% Auto-generated from data/segments.json + layout.json (+ transforms) — do not edit by hand.
% volume: 07Di02
% mode: reading
% segments: 1428
% transform_rules: 0
% toc_marks: 316

% Volume layout defaults (layout.json ``layout``).
\csromanlayoutapply{1}{1.5}{21.6pt}{5pt}{6.3pt}{65pt}{2.5em}%

% Reading physical-page layout (layout.json ``page_layout_reading_mode``).
\csromanreadingpagelayoutvolume{1}{1.5}{21.6pt}{5pt}{6.3pt}{65pt}{2.5em}%
\csromanreadingpagelayoutenable

\csromanheader{2}{\textbf{ทีฆนิกาย}}
\csromanmatikaanchor{mtk.0}
\csromantocmarknopage{chapter}{มหาวคฺคปาฬิ}{mtk.0}
\bookmark[dest={mtk.0},level=0]{มหาวคฺคปาฬิ}
\gambhira{\textbf{มหาวคฺคปาฬิ}}
\namakkaram{นโม ตสฺส ภควโต อรหโต สมฺมาสมฺพุทฺธสฺส.}
\csromanmatikaanchor{mtk.1}
\csromanmatikaanchor{mtk.2}
\csromantocmarknopage{chapter}{1. มหาปทานสุตฺต}{mtk.1}
\bookmark[dest={mtk.1},level=0]{1. มหาปทานสุตฺต}
\csromantocmarkref{section}{ปุพฺเพนิวาสปฏิสํยุตฺตกถา}{mtk.2}
\bookmark[dest={mtk.2},level=1]{ปุพฺเพนิวาสปฏิสํยุตฺตกถา}
\csromanheader{1}{\textbf{ปุพฺเพนิวาสปฏิสํยุตฺตกถา}}
\proseitem{1}{\csromanfolio{1}เอวํ เม สุตํ\csromandash{}เอกํ สมยํ ภควา สาวตฺถิยํ วิหรติ เชตวเน อนาถปิณฺฑิกสฺส อาราเม กเรริกุฏิกายํ.\csromanspacer{} อถ โข สมฺพหุลานํ ภิกฺขูนํ ปจฺฉาภตฺตํ ปิณฺฑปาตปฏิกฺกนฺตานํ กเรริมณฺฑลมาเฬ สนฺนิสินฺนานํ สนฺนิปติตานํ ปุพฺเพนิวาสปฏิสํยุตฺตา ธมฺมี กถา อุทปาทิ “อิติปิ ปุพฺเพนิวาโส อิติปิ ปุพฺเพนิวาโส”ติ.}

\proseitem{2}{อสฺโสสิ โข ภควา ทิพฺพาย โสตธาตุยา วิสุทฺธาย อติกฺกนฺตมานุสิกาย เตสํ ภิกฺขูนํ อิมํ กถาสลฺลาปํ.\csromanspacer{} อถ โข ภควา อุฏฺฐายาสนา เยน กเรริมณฺฑลมาโฬ เตนุปสงฺกมิ, อุปสงฺกมิตฺวา ปญฺญตฺเต อาสเน นิสีทิ, นิสชฺช โข ภควา ภิกฺขู อามนฺเตสิ “กายนุตฺถ ภิกฺขเว เอตรหิ กถาย สนฺนิสินฺนา, กา จ ปน โว อนฺตรากถา วิปฺปกตา”ติ?}

\prose{เอวํ วุตฺเต เต ภิกฺขู ภควนฺตํ เอตทโวจุํ “อิธ ภนฺเต อมฺหากํ ปจฺฉาภตฺตํ ปิณฺฑปาตปฏิกฺกนฺตานํ กเรริมณฺฑลมาเฬ สนฺนิสินฺนานํ สนฺนิปติตานํ ปุพฺเพนิวาสปฏิสํยุตฺตา ธมฺมี กถา อุทปาทิ ‘อิติปิ ปุพฺเพนิวาโส อิติปิ ปุพฺเพนิวาโส’ติ.\csromanspacer{} อยํ โข โน ภนฺเต อนฺตรากถา วิปฺปกตา.\csromanspacer{} อถ ภควา อนุปฺปตฺโต”ติ.}

\proseitem{3}{\csromanfolio{2}อิจฺเฉยฺยาถ โน ตุเมฺห ภิกฺขเว ปุพฺเพนิวาสปฏิสํยุตฺตํ ธมฺมิํ กถํ โสตุนฺติ? เอตสฺส ภควา กาโล, เอตสฺส สุคต กาโล, ยํ ภควา ปุพฺเพนิวาสปฏิสํยุตฺตํ ธมฺมิํ กถํ กเรยฺย, ภควโต สุตฺวา\footnote{ภควโต วจนํ สุตฺวา (สฺยา)} ภิกฺขู ธาเรสฺสนฺตีติ.\csromanspacer{} เตน หิ ภิกฺขเว สุณาถ สาธุกํ มนสิ กโรถ ภาสิสฺสามีติ. “เอวํ ภนฺเต”ติ โข เต ภิกฺขู ภควโต ปจฺจสฺโสสุํ.\csromanspacer{} ภควา เอตทโวจ\csromandash{}}

\proseitem{4}{อิโต โส ภิกฺขเว เอกนวุติกปฺเป\footnote{เอกนวุติกปฺโป (สี), เอกนวุโต กปฺโป (สฺยา, กํ, อิ)} ยํ วิปสฺสี ภควา อรหํ สมฺมาสมฺพุทฺโธ โลเก อุทปาทิ.\csromanspacer{} อิโต โส ภิกฺขเว เอกติํเส กปฺเป\footnote{เอกติํสกปฺโป (สี), เอกติํโส กปฺโป (สฺยา, กํ, อิ)} ยํ สิขี ภควา อรหํ สมฺมาสมฺพุทฺโธ โลเก อุทปาทิ.\csromanspacer{} ตสฺมิญฺเญว โข ภิกฺขเว เอกติํเส กปฺเป เวสฺสภู ภควา อรหํ สมฺมาสมฺพุทฺโธ โลเก อุทปาทิ.\csromanspacer{} อิมสฺมิญฺเญว\footnote{อิมสฺมิํ (กตฺถจิ)} โข ภิกฺขเว ภทฺทกปฺเป กกุสนฺโธ ภควา อรหํ สมฺมาสมฺพุทฺโธ โลเก อุทปาทิ.\csromanspacer{} อิมสฺมิญฺเญว โข ภิกฺขเว ภทฺทกปฺเป โกณาคมโน ภควา อรหํ สมฺมาสมฺพุทฺโธ โลเก อุทปาทิ.\csromanspacer{} อิมสฺมิญฺเญว โข ภิกฺขเว ภทฺทกปฺเป กสฺสโป ภควา อรหํ สมฺมาสมฺพุทฺโธ โลเก อุทปาทิ.\csromanspacer{} อิมสฺมิญฺเญว โข ภิกฺขเว ภทฺทกปฺเป อหํ เอตรหิ อรหํ สมฺมาสมฺพุทฺโธ โลเก อุปฺปนฺโน.}

\proseitem{5}{วิปสฺสี ภิกฺขเว ภควา อรหํ สมฺมาสมฺพุทฺโธ ขตฺติโย ชาติยา อโหสิ, ขตฺติยกุเล อุทปาทิ.\csromanspacer{} สิขี ภิกฺขเว ภควา อรหํ สมฺมาสมฺพุทฺโธ ขตฺติโย ชาติยา อโหสิ, ขตฺติยกุเล อุทปาทิ.\csromanspacer{} เวสฺสภู ภิกฺขเว ภควา อรหํ สมฺมาสมฺพุทฺโธ ขตฺติโย ชาติยา อโหสิ, ขตฺติยกุเล อุทปาทิ.\csromanspacer{} กกุสนฺโธ ภิกฺขเว ภควา อรหํ สมฺมาสมฺพุทฺโธ พฺราหฺมโณ ชาติยา อโหสิ, พฺราหฺมณกุเล อุทปาทิ.\csromanspacer{} โกณาคมโน ภิกฺขเว ภควา อรหํ สมฺมาสมฺพุทฺโธ พฺราหฺมโณ ชาติยา อโหสิ, พฺราหฺมณกุเล อุทปาทิ.\csromanspacer{} กสฺสโป ภิกฺขเว ภควา อรหํ สมฺมาสมฺพุทฺโธ พฺราหฺมโณ ชาติยา อโหสิ, พฺราหฺมณกุเล อุทปาทิ.\csromanspacer{} อหํ ภิกฺขเว เอตรหิ อรหํ สมฺมาสมฺพุทฺโธ ขตฺติโย ชาติยา อโหสิํ, ขตฺติยกุเล อุปฺปนฺโน.}

\proseitem{6}{\csromanfolio{3}วิปสฺสี ภิกฺขเว ภควา อรหํ สมฺมาสมฺพุทฺโธ โกณฺฑญฺโญ โคตฺเตน อโหสิ.\csromanspacer{} สิขี ภิกฺขเว ภควา อรหํ สมฺมาสมฺพุทฺโธ โกณฺฑญฺโญ โคตฺเตน อโหสิ.\csromanspacer{} เวสฺสภู ภิกฺขเว ภควา อรหํ สมฺมาสมฺพุทฺโธ โกณฺฑญฺโญ โคตฺเตน อโหสิ.\csromanspacer{} กกุสนฺโธ ภิกฺขเว ภควา อรหํ สมฺมาสมฺพุทฺโธ กสฺสโป โคตฺเตน อโหสิ.\csromanspacer{} โกณาคมโน ภิกฺขเว ภควา อรหํ สมฺมาสมฺพุทฺโธ กสฺสโป โคตฺเตน อโหสิ.\csromanspacer{} กสฺสโป ภิกฺขเว ภควา อรหํ สมฺมาสมฺพุทฺโธ กสฺสโป โคตฺเตน อโหสิ.\csromanspacer{} อหํ ภิกฺขเว เอตรหิ อรหํ สมฺมาสมฺพุทฺโธ โคตโม โคตฺเตน อโหสิํ.}

\proseitem{7}{วิปสฺสิสฺส ภิกฺขเว ภควโต อรหโต สมฺมาสมฺพุทฺธสฺส อสีติวสฺสสหสฺสานิ อายุปฺปมาณํ อโหสิ.\csromanspacer{} สิขิสฺส ภิกฺขเว ภควโต อรหโต สมฺมาสมฺพุทฺธสฺส สตฺตติวสฺสสหสฺสานิ อายุปฺปมาณํ อโหสิ.\csromanspacer{} เวสฺสภุสฺส ภิกฺขเว ภควโต อรหโต สมฺมาสมฺพุทฺธสฺส สฏฺฐิวสฺสสหสฺสานิ อายุปฺปมาณํ อโหสิ.\csromanspacer{} กกุสนฺธสฺส ภิกฺขเว ภควโต อรหโต สมฺมาสมฺพุทฺธสฺส จตฺตาลีสวสฺสสหสฺสานิ อายุปฺปมาณํ อโหสิ.\csromanspacer{} โกณาคมนสฺส ภิกฺขเว ภควโต อรหโต สมฺมาสมฺพุทฺธสฺส ติํสวสฺสสหสฺสานิ อายุปฺปมาณํ อโหสิ.\csromanspacer{} กสฺสปสฺส ภิกฺขเว ภควโต อรหโต สมฺมาสมฺพุทฺธสฺส วีสติวสฺสสหสฺสานิ อายุปฺปมาณํ อโหสิ.\csromanspacer{} มยฺหํ ภิกฺขเว เอตรหิ อปฺปกํ อายุปฺปมาณํ ปริตฺตํ ลหุกํ, โย จิรํ ชีวติ, โส วสฺสสตํ อปฺปํ วา ภิยฺโย.}

\proseitem{8}{วิปสฺสี ภิกฺขเว ภควา อรหํ สมฺมาสมฺพุทฺโธ ปาฏลิยา มูเล อภิสมฺพุทฺโธ.\csromanspacer{} สิขี ภิกฺขเว ภควา อรหํ สมฺมาสมฺพุทฺโธ ปุณฺฑรีกสฺส มูเล อภิสมฺพุทฺโธ.\csromanspacer{} เวสฺสภู ภิกฺขเว ภควา อรหํ สมฺมาสมฺพุทฺโธ สาลสฺส มูเล อภิสมฺพุทฺโธ.\csromanspacer{} กกุสนฺโธ ภิกฺขเว ภควา อรหํ สมฺมาสมฺพุทฺโธ สิรีสสฺส มูเล อภิสมฺพุทฺโธ.\csromanspacer{} โกณาคมโน ภิกฺขเว ภควา อรหํ สมฺมาสมฺพุทฺโธ อุทุมฺพรสฺส มูเล อภิสมฺพุทฺโธ.\csromanspacer{} กสฺสโป ภิกฺขเว ภควา อรหํ สมฺมาสมฺพุทฺโธ นิโคฺรธสฺส มูเล อภิสมฺพุทฺโธ.\csromanspacer{} อหํ ภิกฺขเว เอตรหิ อรหํ สมฺมาสมฺพุทฺโธ อสฺสตฺถสฺส มูเล อภิสมฺพุทฺโธ.}

\proseitem{9}{\csromanfolio{4}วิปสฺสิสฺส ภิกฺขเว ภควโต อรหโต สมฺมาสมฺพุทฺธสฺส ขณฺฑติสฺสํ นาม สาวกยุคํ อโหสิ อคฺคํ ภทฺทยุคํ.\csromanspacer{} สิขิสฺส ภิกฺขเว ภควโต อรหโต สมฺมาสมฺพุทฺธสฺส อภิภูสมฺภวํ นาม สาวกยุคํ อโหสิ อคฺคํ ภทฺทยุคํ.\csromanspacer{} เวสฺสภุสฺส ภิกฺขเว ภควโต อรหโต สมฺมาสมฺพุทฺธสฺส โสณุตฺตรํ นาม สาวกยุคํ อโหสิ อคฺคํ ภทฺทยุคํ.\csromanspacer{} กกุสนฺธสฺส ภิกฺขเว ภควโต อรหโต สมฺมาสมฺพุทฺธสฺส วิธุรสญฺชีวํ นาม สาวกยุคํ อโหสิ อคฺคํ ภทฺทยุคํ.\csromanspacer{} โกณาคมนสฺส ภิกฺขเว ภควโต อรหโต สมฺมาสมฺพุทฺธสฺส ภิยฺโยสุตฺตรํ นาม สาวกยุคํ อโหสิ อคฺคํ ภทฺทยุคํ.\csromanspacer{} กสฺสปสฺส ภิกฺขเว ภควโต อรหโต สมฺมาสมฺพุทฺธสฺส ติสฺสภารทฺวาชํ นาม สาวกยุคํ อโหสิ อคฺคํ ภทฺทยุคํ.\csromanspacer{} มยฺหํ ภิกฺขเว เอตรหิ สาริปุตฺตโมคฺคลฺลานํ นาม สาวกยุคํ อโหสิ อคฺคํ ภทฺทยุคํ.}

\proseitem{10}{วิปสฺสิสฺส ภิกฺขเว ภควโต อรหโต สมฺมาสมฺพุทฺธสฺส ตโย สาวกานํ สนฺนิปาตา อเหสุํ.\csromanspacer{} เอโก สาวกานํ สนฺนิปาโต อโหสิ อฏฺฐสฏฺฐิภิกฺขุสตสหสฺสํ, เอโก สาวกานํ สนฺนิปาโต อโหสิ ภิกฺขุสตสหสฺสํ, เอโก สาวกานํ สนฺนิปาโต อโหสิ อสีติภิกฺขุสหสฺสานิ.\csromanspacer{} วิปสฺสิสฺส ภิกฺขเว ภควโต อรหโต สมฺมาสมฺพุทฺธสฺส อิเม ตโย สาวกานํ สนฺนิปาตา อเหสุํ สพฺเพสํเยว}

\prose{สิขิสฺส ภิกฺขเว ภควโต อรหโต สมฺมาสมฺพุทฺธสฺส ตโย สาวกานํ สนฺนิปาตา อเหสุํ.\csromanspacer{} เอโก สาวกานํ สนฺนิปาโต อโหสิ ภิกฺขุสตสหสฺสํ, เอโก สาวกานํ สนฺนิปาโต อโหสิ อสีติภิกฺขุสหสฺสานิ, เอโก สาวกานํ สนฺนิปาโต อโหสิ สตฺตติภิกฺขุสหสฺสานิ.\csromanspacer{} สิขิสฺส ภิกฺขเว ภควโต อรหโต สมฺมาสมฺพุทฺธสฺส อิเม ตโย สาวกานํ สนฺนิปาตา อเหสุํ สพฺเพสํเยว}

\prose{เวสฺสภุสฺส ภิกฺขเว ภควโต อรหโต สมฺมาสมฺพุทฺธสฺส ตโย สาวกานํ สนฺนิปาตา อเหสุํ.\csromanspacer{} เอโก สาวกานํ สนฺนิปาโต อโหสิ อสีติภิกฺขุสหสฺสานิ, เอโก สาวกานํ สนฺนิปาโต อโหสิ สตฺตติภิกฺขุสหสฺสานิ, เอโก สาวกานํ สนฺนิปาโต \csromanfolio{5}อโหสิ สฏฺฐิภิกฺขุสหสฺสานิ.\csromanspacer{} เวสฺสภุสฺส ภิกฺขเว ภควโต อรหโต สมฺมาสมฺพุทฺธสฺส อิเม ตโย สาวกานํ สนฺนิปาตา อเหสุํ สพฺเพสํเยว}

\prose{กกุสนฺธสฺส ภิกฺขเว ภควโต อรหโต สมฺมาสมฺพุทฺธสฺส เอโก สาวกานํ สนฺนิปาโต อโหสิ จตฺตาลีสภิกฺขุสหสฺสานิ.\csromanspacer{} กกุสนฺธสฺส ภิกฺขเว ภควโต อรหโต สมฺมาสมฺพุทฺธสฺส อยํ เอโก สาวกานํ สนฺนิปาโต อโหสิ สพฺเพสํเยว ขีณาสวานํ.}

\prose{โกณาคมนสฺส ภิกฺขเว ภควโต อรหโต สมฺมาสมฺพุทฺธสฺส เอโก สาวกานํ สนฺนิปาโต อโหสิ ติํสภิกฺขุสหสฺสานิ.\csromanspacer{} โกณาคมนสฺส ภิกฺขเว ภควโต อรหโต สมฺมาสมฺพุทฺธสฺส อยํ เอโก สาวกานํ สนฺนิปาโต อโหสิ สพฺเพสํเยว ขีณาสวานํ.}

\prose{กสฺสปสฺส ภิกฺขเว ภควโต อรหโต สมฺมาสมฺพุทฺธสฺส เอโก สาวกานํ สนฺนิปาโต อโหสิ วีสติภิกฺขุสหสฺสานิ.\csromanspacer{} กสฺสปสฺส ภิกฺขเว ภควโต อรหโต สมฺมาสมฺพุทฺธสฺส อยํ เอโก สาวกานํ สนฺนิปาโต อโหสิ สพฺเพสํเยว ขีณาสวานํ.}

\prose{มยฺหํ ภิกฺขเว เอตรหิ เอโก สาวกานํ สนฺนิปาโต อโหสิ อฑฺฒเตฬสานิ ภิกฺขุสตานิ.\csromanspacer{} มยฺหํ ภิกฺขเว อยํ เอโก สาวกานํ สนฺนิปาโต อโหสิ สพฺเพสํเยว ขีณาสวานํ.}

\proseitem{11}{วิปสฺสิสฺส ภิกฺขเว ภควโต อรหโต สมฺมาสมฺพุทฺธสฺส อโสโก นาม ภิกฺขุ อุปฏฺฐาโก อโหสิ อคฺคุปฏฺฐาโก.\csromanspacer{} สิขิสฺส ภิกฺขเว ภควโต อรหโต สมฺมาสมฺพุทฺธสฺส เขมงฺกโร นาม ภิกฺขุ อุปฏฺฐาโก อโหสิ อคฺคุปฏฺฐาโก.\csromanspacer{} เวสฺสภุสฺส ภิกฺขเว ภควโต อรหโต สมฺมาสมฺพุทฺธสฺส อุปสนฺโต นาม ภิกฺขุ อุปฏฺฐาโก อโหสิ อคฺคุปฏฺฐาโก.\csromanspacer{} กกุสนฺธสฺส ภิกฺขเว ภควโต อรหโต สมฺมาสมฺพุทฺธสฺส พุทฺธิโช นาม ภิกฺขุ อุปฏฺฐาโก อโหสิ อคฺคุปฏฺฐาโก.\csromanspacer{} โกณาคมนสฺส ภิกฺขเว ภควโต อรหโต สมฺมาสมฺพุทฺธสฺส โสตฺถิโช นาม ภิกฺขุ อุปฏฺฐาโก อโหสิ อคฺคุปฏฺฐาโก.\csromanspacer{} กสฺสปสฺส ภิกฺขเว ภควโต อรหโต สมฺมาสมฺพุทฺธสฺส สพฺพมิตฺโต \csromanfolio{6}นาม ภิกฺขุ อุปฏฺฐาโก อโหสิ อคฺคุปฏฺฐาโก.\csromanspacer{} มยฺหํ ภิกฺขเว เอตรหิ อานนฺโท นาม ภิกฺขุ อุปฏฺฐาโก อโหสิ อคฺคุปฏฺฐาโก.}

\proseitem{12}{วิปสฺสิสฺส ภิกฺขเว ภควโต อรหโต สมฺมาสมฺพุทฺธสฺส พนฺธุมา นาม ราชา ปิตา อโหสิ.\csromanspacer{} พนฺธุมตี นาม เทวี มาตา อโหสิ ชเนตฺติ\footnote{ชเนตฺตี (สฺยา)}.\csromanspacer{} พนฺธุมสฺส รญฺโญ พนฺธุมตี นาม นครํ ราชธานี อโหสิ.}

\prose{สิขิสฺส ภิกฺขเว ภควโต อรหโต สมฺมาสมฺพุทฺธสฺส อรุโณ นาม ราชา ปิตา อโหสิ.\csromanspacer{} ปภาวตี นาม เทวี มาตา อโหสิ ชเนตฺติ.\csromanspacer{} อรุณสฺส รญฺโญ อรุณวตี นาม นครํ ราชธานี อโหสิ.}

\prose{เวสฺสภุสฺส ภิกฺขเว ภควโต อรหโต สมฺมาสมฺพุทฺธสฺส สุปฺปติโต นาม\footnote{สุปฺปตีโต นาม (สฺยา) 3.ยสวตี นาม (สฺยา, อิ)} ราชา ปิตา อโหสิ.\csromanspacer{} วสฺสวตี นาม3 เทวี มาตา อโหสิ ชเนตฺติ.\csromanspacer{} สุปฺปติตสฺส รญฺโญ อโนมํ นาม นครํ ราชธานี อโหสิ.}

\prose{กกุสนฺธสฺส ภิกฺขเว ภควโต อรหโต สมฺมาสมฺพุทฺธสฺส อคฺคิทตฺโต นาม พฺราหฺมโณ ปิตา อโหสิ.\csromanspacer{} วิสาขา นาม พฺราหฺมณี มาตา อโหสิ ชเนตฺติ.\csromanspacer{} เตน โข ปน ภิกฺขเว สมเยน เขโม นาม ราชา อโหสิ.\csromanspacer{} เขมสฺส รญฺโญ เขมวตี นาม นครํ ราชธานี อโหสิ.}

\prose{โกณาคมนสฺส ภิกฺขเว ภควโต อรหโต สมฺมาสมฺพุทฺธสฺส ยญฺญทตฺโต นาม พฺราหฺมโณ ปิตา อโหสิ.\csromanspacer{} อุตฺตรา นาม พฺราหฺมณี มาตา อโหสิ ชเนตฺติ.\csromanspacer{} เตน โข ปน ภิกฺขเว สมเยน โสโภ นาม ราชา อโหสิ.\csromanspacer{} โสภสฺส รญฺโญ โสภวตี นาม นครํ ราชธานี อโหสิ.}

\prose{กสฺสปสฺส ภิกฺขเว ภควโต อรหโต สมฺมาสมฺพุทฺธสฺส พฺรหฺมทตฺโต นาม พฺราหฺมโณ ปิตา อโหสิ.\csromanspacer{} ธนวตี นาม พฺราหฺมนี มาตา อโหสิ ชเนตฺติ.\csromanspacer{} เตน โข ปน ภิกฺขเว สมเยน \csromanfolio{7}กิกี นาม\footnote{กิํกี นาม (สฺยา)} ราชา อโหสิ.\csromanspacer{} กิกิสฺส รญฺโญ พาราณสี นาม นครํ ราชธานี อโหสิ.}

\prose{มยฺหํ ภิกฺขเว เอตรหิ สุทฺโธทโน นาม ราชา ปิตา อโหสิ.\csromanspacer{} มายา นาม เทวี มาตา อโหสิ ชเนตฺติ.\csromanspacer{} กปิลวตฺถุ นาม นครํ ราชธานี อโหสีติ.\csromanspacer{} อิทมโวจ ภควา, อิทํ วตฺวาน สุคโต อุฏฺฐายาสนา วิหารํ ปาวิสิ.}

\proseitem{13}{อถ โข เตสํ ภิกฺขูนํ อจิรปกฺกนฺตสฺส ภควโต อยมนฺตรากถา อุทปาทิ “อจฺฉริยํ อาวุโส อพฺภุตํ อาวุโส ตถาคตสฺส มหิทฺธิกตา มหานุภาวตา.\csromanspacer{} ยตฺร หิ นาม ตถาคโต อตีเต พุทฺเธ ปรินิพฺพุเต ฉินฺนปปญฺเจ ฉินฺนวฏุเม ปริยาทินฺนวฏฺเฏ สพฺพทุกฺขวีติวตฺเต ชาติโตปิ อนุสฺสริสฺสติ, นามโตปิ อนุสฺสริสฺสติ, โคตฺตโตปิ อนุสฺสริสฺสติ, อายุปฺปมาณโตปิ อนุสฺสริสฺสติ, สาวกยุคโตปิ อนุสฺสริสฺสติ, สาวกสนฺนิปาตโตปิ อนุสฺสริสฺสติ ‘เอวํชจฺจา เต ภควนฺโต อเหสุํ อิติปิ, เอวํนามา เอวํโคตฺตา เอวํสีลา เอวํธมฺมา เอวํปญฺญา เอวํวิหารี เอวํวิมุตฺตา เต ภควนฺโต อเหสุํ อิติปี’ติ.}

\prose{กิํ นุ โข อาวุโส ตถาคตสฺเสว นุ โข เอสา ธมฺมธาตุ สุปฺปฏิวิทฺธา, ยสฺสา ธมฺมธาตุยา สุปฺปฏิวิทฺธตฺตา ตถาคโต อตีเต พุทฺเธ ปรินิพฺพุเต ฉินฺนปปญฺเจ ฉินฺนวฏุเม ปริยาทินฺนวฏฺเฏ สพฺพทุกฺขวีติวตฺเต ชาติโตปิ อนุสฺสรติ, นามโตปิ อนุสฺสรติ, โคตฺตโตปิ อนุสฺสรติ.\csromanspacer{} อายุปฺปมาณโตปิ อนุสฺสรติ, สาวกยุคโตปิ อนุสฺสรติ, สาวกสนฺนิปาตโตปิ อนุสฺสรติ ‘เอวํชจฺจา เต ภควนฺโต อเหสุํ อิติปิ, เอวํนามา เอวํโคตฺตา เอวํสีลา เอวํธมฺมา เอวํปญฺญา เอวํวิหารี เอวํวิมุตฺตา เต ภควนฺโต อเหสุํ อิติปี’ติ, อุทาหุ เทวตา ตถาคตสฺส เอตมตฺถํ อาโรเจสุํ, เยน ตถาคโต อตีเต พุทฺเธ ปรินิพฺพุเต ฉินฺนปปญฺเจ ฉินฺนวฏุเม ปริยาทินฺนวฏฺเฏ สพฺพทุกฺขวีติวตฺเต ชาติโตปิ อนุสฺสรติ, นามโตปิ อนุสฺสรติ, โคตฺตโตปิ อนุสฺสรติ, อายุปฺปมาณโตปิ อนุสฺสรติ, สาวกยุคโตปิ อนุสฺสรติ, สาวกสนฺนิปาตโตปิ อนุสฺสรติ ‘เอวํชจฺจา เต ภควนฺโต อเหสุํ อิติปิ, เอวํนามา เอวํโคตฺตา เอวํสีลา เอวํธมฺมา เอวํปญฺญา เอวํวิหารี \csromanfolio{8}เอวํวิมุตฺตา เต ภควนฺโต อเหสุํ อิติปี”ติ.\csromanspacer{} อยญฺจ หิทํ เตสํ ภิกฺขูนํ อนฺตรากถา วิปฺปกตา โหติ.}

\proseitem{14}{อถ โข ภควา สายนฺหสมยํ ปฏิสลฺลานา วุฏฺฐิโต เยน กเรริมณฺฑลมาโฬ เตนุปสงฺกมิ, อุปสงฺกมิตฺวา ปญฺญตฺเต อาสเน นิสีทิ, นิสชฺช โข ภควา ภิกฺขู อามนฺเตสิ “กายนุตฺถ ภิกฺขเว เอตรหิ กถาย สนฺนิสินฺนา.\csromanspacer{} กา จ ปน โว อนฺตรากถา วิปฺปกตา”ติ.}

\prose{เอวํ วุตฺเต เต ภิกฺขู ภควนฺตํ เอตทโวจุํ “อิธ ภนฺเต อมฺหากํ อจิรปกฺกนฺตสฺส ภควโต อยํ อนฺตรากถา อุทปาทิ ‘อจฺฉริยํ อาวุโส อพฺภุตํ อาวุโส ตถาคตสฺส มหิทฺธิกตา มหานุภาวตา, ยตฺร หิ นาม ตถาคโต อตีเต พุทฺเธ ปรินิพฺพุเต ฉินฺนปปญฺเจ ฉินฺนวฏุเม ปริยาทินฺนวฏฺเฏ สพฺพทุกฺขวีติวตฺเต ชาติโตปิ อนุสฺสริสฺสติ, นามโตปิ อนุสฺสริสฺสติ, โคตฺตโตปิ อนุสฺสริสฺสติ, อายุปฺปมาณโตปิ อนุสฺสริสฺสติ, สาวกยุคโตปิ อนุสฺสริสฺสติ, สาวกสนฺนิปาตโตปิ อนุสฺสริสฺสติ ‘เอวํชจฺจา เต ภควนฺโต อเหสุํ อิติปิ, เอวํนามา เอวํโคตฺตา เอวํสีลา เอวํธมฺมา เอวํปญฺญา เอวํวิหารี เอวํวิมุตฺตา เต ภควนฺโต อเหสุํ อิติปี’ติ.\csromanspacer{} กิํ นุ โข อาวุโส ตถาคตสฺเสว นุ โข เอสา ธมฺมธาตุ สุปฺปฏิวิทฺธา, ยสฺสา ธมฺมธาตุยา สุปฺปฏิวิทฺธตฺตา ตถาคโต อตีเต พุทฺเธ ปรินิพฺพุเต ฉินฺนปปญฺเจ ฉินฺนวฏุเม ปริยาทินฺนวฏฺเฏ สพฺพทุกฺขวีติวตฺเต ชาติโตปิ อนุสฺสรติ, นามโตปิ อนุสฺสรติ, โคตฺตโตปิ อนุสฺสรติ, อายุปฺปมาณโตปิ อนุสฺสรติ, สาวกยุคโตปิ อนุสฺสรติ, สาวกสนฺนิปาตโตปิ อนุสฺสรติ ‘เอวํชจฺจา เต ภควนฺโต อเหสุํ อิติปิ, เอวํนามา เอวํโคตฺตา เอวํสีลา เอวํธมฺมา เอวํปญฺญา เอวํวิหารี เอวํวิมุตฺตา เต ภควนฺโต อเหสุํ อิติปี’ติ, อุทาหุ เทวตา ตถาคตสฺส เอตมตฺถํ อาโรเจสุํ, เยน ตถาคโต อตีเต พุทฺเธ ปรินิพฺพุเต ฉินฺนปปญฺเจ ฉินฺนวฏุเม ปริยาทินฺนวฏฺเฏ สพฺพทุกฺขวีติวตฺเต ชาติโตปิ อนุสฺสรติ, นามโตปิ อนุสฺสรติ, โคตฺตโตปิ อนุสฺสรติ, อายุปฺปมาณโตปิ อนุสฺสรติ, สาวกยุคโตปิ อนุสฺสรติ, สาวกสนฺนิปาตโตปิ อนุสฺสรติ ‘เอวํชจฺจา เต ภควนฺโต อเหสุํ อิติปิ, เอวํนามา เอวํโคตฺตา เอวํสีลา เอวํธมฺมา เอวํปญฺญา เอวํวิหารี \csromanfolio{9}เอวํวิมุตฺตา เต ภควนฺโต อเหสุํ อิติปี’ติ.\csromanspacer{} อยํ โข โน ภนฺเต อนฺตรากถา วิปฺปกตา, อถ ภควา อนุปฺปตฺโต”ติ.}

\proseitem{15}{ตถาคตสฺเสเวสา ภิกฺขเว ธมฺมธาตุ สุปฺปฏิวิทฺธา, ยสฺสา ธมฺมธาตุยา สุปฺปฏิวิทฺธตฺตา ตถาคโต อตีเต พุทฺเธ ปรินิพฺพุเต ฉินฺนปปญฺเจ ฉินฺนวฏุเม ปริยาทินฺนวฏฺเฏ สพฺพทุกฺขวีติวตฺเต ชาติโตปิ อนุสฺสรติ, นามโตปิ อนุสฺสรติ, โคตฺตโตปิ อนุสฺสรติ, อายุปฺปมาณโตปิ อนุสฺสรติ, สาวกยุคโตปิ อนุสฺสรติ, สาวกสนฺนิปาตโตปิ อนุสฺสรติ “เอวํชจฺจา เต ภควนฺโต อเหสุํ อิติปิ, เอวํนามา เอวํโคตฺตา เอวํสีลา เอวํธมฺมา เอวํปญฺญา เอวํวิหารี เอวํวิมุตฺตา เต ภควนฺโต อเหสุํ อิติปี”ติ.\csromanspacer{} เทวตาปิ ตถาคตสฺส เอตมตฺถํ อาโรเจสุํ, เยน ตถาคโต อตีเต พุทฺเธ ปรินิพฺพุเต ฉินฺนปปญฺเจ ฉินฺนวฏุเม ปริยาทินฺนวฏฺเฏ สพฺพทุกฺขวีติวตฺเต ชาติโตปิ อนุสฺสรติ, นามโตปิ อนุสฺสรติ, โคตฺตโตปิ อนุสฺสรติ, อายุปฺปมาณโตปิ อนุสฺสรติ, สาวกยุคโตปิ อนุสฺสรติ, สาวกสนฺนิปาตโตปิ อนุสฺสรติ “เอวํชจฺจา เต ภควนฺโต อเหสุํ อิติปิ, เอวํนามา เอวํโคตฺตา เอวํสีลา เอวํธมฺมา เอวํปญฺญา เอวํวิหารี เอวํวิมุตฺตา เต ภควนฺโต อเหสุํ อิติปี”ติ.}

\prose{อิจฺเฉยฺยาถ โน ตุเมฺห ภิกฺขเว ภิยฺโยโส มตฺตาย ปุพฺเพนิวาสปฏิสํยุตฺตํ ธมฺมิํ กถํ โสตุนฺติ? เอตสฺส ภควา กาโล, เอตสฺส สุคต กาโล, ยํ ภควา ภิยฺโยโส มตฺตาย ปุพฺเพนิวาสปฏิสํยุตฺตํ ธมฺมิํ กถํ กเรยฺย, ภควโต สุตฺวา ภิกฺขู ธาเรสฺสนฺตีติ.\csromanspacer{} เตน หิ ภิกฺขเว สุณาถ สาธุกํ มนสิ กโรถ ภาสิสฺสามีติ. “เอวํ ภนฺเต”ติ โข เต ภิกฺขู ภควโต ปจฺจสฺโสสุํ.\csromanspacer{} ภควา เอตทโวจ\csromandash{}}

\proseitem{16}{อิโต โส ภิกฺขเว เอกนวุติกปฺเป ยํ วิปสฺสี ภควา อรหํ สมฺมาสมฺพุทฺโธ โลเก อุทปาทิ.\csromanspacer{} วิปสฺสี ภิกฺขเว ภควา อรหํ สมฺมาสมฺพุทฺโธ ขตฺติโย ชาติยา อโหสิ, ขตฺติยกุเล อุทปาทิ.\csromanspacer{} วิปสฺสี ภิกฺขเว ภควา อรหํ สมฺมาสมฺพุทฺโธ โกณฺฑญฺโญ โคตฺเตน อโหสิ.\csromanspacer{} วิปสฺสิสฺส ภิกฺขเว ภควโต อรหโต สมฺมาสมฺพุทฺธสฺส อสีติวสฺสสหสฺสานิ อายุปฺปมาณํ อโหสิ.\csromanspacer{} วิปสฺสี ภิกฺขเว ภควา \csromanfolio{10}อรหํ สมฺมาสมฺพุทฺโธ ปาฏลิยา มูเล อภิสมฺพุทฺโธ.\csromanspacer{} วิปสฺสิสฺส ภิกฺขเว ภควโต อรหโต สมฺมาสมฺพุทฺธสฺส ขณฺฑติสฺสํ นาม สาวกยุคํ อโหสิ อคฺคํ ภทฺทยุคํ.\csromanspacer{} วิปสฺสิสฺส ภิกฺขเว ภควโต อรหโต สมฺมาสมฺพุทฺธสฺส ตโย สาวกานํ สนฺนิปาตา อเหสุํ.\csromanspacer{} เอโก สาวกานํ สนฺนิปาโต อโหสิ อฏฺฐสฏฺฐิภิกฺขุสตสหสฺสํ, เอโก สาวกานํ สนฺนิปาโต อโหสิ ภิกฺขุสตสหสฺสํ, เอโก สาวกานํ สนฺนิปาโต อโหสิ อสีติภิกฺขุสหสฺสานิ.\csromanspacer{} วิปสฺสิสฺส ภิกฺขเว ภควโต อรหโต สมฺมาสมฺพุทฺธสฺส อิเม ตโย สาวกานํ สนฺนิปาตา อเหสุํ สพฺเพสํเยว ขีณาสวานํ.\csromanspacer{} วิปสฺสิสฺส ภิกฺขเว ภควโต อรหโต สมฺมาสมฺพุทฺธสฺส อโสโก นาม ภิกฺขุ อุปฏฺฐาโก อโหสิ อคฺคุปฏฺฐาโก.\csromanspacer{} วิปสฺสิสฺส ภิกฺขเว ภควโต อรหโต สมฺมาสมฺพุทฺธสฺส พนฺธุมา นาม ราชา ปิตา อโหสิ.\csromanspacer{} พนฺธุมตี นาม เทวี มาตา อโหสิ ชเนตฺติ.\csromanspacer{} พนฺธุมสฺส รญฺโญ พนฺธุมตี นาม นครํ ราชธานี อโหสิ.}

\csromanmatikaanchor{mtk.3}
\csromantocmarkref{section}{โพธิสตฺตธมฺมตา}{mtk.3}
\bookmark[dest={mtk.3},level=1]{โพธิสตฺตธมฺมตา}
\csromanheader{1}{\textbf{โพธิสตฺตธมฺมตา}}
\proseitem{17}{อถ โข ภิกฺขเว วิปสฺสี โพธิสตฺโต ตุสิตา กายา จวิตฺวา สโต สมฺปชาโน มาตุกุจฺฉิํ โอกฺกมิ.\csromanspacer{} อยเมตฺถ ธมฺมตา. (1)}

\proseitem{18}{ธมฺมตา เอสา ภิกฺขเว, ยทา โพธิสตฺโต ตุสิตา กายา จวิตฺวา มาตุกุจฺฉิํ โอกฺกมติ.\csromanspacer{} อถ สเทวเก โลเก สมารเก สพฺรหฺมเก สสฺสมณพฺราหฺมณิยา ปชาย สเทวมนุสฺสาย อปฺปมาโณ อุฬาโร โอภาโส ปาตุ ภวติ อติกฺกมฺเมว เทวานํ เทวานุภาวํ.\csromanspacer{} ยาปิ ตา โลกนฺตริกา อฆา อสํวุตา อนฺธการา อนฺธการติมิสา, ยตฺถปิเม จนฺทิมสูริยา เอวํมหิทฺธิกา เอวํมหานุภาวา อาภาย นานุโภนฺติ, ตตฺถปิ อปฺปมาโณ อุฬาโร โอภาโส ปาตุ ภวติ อติกฺกมฺเมว เทวานํ เทวานุภาวํ.\csromanspacer{} เยปิ ตตฺถ สตฺตา อุปปนฺนา, เตปิ เตโนภาเสน อญฺญมญฺญํ สญฺชานนฺติ “อญฺเญปิ กิร โภ สนฺติ สตฺตา อิธูปปนฺนา”ติ.\csromanspacer{} อยญฺจ ทสสหสฺสี โลกธาตุ สงฺกมฺปติ สมฺปกมฺปติ สมฺปเวธติ.\csromanspacer{} อปฺปมาโณ จ อุฬาโร โอภาโส โลเก ปาตุ ภวติ อติกฺกมฺเมว เทวานํ เทวานุภาวํ.\csromanspacer{} อยเมตฺถ ธมฺมตา. (2)}

\proseitem{19}{\csromanfolio{11}ธมฺมตา เอสา ภิกฺขเว, ยทา โพธิสตฺโต มาตุกุจฺฉิํ โอกฺกนฺโต โหติ.\csromanspacer{} จตฺตาโร นํ เทวปุตฺตา จตุทฺทิสํ\footnote{จาตุทฺทิสํ (สฺยา)} รกฺขาย อุปคจฺฉนฺติ “มา นํ โพธิสตฺตํ วา โพธิสตฺตมาตรํ วา มนุสฺโส วา อมนุสฺโส วา โกจิ วา วิเหเฐสี”ติ.\csromanspacer{} อยเมตฺถ ธมฺมตา. (3)}

\proseitem{20}{ธมฺมตา เอสา ภิกฺขเว, ยทา โพธิสตฺโต มาตุกุจฺฉิํ โอกฺกนฺโต โหติ.\csromanspacer{} ปกติยา สีลวตี โพธิสตฺตมาตา โหติ, วิรตา ปาณาติปาตา, วิรตา อทินฺนาทานา, วิรตา กาเมสุมิจฺฉาจารา, วิรตา มุสาวาทา, วิรตา สุราเมรยมชฺชปฺปมาทฏฺฐานา.\csromanspacer{} อยเมตฺถ ธมฺมตา. (4)}

\proseitem{21}{ธมฺมตา เอสา ภิกฺขเว, ยทา โพธิสตฺโต มาตุกุจฺฉิํ โอกฺกนฺโต โหติ.\csromanspacer{} น โพธิสตฺตมาตุ ปุริเสสุ มานสํ อุปฺปชฺชติ กามคุณูปสํหิตํ, อนติกฺกมนียา จ โพธิสตฺตมาตา โหติ เกนจิ ปุริเสน รตฺตจิตฺเตน.\csromanspacer{} อยเมตฺถ ธมฺมตา. (5)}

\proseitem{22}{ธมฺมตา เอสา ภิกฺขเว, ยทา โพธิสตฺโต มาตุกุจฺฉิํ โอกฺกนฺโต โหติ.\csromanspacer{} ลาภินี โพธิสตฺตมาตา โหติ ปญฺจนฺนํ กามคุณานํ, สา ปญฺจหิ กามคุเณหิ สมปฺปิตา สมงฺคีภูตา ปริจาเรติ.\csromanspacer{} อยเมตฺถ ธมฺมตา. (6)}

\proseitem{23}{ธมฺมตา เอสา ภิกฺขเว, ยทา โพธิสตฺโต มาตุกุจฺฉิํ โอกฺกนฺโต โหติ.\csromanspacer{} น โพธิสตฺตมาตุ โกจิเทว อาพาโธ อุปฺปชฺชติ, สุขินี โพธิสตฺตมาตา โหติ อกิลนฺตกายา, โพธิสตฺตญฺจ โพธิสตฺตมาตา ติโรกุจฺฉิคตํ ปสฺสติ สพฺพงฺคปจฺจงฺคิํ อหีนินฺทฺริยํ.\csromanspacer{} เสยฺยถาปิ ภิกฺขเว มณิ เวฬุริโย สุโภ ชาติมา อฏฺฐํโส สุปริกมฺมกโต อจฺโฉ วิปฺปสนฺโน อนาวิโล สพฺพาการสมฺปนฺโน.\csromanspacer{} ตตฺราสฺส\footnote{ตตฺรสฺส (สฺยา)} สุตฺตํ อาวุตํ นีลํ วา ปีตํ วา โลหิตํ วา โอทาตํ วา ปณฺฑุสุตฺตํ วา, ตเมนํ จกฺขุมา ปุริโส หตฺเถ กริตฺวา ปจฺจเวกฺเขยฺย “อยํ โข มณิ เวฬุริโย สุโภ ชาติมา อฏฺฐํโส สุปริกมฺมกโต อจฺโฉ วิปฺปสนฺโน อนาวิโล สพฺพาการสมฺปนฺโน.\csromanspacer{} ตตฺริทํ สุตฺตํ อาวุตํ นีลํ วา ปีตํ วา โลหิตํ วา โอทาตํ วา ปณฺฑุสุตฺตํ วา”ติ.\csromanspacer{} เอวเมว โข ภิกฺขเว ยทา \csromanfolio{12}โพธิสตฺโต มาตุกุจฺฉิํ โอกฺกนฺโต โหติ.\csromanspacer{} น โพธิสตฺตมาตุ โกจิเทว อาพาโธ อุปฺปชฺชติ, สุขินี โพธิสตฺตมาตา โหติ อกิลนฺตกายา, โพธิสตฺตญฺจ โพธิสตฺตมาตา ติโรกุจฺฉิคตํ ปสฺสติ สพฺพงฺคปจฺจงฺคิํ อหีนินฺทฺริยํ.\csromanspacer{} อยเมตฺถ ธมฺมตา. (7)}

\proseitem{24}{ธมฺมตา เอสา ภิกฺขเว, สตฺตาหชาเต โพธิสตฺเต โพธิสตฺตมาตา กาลํ กโรติ ตุสิตํ กายํ อุปปชฺชติ.\csromanspacer{} อยเมตฺถ ธมฺมตา. (8)}

\proseitem{25}{ธมฺมตา เอสา ภิกฺขเว, ยถา อญฺญา อิตฺถิกา นว วา ทส วา มาเส คพฺภํ กุจฺฉินา ปริหริตฺวา วิชายนฺติ, น เหวํ โพธิสตฺตํ โพธิสตฺตมาตา วิชายติ, ทเสว มาสานิ โพธิสตฺตํ โพธิสตฺตมาตา กุจฺฉินา ปริหริตฺวา วิชายติ.\csromanspacer{} อยเมตฺถ ธมฺมตา. (9)}

\proseitem{26}{ธมฺมตา เอสา ภิกฺขเว, ยถา อญฺญา อิตฺถิกา นิสินฺนา วา นิปนฺนา วา วิชายนฺติ, น เหวํ โพธิสตฺตํ โพธิสตฺตมาตา วิชายติ, ฐิตาว โพธิสตฺตํ โพสตฺตมาตา วิชายติ.\csromanspacer{} อยเมตฺถ ธมฺมตา. (10)}

\proseitem{27}{ธมฺมตา เอสา ภิกฺขเว, ยทา โพธิสตฺโต มาตุกุจฺฉิมฺหา นิกฺขมติ.\csromanspacer{} เทวา ปฐมํ ปฏิคฺคณฺหนฺติ, ปจฺฉา มนุสฺสา.\csromanspacer{} อยเมตฺถ ธมฺมตา. (11)}

\proseitem{28}{ธมฺมตา เอสา ภิกฺขเว, ยทา โพธิสตฺโต มาตุกุจฺฉิมฺหา นิกฺขมติ, อปฺปตฺโตว โพธิสตฺโต ปถวิํ โหติ, จตฺตาโร นํ เทวปุตฺตา ปฏิคฺคเหตฺวา มาตุ ปุรโต ฐเปนฺติ “อตฺตมนา เทวิ โหหิ, มเหสกฺโข เต ปุตฺโต อุปฺปนฺโน”ติ.\csromanspacer{} อยเมตฺถ ธมฺมตา. (12)}

\proseitem{29}{ธมฺมตา เอสา ภิกฺขเว ยทา โพธิสตฺโต มาตุกุจฺฉิมฺหา นิกฺขมติ.\csromanspacer{} วิสโทว นิกฺขมติ อมกฺขิโต อุเทน\footnote{อุทฺเทน (สฺยา), อุทเรน (กตฺถจิ)} อมกฺขิโต เสเมฺหน อมกฺขิโต รุหิเรน อมกฺขิโต เกนจิ อสุจินา สุทฺโธ\footnote{วิสุทฺโธ (สฺยา)} วิสโท.\csromanspacer{} เสยฺยถาปิ ภิกฺขเว มณิรตนํ กาสิเก วตฺเถ นิกฺขิตฺตํ เนว มณิรตนํ กาสิกํ วตฺถํ มกฺเขติ, นาปิ กาสิกํ วตฺถํ มณิรตนํ มกฺเขติ.\csromanspacer{} ตํ กิสฺส เหตุ, อุภินฺนํ สุทฺธตฺตา.\csromanspacer{} เอวเมว โข ภิกฺขเว ยทา โพธิสตฺโต มาตุกุจฺฉิมฺหา นิกฺขมติ, วิสโทว นิกฺขมติ อมกฺขิโต อุเทน \csromanfolio{13}อมกฺขิโต เสเมฺหน อมกฺขิโต รุหิเรน อมกฺขิโต เกนจิ อสุจินา สุทฺโธ วิสโท.\csromanspacer{} อยเมตฺถ ธมฺมตา. (13)}

\proseitem{30}{ธมฺมตา เอสา ภิกฺขเว, ยทา โพธิสตฺโต มาตุกุจฺฉิมฺหา นิกฺขมติ.\csromanspacer{} เทฺว อุทกสฺส ธารา อนฺตลิกฺขา ปาตุ ภวนฺติ เอกา สีตสฺส เอกา อุณฺหสฺส, เยน โพธิสตฺตสฺส อุทกกิจฺจํ กโรนฺติ มาตุ จ.\csromanspacer{} อยเมตฺถ ธมฺมตา. (14)}

\proseitem{31}{ธมฺมตา เอสา ภิกฺขเว, สมฺปติชาโต โพธิสตฺโต สเมหิ ปาเทหิ ปติฏฺฐหิตฺวา อุตฺตราภิมุโข\footnote{อุตฺตเรนาภิมุโข (สฺยา) อุตฺตเรนมุโข (ก)} สตฺตปทวีติหาเรน คจฺฉติ เสตมฺหิ ฉตฺเต อนุธาริยมาเน, สพฺพา จ ทิสา อนุวิโลเกติ, อาสภิํ วาจํ ภาสติ “อคฺโคหมสฺมิ โลกสฺส, เชฏฺโฐหมสฺมิ โลกสฺส, เสฏฺโฐหมสฺมิ โลกสฺส, อยมนฺติมา ชาติ, นตฺถิทานิ ปุนพฺภโว”ติ.\csromanspacer{} อยเมตฺถ ธมฺมตา. (15)}

\proseitem{32}{ธมฺมตา เอสา ภิกฺขเว, ยทา โพธิสตฺโต มาตุกุจฺฉิมฺหา นิกฺขมติ.\csromanspacer{} อถ สเทวเก โลเก สมารเก สพฺรหฺมเก สสฺสมณพฺราหฺมณิยา ปชาย สเทวมนุสฺสาย อปฺปมาโณ อุฬาโร โอภาโส ปาตุ ภวติ อติกฺกมฺเมว เทวานํ เทวานุภาวํ.\csromanspacer{} ยาปิ ตา โลกนฺตริกา อฆา อสํวุตา อนฺธการา อนฺธการติมิสา, ยตฺถปิเม จนฺทิมสูริยา เอวํมหิทฺธิกา เอวํมหานุภาวา อาภาย นานุโภนฺติ, ตตฺถปิ อปฺปมาโณ อุฬาโร โอภาโส ปาตุ ภวติ อติกฺกมฺเมว เทวานํ เทวานุภาวํ.\csromanspacer{} เยปิ ตตฺถ สตฺตา อุปปนฺนา, เตปิ เตโนภาเสน อญฺญมญฺญํ สญฺชานนฺติ “อญฺเญปิ กิร โภ สนฺติ สตฺตา อิธูปปนฺนา”ติ.\csromanspacer{} อยญฺจ ทสสหสฺสี โลกธาตุ สงฺกมฺปติ สมฺปกมฺปติ สมฺปเวธติ.\csromanspacer{} อปฺปมาโณ จ อุฬาโร โอภาโส โลเก ปาตุ ภวติ อติกฺกมฺเมว เทวานํ เทวานุภาวํ.\csromanspacer{} อยเมตฺถ ธมฺมตา. (16)}

\csromanmatikaanchor{mtk.4}
\csromantocmarkref{section}{ทฺวตฺติํสมหาปุริสลกฺขณา}{mtk.4}
\bookmark[dest={mtk.4},level=1]{ทฺวตฺติํสมหาปุริสลกฺขณา}
\csromanheader{1}{\textbf{ทฺวตฺติํสมหาปุริสลกฺขณา}}
\proseitem{33}{ชาเต โข ปน ภิกฺขเว วิปสฺสิมฺหิ กุมาเร พนฺธุมโต รญฺโญ ปฏิเวเทสุํ “ปุตฺโต เต เทว\footnote{เทว เต (ก)} ชาโต, ตํ เทโว ปสฺสตู”ติ.\csromanspacer{} อทฺทสา โข ภิกฺขเว พนฺธุมา ราชา วิปสฺสิํ กุมารํ, ทิสฺวา เนมิตฺเต \csromanfolio{14}พฺราหฺมเณ อามนฺตาเปตฺวา เอตทโวจ “ปสฺสนฺตุ โภนฺโต เนมิตฺตา พฺราหฺมณา กุมารนฺ”ติ.\csromanspacer{} อทฺทสํสุ โข ภิกฺขเว เนมิตฺตา พฺราหฺมณา วิปสฺสิํ กุมารํ, ทิสฺวา พนฺธุมนฺตํ ราชานํ เอตทโวจุํ “อตฺตมโน เทว โหหิ, มเหสกฺโข เต ปุตฺโต อุปฺปนฺโน, ลาภา เต มหาราช, สุลทฺธํ เต มหาราช, ยสฺส เต กุเล เอวรูโป ปุตฺโต อุปฺปนฺโน, อยํ หิ เทว กุมาโร ทฺวตฺติํสมหาปุริสลกฺขเณหิ สมนฺนาคโต, เยหิ สมนฺนาคตสฺส มหาปุริสสฺส เทฺวว คติโย ภวนฺติ อนญฺญา.\csromanspacer{} สเจ อคารํ อชฺฌาวสติ, ราชา โหติ จกฺกวตฺตี ธมฺมิโก ธมฺมราชา จาตุรนฺโต วิชิตาวี ชนปทตฺถาวริยปฺปตฺโต สตฺตรตนสมนฺนาคโต.\csromanspacer{} ตสฺสิมานิ สตฺตรตนานิ ภวนฺติ.\csromanspacer{} เสยฺยถิทํ, จกฺกรตนํ หตฺถิรตนํ อสฺสรตนํ มณิรตนํ อิตฺถิรตนํ คหปติรตนํ ปริณายกรตนเมว สตฺตมํ.\csromanspacer{} ปโรสหสฺสํ โข ปนสฺส ปุตฺตา ภวนฺติ สูรา วีรงฺครูปา ปรเสนปฺปมทฺทนา, โส อิมํ ปถวิํ สาครปริยนฺตํ อทณฺเฑน อสตฺเถน ธมฺเมน อภิวิชิย อชฺฌาวสติ.\csromanspacer{} สเจ โข ปน อคารสฺมา อนคาริยํ ปพฺพชติ, อรหํ โหติ สมฺมาสมฺพุทฺโธ โลเก วิวฏจฺฉโท.}

\proseitem{34}{กตเมหิ จายํ เทว กุมาโร ทฺวตฺติํสมหาปุริสลกฺขเณหิ สมนฺนาคโต, เยหิ สมนฺนาคตสฺส มหาปุริสสฺส เทฺวว คติโย ภวนฺติ อนญฺญา.\csromanspacer{} สเจ อคารํ อชฺฌาวสติ, ราชา โหติ จกฺกวตฺตี ธมฺมิโก ธมฺมราชา จาตุรนฺโต วิชิตาวี ชนปทตฺถาวริยปฺปตฺโต สตฺตรตนสมนฺนาคโต.\csromanspacer{} ตสฺสิมานิ สตฺตรตนานิ ภวนฺติ.\csromanspacer{} เสยฺยถิทํ, จกฺกรตนํ หตฺถิรตนํ อสฺสรตนํ มณิรตนํ อิตฺถิรตนํ คหปติรตนํ ปริณายกรตนเมว สตฺตมํ.\csromanspacer{} ปโรสหสฺสํ โข ปนสฺส ปุตฺตา ภวนฺติ สูรา วีรงฺครูปา ปรเสนปฺปมทฺทนา, โส อิมํ ปถวิํ สาครปริยนฺตํ อทณฺเฑน อสตฺเถน ธมฺเมน อภิวิชิย อชฺฌาวสติ.\csromanspacer{} สเจ โข ปน อคารสฺมา อนคาริยํ ปพฺพชติ, อรหํ โหติ สมฺมาสมฺพุทฺโธ โลเก วิวฏจฺฉโท.}

\proseitem{35}{อยญฺหิ เทว กุมาโร สุปฺปติฏฺฐิตปาโท, ยํปายํ เทว กุมาโร สุปฺปติฏฺฐิตปาโท, อิทมฺปิมสฺส มหาปุริสสฺส มหาปุริสลกฺขณํ ภวติ. (1)}

\prose{อิมสฺส เทว\footnote{อิมสฺส หิ เทว (?)} กุมารสฺส เหฏฺฐา ปาทตเลสุ จกฺกานิ ชาตานิ สหสฺสารานิ สเนมิกานิ สนาภิกานิ สพฺพาการปริปูรานิ, ยมฺปิ \csromanfolio{15}อิมสฺส เทว กุมารสฺส เหฏฺฐา ปาทตเลสุ จกฺกานิ ชาตานิ สหสฺสารานิ สเนมิกานิ สนาภิกานิ สพฺพาการปริปูรานิ, อิทมฺปิมสฺส มหาปุริสสฺส มหาปุริสลกฺขณํ ภวติ. (2)}

\csromanheader{2}{อยญฺหิ เทว กุมาโร อายตปณฺหี ฯเปฯ. (3)}
\csromanheader{2}{อยญฺหิ เทว กุมาโร ทีฆงฺคุลี ฯเปฯ. (4)}
\csromanheader{2}{อยญฺหิ เทว กุมาโร มุทุตลุนหตฺถปาโท ฯเปฯ. (5)}
\csromanheader{2}{อยญฺหิ เทว กุมาโร ชาลหตฺถปาโท ฯเปฯ. (6)}
\csromanheader{2}{อยญฺหิ เทว กุมาโร อุสฺสงฺขปาโท ฯเปฯ. (7)}
\csromanheader{2}{อยญฺหิ เทว กุมาโร เอณิชงฺโฆ ฯเปฯ. (8)}
\prose{อยญฺหิ เทว กุมาโร ฐิตโกว อโนนมนฺโต อุโภหิ ปาณิตเลหิ ชณฺณุกานิ ปริมสติ\footnote{ปรามสติ (ก)} ปริมชฺชติ ฯเปฯ. (9)}

\csromanheader{2}{อยญฺหิ เทว กุมาโร โกโสหิตวตฺถคุโยฺห ฯเปฯ. (10)}
\prose{อยญฺหิ เทว กุมาโร สุวณฺณวณฺโณ กญฺจนสนฺนิภตฺตโจ ฯเปฯ. (11)}

\prose{อยญฺหิ เทว กุมาโร สุขุมจฺฉวี, สุขุมตฺตา ฉวิยา รโชชลฺลํ กาเย น อุปลิมฺปติ\footnote{อุปลิปฺปติ (สฺยา) 3. ปิตนฺตรํโส (สฺยา)} ฯเปฯ. (12)}

\prose{อยญฺหิ เทว กุมาโร เอเกกโลโม, เอเกกานิ โลมานิ โลมกูเปสุ ชาตานิ ฯเปฯ. (13)}

\prose{อยญฺหิ เทว กุมาโร อุทฺธคฺคโลโม, อุทฺธคฺคานิ โลมานิ ชาตานิ นีลานิ อญฺชนวณฺณานิ กุณฺฑลาวฏฺฏานิ ทกฺขิณาวฏฺฏกชาตานิ ฯเปฯ. (14)}

\csromanheader{2}{อยญฺหิ เทว กุมาโร พฺรหฺมุชุคตฺโต ฯเปฯ. (15)}
\csromanheader{2}{อยญฺหิ เทว กุมาโร สตฺตุสฺสโท ฯเปฯ. (16)}
\csromanheader{2}{อยญฺหิ เทว กุมาโร สีหปุพฺพทฺธกาโย ฯเปฯ. (17)}
\csromanheader{2}{อยญฺหิ เทว กุมาโร จิตนฺตรํโส3 ฯเปฯ. (18)}
\prose{อยญฺหิ เทว กุมาโร นิโคฺรธปริมณฺฑโล ยาวตกฺวสฺส กาโย, ตาวตกฺวสฺส พฺยาโม, ยาวตกฺวสฺส พฺยาโม, ตาวตกฺวสฺส กาโย ฯเปฯ. (19)}

\csromanheader{2}{อยญฺหิ เทว กุมาโร สมวฏฺฏกฺขนฺโธ ฯเปฯ. (20)}
\csromanheader{2}{อยญฺหิ เทว กุมาโร รสคฺคสคฺคี ฯเปฯ. (21)}
\csromanheader{2}{อยญฺหิ เทว กุมาโร สีหหนุ ฯเปฯ. (22)}
\csromanheader{2}{อยญฺหิ เทว กุมาโร จตฺตาลีสทนฺโต ฯเปฯ. (23)}
\csromanheader{2}{อยญฺหิ เทว กุมาโร สมทนฺโต ฯเปฯ. (24)}
\csromanheader{2}{อยญฺหิ เทว กุมาโร อวิรฬทนฺโต ฯเปฯ. (25)}
\csromanheader{2}{อยญฺหิ เทว กุมาโร สุสุกฺกทาโฐ ฯเปฯ. (26)}
\csromanheader{2}{อยญฺหิ เทว กุมาโร ปหูตชิโวฺห ฯเปฯ. (27)}
\csromanheader{2}{อยญฺหิ เทว กุมาโร พฺรหฺมสฺสโร กรวีกภาณี ฯเปฯ. (28)}
\csromanheader{2}{อยญฺหิ เทว กุมาโร อภินีลเนตฺโต ฯเปฯ. (29)}
\csromanheader{2}{อยญฺหิ เทว กุมาโร โคปขุโม ฯเปฯ. (30)}
\prose{\csromanfolio{16}อิมสฺส เทว กุมารสฺส อุณฺณา ภมุกนฺตเร ชาตา โอทาตา มุทุตูลสนฺนิภา, ยมฺปิ อิมสฺส เทว กุมารสฺส อุณฺณา ภมุกนฺตเร ชาตา โอทาตา มุทุตูลสนฺนิภา, อิทมฺปิมสฺส มหาปุริสสฺส มหาปุริสลกฺขณํ ภวติ. (31)}

\prose{อยญฺหิ เทว กุมาโร อุณฺหีสสีโส, ยํปายํ เทว กุมาโร อุณฺหีสสีโส, อิทมฺปิมสฺส มหาปุริสสฺส มหาปุริสลกฺขณํ ภวติ. (32)}

\proseitem{36}{อิเมหิ โข อยํ เทว กุมาโร ทฺวตฺติํสมหาปุริสลกฺขเณหิ สมนฺนาคโต, เยหิ สมนฺนาคตสฺส มหาปุริสสฺส เทฺวว คติโย ภวนฺติ อนญฺญา.\csromanspacer{} สเจ อคารํ อชฺฌาวสติ, ราชา โหติ จกฺกวตฺตี ธมฺมิโก ธมฺมราชา จาตุรนฺโต วิชิตาวี ชนปทตฺถาวริยปฺปตฺโต สตฺตรตนสมนฺนาคโต.\csromanspacer{} ตสฺสิมานิ สตฺตรตนานิ ภวนฺติ.\csromanspacer{} เสยฺยถิทํ, จกฺกรตนํ หตฺถิรตนํ อสฺสรตนํ มณิรตนํ อิตฺถิรตนํ คหปติรตนํ ปริณายกรตนเมว สตฺตมํ.\csromanspacer{} ปโรสหสฺสํ โข ปนสฺส ปุตฺตา ภวนฺติ สูรา วีรงฺครูปา ปรเสนปฺปมทฺทนา, โส อิมํ ปถวิํ สาครปริยนฺตํ อทณฺเฑน อสตฺเถน ธมฺเมน\footnote{ธมฺเมน สเมน (สฺยา)} อภิวิชิย อชฺฌาวสติ.\csromanspacer{} สเจ โข \csromanfolio{17}ปน อคารสฺมา อนคาริยํ ปพฺพชติ, อรหํ โหติ สมฺมาสมฺพุทฺโธ โลเก วิวฏจฺฉโทติ.}

\csromanmatikaanchor{mtk.5}
\csromantocmarkref{section}{วิปสฺสีสมญฺญา}{mtk.5}
\bookmark[dest={mtk.5},level=1]{วิปสฺสีสมญฺญา}
\csromanheader{1}{\textbf{วิปสฺสีสมญฺญา}}
\proseitem{37}{อถ โข ภิกฺขเว พนฺธุมา ราชา เนมิตฺเต พฺราหฺมเณ อหเตหิ วตฺเถหิ อจฺฉาทาเปตฺวา\footnote{อจฺฉาเทตฺวา (สฺยา)} สพฺพกาเมหิ สนฺตปฺเปสิ.\csromanspacer{} อถ โข ภิกฺขเว พนฺธุมา ราชา วิปสฺสิสฺส กุมารสฺส ธาติโย อุปฏฺฐาปฺเปสิ.\csromanspacer{} อญฺญา ขีรํ ปาเยนฺติ, อญฺญา นฺหาเปนฺติ, อญฺญา ธาเรนฺติ, อญฺญา องฺเกน ปริหรนฺติ.\csromanspacer{} ชาตสฺส โข ปน ภิกฺขเว วิปสฺสิสฺส กุมารสฺส เสตจฺฉตฺตํ ธารยิตฺถ ทิวา เจว รตฺติํ จ “มา นํ สีตํ วา อุณฺหํ วา ติณํ วา รโช วา อุสฺสาโว วา พาธยิตฺถา”ติ.\csromanspacer{} ชาโต โข ปน ภิกฺขเว วิปสฺสี กุมาโร พหุโน ชนสฺส ปิโย อโหสิ มนาโป.\csromanspacer{} เสยฺยถาปิ ภิกฺขเว อุปฺปลํ วา ปทุมํ วา ปุณฺฑรีกํ วา พหุโน ชนสฺส ปิยํ มนาปํ, เอวเมว โข ภิกฺขเว วิปสฺสี กุมาโร พหุโน ชนสฺส ปิโย อโหสิ มนาโป.\csromanspacer{} สฺวาสฺสุทํ องฺเกเนว องฺกํ ปริหริยติ.}

\proseitem{38}{ชาโต โข ปน ภิกฺขเว วิปสฺสี กุมาโร มญฺชุสฺสโร จ\footnote{กุมาโร พฺรหฺมสฺสโร มญฺชุสฺสโร จ (สี, ก)} อโหสิ วคฺคุสฺสโร จ มธุรสฺสโร จ เปมนิยสฺสโร จ.\csromanspacer{} เสยฺยถาปิ ภิกฺขเว หิมวนฺเต ปพฺพเต กรวีกา นาม สกุณชาติ มญฺชุสฺสรา จ วคฺคุสฺสรา จ มธุรสฺสรา จ เปมนิยสฺสรา จ, เอวเมว โข ภิกฺขเว วิปสฺสี กุมาโร มญฺชุสฺสโร จ อโหสิ วคฺคุสฺสโร จ มธุรสฺสโร จ เปมนิยสฺสโร จ.}

\proseitem{39}{ชาตสฺส โข ปน ภิกฺขเว วิปสฺสิสฺส กุมารสฺส กมฺมวิปากชํ ทิพฺพจกฺขุ ปาตุรโหสิ.\csromanspacer{} เยน สุทํ\footnote{เยน ทูรํ (สฺยา)} สมนฺตา โยชนํ ปสฺสติ ทิวา เจว รตฺติํ จ.}

\proseitem{40}{ชาโต โข ปน ภิกฺขเว วิปสฺสี กุมาโร อนิมิสนฺโต เปกฺขติ เสยฺยถาปิ เทวา ตาวติํสา.\csromanspacer{} อนิมิสนฺโต กุมาโร เปกฺขตีติ โข ภิกฺขเว\footnote{อนิมิสนฺโต เปกฺขติ, ชาตสฺส โข ปน ภิกฺขเว (ก)} วิปสฺสิสฺส กุมารสฺส “วิปสฺสี วิปสฺสี” เตฺวว สมญฺญา อุทปาทิ.}

\proseitem{41}{\csromanfolio{18}อถ โข ภิกฺขเว พนฺธุมา ราชา อตฺถกรเณ\footnote{อฏฺฏ กรเณ (สฺยา)} นิสินฺโน วิปสฺสิํ กุมารํ องฺเก นิสีทาเปตฺวา อตฺเถ อนุสาสติ.\csromanspacer{} ตตฺร สุทํ ภิกฺขเว วิปสฺสี กุมาโร ปิตุ-องฺเก นิสินฺโน วิเจยฺย วิเจยฺย อตฺเถ ปนายติ ญาเยน\footnote{อฏฺเฏ ปนายติ ญาเณน (สฺยา)}.\csromanspacer{} วิเจยฺย วิเจยฺย กุมาโร อตฺเถ ปนายติ ญาเยนาติ โข ภิกฺขเว วิปสฺสิสฺส กุมารสฺส ภิยฺโยโส มตฺตาย “วิปสฺสี วิปสฺสี” เตฺวว สมญฺญา อุทปาทิ.}

\proseitem{42}{อถ โข ภิกฺขเว พนฺธุมา ราชา วิปสฺสิสฺส กุมารสฺส ตโย ปาสาเท การาเปสิ, เอกํ วสฺสิกํ เอกํ เหมนฺติกํ เอกํ คิมฺหิกํ, ปญฺจ กามคุณานิ อุปฏฺฐาเปสิ.\csromanspacer{} ตตฺร สุทํ ภิกฺขเว วิปสฺสี กุมาโร วสฺสิเก ปาสาเท จตฺตาโร มาเส\footnote{วสฺสิเก ปาสาเท วสฺสิเก จตฺตาโร มาเส (สี, อิ)} นิปฺปุริเสหิ ตูริเยหิ ปริจารยมาโน น เหฏฺฐาปาสาทํ โอโรหตีติ.}

\csromanheader{2}{ปฐมภาณวาโร.}
\csromansectionrule
\csromanmatikaanchor{mtk.6}
\csromantocmarkref{section}{ชิณฺณปุริส}{mtk.6}
\bookmark[dest={mtk.6},level=1]{ชิณฺณปุริส}
\csromanheader{1}{\textbf{ชิณฺณปุริส}}
\proseitem{43}{อถ โข ภิกฺขเว วิปสฺสี กุมาโร พหูนํ วสฺสานํ พหูนํ วสฺสสตานํ พหูนํ วสฺสสหสฺสานํ อจฺจเยน สารถิํ อามนฺเตสิ “โยเชหิ สมฺม สารถิ ภทฺทานิ ภทฺทานิ ยานานิ อุยฺยานภูมิํ คจฺฉาม สุภูมิทสฺสนายา”ติ. “เอวํ เทวา”ติ โข ภิกฺขเว สารถิ วิปสฺสิสฺส กุมารสฺส ปฏิสฺสุตฺวา ภทฺทานิ ภทฺทานิ ยานานิ โยเชตฺวา วิปสฺสิสฺส กุมารสฺส ปฏิเวเทสิ “ยุตฺตานิ โข เต เทว ภทฺทานิ ภทฺทานิ ยานานิ, ยสฺสทานิ กาลํ มญฺญสี”ติ.\csromanspacer{} อถ โข ภิกฺขเว วิปสฺสี กุมาโร ภทฺทํ ภทฺทํ ยานํ\footnote{ภทฺรํ ยานํ (สฺยา), ภทฺทํ ยานํ (อิ)} อภิรุหิตฺวา ภทฺเทหิ ภทฺเทหิ ยาเนหิ อุยฺยานภูมิํ นิยฺยาสิ.}

\proseitem{44}{อทฺทสา โข ภิกฺขเว วิปสฺสี กุมาโร อุยฺยานภูมิํ นิยฺยนฺโต ปุริสํ ชิณฺณํ โคปานสิวงฺกํ โภคฺคํ\footnote{ภคฺคํ (สฺยา)} ทณฺฑปรายนํ ปเวธมานํ คจฺฉนฺตํ \csromanfolio{19}อาตุรํ คตโยพฺพนํ, ทิสฺวา สารถิํ อามนฺเตสิ “อยํ ปน สมฺม สารถิ ปุริโส กิํกโต, เกสาปิสฺส น ยถา อญฺเญสํ, กาโยปิสฺส น ยถา อญฺเญสนฺ”ติ.\csromanspacer{} เอโส โข เทว ชิณฺโณ นามาติ.\csromanspacer{} กิํ ปเนโส สมฺม สารถิ ชิณฺโณ นามาติ? เอโส โข เทว ชิณฺโณ นาม น ทานิ เตน จิรํ ชีวิตพฺพํ ภวิสฺสตีติ.\csromanspacer{} กิํ ปน สมฺม สารถิ อหํปิ ชราธมฺโม ชรํ อนตีโตติ? ตฺวญฺจ เทว มยญฺจมฺห สพฺเพ ชราธมฺมา ชรํ อนตีตาติ.\csromanspacer{} เตน หิ สมฺม สารถิ อลํ ทานชฺช อุยฺยานภูมิยา, อิโตว อนฺเตปุรํ ปจฺจนิยฺยาหีติ. “เอวํ เทวา”ติ โข ภิกฺขเว สารถิ วิปสฺสิสฺส กุมารสฺส ปฏิสฺสุตฺวา ตโตว อนฺเตปุรํ ปจฺจนิยฺยาสิ.\csromanspacer{} ตตฺร สุทํ ภิกฺขเว วิปสฺสี กุมาโร อนฺเตปุรํ คโต ทุกฺขี ทุมฺมโน ปชฺฌายติ “ธิรตฺถุ กิร โภ ชาติ นาม, ยตฺร หิ นาม ชาตสฺส ชรา ปญฺญายิสฺสตี”ติ.}

\proseitem{45}{อถ โข ภิกฺขเว พนฺธุมา ราชา สารถิํ อามนฺตาเปตฺวา เอตทโวจ “กจฺจิ สมฺม สารถิ กุมาโร อุยฺยานภูมิยา อภิรมิตฺถ, กจฺจิ สมฺม สารถิ กุมาโร อุยฺยานภูมิยา อตฺตมโน อโหสี”ติ.\csromanspacer{} น โข เทว กุมาโร อุยฺยานภูมิยา อภิรมิตฺถ, น โข เทว กุมาโร อุยฺยานภูมิยา อตฺตมโน อโหสีติ.\csromanspacer{} กิํ ปน สมฺม สารถิ อทฺทส กุมาโร อุยฺยานภูมิํ นิยฺยนฺโตติ? อทฺทสา โข เทว กุมาโร อุยฺยานภูมิํ นิยฺยนฺโต ปุริสํ ชิณฺณํ โคปานสิวงฺกํ โภคฺคํ ทณฺฑปรายนํ ปเวธมานํ คจฺฉนฺตํ อาตุรํ คตโยพฺพนํ, ทิสฺวา มํ เอตทโวจ “อยํ ปน สมฺม สารถิ ปุริโส กิํกโต, เกสาปิสฺส น ยถา อญฺเญสํ, กาโยปิสฺส น ยถา อญฺเญสนฺ”ติ.\csromanspacer{} เอโส โข เทว ชิณฺโณ นามาติ.\csromanspacer{} กิํ ปเนโส สมฺม สารถิ ชิณฺโณ นามาติ? เอโส โข เทว ชิณฺโณ นาม น ทานิ เตน จิรํ ชีวิตพฺพํ ภวิสฺสตีติ.\csromanspacer{} กิํ ปน สมฺม สารถิ อหํปิ ชราธมฺโม ชรํ อนตีโตติ? ตฺวญฺจ เทว มยญฺจมฺห สพฺเพ ชราธมฺมา ชรํ อนตีตาติ.\csromanspacer{} เตน หิ สมฺม สารถิ อลํ ทานชฺช อุยฺยานภูมิยา, อิโตว อนฺเตปุรํ ปจฺจนิยฺยาหีติ. “เอวํ เทวา”ติ โข อหํ เทว วิปสฺสิสฺส กุมารสฺส ปฏิสฺสุตฺวา ตโตว อนฺเตปุรํ ปจฺจนิยฺยาสิํ.\csromanspacer{} โส โข เทว กุมาโร อนฺเตปุรํ คโต ทุกฺขี ทุมฺมโน ปชฺฌายติ “ธิรตฺถุ กิร โภ ชาติ นาม, ยตฺร หิ นาม ชาตสฺส ชรา ปญฺญายิสฺสตี”ติ.}

\csromanmatikaanchor{mtk.7}
\csromantocmarkref{section}{พฺยาธิตปุริส}{mtk.7}
\bookmark[dest={mtk.7},level=1]{พฺยาธิตปุริส}
\csromanheader{1}{\textbf{พฺยาธิตปุริส}}
\proseitem{46}{\csromanfolio{20}อถ โข ภิกฺขเว พนฺธุมสฺส รญฺโญ เอตทโหสิ “มาเหว โข วิปสฺสี กุมาโร น รชฺชํ กาเรสิ, มาเหว วิปสฺสี กุมาโร อคารสฺมา อนคาริยํ ปพฺพชิ, มาเหว เนมิตฺตานํ พฺราหฺมณานํ สจฺจํ อสฺส วจนนฺ”ติ.\csromanspacer{} อถ โข ภิกฺขเว พนฺธุมา ราชา วิปสฺสิสฺส กุมารสฺส ภิยฺโยโส มตฺตาย ปญฺจ กามคุณานิ อุปฏฺฐาเปสิ, ยถา วิปสฺสี กุมาโร รชฺชํ กเรยฺย, ยถา วิปสฺสี กุมาโร น อคารสฺมา อนคาริยํ ปพฺพเชยฺย, ยถา เนมิตฺตานํ พฺราหฺมณานํ มิจฺฉา อสฺส วจนนฺติ.}

\prose{ตตฺร สุทํ ภิกฺขเว วิปสฺสี กุมาโร ปญฺจหิ กามคุเณหิ สมปฺปิโต สมงฺคีภูโต ปริจาเรติ.\csromanspacer{} อถ โข ภิกฺขเว วิปสฺสี กุมาโร พหูนํ วสฺสานํ ฯเปฯ.}

\proseitem{47}{อทฺทสา โข ภิกฺขเว วิปสฺสี กุมาโร อุยฺยานภูมิํ นิยฺยนฺโต ปุริสํ อาพาธิกํ ทุกฺขิตํ พาฬฺหคิลานํ สเก มุตฺตกรีเส ปลิปนฺนํ เสมานํ\footnote{สยมานํ (สฺยา, ก)} อญฺเญหิ วุฏฺฐาปิยมานํ อญฺเญหิ สํเวสิยมานํ, ทิสฺวา สารถิํ อามนฺเตสิ “อยํ ปน สมฺม สารถิ ปุริโส กิํกโต, อกฺขีนิปิสฺส น ยถา อญฺเญสํ, สโรปิสฺส\footnote{สิโรปิสฺส (สฺยา)} น ยถา อญฺเญสนฺ”ติ.\csromanspacer{} เอโส โข เทว พฺยาธิโต นามาติ.\csromanspacer{} กิํ ปเนโส สมฺม สารถิ พฺยาธิโต นามาติ? เอโส โข เทว พฺยาธิโต นาม อปฺเปว นาม ตมฺหา อาพาธา วุฏฺฐเหยฺยาติ.\csromanspacer{} กิํ ปน สมฺม สารถิ อหํปิ พฺยาธิธมฺโม พฺยาธิํ อนตีโตติ? ตฺวญฺจ เทว มยญฺจมฺห สพฺเพ พฺยาธิธมฺมา พฺยาธิํ อนตีตาติ.\csromanspacer{} เตนหิ สมฺม สารถิ อลํ ทานชฺช อุยฺยานภูมิยา, อิโตว อนฺเตปุรํ ปจฺจนิยฺยาหีติ. “เอวํ เทวา”ติ โข ภิกฺขเว สารถิ วิปสฺสิสฺส กุมารสฺส ปฏิสฺสุตฺวา ตโตว อนฺเตปุรํ ปจฺจนิยฺยาสิ.\csromanspacer{} ตตฺร สุทํ ภิกฺขเว วิปสฺสี กุมาโร อนฺเตปุรํ คโต ทุกฺขี ทุมฺมโน ปชฺฌายติ “ธิรตฺถุ กิร โภ ชาติ นาม, ยตฺร หิ นาม ชาตสฺส ชรา ปญฺญายิสฺสติ, พฺยาธิ ปญฺญายิสฺสตี”ติ.}

\proseitem{48}{อถ โข ภิกฺขเว พนฺธุมา ราชา สารถิํ อามนฺตาเปตฺวา เอตทโวจ “กจฺจิ สมฺม สารถิ กุมาโร อุยฺยานภูมิยา อภิรมิตฺถ, กจฺจิ \csromanfolio{21}สมฺม สารถิ กุมาโร อุยฺยานภูมิยา อตฺตมโน อโหสี”ติ.\csromanspacer{} น โข เทว กุมาโร อุยฺยานภูมิยา อภิรมิตฺถ, น โข เทว กุมาโร อุยฺยานภูมิยา อตฺตมโน อโหสีติ.\csromanspacer{} กิํ ปน สมฺม สารถิ อทฺทส กุมาโร อุยฺยานภูมิํ นิยฺยนฺโตติ? อทฺทสา โข เทว กุมาโร อุยฺยานภูมิํ นิยฺยนฺโต ปุริสํ อาพาธิกํ ทุกฺขิตํ พาฬฺหคิลานํ สเก มุตฺตกรีเส ปลิปนฺนํ เสมานํ อญฺเญหิ วุฏฺฐาปิยมานํ อญฺเญหิ สํเวสิยมานํ, ทิสฺวา มํ เอตทโวจ “อยํ ปน สมฺม สารถิ ปุริโส กิํกโต, อกฺขีนิปิสฺส น ยถา อญฺเญสํ, สโรปิสฺส น ยถา อญฺเญสนฺ”ติ.\csromanspacer{} เอโส โข เทว พฺยาธิโต นามาติ.\csromanspacer{} กิํ ปเนโส สมฺม สารถิ พฺยาธิโต นามาติ? เอโส โข เทว พฺยาธิโต นาม อปฺเปว นาม ตมฺหา อาพาธา วุฏฺฐเหยฺยาติ.\csromanspacer{} กิํ ปน สมฺม สารถิ อหํปิ พฺยาธิธมฺโม พฺยาธิํ อนตีโตติ? ตฺวญฺจ เทว มยญฺจมฺห สพฺเพ พฺยาธิธมฺมา พฺยาธิํ อนตีตาติ.\csromanspacer{} เตน หิ สมฺม สารถิ อลํ ทานชฺช อุยฺยานภูมิยา, อิโตว อนฺเตปุรํ ปจฺจนิยฺยาหีติ. “เอวํ เทวา”ติ โข อหํ เทว วิปสฺสิสฺส กุมารสฺส ปฏิสฺสุตฺวา ตโตว อนฺเตปุรํ ปจฺจนิยฺยาสิํ.\csromanspacer{} โส โข เทว กุมาโร อนฺเตปุรํ คโต ทุกฺขี ทุมฺมโน ปชฺฌายติ “ธิรตฺถุ กิร โภ ชาติ นาม, ยตฺร หิ นาม ชาตสฺส ชรา ปญฺญายิสฺสติ, พฺยาธิ ปญฺญายิสฺสตี”ติ.}

\csromanmatikaanchor{mtk.8}
\csromantocmarkref{section}{กาลงฺกตปุริส}{mtk.8}
\bookmark[dest={mtk.8},level=1]{กาลงฺกตปุริส}
\csromanheader{1}{\textbf{กาลงฺกตปุริส}}
\proseitem{49}{อถ โข ภิกฺขเว พนฺธุมสฺส รญฺโญ เอตทโหสิ “มาเหว โข วิปสฺสี กุมาโร น รชฺชํ กาเรสิ, มาเหว วิปสฺสี กุมาโร อคารสฺมา อนคาริยํ ปพฺพชิ, มาเหว เนมิตฺตานํ พฺราหฺมณานํ สจฺจํ อสฺส วจนนฺ”ติ.\csromanspacer{} อถ โข ภิกฺขเว พนฺธุมา ราชา วิปสฺสิสฺส กุมารสฺส ภิยฺโยโส มตฺตาย ปญฺจ กามคุณานิ อุปฏฺฐาเปสิ, ยถา วิปสฺสี กุมาโร รชฺชํ กเรยฺย, ยถา วิปสฺสี กุมาโร น อคารสฺมา อนคาริยํ ปพฺพเชยฺย, ยถา เนมิตฺตานํ พฺราหฺมณานํ มิจฺฉา อสฺส วจนนฺติ.}

\prose{ตตฺร สุทํ ภิกฺขเว วิปสฺสี กุมาโร ปญฺจหิ กามคุเณหิ สมปฺปิโต สมงฺคีภูโต ปริจาเรติ.\csromanspacer{} อถ โข ภิกฺขเว วิปสฺสี กุมาโร พหูนํ วสฺสานํ ฯเปฯ.}

\proseitem{50}{\csromanfolio{22}อทฺทสา โข ภิกฺขเว วิปสฺสี กุมาโร อุยฺยานภูมิํ นิยฺยนฺโต มหาชนกายํ สนฺนิปติตํ นานารตฺตานํ จ ทุสฺสานํ วิลาตํ กยิรมานํ, ทิสฺวา สารถิํ อามนฺเตสิ “กิํ นุ โข โส สมฺม สารถิ มหาชนกาโย สนฺนิปติโต นานารตฺตานํ จ ทุสฺสานํ วิลาตํ กยิรตี”ติ.\csromanspacer{} เอโส โข เทว กาลงฺกโต นามาติ.\csromanspacer{} เตน หิ สมฺม สารถิ เยน โส กาลงฺกโต เตน รถํ เปเสหีติ. “เอวํ เทวา”ติ โข ภิกฺขเว สารถิ วิปสฺสิสฺส กุมารสฺส ปฏิสฺสุตฺวา เยน โส กาลงฺกโต เตน รถํ เปเสสิ.\csromanspacer{} อทฺทสา โข ภิกฺขเว วิปสฺสี กุมาโร เปตํ กาลงฺกตํ, ทิสฺวา สารถิํ อามนฺเตสิ “กิํ ปนายํ สมฺม สารถิ กาลงฺกโต นามา”ติ.\csromanspacer{} เอโส โข เทว กาลงฺกโต นาม น ทานิ ตํ ทกฺขนฺติ มาตา วา ปิตา วา อญฺเญ วา ญาติสาโลหิตา, โสปิ น ทกฺขิสฺสติ มาตรํ วา ปิตรํ วา อญฺเญ วา ญาติสาโลหิเตติ. “กิํ ปน สมฺม สารถิ อหํปิ มรณธมฺโม มรณํ อนตีโต, มํปิ น ทกฺขนฺติ เทโว วา เทวี วา อญฺเญ วา ญาติสาโลหิตา, อหํปิ น ทกฺขิสฺสามิ เทวํ วา เทวิํ วา อญฺเญ วา ญาติสาโลหิเต”ติ.\csromanspacer{} ตฺวญฺจ เทว มยญฺจมฺห สพฺเพ มรณธมฺมา มรณํ อนตีตา, ตํปิ น ทกฺขนฺติ เทโว วา เทวี วา อญฺเญ วา ญาติสาโลหิตา, ตฺวํปิ น ทกฺขิสฺสสิ เทวํ วา เทวิํ วา อญฺเญ วา ญาติสาโลหิเตติ.\csromanspacer{} เตน หิ สมฺม สารถิ อลํ ทานชฺช อุยฺยานภูมิยา, อิโตว อนฺเตปุรํ ปจฺจนิยฺยาหีติ. “เอวํ เทวา”ติ โข ภิกฺขเว สารถิ วิปสฺสิสฺส กุมารสฺส ปฏิสฺสุตฺวา ตโตว อนฺเตปุรํ ปจฺจนิยฺยาสิ.\csromanspacer{} ตตฺร สุทํ ภิกฺขเว วิปสฺสี กุมาโร อนฺเตปุรํ คโต ทุกฺขี ทุมฺมโน ปชฺฌายติ “ธิรตฺถุ กิร โภ ชาติ นาม, ยตฺร หิ นาม ชาตสฺส ชรา ปญฺญายิสฺสติ, พฺยาธิ ปญฺญายิสฺสติ, มรณํ ปญฺญายิสฺสตี”ติ.}

\proseitem{51}{อถ โข ภิกฺขเว พนฺธุมา ราชา สารถิํ อามนฺตาเปตฺวา เอตทโวจ “กจฺจิ สมฺม สารถิ กุมาโร อุยฺยานภูมิยา อภิรมิตฺถ, กจฺจิ สมฺม สารถิ กุมาโร อุยฺยานภูมิยา อตฺตมโน อโหสี”ติ.\csromanspacer{} น โข เทว กุมาโร อุยฺยานภูมิยา อภิรมิตฺถ, น โข เทว กุมาโร อุยฺยานภูมิยา อตฺตมโน อโหสีติ.\csromanspacer{} กิํ ปน สมฺม สารถิ อทฺทส กุมาโร อุยฺยานภูมิํ นิยฺยนฺโตติ.\csromanspacer{} อทฺทสา โข เทว กุมาโร อุยฺยานภูมิํ นิยฺยนฺโต มหาชนกายํ สนฺนิปติตํ นานารตฺตานํ จ ทุสฺสานํ วิลาตํ กยิรมานํ, ทิสฺวา มํ เอตทโวจ “กิํ นุ โข \csromanfolio{23}โส สมฺม สารถิ มหาชนกาโย สนฺนิปติโต นานารตฺตานํ จ ทุสฺสานํ วิลาตํ กยิรตี”ติ.\csromanspacer{} เอโส โข เทว กาลงฺกโต นามาติ.\csromanspacer{} เตน หิ สมฺม สารถิ เยน โส กาลงฺกโต เตน รถํ เปเสหีติ. “เอวํ เทวา”ติ โข อหํ เทว วิปสฺสิสฺส กุมารสฺส ปฏิสฺสุตฺวา เยน โส กาลงฺกโต เตน รถํ เปเสสิํ.\csromanspacer{} อทฺทสา โข เทว กุมาโร เปตํ กาลงฺกตํ, ทิสฺวา มํ เอตทโวจ “กิํ ปนายํ สมฺม สารถิ กาลงฺกโต นามา”ติ.\csromanspacer{} เอโส โข เทว กาลงฺกโต นาม น ทานิ ตํ ทกฺขนฺติ มาตา วา ปิตา วา อญฺเญ วา ญาติสาโลหิตา, โสปิ น ทกฺขิสฺสติ มาตรํ วา ปิตรํ วา อญฺเญ วา ญาติสาโลหิเตติ.\csromanspacer{} กิํ ปน สมฺม สารถิ อหํปิ มรณธมฺโม มรณํ อนตีโต, มํปิ น ทกฺขนฺติ เทโว วา เทวี วา อญฺเญ วา ญาติสาโลหิตา, อหํปิ น ทกฺขิสฺสามิ เทวํ วา เทวิํ วา อญฺเญ วา ญาติสาโลหิเตติ.\csromanspacer{} ตฺวญฺจ เทว มยญฺจมฺห สพฺเพ มรณธมฺมา มรณํ อนตีตา, ตํปิ น ทกฺขนฺติ เทโว วา เทวี วา อญฺเญ วา ญาติสาโลหิตา, ตฺวํปิ น ทกฺขิสฺสสิ เทวํ วา เทวิํ วา อญฺเญ วา ญาติสาโลหิเตติ.\csromanspacer{} เตน หิ สมฺม สารถิ อลํ ทานชฺช อุยฺยานภูมิยา, อิโตว อนฺเตปุรํ ปจฺจนิยฺยาหีติ. “เอวํ เทวา”ติ โข อหํ เทว วิปสฺสิสฺส กุมารสฺส ปฏิสฺสุตฺวา ตโตว อนฺเตปุรํ ปจฺจนิยฺยาสิํ.\csromanspacer{} โส โข เทว กุมาโร อนฺเตปุรํ คโต ทุกฺขี ทุมฺมโน ปชฺฌายติ “ธิรตฺถุ กิร โภ ชาติ นาม, ยตฺร หิ นาม ชาตสฺส ชรา ปญฺญายิสฺสติ, พฺยาธิ ปญฺญายิสฺสติ, มรณํ ปญฺญายิสฺสตี”ติ.}

\csromanmatikaanchor{mtk.9}
\csromantocmarkref{section}{ปพฺพชิต}{mtk.9}
\bookmark[dest={mtk.9},level=1]{ปพฺพชิต}
\csromanheader{1}{\textbf{ปพฺพชิต}}
\proseitem{52}{อถ โข ภิกฺขเว พนฺธุมสฺส รญฺโญ เอตทโหสิ “มาเหว โข วิปสฺสี กุมาโร น รชฺชํ กาเรสิ, มาเหว วิปสฺสี กุมาโร อคารสฺมา อนคาริยํ ปพฺพชิ, มาเหว เนมิตฺตานํ พฺราหฺมณานํ สจฺจํ อสฺส วจนนฺ”ติ.\csromanspacer{} อถ โข ภิกฺขเว พนฺธุมา ราชา วิปสฺสิสฺส กุมารสฺส ภิยฺโยโส มตฺตาย ปญฺจ กามคุณานิ อุปฏฺฐาเปสิ, ยถา วิปสฺสี กุมาโร รชฺชํ กเรยฺย, ยถา วิปสฺสี กุมาโร น อคารสฺมา อนคาริยํ ปพฺพเชยฺย, ยถา เนมิตฺตานํ พฺราหฺมณานํ มิจฺฉา อสฺส วจนนฺติ.}

\prose{ตตฺร สุทํ ภิกฺขเว วิปสฺสี กุมาโร ปญฺจหิ กามคุเณหิ สมปฺปิโต สมงฺคีภูโต ปริจาเรติ.\csromanspacer{} อถ โข ภิกฺขเว วิปสฺสี กุมาโร พหูนํ \csromanfolio{24}วสฺสานํ พหูนํ วสฺสสตานํ พหูนํ วสฺสสหสฺสานํ อจฺจเยน สารถิํ อามนฺเตสิ “โยเชหิ สมฺม สารถิ ภทฺทานิ ภทฺทานิ ยานานิ, อุยฺยานภูมิํ คจฺฉาม สุภูมิทสฺสนายา”ติ. “เอวํ เทวา”ติ โข ภิกฺขเว สารถิ วิปสฺสิสฺส กุมารสฺส ปฏิสฺสุตฺวา ภทฺทานิ ภทฺทานิ ยานานิ โยเชตฺวา วิปสฺสิสฺส กุมารสฺส ปฏิเวเทสิ “ยุตฺตานิ โข เต เทว ภทฺทานิ ภทฺทานิ ยานานิ, ยสฺสทานิ กาลํ มญฺญสี”ติ.\csromanspacer{} อถ โข ภิกฺขเว วิปสฺสี กุมาโร ภทฺทํ ภทฺทํ ยานํ อภิรุหิตฺวา ภทฺเทหิ ภทฺเทหิ ยาเนหิ อุยฺยานภูมิํ นิยฺยาสิ.}

\proseitem{53}{อทฺทสา โข ภิกฺขเว วิปสฺสี กุมาโร อุยฺยานภูมิํ นิยฺยนฺโต ปุริสํ ภณฺฑุํ ปพฺพชิตํ กาสายวสนํ, ทิสฺวา สารถิํ อามนฺเตสิ “อยํ ปน สมฺม สารถิ ปุริโส กิํกโต, สีสํปิสฺส น ยถา อญฺเญสํ, วตฺถานิปิสฺส น ยถา อญฺเญสนฺ”ติ.\csromanspacer{} เอโส โข เทว ปพฺพชิโต นามาติ.\csromanspacer{} กิํ ปเนโส สมฺม สารถิ ปพฺพชิโต นามาติ.\csromanspacer{} เอโส โข เทว ปพฺพชิโต นาม สาธุ ธมฺมจริยา สาธุ สมจริยา\footnote{สมฺมจริยา (ก)} สาธุ กุสลกิริยา\footnote{กุสลจริยา (สฺยา)} สาธุ ปุญฺญกิริยา สาธุ อวิหิํสา สาธุ ภูตานุกมฺปาติ.\csromanspacer{} สาธุ โข โส สมฺม สารถิ ปพฺพชิโต นาม สาธุ ธมฺมจริยา สาธุ สมจริยา สาธุ กุสลกิริยา สาธุ ปุญฺญกิริยา สาธุ อวิหิํสา สาธุ ภูตานุกมฺปา, เตน หิ สมฺม สารถิ เยน โส ปพฺพชิโต เตน รถํ เปเสหีติ. “เอวํ เทวา”ติ โข ภิกฺขเว สารถิ วิปสฺสิสฺส กุมารสฺส ปฏิสฺสุตฺวา เยน โส ปพฺพชิโต เตน รถํ เปเสสิ.\csromanspacer{} อถ โข ภิกฺขเว วิปสฺสี กุมาโร ตํ ปพฺพชิตํ เอตทโวจ “ตฺวํ ปน สมฺม กิํกโต, สีสํปิ เต น ยถา อญฺเญสํ, วตฺถานิปิ เต น ยถา อญฺเญสนฺ”ติ.\csromanspacer{} อหํ โข เทว ปพฺพชิโต นามาติ.\csromanspacer{} กิํ ปน ตฺวํ สมฺม ปพฺพชิโต นามาติ.\csromanspacer{} อหํ โข เทว ปพฺพชิโต นาม สาธุ ธมฺมจริยา สาธุ สมจริยา สาธุ กุสลกิริยา สาธุ ปุญฺญกิริยา สาธุ อวิหิํสา สาธุ ภูตานุกมฺปาติ.\csromanspacer{} สาธุ โข ตฺวํ สมฺม ปพฺพชิโต นาม สาธุ ธมฺมจริยา สาธุ สมจริยา สาธุ กุสลกิริยา สาธุ ปุญฺญกิริยา สาธุ อวิหิํสา สาธุ ภูตานุกมฺปาติ.}

\csromanmatikaanchor{mtk.10}
\csromantocmarkref{section}{โพธิสตฺตปพฺพชฺชา}{mtk.10}
\bookmark[dest={mtk.10},level=1]{โพธิสตฺตปพฺพชฺชา}
\csromanheader{1}{\textbf{โพธิสตฺตปพฺพชฺชา}}
\proseitem{54}{\csromanfolio{25}อถ โข ภิกฺขเว วิปสฺสี กุมาโร สารถิํ อามนฺเตสิ “เตน หิ สมฺม สารถิ รถํ อาทาย อิโตว อนฺเตปุรํ ปจฺจนิยฺยาหิ.\csromanspacer{} อหํ ปน อิเธว เกสมสฺสุํ โอหาเรตฺวา กาสายานิ วตฺถานิ อจฺฉาเทตฺวา อคารสฺมา อนคาริยํ ปพฺพชิสฺสามี”ติ. “เอวํ เทวา”ติ โข ภิกฺขเว สารถิ วิปสฺสิสฺส กุมารสฺส ปฏิสฺสุตฺวา รถํ อาทาย ตโตว อนฺเตปุรํ ปจฺจนิยฺยาสิ.\csromanspacer{} วิปสฺสี ปน กุมาโร ตตฺเถว เกสมสฺสุํ โอหาเรตฺวา กาสายานิ วตฺถานิ อจฺฉาเทตฺวา อคารสฺมา อนคาริยํ ปพฺพชิ.}

\csromanmatikaanchor{mtk.11}
\csromantocmarkref{section}{มหาชนกาย-อนุปพฺพชฺชา}{mtk.11}
\bookmark[dest={mtk.11},level=1]{มหาชนกาย-อนุปพฺพชฺชา}
\csromanheader{1}{\textbf{มหาชนกาย-อนุปพฺพชฺชา}}
\proseitem{55}{อสฺโสสิ โข ภิกฺขเว พนฺธุมติยา ราชธานิยา มหาชนกาโย จตุราสีติ ปาณสหสฺสานิ “วิปสฺสี กิร กุมาโร เกสมสฺสุํ โอหาเรตฺวา กาสายานิ วตฺถานิ อจฺฉาเทตฺวา อคารสฺมา อนคาริยํ ปพฺพชิโต”ติ, สุตฺวาน เตสํ เอตทโหสิ “น หิ นูน โส โอรโก ธมฺมวินโย, น สา โอรกา\footnote{โอริกา (สี, สฺยา)} ปพฺพชฺชา, ยตฺถ วิปสฺสี กุมาโร เกสมสฺสุํ โอหาเรตฺวา กาสายานิ วตฺถานิ อจฺฉาเทตฺวา อคารสฺมา อนคาริยํ ปพฺพชิโต, วิปสฺสีปิ นาม กุมาโร เกสมสฺสุํ โอหาเรตฺวา กาสายานิ วตฺถานิ อจฺฉาเทตฺวา อคารสฺมา อนคาริยํ ปพฺพชิสฺสติ, กิมงฺคํ\footnote{กิมงฺค (สี)} ปน มยนฺ”ติ.}

\prose{อถ โข โส ภิกฺขเว มหาชนกาโย\footnote{มหาชนกายา (สฺยา)} จตุราสีติ ปาณสหสฺสานิ เกสมสฺสุํ โอหาเรตฺวา กาสายานิ วตฺถานิ อจฺฉาเทตฺวา วิปสฺสิํ โพธิสตฺตํ อคารสฺมา อนคาริยํ ปพฺพชิตํ อนุปพฺพชิํสุ.\csromanspacer{} ตาย สุทํ ภิกฺขเว ปริสาย ปริวุโต วิปสฺสี โพธิสตฺโต คามนิคมชนปทราชธานีสุ จาริกํ จรติ.}

\proseitem{56}{อถ โข ภิกฺขเว วิปสฺสิสฺส โพธิสตฺตสฺส รโหคตสฺส ปฏิสลฺลีนสฺส เอวํ เจตโส ปริวิตกฺโก อุทปาทิ “น โข เมตํ\footnote{น โข ปเนตํ (สฺยา)} ปติรูปํ โยหํ อากิณฺโณ วิหรามิ, ยํนูนาหํ เอโก คณมฺหา วูปกฏฺโฐ วิหเรยฺยนฺ”ติ.\csromanspacer{} อถ โข ภิกฺขเว วิปสฺสี โพธิสตฺโต อปเรน สมเยน \csromanfolio{26}เอโก คณมฺหา วูปกฏฺโฐ วิหาสิ, อญฺเญเนว ตานิ จตุราสีติ ปพฺพชิตสหสฺสานิ อคมํสุ, อญฺเญน มคฺเคน วิปสฺสี โพธิสตฺโต.}

\csromanmatikaanchor{mtk.12}
\csromantocmarkref{section}{โพธิสตฺต-อภินิเวส}{mtk.12}
\bookmark[dest={mtk.12},level=1]{โพธิสตฺต-อภินิเวส}
\csromanheader{1}{\textbf{โพธิสตฺต-อภินิเวส}}
\proseitem{57}{อถ โข ภิกฺขเว วิปสฺสิสฺส โพธิสตฺตสฺส วาสูปคตสฺส รโหคตสฺส ปฏิสลฺลีนสฺส เอวํ เจตโส ปริวิตกฺโก อุทปาทิ “กิจฺฉํ วตายํ โลโก อาปนฺโน ชายติ จ ชียติ จ มียติ จ\footnote{ชิยฺยติ จ มิยฺยติ จ (ก)} จวติ จ อุปปชฺชติ จ, อถ จ ปนิมสฺส ทุกฺขสฺส นิสฺสรณํ นปฺปชานาติ ชรามรณสฺส, กุทาสฺสุ นาม อิมสฺส ทุกฺขสฺส นิสฺสรณํ ปญฺญายิสฺสติ ชรามรณสฺสา”ติ.}

\prose{อถ โข ภิกฺขเว วิปสฺสิสฺส โพธิสตฺตสฺส เอตทโหสิ “กิมฺหิ นุ โข สติ ชรามรณํ โหติ กิํ ปจฺจยา ชรามรณนฺ”ติ.\csromanspacer{} อถ โข ภิกฺขเว วิปสฺสิสฺส โพธิสตฺตสฺส โยนิโส มนสิการา อหุ ปญฺญาย อภิสมโย “ชาติยา โข สติ ชรามรณํ โหติ, ชาติปจฺจยา ชรามรณนฺ”ติ.}

\prose{อถ โข ภิกฺขเว วิปสฺสิสฺส โพธิสตฺตสฺส เอตทโหสิ “กิมฺหิ นุ โข สติ ชาติ โหติ กิํ ปจฺจยา ชาตี”ติ.\csromanspacer{} อถ โข ภิกฺขเว วิปสฺสิสฺส โพธิสตฺตสฺส โยนิโส มนสิการา อหุ ปญฺญาย อภิสมโย “ภเว โข สติ ชาติ โหติ, ภวปจฺจยา ชาตี”ติ.}

\prose{อถ โข ภิกฺขเว วิปสฺสิสฺส โพธิสตฺตสฺส เอตทโหสิ “กิมฺหิ นุ โข สติ ภโว โหติ กิํ ปจฺจยา ภโว”ติ.\csromanspacer{} อถ โข ภิกฺขเว วิปสฺสิสฺส โพธิสตฺตสฺส โยนิโส มนสิการา อหุ ปญฺญาย อภิสมโย “อุปาทาเน โข สติ ภโว โหติ, อุปาทานปจฺจยา ภโว”ติ.}

\prose{อถ โข ภิกฺขเว วิปสฺสิสฺส โพธิสตฺตสฺส เอตทโหสิ “กิมฺหิ นุ โข สติ อุปาทานํ โหติ กิํ ปจฺจยา อุปาทานนฺ”ติ.\csromanspacer{} อถ โข ภิกฺขเว วิปสฺสิสฺส โพธิสตฺตสฺส โยนิโส มนสิการา อหุ ปญฺญาย อภิสมโย “ตณฺหาย โข สติ อุปาทานํ โหติ, ตณฺหาปจฺจยา อุปาทานนฺ”ติ.}

\prose{อถ โข ภิกฺขเว วิปสฺสิสฺส โพธิสตฺตสฺส เอตทโหสิ “กิมฺหิ นุ โข สติ ตณฺหา โหติ กิํ ปจฺจยา ตณฺหา”ติ.\csromanspacer{} อถ โข ภิกฺขเว \csromanfolio{27}วิปสฺสิสฺส โพธิสตฺตสฺส โยนิโส มนสิการา อหุ ปญฺญาย อภิสมโย “เวทนาย โข สติ ตณฺหา โหติ, เวทนาปจฺจยา ตณฺหา”ติ.}

\prose{อถ โข ภิกฺขเว วิปสฺสิสฺส โพธิสตฺตสฺส เอตทโหสิ “กิมฺหิ นุ โข สติ เวทนา โหติ กิํ ปจฺจยา เวทนา”ติ.\csromanspacer{} อถ โข ภิกฺขเว วิปสฺสิสฺส โพธิสตฺตสฺส โยนิโส มนสิการา อหุ ปญฺญาย อภิสมโย “ผสฺเส โข สติ เวทนา โหติ, ผสฺสปจฺจยา เวทนา”ติ.}

\prose{อถ โข ภิกฺขเว วิปสฺสิสฺส โพธิสตฺตสฺส เอตทโหสิ “กิมฺหิ นุ โข สติ ผสฺโส โหติ กิํ ปจฺจยา ผสฺโส”ติ.\csromanspacer{} อถ โข ภิกฺขเว วิปสฺสิสฺส โพธิสตฺตสฺส โยนิโส มนสิการา อหุ ปญฺญาย อภิสมโย “สฬายตเน โข สติ ผสฺโส โหติ, สฬายตนปจฺจยา ผสฺโส”ติ.}

\prose{อถ โข ภิกฺขเว วิปสฺสิสฺส โพธิสตฺตสฺส เอตทโหสิ “กิมฺหิ นุ โข สติ สฬายตนํ โหติ กิํ ปจฺจยา สฬายตนนฺ”ติ.\csromanspacer{} อถ โข ภิกฺขเว วิปสฺสิสฺส โพธิสตฺตสฺส โยนิโส มนสิการา อหุ ปญฺญาย อภิสมโย “นามรูเป โข สติ สฬายตนํ โหติ, นามรูปปจฺจยา สฬายตนนฺ”ติ.}

\prose{อถ โข ภิกฺขเว วิปสฺสิสฺส โพธิสตฺตสฺส เอตทโหสิ “กิมฺหิ นุ โข สติ นามรูปํ โหติ กิํ ปจฺจยา นามรูปนฺ”ติ.\csromanspacer{} อถ โข ภิกฺขเว วิปสฺสิสฺส โพธิสตฺตสฺส โยนิโส มนสิการา อหุ ปญฺญาย อภิสมโย “วิญฺญาเณ โข สติ นามรูปํ โหติ, วิญฺญาณปจฺจยา นามรูปนฺ”ติ.}

\prose{อถ โข ภิกฺขเว วิปสฺสิสฺส โพธิสตฺตสฺส เอตทโหสิ “กิมฺหิ นุ โข สติ วิญฺญาณํ โหติ กิํ ปจฺจยา วิญฺญาณนฺ”ติ.\csromanspacer{} อถ โข ภิกฺขเว วิปสฺสิสฺส โพธิสตฺตสฺส โยนิโส มนสิการา อหุ ปญฺญาย อภิสมโย “นามรูเป โข สติ วิญฺญาณํ โหติ, นามรูปปจฺจยา วิญฺญาณนฺ”ติ.}

\proseitem{58}{อถ โข ภิกฺขเว วิปสฺสิสฺส โพธิสตฺตสฺส เอตทโหสิ, “ปจฺจุทาวตฺตติ โข อิทํ วิญฺญาณํ นามรูปมฺหา, นาปรํ คจฺฉติ.\csromanspacer{} เอตฺตาวตา \csromanfolio{28}ชาเยถ วา ชิยฺเยถ วา มิยฺเยถ วา จเวถ วา อุปปชฺเชถ วา, ยทิทํ นามรูปปจฺจยา วิญฺญาณํ, วิญฺญาณปจฺจยา นามรูปํ, นามรูปปจฺจยา สฬายตนํ, สฬายตนปจฺจยา ผสฺโส, ผสฺสปจฺจยา เวทนา, เวทนาปจฺจยา ตณฺหา, ตณฺหาปจฺจยา อุปาทานํ, อุปาทานปจฺจยา ภโว, ภวปจฺจยา ชาติ, ชาติปจฺจยา ชรามรณํ โสกปริเทวทุกฺขโทมนสฺสุปายาสา สมฺภวนฺติ.\csromanspacer{} เอวเมตสฺส เกวลสฺส ทุกฺขกฺขนฺธสฺส สมุทโย โหติ”.}

\proseitem{59}{“สมุทโย สมุทโย”ติ โข ภิกฺขเว วิปสฺสิสฺส โพธิสตฺตสฺส ปุพฺเพ อนนุสฺสุเตสุ ธมฺเมสุ จกฺขุํ อุทปาทิ, ญาณํ อุทปาทิ, ปญฺญา อุทปาทิ, วิชฺชา อุทปาทิ, อาโลโก อุทปาทิ.}

\proseitem{60}{อถ โข ภิกฺขเว วิปสฺสิสฺส โพธิสตฺตสฺส เอตทโหสิ “กิมฺหิ นุ โข อสติ ชรามรณํ น โหติ กิสฺส นิโรธา ชรามรณนิโรโธ”ติ.\csromanspacer{} อถ โข ภิกฺขเว วิปสฺสิสฺส โพธิสตฺตสฺส โยนิโส มนสิการา อหุ ปญฺญาย อภิสมโย “ชาติยา โข อสติ ชรามรณํ น โหติ, ชาตินิโรธา ชรามรณนิโรโธ”ติ.}

\prose{อถ โข ภิกฺขเว วิปสฺสิสฺส โพธิสตฺตสฺส เอตทโหสิ “กิมฺหิ นุ โข อสติ ชาติ น โหติ กิสฺส นิโรธา ชาตินิโรโธ”ติ.\csromanspacer{} อถ โข ภิกฺขเว วิปสฺสิสฺส โพธิสตฺตสฺส โยนิโส มนสิการา อหุ ปญฺญาย อภิสมโย “ภเว โข อสติ ชาติ น โหติ, ภวนิโรธา ชาตินิโรโธ”ติ.}

\prose{อถ โข ภิกฺขเว วิปสฺสิสฺส โพธิสตฺตสฺส เอตทโหสิ “กิมฺหิ นุ โข อสติ ภโว น โหติ กิสฺส นิโรธา ภวนิโรโธ”ติ.\csromanspacer{} อถ โข ภิกฺขเว วิปสฺสิสฺส โพธิสตฺตสฺส โยนิโส มนสิการา อหุ ปญฺญาย อภิสมโย “อุปาทาเน โข อสติ ภโว น โหติ, อุปาทานนิโรธา ภวนิโรโธ”ติ.}

\prose{อถ โข ภิกฺขเว วิปสฺสิสฺส โพธิสตฺตสฺส เอตทโหสิ “กิมฺหิ นุ โข อสติ อุปาทานํ น โหติ กิสฺส นิโรธา อุปาทานนิโรโธ”ติ, อถ โข ภิกฺขเว วิปสฺสิสฺส โพธิสตฺตสฺส โยนิโส มนสิการา \csromanfolio{29}อหุ ปญฺญาย อภิสมโย “ตณฺหาย โข อสติ อุปาทานํ น โหติ, ตณฺหานิโรธา อุปาทานนิโรโธ”ติ.}

\prose{อถ โข ภิกฺขเว วิปสฺสิสฺส โพธิสตฺตสฺส เอตทโหสิ “กิมฺหิ นุ โข อสติ ตณฺหา น โหติ กิสฺส นิโรธา ตณฺหานิโรโธ”ติ.\csromanspacer{} อถ โข ภิกฺขเว วิปสฺสิสฺส โพธิสตฺตสฺส โยนิโส มนสิการา อหุ ปญฺญาย อภิสมโย “เวทนาย โข อสติ ตณฺหา น โหติ เวทนานิโรธา ตณฺหานิโรโธ”ติ.}

\prose{อถ โข ภิกฺขเว วิปสฺสิสฺส โพธิสตฺตสฺส เอตทโหสิ “กิมฺหิ นุ โข อสติ เวทนา น โหติ กิสฺส นิโรธา เวทนานิโรโธ”ติ.\csromanspacer{} อถ โข ภิกฺขเว วิปสฺสิสฺส โพธิสตฺตสฺส โยนิโส มนสิการา อหุ ปญฺญาย อภิสมโย “ผสฺเส โข อสติ เวทนา น โหติ ผสฺสนิโรธา เวทนานิโรโธ”ติ.}

\prose{อถ โข ภิกฺขเว วิปสฺสิสฺส โพธิสตฺตสฺส เอตทโหสิ “กิมฺหิ นุ โข อสติ ผสฺโส น โหติ กิสฺส นิโรธา ผสฺสนิโรโธ”ติ.\csromanspacer{} อถ โข ภิกฺขเว วิปสฺสิสฺส โพธิสตฺตสฺส โยนิโส มนสิการา อหุ ปญฺญาย อภิสมโย “สฬายตเน โข อสติ ผสฺโส น โหติ สฬายตนนิโรธา ผสฺสนิโรโธ”ติ.}

\prose{อถ โข ภิกฺขเว วิปสฺสิสฺส โพธิสตฺตสฺส เอตทโหสิ “กิมฺหิ นุ โข อสติ สฬายตนํ น โหติ กิสฺส นิโรธา สฬายตนนิโรโธ”ติ.\csromanspacer{} อถ โข ภิกฺขเว วิปสฺสิสฺส โพธิสตฺตสฺส โยนิโส มนสิการา อหุ ปญฺญาย อภิสมโย “นามรูเป โข อสติ สฬายตนํ น โหติ นามรูปนิโรธา สฬายตนนิโรโธ”ติ.}

\prose{อถ โข ภิกฺขเว วิปสฺสิสฺส โพธิสตฺตสฺส เอตทโหสิ “กิมฺหิ นุ โข อสติ นามรูปํ น โหติ กิสฺส นิโรธา นามรูปนิโรโธ”ติ.\csromanspacer{} อถ โข ภิกฺขเว วิปสฺสิสฺส โพธิสตฺตสฺส โยนิโส มนสิการา อหุ ปญฺญาย อภิสมโย “วิญฺญาเณ โข อสติ นามรูปํ น โหติ วิญฺญาณนิโรธา นามรูปนิโรโธ”ติ.}

\prose{อถ โข ภิกฺขเว วิปสฺสิสฺส โพธิสตฺตสฺส เอตทโหสิ “กิมฺหิ นุ โข อสติ วิญฺญาณํ น โหติ กิสฺส นิโรธา วิญฺญาณนิโรโธ”ติ.\csromanspacer{} อถ โข ภิกฺขเว วิปสฺสิสฺส โพธิสตฺตสฺส โยนิโส มนสิการา \csromanfolio{30}อหุ ปญฺญาย อภิสมโย “นามรูเป โข อสติ วิญฺญาณํ น โหติ นามรูปนิโรธา วิญฺญาณนิโรโธ”ติ.}

\proseitem{61}{อถ โข ภิกฺขเว วิปสฺสิสฺส โพธิสตฺตสฺส เอตทโหสิ, “อธิคโต โข มฺยายํ มคฺโค สมฺโพธาย, ยทิทํ นามรูปนิโรธา วิญฺญาณนิโรโธ, วิญฺญาณนิโรธา นามรูปนิโรโธ, นามรูปนิโรธา สฬายตนนิโรโธ, สฬายตนนิโรธา ผสฺสนิโรโธ, ผสฺสนิโรธา เวทนานิโรโธ, เวทนานิโรธา ตณฺหานิโรโธ, ตณฺหานิโรธา อุปาทานนิโรโธ, อุปาทานนิโรธา ภวนิโรโธ, ภวนิโรธา ชาตินิโรโธ, ชาตินิโรธา ชรามรณํ โสกปริเทวทุกฺขโทมนสฺสุปายาสา นิรุชฺฌนฺติ.\csromanspacer{} เอวเมตสฺส เกวลสฺส ทุกฺขกฺขนฺธสฺส นิโรโธ โหติ”.}

\proseitem{62}{“นิโรโธ นิโรโธ”ติ โข ภิกฺขเว วิปสฺสิสฺส โพธิสตฺตสฺส ปุพฺเพ อนนุสฺสุเตสุ ธมฺเมสุ จกฺขุํ อุทปาทิ, ญาณํ อุทปาทิ, ปญฺญา อุทปาทิ, วิชฺชา อุทปาทิ, อาโลโก อุทปาทิ.}

\proseitem{63}{อถ โข ภิกฺขเว วิปสฺสี โพธิสตฺโต อปเรน สมเยน ปญฺจสุ อุปาทานกฺขนฺเธสุ อุทยพฺพยานุปสฺสี วิหาสิ “อิติ รูปํ อิติ รูปสฺส สมุทโย อิติ รูปสฺส อตฺถงฺคโม, อิติ เวทนา อิติ เวทนาย สมุทโย อิติ เวทนาย อตฺถงฺคโม, อิติ สญฺญา อิติ สญฺญาย สมุทโย อิติ สญฺญาย อตฺถงฺคโม, อิติ สงฺขารา อิติ สงฺขารานํ สมุทโย อิติ สงฺขารานํ อตฺถงฺคโม, อิติ วิญฺญาณํ อิติ วิญฺญาณสฺส สมุทโย อิติ วิญฺญาณสฺส อตฺถงฺคโม”ติ, ตสฺส ปญฺจสุ อุปาทานกฺขนฺเธสุ อุทยพฺพยานุปสฺสิโน วิหรโต น จิรสฺเสว อนุปาทาย อาสเวหิ จิตฺตํ วิมุจฺจีติ.}

\csromanheader{2}{ทุติยภาณวาโร.}
\csromansectionrule
\csromanmatikaanchor{mtk.13}
\csromantocmarkref{section}{พฺรหฺมยาจนกถา}{mtk.13}
\bookmark[dest={mtk.13},level=1]{พฺรหฺมยาจนกถา}
\csromanheader{1}{\textbf{พฺรหฺมยาจนกถา}}
\proseitem{64}{อถ โข ภิกฺขเว วิปสฺสิสฺส ภควโต อรหโต สมฺมาสมฺพุทฺธสฺส เอตทโหสิ “ยํนูนาหํ ธมฺมํ เทเสยฺยนฺ”ติ.\csromanspacer{} อถ โข ภิกฺขเว \csromanfolio{31}วิปสฺสิสฺส ภควโต อรหโต สมฺมาสมฺพุทฺธสฺส เอตทโหสิ “อธิคโต โข มฺยายํ ธมฺโม คมฺภีโร ทุทฺทโส ทุรนุโพโธ สนฺโต ปณีโต อตกฺกาวจโร นิปุโณ ปณฺฑิตเวทนีโย, อาลยรามา โข ปนายํ ปชา อาลยรตา อาลยสมฺมุทิตา, อาลยรามาย โข ปน ปชาย อาลยรตาย อาลยสมฺมุทิตาย ทุทฺทสํ อิทํ ฐานํ ยทิทํ อิทปฺปจฺจยตาปฏิจฺจสมุปฺปาโท, อิทมฺปิ โข ฐานํ ทุทฺทสํ ยทิทํ สพฺพสงฺขารสมโถ สพฺพูปธิปฏินิสฺสคฺโค ตณฺหากฺขโย วิราโค นิโรโธ นิพฺพานํ.\csromanspacer{} อหญฺเจว โข ปน ธมฺมํ เทเสยฺยํ, ปเร จ เม น อาชาเนยฺยุํ, โส มมสฺส กิลมโถ, สา มมสฺส วิเหสา”ติ.}

\proseitem{65}{อปิสฺสุ ภิกฺขเว วิปสฺสิํ ภควนฺตํ อรหนฺตํ สมฺมาสมฺพุทฺธํ อิมา อนจฺฉริยา คาถาโย ปฏิภํสุ ปุพฺเพ อสฺสุตปุพฺพา\csromandash{}}

\csromangathameasure{“กิจฺเฉน เม อธิคตํ, หลํ ทานิ ปกาสิตุํ. \\ ราคโทสปเรเตหิ, นายํ ธมฺโม สุสมฺพุโธ. \\ ปฏิโสตคามิํ นิปุณํ, คมฺภีรํ ทุทฺทสํ อณุํ. \\ ราครตฺตา น ทกฺขนฺติ, ตโมขนฺเธน อาวุฏา”ติ.}
\csromangathagroup{\csromangathastanza{“กิจฺเฉน เม อธิคตํ, หลํ ทานิ ปกาสิตุํ. \\ ราคโทสปเรเตหิ, นายํ ธมฺโม สุสมฺพุโธ.}\csromangathastanza{ปฏิโสตคามิํ นิปุณํ, คมฺภีรํ ทุทฺทสํ อณุํ. \\ ราครตฺตา น ทกฺขนฺติ, ตโมขนฺเธน อาวุฏา”ติ.}}

\prose{อิติห ภิกฺขเว วิปสฺสิสฺส ภควโต อรหโต สมฺมาสมฺพุทฺธสฺส ปฏิสญฺจิกฺขโต อปฺโปสฺสุกฺกตาย จิตฺตํ นมิ, โน ธมฺมเทสนาย.}

\proseitem{66}{อถ โข ภิกฺขเว อญฺญตรสฺส มหาพฺรหฺมุโน วิปสฺสิสฺส ภควโต อรหโต สมฺมาสมฺพุทฺธสฺส เจตสา เจโตปริวิตกฺกมญฺญาย เอตทโหสิ “นสฺสติ วต โภ โลโก, วินสฺสติ วต โภ โลโก, ยตฺร หิ นาม วิปสฺสิสฺส ภควโต อรหโต สมฺมาสมฺพุทฺธสฺส อปฺโปสฺสุกฺกตาย จิตฺตํ นมติ\footnote{นมิ (สฺยา, ก), นมิสฺสติ (?)}, โน ธมฺมเทสนายา”ติ.\csromanspacer{} อถ โข โส ภิกฺขเว มหาพฺรหฺมา เสยฺยถาปิ นาม พลวา ปุริโส สมิญฺชิตํ วา พาหํ ปสาเรยฺย, ปสาริตํ วา พาหํ สมิญฺเชยฺย, เอวเมว พฺรหฺมโลเก อนฺตรหิโต วิปสฺสิสฺส ภควโต อรหโต สมฺมาสมฺพุทฺธสฺส ปุรโต ปาตุรโหสิ.\csromanspacer{} อถ โข โส ภิกฺขเว มหาพฺรหฺมา เอกํสํ อุตฺตราสงฺคํ กริตฺวา ทกฺขิณํ ชาณุมณฺฑลํ \csromanfolio{32}ปถวิยํ นิหนฺตฺวา\footnote{นิทหนฺโต (สฺยา)} เยน วิปสฺสี ภควา อรหํ สมฺมาสมฺพุทฺโธ เตนญฺชลิํ ปณาเมตฺวา วิปสฺสิํ ภควนฺตํ อรหนฺตํ สมฺมาสมฺพุทฺธํ เอตทโวจ “เทเสตุ ภนฺเต ภควา ธมฺมํ, เทเสตุ สุคโต ธมฺมํ, สนฺติ\footnote{สนฺตีธ (สฺยา)} สตฺตา อปฺปรชกฺขชาติกา, อสฺสวนตา ธมฺมสฺส ปริหายนฺติ, ภวิสฺสนฺติ ธมฺมสฺส อญฺญาตาโร”ติ.}

\proseitem{67}{เอวํ วุตฺเต\footnote{อถ โข (ก)} ภิกฺขเว วิปสฺสี ภควา อรหํ สมฺมาสมฺพุทฺโธ ตํ มหาพฺรหฺมานํ เอตทโวจ “มยฺหมฺปิ โข พฺรหฺเม เอตทโหสิ ‘ยํนูนาหํ ธมฺมํ เทเสยฺยนฺ’ติ.\csromanspacer{} ตสฺส มยฺหํ พฺรหฺเม เอตทโหสิ ‘อธิคโต โข มฺยายํ ธมฺโม คมฺภีโร ทุทฺทโส ทุรนุโพโธ สนฺโต ปณีโต อตกฺกาวจโร นิปุโณ ปณฺฑิตเวทนีโย, อาลยรามา โข ปนายํ ปชา อาลยรตา อาลยสมฺมุทิตา, อาลยรามาย โข ปน ปชาย อาลยรตาย อาลยสมฺมุทิตาย ทุทฺทสํ อิทํ ฐานํ ยทิทํ อิทปฺปจฺจยตาปฏิจฺจสมุปฺปาโท, อิทมฺปิ โข ฐานํ ทุทฺทสํ ยทิทํ สพฺพสงฺขารสมโถ สพฺพูปธิปฏินิสฺสคฺโค ตณฺหกฺขโย วิราโค นิโรโธ นิพฺพานํ.\csromanspacer{} อหญฺเจว โข ปน ธมฺมํ เทเสยฺยํ, ปเร จ เม น อาชาเนยฺยุํ, โส มมสฺส กิลมโถ, สา มมสฺส วิเหสา’ติ.\csromanspacer{} อปิสฺสุ มํ พฺรหฺเม อิมา อนจฺฉริยา คาถาโย ปฏิภํสุ ปุพฺเพ อสฺสุตปุพฺพา\csromandash{}}

\csromangathameasure{‘กิจฺเฉน เม อธิคตํ, หลํ ทานิ ปกาสิตุํ. \\ ราคโทสปเรเตหิ, นายํ ธมฺโม สุสมฺพุโธ. \\ ปฏิโสตคามิํ นิปุณํ, คมฺภีรํ ทุทฺทสํ อณุํ. \\ ราครตฺตา น ทกฺขนฺติ, ตโมขนฺเธน อาวุฏา’ติ.}
\csromangathagroup{\csromangathastanza{‘กิจฺเฉน เม อธิคตํ, หลํ ทานิ ปกาสิตุํ. \\ ราคโทสปเรเตหิ, นายํ ธมฺโม สุสมฺพุโธ.}\csromangathastanza{ปฏิโสตคามิํ นิปุณํ, คมฺภีรํ ทุทฺทสํ อณุํ. \\ ราครตฺตา น ทกฺขนฺติ, ตโมขนฺเธน อาวุฏา’ติ.}}

\prose{อิติห เม พฺรหฺเม ปฏิสญฺจิกฺขโต อปฺโปสฺสุกฺกตาย จิตฺตํ นมิ, โน ธมฺมเทสนายา”ติ.}

\proseitem{68}{ทุติยมฺปิ โข ภิกฺขเว โส มหาพฺรหฺมา ฯเปฯ.\csromanspacer{} ตติยมฺปิ โข ภิกฺขเว โส มหาพฺรหฺมา วิปสฺสิํ ภควนฺตํ อรหนฺตํ สมฺมาสมฺพุทฺธํ เอตทโวจ “เทเสตุ ภนฺเต ภควา ธมฺมํ, เทเสตุ สุคโต ธมฺมํ, สนฺติ สตฺตา อปฺปรชกฺขชาติกา, อสฺสวนตา ธมฺมสฺส ปริหายนฺติ, ภวิสฺสนฺติ ธมฺมสฺส อญฺญาตาโร”ติ.}

\proseitem{69}{\csromanfolio{33}อถ โข ภิกฺขเว วิปสฺสี ภควา อรหํ สมฺมาสมฺพุทฺโธ พฺรหฺมุโน จ อชฺเฌสนํ วิทิตฺวา สตฺเตสุ จ การุญฺญตํ ปฏิจฺจ พุทฺธจกฺขุนา โลกํ โวโลเกสิ, อทฺทสา โข ภิกฺขเว วิปสฺสี ภควา อรหํ สมฺมาสมฺพุทฺโธ พุทฺธจกฺขุนา โลกํ โวโลเกนฺโต สตฺเต อปฺปรชกฺเข มหารชกฺเข ติกฺขินฺทฺริเย มุทินฺทฺริเย สฺวากาเร ทฺวากาเร สุวิญฺญาปเย ทุวิญฺญาปเย\footnote{ทุวิญฺญาปเย ภพฺเพ อภพฺเพ (สฺยา)} อปฺเปกจฺเจ ปรโลกวชฺชภยทสฺสาวิเน\footnote{ทสฺสาวิโน (สี, สฺยา, กํ, ก)} วิหรนฺเต อปฺเปกจฺเจ น ปรโลกวชฺชภยทสฺสาวิเน2 วิหรนฺเต.\csromanspacer{} เสยฺยถาปิ นาม อุปฺปลินิยํ วา ปทุมินิยํ วา ปุณฺฑรีกินิยํ วา อปฺเปกจฺจานิ อุปฺปลานิ วา ปทุมานิ วา ปุณฺฑรีกานิ วา อุทเก ชาตานิ อุทเก สํวฑฺฒานิ อุทกานุคฺคตานิ อนฺโต นิมุคฺคโปสีนิ.\csromanspacer{} อปฺเปกจฺจานิ อุปฺปลานิ วา ปทุมานิ วา ปุณฺฑรีกานิ วา อุทเก ชาตานิ อุทเก สํวฑฺฒานิ สโมทกํ ฐิตานิ.\csromanspacer{} อปฺเปกจฺจานิ อุปฺปลานิ วา ปทุมานิ วา ปุณฺฑรีกานิ วา อุทเก ชาตานิ อุทเก สํวฑฺฒานิ อุทกา อจฺจุคฺคมฺม ฐิตานิ อนุปลิตฺตานิ อุทเกน.\csromanspacer{} เอวเมว โข ภิกฺขเว วิปสฺสี ภควา อรหํ สมฺมาสมฺพุทฺโธ พุทฺธจกฺขุนา โลกํ โวโลเกนฺโต อทฺทส สตฺเต อปฺปรชกฺเข มหารชกฺเข ติกฺขินฺทฺริเย มุทินฺทฺริเย สฺวากาเร ทฺวากาเร สุวิญฺญาปเย ทุวิญฺญาปเย อปฺเปกจฺเจ ปรโลกวชฺชภยทสฺสาวิเน วิหรนฺเต อปฺเปกจฺเจ น ปรโลกวชฺชภยทสฺสาวิเน วิหรนฺเต.}

\proseitem{70}{อถ โข โส ภิกฺขเว มหาพฺรหฺมา วิปสฺสิสฺส ภควโต อรหโต สมฺมาสมฺพุทฺธสฺส เจตสา เจโตปริวิตกฺกมญฺญาย วิปสฺสิํ ภควนฺตํ อรหนฺตํ สมฺมาสมฺพุทฺธํ คาถาหิ อชฺฌภาสิ\csromandash{}}

\nitthitam{“เสเล ยถา ปพฺพตมุทฺธนิฏฺฐิโต, ยถาปิ ปสฺเส ชนตํ สมนฺตโต.}
\prose{ตถูปมํ ธมฺมมยํ สุเมธ, ปาสาทมารุยฺห สมนฺตจกฺขุ.}

\prose{โสกาวติณฺณํ\footnote{โสกาวกิณฺณํ (สฺยา)} ชนตมเปตโสโก,}

\prose{อเวกฺขสฺสุ ชาติชราภิภูตํ.}

\prose{อุฏฺเฐหิ วีร วิชิตสงฺคาม,}

\prose{สตฺถวาห อณณ วิจร โลเก.}

\prose{เทสสฺสุ\footnote{เทเสตุ (สฺยา, อิ)} ภควา ธมฺมํ,}

\prose{อญฺญาตาโร ภวิสฺสนฺตี”ติ.}

\proseitem{71}{\csromanfolio{34}อถ โข ภิกฺขเว วิปสฺสี ภควา อรหํ สมฺมาสมฺพุทฺโธ ตํ มหาพฺรหฺมานํ คาถาย อชฺฌภาสิ\csromandash{}}

\csromangathameasure{“อปารุตา เตสํ อมตสฺส ทฺวารา, \\ เย โสตวนฺโต ปมุญฺจนฺตุ สทฺธํ. \\ วิหิํสสญฺญี ปคุณํ น ภาสิํ, \\ ธมฺมํ ปณีตํ มนุเชสุ พฺรหฺเม”ติ.}
\csromangathagroup{\csromangathastanza{“อปารุตา เตสํ อมตสฺส ทฺวารา, \\ เย โสตวนฺโต ปมุญฺจนฺตุ สทฺธํ. \\ วิหิํสสญฺญี ปคุณํ น ภาสิํ, \\ ธมฺมํ ปณีตํ มนุเชสุ พฺรหฺเม”ติ.}}

\prose{อถ โข โส ภิกฺขเว มหาพฺรหฺมา “กตาวกาโส โขมฺหิ วิปสฺสินา ภควตา อรหตา สมฺมาสมฺพุทฺเธน ธมฺมเทสนายา”ติ วิปสฺสิํ ภควนฺตํ อรหนฺตํ สมฺมาสมฺพุทฺธํ อภิวาเทตฺวา ปทกฺขิณํ กตฺวา ตตฺเถว อนฺตรธายิ.}

\csromanmatikaanchor{mtk.14}
\csromantocmarkref{section}{อคฺคสาวกยุค}{mtk.14}
\bookmark[dest={mtk.14},level=1]{อคฺคสาวกยุค}
\csromanheader{1}{\textbf{อคฺคสาวกยุค}}
\proseitem{72}{อถ โข ภิกฺขเว วิปสฺสิสฺส ภควโต อรหโต สมฺมาสมฺพุทฺธสฺส เอตทโหสิ “กสฺส นุ โข อหํ ปฐมํ ธมฺมํ เทเสยฺยํ, โก อิมํ ธมฺมํ ขิปฺปเมว อาชานิสฺสตี”ติ.\csromanspacer{} อถ โข ภิกฺขเว วิปสฺสิสฺส ภควโต อรหโต สมฺมาสมฺพุทฺธสฺส เอตทโหสิ “อยํ โข ขณฺโฑ จ ราชปุตฺโต ติสฺโส จ ปุโรหิตปุตฺโต พนฺธุมติยา ราชธานิยา ปฏิวสนฺติ ปณฺฑิตา วิยตฺตา เมธาวิโน ทีฆรตฺตํ อปฺปรชกฺขชาติกา.\csromanspacer{} ยํนูนาหํ ขณฺฑสฺส จ ราชปุตฺตสฺส ติสฺสสฺส จ ปุโรหิตปุตฺตสฺส ปฐมํ ธมฺมํ เทเสยฺยํ, เต อิมํ ธมฺมํ ขิปฺปเมว อาชานิสฺสนฺตี”ติ.}

\proseitem{73}{อถ โข ภิกฺขเว วิปสฺสี ภควา อรหํ สมฺมาสมฺพุทฺโธ เสยฺยถาปิ นาม พลวา ปุริโส สมิญฺชิตํ วา พาหํ ปสาเรยฺย, ปสาริตํ วา พาหํ สมิญฺเชยฺย, เอวเมว โพธิรุกฺขมูเล อนฺตรหิโต พนฺธุมติยา ราชธานิยา เขเม มิคทาเย ปาตุรโหสิ.\csromanspacer{} อถ โข ภิกฺขเว วิปสฺสี ภควา อรหํ สมฺมาสมฺพุทฺโธ ทายปาลํ\footnote{มิคทายปาลํ (สฺยา)} อามนฺเตสิ “เอหิ ตฺวํ สมฺม ทายปาล พนฺธุมติํ ราชธานิํ ปวิสิตฺวา ขณฺฑญฺจ ราชปุตฺตํ ติสฺสญฺจ ปุโรหิตปุตฺตํ เอวํ วเทหิ ‘วิปสฺสี ภนฺเต ภควา อรหํ สมฺมาสมฺพุทฺโธ พนฺธุมติํ ราชธานิํ อนุปฺปตฺโต เขเม มิคทาเย วิหรติ, \csromanfolio{35}โส ตุมฺหากํ ทสฺสนกาโม’ติ”. “เอวํ ภนฺเต”ติ โข ภิกฺขเว ทายปาโล วิปสฺสิสฺส ภควโต อรหโต สมฺมาสมฺพุทฺธสฺส ปฏิสฺสุตฺวา พนฺธุมติํ ราชธานิํ ปวิสิตฺวา ขณฺฑญฺจ ราชปุตฺตํ ติสฺสญฺจ ปุโรหิตปุตฺตํ เอตทโวจ “วิปสฺสี ภนฺเต ภควา อรหํ สมฺมาสมฺพุทฺโธ พนฺธุมติํ ราชธานิํ อนุปฺปตฺโต เขเม มิคทาเย วิหรติ, โส ตุมฺหากํ ทสฺสนกาโม”ติ.}

\proseitem{74}{อถ โข ภิกฺขเว ขณฺโฑ จ ราชปุตฺโต ติสฺโส จ ปุโรหิตปุตฺโต ภทฺทานิ ภทฺทานิ ยานานิ โยชาเปตฺวา ภทฺทํ ภทฺทํ ยานํ อภิรุหิตฺวา ภทฺเทหิ ภทฺเทหิ ยาเนหิ พนฺธุมติยา ราชธานิยา นิยฺยิํสุ.\csromanspacer{} เยน เขโม มิคทาโย เตน ปายิํสุ, ยาวติกา ยานสฺส ภูมิ, ยาเนน คนฺตฺวา ยานา ปจฺโจโรหิตฺวา ปตฺติกาว\footnote{ปทิกาว (สฺยา)} เยน วิปสฺสี ภควา อรหํ สมฺมาสมฺพุทฺโธ เตนุปสงฺกมิํสุ, อุปสงฺกมิตฺวา วิปสฺสิํ ภควนฺตํ อรหนฺตํ สมฺมาสมฺพุทฺธํ อภิวาเทตฺวา เอกมนฺตํ นิสีทิํสุ.}

\proseitem{75}{เตสํ วิปสฺสี ภควา อรหํ สมฺมาสมฺพุทฺโธ อนุปุพฺพิํ กถํ\footnote{อานุปุพฺพิกถํ (สี, อิ)} กเถสิ, เสยฺยถิทํ, ทานกถํ สีลกถํ สคฺคกถํ กามานํ อาทีนวํ โอการํ สํกิเลสํ เนกฺขมฺเม อานิสํสํ ปกาเสสิ.\csromanspacer{} ยทา เต ภควา อญฺญาสิ กลฺลจิตฺเต มุทุจิตฺเต วินีวรณจิตฺเต อุทคฺคจิตฺเต ปสนฺนจิตฺเต, อถ ยา พุทฺธานํ สามุกฺกํสิกา ธมฺมเทสนา, ตํ ปกาเสสิ ทุกฺขํ สมุทยํ นิโรธํ มคฺคํ.\csromanspacer{} เสยฺยถาปิ นาม สุทฺธํ วตฺถํ อปคตกาฬกํ สมฺมเทว รชนํ ปฏิคฺคเณฺหยฺย, เอวเมว ขณฺฑสฺส จ ราชปุตฺตสฺส ติสฺสสฺส จ ปุโรหิตปุตฺตสฺส ตสฺมิํเยว อาสเน วิรชํ วีตมลํ ธมฺมจกฺขุํ อุทปาทิ “ยํ กิญฺจิ สมุทยธมฺมํ, สพฺพํ ตํ นิโรธธมฺมนฺ”ติ.}

\proseitem{76}{เต ทิฏฺฐธมฺมา ปตฺตธมฺมา วิทิตธมฺมา ปริโยคาฬฺหธมฺมา ติณฺณวิจิกิจฺฉา วิคตกถํกถา เวสารชฺชปฺปตฺตา อปรปฺปจฺจยา สตฺถุสาสเน วิปสฺสิํ ภควนฺตํ อรหนฺตํ สมฺมาสมฺพุทฺธํ เอตทโวจุํ “อภิกฺกนฺตํ ภนฺเต, อภิกฺกนฺตํ ภนฺเต, เสยฺยถาปิ ภนฺเต นิกฺกุชฺชิตํ วา อุกฺกุชฺเชยฺย, ปฏิจฺฉนฺนํ วา วิวเรยฺย, มูฬฺหสฺส วา มคฺคํ อาจิกฺเขยฺย, อนฺธกาเร วา เตลปชฺโชตํ ธาเรยฺย ‘จกฺขุมนฺโต รูปานิ ทกฺขนฺตี’ติ.\csromanspacer{} เอวเมวํ ภควตา อเนกปริยาเยน \csromanfolio{36}ธมฺโม ปกาสิโต.\csromanspacer{} เอเต มยํ ภนฺเต ภควนฺตํ สรณํ คจฺฉาม ธมฺมญฺจ, ลเภยฺยาม มยํ ภนฺเต ภควโต สนฺติเก ปพฺพชฺชํ ลเภยฺยาม อุปสมฺปทนฺ”ติ.}

\proseitem{77}{อลตฺถุํ โข ภิกฺขเว ขณฺโฑ จ ราชปุตฺโต ติสฺโส จ ปุโรหิตปุตฺโต วิปสฺสิสฺส ภควโต อรหโต สมฺมาสมฺพุทฺธสฺส สนฺติเก ปพฺพชฺชํ อลตฺถุํ อุปสมฺปทํ.\csromanspacer{} เต วิปสฺสี ภควา อรหํ สมฺมาสมฺพุทฺโธ ธมฺมิยา กถาย สนฺทสฺเสสิ สมาทเปสิ สมุตฺเตเชสิ สมฺปหํเสสิ, สงฺขารานํ อาทีนวํ โอการํ สํกิเลสํ นิพฺพาเน\footnote{เนกฺขมฺเม (สฺยา)} อานิสํสํ ปกาเสสิ.\csromanspacer{} เตสํ วิปสฺสินา ภควตา อรหตา สมฺมาสมฺพุทฺเธน ธมฺมิยา กถาย สนฺทสฺสิยมานานํ สมาทปิยมานานํ สมุตฺเตชิยมานานํ สมฺปหํสิยมานานํ นจิรสฺเสว อนุปาทาย อาสเวหิ จิตฺตานิ วิมุจฺจิํสุ.}

\csromanmatikaanchor{mtk.15}
\csromantocmarkref{section}{มหาชนกายปพฺพชฺชา}{mtk.15}
\bookmark[dest={mtk.15},level=1]{มหาชนกายปพฺพชฺชา}
\csromanheader{1}{\textbf{มหาชนกายปพฺพชฺชา}}
\proseitem{78}{อสฺโสสิ โข ภิกฺขเว พนฺธุมติยา ราชธานิยา มหาชนกาโย จตุราสีติปาณสหสฺสานิ “วิปสฺสี กิร ภควา อรหํ สมฺมาสมฺพุทฺโธ พนฺธุมติํ ราชธานิํ อนุปฺปตฺโต เขเม มิคทาเย วิหรติ, ขณฺโฑ จ กิร ราชปุตฺโต ติสฺโส จ ปุโรหิตปุตฺโต วิปสฺสิสฺส ภควโต อรหโต สมฺมาสมฺพุทฺธสฺส สนฺติเก เกสมสฺสุํ โอหาเรตฺวา กาสายานิ วตฺถานิ อจฺฉาเทตฺวา อคารสฺมา อนคาริยํ ปพฺพชิตา”ติ.\csromanspacer{} สุตฺวาน เนสํ เอตทโหสิ “น หิ นูน โส โอรโก ธมฺมวินโย, น สา โอรกา ปพฺพชฺชา, ยตฺถ ขณฺโฑ จ ราชปุตฺโต ติสฺโส จ ปุโรหิตปุตฺโต เกสมสฺสุํ โอหาเรตฺวา กาสายานิ วตฺถานิ อจฺฉาเทตฺวา อคารสฺมา อนคาริยํ ปพฺพชิตา, ขณฺโฑ จ ราชปุตฺโต ติสฺโส จ ปุโรหิตปุตฺโต เกสมสฺสุํ โอหาเรตฺวา กาสายานิ วตฺถานิ อจฺฉาเทตฺวา อคารสฺมา อนคาริยํ ปพฺพชิสฺสนฺติ, กิมงฺคํ ปน มยนฺ”ติ.\csromanspacer{} อถ โข โส ภิกฺขเว มหาชนกาโย จตุราสีติปาณสหสฺสานิ พนฺธุมติยา ราชธานิยา นิกฺขมิตฺวา เยน เขโม มิคทาโย, เยน วิปสฺสี ภควา อรหํ สมฺมาสมฺพุทฺโธ เตนุปสงฺกมิํสุ, อุปสงฺกมิตฺวา วิปสฺสิํ ภควนฺตํ อรหนฺตํ สมฺมาสมฺพุทฺธํ อภิวาเทตฺวา เอกมนฺตํ นิสีทิํสุ.}

\proseitem{79}{\csromanfolio{37}เตสํ วิปสฺสี ภควา อรหํ สมฺมาสมฺพุทฺโธ อนุปุพฺพิํ กถํ กเถสิ.\csromanspacer{} เสยฺยถิทํ, ทานกถํ สีลกถํ สคฺคกถํ กามานํ อาทีนวํ โอการํ สํกิเลสํ เนกฺขมฺเม อานิสํสํ ปกาเสสิ.\csromanspacer{} ยทา เต ภควา อญฺญาสิ กลฺลจิตฺเต มุทุจิตฺเต วินีวรณจิตฺเต อุทคฺคจิตฺเต ปสนฺนจิตฺเต, อถ ยา พุทฺธานํ สามุกฺกํสิกา ธมฺมเทสนา, ตํ ปกาเสสิ ทุกฺขํ สมุทยํ นิโรธํ มคฺคํ.\csromanspacer{} เสยฺยถาปิ นาม สุทฺธํ วตฺถํ อปคตกาฬกํ สมฺมเทว รชนํ ปฏิคฺคเณฺหยฺย, เอวเมว เตสํ จตุราสีติปาณสหสฺสานํ ตสฺมิํเยว อาสเน วิรชํ วีตมลํ ธมฺมจกฺขุํ อุทปาทิ “ยํ กิญฺจิ สมุทยธมฺมํ, สพฺพํ ตํ นิโรธธมฺมนฺ”ติ.}

\proseitem{80}{เต ทิฏฺฐธมฺมา ปตฺตธมฺมา วิทิตธมฺมา ปริโยคาฬฺหธมฺมา ติณฺณวิจิกิจฺฉา วิคตกถํกถา เวสารชฺชปฺปตฺตา อปรปฺปจฺจยา สตฺถุสาสเน วิปสฺสิํ ภควนฺตํ อรหนฺตํ สมฺมาสมฺพุทฺธํ เอตทโวจุํ “อภิกฺกนฺตํ ภนฺเต, อภิกฺกนฺตํ ภนฺเต, เสยฺยถาปิ ภนฺเต นิกฺกุชฺชิตํ วา อุกฺกุชฺเชยฺย, ปฏิจฺฉนฺนํ วา วิวเรยฺย, มูฬฺหสฺส วา มคฺคํ อาจิกฺเขยฺย, อนฺธกาเร วา เตลปชฺโชตํ ธาเรยฺย “จกฺขุมนฺโต รูปานิ ทกฺขนฺตี”ติ.\csromanspacer{} เอวเมวํ ภควตา อเนกปริยาเยน ธมฺโม ปกาสิโต.\csromanspacer{} เอเต มยํ ภนฺเต ภควนฺตํ สรณํ คจฺฉาม ธมฺมญฺจ (ภิกฺขุสํฆญฺจ)\footnote{( ) นตฺถิ อฏฺฐกถายํ, ปาฬิยํ ปน สพฺพตฺถปิ ทิสฺสติ.}, ลเภยฺยาม มยํ ภนฺเต ภควโต สนฺติเก ปพฺพชฺชํ ลเภยฺยาม อุปสมฺปทนฺ”ติ.}

\proseitem{81}{อลตฺถุํ โข ภิกฺขเว ตานิ จตุราสีติปาณสหสฺสานิ วิปสฺสิสฺส ภควโต อรหโต สมฺมาสมฺพุทฺธสฺส สนฺติเก ปพฺพชฺชํ อลตฺถุํ อุปสมฺปทํ.\csromanspacer{} เต วิปสฺสี ภควา อรหํ สมฺมาสมฺพุทฺโธ ธมฺมิยา กถาย สนฺทสฺเสสิ สมาทเปสิ สมุตฺเตเชสิ สมฺปหํเสสิ, สงฺขารานํ อาทีนวํ โอการํ สํกิเลสํ นิพฺพาเน อานิสํสํ ปกาเสสิ.\csromanspacer{} เตสํ วิปสฺสินา ภควตา อรหตา สมฺมาสมฺพุทฺเธน ธมฺมิยา กถาย สนฺทสฺสิยมานานํ สมาทปิยมานานํ สมุตฺเตชิยมานานํ สมฺปหํสิยมานานํ นจิรสฺเสว อนุปาทาย อาสเวหิ จิตฺตานิ วิมุจฺจิํสุ.}

\csromanmatikaanchor{mtk.16}
\csromantocmarkref{section}{ปุริมปพฺพชิตานํ ธมฺมาภิสมย}{mtk.16}
\bookmark[dest={mtk.16},level=1]{ปุริมปพฺพชิตานํ ธมฺมาภิสมย}
\csromanheader{1}{\textbf{ปุริมปพฺพชิตานํ ธมฺมาภิสมย}}
\proseitem{82}{อสฺโสสุํ โข ภิกฺขเว ตานิ ปุริมานิ จตุราสีติปพฺพชิตสหสฺสานิ “วิปสฺสี กิร ภควา อรหํ สมฺมาสมฺพุทฺโธ พนฺธุมติํ \csromanfolio{38}ราชธานิํ อนุปฺปตฺโต เขเม มิคทาเย วิหรติ, ธมฺมญฺจ กิร เทเสตี”ติ.\csromanspacer{} อถ โข ภิกฺขเว ตานิ จตุราสีติปพฺพชิตสหสฺสานิ เยน พนฺธุมตี ราชธานี, เยน เขโม มิคทาโย, เยน วิปสฺสี ภควา อรหํ สมฺมาสมฺพุทฺโธ เตนุปสงฺกมิํสุ, อุปสงฺกมิตฺวา วิปสฺสิํ ภควนฺตํ อรหนฺตํ สมฺมาสมฺพุทฺธํ อภิวาเทตฺวา เอกมนฺตํ นิสีทิํสุ.}

\proseitem{83}{เตสํ วิปสฺสี ภควา อรหํ สมฺมาสมฺพุทฺโธ อนุปุพฺพิํ กถํ กเถสิ.\csromanspacer{} เสยฺยถิทํ, ทานกถํ สีลกถํ สคฺคกถํ กามานํ อาทีนวํ โอการํ สํกิเลสํ เนกฺขมฺเม อานิสํสํ ปกาเสสิ.\csromanspacer{} ยทา เต ภควา อญฺญาสิ กลฺลจิตฺเต มุทุจิตฺเต วินีวรณจิตฺเต อุทคฺคจิตฺเต ปสนฺนจิตฺเต, อถ ยา พุทฺธานํ สามุกฺกํสิกา ธมฺมเทสนา, ตํ ปกาเสสิ ทุกฺขํ สมุทยํ นิโรธํ มคฺคํ.\csromanspacer{} เสยฺยถาปิ นาม สุทฺธํ วตฺถํ อปคตกาฬกํ สมฺมเทว รชนํ ปฏิคฺคเณฺหยฺย, เอวเมว เตสํ จตุราสีติปพฺพชิตสหสฺสานํ ตสฺมิํเยว อาสเน วิรชํ วีตมลํ ธมฺมจกฺขุํ อุทปาทิ “ยํ กิญฺจิ สมุทยธมฺมํ, สพฺพํ ตํ นิโรธธมฺมนฺ”ติ.}

\proseitem{84}{เต ทิฏฺฐธมฺมา ปตฺตธมฺมา วิทิตธมฺมา ปริโยคาฬฺหธมฺมา ติณฺณวิจิกิจฺฉา วิคตกถํกถา เวสารชฺชปฺปตฺตา อปรปฺปจฺจยา สตฺถุสาสเน วิปสฺสิํ ภควนฺตํ อรหนฺตํ สมฺมาสมฺพุทฺธํ เอตทโวจุํ “อภิกฺกนฺตํ ภนฺเต, อภิกฺกนฺตํ ภนฺเต, เสยฺยถาปิ ภนฺเต นิกฺกุชฺชิตํ วา อุกฺกุชฺเชยฺย, ปฏิจฺฉนฺนํ วา วิวเรยฺย, มูฬฺหสฺส วา มคฺคํ อาจิกฺเขยฺย, อนฺธกาเร วา เตลปชฺโชตํ ธาเรยฺย “จกฺขุมนฺโต รูปานิ ทกฺขนฺตี”ติ.\csromanspacer{} เอวเมวํ ภควตา อเนกปริยาเยน ธมฺโม ปกาสิโต.\csromanspacer{} เอเต มยํ ภนฺเต ภควนฺตํ สรณํ คจฺฉาม ธมฺมญฺจ ภิกฺขุสํฆญฺจ, ลเภยฺยาม มยํ ภนฺเต ภควโต สนฺติเก ปพฺพชฺชํ ลเภยฺยาม อุปสมฺปทนฺ”ติ.}

\proseitem{85}{อลตฺถุํ โข ภิกฺขเว ตานิ จตุราสีติปพฺพชิตสหสฺสานิ วิปสฺสิสฺส ภควโต อรหโต สมฺมาสมฺพุทฺธสฺส สนฺติเก ปพฺพชฺชํ อลตฺถุํ อุปสมฺปทํ.\csromanspacer{} เต วิปสฺสี ภควา อรหํ สมฺมาสมฺพุทฺโธ ธมฺมิยา กถาย สนฺทสฺเสสิ สมาทเปสิ สมุตฺเตเชสิ สมฺปหํเสสิ, สงฺขารานํ อาทีนวํ โอการํ สํกิเลสํ นิพฺพาเน อานิสํสํ ปกาเสสิ.\csromanspacer{} เตสํ วิปสฺสินา ภควตา อรหตา สมฺมาสมฺพุทฺเธน ธมฺมิยา กถาย สนฺทสฺสิยมานานํ \csromanfolio{39}สมาทปิยมานานํ สมุตฺเตชิยมานานํ สมฺปหํสิยมานานํ นจิรสฺเสว อนุปาทาย อาสเวหิ จิตฺตานิ วิมุจฺจิํสุ.}

\csromanmatikaanchor{mtk.17}
\csromantocmarkref{section}{จาริกา-อนุชานน}{mtk.17}
\bookmark[dest={mtk.17},level=1]{จาริกา-อนุชานน}
\csromanheader{1}{\textbf{จาริกา-อนุชานน}}
\proseitem{86}{เตน โข ปน ภิกฺขเว สมเยน พนฺธุมติยา ราชธานิยา มหาภิกฺขุสํโฆ ปฏิวสติ อฏฺฐสฏฺฐิภิกฺขุสตสหสฺสํ.\csromanspacer{} อถ โข ภิกฺขเว วิปสฺสิสฺส ภควโต อรหโต สมฺมาสมฺพุทฺธสฺส รโหคตสฺส ปฏิสลฺลีนสฺส เอวํ เจตโส ปริวิตกฺโก อุทปาทิ “มหา โข เอตรหิ ภิกฺขุสํโฆ พนฺธุมติยา ราชธานิยา ปฏิวสติ อฏฺฐสฏฺฐิภิกฺขุสตสหสฺสํ, ยํนูนาหํ ภิกฺขู อนุชาเนยฺยํ ‘จรถ ภิกฺขเว จาริกํ พหุชนหิตาย พหุชนสุขาย โลกานุกมฺปาย อตฺถาย หิตาย สุขาย เทวมนุสฺสานํ, มา เอเกน เทฺว อคมิตฺถ, เทเสถ ภิกฺขเว ธมฺมํ อาทิกลฺยาณํ มชฺเฌกลฺยาณํ ปริโยสานกลฺยาณํ สาตฺถํ สพฺยญฺชนํ เกวลปริปุณฺณํ ปริสุทฺธํ พฺรหฺมจริยํ ปกาเสถ, สนฺติ สตฺตา อปฺปรชกฺขชาติกา, อสฺสวนตา ธมฺมสฺส ปริหายนฺติ, ภวิสฺสนฺติ ธมฺมสฺส อญฺญาตาโร.\csromanspacer{} อปิ จ ฉนฺนํ ฉนฺนํ วสฺสานํ อจฺจเยน พนฺธุมตี ราชธานี อุปสงฺกมิตพฺพา ปาติโมกฺขุทฺเทสายา’ติ”.}

\proseitem{87}{อถ โข ภิกฺขเว อญฺญตโร มหาพฺรหฺมา วิปสฺสิสฺส ภควโต อรหโต สมฺมาสมฺพุทฺธสฺส เจตสา เจโตปริวิตกฺกมญฺญาย เสยฺยถาปิ นาม พลวา ปุริโส สมิญฺชิตํ วา พาหํ ปสาเรยฺย, ปสาริตํ วา พาหํ สมิญฺเชยฺย, เอวเมว พฺรหฺมโลเก อนฺตรหิโต วิปสฺสิสฺส ภควโต อรหโต สมฺมาสมฺพุทฺธสฺส ปุรโต ปาตุรโหสิ.\csromanspacer{} อถ โข โส ภิกฺขเว มหาพฺรหฺมา เอกํสํ อุตฺตราสงฺคํ กริตฺวา เยน วิปสฺสี ภควา อรหํ สมฺมาสมฺพุทฺโธ, เตนญฺชลิํ ปณาเมตฺวา วิปสฺสิํ ภควนฺตํ อรหนฺตํ สมฺมาสมฺพุทฺธํ เอตทโวจ “เอวเมตํ ภควา, เอวเมตํ สุคต.\csromanspacer{} มหา โข ภนฺเต เอตรหิ ภิกฺขุสํโฆ พนฺธุมติยา ราชธานิยา ปฏิวสติ อฏฺฐสฏฺฐิภิกฺขุสตสหสฺสํ, อนุชานาตุ ภนฺเต ภควา ภิกฺขู ‘จรถ ภิกฺขเว จาริกํ พหุชนหิตาย พหุชนสุขาย โลกานุกมฺปาย อตฺถาย หิตาย สุขาย เทวมนุสฺสานํ, มา เอเกน เทฺว อคมิตฺถ, เทเสถ ภิกฺขเว ธมฺมํ \csromanfolio{40}อาทิกลฺยาณํ มชฺเฌกลฺยาณํ ปริโยสานกลฺยาณํ สาตฺถํ สพฺยญฺชนํ เกวลปริปุณฺณํ ปริสุทฺธํ พฺรหฺมจริยํ ปกาเสถ, สนฺติ สตฺตา อปฺปรชกฺขชาติกา, อสฺสวนตา ธมฺมสฺส ปริหายนฺติ, ภวิสฺสนฺติ ธมฺมสฺส อญฺญาตาโร’ติ\footnote{อญฺญาตาโร (สพฺพตฺถ)}.\csromanspacer{} อปิ จ ภนฺเต มยํ ตถา กริสฺสาม, ยถา ภิกฺขู ฉนฺนํ ฉนฺนํ วสฺสานํ อจฺจเยน พนฺธุมติํ ราชธานิํ อุปสงฺกมิสฺสนฺติ ปาติโมกฺขุทฺเทสายา”ติ.\csromanspacer{} อิทมโวจ ภิกฺขเว โส มหาพฺรหฺมา, อิทํ วตฺวา วิปสฺสิํ ภควนฺตํ อรหนฺตํ สมฺมาสมฺพุทฺธํ อภิวาเทตฺวา ปทกฺขิณํ กตฺวา ตตฺเถว อนฺตรธายิ.}

\proseitem{88}{อถ โข ภิกฺขเว วิปสฺสี ภควา อรหํ สมฺมาสมฺพุทฺโธ สายนฺหสมยํ ปฏิสลฺลานา วุฏฺฐิโต ภิกฺขู อามนฺเตสิ\csromandash{}อิธ มยฺหํ ภิกฺขเว รโหคตสฺส ปฏิสลฺลีนสฺส เอวํ เจตโส ปริวิตกฺโก อุทปาทิ “มหา โข เอตรหิ ภิกฺขุสํโฆ พนฺธุมติยา ราชธานิยา ปฏิวสติ อฏฺฐสฏฺฐิภิกฺขุสตสหสฺสํ.\csromanspacer{} ยํนูนาหํ ภิกฺขู อนุชาเนยฺยํ ‘จรถ ภิกฺขเว จาริกํ พหุชนหิตาย พหุชนสุขาย โลกานุกมฺปาย อตฺถาย หิตาย สุขาย เทวมนุสฺสานํ, มา เอเกน เทฺว อคมิตฺถ, เทเสถ ภิกฺขเว ธมฺมํ อาทิกลฺยาณํ มชฺเฌกลฺยาณํ ปริโยสานกลฺยาณํ สาตฺถํ สพฺยญฺชนํ เกวลปริปุณฺณํ ปริสุทฺธํ พฺรหฺมจริยํ ปกาเสถ, สนฺติ สตฺตา อปฺปรชกฺขชาติกา, อสฺสวนตา ธมฺมสฺส ปริหายนฺติ, ภวิสฺสนฺติ ธมฺมสฺส อญฺญาตาโร.\csromanspacer{} อปิ จ ฉนฺนํ ฉนฺนํ วสฺสานํ อจฺจเยน พนฺธุมตี ราชธานี อุปสงฺกมิตพฺพา ปาติโมกฺขุทฺเทสายา’ติ”.}

\prose{อถ โข ภิกฺขเว อญฺญตโร มหาพฺรหฺมา มม เจตสา เจโตปริวิตกฺกมญฺญาย เสยฺยถาปิ นาม พลวา ปุริโส สมิญฺชิตํ วา พาหํ ปสาเรยฺย, ปสาริตํ วา พาหํ สมิญฺเชยฺย, เอวเมว พฺรหฺมโลเก อนฺตรหิโต มม ปุรโต ปาตุรโหสิ.\csromanspacer{} อถ โข โส ภิกฺขเว มหาพฺรหฺมา เอกํสํ อุตฺตราสงฺคํ กริตฺวา เยนาหํ เตนญฺชลิํ ปณาเมตฺวา มํ เอตทโวจ “เอวเมตํ ภควา, เอวเมตํ สุคต.\csromanspacer{} มหา โข ภนฺเต เอตรหิ ภิกฺขุสํโฆ พนฺธุมติยา ราชธานิยา ปฏิวสติ อฏฺฐสฏฺฐิภิกฺขุสตสหสฺสํ, อนุชานาตุ ภนฺเต ภควา ภิกฺขู ‘จรถ ภิกฺขเว จาริกํ พหุชนหิตาย พหุชนสุขาย โลกานุกมฺปาย อตฺถาย \csromanfolio{41}หิตาย สุขาย เทวมนุสฺสานํ, มา เอเกน เทฺว อคมิตฺถ, เทเสถ ภิกฺขเว ธมฺมํ ฯเปฯ สนฺติ สตฺตา อปฺปรชกฺขชาติกา, อสฺสวนตา ธมฺมสฺส ปริหายนฺติ, ภวิสฺสนฺติ ธมฺมสฺส อญฺญาตาโร’ติ.\csromanspacer{} อปิ จ ภนฺเต มยํ ตถา กริสฺสาม, ยถา ภิกฺขู ฉนฺนํ ฉนฺนํ วสฺสานํ อจฺจเยน พนฺธุมติํ ราชธานิํ อุปสงฺกมิสฺสนฺติ ปาติโมกฺขุทฺเทสายา”ติ.\csromanspacer{} อิทมโวจ ภิกฺขเว โส มหาพฺรหฺมา, อิทํ วตฺวา มํ อภิวาเทตฺวา ปทกฺขิณํ กตฺวา ตตฺเถว อนฺตรธายิ.}

\prose{อนุชานามิ ภิกฺขเว, จรถ จาริกํ พหุชนหิตาย พหุชนสุขาย โลกานุกมฺปาย อตฺถาย หิตาย สุขาย เทวมนุสฺสานํ, มา เอเกน เทฺว อคมิตฺถ, เทเสถ ภิกฺขเว ธมฺมํ อาทิกลฺยาณํ มชฺเฌกลฺยาณํ ปริโยสานกลฺยาณํ สาตฺถํ สพฺยญฺชนํ เกวลปริปุณฺณํ ปริสุทฺธํ พฺรหฺมจริยํ ปกาเสถ, สนฺติ สตฺตา อปฺปรชกฺขชาติกา, อสฺสวนตา ธมฺมสฺส ปริหายนฺติ, ภวิสฺสนฺติ ธมฺมสฺส อญฺญาตาโร.\csromanspacer{} อปิ จ ภิกฺขเว ฉนฺนํ ฉนฺนํ วสฺสานํ อจฺจเยน พนฺธุมตี ราชธานี อุปสงฺกมิตพฺพา ปาติโมกฺขุทฺเทสายาติ.\csromanspacer{} อถ โข ภิกฺขเว ภิกฺขู เยภุยฺเยน เอกาเหเนว ชนปทจาริกํ ปกฺกมิํสุ.}

\proseitem{89}{เตน โข ปน สมเยน ชมฺพุทีเป จตุราสีติ อาวาสสหสฺสานิ โหนฺติ.\csromanspacer{} เอกมฺหิ หิ วสฺเส นิกฺขนฺเต เทวตา สทฺทมนุสฺสาเวสุํ “นิกฺขนฺตํ โข มาริสา เอกํ วสฺสํ, ปญฺจ ทานิ วสฺสานิ เสสานิ, ปญฺจนฺนํ วสฺสานํ อจฺจเยน พนฺธุมตี ราชธานี อุปสงฺกมิตพฺพา ปาติโมกฺขุทฺเทสายา”ติ.\csromanspacer{} ทฺวีสุ วสฺเสสุ นิกฺขนฺเตสุ.\csromanspacer{} ตีสุ วสฺเสสุ นิกฺขนฺเตสุ.\csromanspacer{} จตูสุ วสฺเสสุ นิกฺขนฺเตสุ.\csromanspacer{} ปญฺจสุ วสฺเสสุ นิกฺขนฺเตสุ เทวตา สทฺทมนุสฺสาเวสุํ “นิกฺขนฺตานิ โข มาริสา ปญฺจวสฺสานิ, เอกํ ทานิ วสฺสํ เสสํ, เอกสฺส วสฺสสฺส อจฺจเยน พนฺธุมตี ราชธานี อุปสงฺกมิตพฺพา ปาติโมกฺขุทฺเทสายา”ติ.\csromanspacer{} ฉสุ วสฺเสสุ นิกฺขนฺเตสุ เทวตา สทฺทมนุสฺสาเวสุํ “นิกฺขนฺตานิ โข มาริสา ฉพฺพสฺสานิ, สมโย ทานิ พนฺธุมติํ ราชธานิํ อุปสงฺกมิตุํ ปาติโมกฺขุทฺเทสายา”ติ.\csromanspacer{} อถ โข เต ภิกฺขเว ภิกฺขู อปฺเปกจฺเจ สเกน อิทฺธานุภาเวน อปฺเปกจฺเจ เทวตานํ อิทฺธานุภาเวน เอกาเหเนว พนฺธุมติํ ราชธานิํ อุปสงฺกมิํสุ ปาติโมกฺขุทฺเทสายาติ\footnote{ปาติโมกฺขุทฺเทสาย (?)}.}

\proseitem{90}{\csromanfolio{42}ตตฺร สุทํ ภิกฺขเว วิปสฺสี ภควา อรหํ สมฺมาสมฺพุทฺโธ ภิกฺขุสํเฆ เอวํ ปาติโมกฺขํ อุทฺทิสติ.}

\csromangathameasure{“ขนฺตี ปรมํ ตโป ติติกฺขา, \\ นิพฺพานํ ปรมํ วทนฺติ พุทฺธา. \\ น หิ ปพฺพชิโต ปรูปฆาตี, \\ น สมโณ โหติ ปรํ วิเหฐยนฺโต. \\ สพฺพปาปสฺส อกรณํ, กุสลสฺส อุปสมฺปทา. \\ สจิตฺตปริโยทปนํ, เอตํ พุทฺธานสาสนํ. \\ อนูปวาโท อนูปฆาโต, ปาติโมกฺเข จ สํวโร. \\ มตฺตญฺญุตา จ ภตฺตสฺมิํ, ปนฺตญฺจ สยนาสนํ. \\ อธิจิตฺเต จ อาโยโค, เอตํ พุทฺธานสาสนนฺ”ติ.}
\csromangathagroup{\csromangathastanza{“ขนฺตี ปรมํ ตโป ติติกฺขา, \\ นิพฺพานํ ปรมํ วทนฺติ พุทฺธา. \\ น หิ ปพฺพชิโต ปรูปฆาตี, \\ น สมโณ\footnote{สมโณ (สี, สฺยา, อิ)} โหติ ปรํ วิเหฐยนฺโต. \\ สพฺพปาปสฺส อกรณํ, กุสลสฺส อุปสมฺปทา. \\ สจิตฺตปริโยทปนํ, เอตํ พุทฺธานสาสนํ.}\csromangathastanza{อนูปวาโท อนูปฆาโต\footnote{อนุปวาโท อนุปฆาโต (อิ, ก)}, ปาติโมกฺเข จ สํวโร. \\ มตฺตญฺญุตา จ ภตฺตสฺมิํ, ปนฺตญฺจ สยนาสนํ. \\ อธิจิตฺเต จ อาโยโค, เอตํ พุทฺธานสาสนนฺ”ติ.}}

\csromanmatikaanchor{mtk.18}
\csromantocmarkref{section}{เทวตาโรจน}{mtk.18}
\bookmark[dest={mtk.18},level=1]{เทวตาโรจน}
\csromanheader{1}{\textbf{เทวตาโรจน}}
\proseitem{91}{เอกมิทาหํ ภิกฺขเว สมยํ อุกฺกฏฺฐายํ วิหรามิ สุภควเน สาลราชมูเล, ตสฺส มยฺหํ ภิกฺขเว รโหคตสฺส ปฏิสลฺลีนสฺส เอวํ เจตโส ปริวิตกฺโก อุทปาทิ “น โข โส สตฺตาวาโส สุลภรูโป, โย มยา อนาวุตฺถปุพฺโพ\footnote{อนชฺฌาวุฏฺฐปุพฺโพ (ก-สี, ก)} อิมินา ทีเฆน อทฺธุนา อญฺญตฺร สุทฺธาวาเสหิ เทเวหิ.\csromanspacer{} ยํนูนาหํ เยน สุทฺธาวาสา เทวา เตนุปสงฺกเมยฺยนฺ”ติ.\csromanspacer{} อถ ขฺวาหํ ภิกฺขเว เสยฺยถาปิ นาม พลวา ปุริโส สมิญฺชิตํ วา พาหํ ปสาเรยฺย, ปสาริตํ วา พาหํ สมิญฺเชยฺย, เอวเมว อุกฺกฏฺฐายํ สุภควเน สาลราชมูเล อนฺตรหิโต อวิเหสุ เทเวสุ ปาตุรโหสิํ.\csromanspacer{} ตสฺมิํ ภิกฺขเว เทวนิกาเย อเนกานิ เทวตาสหสฺสานิ อเนกานิ เทวตาสตสหสฺสานิ\footnote{อเนกานิ เทวตาสตานิ อเนกานิ เทวตาสหสฺสานิ (สฺยา)} เยนาหํ เตนุปสงฺกมิํสุ, อุปสงฺกมิตฺวา มํ อภิวาเทตฺวา เอกมนฺตํ อฏฺฐํสุ.\csromanspacer{} เอกมนฺตํ ฐิตา โข ภิกฺขเว ตา เทวตา มํ เอตทโวจุํ “อิโต โส มาริสา เอกนวุติกปฺเป ยํ วิปสฺสี ภควา อรหํ สมฺมาสมฺพุทฺโธ โลเก อุทปาทิ.\csromanspacer{} วิปสฺสี มาริสา ภควา อรหํ สมฺมาสมฺพุทฺโธ ขตฺติโย ชาติยา อโหสิ, ขตฺติยกุเล อุทปาทิ.\csromanspacer{} วิปสฺสี มาริสา ภควา อรหํ สมฺมาสมฺพุทฺโธ โกณฺฑญฺโญ โคตฺเตน \csromanfolio{43}อโหสิ.\csromanspacer{} วิปสฺสิสฺส มาริสา ภควโต อรหโต สมฺมาสมฺพุทฺธสฺส อสีติวสฺสสหสฺสานิ อายุปฺปมาณํ อโหสิ.\csromanspacer{} วิปสฺสี มาริสา ภควา อรหํ สมฺมาสมฺพุทฺโธ ปาฏลิยา มูเล อภิสมฺพุทฺโธ.\csromanspacer{} วิปสฺสิสฺส มาริสา ภควโต อรหโต สมฺมาสมฺพุทฺธสฺส ขณฺฑติสฺสํ นาม สาวกยุคํ อโหสิ อคฺคํ ภทฺทยุคํ.\csromanspacer{} วิปสฺสิสฺส มาริสา ภควโต อรหโต สมฺมาสมฺพุทฺธสฺส ตโย สาวกานํ สนฺนิปาตา อเหสุํ.\csromanspacer{} เอโก สาวกานํ สนฺนิปาโต อโหสิ อตฺฏฺฐสฏฺฐิภิกฺขุสตสหสฺสํ.\csromanspacer{} เอโก สาวกานํ สนฺนิปาโต อโหสิ ภิกฺขุสตสหสฺสํ.\csromanspacer{} เอโก สาวกานํ สนฺนิปาโต อโหสิ อสีติภิกฺขุสหสฺสานิ.\csromanspacer{} วิปสฺสิสฺส มาริสา ภควโต อรหโต สมฺมาสมฺพุทฺธสฺส อิเม ตโย สาวกานํ สนฺนิปาตา อเหสุํ สพฺเพสํเยว ขีณาสวานํ.\csromanspacer{} วิปสฺสิสฺส มาริสา ภควโต อรหโต สมฺมาสมฺพุทฺธสฺส อโสโก นาม ภิกฺขุ อุปฏฺฐาโก อโหสิ อคฺคุปฏฺฐาโก.\csromanspacer{} วิปสฺสิสฺส มาริสา ภควโต อรหโต สมฺมาสมฺพุทฺธสฺส พนฺธุมา นาม ราชา ปิตา อโหสิ.\csromanspacer{} พนฺธุมตี นาม เทวี มาตา อโหสิ ชเนตฺติ.\csromanspacer{} พนฺธุมสฺส รญฺโญ พนฺธุมตี นาม นครํ ราชธานี อโหสิ.\csromanspacer{} วิปสฺสิสฺส มาริสา ภควโต อรหโต สมฺมาสมฺพุทฺธสฺส เอวํ อภินิกฺขมนํ อโหสิ เอวํ ปพฺพชฺชา เอวํ ปธานํ เอวํ อภิสมฺโพธิ เอวํ ธมฺมจกฺกปฺปวตฺตนํ.\csromanspacer{} เต มยํ มาริสา วิปสฺสิมฺหิ ภควติ พฺรหฺมจริยํ จริตฺวา กาเมสุ กามจฺฉนฺทํ วิราเชตฺวา อิธูปปนฺนา”ติ ฯเปฯ.}

\prose{ตสฺมิํเยว โข ภิกฺขเว เทวนิกาเย อเนกานิ เทวตาสหสฺสานิ อเนกานิ เทวตาสตสหสฺสานิ\footnote{อเนกานิ เทวตาสตานิ อเนกานิ เทวตาสหสฺสานิ (สฺยา, เอวมุปริปิ)} เยนาหํ เตนุปสงฺกมิํสุ, อุปสงฺกมิตฺวา มํ อภิวาเทตฺวา เอกมนฺตํ อฏฺฐํสุ.\csromanspacer{} เอกมนฺตํ ฐิตา โข ภิกฺขเว ตา เทวตา มํ เอตทโวจุํ “อิมสฺมิํเยว โข มาริสา ภทฺทกปฺเป ภควา เอตรหิ อรหํ สมฺมาสมฺพุทฺโธ โลเก อุปฺปนฺโน.\csromanspacer{} ภควา มาริสา ขตฺติโย ชาติยา ขตฺติยกุเล อุปฺปนฺโน.\csromanspacer{} ภควา มาริสา โคตโม โคตฺเตน.\csromanspacer{} ภควโต มาริสา อปฺปกํ อายุปฺปมาณํ ปริตฺตํ ลหุกํ โย จิรํ ชีวติ, โส วสฺสสตํ อปฺปํ วา ภิยฺโย.\csromanspacer{} ภควา มาริสา อสฺสตฺถสฺส มูเล อภิสมฺพุทฺโธ.\csromanspacer{} ภควโต มาริสา สาริปุตฺตโมคฺคลฺลานํ นาม สาวกยุคํ อโหสิ อคฺคํ \csromanfolio{44}ภทฺทยุคํ.\csromanspacer{} ภควโต มาริสา เอโก สาวกานํ สนฺนิปาโต อโหสิ อฑฺฒเตฬสานิ ภิกฺขุสตานิ.\csromanspacer{} ภควโต มาริสา อยํ เอโก สาวกานํ สนฺนิปาโต อโหสิ สพฺเพสํเยว ขีณาสวานํ.\csromanspacer{} ภควโต มาริสา อานนฺโท นาม ภิกฺขุ อุปฏฺฐาโก อโหสิ อคฺคุปฏฺฐาโก.\csromanspacer{} ภควโต มาริสา สุทฺโธทโน นาม ราชา ปิตา อโหสิ.\csromanspacer{} มายา นาม เทวี มาตา อโหสิ ชเนตฺติ.\csromanspacer{} กปิลวตฺถุ นาม นครํ ราชธานี อโหสิ.\csromanspacer{} ภควโต มาริสา เอวํ อภินิกฺขมนํ อโหสิ เอวํ ปพฺพชฺชา เอวํ ปธานํ เอวํ อภิสมฺโพธิ เอวํ ธมฺมจกฺกปฺปวตฺตนํ.\csromanspacer{} เต มยํ มาริสา ภควติ พฺรหฺมจริยํ จริตฺวา กาเมสุ กามจฺฉนฺทํ วิราเชตฺวา อิธูปปนฺนา”ติ.}

\proseitem{92}{อถ ขฺวาหํ ภิกฺขเว อวิเหหิ เทเวหิ สทฺธิํ เยน อตปฺปา เทวา เตนุปสงฺกมิํ ฯเปฯ.\csromanspacer{} อถ ขฺวาหํ ภิกฺขเว อวิเหหิ จ เทเวหิ อตปฺเปหิ จ เทเวหิ สทฺธิํ เยน สุทสฺสา เทวา เตนุปสงฺกมิํ.\csromanspacer{} อถ ขฺวาหํ ภิกฺขเว อวิเหหิ จ เทเวหิ อตปฺเปหิ จ เทเวหิ สุทสฺเสหิ จ เทเวหิ สทฺธิํ เยน สุทสฺสี เทวา เตนุปสงฺกมิํ.\csromanspacer{} อถ ขฺวาหํ ภิกฺขเว อวิเหหิ จ เทเวหิ อตปฺเปหิ จ เทเวหิ สุทสฺเสหิ จ เทเวหิ สุทสฺสีหิ จ เทเวหิ สทฺธิํ เยน อกนิฏฺฐา เทวา เตนุปสงฺกมิํ.\csromanspacer{} ตสฺมิํ ภิกฺขเว เทวนิกาเย อเนกานิ เทวตาสหสฺสานิ อเนกานิ เทวตาสตสหสฺสานิ เยนาหํ เตนุปสงฺกมิํสุ, อุปสงฺกมิตฺวา มํ อภิวาเทตฺวา เอกมนฺตํ อฏฺฐํสุ.}

\prose{เอกมนฺตํ ฐิตา โข ภิกฺขเว ตา เทวตา มํ เอตทโวจุํ “อิโต โส มาริสา เอกนวุติกปฺเป ยํ วิปสฺสี ภควา อรหํ สมฺมาสมฺพุทฺโธ โลเก อุทปาทิ.\csromanspacer{} วิปสฺสี มาริสา ภควา อรหํ สมฺมาสมฺพุทฺโธ ขตฺติโย ชาติยา อโหสิ, ขตฺติยกุเล อุทปาทิ.\csromanspacer{} วิปสฺสี มาริสา ภควา อรหํ สมฺมาสมฺพุทฺโธ โกณฺฑญฺโญ โคตฺเตน อโหสิ.\csromanspacer{} วิปสฺสิสฺส มาริสา ภควโต อรหโต สมฺมาสมฺพุทฺธสฺส อสีติวสฺสสหสฺสานิ อายุปฺปมาณํ อโหสิ.\csromanspacer{} วิปสฺสี มาริสา ภควา อรหํ สมฺมาสมฺพุทฺโธ ปาฏลิยา มูเล อภิสมฺพุทฺโธ.\csromanspacer{} วิปสฺสิสฺส มาริสา ภควโต อรหโต สมฺมาสมฺพุทฺธสฺส ขณฺฑติสฺสํ นาม สาวกยุคํ อโหสิ อคฺคํ ภทฺทยุคํ.\csromanspacer{} วิปสฺสิสฺส มาริสา ภควโต อรหโต สมฺมาสมฺพุทฺธสฺส ตโย สาวกานํ สนฺนิปาตา อเหสุํ.\csromanspacer{} เอโก สาวกานํ สนฺนิปาโต อโหสิ อฏฺฐสฏฺฐิภิกฺขุสตสหสฺสํ.\csromanspacer{} เอโก สาวกานํ สนฺนิปาโต \csromanfolio{45}อโหสิ ภิกฺขุสตสหสฺสํ.\csromanspacer{} เอโก สาวกานํ สนฺนิปาโต อโหสิ อสีติภิกฺขุสหสฺสานิ.\csromanspacer{} วิปสฺสิสฺส มาริสา ภควโต อรหโต สมฺมาสมฺพุทฺธสฺส อิเม ตโย สาวกานํ สนฺนิปาตา อเหสุํ สพฺเพสํเยว ขีณาสวานํ.\csromanspacer{} วิปสฺสิสฺส มาริสา ภควโต อรหโต สมฺมาสมฺพุทฺธสฺส อโสโก นาม ภิกฺขุ อุปฏฺฐาโก อโหสิ อคฺคุปฏฺฐาโก.\csromanspacer{} วิปสฺสิสฺส มาริสา ภควโต อรหโต สมฺมาสมฺพุทฺธสฺส พนฺธุมา นาม ราชา ปิตา อโหสิ.\csromanspacer{} พนฺธุมตี นาม เทวี มาตา อโหสิ ชเนตฺติ.\csromanspacer{} พนฺธุมสฺส รญฺโญ พนฺธุมตี นาม นครํ ราชธานี อโหสิ.\csromanspacer{} วิปสฺสิสฺส มาริสา ภควโต อรหโต สมฺมาสมฺพุทฺธสฺส เอวํ อภินิกฺขมนํ อโหสิ เอวํ ปพฺพชฺชา เอวํ ปธานํ เอวํ อภิสมฺโพธิ เอวํ ธมฺมจกฺกปฺปวตฺตนํ.\csromanspacer{} เต มยํ มาริสา วิปสฺสิมฺหิ ภควติ พฺรหฺมจริยํ จริตฺวา กาเมสุ กามจฺฉนฺทํ วิราเชตฺวา อิธูปปนฺนา”ติ.\csromanspacer{} ตสฺมิํเยว โข ภิกฺขเว เทวนิกาเย อเนกานิ เทวตาสหสฺสานิ อเนกานิ เทวตาสตสหสฺสานิ เยนาหํ เตนุปสงฺกมิํสุ, อุปสงฺกมิตฺวา มํ อภิวาเทตฺวา เอกมนฺตํ อฏฺฐํสุ.\csromanspacer{} เอกมนฺตํ ฐิตา โข ภิกฺขเว ตา เทวตา มํ เอตทโวจุํ “อิโต โส มาริสา เอกติํเส กปฺเป ยํ สิขี ภควา ฯเปฯ.\csromanspacer{} เต มยํ มาริสา สิขิมฺหิ ภควติ.\csromanspacer{} ตสฺมิํเยว โข มาริสา เอกติํเส กปฺเป ยํ เวสฺสภู ภควา ฯเปฯ.\csromanspacer{} เต มยํ มาริสา เวสฺสภุมฺหิ ภควติ ฯเปฯ.\csromanspacer{} อิมสฺมิํเยว โข มาริสา ภทฺทกปฺเป กกุสนฺโธ โกณาคมโน กสฺสโป ภควา ฯเปฯ.\csromanspacer{} เต มยํ มาริสา กกุสนฺธมฺหิ โกณาคมนมฺหิ กสฺสปมฺหิ ภควติ พฺรหฺมจริยํ จริตฺวา กาเมสุ กามจฺฉนฺทํ วิราเชตฺวา อิธูปปนฺนา”ติ.}

\proseitem{93}{ตสฺมิํเยว โข ภิกฺขเว เทวนิกาเย อเนกานิ เทวตาสหสฺสานิ อเนกานิ เทวตาสตสหสฺสานิ เยนาหํ เตนุปสงฺกมิํสุ, อุปสงฺกมิตฺวา มํ อภิวาเทตฺวา เอกมนฺตํ อฏฺฐํสุ.\csromanspacer{} เอกมนฺตํ ฐิตา โข ภิกฺขเว ตา เทวตา มํ เอตทโวจุํ “อิมสฺมิํเยว โข มาริสา ภทฺทกปฺเป ภควา เอตรหิ อรหํ สมฺมาสมฺพุทฺโธ โลเก อุปฺปนฺโน.\csromanspacer{} ภควา มาริสา ขตฺติโย ชาติยา, ขตฺติยกุเล อุปฺปนฺโน.\csromanspacer{} ภควา มาริสา โคตโม โคตฺเตน.\csromanspacer{} ภควโต มาริสา อปฺปกํ อายุปฺปมาณํ ปริตฺตํ ลหุกํ โย จิรํ ชีวติ, โส วสฺสสตํ อปฺปํ วา ภิยฺโย.\csromanspacer{} ภควา มาริสา อสฺสตฺถสฺส มูเล อภิสมฺพุทฺโธ.\csromanspacer{} ภควโต มาริสา สาริปุตฺตโมคฺคลฺลานํ นาม สาวกยุคํ อโหสิ อคฺคํ ภทฺทยุคํ. \csromanfolio{46}ภควโต มาริสา เอโก สาวกานํ สนฺนิปาโต อโหสิ อฑฺฒเตฬสานิ ภิกฺขุสตานิ.\csromanspacer{} ภควโต มาริสา อยํ เอโก สาวกานํ สนฺนิปาโต อโหสิ สพฺเพสํเยว ขีณาสวานํ.\csromanspacer{} ภควโต มาริสา อานนฺโท นาม ภิกฺขุ อุปฏฺฐาโก อคฺคุปฏฺฐาโก อโหสิ.\csromanspacer{} ภควโต มาริสา สุทฺโธทโน นาม ราชา ปิตา อโหสิ.\csromanspacer{} มายา นาม เทวี มาตา อโหสิ ชเนตฺติ.\csromanspacer{} กปิลวตฺถุ นาม นครํ ราชธานี อโหสิ.\csromanspacer{} ภควโต มาริสา เอวํ อภินิกฺขมนํ อโหสิ เอวํ ปพฺพชฺชา เอวํ ปธานํ เอวํ อภิสมฺโพธิ เอวํ ธมฺมจกฺกปฺปวตฺตนํ.\csromanspacer{} เต มยํ มาริสา ภควติ พฺรหฺมจริยํ จริตฺวา กาเมสุ กามจฺฉนฺทํ วิราเชตฺวา อิธูปปนฺนา”ติ.}

\proseitem{94}{อิติ โข ภิกฺขเว ตถาคตสฺเสเวสา ธมฺมธาตุ สุปฺปฏิวิทฺธา, ยสฺสา ธมฺมธาตุยา สุปฺปฏิวิทฺธตฺตา ตถาคโต อตีเต พุทฺเธ ปรินิพฺพุเต ฉินฺนปปญฺเจ ฉินวฏุเม ปริยาทินฺนวฏฺเฏ สพฺพทุกฺขวีติวตฺเต ชาติโตปิ อนุสฺสรติ, นามโตปิ อนุสฺสรติ, โคตฺตโตปิ อนุสฺสรติ, อายุปฺปมาณโตปิ อนุสฺสรติ, สาวกยุคโตปิ อนุสฺสรติ, สาวกสนฺนิปาตโตปิ อนุสฺสรติ “เอวํชจฺจา เต ภควนฺโต อเหสุํ" อิติปิ. “เอวํนามา เอวํโคตฺตา เอวํสีลา เอวํธมฺมา เอวํปญฺญา เอวํวิหารี เอวํวิมุตฺตา เต ภควนฺโต อเหสุํ” อิติปีติ.}

\prose{เทวตาปิ ตถาคตสฺส เอตมตฺถํ อาโรเจสุํ, เยน ตถาคโต อตีเต พุทฺเธ ปรินิพฺพุเต ฉินฺนปปญฺเจ ฉินฺนวฏุเม ปริยาทินฺนวฏฺเฏ สพฺพทุกฺขวีติวตฺเต ชาติโตปิ อนุสฺสรติ, นามโตปิ อนุสฺสรติ, โคตฺตโตปิ อนุสฺสรติ, อายุปฺปมาณโตปิ อนุสฺสรติ, สาวกยุคโตปิ อนุสฺสรติ, สาวกสนฺนิปาตโตปิ อนุสฺสรติ “เอวํ ชจฺจา เต ภควนฺโต อเหสุํ" อิติปิ. “เอวํนามา เอวํโคตฺตา เอวํสีลา เอวํธมฺมา เอวํปญฺญา เอวํวิหารี เอวํวิมุตฺตา เต ภควนฺโต อเหสุํ" อิติปีติ.}

\prose{อิทมโวจ ภควา.\csromanspacer{} อตฺตมนา เต ภิกฺขู ภควโต ภาสิตํ อภินนฺทุนฺติ.}

\nitthitam{\textbf{มหาปทานสุตฺตํ นิฏฺฐิตํ ปฐมํ.}}
\csromanmatikaanchor{mtk.19}
\csromantocmarknopage{chapter}{2. มหานิทานสุตฺต}{mtk.19}
\bookmark[dest={mtk.19},level=0]{2. มหานิทานสุตฺต}
\chapterheadpageread{2. มหานิทานสุตฺต}
\csromanmatikaanchor{mtk.20}
\csromantocmarkref{section}{ปฏิจฺจสมุปฺปาท}{mtk.20}
\bookmark[dest={mtk.20},level=1]{ปฏิจฺจสมุปฺปาท}
\csromanheader{1}{\textbf{ปฏิจฺจสมุปฺปาท}}
\proseitem{95}{\csromanfolio{47}เอวํ เม สุตํ\csromandash{}เอกํ สมยํ ภควา กุรูสุ วิหรติ กมฺมาสธมฺมํ นาม\footnote{กมฺมาสทมฺมํ นาม (สฺยา)} กุรูนํ นิคโม.\csromanspacer{} อถ โข อายสฺมา อานนฺโท เยน ภควา เตนุปสงฺกมิ, อุปสงฺกมิตฺวา ภควนฺตํ อภิวาเทตฺวา เอกมนฺตํ นิสีทิ.\csromanspacer{} เอกมนฺตํ นิสินฺโน โข อายสฺมา อานนฺโท ภควนฺตํ เอตทโวจ “อจฺฉริยํ ภนฺเต, อพฺภุตํ ภนฺเต, ยาว คมฺภีโร จายํ ภนฺเต ปฏิจฺจสมุปฺปาโท คมฺภีราวภาโส จ, อถ จ ปน เม อุตฺตานกุตฺตานโก วิย ขายตี”ติ.\csromanspacer{} มา เหวํ อานนฺท อวจ, มา เหวํ อานนฺท อวจ.\csromanspacer{} คมฺภีโร จายํ อานนฺท ปฏิจฺจสมุปฺปาโท คมฺภีราวภาโส จ, เอตสฺส อานนฺท ธมฺมสฺส อนนุโพธา อปฺปฏิเวธา เอวมยํ ปชา ตนฺตากุลกชาตา กุลคณฺฐิกชาตา\footnote{คุลาคุณฺฐิกชาตา (สี, อิ), คุณคณฺฐิกชาตา (สฺยา)} มุญฺชปพฺพชภูตา อปายํ ทุคฺคติํ วินิปาตํ สํสารํ นาติวตฺตติ.}

\proseitem{96}{อตฺถิ อิทปฺปจฺจยา ชรามรณนฺติ อิติ ปุฏฺเฐน สตา อานนฺท อตฺถีติสฺส วจนียํ.\csromanspacer{} กิํ ปจฺจยา ชรามรณนฺติ อิติ เจ วเทยฺย, ชาติปจฺจยา ชรามรณนฺติ อิจฺจสฺส วจนียํ.}

\prose{อตฺถิ อิทปฺปจฺจยา ชาตีติ อิติ ปุฏฺเฐน สตา อานนฺท อตฺถีติสฺส วจนียํ.\csromanspacer{} กิํ ปจฺจยา ชาตีติ อิติ เจ วเทยฺย, ภวปจฺจยา ชาตีติ อิจฺจสฺส วจนียํ.}

\prose{อตฺถิ อิทปฺปจฺจยา ภโวติ อิติ ปุฏฺเฐน สตา อานนฺท อตฺถีติสฺส วจนียํ.\csromanspacer{} กิํ ปจฺจยา ภโวติ อิติ เจ วเทยฺย, อุปาทานปจฺจยา ภโวติ อิจฺจสฺส วจนียํ.}

\prose{อตฺถิ อิทปฺปจฺจยา อุปาทานนฺติ อิติ ปุฏฺเฐน สตา อานนฺท อตฺถีติสฺส วจนียํ.\csromanspacer{} กิํ ปจฺจยา อุปาทานนฺติ อิติ เจ วเทยฺย, ตณฺหาปจฺจยา อุปาทานนฺติ อิจฺจสฺส วจนียํ.}

\prose{\csromanfolio{48}อตฺถิ อิทปฺปจฺจยา ตณฺหาติ อิติ ปุฏฺเฐน สตา อานนฺท อตฺถีติสฺส วจนียํ.\csromanspacer{} กิํ ปจฺจยา ตณฺหาติ อิติ เจ วเทยฺย, เวทนาปจฺจยา ตณฺหาติ อิจฺจสฺส วจนียํ.}

\prose{อตฺถิ อิทปฺปจฺจยา เวทนาติ อิติ ปุฏฺเฐน สตา อานนฺท อตฺถีติสฺส วจนียํ.\csromanspacer{} กิํ ปจฺจยา เวทนาติ อิติ เจ วเทยฺย, ผสฺสปจฺจยา เวทนาติ อิจฺจสฺส วจนียํ.}

\prose{อตฺถิ อิทปฺปจฺจยา ผสฺโสติ อิติ ปุฏฺเฐน สตา อานนฺท อตฺถีติสฺส วจนียํ.\csromanspacer{} กิํ ปจฺจยา ผสฺโสติ อิติ เจ วเทยฺย, นามรูปปจฺจยา ผสฺโสติ อิจฺจสฺส วจนียํ.}

\prose{อตฺถิ อิทปฺปจฺจยา นามรูปนฺติ อิติ ปุฏฺเฐน สตา อานนฺท อตฺถีติสฺส วจนียํ.\csromanspacer{} กิํ ปจฺจยา นามรูปนฺติ อิติ เจ วเทยฺย, วิญฺญาณปจฺจยา นามรูปนฺติ อิจฺจสฺส วจนียํ.}

\prose{อตฺถิ อิทปฺปจฺจยา วิญฺญาณนฺติ อิติ ปุฏฺเฐน สตา อานนฺท อตฺถีติสฺส วจนียํ.\csromanspacer{} กิํ ปจฺจยา วิญฺญาณนฺติ อิติ เจ วเทยฺย, นามรูปปจฺจยา วิญฺญาณนฺติ อิจฺจสฺส วจนียํ.}

\proseitem{97}{อิติ โข อานนฺท นามรูปปจฺจยา วิญฺญาณํ, วิญฺญาณปจฺจยา นามรูปํ, นามรูปปจฺจยา ผสฺโส, ผสฺสปจฺจยา เวทนา, เวทนาปจฺจยา ตณฺหา, ตณฺหาปจฺจยา อุปาทานํ, อุปาทานปจฺจยา ภโว, ภวปจฺจยา ชาติ, ชาติปจฺจยา ชรามรณํ โสกปริเทวทุกฺขโทมนสฺสุปายาสา สมฺภวนฺติ, เอวเมตสฺส เกวลสฺส ทุกฺขกฺขนฺธสฺส สมุทโย โหติ.}

\proseitem{98}{ชาติปจฺจยา ชรามรณนฺติ อิติ โข ปเนตํ วุตฺตํ, ตทานนฺท อิมินาเปตํ ปริยาเยน เวทิตพฺพํ, ยถา ชาติปจฺจยา ชรามรณํ.\csromanspacer{} ชาติ จ หิ อานนฺท นาภวิสฺส สพฺเพนสพฺพํ สพฺพถาสพฺพํ กสฺสจิ กิมฺหิจิ, เสยฺยถิทํ, เทวานํ วา เทวตฺตาย, คนฺธพฺพานํ วา คนฺธพฺพตฺตาย, ยกฺขานํ วา ยกฺขตฺตาย, ภูตานํ วา ภูตตฺตาย, มนุสฺสานํ วา มนุสฺสตฺตาย, จตุปฺปทานํ วา จตุปฺปทตฺตาย, ปกฺขีนํ วา ปกฺขิตฺตาย, สริสปานํ วา สรีสปตฺตาย\footnote{สิริํ สปานํ สิริํ สปตฺตาย (สี, สฺยา)}, เตสํ เตสํ จ หิ อานนฺท สตฺตานํ}

\vspace{\tipitakaparskip}
\csromancenter{มหานิทานสุตฺต ตทตฺตาย ชาติ นาภวิสฺส, สพฺพโส ชาติยา อสติ ชาตินิโรธา อปิ นุ โข ชรามรณํ ปญฺญาเยถาติ.\csromanspacer{} โน เหตํ ภนฺเต.\csromanspacer{} ตสฺมาติหานนฺท เอเสว เหตุ เอตํ นิทานํ เอส สมุทโย เอส ปจฺจโย ชรามรณสฺส, ยทิทํ ชาติ.}
\proseitem{99}{\csromanfolio{49}ภวปจฺจยา ชาตีติ อิติ โข ปเนตํ วุตฺตํ, ตทานนฺท อิมินาเปตํ ปริยาเยน เวทิตพฺพํ, ยถา ภวปจฺจยา ชาติ.\csromanspacer{} ภโว จ หิ อานนฺท นาภวิสฺส สพฺเพนสพฺพํ สพฺพถาสพฺพํ กสฺสจิ กิมฺหิจิ, เสยฺยถิทํ, กามภโว วา รูปภโว วา อรูปภโว วา.\csromanspacer{} สพฺพโส ภเว อสติ ภวนิโรธา อปิ นุ โข ชาติ ปญฺญาเยถาติ.\csromanspacer{} โน เหตํ ภนฺเต.\csromanspacer{} ตสฺมาติหานนฺท เอเสว เหตุ เอตํ นิทานํ เอส สมุทโย เอส ปจฺจโย ชาติยา, ยทิทํ ภโว.}

\proseitem{100}{อุปาทานปจฺจยา ภโวติ อิติ โข ปเนตํ วุตฺตํ, ตทานนฺท อิมินาเปตํ ปริยาเยน เวทิตพฺพํ, ยถา อุปาทานปจฺจยา ภโว.\csromanspacer{} อุปาทานํ จ หิ อานนฺท นาภวิสฺส สพฺเพนสพฺพํ สพฺพถาสพฺพํ กสฺสจิ กิมฺหิจิ, เสยฺยถิทํ, กามุปาทานํ วา ทิฏฺฐุปาทานํ วา สีลพฺพตุปาทานํ วา อตฺตวาทุปาทานํ วา.\csromanspacer{} สพฺพโส อุปาทาเน อสติ อุปาทานนิโรธา อปิ นุ โข ภโว ปญฺญาเยถาติ.\csromanspacer{} โน เหตํ ภนฺเต.\csromanspacer{} ตสฺมาติหานนฺท เอเสว เหตุ เอตํ นิทานํ เอส สมุทโย เอส ปจฺจโย ภวสฺส, ยทิทํ อุปาทานํ.}

\proseitem{101}{ตณฺหาปจฺจยา อุปาทานนฺติ อิติ โข ปเนตํ วุตฺตํ, ตทานนฺท อิมินาเปตํ ปริยาเยน เวทิตพฺพํ, ยถา ตณฺหาปจฺจยา อุปาทานํ.\csromanspacer{} ตณฺหา จ หิ อานนฺท นาภวิสฺส สพฺเพนสพฺพํ สพฺพถาสพฺพํ กสฺสจิ กิมฺหิจิ, เสยฺยถิทํ, รูปตณฺหา สทฺทตณฺหา คนฺธตณฺหา รสตณฺหา โผฏฺฐพฺพตณฺหา ธมฺมตณฺหา.\csromanspacer{} สพฺพโส ตณฺหาย อสติ ตณฺหานิโรธา อปิ นุ โข อุปาทานํ ปญฺญาเยถาติ.\csromanspacer{} โน เหตํ ภนฺเต.\csromanspacer{} ตสฺมาติหานนฺท เอเสว เหตุ เอตํ นิทานํ เอส สมุทโย เอส ปจฺจโย อุปาทานสฺส, ยทิทํ ตณฺหา.}

\proseitem{102}{เวทนาปจฺจยา ตณฺหาติ อิติ โข ปเนตํ วุตฺตํ, ตทานนฺท อิมินาเปตํ ปริยาเยน เวทิตพฺพํ, ยถา เวทนาปจฺจยา ตณฺหา. \csromanfolio{50}เวทนา จ หิ อานนฺท นาภวิสฺส สพฺเพนสพฺพํ สพฺพถาสพฺพํ กสฺสจิ กิมฺหิจิ, เสยฺยถิทํ, จกฺขุสมฺผสฺสชา เวทนา โสตสมฺผสฺสชา เวทนา ฆานสมฺผสฺสชา เวทนา ชิวฺหาสมฺผสฺสชา เวทนา กายสมฺผสฺสชา เวทนา มโนสมฺผสฺสชา เวทนา.\csromanspacer{} สพฺพโส เวทนาย อสติ เวทนานิโรธา อปิ นุ โข ตณฺหา ปญฺญาเยถาติ.\csromanspacer{} โน เหตํ ภนฺเต.\csromanspacer{} ตสฺมาติหานนฺท เอเสว เหตุ เอตํ นิทานํ เอส สมุทโย เอส ปจฺจโย ตณฺหาย, ยทิทํ เวทนา.}

\proseitem{103}{อิติ โข ปเนตํ อานนฺท เวทนํ ปฏิจฺจ ตณฺหา, ตณฺหํ ปฏิจฺจ ปริเยสนา, ปริเยสนํ ปฏิจฺจ ลาโภ, ลาภํ ปฏิจฺจ วินิจฺฉโย, วินิจฺฉยํ ปฏิจฺจ ฉนฺทราโค, ฉนฺทราคํ ปฏิจฺจ อชฺโฌสานํ, อชฺโฌสานํ ปฏิจฺจ ปริคฺคโห, ปริคฺคหํ ปฏิจฺจ มจฺฉริยํ, มจฺฉริยํ ปฏิจฺจ อารกฺโข.\csromanspacer{} อารกฺขาธิกรณํ ทณฺฑาทาน สตฺถาทาน กลห วิคฺคห วิวาท ตุวํตุวํ เปสุญฺญมุสาวาทา อเนเก ปาปกา อกุสลา ธมฺมา สมฺภวนฺติ.}

\proseitem{104}{อารกฺขาธิกรณํ\footnote{อารกฺขํ ปฏิจฺจ อารกฺขาธิกรณํ (สฺยา)} ทณฺฑาทาน สตฺถาทาน กลห วิคฺคห วิวาทตุวํตุวํ เปสุญฺญ มุสาวาทา อเนเก ปาปกา อกุสลา ธมฺมา สมฺภวนฺตีติ อิติ โข ปเนตํ วุตฺตํ, ตทานนฺท อิมินาเปตํ ปริยาเยน เวทิตพฺพํ, ยถา อารกฺขาธิกรณํ ทณฺฑาทาน สตฺถาทาน กลห วิคฺคห วิวาท ตุวํตุวํเปสุญฺญ มุสาวาทา อเนเก ปาปกา อกุสลา ธมฺมา สมฺภวนฺติ.\csromanspacer{} อารกฺโข จ หิ อานนฺท นาภวิสฺส สพฺเพนสพฺพํ สพฺพถาสพฺพํ กสฺสจิ กิมฺหิจิ, สพฺพโส อารกฺเข อสติ อารกฺขนิโรธา อปิ นุ โข ทณฺฑาทาน สตฺถาทาน กลห วิคฺคห วิวาท ตุวํ ตุวํ เปสุญฺญ มุสาวาทา อเนเก ปาปกา อกุสลา ธมฺมา สมฺภเวยฺยุนฺติ.\csromanspacer{} โน เหตํ ภนฺเต.\csromanspacer{} ตสฺมาติหานนฺท เอเสว เหตุ เอตํ นิทานํ เอส สมุทโย เอส ปจฺจโย ทณฺฑาทาน สตฺถาทาน กลห วิคฺคห วิวาท ตุวํตุวํ เปสุญฺญ มุสาวาทานํ อเนเกสํ ปาปกานํ อกุสลานํ ธมฺมานํ สมฺภวาย, ยทิทํ อารกฺโข.}

\proseitem{105}{มจฺฉริยํ ปฏิจฺจ อารกฺโขติ อิติ โข ปเนตํ วุตฺตํ, ตทานนฺท อิมินาเปตํ ปริยาเยน เวทิตพฺพํ, ยถา มจฺฉริยํ ปฏิจฺจ อารกฺโข.\csromanspacer{} มจฺฉริยํ จ หิ อานนฺท นาภวิสฺส สพฺเพนสพฺพํ สพฺพถาสพฺพํ กสฺสจิ}

\vspace{\tipitakaparskip}
\csromancenter{มหานิทานสุตฺต กิมฺหิจิ, สพฺพโส มจฺฉริเย อสติ มจฺฉริยนิโรธา อปิ นุ โข อารกฺโข ปญฺญาเยถาติ.\csromanspacer{} โน เหตํ ภนฺเต.\csromanspacer{} ตสฺมาติหานนฺท เอเสว เหตุ เอตํ นิทานํ เอส สมุทโย เอส ปจฺจโย อารกฺขสฺส, ยทิทํ มจฺฉริยํ.}
\proseitem{106}{\csromanfolio{51}ปริคฺคหํ ปฏิจฺจ มจฺฉริยนฺติ อิติ โข ปเนตํ วุตฺตํ, ตทานนฺท อิมินาเปตํ ปริยาเยน เวทิตพฺพํ, ยถา ปริคฺคหํ ปฏิจฺจ มจฺฉริยํ.\csromanspacer{} ปริคฺคโห จ หิ อานนฺท นาภวิสฺส สพฺเพนสพฺพํ สพฺพถาสพฺพํ กสฺสจิ กิมฺหิจิ, สพฺพโส ปริคฺคเห อสติ ปริคฺคหนิโรธา อปิ นุ โข มจฺฉริยํ ปญฺญาเยถาติ.\csromanspacer{} โน เหตํ ภนฺเต.\csromanspacer{} ตสฺมาติหานนฺท เอเสว เหตุ เอตํ นิทานํ เอส สมุทโย เอส ปจฺจโย มจฺฉริยสฺส, ยทิทํ ปริคฺคโห.}

\proseitem{107}{อชฺโฌสานํ ปฏิจฺจ ปริคฺคโหติ อิติ โข ปเนตํ วุตฺตํ, ตทานนฺท อิมินาเปตํ ปริยาเยน เวทิตพฺพํ, ยถา อชฺโฌสานํ ปฏิจฺจ ปริคฺคโห.\csromanspacer{} อชฺโฌสานํ จ หิ อานนฺท นาภวิสฺส สพฺเพนสพฺพํ สพฺพถาสพฺพํ กสฺสจิ กิมฺหิจิ, สพฺพโส อชฺโฌสาเน อสติ อชฺโฌสานนิโรธา อปิ นุ โข ปริคฺคโห ปญฺญาเยถาติ.\csromanspacer{} โน เหตํ ภนฺเต.\csromanspacer{} ตสฺมาติหานนฺท เอเสว เหตุ เอตํ นิทานํ เอส สมุทโย เอส ปจฺจโย ปริคฺคหสฺส, ยทิทํ อชฺโฌสานํ.}

\proseitem{108}{ฉนฺทราคํ ปฏิจฺจ อชฺโฌสานนฺติ อิติ โข ปเนตํ วุตฺตํ, ตทานนฺท อิมินาเปตํ ปริยาเยน เวทิตพฺพํ, ยถา ฉนฺทราคํ ปฏิจฺจ อชฺโฌสานํ.\csromanspacer{} ฉนฺทราโค จ หิ อานนฺท นาภวิสฺส สพฺเพนสพฺพํ สพฺพถาสพฺพํ กสฺสจิ กิมฺหิจิ, สพฺพโส ฉนฺทราเค อสติ ฉนฺทราคนิโรธา อปิ นุ โข อชฺโฌสานํ ปญฺญาเยถาติ.\csromanspacer{} โน เหตํ ภนฺเต.\csromanspacer{} ตสฺมาติหานนฺท เอเสว เหตุ เอตํ นิทานํ เอส สมุทโย เอส ปจฺจโย อชฺโฌสานสฺส, ยทิทํ ฉนฺทราโค.}

\proseitem{109}{วินิจฺฉยํ ปฏิจฺจ ฉนฺทราโคติ อิติ โข ปเนตํ วุตฺตํ, ตทานนฺท อิมินาเปตํ ปริยาเยน เวทิตพฺพํ, ยถา วินิจฺฉยํ ปฏิจฺจ ฉนฺทราโค.\csromanspacer{} วินิจฺฉโย จ หิ อานนฺท นาภวิสฺส สพฺเพนสพฺพํ สพฺพถาสพฺพํ กสฺสจิ กิมฺหิจิ, สพฺพโส วินิจฺฉเย อสติ วินิจฺฉยนิโรธา อปิ นุ โข ฉนฺทราโค ปญฺญาเยถาติ.\csromanspacer{} โน เหตํ ภนฺเต.\csromanspacer{} ตสฺมาติหานนฺท เอเสว เหตุ เอตํ นิทานํ เอส สมุทโย เอส ปจฺจโย ฉนฺทราคสฺส, ยทิทํ วินิจฺฉโย.}

\proseitem{110}{\csromanfolio{52}ลาภํ ปฏิจฺจ วินิจฺฉโยติ อิติ โข ปเนตํ วุตฺตํ, ตทานนฺท อิมินาเปตํ ปริยาเยน เวทิตพฺพํ, ยถา ลาภํ ปฏิจฺจ วินิจฺฉโย.\csromanspacer{} ลาโภ จ หิ อานนฺท นาภวิสฺส สพฺเพนสพฺพํ สพฺพถาสพฺพํ กสฺสจิ กิมฺหิจิ, สพฺพโส ลาเภ อสติ ลาภนิโรธา อปิ นุ โข วินิจฺฉโย ปญฺญาเยถาติ.\csromanspacer{} โน เหตํ ภนฺเต.\csromanspacer{} ตสฺมาติหานนฺท เอเสว เหตุ เอตํ นิทานํ เอส สมุทโย เอส ปจฺจโย วินิจฺฉยสฺส, ยทิทํ ลาโภ.}

\proseitem{111}{ปริเยสนํ ปฏิจฺจ ลาโภติ อิติ โข ปเนตํ วุตฺตํ, ตทานนฺท อิมินาเปตํ ปริยาเยน เวทิตพฺพํ, ยถา ปริเยสนํ ปฏิจฺจ ลาโภ.\csromanspacer{} ปริเยสนา จ หิ อานนฺท นาภวิสฺส สพฺเพนสพฺพํ สพฺพถาสพฺพํ กสฺสจิ กิมฺหิจิ, สพฺพโส ปริเยสนาย อสติ ปริเยสนานิโรธา อปิ นุ โข ลาโภ ปญฺญาเยถาติ.\csromanspacer{} โน เหตํ ภนฺเต.\csromanspacer{} ตสฺมาติหานนฺท เอเสว เหตุ เอตํ นิทานํ เอส สมุทโย เอส ปจฺจโย ลาภสฺส, ยทิทํ ปริเยสนา.}

\proseitem{112}{ตณฺหํ ปฏิจฺจ ปริเยสนาติ อิติ โข ปเนตํ วุตฺตํ, ตทานนฺท อิมินาเปตํ ปริยาเยน เวทิตพฺพํ, ยถา ตณฺหํ ปฏิจฺจ ปริเยสนา.\csromanspacer{} ตณฺหา จ หิ อานนฺท นาภวิสฺส สพฺเพนสพฺพํ สพฺพถาสพฺพํ กสฺสจิ กิมฺหิจิ, เสยฺยถิทํ, กามตณฺหา ภวตณฺหา วิภวตณฺหา.\csromanspacer{} สพฺพโส ตณฺหาย อสติ ตณฺหานิโรธา อปิ นุ โข ปริเยสนา ปญฺญาเยถาติ.\csromanspacer{} โน เหตํ ภนฺเต.\csromanspacer{} ตสฺมาติหานนฺท เอเสว เหตุ เอตํ นิทานํ เอส สมุทโย เอส ปจฺจโย ปริเยสนาย, ยทิทํ ตณฺหา.\csromanspacer{} อิติ โข อานนฺท อิเม เทฺว ธมฺมา\footnote{อิเม ธมฺมา (ก)} ทฺวเยน เวทนาย เอกสโมสรณา ภวนฺติ.}

\proseitem{113}{ผสฺสปจฺจยา เวทนาติ อิติ โข ปเนตํ วุตฺตํ, ตทานนฺท อิมินาเปตํ ปริยาเยน เวทิตพฺพํ, ยถา ผสฺสปจฺจยา เวทนา.\csromanspacer{} ผสฺโส จ หิ อานนฺท นาภวิสฺส สพฺเพนสพฺพํ สพฺพถาสพฺพํ กสฺสจิ กิมฺหิจิ, เสยฺยถิทํ, จกฺขุสมฺผสฺโส โสตสมฺผสฺโส ฆานสมฺผสฺโส ชิวฺหาสมฺผสฺโส กายสมฺผสฺโส มโนสมฺผสฺโส.\csromanspacer{} สพฺพโส ผสฺเส อสติ ผสฺสนิโรธา อปิ นุ โข เวทนา ปญฺญาเยถาติ.\csromanspacer{} โน เหตํ ภนฺเต.}

\vspace{\tipitakaparskip}
\csromancenter{มหานิทานสุตฺต ตสฺมาติหานนฺท เอเสว เหตุ เอตํ นิทานํ เอส สมุทโย เอส ปจฺจโย เวทนาย, ยทิทํ ผสฺโส.}
\proseitem{114}{\csromanfolio{53}นามรูปปจฺจยา ผสฺโสติ อิติ โข ปเนตํ วุตฺตํ, ตทานนฺท อิมินาเปตํ ปริยาเยน เวทิตพฺพํ, ยถา นามรูปปจฺจยา ผสฺโส.\csromanspacer{} เยหิ อานนฺท อากาเรหิ เยหิ ลิงฺเคหิ เยหิ นิมิตฺเตหิ เยหิ อุทฺเทเสหิ นามกายสฺส ปญฺญตฺติ โหติ, เตสุ อากาเรสุ เตสุ ลิงฺเคสุ เตสุ นิมิตฺเตสุ เตสุ อุทฺเทเสสุ อสติ อปิ นุ โข รูปกาเย อธิวจนสมฺผสฺโส ปญฺญาเยถาติ.\csromanspacer{} โน เหตํ ภนฺเต.\csromanspacer{} เยหิ อานนฺท อากาเรหิ เยหิ ลิงฺเคหิ เยหิ นิมิตฺเตหิ เยหิ อุทฺเทเสหิ รูปกายสฺส ปญฺญตฺติ โหติ, เตสุ อากาเรสุ ฯเปฯ เตสุ อุทฺเทเสสุ อสติ อปิ นุ โข นามกาเย ปฏิฆสมฺผสฺโส ปญฺญาเยถาติ.\csromanspacer{} โน เหตํ ภนฺเต.\csromanspacer{} เยหิ อานนฺท อากาเรหิ ฯเปฯ เยหิ อุทฺเทเสหิ นามกายสฺส จ รูปกายสฺส จ ปญฺญตฺติ โหติ, เตสุ อากาเรสุ ฯเปฯ เตสุ อุทฺเทเสสุ อสติ อปิ นุ โข อธิวจนสมฺผสฺโส วา ปฏิฆสมฺผสฺโส วา ปญฺญาเยถาติ.\csromanspacer{} โน เหตํ ภนฺเต.\csromanspacer{} เยหิ อานนฺท อากาเรหิ ฯเปฯ เยหิ อุทฺเทเสหิ นามรูปสฺส ปญฺญตฺติ โหติ, เตสุ อากาเรสุ ฯเปฯ เตสุ อุทฺเทเสสุ อสติ อปิ นุ โข ผสฺโส ปญฺญาเยถาติ.\csromanspacer{} โน เหตํ ภนฺเต.\csromanspacer{} ตสฺมาติหานนฺท เอเสว เหตุ เอตํ นิทานํ เอส สมุทโย เอส ปจฺจโย ผสฺสสฺส, ยทิทํ นามรูปํ.}

\proseitem{115}{วิญฺญาณปจฺจยา นามรูปนฺติ อิติ โข ปเนตํ วุตฺตํ, ตทานนฺท อิมินาเปตํ ปริยาเยน เวทิตพฺพํ, ยถา วิญฺญาณปจฺจยา นามรูปํ.\csromanspacer{} วิญฺญาณญฺจ หิ อานนฺท มาตุกุจฺฉิสฺมิํ น โอกฺกมิสฺสถ, อปิ นุ โข นามรูปํ มาตุกุจฺฉิสฺมิํ สมุจฺจิสฺสถาติ.\csromanspacer{} โน เหตํ ภนฺเต.\csromanspacer{} วิญฺญาณญฺจ หิ อานนฺท มาตุกุจฺฉิสฺมิํ โอกฺกมิตฺวา โวกฺกมิสฺสถ, อปิ นุ โข นามรูปํ อิตฺถตฺตาย อภินิพฺพตฺติสฺสถาติ.\csromanspacer{} โน เหตํ ภนฺเต.\csromanspacer{} วิญฺญาณญฺจ หิ อานนฺท ทหรสฺเสว สโต โวจฺฉิชฺชิสฺสถ กุมารกสฺส วา กุมาริกาย วา, อปิ นุ โข นามรูปํ วุทฺธิํ วิรูฬฺหิํ เวปุลฺลํ อาปชฺชิสฺสถาติ.\csromanspacer{} โน เหตํ ภนฺเต.\csromanspacer{} ตสฺมาติหานนฺท เอเสว เหตุ เอตํ นิทานํ เอส สมุทโย เอส ปจฺจโย นามรูปสฺส, ยทิทํ วิญฺญาณํ.}

\proseitem{116}{\csromanfolio{54}นามรูปปจฺจยา วิญฺญาณนฺติ อิติ โข ปเนตํ วุตฺตํ, ตทานนฺท อิมินาเปตํ ปริยาเยน เวทิตพฺพํ, ยถา นามรูปปจฺจยา วิญฺญาณํ.\csromanspacer{} วิญฺญาณญฺจ หิ อานนฺท นามรูเป ปติฏฺฐํ น ลภิสฺสถ, อปิ นุ โข อายติํ ชาติชรามรณํ ทุกฺขสมุทยสมฺภโว\footnote{ชาติชรามรณทุกฺขสมุทยสมฺภโว (สี, สฺยา, อิ)} ปญฺญาเยถาติ.\csromanspacer{} โน เหตํ ภนฺเต.\csromanspacer{} ตสฺมาติหานนฺท เอเสว เหตุ เอตํ นิทานํ เอส สมุทโย เอส ปจฺจโย วิญฺญาณสฺส ยทิทํ นามรูปํ, เอตฺตาวตา โข อานนฺท ชาเยถ วา ชีเยถ\footnote{ชิยฺเยถ (ก)} วา มีเยถ\footnote{มิยฺเยถ (ก)} วา จเวถ วา อุปปชฺเชถ วา.\csromanspacer{} เอตฺตาวตา อธิวจนปโถ, เอตฺตาวตา นิรุตฺติปโถ, เอตฺตาวตา ปญฺญตฺติปโถ, เอตฺตาวตา ปญฺญาวจรํ, เอตฺตาวตา วฏฺฏํ วตฺตติ อิตฺถตฺตํ ปญฺญาปนาย ยทิทํ นามรูปํ สห วิญฺญาเณน อญฺญมญฺญปจฺจยตา ปวตฺตติ.}

\csromanmatikaanchor{mtk.21}
\csromantocmarkref{section}{อตฺตปญฺญตฺติ}{mtk.21}
\bookmark[dest={mtk.21},level=1]{อตฺตปญฺญตฺติ}
\csromanheader{1}{\textbf{อตฺตปญฺญตฺติ}}
\proseitem{117}{กิตฺตาวตา จ อานนฺท อตฺตานํ ปญฺญเปนฺโต ปญฺญเปติ.\csromanspacer{} รูปิํ วา หิ อานนฺท ปริตฺตํ อตฺตานํ ปญฺญเปนฺโต ปญฺญเปติ “รูปี เม ปริตฺโต อตฺตา”ติ.\csromanspacer{} รูปิํ วา หิ อานนฺท อนนฺตํ อตฺตานํ ปญฺญเปนฺโต ปญฺญเปติ “รูปี เม อนนฺโต อตฺตา”ติ.\csromanspacer{} อรูปิํ วา หิ อานนฺท ปริตฺตํ อตฺตานํ ปญฺญเปนฺโต ปญฺญเปติ “อรูปี เม ปริตฺโต อตฺตา”ติ.\csromanspacer{} อรูปิํ วา หิ อานนฺท อนนฺตํ อตฺตานํ ปญฺญเปนฺโต ปญฺญเปติ “อรูปี เม อนนฺโต อตฺตา”ติ.}

\proseitem{118}{ตตฺรานนฺท โย โส รูปิํ ปริตฺตํ อตฺตานํ ปญฺญเปนฺโต ปญฺญเปติ, เอตรหิ วา โส รูปิํ ปริตฺตํ อตฺตานํ ปญฺญเปนฺโต ปญฺญเปติ, ตตฺถ ภาวิํ วา โส รูปิํ ปริตฺตํ อตฺตานํ ปญฺญเปนฺโต ปญฺญเปติ, “อตถํ วา ปน สนฺตํ ตถตฺตาย อุปกปฺเปสฺสามี”ติ อิติ วา ปนสฺส โหติ.\csromanspacer{} เอวํ สนฺตํ โข อานนฺท รูปิํ\footnote{รูปี (ก)} ปริตฺตตฺตานุทิฏฺฐิ อนุเสตีติ อิจฺจาลํ วจนาย.}

\prose{ตตฺรานนฺท โย โส รูปิํ อนนฺตํ อตฺตานํ ปญฺญเปนฺโต ปญฺญเปติ, เอตรหิ วา โส รูปิํ อนนฺตํ อตฺตานํ ปญฺญเปนฺโต ปญฺญเปติ, ตตฺถ ภาวิํ วา โส รูปิํ อนนฺตํ อตฺตานํ ปญฺญเปนฺโต ปญฺญเปติ, “อตถํ วา}

\vspace{\tipitakaparskip}
\csromancenter{มหานิทานสุตฺต ปน สนฺตํ ตถตฺตาย อุปกปฺเปสฺสามี”ติ อิติ วา ปนสฺส โหติ.\csromanspacer{} เอวํ สนฺตํ โข อานนฺท รูปิํ\footnote{รูปี (ก)} อนนฺตตฺตานุทิฏฺฐิ อนุเสตีติ อิจฺจาลํ วจนาย.}
\prose{\csromanfolio{55}ตตฺรานนฺท โย โส อรูปิํ ปริตฺตํ อตฺตานํ ปญฺญเปนฺโต ปญฺญเปติ, เอตรหิ วา โส อรูปิํ ปริตฺตํ อตฺตานํ ปญฺญาเปนฺโต ปญฺญเปติ, ตตฺถ ภาวิํ วา โส อรูปิํ ปริตฺตํ อตฺตานํ ปญฺญเปนฺโต ปญฺญเปติ, “อตถํ วา ปน สนฺตํ ตถตฺตาย อุปกปฺเปสฺสามี”ติ อิติ วา ปนสฺส โหติ.\csromanspacer{} เอวํ สนฺตํ โข อานนฺท อรูปิํ\footnote{อรูปี (ก)} ปริตฺตตฺตานุทิฏฺฐิ อนุเสตีติ อิจฺจาลํ วจนาย.}

\prose{ตตฺรานนฺท โย โส อรูปิํ อนนฺตํ อตฺตานํ ปญฺญเปนฺโต ปญฺญเปติ, เอตรหิ วา โส อรูปิํ อนนฺตํ อตฺตานํ ปญฺญเปนฺโต ปญฺญเปติ, ตตฺถ ภาวิํ วา โส อรูปิํ อนนฺตํ อตฺตานํ ปญฺญเปนฺโต ปญฺญเปติ, “อตถํ วา ปน สนฺตํ ตถตฺตาย อุปกปฺเปสฺสามี”ติ อิติ วา ปนสฺส โหติ.\csromanspacer{} เอวํ สนฺตํ โข อานนฺท อรูปิํ2 อนนฺตตฺตานุทิฏฺฐิ อนุเสตีติ อิจฺจาลํ วจนาย.\csromanspacer{} เอตฺตาวตา โข อานนฺท อตฺตานํ ปญฺญเปนฺโต ปญฺญเปติ.}

\csromanmatikaanchor{mtk.22}
\csromantocmarkref{section}{น-อตฺตปญฺญตฺติ}{mtk.22}
\bookmark[dest={mtk.22},level=1]{น-อตฺตปญฺญตฺติ}
\csromanheader{1}{\textbf{น อตฺตปญฺญตฺติ}}
\proseitem{119}{กิตฺตาวตา จ อานนฺท อตฺตานํ นปญฺญเปนฺโต น ปญฺญเปติ.\csromanspacer{} รูปิํ วา หิ อานนฺท ปริตฺตํ อตฺตานํ นปญฺญเปนฺโต น ปญฺญเปติ “รูปี เม ปริตฺโต อตฺตา”ติ.\csromanspacer{} รูปิํ วา หิ อานนฺท อนนฺตํ อตฺตานํ นปญฺญเปนฺโต น ปญฺญเปติ “รูปี เม อนนฺโต อตฺตา”ติ.\csromanspacer{} อรูปิํ วา หิ อานนฺท ปริตฺตํ อตฺตานํ นปญฺญเปนฺโต น ปญฺญเปติ “อรูปี เม ปริตฺโต อตฺตา”ติ.\csromanspacer{} อรูปิํ วา หิ อานนฺท อนนฺตํ อตฺตานํ นปญฺญเปนฺโต น ปญฺญเปติ “อรูปี เม อนนฺโต อตฺตา”ติ.}

\proseitem{120}{ตตฺรานนฺท โย โส รูปิํ ปริตฺตํ อตฺตานํ นปญฺญเปนฺโต น ปญฺญเปติ, เอตรหิ วา โส รูปิํ ปริตฺตํ อตฺตานํ นปญฺญเปนฺโต น ปญฺญเปติ, ตตฺถ ภาวิํ วา โส รูปิํ ปริตฺตํ อตฺตานํ นปญฺญเปนฺโต น ปญฺญเปติ, “อตถํ วา ปน สนฺตํ ตถตฺตาย อุปกปฺเปสฺสามี”ติ อิติ วา ปนสฺส น โหติ.\csromanspacer{} เอวํ สนฺตํ โข อานนฺท รูปิํ ปริตฺตตฺตานุทิฏฺฐิ นานุเสตี”ติ อิจฺจาลํ วจนาย.}

\prose{\csromanfolio{56}ตตฺรานนฺท โย โส รูปิํ อนนฺตํ อตฺตานํ นปญฺญเปนฺโต น ปญฺญเปติ, เอตรหิ วา โส รูปิํ อนนฺตํ อตฺตานํ นปญฺญเปนฺโต น ปญฺญเปติ, ตตฺถ ภาวิํ วา โส รูปิํ อนนฺตํ อตฺตานํ นปญฺญเปนฺโต น ปญฺญเปติ, “อตถํ วา ปน สนฺตํ ตถตฺตาย อุปกปฺเปสฺสามี”ติ อิติ วา ปนสฺส น โหติ.\csromanspacer{} เอวํ สนฺตํ โข อานนฺท รูปิํ อนนฺตตฺตานุทิฏฺฐิ นานุเสตี”ติ อิจฺจาลํ วจนาย.}

\prose{ตตฺรานนฺท โย โส อรูปิํ ปริตฺตํ อตฺตานํ นปญฺญเปนฺโต น ปญฺญเปติ.\csromanspacer{} เอตรหิ วา โส อรูปิํ ปริตฺตํ อตฺตานํ นปญฺญเปนฺโต น ปญฺญเปติ, ตตฺถ ภาวิํ วา โส อรูปิํ ปริตฺตํ อตฺตานํ นปญฺญเปนฺโต น ปญฺญเปติ, “อตถํ วา ปน สนฺตํ ตถตฺตาย อุปกปฺเปสฺสามี”ติ อิติ วา ปนสฺส น โหติ.\csromanspacer{} เอวํ สนฺตํ โข อานนฺท อรูปิํ ปริตฺตตฺตานุทิฏฺฐิ นานุเสตีติ อิจฺจาลํ วจนาย.}

\prose{ตตฺรานนฺท โย โส อรูปิํ อนนฺตํ อตฺตานํ นปญฺญเปนฺโต น ปญฺญเปติ, เอตรหิ วา โส อรูปิํ อนนฺตํ อตฺตานํ นปญฺญเปนฺโต น ปญฺญเปติ, ตตฺถ ภาวิํ วา โส อรูปิํ อนนฺตํ อตฺตานํ นปญฺญเปนฺโต น ปญฺญเปติ, “อตถํ วา ปน สนฺตํ ตถตฺตาย อุปกปฺเปสฺสามี”ติ อิติ วา ปนสฺส น โหติ.\csromanspacer{} เอวํ สนฺตํ โข อานนฺท อรูปิํ อนนฺตตฺตานุทิฏฺฐิ นานุเสตีติ อิจฺจาลํ วจนาย.\csromanspacer{} เอตฺตาวตา โข อานนฺท อตฺตานํ นปญฺญเปนฺโต น ปญฺญเปติ.}

\csromanmatikaanchor{mtk.23}
\csromantocmarkref{section}{อตฺตสมนุปสฺสนา}{mtk.23}
\bookmark[dest={mtk.23},level=1]{อตฺตสมนุปสฺสนา}
\csromanheader{1}{\textbf{อตฺตสมนุปสฺสนา}}
\proseitem{121}{กิตฺตาวตา จ อานนฺท อตฺตานํ สมนุปสฺสมาโน สมนุปสฺสติ.\csromanspacer{} เวทนํ วา หิ อานนฺท อตฺตานํ สมนุปสฺสมาโน สมนุปสฺสติ “เวทนา เม อตฺตา”ติ. “นเหว โข เม เวทนา อตฺตา อปฺปฏิสํเวทโน เม อตฺตา”ติ, อิติ วา หิ อานนฺท อตฺตานํ สมนุปสฺสมาโน สมนุปสฺสติ. “น เหว โข เม เวทนา อตฺตา, โนปิ อปฺปฏิสํเวทโน เม อตฺตา, อตฺตา เม เวทิยติ, เวทนาธมฺโม หิ เม อตฺตา”ติ, อิติ วา หิ อานนฺท อตฺตานํ สมนุปสฺสมาโน สมนุปสฺสติ.}

\proseitem{122}{ตตฺรานนฺท โย โส เอวมาห “เวทนา เม อตฺตา”ติ.\csromanspacer{} โส เอวมสฺส วจนีโย “ติสฺโส โข อิมา อาวุโส เวทนา สุขา เวทนา ทุกฺขา เวทนา อทุกฺขมสุขา เวทนา.\csromanspacer{} อิมาสํ โข ตฺวํ}

\vspace{\tipitakaparskip}
\csromancenter{มหานิทานสุตฺต ติสฺสนฺนํ เวทนานํ กตมํ อตฺตโต สมนุปสฺสสี”ติ.\csromanspacer{} ยสฺมิํ อานนฺท สมเย สุขํ เวทนํ เวเทติ, เนว ตสฺมิํ สมเย ทุกฺขํ เวทนํ เวเทติ, น อทุกฺขมสุขํ เวทนํ เวเทติ, สุขํเยว ตสฺมิํ สมเย เวทนํ เวเทติ.\csromanspacer{} ยสฺมิํ อานนฺท สมเย ทุกฺขํ เวทนํ เวเทติ, เนว ตสฺมิํ สมเย สุขํ เวทนํ เวเทติ, น อทุกฺขมสุขํ เวทนํ เวเทติ, ทุกฺขํเยว ตสฺมิํ สมเย เวทนํ เวเทติ.\csromanspacer{} ยสฺมิํ อานนฺท สมเย อทุกฺขมสุขํ เวทนํ เวเทติ, เนว ตสฺมิํ สมเย สุขํ เวทนํ เวเทติ, น ทุกฺขํ เวทนํ เวเทติ, อทุกฺขมสุขํเยว ตสฺมิํ สมเย เวทนํ เวเทติ.}
\proseitem{123}{\csromanfolio{57}สุขาปิ โข อานนฺท เวทนา อนิจฺจา สงฺขตา ปฏิจฺจสมุปฺปนฺนา ขยธมฺมา วยธมฺมา วิราคธมฺมา นิโรธธมฺมา.\csromanspacer{} ทุกฺขาปิ โข อานนฺท เวทนา อนิจฺจา สงฺขตา ปฏิจฺจสมุปฺปนฺนา ขยธมฺมา วยธมฺมา วิราคธมฺมา นิโรธธมฺมา.\csromanspacer{} อทุกฺขมสุขาปิ โข อานนฺท เวทนา อนิจฺจา สงฺขตา ปฏิจฺจสมุปฺปนฺนา ขยธมฺมา วยธมฺมา วิราคธมฺมา นิโรธธมฺมา.\csromanspacer{} ตสฺส สุขํ เวทนํ เวทิยมานสฺส “เอโส เม อตฺตา”ติ โหติ.\csromanspacer{} ตสฺสาเยว สุขาย เวทนาย นิโรธา “พฺยคา\footnote{พฺยคฺคา (สี, ก)} เม อตฺตา”ติ โหติ.\csromanspacer{} ทุกฺขํ เวทนํ เวทิยมานสฺส “เอโส เม อตฺตา”ติ โหติ.\csromanspacer{} ตสฺสาเยว ทุกฺขาย เวทนาย นิโรธา “พฺยคา เม อตฺตา”ติ โหติ.\csromanspacer{} อทุกฺขมสุขํ เวทนํ เวทิยมานสฺส “เอโส เม อตฺตา”ติ โหติ.\csromanspacer{} ตสฺสาเยว อทุกฺขมสุขาย เวทนาย นิโรธา “พฺยคา เม อตฺตา”ติ โหติ.\csromanspacer{} อิติ โส ทิฏฺเฐวธมฺเม อนิจฺจสุขทุกฺขโวกิณฺณํ อุปฺปาทวยธมฺมํ อตฺตานํ สมนุปสฺสมาโน สมนุปสฺสติ, โย โส เอวมาห “เวทนา เม อตฺตา”ติ.\csromanspacer{} ตสฺมาติหานนฺท เอเตนเปตํ นกฺขมติ “เวทนา เม อตฺตา”ติ สมนุปสฺสิตุํ.}

\proseitem{124}{ตตฺรานนฺท โย โส เอวมาห “น เหว โข เม เวทนา อตฺตา, อปฺปฏิสํเวทโน เม อตฺตา”ติ.\csromanspacer{} โส เอวมสฺส วจนีโย “ยตฺถ ปนาวุโส สพฺพโส เวทยิตํ นตฺถิ, อปิ นุ โข ตตฺถ ‘อยมหมสฺมี’ติ สิยา”ติ.\csromanspacer{} โน เหตํ ภนฺเต.\csromanspacer{} ตสฺมาติหานนฺท เอเตนเปตํ นกฺขมติ “น เหว โข เม เวทนา อตฺตา, อปฺปฏิสํเวทโน เม อตฺตา”ติ สมนุปสฺสิตุํ.}

\proseitem{125}{\csromanfolio{58}ตตฺรานนฺท โย โส เอวมาห “น เหว โข เม เวทนา อตฺตา, โนปิ อปฺปฏิสํเวทโน เม อตฺตา, อตฺตา เม เวทิยติ, เวทนาธมฺโม หิ เม อตฺตา”ติ.\csromanspacer{} โส เอวมสฺส วจนีโย “เวทนา จ หิ อาวุโส สพฺเพนสพฺพํ สพฺพถาสพฺพํ อปริเสสา นิรุชฺเฌยฺยุํ.\csromanspacer{} สพฺพโส เวทนาย อสติ เวทนานิโรธา อปิ นุ โข ตตฺถ ‘อยมหมสฺมี’ติ สิยา”ติ.\csromanspacer{} โน เหตํ ภนฺเต.\csromanspacer{} ตสฺมาติหานนฺท เอเตนเปตํ นกฺขมติ “น เหว โข เม เวทนา อตฺตา, โนปิ อปฺปฏิสํเวทโน เม อตฺตา, อตฺตา เม เวทิยติ, เวทนาธมฺโม หิ เม อตฺตา”ติ สมนุปสฺสิตุํ.}

\proseitem{126}{ยโต โข อานนฺท ภิกฺขุ เนว เวทนํ อตฺตานํ สมนุปสฺสติ, โนปิ อปฺปฏิสํเวทนํ อตฺตานํ สมนุปสฺสติ, โนปิ “อตฺตา เม เวทิยติ, เวทนาธมฺโม หิ เม อตฺตา”ติ สมนุปสฺสติ.\csromanspacer{} โส เอวํ นสมนุปสฺสนฺโต น จ กิญฺจิ โลเก อุปาทิยติ, อนุปาทิยํ น ปริตสฺสติ, อปริตสฺสํ\footnote{อปริตสฺสนํ (ก)} ปจฺจตฺตญฺเญว ปรินิพฺพายติ, “ขีณา ชาติ, วุสิตํ พฺรหฺมจริยํ, กตํ กรณียํ, นาปรํ อิตฺถตฺตายา”ติ ปชานาติ.\csromanspacer{} เอวํ วิมุตฺตจิตฺตํ โข อานนฺท ภิกฺขุํ โย เอวํ วเทยฺย “โหติ ตถาคโต ปรํ มรณา อิติสฺส\footnote{อิติ สา (อฏฺฐกถายํ ปาฐนฺตรํ)} ทิฏฺฐี”ติ, ตทกลฺลํ. “น โหติ ตถาคโต ปรํ มรณา อิติสฺส ทิฏฺฐี”ติ, ตทกลฺลํ. “โหติ จ น จ โหติ ตถาคโต ปรํ มรณา อิติสฺส ทิฏฺฐี”ติ, ตทกลฺลํ. “เนว โหติ น น โหติ ตถาคโต ปรํ มรณา อิติสฺส ทิฏฺฐี”ติ, ตทกลฺลํ.\csromanspacer{} ตํ กิสฺส เหตุ.\csromanspacer{} ยาวตา อานนฺท อธิวจนํ ยาวตา อธิวจนปโถ, ยาวตา นิรุตฺติ ยาวตา นิรุตฺติปโถ, ยาวตา ปญฺญตฺติ ยาวตา ปญฺญตฺติปโถ, ยาวตา ปญฺญา ยาวตา ปญฺญาวจรํ, ยาวตา วฏฺฏํ3 ยาวตา วฏฺฏติ3, ตทภิญฺญาวิมุตฺโต ภิกฺขุ, ตทภิญฺญาวิมุตฺตํ ภิกฺขุํ “น ชานาติ น ปสฺสติ อิติสฺส ทิฏฺฐี”ติ, ตทกลฺลํ.}

\csromanmatikaanchor{mtk.24}
\csromantocmarkref{section}{สตฺตวิญฺญาณฏฺฐิติ}{mtk.24}
\bookmark[dest={mtk.24},level=1]{สตฺตวิญฺญาณฏฺฐิติ}
\csromanheader{1}{\textbf{สตฺตวิญฺญาณฏฺฐิติ}}
\proseitem{127}{สตฺต โข อานนฺท\footnote{สตฺต โข อิมา อานนฺท (ก-สี, สฺยา)} วิญฺญาณฏฺฐิติโย.\csromanspacer{} เทฺว อายตนานิ.\csromanspacer{} กตมา สตฺต.\csromanspacer{} สนฺตานนฺท สตฺตา นานตฺตกายา นานตฺตสญฺญิโน, เสยฺยถาปิ}

\vspace{\tipitakaparskip}
\csromancenter{มหานิทานสุตฺต มนุสฺสา เอกจฺเจ จ เทวา เอกจฺเจ จ วินิปาติกา, อยํ ปฐมา วิญฺญาณฏฺฐิติ.\csromanspacer{} สนฺตานนฺท สตฺตา นานตฺตกายา เอกตฺตสญฺญิโน, เสยฺยถาปิ เทวา พฺรหฺมกายิกา ปฐมาภินิพฺพตฺตา, อยํ ทุติยา วิญฺญาณฏฺฐิติ.\csromanspacer{} สนฺตานนฺท สตฺตา เอกตฺตกายา นานตฺตสญฺญิโน, เสยฺยถาปิ เทวา อาภสฺสรา, อยํ ตติยา วิญฺญาณฏฺฐิติ.\csromanspacer{} สนฺตานนฺท สตฺตา เอกตฺตกายา เอกตฺตสญฺญิโน, เสยฺยถาปิ เทวา สุภกิณฺหา, อยํ จตุตฺถี วิญฺญาณฏฺฐิติ.\csromanspacer{} สนฺตานนฺท สตฺตา สพฺพโส รูปสญฺญานํ สมติกฺกมา ปฏิฆสญฺญานํ อตฺถงฺคมา นานตฺตสญฺญานํ อมนสิการา “อนนฺโต อากาโส”ติ อากาสานญฺจายตนูปคา, อยํ ปญฺจมี วิญฺญาณฏฺฐิติ.\csromanspacer{} สนฺตานนฺท สตฺตา สพฺพโส อากาสานญฺจายตนํ สมติกฺกมฺม “อนนฺตํ วิญฺญาณนฺ”ติ วิญฺญาณญฺจายตนูปคา, อยํ ฉฏฺฐี วิญฺญาณฏฺฐิติ.\csromanspacer{} สนฺตานนฺท สตฺตา สพฺพโส วิญฺญาณญฺจายตนํ สมติกฺกมฺม “นตฺถิ กิญฺจี”ติ อากิญฺจญฺญายตนูปคา, อยํ สตฺตมี วิญฺญาณฏฺฐิติ.\csromanspacer{} อสญฺญสตฺตายตนํ เนวสญฺญานาสญฺญายตนเมว ทุติยํ.}
\proseitem{128}{\csromanfolio{59}ตตฺรานนฺท ยายํ ปฐมา วิญฺญาณฏฺฐิติ นานตฺตกายา นานตฺตสญฺญิโน, เสยฺยถาปิ มนุสฺสา เอกจฺเจ จ เทวา เอกจฺเจ จ วินิปาติกา.\csromanspacer{} โย นุ โข อานนฺท ตญฺจ ปชานาติ, ตสฺสา จ สมุทยํ ปชานาติ, ตสฺสา จ อตฺถงฺคมํ ปชานาติ, ตสฺสา จ อสฺสาทํ ปชานาติ, ตสฺสา จ อาทีนวํ ปชานาติ, ตสฺสา จ นิสฺสรณํ ปชานาติ, กลฺลํ นุ เตน ตทภินนฺทิตุนฺติ.\csromanspacer{} โน เหตํ ภนฺเต ฯเปฯ.\csromanspacer{} ตตฺรานนฺท ยมิทํ อสญฺญสตฺตายตนํ.\csromanspacer{} โย นุ โข อานนฺท ตญฺจ ปชานาติ, ตสฺส จ สมุทยํ ปชานาติ, ตสฺส จ อตฺถงฺคมํ ปชานาติ, ตสฺส จ อสฺสาทํ ปชานาติ, ตสฺส จ อาทีนวํ ปชานาติ, ตสฺส จ นิสฺสรณํ ปชานาติ, กลฺลํ นุ เตน ตทภินนฺทิตุนฺติ.\csromanspacer{} โน เหตํ ภนฺเต.\csromanspacer{} ตตฺรานนฺท ยมิทํ เนวสญฺญานาสญฺญายตนํ.\csromanspacer{} โย นุ โข อานนฺท ตญฺจ ปชานาติ, ตสฺส จ สมุทยํ ปชานาติ, ตสฺส จ อตฺถงฺคมํ ปชานาติ, ตสฺส จ อสฺสาทํ ปชานาติ, ตสฺส จ อาทีนวํ ปชานาติ, ตสฺส จ นิสฺสรณํ ปชานาติ, กลฺลํ นุ เตน ตทภินนฺทิตุนฺติ.\csromanspacer{} โน เหตํ ภนฺเต.\csromanspacer{} ยโต โข อานนฺท ภิกฺขุ อิมาสญฺจ สตฺตนฺนํ วิญฺญาณฏฺฐิตีนํ อิเมสญฺจ ทฺวินฺนํ อายตนานํ สมุทยญฺจ อตฺถงฺคมญฺจ อสฺสาทญฺจ อาทีนวญฺจ นิสฺสรณญฺจ ยถาภูตํ วิทิตฺวา อนุปาทา วิมุตฺโต โหติ.\csromanspacer{} อยํ วุจฺจตานนฺท ภิกฺขุ ปญฺญาวิมุตฺโต.}

\csromanmatikaanchor{mtk.25}
\csromantocmarkref{section}{อฏฺฐวิโมกฺขา}{mtk.25}
\bookmark[dest={mtk.25},level=1]{อฏฺฐวิโมกฺขา}
\csromanheader{1}{\textbf{อฏฺฐวิโมกฺขา}}
\proseitem{129}{\csromanfolio{60}อฏฺฐ โข อิเม อานนฺท วิโมกฺขา.\csromanspacer{} กตเม อฏฺฐ.\csromanspacer{} รูปี รูปานิ ปสฺสติ, อยํ ปฐโม วิโมกฺโข.\csromanspacer{} อชฺฌตฺตํ อรูปสญฺญี พหิทฺธา รูปานิ ปสฺสติ, อยํ ทุติโย วิโมกฺโข.\csromanspacer{} สุภนฺเตว อธิมุตฺโต โหติ, อยํ ตติโย วิโมกฺโข.\csromanspacer{} สพฺพโส รูปสญฺญานํ สมติกฺกมา ปฏิฆสญฺญานํ อตฺถงฺคมา นานตฺตสญฺญานํ อมนสิการา “อนนฺโต อากาโส”ติ อากาสานญฺจายตนํ อุปสมฺปชฺช วิหรติ, อยํ จตุตฺโถ วิโมกฺโข.\csromanspacer{} สพฺพโส อากาสานญฺจายตนํ สมติกฺกมฺม “อนนฺตํ วิญฺญาณนฺ”ติ วิญฺญาณญฺจายตนํ อุปสมฺปชฺช วิหรติ, อยํ ปญฺจโม วิโมกฺโข.\csromanspacer{} สพฺพโส วิญฺญาณญฺจายตนํ สมติกฺกมฺม “นตฺถิ กิญฺจี”ติ อากิญฺจญฺญายตนํ อุปสมฺปชฺช วิหรติ, อยํ ฉฏฺโฐ วิโมกฺโข.\csromanspacer{} สพฺพโส อากิญฺจญฺญายตนํ สมติกฺกมฺม เนวสญฺญานาสญฺญายตนํ อุปสมฺปชฺช วิหรติ, อยํ สตฺตโม วิโมกฺโข.\csromanspacer{} สพฺพโส เนวสญฺญานาสญฺญายตนํ สมติกฺกมฺม สญฺญาเวทยิตนิโรธํ อุปสมฺปชฺช วิหรติ, อยํ อฏฺฐโม วิโมกฺโข.\csromanspacer{} อิเม โข อานนฺท อฏฺฐ วิโมกฺขา.}

\proseitem{130}{ยโต โข อานนฺท ภิกฺขุ อิเม อฏฺฐ วิโมกฺเข อนุโลมํปิ สมาปชฺชติ, ปฏิโลมํปิ สมาปชฺชติ, อนุโลมปฏิโลมํปิ สมาปชฺชติ, ยตฺถิจฺฉกํ ยทิจฺฉกํ ยาวติจฺฉกํ สมาปชฺชติปิ วุฏฺฐาติปิ, อาสวานญฺจ ขยา อนาสวํ เจโตวิมุตฺติํ ปญฺญาวิมุตฺติํ ทิฏฺเฐวธมฺเม สยํ อภิญฺญา สจฺฉิกตฺวา อุปสมฺปชฺช วิหรติ.\csromanspacer{} อยํ วุจฺจตานนฺท ภิกฺขุ อุภโตภาควิมุตฺโต.\csromanspacer{} อิมาย จ อานนฺท อุภโตภาควิมุตฺติยา อญฺญา อุภโตภาควิมุตฺติ อุตฺตริตรา วา ปณีตตรา วา นตฺถีติ.\csromanspacer{} อิทมโวจ ภควา.\csromanspacer{} อตฺตมโน อายสฺมา อานนฺโท ภควโต ภาสิตํ อภินนฺทีติ.}

\nitthitam{\textbf{มหานิทานสุตฺตํ นิฏฺฐิตํ ทุติยํ.}}
\csromanmatikaanchor{mtk.26}
\csromantocmarknopage{chapter}{3. มหาปรินิพฺพานสุตฺต}{mtk.26}
\bookmark[dest={mtk.26},level=0]{3. มหาปรินิพฺพานสุตฺต}
\chapterheadpageread{3. มหาปรินิพฺพานสุตฺต}
\proseitem{131}{\csromanfolio{61}เอวํ เม สุตํ\csromandash{}เอกํ สมยํ ภควา ราชคเห วิหรติ คิชฺฌกูเฏ ปพฺพเต.\csromanspacer{} เตน โข ปน สมเยน ราชา มาคโธ อชาตสตฺตุ เวเทหิปุตฺโต วชฺชี อภิยาตุกาโม โหติ.\csromanspacer{} โส เอวมาห “อหํ หิเม วชฺชี เอวํมหิทฺธิเก เอวํมหานุภาเว อุจฺเฉจฺฉามิ\footnote{อุจฺเฉชฺชามิ (สฺยา, อิ), อุจฺฉิชฺชามิ (ก)} วชฺชี วินาเสสฺสามิ วชฺชี อนยพฺยสนํ อาปาเทสฺสามี”ติ\footnote{อาปาเทสฺสามิ วชฺชีติ (สพฺพตฺถ) อํ 2. 409 ปสฺสิตพฺพํ.}.}

\proseitem{132}{อถ โข ราชา มาคโธ อชาตสตฺตุ เวเทหิปุตฺโต วสฺสการํ พฺราหฺมณํ มคธมหามตฺตํ อามนฺเตสิ “เอหิ ตฺวํ พฺราหฺมณ เยน ภควา เตนุปสงฺกม, อุปสงฺกมิตฺวา มม วจเนน ภควโต ปาเท สิรสา วนฺทาหิ, อปฺปาพาธํ อปฺปาตงฺกํ ลหุฏฺฐานํ พลํ ผาสุวิหารํ ปุจฺฉ, ‘ราชา ภนฺเต มาคโธ อชาตสตฺตุ เวเทหิปุตฺโต ภควโต ปาเท สิรสา วนฺทติ, อปฺปาพาธํ อปฺปาตงฺกํ ลหุฏฺฐานํ พลํ ผาสุวิหารํ ปุจฺฉตี’ติ.\csromanspacer{} เอวญฺจ วเทหิ ‘ราชา ภนฺเต มาคโธ อชาตสตฺตุ เวเทหิปุตฺโต วชฺชี อภิยาตุกาโม, โส เอวมาห อหํ หิเม วชฺชี เอวํมหิทฺธิเก เอวํมหานุภาเว อุจฺเฉจฺฉามิ วชฺชี วินาเสสฺสามิ วชฺชี อนยพฺยสนํ อาปาเทสฺสามี’ติ.\csromanspacer{} ยถา เต ภควา พฺยากโรติ, ตํ สาธุกํ อุคฺคเหตฺวา มม อาโรเจยฺยาสิ.\csromanspacer{} น หิ ตถาคตา วิตถํ ภณนฺตี”ติ.}

\csromanmatikaanchor{mtk.27}
\csromantocmarkref{section}{วสฺสการพฺราหฺมณ}{mtk.27}
\bookmark[dest={mtk.27},level=1]{วสฺสการพฺราหฺมณ}
\csromanheader{1}{\textbf{วสฺสการพฺราหฺมณ}}
\proseitem{133}{“เอวํ โภ”ติ โข วสฺสกาโร พฺราหฺมโณ มคธมหามตฺโต รญฺโญ มาคธสฺส อชาตสตฺตุสฺส เวเทหิปุตฺตสฺส ปฏิสฺสุตฺวา ภทฺทานิ ภทฺทานิ ยานานิ โยเชตฺวา ภทฺทํ ภทฺทํ ยานํ อภิรุหิตฺวา ภทฺเทหิ ภทฺเทหิ ยาเนหิ ราชคหมฺหา นิยฺยาสิ, เยน คิชฺฌกูโฏ ปพฺพโต เตน ปายาสิ.\csromanspacer{} ยาวติกา ยานสฺส ภูมิ, ยาเนน คนฺตฺวา ยานา ปจฺโจโรหิตฺวา ปตฺติโกว เยน ภควา เตนุปสงฺกมิ, อุปสงฺกมิตฺวา ภควตา สทฺธิํ สมฺโมทิ, สมฺโมทนียํ กถํ สารณียํ วีติสาเรตฺวา \csromanfolio{62}เอกมนฺตํ นิสีทิ, เอกมนฺตํ นิสินฺโน โข วสฺสกาโร พฺราหฺมโณ มคธมหามตฺโต ภควนฺตํ เอตทโวจ “ราชา โภ โคตม มาคโธ อชาตสตฺตุ เวเทหิปุตฺโต โภโต โคตมสฺส ปาเท สิรสา วนฺทติ, อปฺปาพาธํ อปฺปาตงฺกํ ลหุฏฺฐานํ พลํ ผาสุวิหารํ ปุจฺฉติ, ราชา\footnote{เอวญฺจ วเทติ ราชา (ก)} โภ โคตม มาคโธ อชาตสตฺตุ เวเทหิปุตฺโต วชฺชี อภิยาตุกาโม, โส เอวมาห อหํ หิเม วชฺชี เอวํมหิทฺธิเก เอวํมหานุภาเว อุจฺเฉจฺฉามิ วชฺชี วินาเสสฺสามิ วชฺชี อนยพฺยสนํ อาปาเทสฺสามี”ติ.}

\csromanmatikaanchor{mtk.28}
\csromantocmarkref{section}{ราช-อปริหานิยธมฺม}{mtk.28}
\bookmark[dest={mtk.28},level=1]{ราช-อปริหานิยธมฺม}
\csromanheader{1}{\textbf{ราช อปริหานิยธมฺม}}
\proseitem{134}{เตน โข ปน สมเยน อายสฺมา อานนฺโท ภควโต ปิฏฺฐิโต ฐิโต โหติ ภควนฺตํ พีชยมาโน\footnote{วีชยมาโน (สี), วีชิยมาโน (สฺยา)}.\csromanspacer{} อถ โข ภควา อายสฺมนฺตํ อานนฺทํ อามนฺเตสิ “กินฺติ เต อานนฺท สุตํ, วชฺชี อภิณฺหํ สนฺนิปาตา สนฺนิปาตพหุลา”ติ.\csromanspacer{} สุตํ เมตํ ภนฺเต “วชฺชี อภิณฺหํ สนฺนิปาตา สนฺนิปาตพหุลา”ติ.\csromanspacer{} ยาวกีวญฺจ อานนฺท วชฺชี อภิณฺหํ สนฺนิปาตา สนฺนิปาตพหุลา ภวิสฺสนฺติ, วุทฺธิเยว อานนฺท วชฺชีนํ ปาฏิกงฺขา, โน ปริหานิ. (1)}

\prose{กินฺติ เต อานนฺท สุตํ, วชฺชี สมคฺคา สนฺนิปตนฺติ, สมคฺคา วุฏฺฐหนฺติ, สมคฺคา วชฺชิกรณียานิ กโรนฺตีติ.\csromanspacer{} สุตํ เมตํ ภนฺเต “วชฺชี สมคฺคา สนฺนิปตนฺติ, สมคฺคา วุฏฺฐหนฺติ, สมคฺคา วชฺชิกรณียานิ กโรนฺตี”ติ.\csromanspacer{} ยาวกีวญฺจ อานนฺท วชฺชี สมคฺคา สนฺนิปติสฺสนฺติ, สมคฺคา วุฏฺฐหิสฺสนฺติ, สมคฺคา วชฺชิกรณียานิ กริสฺสนฺติ, วุทฺธิเยว อานนฺท วชฺชีนํ ปาฏิกงฺขา, โน ปริหานิ. (2)}

\prose{กินฺติ เต อานนฺท สุตํ, วชฺชี อปญฺญตฺตํ น ปญฺญเปนฺติ, ปญฺญตฺตํ น สมุจฺฉินฺทนฺติ, ยถาปญฺญตฺเต โปราเณ วชฺชิธมฺเม สมาทาย วตฺตนฺตีติ.\csromanspacer{} สุตํ เมตํ ภนฺเต “วชฺชี อปญฺญตฺตํ น ปญฺญเปนฺติ, ปญฺญตฺตํ น สมุจฺฉินฺทนฺติ, ยถาปญฺญตฺเต โปราเณ วชฺชิธมฺเม สมาทาย วตฺตนฺตี”ติ.\csromanspacer{} ยาวกีวญฺจ อานนฺท วชฺชี อปญฺญตฺตํ น ปญฺญเปสฺสนฺติ, ปญฺญตฺตํ น สมุจฺฉินฺทิสฺสนฺติ, ยถาปญฺญตฺเต โปราเณ วชฺชิธมฺเม สมาทาย วตฺติสฺสนฺติ, วุทฺธิเยว อานนฺท วชฺชีนํ ปาฏิกงฺขา, โน ปริหานิ. (3)}

\csromanheader{2}{3. มหาปรินิพฺพานสุตฺต}
\prose{\csromanfolio{63}กินฺติ เต อานนฺท สุตํ, วชฺชี เย เต วชฺชีนํ วชฺชิมหลฺลกา, เต สกฺกโรนฺติ ครุํ กโรนฺติ\footnote{ครุกโรนฺติ (สี, สฺยา, อิ)} มาเนนฺติ ปูเชนฺติ, เตสญฺจ โสตพฺพํ มญฺญนฺตีติ.\csromanspacer{} สุตํ เมตํ ภนฺเต “วชฺชี เย เต วชฺชีนํ วชฺชิมหลฺลกา, เต สกฺกโรนฺติ ครุํ กโรนฺติ มาเนนฺติ ปูเชนฺติ, เตสญฺจ โสตพฺพํ มญฺญนฺตี”ติ.\csromanspacer{} ยาวกีวญฺจ อานนฺท วชฺชี เย เต วชฺชีนํ วชฺชิมหลฺลกา, เต สกฺกริสฺสนฺติ ครุํ กริสฺสนฺติ มาเนสฺสนฺติ ปูเชสฺสนฺติ, เตสญฺจ โสตพฺพํ มญฺญิสฺสนฺติ, วุทฺธิเยว อานนฺท วชฺชีนํ ปาฏิกงฺขา, โน ปริหานิ. (4)}

\prose{กินฺติ เต อานนฺท สุตํ, วชฺชี ยา ตา กุลิตฺถิโย กุลกุมาริโย, ตา น โอกฺกสฺส ปสยฺห วาเสนฺตีติ.\csromanspacer{} สุตํ เมตํ ภนฺเต “วชฺชี ยา ตา กุลิตฺถิโย กุลกุมาริโย, ตา น โอกฺกสฺส ปสยฺห วาเสนฺตี”ติ.\csromanspacer{} ยาวกีวญฺจ อานนฺท วชฺชี ยา ตา กุลิตฺถิโย กุลกุมาริโย, ตา น โอกฺกสฺส ปสยฺห วาเสสฺสนฺติ, วุทฺธิเยว อานนฺท วชฺชีนํ ปาฏิกงฺขา, โน ปริหานิ. (5)}

\prose{กินฺติ เต อานนฺท สุตํ, วชฺชี ยานิ ตานิ วชฺชีนํ วชฺชิเจติยานิ อพฺภนฺตรานิ เจว พาหิรานิ จ, ตานิ สกฺกโรนฺติ ครุํ กโรนฺติ มาเนนฺติ ปูเชนฺติ, เตสญฺจ ทินฺนปุพฺพํ กตปุพฺพํ ธมฺมิกํ พลิํ โน ปริหาเปนฺตีติ.\csromanspacer{} สุตํ เมตํ ภนฺเต “วชฺชี ยานิ ตานิ วชฺชีนํ วชฺชิเจติยานิ อพฺภนฺตรานิ เจว พาหิรานิ จ, ตานิ สกฺกโรนฺติ ครุํ กโรนฺติ มาเนนฺติ ปูเชนฺติ, เตสญฺจ ทินฺนปุพฺพํ กตปุพฺพํ ธมฺมิกํ พลิํ โน ปริหาเปนฺตี”ติ.\csromanspacer{} ยาวกีวญฺจ อานนฺท วชฺชี ยานิ ตานิ วชฺชีนํ วชฺชิเจติยานิ อพฺภนฺตรานิ เจว พาหิรานิ จ, ตานิ สกฺกริสฺสนฺติ ครุํ กริสฺสนฺติ มาเนสฺสนฺติ ปูเชสฺสนฺติ, เตสญฺจ ทินฺนปุพฺพํ กตปุพฺพํ ธมฺมิกํ พลิํ โน ปริหาเปสฺสนฺติ, วุทฺธิเยว อานนฺท วชฺชีนํ ปาฏิกงฺขา, โน ปริหานิ. (6)}

\prose{กินฺติ เต อานนฺท สุตํ, วชฺชีนํ อรหนฺเตสุ ธมฺมิกา รกฺขาวรณคุตฺติ สุสํวิหิตา “กินฺติ อนาคตา จ อรหนฺโต วิชิตํ อาคจฺเฉยฺยุํ, อาคตา จ อรหนฺโต วิชิเต ผาสุ วิหเรยฺยุนฺ”ติ.\csromanspacer{} สุตํ เมตํ ภนฺเต “วชฺชีนํ อรหนฺเตสุ ธมฺมิกา รกฺขาวรณคุตฺติ สุสํวิหิตา ‘กินฺติ อนาคตา จ อรหนฺโต วิชิตํ อาคจฺเฉยฺยุํ, อาคตา จ อรหนฺโต วิชิเต ผาสุ วิหเรยฺยุนฺ’ติ”.\csromanspacer{} ยาวกีวญฺจ อานนฺท วชฺชีนํ อรหนฺเตสุ ธมฺมิกา รกฺขาวรณคุตฺติ สุสํวิหิตา ภวิสฺสติ.\csromanspacer{} กินฺติ อนาคตา จ อรหนฺโต}

\prose{\csromanfolio{64}วิชิตํ อาคจฺเฉยฺยุํ, อาคตา จ อรหนฺโต วิชิเต ผาสุ วิหเรยฺยุนฺติ, วุทฺธิเยว อานนฺท วชฺชีนํ ปาฏิกงฺขา, โน ปริหานีติ. (7)}

\proseitem{135}{อถ โข ภควา วสฺสการํ พฺราหฺมณํ มคธมหามตฺตํ อามนฺเตสิ “เอกมิทาหํ พฺราหฺมณ สมยํ เวสาลิยํ วิหรามิ สารนฺทเท\footnote{สานนฺทเร (ก)} เจติเย.\csromanspacer{} ตตฺราหํ วชฺชีนํ อิเม สตฺต อปริหานิเย ธมฺเม เทเสสิํ.\csromanspacer{} ยาวกีวญฺจ พฺราหฺมณ อิเม สตฺต อปริหานิยา ธมฺมา วชฺชีสุ ฐสฺสนฺติ, อิเมสุ จ สตฺตสุ อปริหานิเยสุ ธมฺเมสุ วชฺชี สนฺทิสฺสิสฺสนฺติ, วุทฺธิเยว พฺราหฺมณ วชฺชีนํ ปาฏิกงฺขา, โน ปริหานี”ติ.}

\prose{เอวํ วุตฺเต วสฺสกาโร พฺราหฺมโณ มคธมหามตฺโต ภควนฺตํ เอตทโวจ “เอกเมเกนปิ โภ โคตม อปริหานิเยน ธมฺเมน สมนฺนาคตานํ วชฺชีนํ วุทฺธิเยว ปาฏิกงฺขา, โน ปริหานิ.\csromanspacer{} โก ปน วาโท สตฺตหิ อปริหานิเยหิ ธมฺเมหิ.\csromanspacer{} อกรณียาว\footnote{อกรณียา จ (สฺยา, ก)} โภ โคตม วชฺชี\footnote{วชฺชีนํ (ก)} รญฺญา มาคเธน อชาตสตฺตุนา เวเทหิปุตฺเตน ยทิทํ ยุทฺธสฺส อญฺญตฺร อุปลาปนาย อญฺญตฺร มิถุเภทา.\csromanspacer{} หนฺท จ ทานิ มยํ โภ โคตม คจฺฉาม, พหุกิจฺจา มยํ พหุกรณียา”ติ.\csromanspacer{} ยสฺสทานิ ตฺวํ พฺราหฺมณ กาลํ มญฺญสีติ.\csromanspacer{} อถ โข วสฺสกาโร พฺราหฺมโณ มคธมหามตฺโต ภควโต ภาสิตํ อภินนฺทิตฺวา อนุโมทิตฺวา อุฏฺฐายาสนา ปกฺกามิ.}

\csromanmatikaanchor{mtk.29}
\csromantocmarkref{section}{ภิกฺขุ-อปริหานิยธมฺม}{mtk.29}
\bookmark[dest={mtk.29},level=1]{ภิกฺขุ-อปริหานิยธมฺม}
\csromanheader{1}{\textbf{ภิกฺขุ อปริหานิยธมฺม}}
\proseitem{136}{อถ โข ภควา อจิรปกฺกนฺเต วสฺสกาเร พฺราหฺมเณ มคธมหามตฺเต อายสฺมนฺตํ อานนฺทํ อามนฺเตสิ “คจฺฉ ตฺวํ อานนฺท ยาวติกา ภิกฺขู ราชคหํ อุปนิสฺสาย วิหรนฺติ, เต สพฺเพ อุปฏฺฐานสาลายํ สนฺนิปาเตหี”ติ. “เอวํ ภนฺเต”ติ โข อายสฺมา อานนฺโท ภควโต ปฏิสฺสุตฺวา ยาวติกา ภิกฺขู ราชคหํ อุปนิสฺสาย วิหรนฺติ, เต สพฺเพ อุปฏฺฐานสาลายํ สนฺนิปาเตตฺวา เยน ภควา เตนุปสงฺกมิ, อุปสงฺกมิตฺวา ภควนฺตํ อภิวาเทตฺวา เอกมนฺตํ อฏฺฐาสิ.\csromanspacer{} เอกมนฺตํ ฐิโต โข อายสฺมา อานนฺโท ภควนฺตํ เอตทโวจ “สนฺนิปติโต ภนฺเต ภิกฺขุสํโฆ, ยสฺสทานิ ภนฺเต ภควา กาลํ มญฺญตี”ติ.}

\csromanheader{2}{3. มหาปรินิพฺพานสุตฺต}
\prose{\csromanfolio{65}อถ โข ภควา อุฏฺฐายาสนา เยน อุปฏฺฐานสาลา เตนุปสงฺกมิ, อุปสงฺกมิตฺวา ปญฺญตฺเต อาสเน นิสีทิ, นิสชฺช โข ภควา ภิกฺขู อามนฺเตสิ “สตฺต โว ภิกฺขเว อปริหานิเย ธมฺเม เทเสสฺสามิ, ตํ สุณาถ สาธุกํ มนสิกโรถ ภาสิสฺสามี”ติ. “เอวํ ภนฺเต”ติ โข เต ภิกฺขู ภควโต ปจฺจสฺโสสุํ.\csromanspacer{} ภควา เอตทโวจ\csromandash{}}

\prose{ยาวกีวญฺจ ภิกฺขเว ภิกฺขู อภิณฺหํ สนฺนิปาตา สนฺนิปาตพหุลา ภวิสฺสนฺติ, วุทฺธิเยว ภิกฺขเว ภิกฺขูนํ ปาฏิกงฺขา, โน ปริหานิ. (1)}

\prose{ยาวกีวญฺจ ภิกฺขเว ภิกฺขู สมคฺคา สนฺนิปติสฺสนฺติ, สมคฺคา วุฏฺฐหิสฺสนฺติ, สมคฺคา สํฆกรณียานิ กริสฺสนฺติ, วุทฺธิเยว ภิกฺขเว ภิกฺขูนํ ปาฏิกงฺขา, โน ปริหานิ. (2)}

\prose{ยาวกีวญฺจ ภิกฺขเว ภิกฺขู อปญฺญตฺตํ น ปญฺญเปสฺสนฺติ, ปญฺญตฺตํ น สมุจฺฉินฺทิสฺสนฺติ, ยถาปญฺญตฺเตสุ สิกฺขาปเทสุ สมาทาย วตฺติสฺสนฺติ, วุทฺธิเยว ภิกฺขเว ภิกฺขูนํ ปาฏิกงฺขา, โน ปริหานิ. (3)}

\prose{ยาวกีวญฺจ ภิกฺขเว ภิกฺขู เย เต ภิกฺขู เถรา รตฺตญฺญู จิรปพฺพชิตา สํฆปิตโร สํฆปริณายกา, เต สกฺกริสฺสนฺติ ครุํ กริสฺสนฺติ มาเนสฺสนฺติ ปูเชสฺสนฺติ, เตสญฺจ โสตพฺพํ มญฺญิสฺสนฺติ, วุทฺธิเยว ภิกฺขเว ภิกฺขูนํ ปาฏิกงฺขา, โน ปริหานิ. (4)}

\prose{ยาวกีวญฺจ ภิกฺขเว ภิกฺขู อุปฺปนฺนาย ตณฺหาย โปโนพฺภวิกาย น วสํ คจฺฉิสฺสนฺติ, วุทฺธิเยว ภิกฺขเว ภิกฺขูนํ ปาฏิกงฺขา, โน ปริหานิ. (5)}

\prose{ยาวกีวญฺจ ภิกฺขเว ภิกฺขู อารญฺญเกสุ เสนาสเนสุ สาเปกฺขา ภวิสฺสนฺติ, วุทฺธิเยว ภิกฺขเว ภิกฺขูนํ ปาฏิกงฺขา, โน ปริหานิ. (6)}

\prose{ยาวกีวญฺจ ภิกฺขเว ภิกฺขู ปจฺจตฺตญฺเญว สติํ อุปฏฺฐเปสฺสนฺติ “กินฺติ อนาคตา จ เปสลา สพฺรหฺมจารี อาคจฺเฉยฺยุํ, อาคตา จ เปสลา สพฺรหฺมจารี ผาสุ\footnote{ผาสุํ (สี, สฺยา, อิ)} วิหเรยฺยุนฺ”ติ, วุทฺธิเยว ภิกฺขเว ภิกฺขูนํ ปาฏิกงฺขา, โน ปริหานิ. (7)}

\prose{ยาวกีวญฺจ ภิกฺขเว อิเม สตฺต อปริหานิยา ธมฺมา ภิกฺขูสุ ฐสฺสนฺติ, อิเมสุ จ สตฺตสุ อปริหานิเยสุ ธมฺเมสุ ภิกฺขู สนฺทิสฺสิสฺสนฺติ, วุทฺธิเยว ภิกฺขเว ภิกฺขูนํ ปาฏิกงฺขา, โน ปริหานิ.}

\proseitem{137}{\csromanfolio{66}อปเรปิ โว ภิกฺขเว สตฺต อปริหานิเย ธมฺเม เทเสสฺสามิ, ตํ สุณาถ สาธุกํ มนสิกโรถ ภาสิสฺสามีติ. “เอวํ ภนฺเต”ติ โข เต ภิกฺขู ภควโต ปจฺจสฺโสสุํ.\csromanspacer{} ภควา เอตทโวจ\csromandash{}}

\prose{ยาวกีวญฺจ ภิกฺขเว ภิกฺขู น กมฺมารามา ภวิสฺสนฺติ น กมฺมรตา น กมฺมารามตมนุยุตฺตา, วุทฺธิเยว ภิกฺขเว ภิกฺขูนํ ปาฏิกงฺขา, โน ปริหานิ. (1)}

\prose{ยาวกีวญฺจ ภิกฺขเว ภิกฺขู น ภสฺสารามา ภวิสฺสนฺติ น ภสฺสรตา น ภสฺสารามตมนุยุตฺตา, วุทฺธิเยว ภิกฺขเว ภิกฺขูนํ ปาฏิกงฺขา, โน ปริหานิ. (2)}

\prose{ยาวกีวญฺจ ภิกฺขเว ภิกฺขู น นิทฺทารามา ภวิสฺสนฺติ น นิทฺทารตา น นิทฺทารามตมนุยุตฺตา, วุทฺธิเยว ภิกฺขเว ภิกฺขูนํ ปาฏิกงฺขา, โน ปริหานิ. (3)}

\prose{ยาวกีวญฺจ ภิกฺขเว ภิกฺขู น สงฺคณิการามา ภวิสฺสนฺติ น สงฺคณิกรตา น สงฺคณิการามตมนุยุตฺตา, วุทฺธิเยว ภิกฺขเว ภิกฺขูนํ ปาฏิกงฺขา, โน ปริหานิ. (4)}

\prose{ยาวกีวญฺจ ภิกฺขเว ภิกฺขู น ปาปิจฺฉา ภวิสฺสนฺติ น ปาปิกานํ อิจฺฉานํ วสํ คตา, วุทฺธิเยว ภิกฺขเว ภิกฺขูนํ ปาฏิกงฺขา, โน ปริหานิ. (5)}

\prose{ยาวกีวญฺจ ภิกฺขเว ภิกฺขู น ปาปมิตฺตา ภวิสฺสนฺติ น ปาปสหายา น ปาปสมฺปวงฺกา, วุทฺธิเยว ภิกฺขเว ภิกฺขูนํ ปาฏิกงฺขา, โน ปริหานิ. (6)}

\prose{ยาวกีวญฺจ ภิกฺขเว ภิกฺขู น โอรมตฺตเกน วิเสสาธิคเมน อนฺตราโวสานํ อาปชฺชิสฺสนฺติ, วุทฺธิเยว ภิกฺขเว ภิกฺขูนํ ปาฏิกงฺขา, โน ปริหานิ. (7)}

\prose{ยาวกีวญฺจ ภิกฺขเว อิเม สตฺต อปริหานิยา ธมฺมา ภิกฺขูสุ ฐสฺสนฺติ, อิเมสุ จ สตฺตสุ อปริหานิเยสุ ธมฺเมสุ ภิกฺขู สนฺทิสฺสิสฺสนฺติ, วุทฺธิเยว ภิกฺขเว ภิกฺขูนํ ปาฏิกงฺขา, โน ปริหานิ.}

\proseitem{138}{อปเรปิ โว ภิกฺขเว สตฺต อปริหานิเย ธมฺเม เทเสสฺสามิ ฯเปฯ.\csromanspacer{} ยาวกีวญฺจ ภิกฺขเว ภิกฺขู สทฺธา ภวิสฺสนฺติ ฯเปฯ หิริมนา ภวิสฺสนฺติ.\csromanspacer{} โอตฺตปฺปี ภวิสฺสนฺติ.\csromanspacer{} พหุสฺสุตา ภวิสฺสนฺติ.\csromanspacer{} อารทฺธวีริยา ภวิสฺสนฺติ.}

\vspace{\tipitakaparskip}
\csromancenter{มหาปรินิพฺพานสุตฺต อุปฏฺฐิตสฺสตี ภวิสฺสนฺติ.\csromanspacer{} ปญฺญวนฺโต ภวิสฺสนฺติ, วุทฺธิเยว ภิกฺขเว ภิกฺขูนํ ปาฏิกงฺขา, โน ปริหานิ.\csromanspacer{} ยาวกีวญฺจ ภิกฺขเว อิเม สตฺต อปริหานิยา ธมฺมา ภิกฺขูสุ ฐสฺสนฺติ, อิเมสุ จ สตฺตสุ อปริหานิเยสุ ธมฺเมสุ ภิกฺขู สนฺทิสฺสิสฺสนฺติ, วุทฺธิเยว ภิกฺขเว ภิกฺขูนํ ปาฏิกงฺขา, โน ปริหานิ.}
\proseitem{139}{\csromanfolio{67}อปเรปิ โว ภิกฺขเว สตฺต อปริหานิเย ธมฺเม เทเสสฺสามิ, ตํ สุณาถ สาธุกํ มนสิกโรถ ภาสิสฺสามีติ. “เอวํ ภนฺเต”ติ โข เต ภิกฺขู ภควโต ปจฺจสฺโสสุํ.\csromanspacer{} ภควา เอตทโวจ\csromandash{}}

\prose{ยาวกีวญฺจ ภิกฺขเว ภิกฺขู สติสมฺโพชฺฌงฺคํ ภาเวสฺสนฺติ ฯเปฯ ธมฺมวิจยสมฺโพชฺฌงฺคํ ภาเวสฺสนฺติ.\csromanspacer{} วีริยสมฺโพชฺฌงฺคํ ภาเวสฺสนฺติ.\csromanspacer{} ปีติสมฺโพชฺฌงฺคํ ภาเวสฺสนฺติ.\csromanspacer{} ปสฺสทฺธิสมฺโพชฺฌงฺคํ ภาเวสฺสนฺติ.\csromanspacer{} สมาธิสมฺโพชฺฌงฺคํ ภาเวสฺสนฺติ.\csromanspacer{} อุเปกฺขาสมฺโพชฺฌงฺคํ ภาเวสฺสนฺติ, วุทฺธิเยว ภิกฺขเว ภิกฺขูนํ ปาฏิกงฺขา, โน ปริหานิ.}

\prose{ยาวกีวญฺจ ภิกฺขเว อิเม สตฺต อปริหานิยา ธมฺมา ภิกฺขูสุ ฐสฺสนฺติ, อิเมสุ จ สตฺตสุ อปริหานิเยสุ ธมฺเมสุ ภิกฺขู สนฺทิสฺสิสฺสนฺติ, วุทฺธิเยว ภิกฺขเว ภิกฺขูนํ ปาฏิกงฺขา, โน ปริหานิ.}

\proseitem{140}{อปเรปิ โว ภิกฺขเว สตฺต อปริหานิเย ธมฺเม เทเสสฺสามิ, ตํ สุณาถ สาธุกํ มนสิกโรถ ภาสิสฺสามีติ. “เอวํ ภนฺเต”ติ โข เต ภิกฺขู ภควโต ปจฺจสฺโสสุํ.\csromanspacer{} ภควา เอตทโวจ\csromandash{}}

\prose{ยาวกีวญฺจ ภิกฺขเว ภิกฺขู อนิจฺจสญฺญํ ภาเวสฺสนฺติ ฯเปฯ อนตฺตสญฺญํ ภาเวสฺสนฺติ.\csromanspacer{} อสุภสญฺญํ ภาเวสฺสนฺติ.\csromanspacer{} อาทีนวสญฺญํ ภาเวสฺสนฺติ.\csromanspacer{} ปหานสญฺญํ ภาเวสฺสนฺติ.\csromanspacer{} วิราคสญฺญํ ภาเวสฺสนฺติ.\csromanspacer{} นิโรธสญฺญํ ภาเวสฺสนฺติ, วุทฺธิเยว ภิกฺขเว ภิกฺขูนํ ปาฏิกงฺขา, โน ปริหานิ.}

\prose{ยาวกีวญฺจ ภิกฺขเว อิเม สตฺต อปริหานิยา ธมฺมา ภิกฺขูสุ ฐสฺสนฺติ, อิเมสุ จ สตฺตสุ อปริหานิเยสุ ธมฺเมสุ ภิกฺขู สนฺทิสฺสิสฺสนฺติ, วุทฺธิเยว ภิกฺขเว ภิกฺขูนํ ปาฏิกงฺขา, โน ปริหานิ.}

\proseitem{141}{ฉ โว ภิกฺขเว อปริหานิเย ธมฺเม เทเสสฺสามิ, ตํ สุณาถ สาธุกํ มนสิกโรถ ภาสิสฺสามีติ. “เอวํ ภนฺเต”ติ โข เต ภิกฺขู ภควโต ปจฺจสฺโสสุํ.\csromanspacer{} ภควา เอตทโวจ\csromandash{}}

\prose{\csromanfolio{68}ยาวกีวญฺจ ภิกฺขเว ภิกฺขู เมตฺตํ กายกมฺมํ ปจฺจุปฏฺฐาเปสฺสนฺติ สพฺรหฺมจารีสุ อาวิ เจว รโห จ, วุทฺธิเยว ภิกฺขเว ภิกฺขูนํ ปาฏิกงฺขา, โน ปริหานิ. (1)}

\prose{ยาวกีวญฺจ ภิกฺขเว ภิกฺขู เมตฺตํ วจีกมฺมํ ปจฺจุปฏฺฐาเปสฺสนฺติ ฯเปฯ เมตฺตํ มโนกมฺมํ ปจฺจุปฏฺฐาเปสฺสนฺติ สพฺรหฺมจารีสุ อาวิ เจว รโห จ, วุทฺธิเยว ภิกฺขเว ภิกฺขูนํ ปาฏิกงฺขา, โน ปริหานิ. (2-3)}

\prose{ยาวกีวญฺจ ภิกฺขเว ภิกฺขู เย เต ลาภา ธมฺมิกา ธมฺมลทฺธา อนฺตมโส ปตฺตปริยาปนฺนมตฺตมฺปิ, ตถารูเปหิ ลาเภหิ อปฺปฏิวิภตฺตโภคี ภวิสฺสนฺติ สีลวนฺเตหิ สพฺรหฺมจารีหิ สาธารณโภคี, วุทฺธิเยว ภิกฺขเว ภิกฺขูนํ ปาฏิกงฺขา, โน ปริหานิ. (4)}

\prose{ยาวกีวญฺจ ภิกฺขเว ภิกฺขู ยานิ ตานิ สีลานิ อขณฺฑานิ อจฺฉิทฺทานิ อสพลานิ อกมฺมาสานิ ภุชิสฺสานิ วิญฺญูปสตฺถานิ\footnote{วิญฺญุปฺปสตฺถานิ (สี)} อปรามฏฺฐานิ สมาธิสํวตฺตนิกานิ, ตถารูเปสุ สีเลสุ สีลสามญฺญคตา วิหริสฺสนฺติ สพฺรหฺมจารีหิ อาวิ เจว รโห จ, วุทฺธิเยว ภิกฺขเว ภิกฺขูนํ ปาฏิกงฺขา, โน ปริหานิ. (5)}

\prose{ยาวกีวญฺจ ภิกฺขเว ภิกฺขู ยายํ ทิฏฺฐิ อริยา นิยฺยานิกา นิยฺยาติ ตกฺกรสฺส สมฺมา ทุกฺขกฺขยาย, ตถารูปาย ทิฏฺฐิยา ทิฏฺฐิสามญฺญคตา วิหริสฺสนฺติ สพฺรหฺมจารีหิ อาวิ เจว รโห จ, วุทฺธิเยว ภิกฺขเว ภิกฺขูนํ ปาฏิกงฺขา, โน ปริหานิ. (6)}

\prose{ยาวกีวญฺจ ภิกฺขเว อิเม ฉ อปริหานิยา ธมฺมา ภิกฺขูสุ ฐสฺสนฺติ, อิเมสุ จ ฉสุ อปริหานิเยสุ ธมฺเมสุ ภิกฺขู สนฺทิสฺสิสฺสนฺติ, วุทฺธิเยว ภิกฺขเว ภิกฺขูนํ ปาฏิกงฺขา, โน ปริหานีติ.}

\proseitem{142}{ตตฺร สุทํ ภควา ราชคเห วิหรนฺโต คิชฺฌกูเฏ ปพฺพเต เอตเทว พหุลํ ภิกฺขูนํ ธมฺมิํ กถํ กโรติ “อิติ สีลํ อิติ สมาธิ อิติ ปญฺญา.\csromanspacer{} สีลปริภาวิโต สมาธิ มหปฺผโล โหติ มหานิสํโส.\csromanspacer{} สมาธิปริภาวิตา ปญฺญา มหปฺผลา โหติ มหานิสํสา.\csromanspacer{} ปญฺญาปริภาวิตํ จิตฺตํ สมฺมเทว อาสเวหิ วิมุจฺจติ.\csromanspacer{} เสยฺยถิทํ, กามาสวา ภวาสวา อวิชฺชาสวา”ติ.}

\csromanheader{2}{3. มหาปรินิพฺพานสุตฺต}
\proseitem{143}{\csromanfolio{69}อถ โข ภควา ราชคเห ยถาภิรนฺตํ วิหริตฺวา อายสฺมนฺตํ อานนฺทํ อามนฺเตสิ “อายามานนฺท เยน อมฺพลฏฺฐิกา เตนุปสงฺกมิสฺสามา”ติ. “เอวํ ภนฺเต”ติ โข อายสฺมา อานนฺโท ภควโต ปจฺจสฺโสสิ.\csromanspacer{} อถ โข ภควา มหตา ภิกฺขุสํเฆน สทฺธิํ เยน อมฺพลฏฺฐิกา ตทวสริ, ตตฺร สุทํ ภควา อมฺพลฏฺฐิกายํ วิหรติ ราชาคารเก.\csromanspacer{} ตตฺราปิ สุทํ ภควา อมฺพลฏฺฐิกายํ วิหรนฺโต ราชาคารเก เอตเทว พหุลํ ภิกฺขูนํ ธมฺมิํ กถํ กโรติ “อิติ สีลํ อิติ สมาธิ อิติ ปญฺญา.\csromanspacer{} สีลปริภาวิโต สมาธิ มหปฺผโล โหติ มหานิสํโส.\csromanspacer{} สมาธิปริภาวิตา ปญฺญา มหปฺผลา โหติ มหานิสํสา.\csromanspacer{} ปญฺญาปริภาวิตํ จิตฺตํ สมฺมเทว อาสเวหิ วิมุจฺจติ.\csromanspacer{} เสยฺยถิทํ, กามาสวา ภวาสวา อวิชฺชาสวา”ติ.}

\proseitem{144}{อถ โข ภควา อมฺพลฏฺฐิกายํ ยถาภิรนฺตํ วิหริตฺวา อายสฺมนฺตํ อานนฺทํ อามนฺเตสิ “อายามานนฺท เยน นาฬนฺทา เตนุปสงฺกมิสฺสามา”ติ. “เอวํ ภนฺเต”ติ โข อายสฺมา อานนฺโท ภควโต ปจฺจสฺโสสิ.\csromanspacer{} อถ โข ภควา มหตา ภิกฺขุสํเฆน สทฺธิํ เยน นาฬนฺทา ตทวสริ, ตตฺร สุทํ ภควา นาฬนฺทายํ วิหรติ ปาวาริกมฺพวเน.}

\csromanmatikaanchor{mtk.30}
\csromantocmarkref{section}{สาริปุตฺตสีหนาท}{mtk.30}
\bookmark[dest={mtk.30},level=1]{สาริปุตฺตสีหนาท}
\csromanheader{1}{\textbf{สาริปุตฺตสีหนาท}}
\proseitem{145}{อถ โข อายสฺมา สาริปุตฺโต เยน ภควา เตนุปสงฺกมิ อุปสงฺกมิตฺวา ภควนฺตํ อภิวาเทตฺวา เอกมนฺตํ นิสีทิ, เอกมนฺตํ นิสินฺโน โข อายสฺมา สาริปุตฺโต ภควนฺตํ เอตทโวจ “เอวํปสนฺโน อหํ ภนฺเต ภควติ, น จาหุ น จ ภวิสฺสติ น เจตรหิ วิชฺชติ อญฺโญ สมโณ วา พฺราหฺมโณ วา ภควตา ภิยฺโยภิญฺญตโร ยทิทํ สมฺโพธิยนฺ”ติ.\csromanspacer{} อุฬารา โข เต อยํ สาริปุตฺต อาสภี วาจา\footnote{อาสภิวาจา (สฺยา)} ภาสิตา, เอกํโส คหิโต, สีหนาโท นทิโต “เอวํปสนฺโน อหํ ภนฺเต ภควติ, น จาหุ น จ ภวิสฺสติ น เจตรหิ วิชฺชติ อญฺโญ สมโณ วา พฺราหฺมโณ วา ภควตา ภิยฺโยภิญฺญตโร ยทิทํ สมฺโพธิยนฺ”ติ.}

\prose{\csromanfolio{70}กิํ เต\footnote{กิํนุ (สฺยา, อิ, ก)} สาริปุตฺต เย เต อเหสุํ อตีตมทฺธานํ อรหนฺโต สมฺมาสมฺพุทฺธา, สพฺเพเต ภควนฺโต เจตสา เจโต ปริจฺจ วิทิตา “เอวํสีลา เต ภควนฺโต อเหสุํ อิติปิ, เอวํธมฺมา เอวํปญฺญา เอวํวิหารี เอวํวิมุตฺตา เต ภควนฺโต อเหสุํ อิติปี”ติ.\csromanspacer{} โน เหตํ ภนฺเต.}

\prose{กิํ ปน เต\footnote{กิํ ปน (สฺยา, อิ, ก)} สาริปุตฺต เย เต ภวิสฺสนฺติ อนาคตมทฺธานํ อรหนฺโต สมฺมาสมฺพุทฺธา, สพฺเพเต ภควนฺโต เจตสา เจโต ปริจฺจ วิทิตา “เอวํสีลา เต ภควนฺโต ภวิสฺสนฺติ อิติปิ, เอวํธมฺมา เอวํปญฺญา เอวํวิหารี เอวํวิมุตฺตา เต ภควนฺโต ภวิสฺสนฺติ อิติปี”ติ.\csromanspacer{} โน เหตํ ภนฺเต.}

\prose{กิํ ปน เต สาริปุตฺต อหํ เอตรหิ อรหํ สมฺมาสมฺพุทฺโธ เจตสา เจโต ปริจฺจ วิทิโต “เอวํสีโล ภควา อิติปิ, เอวํธมฺโม เอวํปญฺโญ เอวํวิหารี เอวํวิมุตฺโต ภควา อิติปี”ติ.\csromanspacer{} โน เหตํ ภนฺเต.}

\prose{เอตฺถ จ หิ เต สาริปุตฺต อตีตานาคตปจฺจุปฺปนฺเนสุ อรหนฺเตสุ สมฺมาสมฺพุทฺเธสุ เจโตปริยญาณํ\footnote{เจโตปริญฺญายญาณํ (สฺยา), เจตสา เจโตปริยายญาณํ (ก)} นตฺถิ.\csromanspacer{} อถ กิญฺจรหิ เต อยํ สาริปุตฺต อุฬารา อาสภี วาจา ภาสิตา, เอกํโส คหิโต สีหนาโท นทิโต “เอวํปสนฺโน อหํ ภนฺเต ภควติ, น จาหุ น จ ภวิสฺสติ น เจตรหิ วิชฺชติ อญฺโญ สมโณ วา พฺราหฺมโณ วา ภควตา ภิยฺโยภิญฺญตโร ยทิทํ สมฺโพธิยนฺ”ติ.}

\proseitem{146}{น โข เม ภนฺเต อตีตานาคตปจฺจุปฺปนฺเนสุ อรหนฺเตสุ สมฺมาสมฺพุทฺเธสุ เจโตปริยญาณํ อตฺถิ, อปิ จ เม ธมฺมนฺวโย วิทิโต.\csromanspacer{} เสยฺยถาปิ ภนฺเต รญฺโญ ปจฺจนฺติมํ นครํ ทฬฺหุทฺธาปํ ทฬฺหปาการโตรณํ เอกทฺวารํ, ตตฺรสฺส โทเวริโก ปณฺฑิโต วิยตฺโต เมธาวี อญฺญาตานํ นิวาเรตา ญาตานํ ปเวเสตา, โส ตสฺส นครสฺส สมนฺตา อนุปริยายปถํ\footnote{อนุจริยายปถํ (สฺยา)} อนุกฺกมมาโน น ปสฺเสยฺย ปาการสนฺธิํ วา ปาการวิวรํ วา อนฺตมโส พิฬารนิกฺขมนมตฺตมฺปิ.\csromanspacer{} ตสฺส เอวมสฺส\footnote{น ปสฺเสยฺย ตสฺส เอวมสฺส (สฺยา)} “เย โข เกจิ โอฬาริกา ปาณา อิมํ นครํ ปวิสนฺติ วา นิกฺขมนฺติ วา, สพฺเพเต อิมินาว ทฺวาเรน ปวิสนฺติ วา นิกฺขมนฺติ วา”ติ.\csromanspacer{} เอวเมว โข เม ภนฺเต ธมฺมนฺวโย วิทิโต.\csromanspacer{} เย เต ภนฺเต อเหสุํ อตีตมทฺธานํ อรหนฺโต}

\vspace{\tipitakaparskip}
\csromancenter{มหาปรินิพฺพานสุตฺต สมฺมาสมฺพุทฺธา, สพฺเพเต ภควนฺโต ปญฺจ นีวรเณ ปหาย เจตโส อุปกฺกิเลเส ปญฺญาย ทุพฺพลีกรเณ จตูสุ สติปฏฺฐาเนสุ สุปติฏฺฐิตจิตฺตา สตฺต โพชฺฌงฺเค ยถาภูตํ ภาเวตฺวา อนุตฺตรํ สมฺมาสมฺโพธิํ อภิสมฺพุชฺฌิํสุ.\csromanspacer{} เยปิ เต ภนฺเต ภวิสฺสนฺติ อนาคตมทฺธานํ อรหนฺโต สมฺมาสมฺพุทฺธา, สพฺเพเต ภควนฺโต ปญฺจ นีวรเณ ปหาย เจตโส อุปกฺกิเลเส ปญฺญาย ทุพฺพลีกรเณ จตูสุ สติปฏฺฐาเนสุ สุปติฏฺฐิตจิตฺตา สตฺต โพชฺฌงฺเค ยถาภูตํ ภาเวตฺวา อนุตฺตรํ สมฺมาสมฺโพธิํ อภิสมฺพุชฺฌิสฺสนฺติ.\csromanspacer{} ภควาปิ ภนฺเต เอตรหิ อรหํ สมฺมาสมฺพุทฺโธ ปญฺจ นีวรเณ ปหาย เจตโส อุปกฺกิเลเส ปญฺญาย ทุพฺพลีกรเณ จตูสุ สติปฏฺฐาเนสุ สุปติฏฺฐิตจิตฺโต สตฺต โพชฺฌงฺเค ยถาภูตํ ภาเวตฺวา อนุตฺตรํ สมฺมาสมฺโพธิํ อภิสมฺพุทฺโธ”ติ.}
\proseitem{147}{\csromanfolio{71}ตตฺรปิ สุทํ ภควา นาฬนฺทายํ วิหรนฺโต ปาวาริกมฺพวเน เอตเทว พหุลํ ภิกฺขูนํ ธมฺมิํ กถํ กโรติ “อิติ สีลํ อิติ สมาธิ อิติ ปญฺญา.\csromanspacer{} สีลปริภาวิโต สมาธิ มหปฺผโล โหติ มหานิสํโส.\csromanspacer{} สมาธิปริภาวิตา ปญฺญา มหปฺผลา โหติ มหานิสํสา.\csromanspacer{} ปญฺญาปริภาวิตํ จิตฺตํ สมฺมเทว อาสเวหิ วิมุจฺจติ.\csromanspacer{} เสยฺยถิทํ, กามาสวา ภวาสวา อวิชฺชาสวา”ติ.}

\csromanmatikaanchor{mtk.31}
\csromantocmarkref{section}{ทุสฺสีล-อาทีนว}{mtk.31}
\bookmark[dest={mtk.31},level=1]{ทุสฺสีล-อาทีนว}
\csromanheader{1}{\textbf{ทุสฺสีล อาทีนว}}
\proseitem{148}{อถ โข ภควา นาฬนฺทายํ ยถาภิรนฺตํ วิหริตฺวา อายสฺมนฺตํ อานนฺทํ อามนฺเตสิ “อายามานนฺท เยน ปาฏลิคาโม เตนุปสงฺกมิสฺสามา”ติ. “เอวํ ภนฺเต”ติ โข อายสฺมา อานนฺโท ภควโต ปจฺจสฺโสสิ.\csromanspacer{} อถ โข ภควา มหตา ภิกฺขุสํเฆน สทฺธิํ เยน ปาฏลิคาโม ตทวสริ.\csromanspacer{} อสฺโสสุํ โข ปาฏลิคามิกา อุปาสกา “ภควา กิร ปาฏลิคามํ อนุปฺปตฺโต”ติ.\csromanspacer{} อถ โข ปาฏลิคามิกา อุปาสกา เยน ภควา เตนุปสงฺกมิํสุ, อุปสงฺกมิตฺวา ภควนฺตํ อภิวาเทตฺวา เอกมนฺตํ นิสีทิํสุ, เอกมนฺตํ นิสินฺนา โข ปาฏลิคามิกา อุปาสกา ภควนฺตํ เอตทโวจุํ “อธิวาเสตุ โน ภนฺเต ภควา อาวสถาคารนฺ”ติ.\csromanspacer{} อธิวาเสสิ ภควา ตุณฺหีภาเวน.\csromanspacer{} อถ โข ปาฏลิคามิกา อุปาสกา ภควโต อธิวาสนํ วิทิตฺวา อุฏฺฐายาสนา ภควนฺตํ อภิวาเทตฺวา ปทกฺขิณํ กตฺวา เยน อาวสถาคารํ เตนุปสงฺกมิํสุ, อุปสงฺกมิตฺวา \csromanfolio{72}สพฺพสนฺถริํ\footnote{สพฺพสนฺถริตํ สนฺถตํ (สฺยา), สพฺพสนฺถริํ สนฺถตํ (ก)} อาวสถาคารํ สนฺถริตฺวา อาสนานิ ปญฺญาเปตฺวา อุทกมณิกํ ปติฏฺฐาเปตฺวา เตลปทีปํ อาโรเปตฺวา เยน ภควา เตนุปสงฺกมิํสุ, อุปสงฺกมิตฺวา ภควนฺตํ อภิวาเทตฺวา เอกมนฺตํ อฏฺฐํสุ, เอกมนฺตํ ฐิตา โข ปาฏลิคามิกา อุปาสกา ภควนฺตํ เอตทโวจุํ “สพฺพสนฺถริสนฺถตํ\footnote{สพฺพสนฺถริํ สนฺถตํ (สี, สฺยา, อิ, ก)} ภนฺเต อาวสถาคารํ, อาสนานิ ปญฺญตฺตานิ, อุทกมณิโก ปติฏฺฐาปิโต, เตลปทีโป อาโรปิโต, ยสฺสทานิ ภนฺเต ภควา กาลํ มญฺญตี”ติ.\csromanspacer{} อถ โข ภควา สายนฺหสมยํ\footnote{อิทํ ปทํ วินยมหาวคฺเค น ทิสฺสติ.} นิวาเสตฺวา ปตฺตจีวรมาทาย สทฺธิํ ภิกฺขุสํเฆน เยน อาวสถาคารํ เตนุปสงฺกมิ, อุปสงฺกมิตฺวา ปาเท ปกฺขาเลตฺวา อาวสถาคารํ ปวิสิตฺวา มชฺฌิมํ ถมฺภํ นิสฺสาย ปุรตฺถาภิมุโข\footnote{ปุรตฺถิมาภิมุโข (ก)} นิสีทิ.\csromanspacer{} ภิกฺขุสํโฆปิ โข ปาเท ปกฺขาเลตฺวา อาวสถาคารํ ปวิสิตฺวา ปจฺฉิมํ ภิตฺติํ นิสฺสาย ปุรตฺถาภิมุโข นิสีทิ ภควนฺตเมว ปุรกฺขตฺวา.\csromanspacer{} ปาฏลิคามิกาปิ โข อุปาสกา ปาเท ปกฺขาเลตฺวา อาวสถาคารํ ปวิสิตฺวา ปุรตฺถิมํ ภิตฺติํ นิสฺสาย ปจฺฉิมาภิมุขา นิสีทิํสุ ภควนฺตเมว ปุรกฺขตฺวา.}

\proseitem{149}{อถ โข ภควา ปาฏลิคามิเก อุปาสเก อามนฺเตสิ\csromandash{}ปญฺจิเม คหปตโย อาทีนวา ทุสฺสีลสฺส สีลวิปตฺติยา.\csromanspacer{} กตเม ปญฺจ.\csromanspacer{} อิธ คหปตโย ทุสฺสีโล สีลวิปนฺโน ปมาทาธิกรณํ มหติํ โภคชานิํ นิคจฺฉติ, อยํ ปฐโม อาทีนโว ทุสฺสีลสฺส สีลวิปตฺติยา.}

\prose{ปุน จปรํ คหปตโย ทุสฺสีลสฺส สีลวิปนฺนสฺส ปาปโก กิตฺติสทฺโท อพฺภุคฺคจฺฉติ, อยํ ทุติโย อาทีนโว ทุสฺสีลสฺส สีลวิปตฺติยา.}

\prose{ปุน จปรํ คหปตโย ทุสฺสีโล สีลวิปนฺโน ยญฺญเทว ปริสํ อุปสงฺกมติ ยทิ ขตฺติยปริสํ ยทิ พฺราหฺมณปริสํ ยทิ คหปติปริสํ ยทิ สมณปริสํ, อวิสารโท อุปสงฺกมติ มงฺกุภูโต, อยํ ตติโย อาทีนโว ทุสฺสีลสฺส สีลวิปตฺติยา.}

\prose{ปุน จปรํ คหปตโย ทุสฺสีโล สีลวิปนฺโน สมฺมูโฬฺห กาลํ กโรติ, อยํ จตุตฺโถ อาทีนโว ทุสฺสีลสฺส สีลวิปตฺติยา.}

\csromanheader{2}{3. มหาปรินิพฺพานสุตฺต}
\prose{\csromanfolio{73}ปุน จปรํ คหปตโย ทุสฺสีโล สีลวิปนฺโน กายสฺส เภทา ปรํ มรณา อปายํ ทุคฺคติํ วินิปาตํ นิรยํ อุปปชฺชติ, อยํ ปญฺจโม อาทีนโว ทุสฺสีลสฺส สีลวิปตฺติยา.\csromanspacer{} อิเม โข คหปตโย ปญฺจ อาทีนวา ทุสฺสีลสฺส สีลวิปตฺติยา.}

\csromanmatikaanchor{mtk.32}
\csromantocmarkref{section}{สีลวนฺต-อานิสํส}{mtk.32}
\bookmark[dest={mtk.32},level=1]{สีลวนฺต-อานิสํส}
\csromanheader{1}{\textbf{สีลวนฺต อานิสํส}}
\proseitem{150}{ปญฺจิเม คหปตโย อานิสํสา สีลวโต สีลสมฺปทาย.\csromanspacer{} กตเม ปญฺจ.\csromanspacer{} อิธ คหปตโย สีลวา สีลสมฺปนฺโน อปฺปมาทาธิกรณํ มหนฺตํ โภคกฺขนฺธํ อธิคจฺฉติ, อยํ ปฐโม อานิสํโส สีลวโต สีลสมฺปทาย.}

\prose{ปุน จปรํ คหปตโย สีลวโต สีลสมฺปนฺนสฺส กลฺยาโณ กิตฺติสทฺโท อพฺภุคฺคจฺฉติ, อยํ ทุติโย อานิสํโส สีลวโต สีลสมฺปทาย.}

\prose{ปุน จปรํ คหปตโย สีลวา สีลสมฺปนฺโน ยญฺญเทว ปริสํ อุปสงฺกมติ ยทิ ขตฺติยปริสํ ยทิ พฺราหฺมณปริสํ ยทิ คหปติปริสํ ยทิ สมณปริสํ, วิสารโท อุปสงฺกมติ อมงฺกุภูโต, อยํ ตติโย อานิสํโส สีลวโต สีลสมฺปทาย.}

\prose{ปุน จปรํ คหปตโย สีลวา สีลสมฺปนฺโน อสมฺมูโฬฺห กาลํ กโรติ, อยํ จตุตฺโถ อานิสํโส สีลวโต สีลสมฺปทาย.}

\prose{ปุน จปรํ คหปตโย สีลวา สีลสมฺปนฺโน กายสฺส เภทา ปรํ มรณา สุคติํ สคฺคํ โลกํ อุปปชฺชติ, อยํ ปญฺจโม อานิสํโส สีลวโต สีลสมฺปทาย.\csromanspacer{} อิเม โข คหปตโย ปญฺจ อานิสํสา สีลวโต สีลสมฺปทายา”ติ.}

\proseitem{151}{อถ โข ภควา ปาฏลิคามิเก อุปาสเก พหุเทว รตฺติํ ธมฺมิยา กถาย สนฺทสฺเสตฺวา สมาทเปตฺวา สมุตฺเตเชตฺวา สมฺปหํเสตฺวา อุยฺโยเชสิ “อภิกฺกนฺตา โข คหปตโย รตฺติ, ยสฺสทานิ ตุเมฺห กาลํ มญฺญถา”ติ. “เอวํ ภนฺเต”ติ โข ปาฏลิคามิกา อุปาสกา ภควโต ปฏิสฺสุตฺวา อุฏฺฐายาสนา ภควนฺตํ อภิวาเทตฺวา ปทกฺขิณํ กตฺวา ปกฺกมิํสุ.\csromanspacer{} อถ โข ภควา อจิรปกฺกนฺเตสุ ปาฏลิคามิเกสุ อุปาสเกสุ สุญฺญาคารํ ปาวิสิ.}

\csromanmatikaanchor{mtk.33}
\csromantocmarkref{section}{ปาฏลิปุตฺตนครมาปน}{mtk.33}
\bookmark[dest={mtk.33},level=1]{ปาฏลิปุตฺตนครมาปน}
\csromanheader{1}{\textbf{ปาฏลิปุตฺตนครมาปน}}
\proseitem{152}{\csromanfolio{74}เตน โข ปน สมเยน สุนิธวสฺสการา\footnote{สุนีธวสฺสการา (สฺยา, ก)} มคธมหามตฺตา ปาฏลิคาเม นครํ มาเปนฺติ วชฺชีนํ ปฏิพาหาย.\csromanspacer{} เตน สมเยน สมฺพหุลา เทวตาโย สหสฺเสว\footnote{สหสฺสสฺเสว (สี, อิ, ก), สหสฺสเสว (ฏีกายํ ปาฐนฺตรํ), สหสฺสสหสฺเสว (อุทานฏฺฐกถา)} ปาฏลิคาเม วตฺถูนิ ปริคฺคณฺหนฺติ.\csromanspacer{} ยสฺมิํ ปเทเส มเหสกฺขา เทวตา วตฺถูนิ ปริคฺคณฺหนฺติ, มเหสกฺขานํ ตตฺถ รญฺญํ ราชมหามตฺตานํ จิตฺตานิ นมนฺติ นิเวสนานิ มาเปตุํ.\csromanspacer{} ยสฺมิํ ปเทเส มชฺฌิมา เทวตา วตฺถูนิ ปริคฺคณฺหนฺติ, มชฺฌิมานํ ตตฺถ รญฺญํ ราชมหามตฺตานํ จิตฺตานิ นมนฺติ นิเวสนานิ มาเปตุํ.\csromanspacer{} ยสฺมิํ ปเทเส นีจา เทวตา วตฺถูนิ ปริคฺคณฺหนฺติ, นีจานํ ตตฺถ รญฺญํ ราชมหามตฺตานํ จิตฺตานิ นมนฺติ นิเวสนานิ มาเปตุํ.\csromanspacer{} อทฺทสา โข ภควา ทิพฺเพน จกฺขุนา วิสุทฺเธน อติกฺกนฺตมานุสเกน ตา เทวตาโย สหสฺเสว ปาฏลิคาเม วตฺถูนิ ปริคฺคณฺหนฺติโย.\csromanspacer{} อถ โข ภควา รตฺติยา ปจฺจูสสมยํ ปจฺจุฏฺฐาย อายสฺมนฺตํ อานนฺทํ อามนฺเตสิ “เก นุ โข\footnote{โก นุ โข (สี, สฺยา, อิ, ก)} อานนฺท ปาฏลิคาเม นครํ มาเปนฺตี”ติ\footnote{มาเปตีติ (สี, สฺยา, อิ, ก)}.\csromanspacer{} สุนิธวสฺสการา ภนฺเต มคธมหามตฺตา ปาฏลิคาเม นครํ มาเปนฺติ วชฺชีนํ ปฏิพาหายาติ.\csromanspacer{} เสยฺยถาปิ อานนฺท เทเวหิ ตาวติํเสหิ สทฺธิํ มนฺเตตฺวา, เอวเมว โข อานนฺท สุนิธวสฺสการา มคธมหามตฺตา ปาฏลิคาเม นครํ มาเปนฺติ วชฺชีนํ ปฏิพาหาย.\csromanspacer{} อิธาหํ อานนฺท อทฺทสํ ทิพฺเพน จกฺขุนา วิสุทฺเธน อติกฺกนฺตมานุสเกน สมฺพหุลา เทวตาโย สหสฺเสว ปาฏลิคาเม วตฺถูนิ ปริคฺคณฺหนฺติโย.\csromanspacer{} ยสฺมิํ อานนฺท ปเทเส มเหสกฺขา เทวตา วตฺถูนิ ปริคฺคณฺหนฺติ, มเหสกฺขานํ ตตฺถ รญฺญํ ราชมหามตฺตานํ จิตฺตานิ นมนฺติ นิเวสนานิ มาเปตุํ.\csromanspacer{} ยสฺมิํ ปเทเส มชฺฌิมา เทวตา วตฺถูนิ ปริคฺคณฺหนฺติ, มชฺฌิมานํ ตตฺถ รญฺญํ ราชมหามตฺตานํ จิตฺตานิ นมนฺติ นิเวสนานิ มาเปตุํ.\csromanspacer{} ยสฺมิํ ปเทเส นีจา เทวตา วตฺถูนิ ปริคฺคณฺหนฺติ, นีจานํ ตตฺถ รญฺญํ ราชมหามตฺตานํ จิตฺตานิ นมนฺติ นิเวสนานิ มาเปตุํ.\csromanspacer{} ยาวตา อานนฺท อริยํ อายตนํ ยาวตา วณิปฺปโถ อิทํ อคฺคนครํ ภวิสฺสติ ปาฏลิปุตฺตํ ปุฏเภทนํ.\csromanspacer{} ปาฏลิปุตฺตสฺส โข อานนฺท ตโย อนฺตรายา ภวิสฺสนฺติ อคฺคิโต วา อุทกโต วา มิถุเภทา วา”ติ.}

\csromanheader{2}{3. มหาปรินิพฺพานสุตฺต}
\proseitem{153}{\csromanfolio{75}อถ โข สุนิธวสฺสการา มคธมหามตฺตา เยน ภควา เตนุปสงฺกมิํสุ, อุปสงฺกมิตฺวา ภควตา สทฺธิํ สมฺโมทิํสุ, สมฺโมทนียํ กถํ สารณียํ วีติสาเรตฺวา เอกมนฺตํ อฏฺฐํสุ, เอกมนฺตํ ฐิตา โข สุนิธวสฺสการา มคธมหามตฺตา ภควนฺตํ เอตทโวจุํ “อธิวาเสตุ โน ภวํ โคตโม อชฺชตนาย ภตฺตํ สทฺธิํ ภิกฺขุสํเฆนา”ติ.\csromanspacer{} อธิวาเสสิ ภควา ตุณฺหีภาเวน.\csromanspacer{} อถ โข สุนิธวสฺสการา มคธมหามตฺตา ภควโต อธิวาสนํ วิทิตฺวา เยน สโก อาวสโถ เตนุปสงฺกมิํสุ, อุปสงฺกมิตฺวา สเก อาวสเถ ปณีตํ ขาทนียํ โภชนียํ ปฏิยาทาเปตฺวา ภควโต กาลํ อาโรจาเปสุํ “กาโล โภ โคตม นิฏฺฐิตํ ภตฺตนฺ”ติ.}

\prose{อถ โข ภควา ปุพฺพณฺหสมยํ นิวาเสตฺวา ปตฺตจีวรมาทาย สทฺธิํ ภิกฺขุสํเฆน เยน สุนิธวสฺสการานํ มคธมหามตฺตานํ อาวสโถ เตนุปสงฺกมิ, อุปสงฺกมิตฺวา ปญฺญตฺเต อาสเน นิสีทิ.\csromanspacer{} อถ โข สุนิธวสฺสการา มคธมหามตฺตา พุทฺธปฺปมุขํ ภิกฺขุสํฆํ ปณีเตน ขาทนีเยน โภชนีเยน สหตฺถา สนฺตปฺเปสุํ สมฺปวาเรสุํ.\csromanspacer{} อถ โข สุนิธวสฺสการา มคธมหามตฺตา ภควนฺตํ ภุตฺตาวิํ โอนีตปตฺตปาณิํ อญฺญตรํ นีจํ อาสนํ คเหตฺวา เอกมนฺตํ นิสีทิํสุ, เอกมนฺตํ นิสินฺเน โข สุนิธวสฺสกาเร มคธมหามตฺเต ภควา อิมาหิ คาถาหิ อนุโมทิ\csromandash{}}

\csromangathameasure{“ยสฺมิํ ปเทเส กปฺเปติ, วาสํ ปณฺฑิตชาติโย. \\ สีลวนฺเตตฺถ โภเชตฺวา, สญฺญเต พฺรหฺมจารโย. \\ ยา ตตฺถ เทวตา อาสุํ, ตาสํ ทกฺขิณมาทิเส. \\ ตา ปูชิตา ปูชยนฺติ, มานิตา มานยนฺติ นํ. \\ ตโต นํ อนุกมฺปนฺติ, มาตา ปุตฺตํว โอรสํ. \\ เทวตานุกมฺปิโต โปโส, สทา ภทฺรานิ ปสฺสตี”ติ.}
\csromangathagroup{\csromangathastanza{“ยสฺมิํ ปเทเส กปฺเปติ, วาสํ ปณฺฑิตชาติโย. \\ สีลวนฺเตตฺถ โภเชตฺวา, สญฺญเต พฺรหฺมจารโย\footnote{พฺรหฺมจาริโน (สฺยา)}.}\csromangathastanza{ยา ตตฺถ เทวตา อาสุํ, ตาสํ ทกฺขิณมาทิเส. \\ ตา ปูชิตา ปูชยนฺติ\footnote{ปูชิตา ปูชยนฺติ นํ (ก)}, มานิตา มานยนฺติ นํ.}\csromangathastanza{ตโต นํ อนุกมฺปนฺติ, มาตา ปุตฺตํว โอรสํ. \\ เทวตานุกมฺปิโต โปโส, สทา ภทฺรานิ ปสฺสตี”ติ.}}

\prose{อถ โข ภควา สุนิธวสฺสกาเร มคธมหามตฺเต อิมาหิ คาถาหิ อนุโมทิตฺวา อุฏฺฐายาสนา ปกฺกามิ.}

\proseitem{154}{\csromanfolio{76}เตน โข ปน สมเยน สุนิธวสฺสการา มคธมหามตฺตา ภควนฺตํ ปิฏฺฐิโต ปิฏฺฐิโต อนุพนฺธา โหนฺติ “เยนชฺช สมโณ โคตโม ทฺวาเรน นิกฺขมิสฺสติ, ตํ โคตมทฺวารํ นาม ภวิสฺสติ.\csromanspacer{} เยน ติตฺเถน คงฺคํ นทิํ ตริสฺสติ, ตํ โคตมติตฺถํ นาม ภวิสฺสตี”ติ.\csromanspacer{} อถ โข ภควา เยน ทฺวาเรน นิกฺขมิ, ตํ โคตมทฺวารํ นาม อโหสิ.\csromanspacer{} อถ โข ภควา เยน คงฺคา นที เตนุปสงฺกมิ.\csromanspacer{} เตน โข ปน สมเยน คงฺคา นที ปูรา โหติ สมติตฺติกา กากเปยฺยา.\csromanspacer{} อปฺเปกจฺเจ มนุสฺสา นาวํ ปริเยสนฺติ, อปฺเปกจฺเจ อุฬุมฺปํ ปริเยสนฺติ, อปฺเปกจฺเจ กุลฺลํ พนฺธนฺติ อปารา\footnote{ปารา (สี, สฺยา, ก), โอรา (วิ-มหาวคฺค)} ปารํ คนฺตุกามา.\csromanspacer{} อถ โข ภควา เสยฺยถาปิ นาม พลวา ปุริโส สมิญฺชิตํ วา พาหํ ปสาเรยฺย, ปสาริตํ วา พาหํ สมิญฺเชยฺย, เอวเมว คงฺคาย นทิยา โอริมตีเร อนฺตรหิโต ปาริมตีเร ปจฺจุฏฺฐาสิ สทฺธิํ ภิกฺขุสํเฆน.\csromanspacer{} อทฺทสา โข ภควา เต มนุสฺเส อปฺเปกจฺเจ นาวํ ปริเยสนฺเต อปฺเปกจฺเจ อุฬุมฺปํ ปริเยสนฺเต อปฺเปกจฺเจ กุลฺลํ พนฺธนฺเต อปารา ปารํ คนฺตุกาเม.\csromanspacer{} อถ โข ภควา เอตมตฺถํ วิทิตฺวา ตายํ เวลายํ อิมํ อุทานํ อุทาเนสิ\csromandash{}}

\csromangathameasure{“เย ตรนฺติ อณฺณวํ สรํ, เสตุํ กตฺวาน วิสชฺช ปลฺลลานิ. \\ กุลฺลํ หิ ชโน พนฺธติ, ติณฺณา เมธาวิโน ชนา”ติ.}
\csromangathagroup{\csromangathastanza{“เย ตรนฺติ อณฺณวํ สรํ, เสตุํ กตฺวาน วิสชฺช ปลฺลลานิ. \\ กุลฺลํ หิ ชโน พนฺธติ\footnote{กุลฺลํ ชโน จ พนฺธติ (สฺยา), กุลฺลํ หิ ชโน ปพนฺธหิ (สี, อิ, ก)}, ติณฺณา\footnote{นิติณฺณา, น ติณฺณา (ก)} เมธาวิโน ชนา”ติ.}}

\csromanheader{2}{ปฐมภาณวาโร.}
\csromansectionrule
\csromanmatikaanchor{mtk.34}
\csromantocmarkref{section}{อริยสจฺจกถา}{mtk.34}
\bookmark[dest={mtk.34},level=1]{อริยสจฺจกถา}
\csromanheader{1}{\textbf{อริยสจฺจกถา}}
\proseitem{155}{อถ โข ภควา อายสฺมนฺตํ อานนฺทํ อามนฺเตสิ “อายามานนฺท, เยน โกฏิคาโม เตนุปสงฺกมิสฺสามา”ติ. “เอวํ ภนฺเต”ติ โข อายสฺมา อานนฺโท ภควโต ปจฺจสฺโสสิ.\csromanspacer{} อถ โข ภควา มหตา ภิกฺขุสํเฆน สทฺธิํ เยน โกฏิคาโม ตทวสริ, ตตฺร สุทํ ภควา โกฏิคาเม วิหรติ.\csromanspacer{} ตตฺร โข ภควา ภิกฺขู อามนฺเตสิ\csromandash{}}

\csromanheader{2}{3. มหาปรินิพฺพานสุตฺต}
\prose{\csromanfolio{77}“จตุนฺนํ ภิกฺขเว อริยสจฺจานํ อนนุโพธา อปฺปฏิเวธา เอวมิทํ ทีฆมทฺธานํ สนฺธาวิตํ สํสริตํ มมญฺเจว ตุมฺหากญฺจ.\csromanspacer{} กตเมสํ จตุนฺนํ.\csromanspacer{} ทุกฺขสฺส ภิกฺขเว อริยสจฺจสฺส อนนุโพธา อปฺปฏิเวธา เอวมิทํ ทีฆมทฺธานํ สนฺธาวิตํ สํสริตํ มมญฺเจว ตุมฺหากญฺจ, ทุกฺขสมุทยสฺส ภิกฺขเว อริยสจฺจสฺส อนนุโพธา อปฺปฏิเวธา เอวมิทํ ทีฆมทฺธานํ สนฺธาวิตํ สํสริตํ มมญฺเจว ตุมฺหากญฺจ, ทุกฺขนิโรธสฺส ภิกฺขเว อริยสจฺจสฺส อนนุโพธา อปฺปฏิเวธา เอวมิทํ ทีฆมทฺธานํ สนฺธาวิตํ สํสริตํ มมญฺเจว ตุมฺหากญฺจ, ทุกฺขนิโรธคามินิยา ปฏิปทาย ภิกฺขเว อริยสจฺจสฺส อนนุโพธา อปฺปฏิเวธา เอวมิทํ ทีฆมทฺธานํ สนฺธาวิตํ สํสริตํ มมญฺเจว ตุมฺหากญฺจ.\csromanspacer{} ตยิทํ ภิกฺขเว ทุกฺขํ อริยสจฺจํ อนุพุทฺธํ ปฏิวิทฺธํ, ทุกฺขสมุทยํ\footnote{ทุกฺขสมุทโย (สฺยา)} อริยสจฺจํ อนุพุทฺธํ ปฏิวิทฺธํ, ทุกฺขนิโรธํ\footnote{ทุกฺขนิโรโธ (สฺยา)} อริยสจฺจํ อนุพุทฺธํ ปฏิวิทฺธํ, ทุกฺขนิโรธคามินี ปฏิปทา อริยสจฺจํ อนุพุทฺธํ ปฏิวิทฺธํ, อุจฺฉินฺนา ภวตณฺหา, ขีณา ภวเนตฺติ, นตฺถิทานิ ปุนพฺภโว”ติ, อิทมโวจ ภควา, อิทํ วตฺวาน สุคโต อถาปรํ เอตทโวจ สตฺถา\csromandash{}}

\csromangathameasure{“จตุนฺนํ อริยสจฺจานํ, ยถาภูตํ อทสฺสนา. \\ สํสิตํ ทีฆมทฺธานํ, ตาสุ ตาเสฺวว ชาติสุ. \\ ตานิ เอตานิ ทิฏฺฐานิ, ภวเนตฺติ สมูหตา. \\ อุจฺฉินฺนํ มูลํ ทุกฺขสฺส, นตฺถิทานิ ปุนพฺภโว”ติ.}
\csromangathagroup{\csromangathastanza{“จตุนฺนํ อริยสจฺจานํ, ยถาภูตํ อทสฺสนา. \\ สํสิตํ ทีฆมทฺธานํ, ตาสุ ตาเสฺวว ชาติสุ.}\csromangathastanza{ตานิ เอตานิ ทิฏฺฐานิ, ภวเนตฺติ สมูหตา. \\ อุจฺฉินฺนํ มูลํ ทุกฺขสฺส, นตฺถิทานิ ปุนพฺภโว”ติ.}}

\prose{ตตฺรปิ สุทํ ภควา โกฏิคาเม วิหรนฺโต เอตเทว พหุลํ ภิกฺขูนํ ธมฺมิํ กถํ กโรติ “อิติ สีลํ อิติ สมาธิ อิติ ปญฺญา.\csromanspacer{} สีลปริภาวิโต สมาธิ มหปฺผโล โหติ มหานิสํโส.\csromanspacer{} สมาธิปริภาวิตา ปญฺญา มหปฺผลา โหติ มหานิสํสา.\csromanspacer{} ปญฺญาปริภาวิ-ตํ จิตฺตํ สมฺมเทว อาสเวหิ วิมุจฺจติ.\csromanspacer{} เสยฺยถิทํ, กามาสวา ภวาสวา อวิชฺชาสวา”ติ.}

\csromanmatikaanchor{mtk.35}
\csromantocmarkref{section}{อนาวตฺติธมฺมสมฺโพธิปรายณ}{mtk.35}
\bookmark[dest={mtk.35},level=1]{อนาวตฺติธมฺมสมฺโพธิปรายณ}
\csromanheader{1}{\textbf{อนาวตฺติธมฺมสมฺโพธิปรายณ}}
\proseitem{156}{อถ โข ภควา โกฏิคาเม ยถาภิรนฺตํ วิหริตฺวา อายสฺมนฺตํ อานนฺทํ อามนฺเตสิ “อายามานนฺท, เยน นาติกา\footnote{นาทิกา (สฺยา, อิ)} เตนุปสงฺกมิสฺสามา”ติ. “เอวํ ภนฺเต”ติ โข อายสฺมา อานนฺโท \csromanfolio{78}ภควโต ปจฺจสฺโสสิ.\csromanspacer{} อถ โข ภควา มหตา ภิกฺขุสํเฆน สทฺธิํ เยน นาติกา ตทวสริ, ตตฺรปิ สุทํ ภควา นาติเก วิหรติ คิญฺชกาวสเถ.\csromanspacer{} อถ โข อายสฺมา อานนฺโท เยน ภควา เตนุปสงฺกมิ, อุปสงฺกมิตฺวา ภควนฺตํ อภิวาเทตฺวา เอกมนฺตํ นิสีทิ, เอกมนฺตํ นิสินฺโน โข อายสฺมา อานนฺโท ภควนฺตํ เอตทโวจ “สาโฬฺห นาม ภนฺเต ภิกฺขุ นาติเก กาลงฺกโต, ตสฺส กา คติ โก อภิสมฺปราโย.\csromanspacer{} นนฺทา นาม ภนฺเต ภิกฺขุนี นาติเก กาลงฺกตา, ตสฺสา กา คติ โก อภิสมฺปราโย.\csromanspacer{} สุทตฺโต นาม ภนฺเต อุปาสโก นาติเก กาลงฺกโต, ตสฺส กา คติ โก อภิสมฺปราโย.\csromanspacer{} สุชาตา นาม ภนฺเต อุปาสิกา นาติเก กาลงฺกตา, ตสฺสา กา คติ โก อภิสมฺปราโย.\csromanspacer{} กุกฺกุโฏ\footnote{กกุโธ (สฺยา)} นาม ภนฺเต อุปาสโก นาติเก กาลงฺกโต, ตสฺส กา คติ โก อภิสมฺปราโย.\csromanspacer{} กาฬิมฺโพ\footnote{กาลิงฺโค (อิ), การฬิมฺโพ (สฺยา)} นาม ภนฺเต อุปาสโก ฯเปฯ.\csromanspacer{} นิกโฏ นาม ภนฺเต อุปาสโก.\csromanspacer{} กฏิสฺสโห\footnote{กฏิสฺสโภ (สี, อิ)} นาม ภนฺเต อุปาสโก.\csromanspacer{} ตุฏฺโฐ นาม ภนฺเต อุปาสโก.\csromanspacer{} สนฺตุฏฺโฐ นาม ภนฺเต อุปาสโก.\csromanspacer{} ภทฺโท\footnote{ภโฏ (สฺยา)} นาม ภนฺเต อุปาสโก.\csromanspacer{} สุภทฺโท\footnote{สุภโฏ (สฺยา)} นาม ภนฺเต อุปาสโก นาติเก กาลงฺกโต, ตสฺส กา คติ โก อภิสมฺปราโย”ติ.}

\proseitem{157}{สาโฬฺห อานนฺท ภิกฺขุ อาสวานํ ขยา อนาสวํ เจโตวิมุตฺติํ ปญฺญาวิมุตฺติํ ทิฏฺเฐว ธมฺเม สยํ อภิญฺญา สจฺฉิกตฺวา อุปสมฺปชฺช วิหาสิ.\csromanspacer{} นนฺทา อานนฺท ภิกฺขุนี ปญฺจนฺนํ โอรมฺภาคิยานํ สํโยชนานํ ปริกฺขยา โอปปาติกา ตตฺถ ปรินิพฺพายินี อนาวตฺติธมฺมา ตสฺมา โลกา.\csromanspacer{} สุทตฺโต อานนฺท อุปาสโก ติณฺณํ สํโยชนานํ ปริกฺขยา ราคโทสโมหานํ ตนุตฺตา สกทาคามี สกิเทว อิมํ โลกํ อาคนฺตฺวา ทุกฺขสฺสนฺตํ กริสฺสติ.\csromanspacer{} สุชาตา อานนฺท อุปาสิกา ติณฺณํ สํโยชนานํ ปริกฺขยา โสตาปนฺนา อวินิปาตธมฺมา นิยตา สมฺโพธิปรายณา\footnote{ปรายนา (สี, สฺยา, อิ, ก)}.\csromanspacer{} กุกฺกุโฏ อานนฺท อุปาสโก ปญฺจนฺนํ โอรมฺภาคิยานํ สํโยชนานํ ปริกฺขยา โอปปาติโก ตตฺถ ปรินิพฺพายี อนาวตฺติธมฺโม ตสฺมา โลกา.\csromanspacer{} กาฬิมฺโพ อานนฺท อุปาสโก ฯเปฯ.\csromanspacer{} นิกโฏ อานนฺท อุปาสโก.}

\vspace{\tipitakaparskip}
\csromancenter{มหาปรินิพฺพานสุตฺต กฏิสฺสโห อานนฺท อุปาสโก.\csromanspacer{} ตุฏฺโฐ อานนฺท อุปาสโก.\csromanspacer{} สนฺตุฏฺโฐ อานนฺท อุปาสโก.\csromanspacer{} ภทฺโท อานนฺท อุปาสโก.\csromanspacer{} สุภทฺโท อานนฺท อุปาสโก ปญฺจนฺนํ โอรมฺภาคิยานํ สํโยชนานํ ปริกฺขยา โอปปาติโก ตตฺถ ปรินิพฺพายี อนาวตฺติธมฺโม ตสฺมา โลกา.\csromanspacer{} ปโรปญฺญาสํ อานนฺท นาติเก อุปาสกา กาลงฺกตา ปญฺจนฺนํ โอรมฺภาคิยานํ สํโยชนานํ ปริกฺขยา โอปปาติกา ตตฺถ ปรินิพฺพายิโน อนาวตฺติธมฺมา ตสฺมา โลกา.\csromanspacer{} สาธิกา นวุติ\footnote{ฉาธิกา นวุติ (สฺยา)} อานนฺท นาติเก อุปาสกา กาลงฺกตา ติณฺณํ สํโยชนานํ ปริกฺขยา ราคโทสโมหานํ ตนุตฺตา สกทาคามิโน สกิเทว อิมํ โลกํ อาคนฺตฺวา ทุกฺขสฺสนฺตํ กริสฺสนฺติ.\csromanspacer{} สาติเรกานิ\footnote{ทสาติเรกานิ (สฺยา)} อานนฺท ปญฺจสตานิ นาติเก อุปาสกา กาลงฺกตา ติณฺณํ สํโยชนานํ ปริกฺขยา โสตาปนฺนา อวินิปาตธมฺมา นิยตา สมฺโพธิปรายณา.}
\csromanmatikaanchor{mtk.36}
\csromantocmarkref{section}{ธมฺมาทาสธมฺมปริยาย}{mtk.36}
\bookmark[dest={mtk.36},level=1]{ธมฺมาทาสธมฺมปริยาย}
\csromanheader{1}{\textbf{ธมฺมาทาสธมฺมปริยาย}}
\proseitem{158}{\csromanfolio{79}อนจฺฉริยํ โข ปเนตํ อานนฺท, ยํ มนุสฺสภูโต กาลํ กเรยฺย, ตสฺมิํเยว\footnote{ตสฺมิํ ตสฺมิํ เจ (สี, อิ), ตสฺมิํ ตสฺมิํ โข (สฺยา)} กาลงฺกเต ตถาคตํ อุปสงฺกมิตฺวา เอตมตฺถํ ปุจฺฉิสฺสถ, วิเหสา เหสา อานนฺท ตถาคตสฺส.\csromanspacer{} ตสฺมาติหานนฺท ธมฺมาทาสํ นาม ธมฺมปริยายํ เทเสสฺสามิ, เยน สมนฺนาคโต อริยสาวโก อากงฺขมาโน อตฺตนาว อตฺตานํ พฺยากเรยฺย “ขีณนิรโยมฺหิ ขีณติรจฺฉานโยนิ ขีณเปตฺติวิสโย ขีณาปายทุคฺคติวินิปาโต, โสตาปนฺโนหมสฺมิ อวินิปาตธมฺโม นิยโต สมฺโพธิปรายโณ”ติ.}

\proseitem{159}{กตโม จ โส อานนฺท ธมฺมาทาโส ธมฺมปริยาโย, เยน สมนฺนาคโต อริยสาวโก อากงฺขมาโน อตฺตนาว อตฺตานํ พฺยากเรยฺย “ขีณนิรโยมฺหิ ขีณติรจฺฉานโยนิ ขีณเปตฺติวิสโย ขีณาปายทุคฺคติวินิปาโต, โสตาปนฺโนหมสฺมิ อวินิปาตธมฺโม นิยโต สมฺโพธิปรายโณ”ติ.}

\prose{\csromanfolio{80}อิธานนฺท อริยสาวโก พุทฺเธ อเวจฺจปฺปสาเทน สมนฺนาคโต โหติ “อิติปิ โส ภควา อรหํ สมฺมาสมฺพุทฺโธ วิชฺชาจรณสมฺปนฺโน สุคโต โลกวิทู อนุตฺตโร ปุริสทมฺมสารถิ สตฺถา เทวมนุสฺสานํ พุทฺโธ ภควา”ติ.}

\prose{ธมฺเม อเวจฺจปฺปสาเทน สมนฺนาคโต โหติ “สฺวากฺขาโต ภควตา ธมฺโม สนฺทิฏฺฐิโก อกาลิโก เอหิปสฺสิโก โอปเนยฺยิโก ปจฺจตฺตํ เวทิตพฺโพ วิญฺญูหี”ติ.}

\prose{สํเฆ อเวจฺจปฺปสาเทน สมนฺนาคโต โหติ “สุปฺปฏิปนฺโน ภควโต สาวกสํโฆ, อุชุปฺปฏิปนฺโน ภควโต สาวกสํโฆ, ญายปฺปฏิปนฺโน ภควโต สาวกสํโฆ, สามีจิปฺปฏิปนฺโน ภควโต สาวกสํโฆ, ยทิทํ จตฺตาริ ปุริสยุคานิ อฏฺฐ ปุริสปุคฺคลา, เอส ภควโต สาวกสํโฆ อาหุเนยฺโย ปาหุเนยฺโย ทกฺขิเณยฺโย อญฺชลิกรณีโย อนุตฺตรํ ปุญฺญกฺเขตฺตํ โลกสฺสา”ติ.}

\prose{อริยกนฺเตหิ สีเลหิ สมนฺนาคโต โหติ อขณฺเฑหิ อจฺฉิทฺเทหิ อสพเลหิ อกมฺมาเสหิ ภุชิสฺเสหิ วิญฺญูปสตฺเถหิ อปรามฏฺเฐหิ สมาธิสํวตฺตนิเกหิ.}

\prose{อยํ โข โส อานนฺท ธมฺมาทาโส ธมฺมปริยาโย, เยน สมนฺนาคโต อริยสาวโก อากงฺขมาโน อตฺตนาว อตฺตานํ พฺยากเรยฺย “ขีณนิรโยมฺหิ ขีณติรจฺฉานโยนิ ขีณเปตฺติวิสโย ขีณาปายทุคฺคติวินิปาโต, โสตาปนฺโนหมสฺมิ อวินิปาตธมฺโม นิยโต สมฺโพธิปรายโณ”ติ.}

\prose{ตตฺรปิ สุทํ ภควา นาติเก วิหรนฺโต คิญฺชกาวสเถ เอตเทว พหุลํ ภิกฺขูนํ ธมฺมิํ กถํ กโรติ “อิติ สีลํ อิติ สมาธิ อิติ ปญฺญา.\csromanspacer{} สีลปริภาวิโต สมาธิ มหปฺผโล โหติ มหานิสํโส.\csromanspacer{} สมาธิปริภาวิตา ปญฺญา มหปฺผลา โหติ มหานิสํสา.\csromanspacer{} ปญฺญาปริภาวิตํ จิตฺตํ สมฺมเทว อาสเวหิ วิมุจฺจติ.\csromanspacer{} เสยฺยถิทํ, กามาสวา ภวาสวา อวิชฺชาสวา”ติ.}

\csromanheader{2}{3. มหาปรินิพฺพานสุตฺต}
\proseitem{160}{\csromanfolio{81}อถ โข ภควา นาติเก ยถาภิรนฺตํ วิหริตฺวา อายสฺมนฺตํ อานนฺทํ อามนฺเตสิ “อายามานนฺท, เยน เวสาลี เตนุปสงฺกมิสฺสามา”ติ. “เอวํ ภนฺเต”ติ โข อายสฺมา อานนฺโท ภควโต ปจฺจสฺโสสิ.\csromanspacer{} อถ โข ภควา มหตา ภิกฺขุสํเฆน สทฺธิํ เยน เวสาลี ตทวสริ, ตตฺร สุทํ ภควา เวสาลิยํ วิหรติ อมฺพปาลิวเน.\csromanspacer{} ตตฺร โข ภควา ภิกฺขู อามนฺเตสิ\csromandash{}}

\prose{“สโต ภิกฺขเว ภิกฺขุ วิหเรยฺย สมฺปชาโน, อยํ โว อมฺหากํ อนุสาสนี.\csromanspacer{} กถญฺจ ภิกฺขเว ภิกฺขุ สโต โหติ.\csromanspacer{} อิธ ภิกฺขเว ภิกฺขุ กาเย กายานุปสฺสี วิหรติ อาตาปี สมฺปชาโน สติมา วิเนยฺย โลเก อภิชฺฌาโทมนสฺสํ.\csromanspacer{} เวทนาสุ เวทนานุปสฺสี ฯเปฯ.\csromanspacer{} จิตฺเต จิตฺตานุปสฺสี ฯเปฯ.\csromanspacer{} ธมฺเมสุ ธมฺมานุปสฺสี วิหรติ อาตาปี สมฺปชาโน สติมา วิเนยฺย โลเก อภิชฺฌาโทมนสฺสํ.\csromanspacer{} เอวํ โข ภิกฺขเว ภิกฺขุ สโต โหติ.}

\prose{กถญฺจ ภิกฺขเว ภิกฺขุ สมฺปชาโน โหติ.\csromanspacer{} อิธ ภิกฺขเว ภิกฺขุ อภิกฺกนฺเต ปฏิกฺกนฺเต สมฺปชานการี โหติ, อาโลกิเต วิโลกิเต สมฺปชานการี โหติ, สมิญฺชิเต ปสาริเต สมฺปชานการี โหติ, สํฆาฏิปตฺตจีวรธารเณ สมฺปชานการี โหติ, อสิเต ปีเต ขายิเต สายิเต สมฺปชานการี โหติ, อุจฺจารปสฺสาวกมฺเม สมฺปชานการี โหติ, คเต ฐิเต นิสินฺเน สุตฺเต ชาคริเต ภาสิเต ตุณฺหีภาเว สมฺปชานการี โหติ.\csromanspacer{} เอวํ โข ภิกฺขเว ภิกฺขุ สมฺปชาโน โหติ.\csromanspacer{} สโต ภิกฺขเว ภิกฺขุ วิหเรยฺย สมฺปชาโน, อยํ โว อมฺหากํ อนุสาสนี”ติ.}

\csromanmatikaanchor{mtk.37}
\csromantocmarkref{section}{อมฺพปาลีคณิกา}{mtk.37}
\bookmark[dest={mtk.37},level=1]{อมฺพปาลีคณิกา}
\csromanheader{1}{\textbf{อมฺพปาลีคณิกา}}
\proseitem{161}{อสฺโสสิ โข \textbf{อมฺพปาลี คณิกา} “ภควา กิร เวสาลิํ อนุปฺปตฺโต เวสาลิยํ วิหรติ มยฺหํ อมฺพวเน”ติ.\csromanspacer{} อถ โข \textbf{อมฺพปาลี คณิกา} ภทฺทานิ ภทฺทานิ ยานานิ โยชาเปตฺวา ภทฺทํ ภทฺทํ ยานํ อภิรุหิตฺวา ภทฺเทหิ ภทฺเทหิ ยาเนหิ เวสาลิยา นิยฺยาสิ.\csromanspacer{} เยน สโก อาราโม เตน ปายาสิ.\csromanspacer{} ยาวติกา ยานสฺส ภูมิ, ยาเนน คนฺตฺวา ยานา ปจฺโจโรหิตฺวา ปตฺติกาว เยน ภควา เตนุปสงฺกมิ, อุปสงฺกมิตฺวา ภควนฺตํ อภิวาเทตฺวา เอกมนฺตํ นิสีทิ. \csromanfolio{82}เอกมนฺตํ นิสินฺนํ โข อมฺพปาลิํ คณิกํ ภควา ธมฺมิยา กถาย สนฺทสฺเสสิ สมาทเปสิ สมุตฺเตเชสิ สมฺปหํเสสิ.\csromanspacer{} อถ โข อมฺพปาลี คณิกา ภควตา ธมฺมิยา กถาย สนฺทสฺสิตา สมาทปิตา สมุตฺเตชิตา สมฺปหํสิตา ภควนฺตํ เอตทโวจ “อธิวาเสตุ เม ภนฺเต ภควา สฺวาตนาย ภตฺตํ สทฺธิํ ภิกฺขุสํเฆนา”ติ.\csromanspacer{} อธิวาเสสิ ภควา ตุณฺหีภาเวน.\csromanspacer{} อถ โข อมฺพปาลี คณิกา ภควโต อธิวาสนํ วิทิตฺวา อุฏฺฐายาสนา ภควนฺตํ อภิวาเทตฺวา ปทกฺขิณํ กตฺวา ปกฺกามิ.}

\prose{อสฺโสสุํ โข เวสาลิกา ลิจฺฉวี “ภควา กิร เวสาลิํ อนุปฺปตฺโต เวสาลิยํ วิหรติ อมฺพปาลิวเน”ติ.\csromanspacer{} อถ โข เต ลิจฺฉวี ภทฺทานิ ภทฺทานิ ยานานิ โยชาเปตฺวา ภทฺทํ ภทฺทํ ยานํ อภิรุหิตฺวา ภทฺเทหิ ภทฺเทหิ ยาเนหิ เวสาลิยา นิยฺยิํสุ.\csromanspacer{} ตตฺร เอกจฺเจ ลิจฺฉวี นีลา โหนฺติ นีลวณฺณา นีลวตฺถา นีลาลงฺการา, เอกจฺเจ ลิจฺฉวี ปีตา โหนฺติ ปีตวณฺณา ปีตวตฺถา ปิตาลงฺการา, เอกจฺเจ ลิจฺฉวี โลหิตา โหนฺติ โลหิตวณฺณา โลหิตวตฺถา โลหิตาลงฺการา, เอกจฺเจ ลิจฺฉวี โอทาตา โหนฺติ โอทาตวณฺณา โอทาตวตฺถา โอทาตาลงฺการา.\csromanspacer{} อถ โข อมฺพปาลี คณิกา ทหรานํ ทหรานํ ลิจฺฉวีนํ อกฺเขน อกฺขํ จกฺเกน จกฺกํ ยุเคน ยุคํ ปฏิวฏฺเฏสิ\footnote{ปริวตฺเตสิ (วิ-มหาวคฺค)}.\csromanspacer{} อถ โข เต ลิจฺฉวี อมฺพปาลิํ คณิกํ เอตทโวจุํ “กิํ เช อมฺพปาลิ ทหรานํ ทหรานํ ลิจฺฉวีนํ อกฺเขน อกฺขํ จกฺเกน จกฺกํ ยุเคน ยุคํ ปฏิวฏฺเฏสี”ติ.\csromanspacer{} ตถา หิ ปน เม อยฺยปุตฺตา ภควา นิมนฺติโต สฺวาตนาย ภตฺตํ สทฺธิํ ภิกฺขุสํเฆนาติ.\csromanspacer{} เทหิ เช อมฺพปาลิ เอตํ\footnote{เอกํ (ก)} ภตฺตํ สตสหสฺเสนาติ.\csromanspacer{} สเจปิ เม อยฺยปุตฺตา เวสาลิํ สาหารํ ทสฺสถ\footnote{ทชฺเชยฺยาถ (วิ-มหาวคฺค)}, เอวมหํ ตํ\footnote{เอวมฺปิ มหนฺตํ (สฺยา), เอวํ มหนฺตํ (สี, อิ)} ภตฺตํ น ทสฺสามีติ\footnote{เนว ทชฺชาหํ ตํ ภตฺตนฺติ (วิ-มหาวคฺค)}.\csromanspacer{} อถ โข เต ลิจฺฉวี องฺคุลิํ โผเฏสุํ “ชิตมฺห\footnote{ชิตมฺหา (พหูสุ)} วต โภ อมฺพกาย ชิตมฺห วต โภ อมฺพกายา”ติ\footnote{“ชิตมฺหา วต โภ อมฺพปาลิกาย วญฺจิตมฺหา วต โภ อมฺพปาลิกายา”ติ (สฺยา)}.}

\prose{อถ โข เต ลิจฺฉวี เยน อมฺพปาลิวนํ เตน ปายิํสุ.\csromanspacer{} อทฺทสา โข ภควา เต ลิจฺฉวี ทูรโตว อาคจฺฉนฺเต, ทิสฺวาน ภิกฺขู อามนฺเตสิ}

\vspace{\tipitakaparskip}
\csromancenter{มหาปรินิพฺพานสุตฺต “เยสํ\footnote{เยหิ (วิ-มหาวคฺค)} ภิกฺขเว ภิกฺขูนํ เทวา ตาวติํสา อทิฏฺฐปุพฺพา, โอโลเกถ ภิกฺขเว ลิจฺฉวิปริสํ, อปโลเกถ ภิกฺขเว ลิจฺฉวิปริสํ, อุปสํหรถ ภิกฺขเว ลิจฺฉวิปริสํ ตาวติํสสทิสนฺ”ติ.\csromanspacer{} อถ โข เต ลิจฺฉวี ยาวติกา ยานสฺส ภูมิ, ยาเนน คนฺตฺวา ยานา ปจฺโจโรหิตฺวา ปตฺติกาว เยน ภควา เตนุปสงฺกมิํสุ, อุปสงฺกมิตฺวา ภควนฺตํ อภิวาเทตฺวา เอกมนฺตํ นิสีทิํสุ.\csromanspacer{} เอกมนฺตํ นิสินฺเน โข เต ลิจฺฉวี ภควา ธมฺมิยา กถาย สนฺทสฺเสสิ สมาทเปสิ สมุตฺเตเชสิ สมฺปหํเสสิ.\csromanspacer{} อถ โข เต ลิจฺฉวี ภควตา ธมฺมิยา กถาย สนฺทสฺสิตา สมาทปิตา สมุตฺเตชิตา สมฺปหํสิตา ภควนฺตํ เอตทโวจุํ “อธิวาเสตุ โน ภนฺเต ภควา สฺวาตนาย ภตฺตํ สทฺธิํ ภิกฺขุสํเฆนา”ติ.\csromanspacer{} อถ โข ภควา เต ลิจฺฉวี เอตทโวจ “อธิวุตฺถํ\footnote{อธิวาสิตํ (สฺยา)} โข เม ลิจฺฉวี สฺวาตนาย อมฺพปาลิยา คณิกาย ภตฺตนฺ”ติ.\csromanspacer{} อถ โข เต ลิจฺฉวี องฺคุลิํ โผเฏสุํ “ชิตมฺห วต โภ อมฺพกาย ชิตมฺห วต โภ อมฺพกายา”ติ.\csromanspacer{} อถ โข เต ลิจฺฉวี ภควโต ภาสิตํ อภินนฺทิตฺวา อนุโมทิตฺวา อุฏฺฐายาสนา ภควนฺตํ อภิวาเทตฺวา ปทกฺขิณํ กตฺวา ปกฺกมิํสุ.}
\proseitem{162}{\csromanfolio{83}อถ โข อมฺพปาลี คณิกา ตสฺสา รตฺติยา อจฺจเยน สเก อาราเม ปณีตํ ขาทนียํ โภชนียํ ปฏิยาทาเปตฺวา ภควโต กาลํ อาโรจาเปสิ “กาโล ภนฺเต นิฏฺฐิตํ ภตฺตนฺ”ติ.\csromanspacer{} อถ โข ภควา ปุพฺพณฺหสมยํ นิวาเสตฺวา ปตฺตจีวรมาทาย สทฺธิํ ภิกฺขุสํเฆน เยน อมฺพปาลิยา คณิกาย นิเวสนํ เตนุปสงฺกมิ, อุปสงฺกมิตฺวา ปญฺญตฺเต อาสเน นิสีทิ.\csromanspacer{} อถ โข อามฺพปาลี คณิกา พุทฺธปฺปมุขํ ภิกฺขุสํฆํ ปณีเตน ขาทนีเยน โภชนีเยน สหตฺถา สนฺตปฺเปสิ สมฺปวาเรสิ.\csromanspacer{} อถ โข อมฺพปาลี คณิกา ภควนฺตํ ภุตฺตาวิํ โอนีตปตฺตปาณิํ อญฺญตรํ นีจํ อาสนํ คเหตฺวา เอกมนฺตํ นิสีทิ, เอกมนฺตํ นิสินฺนา โข อมฺพปาลี คณิกา ภควนฺตํ เอตทโวจ “อิมาหํ ภนฺเต อารามํ พุทฺธปฺปมุขสฺส ภิกฺขุสํฆสฺส ทมฺมี”ติ.\csromanspacer{} ปฏิคฺคเหสิ ภควา อารามํ.\csromanspacer{} อถ โข ภควา อมฺพปาลิํ คณิกํ ธมฺมิยา กถาย สนฺทสฺเสตฺวา สมาทเปตฺวา สมุตฺเตเชตฺวา สมฺปหํเสตฺวา อุฏฺฐายาสนา ปกฺกามิ.\csromanspacer{} ตตฺรปิ สุทํ ภควา เวสาลิยํ วิหรนฺโต อมฺพปาลิวเน เอตเทว พหุลํ ภิกฺขูนํ ธมฺมิํ กถํ \csromanfolio{84}กโรติ “อิติ สีลํ อิติ สมาธิ อิติ ปญฺญา.\csromanspacer{} สีลปริภาวิโต สมาธิ มหปฺผโล โหติ มหานิสํโส.\csromanspacer{} สมาธิปริภาวิตา ปญฺญา มหปฺผลา โหติ มหานิสํสา.\csromanspacer{} ปญฺญาปริภาวิตํ จิตฺตํ สมฺมเทว อาสเวหิ วิมุจฺจติ.\csromanspacer{} เสยฺยถิทํ, กามาสวา ภวาสวา อวิชฺชาสวา”ติ.}

\csromanmatikaanchor{mtk.38}
\csromantocmarkref{section}{เวฬุวคามวสฺสูปคมน}{mtk.38}
\bookmark[dest={mtk.38},level=1]{เวฬุวคามวสฺสูปคมน}
\csromanheader{1}{\textbf{เวฬุวคามวสฺสูปคมน}}
\proseitem{163}{อถ โข ภควา อมฺพปาลิวเน ยถาภิรนฺตํ วิหริตฺวา อายสฺมนฺตํ อานนฺทํ อามนฺเตสิ “อายามานนฺท, เยน เวฬุวคามโก\footnote{เพฬุว คามโก (สี, อิ)} เตนุปสงฺกมิสฺสามา”ติ. “เอวํ ภนฺเต”ติ โข อายสฺมา อานนฺโท ภควโต ปจฺจสฺโสสิ.\csromanspacer{} อถ โข ภควา มหตา ภิกฺขุสํเฆน สทฺธิํ เยน เวฬุวคามโก ตทวสริ, ตตฺร สุทํ ภควา เวฬุวคามเก วิหรติ.\csromanspacer{} ตตฺร โข ภควา ภิกฺขู อามนฺเตสิ “เอถ ตุเมฺห ภิกฺขเว สมนฺตา เวสาลิํ ยถามิตฺตํ ยถาสนฺทิฏฺฐํ ยถาสมฺภตฺตํ วสฺสํ อุเปถ\footnote{อุปคจฺฉถ (สฺยา)}, อหํ ปน อิเธว เวฬุวคามเก วสฺสํ อุปคจฺฉามี”ติ. “เอวํ ภนฺเต”ติ โข เต ภิกฺขู ภควโต ปฏิสฺสุตฺวา สมนฺตา เวสาลิํ ยถามิตฺตํ ยถาสนฺทิฏฺฐํ ยถาสมฺภตฺตํ วสฺสํ อุปคจฺฉิํสุ.\csromanspacer{} ภควา ปน ตตฺเถว เวฬุวคามเก วสฺสํ อุปคจฺฉิ.}

\proseitem{164}{อถ โข ภควโต วสฺสูปคตสฺส ขโร อาพาโธ อุปฺปชฺชิ, พาฬฺหา เวทนา วตฺตนฺติ มารณนฺติกา.\csromanspacer{} ตา สุทํ ภควา สโต สมฺปชาโน อธิวาเสสิ อวิหญฺญมาโน.\csromanspacer{} อถ โข ภควโต เอตทโหสิ “น โข เมตํ ปติรูปํ, ยฺวาหํ อนามนฺเตตฺวา อุปฏฺฐาเก อนปโลเกตฺวา ภิกฺขุสํฆํ ปรินิพฺพาเยยฺยํ, ยํนูนาหํ อิมํ อาพาธํ วีริเยน ปฏิปณาเมตฺวา ชีวิตสงฺขารํ อธิฏฺฐาย วิหเรยฺยนฺ”ติ.\csromanspacer{} อถ โข ภควา ตํ อาพาธํ วีริเยน ปฏิปณาเมตฺวา ชีวิตสงฺขารํ อธิฏฺฐาย วิหาสิ.\csromanspacer{} อถ โข ภควโต โส อาพาโธ ปฏิปฺปสฺสมฺภิ.\csromanspacer{} อถ โข ภควา คิลานา วุฏฺฐิโต\footnote{คิลานวุฏฺฐิโต (สทฺทนีติ)} อจิรวุฏฺฐิโต เคลญฺญา วิหารา นิกฺขมฺม วิหารปจฺฉายายํ ปญฺญตฺเต อาสเน นิสีทิ.\csromanspacer{} อถ โข อายสฺมา อานนฺโท เยน ภควา เตนุปสงฺกมิ, อุปสงฺกมิตฺวา ภควนฺตํ อภิวาเทตฺวา เอกมนฺตํ นิสีทิ, เอกมนฺตํ นิสินฺโน โข อายสฺมา อานนฺโท}

\vspace{\tipitakaparskip}
\csromancenter{มหาปรินิพฺพานสุตฺต ภควนฺตํ เอตทโวจ “ทิฏฺโฐ เม ภนฺเต ภควโต ผาสุ, ทิฏฺฐํ เม ภนฺเต ภควโต ขมนียํ, อปิ จ เม ภนฺเต มธุรกชาโต วิย กาโย, ทิสาปิ เม น ปกฺขายนฺติ, ธมฺมาปิ มํ น ปฏิภนฺติ ภควโต เคลญฺเญน, อปิ จ เม ภนฺเต อโหสิ กาจิเทว อสฺสาสมตฺตา น ตาว ภควา ปรินิพฺพายิสฺสติ, น ยาว ภควา ภิกฺขุสํฆํ อารพฺภ กิญฺจิเทว อุทาหรตี”ติ.}
\proseitem{165}{\csromanfolio{85}กิํ ปนานนฺท ภิกฺขุสํโฆ มยิ ปจฺจาสีสติ\footnote{ปจฺจาสิํสติ (สี, สฺยา)}, เทสิโต อานนฺท มยา ธมฺโม อนนฺตรํ อพาหิรํ กริตฺวา, นตฺถานนฺท ตถาคตสฺส ธมฺเมสุ อาจริยมุฏฺฐิ.\csromanspacer{} ยสฺส นูน อานนฺท เอวมสฺส “อหํ ภิกฺขุสํฆํ ปริหริสฺสามี”ติ วา “มมุทฺเทสิโก ภิกฺขุสํโฆ”ติ วา, โส นูน อานนฺท ภิกฺขุสํฆํ อารพฺภ กิญฺจิเทว อุทาหเรยฺย.\csromanspacer{} ตถาคตสฺส โข อานนฺท น เอวํ โหติ “อหํ ภิกฺขุสํฆํ ปริหริสฺสามี”ติ วา “มมุทฺเทสิโก ภิกฺขุสํโฆ”ติ วา.\csromanspacer{} สกิํ\footnote{กิํ (สี, อิ)} อานนฺท ตถาคโต ภิกฺขุสํฆํ อารพฺภ กิญฺจิเทว อุทาหริสฺสติ.\csromanspacer{} อหํ โข ปนานนฺท เอตรหิ ชิณฺโณ วุทฺโธ มหลฺลโก อทฺธคโต วโย-อนุปฺปตฺโต, อาสีติโก เม วโย วตฺตติ.\csromanspacer{} เสยฺยถาปิ อานนฺท ชชฺชรสกฏํ เวฐมิสฺสเกน\footnote{เวฬุมิสฺสเกน (สฺยา), เวฆมิสฺสเกน (อิ), เวธมิสฺสเกน, เวขมิสฺสเกน (ก)} ยาเปติ, เอวเมว โข อานนฺท เวฐมิสฺสเกน มญฺเญ ตถาคตสฺส กาโย ยาเปติ.\csromanspacer{} ยสฺมิํ อานนฺท สมเย ตถาคโต สพฺพนิมิตฺตานํ อมนสิการา เอกจฺจานํ เวทนานํ นิโรธา อนิมิตฺตํ เจโตสมาธิํ อุปสมฺปชฺช วิหรติ, ผาสุตโร อานนฺท ตสฺมิํ สมเย ตถาคตสฺส กาโย โหติ.\csromanspacer{} ตสฺมาติหานนฺท อตฺตทีปา วิหรถ อตฺตสรณา อนญฺญสรณา ธมฺมทีปา ธมฺมสรณา อนญฺญสรณา.\csromanspacer{} กถญฺจานนฺท ภิกฺขุ อตฺตทีโป วิหรติ อตฺตสรโณ อนญฺญสรโณ ธมฺมทีโป ธมฺมสรโณ อนญฺญสรโณ.\csromanspacer{} อิธานนฺท ภิกฺขุ กาเย กายานุปสฺสี วิหรติ อาตาปี สมฺปชาโน สติมา วิเนยฺย โลเก อภิชฺฌาโทมนสฺสํ.\csromanspacer{} เวทนาสุ ฯเปฯ.\csromanspacer{} จิตฺเต ฯเปฯ.\csromanspacer{} ธมฺเมสุ ธมฺมานุปสฺสี วิหรติ อาตาปี สมฺปชาโน สติมา วิเนยฺย โลเก อภิชฺฌาโทมนสฺสํ.\csromanspacer{} เอวํ โข อานนฺท ภิกฺขุ อตฺตทีโป วิหรติ อตฺตสรโณ อนญฺญสรโณ ธมฺมทีโป ธมฺมสรโณ \csromanfolio{86}อนญฺญสรโณ.\csromanspacer{} เย หิ เกจิ อานนฺท เอตรหิ วา มม วา อจฺจเยน อตฺตทีปา วิหริสฺสนฺติ อตฺตสรณา อนญฺญสรณา ธมฺมทีปา ธมฺมสรณา อนญฺญสรณา, ตมตคฺเค เม เต อานนฺท ภิกฺขู ภวิสฺสนฺติ เย เกจิ สิกฺขากามาติ.\csromanspacer{} ทุติยภาณวาโร.}

\csromanmatikaanchor{mtk.39}
\csromantocmarkref{section}{นิมิตฺโตภาสกถา}{mtk.39}
\bookmark[dest={mtk.39},level=1]{นิมิตฺโตภาสกถา}
\csromanheader{1}{\textbf{นิมิตฺโตภาสกถา}}
\proseitem{166}{อถ โข ภควา ปุพฺพณฺหสมยํ นิวาเสตฺวา ปตฺตจีวรมาทาย เวสาลิํ ปิณฺฑาย ปาวิสิ.\csromanspacer{} เวสาลิยํ ปิณฺฑาย จริตฺวา ปจฺฉาภตฺตํ ปิณฺฑปาตปฏิกฺกนฺโต อายสฺมนฺตํ อานนฺทํ อามนฺเตสิ “คณฺหาหิ อานนฺท นิสีทนํ, เยน จาปาลํ เจติยํ\footnote{ปาวาลํ เจติยํ (สฺยา)} เตนุปสงฺกมิสฺสาม ทิวา วิหารายา”ติ. “เอวํ ภนฺเต”ติ โข อายสฺมา อานนฺโท ภควโต ปฏิสฺสุตฺวา นิสีทนํ อาทาย ภควนฺตํ ปิฏฺฐิโต ปิฏฺฐิโต อนุพนฺธิ.\csromanspacer{} อถ โข ภควา เยน จาปาลํ เจติยํ เตนุปสงฺกมิ, อุปสงฺกมิตฺวา ปญฺญตฺเต อาสเน นิสีทิ.\csromanspacer{} อายสฺมาปิ โข อานนฺโท ภควนฺตํ อภิวาเทตฺวา เอกมนฺตํ นิสีทิ.}

\proseitem{167}{เอกมนฺตํ นิสินฺนํ โข อายสฺมนฺตํ อานนฺทํ ภควา เอตทโวจ “รมณียา อานนฺท เวสาลี, รมณียํ อุเทนํ เจติยํ, รมณียํ โคตมกํ เจติยํ, รมณียํ สตฺตมฺพํ\footnote{สตฺตมฺพกํ (อิ)} เจติยํ, รมณียํ พหุปุตฺตํ เจติยํ, รมณียํ สารนฺททํ เจติยํ, รมณียํ จาปาลํ เจติยํ.\csromanspacer{} ยสฺส กสฺสจิ อานนฺท จตฺตาโร อิทฺธิปาทา ภาวิตา พหุลีกตา ยานีกตา วตฺถุกตา อนุฏฺฐิตา ปริจิตา สุสมารทฺธา, โส อากงฺขมาโน กปฺปํ วา ติฏฺเฐยฺย กปฺปาวเสสํ วา.\csromanspacer{} ตถาคตสฺส โข อานนฺท จตฺตาโร อิทฺธิปาทา ภาวิตา พหุลีกตา ยานีกตา วตฺถุกตา อนุฏฺฐิตา ปริจิตา สุสมารทฺธา, โส อากงฺขมาโน\footnote{อากงฺขมาโน (?)} อานนฺท ตถาคโต กปฺปํ วา ติฏฺเฐยฺย กปฺปาวเสสํ วา”ติ.\csromanspacer{} เอวมฺปิ โข อายสฺมา อานนฺโท ภควตา โอฬาริเก นิมิตฺเต กยิรมาเน โอฬาริเก โอภาเส กยิรมาเน นาสกฺขิ ปฏิวิชฺฌิตุํ, น ภควนฺตํ ยาจิ “ติฏฺฐตุ ภนฺเต ภควา กปฺปํ, ติฏฺฐตุ สุคโต กปฺปํ พหุชนหิตาย พหุชนสุขาย โลกานุกมฺปาย อตฺถาย หิตาย สุขาย เทวมนุสฺสานนฺ”ติ ยถา ตํ มาเรน ปริยุฏฺฐิตจิตฺโต.}

\vspace{\tipitakaparskip}
\csromancenter{มหาปรินิพฺพานสุตฺต ทุติยมฺปิ โข ภควา ฯเปฯ.\csromanspacer{} ตติยมฺปิ โข ภควา อายสฺมนฺตํ อานนฺทํ อามนฺเตสิ “รมณียา อานนฺท เวสาลี, รมณียํ อุเทนํ เจติยํ, รมณียํ โคตมกํ เจติยํ, รมณียํ สตฺตมฺพํ เจติยํ, รมณียํ พหุปุตฺตํ เจติยํ, รมณียํ สารนฺททํ เจติยํ, รมณียํ จาปาลํ เจติยํ.\csromanspacer{} ยสฺส กสฺสจิ อานนฺท จตฺตาโร อิทฺธิปาทา ภาวิตา พหุลีกตา ยานีกตา วตฺถุกตา อนุฏฺฐิตา ปริจิตา สุสมารทฺธา, โส อากงฺขมาโน กปฺปํ วา ติฏฺเฐยฺย กปฺปาวเสสํ วา.\csromanspacer{} ตถาคตสฺส โข อานนฺท จตฺตาโร อิทฺธิปาทา ภาวิตา พหุลีกตา ยานีกตา วตฺถุกตา อนุฏฺฐิตา ปริจิตา สุสมารทฺธา, โส อากงฺขมาโน อานนฺท ตถาคโต กปฺปํ วา ติฏฺเฐยฺย กปฺปาวเสสํ วา”ติ.\csromanspacer{} เอวมฺปิ โข อายสฺมา อานนฺโท ภควตา โอฬาริเก นิมิตฺเต กยิรมาเน โอฬาริเก โอภาเส กยิรมาเน นาสกฺขิ ปฏิวิชฺฌิตุํ, น ภควนฺตํ ยาจิ “ติฏฺฐตุ ภนฺเต ภควา กปฺปํ, ติฏฺฐตุ สุคโต กปฺปํ พหุชนหิตาย พหุชนสุขาย โลกานุกมฺปาย อตฺถาย หิตาย สุขาย เทวมนุสฺสานนฺ”ติ ยถา ตํ มาเรน ปริยุฏฺฐิตจิตฺโต.\csromanspacer{} อถ โข ภควา อายสฺมนฺตํ อานนฺทํ อามนฺเตสิ “คจฺฉ ตฺวํ อานนฺท, ยสฺสทานิ กาลํ มญฺญสี”ติ. “เอวํ ภนฺเต”ติ โข อายสฺมา อานนฺโท ภควโต ปฏิสฺสุตฺวา อุฏฺฐายาสนา ภควนฺตํ อภิวาเทตฺวา ปทกฺขิณํ กตฺวา อวิทูเร อญฺญตรสฺมิํ รุกฺขมูเล นิสีทิ.}
\csromanmatikaanchor{mtk.40}
\csromantocmarkref{section}{มารยาจนกถา}{mtk.40}
\bookmark[dest={mtk.40},level=1]{มารยาจนกถา}
\csromanheader{1}{\textbf{มารยาจนกถา}}
\proseitem{168}{\csromanfolio{87}อถ โข มาโร ปาปิมา อจิรปกฺกนฺเต อายสฺมนฺเต อานนฺเท เยน ภควา เตนุปสงฺกมิ, อุปสงฺกมิตฺวา เอกมนฺตํ อฏฺฐาสิ, เอกมนฺตํ ฐิโต โข มาโร ปาปิมา ภควนฺตํ เอตทโวจ “ปรินิพฺพาตุทานิ ภนฺเต ภควา, ปรินิพฺพาตุ สุคโต, ปรินิพฺพานกาโลทานิ ภนฺเต ภควโต, ภาสิตา โข ปเนสา ภนฺเต ภควตา วาจา ‘น ตาวาหํ ปาปิม ปรินิพฺพายิสฺสามิ, ยาว เม ภิกฺขู น สาวกา ภวิสฺสนฺติ วิยตฺตา วินีตา วิสารทา พหุสฺสุตา ธมฺมธรา ธมฺมานุธมฺมปฺปฏิปนฺนา สามีจิปฺปฏิปนฺนา อนุธมฺมจาริโน, สกํ อาจริยกํ อุคฺคเหตฺวา อาจิกฺขิสฺสนฺติ เทเสสฺสนฺติ ปญฺญเปสฺสนฺติ ปฏฺฐเปสฺสนฺติ วิวริสฺสนฺติ วิภชิสฺสนฺติ อุตฺตานี\footnote{อุตฺตานิํ (ก), อุตฺตานิ (สี, อิ)} กริสฺสนฺติ, อุปฺปนฺนํ ปรปฺปวาทํ สหธมฺเมน สุนิคฺคหิตํ นิคฺคเหตฺวา สปฺปาฏิหาริยํ ธมฺมํ \csromanfolio{88}เทเสสฺสนฺตี’ติ.\csromanspacer{} เอตรหิ โข ปน ภนฺเต ภิกฺขู ภควโต สาวกา วิยตฺตา วินีตา วิสารทา พหุสฺสุตา ธมฺมธรา ธมฺมานุธมฺมปฺปฏิปนฺนา สามีจิปฺปฏิปนฺนา อนุธมฺมจาริโน, สกํ อาจริยกํ อุคฺคเหตฺวา อาจิกฺขนฺติ เทเสนฺติ ปญฺญเปนฺติ ปฏฺฐเปนฺติ วิวรนฺติ วิภชนฺติ อุตฺตานี กโรนฺติ, อุปฺปนฺนํ ปรปฺปวาทํ สหธมฺเมน สุนิคฺคหิตํ นิคฺคเหตฺวา สปฺปาฏิหาริยํ ธมฺมํ เทเสนฺติ.\csromanspacer{} ปรินิพฺพาตุทานิ ภนฺเต ภควา, ปรินิพฺพาตุ สุคโต, ปรินิพฺพานกาโลทานิ ภนฺเต ภควโต.}

\prose{ภาสิตา โข ปเนสา ภนฺเต ภควตา วาจา ‘น ตาวาหํ ปาปิม ปรินิพฺพายิสฺสามิ, ยาว เม ภิกฺขุนิโย น สาวิกา ภวิสฺสนฺติ วิยตฺตา วินีตา วิสารทา พหุสฺสุตา ธมฺมธรา ธมฺมานุธมฺมปฺปฏิปนฺนา สามีจิปฺปฏิปนฺนา อนุธมฺมจารินิโย, สกํ อาจริยกํ อุคฺคเหตฺวา อาจิกฺขิสฺสนฺติ เทเสสฺสนฺติ ปญฺญเปสฺสนฺติ ปฏฺฐเปสฺสนฺติ วิวริสฺสนฺติ วิภชิสฺสนฺติ อุตฺตานี กริสฺสนฺติ, อุปฺปนฺนํ ปรปฺปวาทํ สหธมฺเมน สุนิคฺคหิตํ นิคฺคเหตฺวา สปฺปาฏิหาริยํ ธมฺมํ เทเสสฺสนฺตี’ติ.\csromanspacer{} เอตรหิ โข ปน ภนฺเต ภิกฺขุนิโย ภควโต สาวิกา วิยตฺตา วินีตา วิสารทา พหุสฺสุตา ธมฺมธรา ธมฺมานุธมฺมปฺปฏิปนฺนา สามีจิปฺปฏิปนฺนา อนุธมฺมจารินิโย, สกํ อาจริยกํ อุคฺคเหตฺวา อาจิกฺขนฺติ เทเสนฺติ ปญฺญเปนฺติ ปฏฺฐเปนฺติ วิวรนฺติ วิภชนฺติ อุตฺตานี กโรนฺติ, อุปฺปนฺนํ ปรปฺปวาทํ สหธมฺเมน สุนิคฺคหิตํ นิคฺคเหตฺวา สปฺปาฏิหาริยํ ธมฺมํ เทเสนฺติ.\csromanspacer{} ปรินิพฺพาตุทานิ ภนฺเต ภควา, ปรินิพฺพาตุ สุคโต, ปรินิพฺพานกาโลทานิ ภนฺเต ภควโต.}

\prose{ภาสิตา โข ปเนสา ภนฺเต ภควตา วาจา ‘น ตาวาหํ ปาปิม ปรินิพฺพายิสฺสามิ, ยาว เม อุปาสกา น สาวกา ภวิสฺสนฺติ วิยตฺตา วินีตา วิสารทา พหุสฺสุตา ธมฺมธรา ธมฺมานุธมฺมปฺปฏิปนฺนา สามีจิปฺปฏิปนฺนา อนุธมฺมจาริโน, สกํ อาจริยกํ อุคฺคเหตฺวา อาจิกฺขิสฺสนฺติ เทเสสฺสนฺติ ปญฺญเปสฺสนฺติ ปฏฺฐเปสฺสนฺติ วิวริสฺสนฺติ วิภชิสฺสนฺติ อุตฺตานี กริสฺสนฺติ, อุปฺปนฺนํ ปรปฺปวาทํ สหธมฺเมน สุนิคฺคหิตํ นิคฺคเหตฺวา สปฺปาฏิหาริยํ ธมฺมํ เทเสสฺสนฺตี’ติ.\csromanspacer{} เอตรหิ โข ปน ภนฺเต อุปาสกา ภควโต สาวกา วิยตฺตา วินีตา วิสารทา พหุสฺสุตา ธมฺมธรา ธมฺมานุธมฺมปฺปฏิปนฺนา สามีจิปฺปฏิปนฺนา อนุธมฺมจาริโน, สกํ อาจริยกํ อุคฺคเหตฺวา อาจิกฺขนฺติ เทเสนฺติ ปญฺญเปนฺติ ปฏฺฐเปนฺติ วิวรนฺติ วิภชนฺติ อุตฺตานี กโรนฺติ, อุปฺปนฺนํ ปรปฺปวาทํ สหธมฺเมน สุนิคฺคหิตํ นิคฺคเหตฺวา สปฺปาฏิหาริยํ ธมฺมํ เทเสนฺติ.}

\vspace{\tipitakaparskip}
\csromancenter{มหาปรินิพฺพานสุตฺต ปรินิพฺพาตุทานิ ภนฺเต ภควา, ปรินิพฺพาตุ สุคโต, ปรินิพฺพานกาโลทานิ ภนฺเต ภควโต.}
\prose{\csromanfolio{89}ภาสิตา โข ปเนสา ภนฺเต ภควตา วาจา ‘น ตาวาหํ ปาปิม ปรินิพฺพายิสฺสามิ, ยาว เม อุปาสิกา น สาวิกา ภวิสฺสนฺติ วิยตฺตา วินีตา วิสารทา พหุสฺสุตา ธมฺมธรา ธมฺมานุธมฺมปฺปฏิปนฺนา สามีจิปฺปฏิปนฺนา อนุธมฺมจารินิโย, สกํ อาจริยกํ อุคฺคเหตฺวา อาจิกฺขิสฺสนฺติ เทเสสฺสนฺติ ปญฺญเปสฺสนฺติ ปฏฺฐเปสฺสนฺติ วิวริสฺสนฺติ วิภชิสฺสนฺติ อุตฺตานี กริสฺสนฺติ, อุปฺปนฺนํ ปรปฺปวาทํ สหธมฺเมน สุนิคฺคหิตํ นิคฺคเหตฺวา สปฺปาฏิหาริยํ ธมฺมํ เทเสสฺสนฺตี’ติ.\csromanspacer{} เอตรหิ โข ปน ภนฺเต อุปาสิกา ภควโต สาวิกา วิยตฺตา วินีตา วิสารทา พหุสฺสุตา ธมฺมธรา ธมฺมานุธมฺมปฺปฏิปนฺนา สามีจิปฺปฏิปนฺนา อนุธมฺมจารินิโย, สกํ อาจริยกํ อุคฺคเหตฺวา อาจิกฺขนฺติ เทเสนฺติ ปญฺญเปนฺติ ปฏฺฐเปนฺติ วิวรนฺติ วิภชนฺติ อุตฺตานี กโรนฺติ, อุปฺปนฺนํ ปรปฺปวาทํ สหธมฺเมน สุนิคฺคหิตํ นิคฺคเหตฺวา สปฺปาฏิหาริยํ ธมฺมํ เทเสนฺติ.\csromanspacer{} ปรินิพฺพาตุทานิ ภนฺเต ภควา, ปรินิพฺพาตุ สุคโต, ปรินิพฺพานกาโลทานิ ภนฺเต ภควโต.}

\prose{ภาสิตา โข ปเนสา ภนฺเต ภควตา วาจา ‘น ตาวาหํ ปาปิม ปรินิพฺพายิสฺสามิ, ยาว เม อิทํ พฺรหฺมจริยํ น อิทฺธญฺเจว ภวิสฺสติ ผีตญฺจ วิตฺถาริกํ พาหุชญฺญํ ปุถุภูตํ ยาว เทวมนุสฺเสหิ สุปฺปกาสิตนฺ’ติ.\csromanspacer{} เอตรหิ โข ปน ภนฺเต ภควโต พฺรหฺมจริยํ อิทฺธญฺเจว ผีตญฺจ วิตฺถาริกํ พาหุชญฺญํ ปุถุภูตํ ยาว เทวมนุสฺเสหิ สุปฺปกาสิตํ.\csromanspacer{} ปรินิพฺพาตุทานิ ภนฺเต ภควา, ปรินิพฺพาตุ สุคโต, ปรินิพฺพานกาโลทานิ ภนฺเต ภควโต”ติ.}

\prose{เอวํ วุตฺเต ภควา มารํ ปาปิมนฺตํ เอตทโวจ “อปฺโปสฺสุกฺโก ตฺวํ ปาปิม โหหิ, น จิรํ ตถาคตสฺส ปรินิพฺพานํ ภวิสฺสติ, อิโต ติณฺณํ มาสานํ อจฺจเยน ตถาคโต ปรินิพฺพายิสฺสตี”ติ.}

\csromanmatikaanchor{mtk.41}
\csromantocmarkref{section}{อายุสงฺขาร-โอสฺสชฺชน}{mtk.41}
\bookmark[dest={mtk.41},level=1]{อายุสงฺขาร-โอสฺสชฺชน}
\csromanheader{1}{\textbf{อายุสงฺขาร โอสฺสชฺชน}}
\proseitem{169}{อถ โข ภควา จาปาเล เจติเย สโต สมฺปชาโน อายุสงฺขารํ โอสฺสชิ, โอสฺสฏฺเฐ จ ภควตา อายุสงฺขาเร มหาภูมิจาโล อโหสิ ภิํสนโก สโลมหํโส\footnote{โลมหํโส (สฺยา)}, เทวทุนฺทุภิโย\footnote{เทวทุทฺรภิโย (ก)} จ \csromanfolio{90}ผลิํสุ.\csromanspacer{} อถ โข ภควา เอตมตฺถํ วิทิตฺวา ตายํ เวลายํ อิมํ อุทานํ อุทาเนสิ\csromandash{}}

\csromangathameasure{“ตุลมตุลญฺจ สมฺภวํ, ภวสงฺขารมวสฺสชิ มุนิ. \\ อชฺฌตฺตรโต สมาหิโต, อภินฺทิ กวจมิวตฺตสมฺภวนฺ”ติ.}
\csromangathagroup{\csromangathastanza{“ตุลมตุลญฺจ สมฺภวํ, ภวสงฺขารมวสฺสชิ มุนิ. \\ อชฺฌตฺตรโต สมาหิโต, อภินฺทิ กวจมิวตฺตสมฺภวนฺ”ติ.}}

\csromanmatikaanchor{mtk.42}
\csromantocmarkref{section}{มหาภูมิจาลเหตุ}{mtk.42}
\bookmark[dest={mtk.42},level=1]{มหาภูมิจาลเหตุ}
\csromanheader{1}{\textbf{มหาภูมิจาลเหตุ}}
\proseitem{170}{อถ โข อายสฺมโต อานนฺทสฺส เอตทโหสิ “อจฺฉริยํ วต โภ อพฺภุตํ วต โภ, มหา วตายํ ภูมิจาโล, สุมหา วตายํ ภูมิจาโล ภิํสนโก สโลมหํโส, เทวทุนฺทุภิโย จ ผลิํสุ.\csromanspacer{} โก นุ โข เหตุ โก ปจฺจโย มหโต ภูมิจาลสฺส ปาตุภาวายา”ติ.}

\prose{อถ โข อายสฺมา อานนฺโท เยน ภควา เตนุปสงฺกมิ, อุปสงฺกมิตฺวา ภควนฺตํ อภิวาเทตฺวา เอกมนฺตํ นิสีทิ, เอกมนฺตํ นิสินฺโน โข อายสฺมา อานนฺโท ภควนฺตํ เอตทโวจ “อจฺฉริยํ ภนฺเต อพฺภุตํ ภนฺเต, มหา วตายํ ภนฺเต ภูมิจาโล, สุมหา วตายํ ภนฺเต ภูมิจาโล ภิํสนโก สโลมหํโส, เทวทุนฺทุภิโย จ ผลิํสุ.\csromanspacer{} โก นุ โข ภนฺเต เหตุ โก ปจฺจโย มหโต ภูมิจาลสฺส ปาตุภาวายา”ติ.}

\proseitem{171}{อฏฺฐ โข อิเม อานนฺท เหตู อฏฺฐ ปจฺจยา มหโต ภูมิจาลสฺส ปาตุภาวาย.\csromanspacer{} กตเม อฏฺฐ.\csromanspacer{} อยํ อานนฺท มหาปถวี อุทเก ปติฏฺฐิตา, อุทกํ วาเต ปติฏฺฐิตํ, วาโต อากาสฏฺโฐ.\csromanspacer{} โหติ โข โส อานนฺท สมโย, ยํ มหาวาตา วายนฺติ, มหาวาตา วายนฺตา อุทกํ กมฺเปนฺติ, อุทกํ กมฺปิตํ ปถวิํ กมฺเปติ.\csromanspacer{} อยํ ปฐโม เหตุ ปฐโม ปจฺจโย มหโต ภูมิจาลสฺส ปาตุภาวาย.}

\prose{ปุน จปรํ อานนฺท สมโณ วา โหติ พฺราหฺมโณ วา อิทฺธิมา เจโตวสิปฺปตฺโต เทโว วา มหิทฺธิโก มหานุภาโว, ตสฺส ปริตฺตา ปถวีสญฺญา ภาวิตา โหติ, อปฺปมาณา อาโปสญฺญา.\csromanspacer{} โส อิมํ ปถวิํ กมฺเปติ สงฺกมฺเปติ สมฺปกมฺเปติ สมฺปเวเธติ.\csromanspacer{} อยํ ทุติโย เหตุ ทุติโย ปจฺจโย มหโต ภูมิจาลสฺส ปาตุภาวาย.}

\csromanheader{2}{3. มหาปรินิพฺพานสุตฺต}
\prose{\csromanfolio{91}ปุน จปรํ อานนฺท ยทา โพธิสตฺโต ตุสิตกายา จวิตฺวา สโต สมฺปชาโน มาตุกุจฺฉิํ โอกฺกมติ, ตทายํ ปถวี กมฺปติ สงฺกมฺปติ สมฺปกมฺปติ สมฺปเวธติ.\csromanspacer{} อยํ ตติโย เหตุ ตติโย ปจฺจโย มหโต ภูมิจาลสฺส ปาตุภาวาย.}

\prose{ปุน จปรํ อานนฺท ยทา โพธิสตฺโต สโต สมฺปชาโน มาตุกุจฺฉิสฺมา นิกฺขมติ, ตทายํ ปถวี กมฺปติ สงฺกมฺปติ สมฺปกมฺปติ สมฺปเวธติ.\csromanspacer{} อยํ จตุตฺโถ เหตุ จตุตฺโถ ปจฺจโย มหโต ภูมิจาลสฺส ปาตุภาวาย.}

\prose{ปุน จปรํ อานนฺท ยทา ตถาคโต อนุตฺตรํ สมฺมาสมฺโพธิํ อภิสมฺพุชฺฌติ, ตทายํ ปถวี กมฺปติ สงฺกมฺปติ สมฺปกมฺปติ สมฺปเวธติ.\csromanspacer{} อยํ ปญฺจโม เหตุ ปญฺจโม ปจฺจโย มหโต ภูมิจาลสฺส ปาตุภาวาย.}

\prose{ปุน จปรํ อานนฺท ยทา ตถาคโต อนุตฺตรํ ธมฺมจกฺกํ ปวตฺเตติ, ตทายํ ปถวี กมฺปติ สงฺกมฺปติ สมฺปกมฺปติ สมฺปเวธติ.\csromanspacer{} อยํ ฉฏฺโฐ เหตุ ฉฏฺโฐ ปจฺจโย มหโต ภูมิจาลสฺส ปาตุภาวาย.}

\prose{ปุน จปรํ อานนฺท ยทา ตถาคโต สโต สมฺปชาโน อายุสงฺขารํ โอสฺสชฺชติ, ตทายํ ปถวี กมฺปติ สงฺกมฺปติ สมฺปกมฺปติ สมฺปเวธติ.\csromanspacer{} อยํ สตฺตโม เหตุ สตฺตโม ปจฺจโย มหโต ภูมิจาลสฺส ปาตุภาวาย.}

\prose{ปุน จปรํ อานนฺท ยทา ตถาคโต อนุปาทิเสสาย นิพฺพานธาตุยา ปรินิพฺพายติ, ตทายํ ปถวี กมฺปติ สงฺกมฺปติ สมฺปกมฺปติ สมฺปเวธติ.\csromanspacer{} อยํ อฏฺฐโม เหตุ อฏฺฐโม ปจฺจโย มหโต ภูมิจาลสฺส ปาตุภาวาย.\csromanspacer{} อิเม โข อานนฺท อฏฺฐ เหตู อฏฺฐ ปจฺจยา มหโต ภูมิจาลสฺส ปาตุภาวาย.}

\csromanmatikaanchor{mtk.43}
\csromantocmarkref{section}{อฏฺฐปริสา}{mtk.43}
\bookmark[dest={mtk.43},level=1]{อฏฺฐปริสา}
\csromanheader{1}{\textbf{อฏฺฐปริสา}}
\proseitem{172}{อฏฺฐ โข อิมา อานนฺท ปริสา.\csromanspacer{} กตมา อฏฺฐ.\csromanspacer{} ขตฺติยปริสา พฺราหฺมณปริสา คหปติปริสา สมณปริสา จาตุมหาราชิกปริสา\footnote{จาตุมฺมหาราชิกปริสา (สี, สฺยา, กํ, อิ)} ตาวติํสปริสา มารปริสา พฺรหฺมปริสา.\csromanspacer{} อภิชานามิ โข ปนาหํ \csromanfolio{92}อานนฺท อเนกสตํ ขตฺติยปริสํ อุปสงฺกมิตา.\csromanspacer{} ตตฺรปิ มยา สนฺนิสินฺนปุพฺพญฺเจว สลฺลปิตปุพฺพญฺจ สากจฺฉา จ สมาปชฺชิตปุพฺพา.\csromanspacer{} ตตฺถ ยาทิสโก เตสํ วณฺโณ โหติ, ตาทิสโก มยฺหํ วณฺโณ โหติ.\csromanspacer{} ยาทิสโก เตสํ สโร โหติ, ตาทิสโก มยฺหํ สโร โหติ.\csromanspacer{} ธมฺมิยา กถาย สนฺทสฺเสมิ สมาทเปมิ สมุตฺเตเชมิ สมฺปหํเสมิ.\csromanspacer{} ภาสมานญฺจ มํ น ชานนฺติ “โก นุ โข อยํ ภาสติ เทโว วา มนุสฺโส วา”ติ.\csromanspacer{} ธมฺมิยา กถาย สนฺทสฺเสตฺวา สมาทเปตฺวา สมุตฺเตเชตฺวา สมฺปหํเสตฺวา อนฺตรธายามิ.\csromanspacer{} อนฺตรหิตญฺจ มํ น ชานนฺติ “โก นุ โข อยํ อนฺตรหิโต เทโว วา มนุสฺโส วา”ติ.\csromanspacer{} อภิชานามิ โข ปนาหํ อานนฺท อเนกสตํ พฺราหฺมณปริสํ ฯเปฯ คหปติปริสํ.\csromanspacer{} สมณปริสํ.\csromanspacer{} จาตุมหาราชิกปริสํ.\csromanspacer{} ตาวติํสปริสํ.\csromanspacer{} มารปริสํ.\csromanspacer{} พฺรหฺมปริสํ อุปสงฺกมิตา.\csromanspacer{} ตตฺรปิ มยา สนฺนิสินฺนปุพฺพญฺเจว สลฺลปิตปุพฺพญฺจ สากจฺฉา จ สมาปชฺชิตปุพฺพา.\csromanspacer{} ตตฺถ ยาทิสโก เตสํ วณฺโณ โหติ, ตาทิสโก มยฺหํ วณฺโณ โหติ.\csromanspacer{} ยาทิสโก เตสํ สโร โหติ, ตาทิสโก มยฺหํ สโร โหติ.\csromanspacer{} ธมฺมิยา กถาย สนฺทสฺเสมิ สมาทเปมิ สมุตฺเตเชมิ สมฺปหํเสมิ.\csromanspacer{} ภาสมานญฺจ มํ น ชานนฺติ “โก นุ โข อยํ ภาสติ เทโว วา มนุสฺโส วา”ติ.\csromanspacer{} ธมฺมิยา กถาย สนฺทสฺเสตฺวา สมาทเปตฺวา สมุตฺเตเชตฺวา สมฺปหํเสตฺวา อนฺตรธายามิ.\csromanspacer{} อนฺตรหิตญฺจ มํ น ชานนฺติ “โก นุ โข อยํ อนฺตรหิโต เทโว วา มนุสฺโส วา”ติ.\csromanspacer{} อิมา โข อานนฺท อฏฺฐ ปริสา.}

\csromanmatikaanchor{mtk.44}
\csromantocmarkref{section}{อฏฺฐ-อภิภายตน}{mtk.44}
\bookmark[dest={mtk.44},level=1]{อฏฺฐ-อภิภายตน}
\csromanheader{1}{\textbf{อฏฺฐ-อภิภายตน}}
\proseitem{173}{อฏฺฐ โข อิมานิ อานนฺท อภิภายตนานิ.\csromanspacer{} กตมานิ อฏฺฐ.\csromanspacer{} อชฺฌตฺตํ รูปสญฺญี เอโก พหิทฺธา รูปานิ ปสฺสติ ปริตฺตานิ สุวณฺณทุพฺพณฺณานิ, “ตานิ อภิภุยฺย ชานามิ ปสฺสามี”ติ เอวํสญฺญี โหติ.\csromanspacer{} อิทํ ปฐมํ อภิภายตนํ.}

\prose{อชฺฌตฺตํ รูปสญฺญี เอโก พหิทฺธา รูปานิ ปสฺสติ อปฺปมาณานิ สุวณฺณทุพฺพณฺณานิ, “ตานิ อภิภุยฺย ชานามิ ปสฺสามี”ติ เอวํสญฺญี โหติ.\csromanspacer{} อิทํ ทุติยํ อภิภายตนํ.}

\prose{อชฺฌตฺตํ อรูปสญฺญี เอโก พหิทฺธา รูปานิ ปสฺสติ ปริตฺตานิ สุวณฺณทุพฺพณฺณานิ, “ตานิ อภิภุยฺย ชานามิ ปสฺสามี”ติ เอวํสญฺญี โหติ.\csromanspacer{} อิทํ ตติยํ อภิภายตนํ.}

\csromanheader{2}{3. มหาปรินิพฺพานสุตฺต}
\prose{\csromanfolio{93}อชฺฌตฺตํ อรูปสญฺญี เอโก พหิทฺธา รูปานิ ปสฺสติ อปฺปมาณานิ สุวณฺณทุพฺพณฺณานิ, “ตานิ อภิภุยฺย ชานามิ ปสฺสามี”ติ เอวํสญฺญี โหติ.\csromanspacer{} อิทํ จตุตฺถํ อภิภายตนํ.}

\prose{อชฺฌตฺตํ อรูปสญฺญี เอโก พหิทฺธา รูปานิ ปสฺสติ นีลานิ นีลวณฺณานิ นีลนิทสฺสนานิ นีลนิภาสานิ.\csromanspacer{} เสยฺยถาปิ นาม อุมาปุปฺผํ นีลํ นีลวณฺณํ นีลนิทสฺสนํ นีลนิภาสํ.\csromanspacer{} เสยฺยถา วา ปน ตํ วตฺถํ พาราณเสยฺยกํ อุภโตภาควิมฏฺฐํ นีลํ นีลวณฺณํ นีลนิทสฺสนํ นิลนิภาสํ.\csromanspacer{} เอวเมว อชฺฌตฺตํ อรูปสญฺญี เอโก พหิทฺธา รูปานิ ปสฺสติ นีลานิ นีลวณฺณานิ นีลนิทสฺสนานิ นีลนิภาสานิ, “ตานิ อภิภุยฺย ชานามิ ปสฺสามี”ติ เอวํ สญฺญี โหติ.\csromanspacer{} อิทํ ปญฺจมํ อภิภายตนํ.}

\prose{อชฺฌตฺตํ อรูปสญฺญี เอโก พหิทฺธา รูปานิ ปสฺสติ ปีตานิ ปีตวณฺณานิ ปีตนิทสฺสนานิ ปีตนิภาสานิ.\csromanspacer{} เสยฺยถาปิ นาม กณิการปุปฺผํ ปีตํ ปีตวณฺณํ ปีตนิทสฺสนํ ปีตนิภาสํ.\csromanspacer{} เสยฺยถา วา ปน ตํ วตฺถํ พาราณเสยฺยกํ อุภโตภาควิมฏฺฐํ ปีตํ ปีตวณฺณํ ปีตนิทสฺสนํ ปีตนิภาสํ.\csromanspacer{} เอวเมว อชฺฌตฺตํ อรูปสญฺญี เอโก พหิทฺธา รูปานิ ปสฺสติ ปีตานิ ปีตวณฺณานิ ปีตนิทสฺสนานิ ปีตนิภาสานิ, “ตานิ อภิภุยฺย ชานามิ ปสฺสามี”ติ เอวํสญฺญี โหติ.\csromanspacer{} อิทํ ฉฏฺฐํ อภิภายตนํ.}

\prose{อชฺฌตฺตํ อรูปสญฺญี เอโก พหิทฺธา รูปานิ ปสฺสติ โลหิตกานิ โลหิตกวณฺณานิ โลหิตกนิทสฺสนานิ โลหิตกนิภาสานิ.\csromanspacer{} เสยฺยถาปิ นาม พนฺธุชีวกปุปฺผํ โลหิตกํ โลหิตกวณฺณํ โลหิตกนิทสฺสนํ โลหิตกนิภาสํ.\csromanspacer{} เสยฺยถา วา ปน ตํ วตฺถํ พาราณเสยฺยกํ อุภโตภาควิมฏฺฐํ โลหิตกํ โลหิตกวณฺณํ โลหิตกนิทสฺสนํ โลหิตกนิภาสํ.\csromanspacer{} เอวเมว อชฺฌตฺตํ อรูปสญฺญี เอโก พหิทฺธา รูปานิ ปสฺสติ โลหิตกานิ โลหิตกวณฺณานิ โลหิตกนิทสฺสนานิ โลหิตกนิภาสานิ, “ตานิ อภิภุยฺย ชานามิ ปสฺสามี”ติ เอวํสญฺญี โหติ.\csromanspacer{} อิทํ สตฺตมํ อภิภายตนํ.}

\prose{อชฺฌตฺตํ อรูปสญฺญี เอโก พหิทฺธา รูปานิ ปสฺสติ โอทาตานิ โอทาตวณฺณานิ โอทาตนิทสฺสนานิ โอทาตนิภาสานิ.\csromanspacer{} เสยฺยถาปิ นาม}

\prose{\csromanfolio{94}โอสธิตารกา โอทาตา โอทาตวณฺณา โอทาตนิทสฺสนา โอทาตนิภาสา.\csromanspacer{} เสยฺยถา วา ปน ตํ วตฺถํ พาราณเสยฺยกํ อุภโตภาควิมฏฺฐํ โอทาตํ โอทาตวณฺณํ โอทาตนิทสฺสนํ โอทาตนิภาสํ.\csromanspacer{} เอวเมว อชฺฌตฺตํ อรูปสญฺญี เอโก พหิทฺธา รูปานิ ปสฺสติ โอทาตานิ โอทาตวณฺณานิ โอทาตนิทสฺสนานิ โอทาตนิภาสานิ, “ตานิ อภิภุยฺย ชานามิ ปสฺสามี”ติ เอวํสญฺญี โหติ.\csromanspacer{} อิทํ อฏฺฐมํ อภิภายตนํ.\csromanspacer{} อิมานิ โข อานนฺท อฏฺฐ อภิภายตนานิ.}

\csromanmatikaanchor{mtk.45}
\csromantocmarkref{section}{อฏฺฐวิโมกฺข}{mtk.45}
\bookmark[dest={mtk.45},level=1]{อฏฺฐวิโมกฺข}
\csromanheader{1}{\textbf{อฏฺฐวิโมกฺข}}
\proseitem{174}{อฏฺฐ โข อิเม อานนฺท วิโมกฺขา.\csromanspacer{} กตเม อฏฺฐ.\csromanspacer{} รูปี รูปานิ ปสฺสติ, อยํ ปฐโม วิโมกฺโข.\csromanspacer{} อชฺฌตฺตํ อรูปสญฺญี พหิทฺธา รูปานิ ปสฺสติ, อยํ ทุติโย วิโมกฺโข.\csromanspacer{} สุภนฺเตว อธิมุตฺโต โหติ, อยํ ตติโย วิโมกฺโข.\csromanspacer{} สพฺพโส รูปสญฺญานํ สมติกฺกมา ปฏิฆสญฺญานํ อตฺถงฺคมา นานตฺตสญฺญานํ อมนสิการา “อนนฺโต อากาโส”ติ อากาสานญฺจายตนํ อุปสมฺปชฺช วิหรติ, อยํ จตุตฺโถ วิโมกฺโข.\csromanspacer{} สพฺพโส อากาสานญฺจายตนํ สมติกฺกมฺม “อนนฺตํ วิญฺญาณนฺ”ติ วิญฺญาณญฺจายตนํ อุปสมฺปชฺช วิหรติ, อยํ ปญฺจโม วิโมกฺโข.\csromanspacer{} สพฺพโส วิญฺญาณญฺจายตนํ สมติกฺกมฺม “นตฺถิ กิญฺจี”ติ อากิญฺจญฺญายตนํ อุปสมฺปชฺช วิหรติ, อยํ ฉฏฺโฐ วิโมกฺโข.\csromanspacer{} สพฺพโส อากิญฺจญฺญายตนํ สมติกฺกมฺม เนวสญฺญานาสญฺญายตนํ อุปสมฺปชฺช วิหรติ, อยํ สตฺตโม วิโมกฺโข.\csromanspacer{} สพฺพโส เนวสญฺญานาสญฺญายตนํ สมติกฺกมฺม สญฺญาเวทยิตนิโรธํ อุปสมฺปชฺช วิหรติ, อยํ อฏฺฐโม วิโมกฺโข.\csromanspacer{} อิเม โข อานนฺท อฏฺฐ วิโมกฺขา.}

\proseitem{175}{เอกมิทาหํ อานนฺท สมยํ อุรุเวลายํ วิหรามิ นชฺชา เนรญฺชราย ตีเร อชปาลนิโคฺรเธ ปฐมาภิสมฺพุทฺโธ.\csromanspacer{} อถ โข อานนฺท มาโร ปาปิมา เยนาหํ เตนุปสงฺกมิ, อุปสงฺกมิตฺวา เอกมนฺตํ อฏฺฐาสิ, เอกมนฺตํ ฐิโต โข อานนฺท มาโร ปาปิมา มํ เอตทโวจ “ปรินิพฺพาตุทานิ ภนฺเต ภควา, ปรินิพฺพาตุ สุคโต, ปรินิพฺพานกาโลทานิ ภนฺเต ภควโต”ติ.\csromanspacer{} เอวํ วุตฺเต อหํ อานนฺท มารํ ปาปิมนฺตํ เอตทโวจํ “น ตาวาหํ ปาปิม ปรินิพฺพายิสฺสามิ, ยาว เม ภิกฺขู น สาวกา ภวิสฺสนฺติ วิยตฺตา วินีตา วิสารทา พหุสฺสุตา ธมฺมธรา}

\vspace{\tipitakaparskip}
\csromancenter{มหาปรินิพฺพานสุตฺต ธมฺมานุธมฺมปฺปฏิปนฺนา สามีจิปฺปฏิปนฺนา อนุธมฺมจาริโน, สกํ อาจริยกํ อุคฺคเหตฺวา อาจิกฺขิสฺสนฺติ เทเสสฺสนฺติ ปญฺญเปสฺสนฺติ ปฏฺฐเปสฺสนฺติ วิวริสฺสนฺติ วิภชิสฺสนฺติ อุตฺตานี กริสฺสนฺติ, อุปฺปนฺนํ ปรปฺปวาทํ สหธมฺเมน สุนิคฺคหิตํ นิคฺคเหตฺวา สปฺปาฏิหาริยํ ธมฺมํ เทเสสฺสนฺติ.}
\prose{\csromanfolio{95}น ตาวาหํ ปาปิม ปรินิพฺพายิสฺสามิ, ยาว เม ภิกฺขุนิโย น สาวิกา ภวิสฺสนฺติ วิยตฺตา วินีตา วิสารทา พหุสฺสุตา ธมฺมธรา ธมฺมานุธมฺมปฺปฏิปนฺนา สามีจิปฺปฏิปนฺนา อนุธมฺมจารินิโย, สกํ อาจริยกํ อุคฺคเหตฺวา อาจิกฺขิสฺสนฺติ เทเสสฺสนฺติ ปญฺญเปสฺสนฺติ ปฏฺฐเปสฺสนฺติ วิวริสฺสนฺติ วิภชิสฺสนฺติ อุตฺตานี กริสฺสนฺติ, อุปฺปนฺนํ ปรปฺปวาทํ สหธมฺเมน สุนิคฺคหิตํ นิคฺคเหตฺวา สปฺปาฏิหาริยํ ธมฺมํ เทเสสฺสนฺติ.}

\prose{น ตาวาหํ ปาปิม ปรินิพฺพายิสฺสามิ, ยาว เม อุปาสกา น สาวกา ภวิสฺสนฺติ วิยตฺตา วินีตา วิสารทา พหุสฺสุตา ธมฺมธรา ธมฺมานุธมฺมปฺปฏิปนฺนา สามีจิปฺปฏิปนฺนา อนุธมฺมจาริโน สกํ อาจริยกํ อุคฺคเหตฺวา อาจิกฺขิสฺสนฺติ เทเสสฺสนฺติ ปญฺญเปสฺสนฺติ ปฏฺฐเปสฺสนฺติ วิวริสฺสนฺติ วิภชิสฺสนฺติ อุตฺตานี กริสฺสนฺติ, อุปฺปนฺนํ ปรปฺปวาทํ สหธมฺเมน สุนิคฺคหิตํ นิคฺคเหตฺวา สปฺปาฏิหาริยํ ธมฺมํ เทเสสฺสนฺติ.}

\prose{น ตาวาหํ ปาปิม ปรินิพฺพายิสฺสามิ, ยาว เม อุปาสิกา น สาวิกา ภวิสฺสนฺติ วิยตฺตา วินีตา วิสารทา พหุสฺสุตา ธมฺมธรา ธมฺมานุธมฺมปฺปฏิปนฺนา สามีจิปฺปฏิปนฺนา อนุธมฺมจารินิโย สกํ อาจริยกํ อุคฺคเหตฺวา อาจิกฺขิสฺสนฺติ เทเสสฺสนฺติ ปญฺญเปสฺสนฺติ ปฏฺฐเปสฺสนฺติ วิวริสฺสนฺติ วิภชิสฺสนฺติ อุตฺตานี กริสฺสนฺติ, อุปฺปนฺนํ ปรปฺปวาทํ สหธมฺเมน สุนิคฺคหิตํ นิคฺคเหตฺวา สปฺปาฏิหาริยํ ธมฺมํ เทเสสฺสนฺติ.}

\prose{น ตาวาหํ ปาปิม ปรินิพฺพายิสฺสามิ, ยาว เม อิทํ พฺรหฺมจริยํ น อิทฺธญฺเจว ภวิสฺสติ ผีตญฺจ วิตฺถาริกํ พาหุชญฺญํ ปุถุภูตํ ยาว เทวมนุสฺเสหิ สุปฺปกาสิตนฺ”ติ.}

\proseitem{176}{อิทาเนว โข อานนฺท อชฺช จาปาเล เจติเย มาโร ปาปิมา เยนาหํ เตนุปสงฺกมิ, อุปสงฺกมิตฺวา เอกมนฺตํ อฏฺฐาสิ, เอกมนฺตํ ฐิโต โข อานนฺท มาโร ปาปิมา มํ เอตทโวจ “ปรินิพฺพาตุทานิ ภนฺเต ภควา, ปรินิพฺพาตุ สุคโต, ปรินิพฺพานกาโลทานิ ภนฺเต ภควโต.\csromanspacer{} ภาสิตา โข ปเนสา ภนฺเต ภควตา วาจา ‘น ตาวาหํ ปาปิม \csromanfolio{96}ปรินิพฺพายิสฺสามิ, ยาว เม ภิกฺขู น สาวกา ภวิสฺสนฺติ ฯเปฯ ยาว เม ภิกฺขุนิโย น สาวิกา ภวิสฺสนฺติ ฯเปฯ ยาว เม อุปาสกา น สาวกา ภวิสฺสนฺติ ฯเปฯ ยาว เม อุปาสิกา น สาวิกา ภวิสฺสนฺติ ฯเปฯ ยาว เม อิทํ พฺรหฺมจริยํ น อิทฺธญฺเจว ภวิสฺสติ ผีตญฺจ วิตฺถาริกํ พาหุชญฺญํ ปุถุภูตํ ยาว เทวมนุสฺเสหิ สุปฺปกาสิตนฺ’ติ.\csromanspacer{} เอตรหิ โข ปน ภนฺเต ภควโต พฺรหฺมจริยํ อิทฺธญฺเจว ผีตญฺจ วิตฺถาริกํ พาหุชญฺญํ ปุถุภูตํ ยาว เทวมนุสฺเสหิ สุปฺปกาสิตํ.\csromanspacer{} ปรินิพฺพาตุทานิ ภนฺเต ภควา, ปรินิพฺพาตุ สุคโต, ปรินิพฺพานกาโลทานิ ภนฺเต ภควโต”ติ.}

\proseitem{177}{เอวํ วุตฺเต อหํ อานนฺท มารํ ปาปิมนฺตํ เอตทโวจํ “อปฺโปสฺสุกฺโก ตฺวํ ปาปิม โหหิ, น จิรํ ตถาคตสฺส ปรินิพฺพานํ ภวิสฺสติ, อิโต ติณฺณํ มาสานํ อจฺจเยน ตถาคโต ปรินิพฺพายิสฺสตี”ติ.\csromanspacer{} อิทาเนว โข อานนฺท อชฺช จาปาเล เจติเย ตถาคเตน สเตน สมฺปชาเนน อายุสงฺขาโร โอสฺสฏฺโฐติ.}

\csromanmatikaanchor{mtk.46}
\csromantocmarkref{section}{อานนฺทยาจนกถา}{mtk.46}
\bookmark[dest={mtk.46},level=1]{อานนฺทยาจนกถา}
\csromanheader{1}{\textbf{อานนฺทยาจนกถา}}
\proseitem{178}{เอวํ วุตฺเต อายสฺมา อานนฺโท ภควนฺตํ เอตทโวจ “ติฏฺฐตุ ภนฺเต ภควา กปฺปํ, ติฏฺฐตุ สุคโต กปฺปํ พหุชนหิตาย พหุชนสุขาย โลกานุกมฺปาย อตฺถาย หิตาย สุขาย เทวมนุสฺสานนฺ”ติ.}

\prose{อลํทานิ อานนฺท มา ตถาคตํ ยาจิ, อกาโลทานิ อานนฺท ตถาคตํ ยาจนายาติ.\csromanspacer{} ทุติยมฺปิ โข อายสฺมา อานนฺโท ฯเปฯ.\csromanspacer{} ตติยมฺปิ โข อายสฺมา อานนฺโท ภควนฺตํ เอตทโวจ “ติฏฺฐตุ ภนฺเต ภควา กปฺปํ, ติฏฺฐตุ สุคโต กปฺปํ พหุชนหิตาย พหุชนสุขาย โลกานุกมฺปาย อตฺถาย หิตาย สุขาย เทวมนุสฺสานนฺ”ติ.}

\prose{สทฺทหสิ ตฺวํ อานนฺท ตถาคตสฺส โพธินฺติ.\csromanspacer{} เอวํ ภนฺเต.\csromanspacer{} อถ กิญฺจรหิ ตฺวํ อานนฺท ตถาคตํ ยาวตติยกํ อภินิปฺปีเฬสีติ.\csromanspacer{} สมฺมุขา เมตํ ภนฺเต ภควโต สุตํ สมฺมุขา ปฏิคฺคหิตํ “ยสฺส กสฺสจิ อานนฺท จตฺตาโร อิทฺธิปาทา ภาวิตา พหุลีกตา ยานีกตา วตฺถุกตา อนุฏฺฐิตา ปริจิตา สุสมารทฺธา, โส อากงฺขมาโน กปฺปํ วา ติฏฺเฐยฺย กปฺปาวเสสํ วา.\csromanspacer{} ตถาคตสฺส โข อานนฺท จตฺตาโร อิทฺธิปาทา ภาวิตา พหุลีกตา ยานีกตา วตฺถุกตา อนุฏฺฐิตา ปริจิตา สุสมารทฺธา,}

\vspace{\tipitakaparskip}
\csromancenter{มหาปรินิพฺพานสุตฺต โส อากงฺขมาโน อานนฺท ตถาคโต กปฺปํ วา ติฏฺเฐยฺย กปฺปาวเสสํ วา”ติ.\csromanspacer{} สทฺทหสิ ตฺวํ อานนฺทาติ.\csromanspacer{} เอวํ ภนฺเต.\csromanspacer{} ตสฺมาติหานนฺท ตุเยฺหเวตํ ทุกฺกฏํ ตุเยฺหเวตํ อปรทฺธํ, ยํ ตฺวํ ตถาคเตน เอวํ โอฬาริเก นิมิตฺเต กยิรมาเน โอฬาริเก โอภาเส กยิรมาเน นาสกฺขิ ปฏิวิชฺฌิตุํ, น ตถาคตํ ยาจิ “ติฏฺฐตุ ภนฺเต ภควา กปฺปํ, ติฏฺฐตุ สุคโต กปฺปํ พหุชนหิตาย พหุชนสุขาย โลกานุกมฺปาย อตฺถาย หิตาย สุขาย เทวมนุสฺสานนฺ”ติ.\csromanspacer{} สเจ ตฺวํ อานนฺท ตถาคตํ ยาเจยฺยาสิ, เทฺวว เต วาจา ตถาคโต ปฏิกฺขิเปยฺย, อถ ตติยกํ อธิวาเสยฺย.\csromanspacer{} ตสฺมาติหานนฺท ตุเยฺหเวตํ ทุกฺกฏํ, ตุเยฺหเวตํ อปรทฺธํ.}
\proseitem{179}{\csromanfolio{97}เอกมิทาหํ อานนฺท สมยํ ราชคเห วิหรามิ คิชฺฌกูเฏ ปพฺพเต.\csromanspacer{} ตตฺราปิ โข ตาหํ อานนฺท อามนฺเตสิํ “รมณียํ อานนฺท ราชคหํ, รมณีโย อานนฺท คิชฺฌกูโฏ ปพฺพโต.\csromanspacer{} ยสฺส กสฺสจิ อานนฺท จตฺตาโร อิทฺธิปาทา ภาวิตา พหุลีกตา ยานีกตา วตฺถุกตา อนุฏฺฐิตา ปริจิตา สุสมารทฺธา, โส อากงฺขมาโน กปฺปํ วา ติฏฺเฐยฺย กปฺปาวเสสํ วา.\csromanspacer{} ตถาคตสฺส โข อานนฺท จตฺตาโร อิทฺธิปาทา ภาวิตา พหุลีกตา ยานีกตา วตฺถุกตา อนุฏฺฐิตา ปริจิตา สุสมารทฺธา, โส อากงฺขมาโน อานนฺท ตถาคโต กปฺปํ วา ติฏฺเฐยฺย กปฺปาวเสสํ วา”ติ.\csromanspacer{} เอวมฺปิ โข ตฺวํ อานนฺท ตถาคเตน โอฬาริเก นิมิตฺเต กยิรมาเน โอฬาริเก โอภาเส กยิรมาเน นาสกฺขิ ปฏิวิชฺฌิตุํ, น ตถาคตํ ยาจิ “ติฏฺฐตุ ภนฺเต ภควา กปฺปํ, ติฏฺฐตุ สุคโต กปฺปํ พหุชนหิตาย พหุชนสุขาย โลกานุกมฺปาย อตฺถาย หิตาย สุขาย เทวมนุสฺสานนฺ”ติ.\csromanspacer{} สเจ ตฺวํ อานนฺท ตถาคตํ ยาเจยฺยาสิ, เทฺวว เต วาจา ตถาคโต ปฏิกฺขิเปยฺย, อถ ตติยกํ อธิวาเสยฺย.\csromanspacer{} ตสฺมาติหานนฺท ตุเยฺหเวตํ ทุกฺกฏํ, ตุเยฺหเวตํ อปรทฺธํ.}

\proseitem{180}{เอกมิทาหํ อานนฺท สมยํ ตตฺเถว ราชคเห วิหรามิ โคตมนิโคฺรเธ ฯเปฯ ตตฺเถว ราชคเห วิหรามิ โจรปปาเต.\csromanspacer{} ตตฺเถว ราชคเห วิหรามิ เวภารปสฺเส สตฺตปณฺณิคุหายํ.\csromanspacer{} ตตฺเถว ราชคเห วิหรามิ อิสิคิลิปสฺเส กาฬสิลายํ.\csromanspacer{} ตตฺเถว ราชคเห วิหรามิ สีตวเน สปฺปโสณฺฑิกปพฺภาเร.\csromanspacer{} ตตฺเถว ราชคเห วิหรามิ ตโปทาราเม.\csromanspacer{} ตตฺเถว ราชคเห วิหรามิ เวฬุวเน กลนฺทกนิวาเป.\csromanspacer{} ตตฺเถว ราชคเห วิหรามิ ชีวกมฺพวเน.\csromanspacer{} ตตฺเถว ราชคเห วิหรามิ มทฺทกุจฺฉิสฺมิํ \csromanfolio{98}มิคทาเย.\csromanspacer{} ตตฺราปิ โข ตาหํ อานนฺท อามนฺเตสิํ “รมณียํ อานนฺท ราชคหํ, รมณีโย คิชฺฌกูโฏ ปพฺพโต, รมณีโย โคตมนิโคฺรโธ, รมณีโย โจรปปาโต, รมณียา เวภารปสฺเส สตฺตปณฺณิคุหา, รมณียา อิสิคิลิปสฺเส กาฬสิลา.\csromanspacer{} รมณีโย สีตวเน สปฺปโสณฺฑิกปพฺภาโร, รมณีโย ตโปทาราโม, รมณีโย เวฬุวเน กลนฺทกนิวาโป, รมณียํ ชีวกมฺพวนํ, รมณีโย มทฺทกุจฺฉิสฺมิํ มิคทาโย.\csromanspacer{} ยสฺส กสฺสจิ อานนฺท จตฺตาโร อิทฺธิปาทา ภาวิตา พหุลีกตา ยานีกตา วตฺถุกตา อนุฏฺฐิตา ปริจิตา สุสมารทฺธา ฯเปฯ อากงฺขมาโน อานนฺท ตถาคโต กปฺปํ วา ติฏฺเฐยฺย กปฺปาวเสสํ วา”ติ.\csromanspacer{} เอวมฺปิ โข ตฺวํ อานนฺท ตถาคเตน โอฬาริเก นิมิตฺเต กยิรมาเน โอฬาริเก โอภาเส กยิรมาเน นาสกฺขิ ปฏิวิชฺฌิตุํ, น ตถาคตํ ยาจิ “ติฏฺฐตุ ภนฺเต ภควา กปฺปํ, ติฏฺฐตุ สุคโต กปฺปํ พหุชนหิตาย พหุชนสุขาย โลกานุกมฺปาย อตฺถาย หิตาย สุขาย เทวมนุสฺสานนฺ”ติ.\csromanspacer{} สเจ ตฺวํ อานนฺท ตถาคตํ ยาเจยฺยาสิ, เทฺวว เต วาจา ตถาคโต ปฏิกฺขิเปยฺย, อถ ตติยกํ อธิวาเสยฺย.\csromanspacer{} ตสฺมาติหานนฺท ตุเยฺหเวตํ ทุกฺกฏํ, ตุเยฺหเวตํ อปรทฺธํ.}

\proseitem{181}{เอกมิทาหํ อานนฺท สมยํ อิเธว เวสาลิยํ วิหรามิ อุเทเน เจติเย.\csromanspacer{} ตตฺราปิ โข ตาหํ อานนฺท อามนฺเตสิํ “รมณียา อานนฺท เวสาลี, รมณียํ อุเทนํ เจติยํ.\csromanspacer{} ยสฺส กสฺสจิ อานนฺท จตฺตาโร อิทฺธิปาทา ภาวิตา พหุลีกตา ยานีกตา วตฺถุกตา อนุฏฺฐิตา ปริจิตา สุสมารทฺธา, โส อากงฺขมาโน กปฺปํ วา ติฏฺเฐยฺย กปฺปาวเสสํ วา.\csromanspacer{} ตถาคตสฺส โข อานนฺท จตฺตาโร อิทฺธิปาทา ภาวิตา พหุลีกตา ยานีกตา วตฺถุกตา อนุฏฺฐิตา ปริจิตา สุสมารทฺธา, โส อากงฺขมาโน อานนฺท ตถาคโต กปฺปํ วา ติฏฺเฐยฺย กปฺปาวเสสํ วา”ติ.\csromanspacer{} เอวมฺปิ โข ตฺวํ อานนฺท ตถาคเตน โอฬาริเก นิมิตฺเต กยิรมาเน โอฬาริเก โอภาเส กยิรมาเน นาสกฺขิ ปฏิวิชฺฌิตุํ, น ตถาคตํ ยาจิ “ติฏฺฐตุ ภนฺเต ภควา กปฺปํ, ติฏฺฐตุ สุคโต กปฺปํ พหุชนหิตาย พหุชนสุขาย โลกานุกมฺปาย อตฺถาย หิตาย สุขาย เทวมนุสฺสานนฺ”ติ.\csromanspacer{} สเจ ตฺวํ อานนฺท ตถาคตํ ยาเจยฺยาสิ, เทฺวว เต วาจา ตถาคโต ปฏิกฺขิเปยฺย, อถ ตติยกํ อธิวาเสยฺย, ตสฺมาติหานนฺท ตุเยฺหเวตํ ทุกฺกฏํ, ตุเยฺหเวตํ อปรทฺธํ.}

\csromanheader{2}{3. มหาปรินิพฺพานสุตฺต}
\proseitem{182}{\csromanfolio{99}เอกมิทาหํ อานนฺท สมยํ อิเธว เวสาลิยํ วิหรามิ โคตมเก เจติเย ฯเปฯ อิเธว เวสาลิยํ วิหรามิ สตฺตมฺเพ เจติเย.\csromanspacer{} อิเธว เวสาลิยํ วิหรามิ พหุปุตฺเต เจติเย.\csromanspacer{} อิเธว เวสาลิยํ วิหรามิ สารนฺทเท เจติเย.\csromanspacer{} อิทาเนว โข ตาหํ อานนฺท อชฺช จาปาเล เจติเย อามนฺเตสิํ “รมณียา อานนฺท เวสาลี, รมณียํ อุเทนํ เจติยํ, รมณียํ โคตมกํ เจติยํ, รมณียํ สตฺตมฺพํ เจติยํ, รมณียํ พหุปุตฺตํ เจติยํ, รมณียํ สารนฺททํ เจติยํ, รมณียํ จาปาลํ เจติยํ.\csromanspacer{} ยสฺส กสฺสจิ อานนฺท จตฺตาโร อิทฺธิปาทา ภาวิตา พหุลีกตา ยานีกตา วตฺถุกตา อนุฏฺฐิตา ปริจิตา สุสมารทฺธา, โส อากงฺขมาโน กปฺปํ วา ติฏฺเฐยฺย กปฺปาวเสสํ วา.\csromanspacer{} ตถาคตสฺส โข อานนฺท จตฺตาโร อิทฺธิปาทา ภาวิตา พหุลีกตา ยานีกตา วตฺถุกตา อนุฏฺฐิตา ปริจิตา สุสมารทฺธา, โส อากงฺขมาโน อานนฺท ตถาคโต กปฺปํ วา ติฏฺเฐยฺย กปฺปาวเสสํ วา”ติ.\csromanspacer{} เอวมฺปิ โข ตฺวํ อานนฺท ตถาคเตน โอฬาริเก นิมิตฺเต กยิรมาเน โอฬาริเก โอภาเส กยิรมาเน นาสกฺขิ ปฏิวิชฺฌิตุํ, น ตถาคตํ ยาจิ “ติฏฺฐตุ ภนฺเต ภควา กปฺปํ, ติฏฺฐตุ สุคโต กปฺปํ พหุชนหิตาย พหุชนสุขาย โลกานุกมฺปาย อตฺถาย หิตาย สุขาย เทวมนุสฺสานนฺ”ติ.\csromanspacer{} สเจ ตฺวํ อานนฺท ตถาคตํ ยาเจยฺยาสิ, เทฺวว เต วาจา ตถาคโต ปฏิกฺขิเปยฺย, อถ ตติยกํ อธิวาเสยฺย.\csromanspacer{} ตสฺมาติหานนฺท ตุเยฺหเวตํ ทุกฺกฏํ, ตุเยฺหเวตํ อปรทฺธํ.}

\proseitem{183}{นนุ เอตํ\footnote{เอวํ (สฺยา, อิ)} อานนฺท มยา ปฏิกจฺเจว\footnote{ปฏิคจฺเจว (สี, อิ)} อกฺขาตํ “สพฺเพเหว ปิเยหิ มนาเปหิ นานาภาโว วินาภาโว อญฺญถาภาโว, ตํ กุเตตฺถ อานนฺท ลพฺภา, ยํ ตํ ชาตํ ภูตํ สงฺขตํ ปโลกธมฺมํ, ตํ วต มาปลุชฺชีติ เนตํ ฐานํ วิชฺชติ”.\csromanspacer{} ยํ โข ปเนตํ อานนฺท ตถาคเตน จตฺตํ วนฺตํ มุตฺตํ ปหีนํ ปฏินิสฺสฏฺฐํ โอสฺสฏฺโฐ อายุสงฺขาโร, เอกํเสน วาจา ภาสิตา “น จิรํ ตถาคตสฺส ปรินิพฺพานํ ภวิสฺสติ, อิโต ติณฺณํ มาสานํ อจฺจเยน ตถาคโต ปรินิพฺพายิสฺสตี”ติ, ตญฺจ\footnote{ตํ วจนํ (สี)} ตถาคโต ชีวิตเหตุ ปุน ปจฺจาวมิสฺสตีติ\footnote{ปจฺจาคมิสฺสตีติ (สฺยา, ก)} เนตํ ฐานํ วิชฺชติ.\csromanspacer{} อายามานนฺท, เยน มหาวนํ กูฏาคารสาลา เตนุปสงฺกมิสฺสามา”ติ. “เอวํ ภนฺเต”ติ \csromanfolio{100}โข อายสฺมา อานนฺโท ภควโต ปจฺจสฺโสสิ.\csromanspacer{} อถ โข ภควา อายสฺมตา อานนฺเทน สทฺธิํ เยน มหาวนํ กูฏาคารสาลา เตนุปสงฺกมิ, อุปสงฺกมิตฺวา อายสฺมนฺตํ อานนฺทํ อามนฺเตสิ “คจฺฉ ตฺวํ อานนฺท ยาวติกา ภิกฺขู เวสาลิํ อุปนิสฺสาย วิหรนฺติ, เต สพฺเพ อุปฏฺฐานสาลายํ สนฺนิปาเตหี”ติ. “เอวํ ภนฺเต”ติ โข อายสฺมา อานนฺโท ภควโต ปฏิสฺสุตฺวา ยาวติกา ภิกฺขู เวสาลิํ อุปนิสฺสาย วิหรนฺติ, เต สพฺเพ อุปฏฺฐานสาลายํ สนฺนิปาเตตฺวา เยน ภควา เตนุปสงฺกมิ, อุปสงฺกมิตฺวา ภควนฺตํ อภิวาเทตฺวา เอกมนฺตํ อฏฺฐาสิ, เอกมนฺตํ ฐิโต โข อายสฺมา อานนฺโท ภควนฺตํ เอตทโวจ “สนฺนิปติโต ภนฺเต ภิกฺขุสํโฆ, ยสฺสทานิ ภนฺเต ภควา กาลํ มญฺญตี”ติ.}

\proseitem{184}{อถ โข ภควา เยนุปฏฺฐานสาลา เตนุปสงฺกมิ, อุปสงฺกมิตฺวา ปญฺญตฺเต อาสเน นิสีทิ, นิสชฺช โข ภควา ภิกฺขู อามนฺเตสิ\csromandash{}ตสฺมาติห ภิกฺขเว เย เต มยา ธมฺมา อภิญฺญา เทสิตา, เต โว สาธุกํ อุคฺคเหตฺวา อาเสวิตพฺพา ภาเวตพฺพา พหุลีกาตพฺพา, ยถยิทํ พฺรหฺมจริยํ อทฺธนิยํ อสฺส จิรฏฺฐิติกํ, ตทสฺส พหุชนหิตาย พหุชนสุขาย โลกานุกมฺปาย อตฺถาย หิตาย สุขาย เทวมนุสฺสานํ.\csromanspacer{} กตเม จ เต ภิกฺขเว ธมฺมา มยา อภิญฺญา เทสิตา, เย โว สาธุกํ อุคฺคเหตฺวา อาเสวิตพฺพา ภาเวตพฺพา พหุลีกาตพฺพา, ยถยิทํ พฺรหฺมจริยํ อทฺธนิยํ อสฺส จิรฏฺฐิติกํ, ตทสฺส พหุชนหิตาย พหุชนสุขาย โลกานุกมฺปาย อตฺถาย หิตาย สุขาย เทวมนุสฺสานํ.\csromanspacer{} เสยฺยถิทํ, จตฺตาโร สติปฏฺฐานา จตฺตาโร สมฺมปฺปธานา จตฺตาโร อิทฺธิปาทา ปญฺจินฺทฺริยานิ ปญฺจ พลานิ สตฺต โพชฺฌงฺคา อริโย อฏฺฐงฺคิโก มคฺโค.\csromanspacer{} อิเม โข เต ภิกฺขเว ธมฺมา มยา อภิญฺญา เทสิตา, เย โว สาธุกํ อุคฺคเหตฺวา อาเสวิตพฺพา ภาเวตพฺพา พหุลีกาตพฺพา, ยถยิทํ พฺรหฺมจริยํ อทฺธนิยํ อสฺส จิรฏฺฐิติกํ, ตทสฺส พหุชนหิตาย พหุชนสุขาย โลกานุกมฺปาย อตฺถาย หิตาย สุขาย เทวมนุสฺสานนฺติ.}

\proseitem{185}{อถ โข ภควา ภิกฺขู อามนฺเตสิ “หนฺท ทานิ ภิกฺขเว อามนฺตยามิ โว, วยธมฺมา สงฺขารา อปฺปมาเทน สมฺปาเทถ.\csromanspacer{} น จิรํ}

\vspace{\tipitakaparskip}
\csromancenter{มหาปรินิพฺพานสุตฺต ตถาคตสฺส ปรินิพฺพานํ ภวิสฺสติ, อิโต ติณฺณํ มาสานํ อจฺจเยน ตถาคโต ปรินิพฺพายิสฺสตี”ติ.\csromanspacer{} อิทมโวจ ภควา, อิทํ วตฺวาน สุคโต อถาปรํ เอตทโวจ สตฺถา\footnote{อิโต ปรํ สฺยามโปตฺถเก เอวํปิ ปาโฐ ทิสฺสติ\csromandash{} ทหราปิ จ เย วุทฺธา, เย พาลา เย จ ปณฺฑิตา. อฑฺฒาเจว ทลิทฺทา จ, สพฺเพ มจฺจุปรายนา. ยถาปิ กุมฺภการสฺส, กตํ มตฺติกภาชนํ. ขุทฺทกญฺจ มหนฺตญฺจ, ยญฺจ ปกฺกํ ยญฺจ อามกํ. สพฺพํ เภทปริยนฺตํ, เอวํ มจฺจาน ชีวิตํ. อถาปรํ เอตทโวจ สตฺถา\csromandash{}}\csromandash{}}
\vspace{0.5\baselineskip}
\csromangathameasure{ปริปกฺโก วโย มยฺหํ, ปริตฺตํ มม ชีวิตํ. \\ ปหาย โว คมิสฺสามิ, กตํ เม สรณมตฺตโน. \\ อปฺปมตฺตา สตีมนฺโต, สุสีลา โหถ ภิกฺขโว. \\ สุสมาหิตสงฺกปฺปา, สจิตฺตมนุรกฺขถ. \\ โย อิมสฺมิํ ธมฺมวินเย, อปฺปมตฺโต วิหสฺสติ. \\ ปหาย ชาติสํสารํ, ทุกฺขสฺสนฺตํ กริสฺสตีติ.}
\csromangathagroup{\csromangathastanza{\csromanfolio{101}ปริปกฺโก วโย มยฺหํ, ปริตฺตํ มม ชีวิตํ. \\ ปหาย โว คมิสฺสามิ, กตํ เม สรณมตฺตโน.}\csromangathastanza{อปฺปมตฺตา สตีมนฺโต, สุสีลา โหถ ภิกฺขโว. \\ สุสมาหิตสงฺกปฺปา, สจิตฺตมนุรกฺขถ.}\csromangathastanza{โย อิมสฺมิํ ธมฺมวินเย, อปฺปมตฺโต วิหสฺสติ\footnote{วิหริสฺสติ (สฺยา), วิเหสฺสติ (สี)}. \\ ปหาย ชาติสํสารํ, ทุกฺขสฺสนฺตํ กริสฺสตีติ.}}

\csromanheader{2}{ตติโย ภาณวาโร.}
\csromansectionrule
\csromanmatikaanchor{mtk.47}
\csromantocmarkref{section}{นาคาปโลกิต}{mtk.47}
\bookmark[dest={mtk.47},level=1]{นาคาปโลกิต}
\csromanheader{1}{\textbf{นาคาปโลกิต}}
\proseitem{186}{อถ โข ภควา ปุพฺพณฺหสมยํ นิวาเสตฺวา ปตฺตจีวรมาทาย เวสาลิํ ปิณฺฑาย ปาวิสิ, เวสาลิยํ ปิณฺฑาย จริตฺวา ปจฺฉาภตฺตํ ปิณฺฑปาตปฏิกฺกนฺโต นาคาปโลกิตํ เวสาลิํ อปโลเกตฺวา อายสฺมนฺตํ อานนฺทํ อามนฺเตสิ “อิทํ ปจฺฉิมกํ อานนฺท ตถาคตสฺส เวสาลิยา ทสฺสนํ ภวิสฺสติ, อายามานนฺท, เยน ภณฺฑคาโม\footnote{ภณฺฑุคาโม(ก)} เตนุปสงฺกมิสฺสามา”ติ. “เอวํ ภนฺเต”ติ โข อายสฺมา อานนฺโท ภควโต ปจฺจสฺโสสิ.}

\prose{\csromanfolio{102}อถ โข ภควา มหตา ภิกฺขุสํเฆน สทฺธิํ เยน ภณฺฑคาโม ตทวสริ, ตตฺร สุทํ ภควา ภณฺฑคาเม วิหรติ.\csromanspacer{} ตตฺร โข ภควา ภิกฺขู อามนฺเตสิ จตุนฺนํ ภิกฺขเว ธมฺมานํ อนนุโพธา อปฺปฏิเวธา เอวมิทํ ทีฆมทฺธานํ สนฺธาวิตํ สํสริตํ มมญฺเจว ตุมฺหากญฺจ.\csromanspacer{} กตเมสํ จตุนฺนํ, อริยสฺส ภิกฺขเว สีลสฺส อนนุโพธา อปฺปฏิเวธา เอวมิทํ ทีฆมทฺธานํ สนฺธาวิตํ สํสริตํ มมญฺเจว ตุมฺหากญฺจ.\csromanspacer{} อริยสฺส ภิกฺขเว สมาธิสฺส อนนุโพธา อปฺปฏิเวธา เอวมิทํ ทิฆมทฺธานํ สนฺธาวิตํ สํสริตํ มมญฺเจว ตุมฺหากญฺจ.\csromanspacer{} อริยาย ภิกฺขเว ปญฺญาย อนนุโพธา อปฺปฏิเวธา เอวมิทํ ทีฆมทฺธานํ สนฺธาวิตํ สํสริตํ มมญฺเจว ตุมฺหากญฺจ.\csromanspacer{} อริยาย ภิกฺขเว วิมุตฺติยา อนนุโพธา อปฺปฏิเวธา เอวมิทํ ทีฆมทฺธานํ สนฺธาวิตํ สํสริตํ มมญฺเจว ตุมฺหากญฺจ.\csromanspacer{} ตยิทํ ภิกฺขเว อริยํ สีลํ อนุพุทฺธํ ปฏิวิทฺธํ, อริโย สเมธิ อนุพุทฺโธ ปฏิวิทฺโธ, อริยา ปญฺญา อนุพุทฺธา ปฏิวิทฺธา, อริยา วิมุตฺติ อนุพุทฺธา ปฏิวิทฺธา, อุจฺฉินฺนา ภวตณฺหา, ขีณา ภวเนตฺติ, นตฺถิ ทานิ ปุนพฺภโวติ.\csromanspacer{} อิทมโวจ ภควา, อิทํ วตฺวาน สุคโต อถาปรํ เอตทโวจ สตฺถา\csromandash{}}

\csromangathameasure{“สีลํ สมาธิ ปญฺญา จ, วิมุตฺติ จ อนุตฺตรา. \\ อนุพุทฺธา อิเม ธมฺมา, โคตเมน ยสสฺสินา. \\ อิติ พุทฺโธ อภิญฺญาย, ธมฺมมกฺขาสิ ภิกฺขุนํ. \\ ทุกฺขสฺสนฺตกโร สตฺถา, จกฺขุมา นิพฺพานํ”ติ.}
\csromangathagroup{\csromangathastanza{“สีลํ สมาธิ ปญฺญา จ, วิมุตฺติ จ อนุตฺตรา. \\ อนุพุทฺธา อิเม ธมฺมา, โคตเมน ยสสฺสินา.}\csromangathastanza{อิติ พุทฺโธ อภิญฺญาย, ธมฺมมกฺขาสิ ภิกฺขุนํ. \\ ทุกฺขสฺสนฺตกโร สตฺถา, จกฺขุมา นิพฺพานํ”ติ.}}

\prose{ตตฺราปิ สุทํ ภควา ภณฺฑคาเม วิหรนฺโต เอตเทว พหุลํ ภิกฺขูนํ ธมฺมิํ กถํ กโรติ “อิติ สีลํ อิติ สมาธิ อิติ ปญฺญา.\csromanspacer{} สีลปริภาวิโต สมาธิ มหปฺผโล โหติ มหานิสํโส.\csromanspacer{} สมาธิปริภาวิตา ปญฺญา มหปฺผลา โหติ มหานิสํสา.\csromanspacer{} ปญฺญาปริภาวิตํ จิตฺตํ สมฺมเทว อาสเวหิ วิมุจฺจติ.\csromanspacer{} เสยฺยถิทํ, กามาสวา ภวาสวา อวิชฺชาสวา”ติ.}

\csromanmatikaanchor{mtk.48}
\csromantocmarkref{section}{จตุมหาปเทสกถา}{mtk.48}
\bookmark[dest={mtk.48},level=1]{จตุมหาปเทสกถา}
\csromanheader{1}{\textbf{จตุมหาปเทสกถา}}
\proseitem{187}{อถ โข ภควา ภณฺฑคาเม ยถาภิรนฺตํ วิหริตฺวา อายสฺมนฺตํ อานนฺทํ อามนฺเตสิ “อายามานนฺท เยน หตฺถิคาโม, เยน อมฺพคาโม, เยน ชมฺพุคาโม, เยน โภคนครํ เตนุปสงฺกมิสฺสามา”ติ. “เอวํ ภนฺเต”ติ โข อายสฺมา อานนฺโท ภควโต ปจฺจสฺโสสิ.\csromanspacer{} อถ}

\vspace{\tipitakaparskip}
\csromancenter{มหาปรินิพฺพานสุตฺต โข ภควา มหตา ภิกฺขุสํเฆน สทฺธิํ เยน โภคนครํ ตทวสริ, ตตฺร สุทํ ภควา โภคนคเร วิหรติ อานนฺเท\footnote{สานนฺทเร (ก)} เจติเย.\csromanspacer{} ตตฺร โข ภควา ภิกฺขู อามนฺเตสิ “จตฺตาโรเม ภิกฺขเว มหาปเทเส เทเสสฺสามิ, ตํ สุณาถ, สาธุกํ มนสิกโรถ, ภาสิสฺสามี”ติ. “เอวํ ภนฺเต”ติ โข เต ภิกฺขู ภควโต ปจฺจสฺโสสุํ.\csromanspacer{} ภควา เอตทโวจ\csromandash{}}
\proseitem{188}{\csromanfolio{103}อิธ ภิกฺขเว ภิกฺขุ เอวํ วเทยฺย “สมฺมุขา เมตํ อาวุโส ภควโต สุตํ สมฺมุขา ปฏิคฺคหิตํ, อยํ ธมฺโม อยํ วินโย อิทํ สตฺถุสาสนนฺ”ติ, ตสฺส ภิกฺขเว ภิกฺขุโน ภาสิตํ เนว อภินนฺทิตพฺพํ นปฺปฏิกฺโกสิตพฺพํ.\csromanspacer{} อนภินนฺทิตฺวา อปฺปฏิกฺโกสิตฺวา ตานิ ปทพฺยญฺชนานิ สาธุกํ อุคฺคเหตฺวา สุตฺเต โอสาเรตพฺพานิ\footnote{โอตาเรตพฺพานิ,}, วินเย สนฺทสฺเสตพฺพานิ.\csromanspacer{} ตานิ เจ สุตฺเต โอสาริยมานานิ\footnote{โอตาริยมานานิ,} วินเย สนฺทสฺสิยมานานิ น เจว สุตฺเต โอสรนฺติ\footnote{โอตรนฺติ (สี, อิ, อํ 1. 487)}, น จ วินเย สนฺทิสฺสนฺติ, นิฏฺฐเมตฺถ คนฺตพฺพํ “อทฺธา อิทํ น เจว ตสฺส ภควโต วจนํ, อิมสฺส จ ภิกฺขุโน ทุคฺคหิตนฺ”ติ, อิติเหตํ ภิกฺขเว ฉฏฺเฏยฺยาถ.\csromanspacer{} ตานิ เจ สุตฺเต โอสาริยมานานิ วินเย สนฺทสฺสิยมานานิ สุตฺเต เจว โอสรนฺติ, วินเย จ สนฺทิสฺสนฺติ, นิฏฺฐเมตฺถ คนฺตพฺพํ “อทฺธา อิทํ ตสฺส ภควโต วจนํ, อิมสฺส จ ภิกฺขุโน สุคฺคหิตนฺ”ติ.\csromanspacer{} อิทํ ภิกฺขเว ปฐมํ มหาปเทสํ ธาเรยฺยาถ.}

\prose{อิธ ปน ภิกฺขเว ภิกฺขุ เอวํ วเทยฺย “อมุกสฺมิํ นาม อาวาเส สํโฆ วิหรติ สเถโร สปาโมกฺโข, ตสฺส เม สํฆสฺส สมฺมุขา สุตํ สมฺมุขา ปฏิคฺคหิตํ, อยํ ธมฺโม อยํ วินโย อิทํ สตฺถุสาสนนฺ”ติ, ตสฺส ภิกฺขเว ภิกฺขุโน ภาสิตํ เนว อภินนฺทิตพฺพํ นปฺปฏิกฺโกสิตพฺพํ.\csromanspacer{} อนภินนฺทิตฺวา อปฺปฏิกฺโกสิตฺวา ตานิ ปทพฺยญฺชนานิ สาธุกํ อุคฺคเหตฺวา สุตฺเต โอสาเรตพฺพานิ, วินเย สนฺทสฺเสตพฺพานิ.\csromanspacer{} ตานิ เจ สุตฺเต โอสาริยมานานิ วินเย สนฺทสฺสิยมานานิ น เจว สุตฺเต โอสรนฺติ, น จ วินเย สนฺทิสฺสนฺติ, นิฏฺฐเมตฺถ คนฺตพฺพํ “อทฺธา อิทํ น เจว ตสฺส ภควโต วจนํ, ตสฺส จ สํฆสฺส ทุคฺคหิตนฺ”ติ, อิติเหตํ ภิกฺขเว ฉฏฺเฏยฺยาถ.\csromanspacer{} ตานิ เจ สุตฺเต โอสาริยมานานิ วินเย สนฺทสฺสิยมานานิ สุตฺเต \csromanfolio{104}เจว โอสรนฺติ, วินเย จ สนฺทิสฺสนฺติ, นิฏฺฐเมตฺถ คนฺตพฺพํ “อทฺธา อิทํ ตสฺส ภควโต วจนํ, ตสฺส จ สํฆสฺส สุคฺคหิตนฺ”ติ.\csromanspacer{} อิทํ ภิกฺขเว ทุติยํ มหาปเทสํ ธาเรยฺยาถ.}

\prose{อิธ ปน ภิกฺขเว ภิกฺขุ เอวํ วเทยฺย “อมุกสฺมิํ นาม อาวาเส สมฺพหุลา เถรา ภิกฺขู วิหรนฺติ พหุสฺสุตา อาคตาคมา ธมฺมธรา วินยธรา มาติกาธรา, เตสํ เม เถรานํ สมฺมุขา สุตํ สมฺมุขา ปฏิคฺคหิตํ, อยํ ธมฺโม อยํ วินโย อิทํ สตฺถุสาสนนฺ”ติ, ตสฺส ภิกฺขเว ภิกฺขุโน ภาสิตํ เนว อภินนฺทิตพฺพํ ฯเปฯ น จ วินเย สนฺทิสฺสนฺติ.\csromanspacer{} นิฏฺฐเมตฺถ คนฺตพฺพํ “อทฺธา อิทํ น เจว ตสฺส ภควโต วจนํ, เตสญฺจ เถรานํ ทุคฺคหิตนฺ”ติ, อิติเหตํ ภิกฺขเว ฉฏฺเฏยฺยาถ.\csromanspacer{} ตานิ เจ สุตฺเต โอสาริยมานานิ ฯเปฯ วินเย จ สนฺทิสฺสนฺติ, นิฏฺฐเมตฺถ คนฺตพฺพํ “อทฺธา อิทํ ตสฺส ภควโต วจนํ, เตสญฺจ เถรานํ สุคฺคหิตนฺ”ติ.\csromanspacer{} อิทํ ภิกฺขเว ตติยํ มหาปเทสํ ธาเรยฺยาถ.}

\prose{อิธ ปน ภิกฺขเว ภิกฺขุ เอวํ วเทยฺย “อมุกสฺมิํ นาม อาวาเส เอโก เถโร ภิกฺขุ วิหรติ พหุสฺสุโต อาคตาคโม ธมฺมธโร วินยธโร มาติกาธโร, ตสฺส เม เถรสฺส สมฺมุขา สุตํ สมฺมุขา ปฏิคฺคหิตํ, อยํ ธมฺโม อยํ วินโย อิทํ สตฺถุสาสนนฺ”ติ, ตสฺส ภิกฺขเว ภิกฺขุโน ภาสิตํ เนว อภินนฺทิตพฺพํ นปฺปฏิกฺโกสิตพฺพํ.\csromanspacer{} อนภินนฺทิตฺวา อปฺปฏิกฺโกสิตฺวา ตานิ ปทพฺยญฺชนานิ สาธุกํ อุคฺคเหตฺวา สุตฺเต โอสาเรตพฺพานิ, วินเย สนฺทสฺเสตพฺพานิ.\csromanspacer{} ตานิ เจ สุตฺเต โอสาริยมานานิ วินเย สนฺทสฺสิยมานานิ น เจว สุตฺเต โอสรนฺติ, น จ วินเย สนฺทิสฺสนฺติ, นิฏฺฐเมตฺถ คนฺตพฺพํ “อทฺธา อิทํ น เจว ตสฺส ภควโต วจนํ, ตสฺส จ เถรสฺส ทุคฺคหิตนฺ”ติ, อิติเหตํ ภิกฺขเว ฉฏฺเฏยฺยาถ.\csromanspacer{} ตานิ เจ สุตฺเต โอสาริยมานานิ วินเย สนฺทสฺสิยมานานิ สุตฺเต เจว โอสรนฺติ, วินเย จ สนฺทิสฺสนฺติ, นิฏฺฐเมตฺถ คนฺตพฺพํ “อทฺธา อิทํ ตสฺส ภควโต วจนํ, ตสฺส จ เถรสฺส สุคฺคหิตนฺ”ติ.\csromanspacer{} อิทํ ภิกฺขเว จตุตฺถํ มหาปเทสํ ธาเรยฺยาถ.\csromanspacer{} อิเม โข ภิกฺขเว จตฺตาโร มหาปเทเส ธาเรยฺยาถาติ.}

\prose{ตตฺรปิ สุทํ ภควา โภคนคเร วิหรนฺโต อานนฺเท เจติเย เอตเทว พหุลํ ภิกฺขูนํ ธมฺมิํ กถํ กโรติ “อิติ สีลํ อิติ สมาธิ อิติ ปญฺญา.\csromanspacer{} สีลปริภาวิโต สมาธิ มหปฺผโล โหติ}

\vspace{\tipitakaparskip}
\csromancenter{มหาปรินิพฺพานสุตฺต มหานิสํโส, สมาธิปริภาวิตา ปญฺญา มหปฺผลา โหติ มหานิสํสา, ปญฺญาปริภาวิตํ จิตฺตํ สมฺมเทว อาสเวหิ วิมุจฺจติ, เสยฺยถิทํ, กามาสวา ภวาสวา อวิชฺชาสวา”ติ.}
\csromanmatikaanchor{mtk.49}
\csromantocmarkref{section}{กมฺมารปุตฺตจุนฺทวตฺถุ}{mtk.49}
\bookmark[dest={mtk.49},level=1]{กมฺมารปุตฺตจุนฺทวตฺถุ}
\csromanheader{1}{\textbf{กมฺมารปุตฺตจุนฺทวตฺถุ}}
\proseitem{189}{\csromanfolio{105}อถ โข ภควา โภคนคเร ยถาภิรนฺตํ วิหริตฺวา อายสฺมนฺตํ อานนฺทํ อามนฺเตสิ “อายามานนฺท, เยน ปาวา เตนุปสงฺกมิสฺสามา”ติ. “เอวํ ภนฺเต”ติ โข อายสฺมา อานนฺโท ภควโต ปจฺจสฺโสสิ.\csromanspacer{} อถ โข ภควา มหตา ภิกฺขุสํเฆน สทฺธิํ เยน ปาวา ตทวสริ, ตตฺร สุทํ ภควา ปาวายํ วิหรติ จุนฺทสฺส กมฺมารปุตฺตสฺส อมฺพวเน.\csromanspacer{} อสฺโสสิ โข จุนฺโท กมฺมารปุตฺโต “ภควา กิร ปาวํ อนุปฺปตฺโต ปาวายํ วิหรติ มยฺหํ อมฺพวเน”ติ.\csromanspacer{} อถ โข จุนฺโท กมฺมารปุตฺโต เยน ภควา เตนุปสงฺกมิ, อุปสงฺกมิตฺวา ภควนฺตํ อภิวาเทตฺวา เอกมนฺตํ นิสีทิ, เอกมนฺตํ นิสินฺนํ โข จุนฺทํ กมฺมารปุตฺตํ ภควา ธมฺมิยา กถาย สนฺทสฺเสสิ สมาทเปสิ สมุตฺเตเชสิ สมฺปหํเสสิ.\csromanspacer{} อถ โข จุนฺโท กมฺมารปุตฺโต ภควตา ธมฺมิยา กถาย สนฺทสฺสิโต สมาทปิโต สมุตฺเตชิโต สมฺปหํสิโต ภควนฺตํ เอตทโวจ “อธิวาเสตุ เม ภนฺเต ภควา สฺวาตนาย ภตฺตํ สทฺธิํ ภิกฺขุสํเฆนา”ติ.\csromanspacer{} อธิวาเสสิ ภควา ตุณฺหีภาเวน.\csromanspacer{} อถ โข จุนฺโท กมฺมารปุตฺโต ภควโต อธิวาสนํ วิทิตฺวา อุฏฺฐายาสนา ภควนฺตํ อภิวาเทตฺวา ปทกฺขิณํ กตฺวา ปกฺกามิ.}

\prose{อถ โข จุนฺโท กมฺมารปุตฺโต ตสฺสา รตฺติยา อจฺจเยน สเก นิเวสเน ปณีตํ ขาทนียํ โภชนียํ ปฏิยาทาเปตฺวา ปหูตญฺจ สูกรมทฺทวํ ภควโต กาลํ อาโรจาเปสิ “กาโล ภนฺเต นิฏฺฐิตํ ภตฺตนฺ”ติ.\csromanspacer{} อถ โข ภควา ปุพฺพณฺหสมยํ นิวาเสตฺวา ปตฺตจีวรมาทาย สทฺธิํ ภิกฺขุสํเฆน เยน จุนฺทสฺส กมฺมารปุตฺตสฺส นิเวสนํ เตนุปสงฺกมิ, อุปสงฺกมิตฺวา ปญฺญตฺเต อาสเน นิสีทิ, นิสชฺช โข ภควา จุนฺทํ กมฺมารปุตฺตํ อามนฺเตสิ “ยํ เต จุนฺท สูกรมทฺทวํ ปฏิยตฺตํ, เตน มํ ปริวิส.\csromanspacer{} ยํ ปนญฺญํ ขาทนียํ โภชนียํ ปฏิยตฺตํ, เตน ภิกฺขุสํฆํ ปริวิสา”ติ. “เอวํ ภนฺเต”ติ โข จุนฺโท กมฺมารปุตฺโต ภควโต ปฏิสฺสุตฺวา ยํ อโหสิ สูกรมทฺทวํ ปฏิยตฺตํ, เตน ภควนฺตํ ปริวิสิ.\csromanspacer{} ยํ ปนญฺญํ ขาทนียํ โภชนียํ \csromanfolio{106}ปฏิยตฺตํ, เตน ภิกฺขุสํฆํ ปริวิสิ.\csromanspacer{} อถ โข ภควา จุนฺทํ กมฺมารปุตฺตํ อามนฺเตสิ “ยํ เต จุนฺท สูกรมทฺทวํ อวสิฏฺฐํ, ตํ โสพฺเภ นิขณาหิ.\csromanspacer{} นาหํ ตํ จุนฺท ปสฺสามิ สเทวเก โลเก สมารเก สพฺรหฺมเก สสฺสมณพฺราหฺมณิยา ปชาย สเทวมนุสฺสาย, ยสฺส ตํ ปริภุตฺตํ สมฺมา ปริณามํ คจฺเฉยฺย อญฺญตฺร ตถาคตสฺสา”ติ. “เอวํ ภนฺเต”ติ โข จุนฺโท กมฺมารปุตฺโต ภควโต ปฏิสฺสุตฺวา ยํ อโหสิ สูกรมทฺทวํ อวสิฏฺฐํ, ตํ โสพฺเภ นิขณิตฺวา เยน ภควา เตนุปสงฺกมิ, อุปสงฺกมิตฺวา ภควนฺตํ อภิวาเทตฺวา เอกมนฺตํ นิสีทิ, เอกมนฺตํ นิสินฺนํ โข จุนฺทํ กมฺมารปุตฺตํ ภควา ธมฺมิยา กถาย สนฺทสฺเสตฺวา สมาทเปตฺวา สมุตฺเตเชตฺวา สมฺปหํเสตฺวา อุฏฺฐายาสนา ปกฺกามิ.}

\proseitem{190}{อถ โข ภควโต จุนฺทสฺส กมฺมารปุตฺตสฺส ภตฺตํ ภุตฺตาวิสฺส ขโร อาพาโธ อุปฺปชฺชิ, โลหิตปกฺขนฺทิกา ปพาฬฺหา เวทนา วตฺตนฺติ มารณนฺติกา, ตา สุทํ ภควา สโต สมฺปชาโน อธิวาเสสิ อวิหญฺญมาโน.\csromanspacer{} อถ โข ภควา อายสฺมนฺตํ อานนฺทํ อามนฺเตสิ “อายามานนฺท เยน กุสินารา เตนุปสงฺกมิสฺสามา”ติ. “เอวํ ภนฺเต”ติ โข อายสฺมา อานนฺโท ภควโต ปจฺจสฺโสสิ.}

\csromangathameasure{จุนฺทสฺส ภตฺตํ ภุญฺชิตฺวา, กมฺมารสฺสาติ เม สุตํ. \\ อาพาธํ สมฺผุสี ธีโร, ปพาฬฺหํ มารณนฺติกํ. \\ ภุตฺตสฺส จ สูกรมทฺทเวน, \\ พฺยาธิปฺปพาโฬฺห อุทปาทิ สตฺถุโน. \\ วิเรจมาโน ภควา อโวจ, \\ คจฺฉามหํ กุสินารํ นครนฺติ.}
\csromangathagroup{\csromangathastanza{จุนฺทสฺส ภตฺตํ ภุญฺชิตฺวา, กมฺมารสฺสาติ เม สุตํ. \\ อาพาธํ สมฺผุสี ธีโร, ปพาฬฺหํ มารณนฺติกํ.}\csromangathastanza{ภุตฺตสฺส จ สูกรมทฺทเวน, \\ พฺยาธิปฺปพาโฬฺห อุทปาทิ สตฺถุโน. \\ วิเรจมาโน\footnote{วิริจฺจมาโน (สี, สฺยา, ก), วิริญฺจมาโน (?)} ภควา อโวจ, \\ คจฺฉามหํ กุสินารํ นครนฺติ.}}

\csromanmatikaanchor{mtk.50}
\csromantocmarkref{section}{ปานียาหรณ}{mtk.50}
\bookmark[dest={mtk.50},level=1]{ปานียาหรณ}
\csromanheader{1}{\textbf{ปานียาหรณ}}
\proseitem{191}{อถ โข ภควา มคฺคา โอกฺกมฺม เยน อญฺญตรํ รุกฺขมูลํ เตนุปสงฺกมิ, อุปสงฺกมิตฺวา อายสฺมนฺตํ อานนฺทํ อามนฺเตสิ “อิงฺฆ เม ตฺวํ อานนฺท จตุคฺคุณํ สํฆาฏิํ ปญฺญเปหิ, กิลนฺโตสฺมิ อานนฺท นิสีทิสฺสามี”ติ. “เอวํ ภนฺเต”ติ โข อายสฺมา อานนฺโท ภควโต ปฏิสฺสุตฺวา จตุคฺคุณํ สํฆาฏิํ ปญฺญเปสิ.\csromanspacer{} นิสีทิ ภควา ปญฺญตฺเต อาสเน, นิสชฺช โข ภควา}

\vspace{\tipitakaparskip}
\csromancenter{มหาปรินิพฺพานสุตฺต อายสฺมนฺตํ อานนฺทํ อามนฺเตสิ “อิงฺฆ เม ตฺวํ อานนฺท ปานียํ อาหร, ปิปาสิโตสฺมิ อานนฺท ปิวิสฺสามี”ติ.\csromanspacer{} เอวํ วุตฺเต อายสฺมา อานนฺโท ภควนฺตํ เอตทโวจ “อิทานิ ภนฺเต ปญฺจมตฺตานิ สกฏสตานิ อติกฺกนฺตานิ, ตํ จกฺกจฺฉินฺนํ อุทกํ ปริตฺตํ ลุฬิตํ อาวิลํ สนฺทติ, อยํ ภนฺเต กกุธา\footnote{กกุถา (สี, อิ)} นที อวิทูเร อจฺโฉทกา สาโตทกา สีโตทกา เสโตทกา\footnote{เสตกา (สี)} สุปฺปติตฺถา รมณียา.\csromanspacer{} เอตฺถ ภควา ปานียญฺจ ปิวิสฺสติ, คตฺตานิ จ สีตี\footnote{สีตํ (สี, อิ, ก)} กริสฺสตี”ติ.}
\prose{\csromanfolio{107}ทุติยมฺปิ โข ภควา อายสฺมนฺตํ อานนฺทํ อามนฺเตสิ “อิงฺฆ เม ตฺวํ อานนฺท ปานียํ อาหร, ปิปาสิโตสฺมิ อานนฺท ปิวิสฺสามี”ติ.\csromanspacer{} ทุติยมฺปิ โข อายสฺมา อานนฺโท ภควนฺตํ เอตทโวจ “อิทานิ ภนฺเต ปญฺจมตฺตานิ สกฏสตานิ อติกฺกนฺตานิ, ตํ จกฺกจฺฉินฺนํ อุทกํ ปริตฺตํ ลุฬิตํ อาวิลํ สนฺทติ, อยํ ภนฺเต กกุธา นที อวิทูเร อจฺโฉทกา สาโตทกา สีโตทกา เสโตทกา สุปฺปติตฺถา รมณียา.\csromanspacer{} เอตฺถ ภควา ปานียญฺจ ปิวิสฺสติ, คตฺตานิ จ สีตีกริสฺสตี”ติ.}

\prose{ตติยมฺปิ โข ภควา อายสฺมนฺตํ อานนฺทํ อามนฺเตสิ “อิงฺฆ เม ตฺวํ อานนฺท ปานียํ อาหร, ปิปาสิโตสฺมิ อานนฺท ปิวิสฺสามี”ติ. “เอวํ ภนฺเต”ติ โข อายสฺมา อานนฺโท ภควโต ปฏิสฺสุตฺวา ปตฺตํ คเหตฺวา เยน สา นทิกา เตนุปสงฺกมิ.\csromanspacer{} อถ โข สา นทิกา จกฺกจฺฉินฺนา ปริตฺตา ลุฬิตา อาวิลา สนฺทมานา อายสฺมนฺเต อานนฺเท อุปสงฺกมนฺเต อจฺฉา วิปฺปสนฺนา อนาวิลา สนฺทิตฺถ\footnote{สนฺทติ (สฺยา)}.\csromanspacer{} อถ โข อายสฺมโต อานนฺทสฺส เอตทโหสิ “อจฺฉริยํ วต โภ อพฺภุตํ วต โภ ตถาคตสฺส มหิทฺธิกตา มหานุภาวตา.\csromanspacer{} อยํ หิ สา นทิกา จกฺกจฺฉินฺนา ปริตฺตา ลุฬิตา อาวิลา สนฺทมานา มยิ อุปสงฺกมนฺเต อจฺฉา วิปฺปสนฺนา อนาวิลา สนฺทตี”ติ.\csromanspacer{} ปตฺเตน ปานียํ อาทาย เยน ภควา เตนุปสงฺกมิ, อุปสงฺกมิตฺวา ภควนฺตํ เอตทโวจ “อจฺฉริยํ ภนฺเต อพฺภุตํ ภนฺเต ตถาคตสฺส มหิทฺธิกตา มหานุภาวตา.\csromanspacer{} อิทานิ สา ภนฺเต นทิกา จกฺกจฺฉินฺนา ปริตฺตา ลุฬิตา อาวิลา สนฺทมานา มยิ อุปสงฺกมนฺเต อจฺฉา วิปฺปสนฺนา อนาวิลา สนฺทิตฺถ.\csromanspacer{} ปิวตุ ภควา ปานียํ ปิวตุ สุคโต ปานียนฺ”ติ.\csromanspacer{} อถ โข ภควา ปานียํ อปายิ.}

\csromanmatikaanchor{mtk.51}
\csromantocmarkref{section}{ปุกฺกุสมลฺลปุตฺตวตฺถุ}{mtk.51}
\bookmark[dest={mtk.51},level=1]{ปุกฺกุสมลฺลปุตฺตวตฺถุ}
\csromanheader{1}{\textbf{ปุกฺกุสมลฺลปุตฺตวตฺถุ}}
\proseitem{192}{\csromanfolio{108}เตน โข ปน สมเยน ปุกฺกุโส มลฺลปุตฺโต อาฬารสฺส กาลามสฺส สาวโก กุสินาราย ปาวํ อทฺธานมคฺคปฺปฏิปนฺโน โหติ.\csromanspacer{} อทฺทสา โข ปุกฺกุโส มลฺลปุตฺโต ภควนฺตํ อญฺญตรสฺมิํ รุกฺขมูเล นิสินฺนํ, ทิสฺวา เยน ภควา เตนุปสงฺกมิ, อุปสงฺกมิตฺวา ภควนฺตํ อภิวาเทตฺวา เอกมนฺตํ นิสีทิ, เอกมนฺตํ นิสินฺโน โข ปุกฺกุโส มลฺลปุตฺโต ภควนฺตํ เอตทโวจ\csromandash{}อจฺฉริยํ ภนฺเต อพฺภุตํ ภนฺเต, สนฺเตน วต ภนฺเต ปพฺพชิตา วิหาเรน วิหรนฺติ.\csromanspacer{} ภูตปุพฺพํ ภนฺเต อาฬาโร กาลาโม อทฺธานมคฺคปฺปฏิปนฺโน มคฺคา โอกฺกมฺม อวิทูเร อญฺญตรสฺมิํ รุกฺขมูเล ทิวาวิหารํ นิสีทิ.\csromanspacer{} อถ โข ภนฺเต ปญฺจมตฺตานิ สกฏสตานิ อาฬารํ กาลามํ นิสฺสาย นิสฺสาย อติกฺกมิํสุ.\csromanspacer{} อถ โข ภนฺเต อญฺญตโร ปุริโส ตสฺส สกฏสตฺถสฺส\footnote{สกฏสตสฺส (ก)} ปิฏฺฐิโต ปิฏฺฐิโต อาคจฺฉนฺโต เยน อาฬาโร กาลาโม เตนุปสงฺกมิ, อุปสงฺกมิตฺวา อาฬารํ กาลามํ เอตทโวจ “อปิ ภนฺเต ปญฺจมตฺตานิ สกฏสตานิ อติกฺกนฺตานิ อทฺทสา”ติ. “น โข อหํ อาวุโส อทฺทสนฺ”ติ.\csromanspacer{} กิํ ปน ภนฺเต สทฺทํ อสฺโสสีติ. “น โข อหํ อาวุโส สทฺทํ อสฺโสสินฺ”ติ.\csromanspacer{} กิํ ปน ภนฺเต สุตฺโต อโหสีติ. “น โข อหํ อาวุโส สุตฺโต อโหสินฺ”ติ.\csromanspacer{} กิํ ปน ภนฺเต สญฺญี อโหสีติ. “เอวมาวุโส”ติ.\csromanspacer{} โสตฺวํ ภนฺเต สญฺญี สมาโน ชาคโร ปญฺจมตฺตานิ สกฏสตานิ นิสฺสาย นิสฺสาย อติกฺกนฺตานิ เนว อทฺทส, น ปน สทฺทํ อสฺโสสิ, อปิสุ\footnote{อปิ หิ (สี, สฺยา, อิ)} เต ภนฺเต สํฆาฏิ รเชน โอกิณฺณาติ. “เอวมาวุโส”ติ.\csromanspacer{} อถ โข ภนฺเต ตสฺส ปุริสสฺส เอตทโหสิ “อจฺฉริยํ วต โภ อพฺภุตํ วต โภ, สนฺเตน วต โภ ปพฺพชิตา วิหาเรน วิหรนฺติ.\csromanspacer{} ยตฺร หิ นาม สญฺญี สมาโน ชาคโร ปญฺจมตฺตานิ สกฏสตานิ นิสฺสาย นิสฺสาย อติกฺกนฺตานิ เนว ทกฺขติ, น ปน สทฺทํ โสสฺสตี”ติ, อาฬาเร กาลาเม อุฬารํ ปสาทํ ปเวเทตฺวา ปกฺกามีติ.}

\proseitem{193}{ตํ กิํ มญฺญสิ ปุกฺกุส, กตมํ นุ โข ทุกฺกรตรํ วา ทุรภิสมฺภวตรํ วา.\csromanspacer{} โย วา สญฺญี สมาโน ชาคโร ปญฺจมตฺตานิ สกฏสตานิ นิสฺสาย นิสฺสาย อติกฺกนฺตานิ เนว ปสฺเสยฺย, น ปน}

\vspace{\tipitakaparskip}
\csromancenter{มหาปรินิพฺพานสุตฺต สทฺทํ สุเณยฺย.\csromanspacer{} โย วา สญฺญี สมาโน ชาคโร เทเว วสฺสนฺเต เทเว คฬคฬายนฺเต วิชฺชุลฺลตาสุ\footnote{วิชฺชุตาสุ (สี, สฺยา, อิ) 2.นว วา สกฏสตานิ ทส วา สกฏสตานิ (สี)} นิจฺฉรนฺตีสุ อสนิยา ผลนฺติยา เนว ปสฺเสยฺย, น ปน สทฺทํ สุเณยฺยาติ.\csromanspacer{} กิญฺหิ ภนฺเต กริสฺสนฺติ ปญฺจ วา สกฏสตานิ ฉ วา สกฏสตานิ สตฺต วา สกฏสตานิ อฏฺฐ วา สกฏสตานิ นว วา สกฏสตานิ2, สกฏสหสฺสํ วา สกฏสตสหสฺสํ วา.\csromanspacer{} อถ โข เอตเทว ทุกฺกรตรญฺเจว ทุรภิสมฺภวตรญฺจ โย สญฺญี สมาโน ชาคโร เทเว วสฺสนฺเต เทเว คฬคฬายนฺเต วิชฺชุลฺลตาสุ นิจฺฉรนฺตีสุ อสนิยา ผลนฺติยา เนว ปสฺเสยฺย, น ปน สทฺทํ สุเณยฺยาติ.}
\prose{\csromanfolio{109}เอกมิทาหํ ปุกฺกุส สมยํ อาตุมายํ วิหรามิ ภุสาคาเร.\csromanspacer{} เตน โข ปน สมเยน เทเว วสฺสนฺเต เทเว คฬคฬายนฺเต วิชฺชุลฺลตาสุ นิจฺฉรนฺตีสุ อสนิยา ผลนฺติยา อวิทูเร ภุสาคารสฺส เทฺว กสฺสกา ภาตโร หตา จตฺตาโร จ พลิพทฺทา\footnote{พลิวทฺทา (สี, อิ)}.\csromanspacer{} อถ โข ปุกฺกุส อาตุมาย มหาชนกาโย นิกฺขมิตฺวา เยน เต เทฺว กสฺสกา ภาตโร หตา จตฺตาโร จ พลิพทฺทา เตนุปสงฺกมิ.\csromanspacer{} เตน โข ปนาหํ ปุกฺกุส สมเยน ภุสาคารา นิกฺขมิตฺวา ภุสาคารทฺวาเร อพฺโภกาเส จงฺกมามิ.\csromanspacer{} อถ โข ปุกฺกุส อญฺญตโร ปุริโส ตมฺหา มหาชนกายา เยนาหํ เตนุปสงฺกมิ, อุปสงฺกมิตฺวา มํ อภิวาเทตฺวา เอกมนฺตํ อฏฺฐาสิ, เอกมนฺตํ ฐิตํ โข อหํ ปุกฺกุส ตํ ปุริสํ เอตทโวจํ “กิํ นุ โข เอโส อาวุโส มหาชนกาโย สนฺนิปติโต”ติ.\csromanspacer{} อิทานิ ภนฺเต เทเว วสฺสนฺเต เทเว คฬคฬายนฺเต วิชฺชุลฺลตาสุ นิจฺฉรนฺตีสุ อสนิยา ผลนฺติยา เทฺว กสฺสกา ภาตโร หตา จตฺตาโร จ พลิพทฺทา, เอตฺเถโส มหาชนกาโย สนฺนิปติโต, ตฺวํ ปน ภนฺเต กฺว อโหสีติ. “อิเธว โข อหํ อาวุโส อโหสินฺ”ติ.\csromanspacer{} กิํ ปน ภนฺเต อทฺทสาติ. “น โข อหํ อาวุโส อทฺทสนฺ”ติ.\csromanspacer{} กิํ ปน ภนฺเต สทฺทํ อสฺโสสีติ. “น โข อหํ อาวุโส สทฺทํ อสฺโสสินฺ”ติ.\csromanspacer{} กิํ ปน ภนฺเต สุตฺโต อโหสีติ. “น โข อหํ อาวุโส สุตฺโต อโหสินฺ”ติ.\csromanspacer{} กิํ ปน ภนฺเต สญฺญี อโหสีติ. “เอวมาวุโส”ติ.\csromanspacer{} โส ตฺวํ ภนฺเต สญฺญี สมาโน ชาคโร เทเว วสฺสนฺเต เทเว \csromanfolio{110}คฬคฬายนฺเต วิชฺชุลฺลตาสุ นิจฺฉรนฺตีสุ อสนิยา ผลนฺติยา เนว อทฺทส, น ปน สทฺทํ อสฺโสสีติ. “เอวมาวุโส”ติ.}

\prose{อถ โข ปุกฺกุส ตสฺส ปุริสสฺส เอตทโหสิ “อจฺฉริยํ วต โภ อพฺภุตํ วต โภ, สนฺเตน วต โภ ปพฺพชิตา วิหาเรน วิหรนฺติ.\csromanspacer{} ยตฺร หิ นาม สญฺญี สมาโน ชาคโร เทเว วสฺสนฺเต เทเว คฬคฬายนฺเต วิชฺชุลฺลตาสุ นิจฺฉรนฺตีสุ อสนิยา ผลนฺติยา เนว ทกฺขติ, น ปน สทฺทํ โสสฺสตี”ติ\footnote{สุณิสฺสติ (สฺยา)}, มยิ อุฬารํ ปสาทํ ปเวเทตฺวา มํ อภิวาเทตฺวา ปทกฺขิณํ กตฺวา ปกฺกามีติ.}

\prose{เอวํ วุตฺเต ปุกฺกุโส มลฺลปุตฺโต ภควนฺตํ เอตทโวจ “เอสาหํ ภนฺเต โย เม อาฬาเร กาลาเม ปสาโท.\csromanspacer{} ตํ มหาวาเต วา โอผุณามิ สีฆโสตาย\footnote{สิงฺฆโสตาย (ก)} วา นทิยา ปวาเหมิ.\csromanspacer{} อภิกฺกนฺตํ ภนฺเต, อภิกฺกนฺตํ ภนฺเต, เสยฺยถาปิ ภนฺเต นิกฺกุชฺชิตํ วา อุกฺกุชฺเชยฺย, ปฏิจฺฉนฺนํ วา วิวเรยฺย, มูฬฺหสฺส วา มคฺคํ อาจิกฺเขยฺย, อนฺธกาเร วา เตลปชฺโชตํ ธาเรยฺย ‘จกฺขุมนฺโต รูปานิ ทกฺขนฺตี’ติ, เอวเมวํ ภควตา อเนกปริยาเยน ธมฺโม ปกาสิโต.\csromanspacer{} เอสาหํ ภนฺเต ภควนฺตํ สรณํ คจฺฉามิ ธมฺมญฺจ ภิกฺขุสํฆญฺจ, อุปาสกํ มํ ภควา ธาเรตุ อชฺชตคฺเค ปาณุเปตํ สรณํ คตนฺ”ติ.}

\proseitem{194}{อถ โข ปุกฺกุโส มลฺลปุตฺโต อญฺญตรํ ปุริสํ อามนฺเตสิ “อิงฺฆ เม ตฺวํ ภเณ สิงฺคีวณฺณํ ยุคมฏฺฐํ ธารณียํ อาหรา”ติ. “เอวํ ภนฺเต”ติ โข โส ปุริโส ปุกฺกุสสฺส มลฺลปุตฺตสฺส ปฏิสฺสุตฺวา ตํ สิงฺคีวณฺณํ ยุคมฏฺฐํ ธารณียํ อาหริ\footnote{อาหรสิ (ก)}.\csromanspacer{} อถ โข ปุกฺกุโส มลฺลปุตฺโต ตํ สิงฺคีวณฺณํ ยุคมฏฺฐํ ธารณียํ ภควโต อุปนาเมสิ “อิทํ ภนฺเต สิงฺคีวณฺณํ ยุคมฏฺฐํ ธารณียํ, ตํ เม ภควา ปฏิคฺคณฺหาตุ อนุกมฺปํ อุปาทายา”ติ.\csromanspacer{} เตน หิ ปุกฺกุส เอเกน มํ อจฺฉาเทหิ, เอเกน อานนฺทนฺติ. “เอวํ ภนฺเต”ติ โข ปุกฺกุโส มลฺลปุตฺโต ภควโต ปฏิสฺสุตฺวา เอเกน ภควนฺตํ อจฺฉาเทติ, เอเกน อายสฺมนฺตํ อานนฺทํ.\csromanspacer{} อถ โข ภควา ปุกฺกุสํ มลฺลปุตฺตํ ธมฺมิยา กถาย สนฺทสฺเสสิ สมาทเปสิ สมุตฺเตเชสิ สมฺปหํเสสิ.\csromanspacer{} อถ โข ปุกฺกุโส มลฺลปุตฺโต ภควตา ธมฺมิยา กถาย สนฺทสฺสิโต สมาทปิโต สมุตฺเตชิโต สมฺปหํสิโต อุฏฺฐายาสนา ภควนฺตํ อภิวาเทตฺวา ปทกฺขิณํ กตฺวา ปกฺกามิ.}

\csromanheader{2}{3. มหาปรินิพฺพานสุตฺต}
\proseitem{195}{\csromanfolio{111}อถ โข อายสฺมา อานนฺโท อจิรปกฺกนฺเต ปุกฺกุเส มลฺลปุตฺเต ตํ สิงฺคีวณฺณํ ยุคมฏฺฐํ ธารณียํ ภควโต กายํ อุปนาเมสิ.\csromanspacer{} ตํ ภควโต กายํ อุปนามิตํ หตจฺจิกํ วิย\footnote{วีตจฺจิกํวิย (สี, อิ)} ขายติ.\csromanspacer{} อถ โข อายสฺมา อานนฺโท ภควนฺตํ เอตทโวจ “อจฺฉริยํ ภนฺเต อพฺภุตํ ภนฺเต, ยาว ปริสุทฺโธ ภนฺเต ตถาคตสฺส ฉวิวณฺโณ ปริโยทาโต, อิทํ ภนฺเต สิงฺคีวณฺณํ ยุคมฏฺฐํ ธารณียํ ภควโต กายํ อุปนามิตํ หตจฺจิกํ วิย ขายตี”ติ.\csromanspacer{} เอวเมตํ อานนฺท เอวเมตํ อานนฺท ทฺวีสุ กาเลสุ อติวิย ตถาคตสฺส กาโย ปริสุทฺโธ โหติ ฉวิวณฺโณ ปริโยทาโต.\csromanspacer{} กตเมสุ ทฺวีสุ, ยญฺจ อานนฺท รตฺติํ ตถาคโต อนุตฺตรํ สมฺมาสมฺโพธิํ อภิสมฺพุชฺฌติ, ยญฺจ รตฺติํ อนุปาทิเสสาย นิพฺพานธาตุยา ปรินิพฺพายติ, อิเมสุ โข อานนฺท ทฺวีสุ กาเลสุ อติวิย ตถาคตสฺส กาโย ปริสุทฺโธ โหติ ฉวิวณฺโณ ปริโยทาโต.\csromanspacer{} อชฺช โข ปเนนนฺท รตฺติยา ปจฺฉิเม ยาเม กุสินารายํ อุปวตฺตเน มลฺลานํ สาลวเน อนฺตเรน\footnote{อนฺตเร (สฺยา)} ยมกสาลานํ ตถาคตสฺส ปรินิพฺพานํ ภวิสฺสติ\footnote{ภวิสฺสตีติ (ก)}.\csromanspacer{} อายามานนฺท, เยน กกุธา นที เตนุปสงฺกมิสฺสามาติ. “เอวํ ภนฺเต”ติ โข อายสฺมา อานนฺโท ภควโต ปจฺจสฺโสสิ.}

\csromangathameasure{สิงฺคีวณฺณํ ยุคมฏฺฐํ, ปุกฺกุโส อภิหารยิ. \\ เตน อจฺฉาทิโต สตฺถา, เหมวณฺโณ อโสภถาติ.}
\csromangathagroup{\csromangathastanza{สิงฺคีวณฺณํ ยุคมฏฺฐํ, ปุกฺกุโส อภิหารยิ. \\ เตน อจฺฉาทิโต สตฺถา, เหมวณฺโณ อโสภถาติ.}}

\proseitem{196}{อถ โข ภควา มหตา ภิกฺขุสํเฆน สทฺธิํ เยน กกุธา นที เตนุปสงฺกมิ, อุปสงฺกมิตฺวา กกุธํ นทิํ อชฺโฌคาเหตฺวา นฺหตฺวา จ ปิวิตฺวา จ ปจฺจุตฺตริตฺวา เยน อมฺพวนํ เตนุปสงฺกมิ, อุปสงฺกมิตฺวา อายสฺมนฺตํ จุนฺทกํ อามนฺเตสิ “อิงฺฆ เม ตฺวํ จุนฺทก จตุคฺคุณํ สํฆาฏิํ ปญฺญเปหิ, กิลนฺโตสฺมิ จุนฺทก นิปชฺชิสฺสามี”ติ.}

\prose{“เอวํ ภนฺเต”ติ โข อายสฺมา จุนฺทโก ภควโต ปฏิสฺสุตฺวา จตุคฺคุณํ สํฆาฏิํ ปญฺญเปสิ.\csromanspacer{} อถ โข ภควา ทกฺขิเณน ปสฺเสน สีหเสยฺยํ กปฺเปสิ ปาเท ปาทํ อจฺจาธาย สโต สมฺปชาโน อุฏฺฐานสญฺญํ มนสิกริตฺวา.\csromanspacer{} อายสฺมา ปน จุนฺทโก ตตฺเถว ภควโต ปุรโต นิสีทิ.}

\csromangathameasure{คนฺตฺวาน พุทฺโธ นทิกํ กกุธํ, \\ อจฺโฉทกํ สาตุทกํ วิปฺปสนฺนํ. \\ โอคาหิ สตฺถา อกิลนฺตรูโป, \\ ตถาคโต อปฺปฏิโม จ โลเก. \\ นฺหตฺวา จ ปิวิตฺวา จุทตาริ สตฺถา, \\ ปุรกฺขโต ภิกฺขุคณสฺส มชฺเฌ. \\ วตฺตา ปวตฺตา ภควา อิธ ธมฺเม, \\ อุปาคมิ อมฺพวนํ มเหสิ. \\ อามนฺตยิ จุนฺทกํ นาม ภิกฺขุํ, \\ จตุคฺคุณํ สนฺถร เม นิปชฺชํ. \\ โส โจทิโต ภาวิตตฺเตน จุนฺโท, \\ จตุคฺคุณํ สนฺถริ ขิปฺปเมว. \\ นิปชฺชิ สตฺถา อกิลนฺตรูโป, \\ จุนฺโทปิ ตตฺถ ปมุเข นิสีทีติ.}
\csromangathagroup{\csromangathastanza{\csromanfolio{112}คนฺตฺวาน พุทฺโธ นทิกํ กกุธํ, \\ อจฺโฉทกํ สาตุทกํ วิปฺปสนฺนํ. \\ โอคาหิ สตฺถา อกิลนฺตรูโป\footnote{สุกิลนฺตรูโป (สี, อิ)}, \\ ตถาคโต อปฺปฏิโม จ\footnote{สุกิลนฺตรูโป (สี, อิ)} โลเก.}\csromangathastanza{นฺหตฺวา จ ปิวิตฺวา จุทตาริ สตฺถา\footnote{ปิวิตฺวา จุนฺทเกน, ปิวิตฺวา จ อุตฺตริ (ก)}, \\ ปุรกฺขโต ภิกฺขุคณสฺส มชฺเฌ. \\ วตฺตา\footnote{ปิวิตฺวา จุนฺทเกน, ปิวิตฺวา จ อุตฺตริ (ก)} ปวตฺตา ภควา อิธ ธมฺเม, \\ อุปาคมิ อมฺพวนํ มเหสิ.}\csromangathastanza{อามนฺตยิ จุนฺทกํ นาม ภิกฺขุํ, \\ จตุคฺคุณํ สนฺถร เม นิปชฺชํ. \\ โส โจทิโต ภาวิตตฺเตน จุนฺโท, \\ จตุคฺคุณํ สนฺถริ ขิปฺปเมว. \\ นิปชฺชิ สตฺถา อกิลนฺตรูโป, \\ จุนฺโทปิ ตตฺถ ปมุเข\footnote{สมุเข (ก)} นิสีทีติ.}}

\proseitem{197}{อถ โข ภควา อายสฺมนฺตํ อานนฺทํ อามนฺเตสิ “สิยา โข\footnote{โย โข (ก)} ปนานนฺท จุนฺทสฺส กมฺมารปุตฺตสฺส โกจิ วิปฺปฏิสารํ อุปฺปาเทยฺย ‘ตสฺส เต อาวุโส จุนฺท อลาภา ตสฺส เต ทุลฺลทฺธํ, ยสฺส เต ตถาคโต ปจฺฉิมํ ปิณฺฑปาตํ ปริภุญฺชิตฺวา นิพฺพานํ’ติ.\csromanspacer{} จุนฺทสฺส อานนฺท กมฺมารปุตฺตสฺส เอวํ วิปฺปฏิสาโร ปฏิวิเนตพฺโพ ‘ตสฺส เต อาวุโส จุนฺท ลาภา ตสฺส เต สุลทฺธํ, ยสฺส เต ตถาคโต ปจฺฉิมํ ปิณฺฑปาตํ ปริภุญฺชิตฺวา นิพฺพานํ.\csromanspacer{} สมฺมุขา เมตํ อาวุโส จุนฺท ภควโต สุตํ สมฺมุขา ปฏิคฺคหิตํ.\csromanspacer{} เทฺว เม ปิณฺฑปาตา สมสมผลา\footnote{สมา สมผลา (ก)} สมวิปากา\footnote{สมสมวิปากา (สี, สฺยา, อิ)}, อติวิย อญฺเญหิ ปิณฺฑปาเตหิ มหปฺผลตรา จ มหานิสํสตรา จ.\csromanspacer{} กตเม เทฺว, ยญฺจ ปิณฺฑปาตํ ปริภุญฺชิตฺวา ตถาคโต อนุตฺตรํ สมฺมาสมฺโพธิํ อภิสมฺพุชฺฌติ, ยญฺจ ปิณฺฑปาตํ ปริภุญฺชิตฺวา ตถาคโต อนุปาทิเสสาย นิพฺพานธาตุยา ปรินิพฺพายติ, อิเม เทฺว ปิณฺฑปาตา สมสมผลา}

\vspace{\tipitakaparskip}
\csromancenter{มหาปรินิพฺพานสุตฺต สมวิปากา, อติวิย อญฺเญหิ ปิณฺฑปาเตหิ มหปฺผลตรา จ มหานิสํสตรา จ, อายุสํวตฺตนิกํ อายสฺมตา จุนฺเทน กมฺมารปุตฺเตน กมฺมํ อุปจิตํ, วณฺณสํวตฺตนิกํ อายสฺมตา จุนฺเทน กมฺมารปุตฺเตน กมฺมํ อุปจิตํ, สุขสํวตฺตนิกํ อายสฺมตา จุนฺเทน กมฺมารปุตฺเตน กมฺมํ อุปจิตํ, ยสสํวตฺตนิกํ อายสฺมตา จุนฺเทน กมฺมารปุตฺเตน กมฺมํ อุปจิตํ, สคฺคสํวตฺตนิกํ อายสฺมตา จุนฺเทน กมฺมารปุตฺเตน กมฺมํ อุปจิตํ, อาธิปเตยฺยสํวตฺตนิกํ อายสฺมตา จุนฺเทน กมฺมารปุตฺเตน กมฺมํ อุปจิตนฺ’ติ.\csromanspacer{} จุนฺทสฺส อานนฺท กมฺมารปุตฺตสฺส เอวํ วิปฺปฏิสาโร ปฏิวิเนตพฺโพ”ติ.\csromanspacer{} อถ โข ภควา เอตมตฺถํ วิทิตฺวา ตายํ เวลายํ อิมํ อุทานํ อุทาเนสิ\csromandash{}}
\vspace{0.5\baselineskip}
\csromangathameasure{“ททโต ปุญฺญํ ปวฑฺฒติ, \\ สํยมโต เวรํ น จียติ. \\ กุสโล จ ชหาติ ปาปกํ, \\ ราคโทสโมหกฺขยา สนิพฺพุโต”ติ.}
\csromangathagroup{\csromangathastanza{\csromanfolio{113}“ททโต ปุญฺญํ ปวฑฺฒติ, \\ สํยมโต เวรํ น จียติ. \\ กุสโล จ ชหาติ ปาปกํ, \\ ราคโทสโมหกฺขยา สนิพฺพุโต”ติ.}}

\csromanheader{2}{จตุตฺโถ ภาณวาโร.}
\csromansectionrule
\csromanmatikaanchor{mtk.52}
\csromantocmarkref{section}{ยมกสาลา}{mtk.52}
\bookmark[dest={mtk.52},level=1]{ยมกสาลา}
\csromanheader{1}{\textbf{ยมกสาลา}}
\proseitem{198}{อถ โข ภควา อายสฺมนฺตํ อานนฺทํ อามนฺเตสิ “อายามานนฺท, เยน หิรญฺญวติยา นทิยา ปาริมํ ตีรํ, เยน กุสินารา อุปวตฺตนํ มลฺลานํ สาลวนํ เตนุปสงฺกมิสฺสามา”ติ. “เอวํ ภนฺเต”ติ โข อายสฺมา อานนฺโท ภควโต ปจฺจสฺโสสิ.\csromanspacer{} อถ โข ภควา มหตา ภิกฺขุสํเฆน สทฺธิํ เยน หิรญฺญวติยา นทิยา ปาริมํ ตีรํ, เยน กุสินารา อุปวตฺตนํ มลฺลานํ สาลวนํ เตนุปสงฺกมิ, อุปสงฺกมิตฺวา อายสฺมนฺตํ อานนฺทํ อามนฺเตสิ “อิงฺฆ เม ตฺวํ อานนฺท อนฺตเรน ยมกสาลานํ อุตฺตรสีสกํ มญฺจกํ ปญฺญเปหิ, กิลนฺโตสฺมิ อานนฺท นิปชฺชิสฺสามี”ติ. “เอวํ ภนฺเต”ติ โข อายสฺมา อานนฺโท ภควโต ปฏิสฺสุตฺวา อนฺตเรน ยมกสาลานํ อุตฺตรสีสกํ มญฺจกํ ปญฺญเปสิ.\csromanspacer{} อถ โข ภควา ทกฺขิเณน ปสฺเสน สีหเสยฺยํ กปฺเปสิ ปาเท ปาทํ อจฺจาธาย สโต สมฺปชาโน.}

\prose{\csromanfolio{114}เตน โข ปน สมเยน ยมกสาลา สพฺพผาลิผุลฺลา โหนฺติ อกาลปุปฺเผหิ, เต ตถาคตสฺส สรีรํ โอกิรนฺติ อชฺโฌกิรนฺติ อภิปฺปกิรนฺติ ตถาคตสฺส ปูชาย.\csromanspacer{} ทิพฺพานิปิ มนฺทารวปุปฺผานิ อนฺตลิกฺขา ปปตนฺติ, ตานิ ตถาคตสฺส สรีรํ โอกิรนฺติ อชฺโฌกิรนฺติ อภิปฺปกิรนฺติ ตถาคตสฺส ปูชาย.\csromanspacer{} ทิพฺพานิปิ จนฺทนจุณฺณานิ อนฺตลิกฺขา ปปตนฺติ, ตานิ ตถาคตสฺส สรีรํ โอกิรนฺติ อชฺโฌกิรนฺติ อภิปฺปกิรนฺติ ตถาคตสฺส ปูชาย.\csromanspacer{} ทิพฺพานิปิ ตูริยานิ อนฺตลิกฺเข วชฺชนฺติ ตถาคตสฺส ปูชาย.\csromanspacer{} ทิพฺพานิปิ สํคีตานิ อนฺตลิกฺเข วตฺตนฺติ ตถาคตสฺส ปูชาย.}

\proseitem{199}{อถ โข ภควา อายสฺมนฺตํ อานนฺทํ อามนฺเตสิ\csromandash{} สพฺพผาลิผุลฺลา โข อานนฺท ยมกสาลา อกาลปุปฺเผหิ, เต ตถาคตสฺส สรีรํ โอกิรนฺติ อชฺโฌกิรนฺติ อภิปฺปกิรนฺติ ตถาคตสฺส ปูชาย.\csromanspacer{} ทิพฺพานิปิ มนฺทารวปุปฺผานิ อนฺตลิกฺขา ปปตนฺติ, ตานิ ตถาคตสฺส สรีรํ โอกิรนฺติ อชฺโฌกิรนฺติ อภิปฺปกิรนฺติ ตถาคตสฺส ปูชาย.\csromanspacer{} ทิพฺพานิปิ จนฺทนจุณฺณานิ อนฺตลิกฺขา ปปตนฺติ, ตานิ ตถาคตสฺส สรีรํ โอกิรนฺติ อชฺโฌกิรนฺติ อภิปฺปกิรนฺติ ตถาคตสฺส ปูชาย.\csromanspacer{} ทิพฺพานิปิ ตูริยานิ อนฺตลิกฺเข วชฺชนฺติ ตถาคตสฺส ปูชาย.\csromanspacer{} ทิพฺพานิปิ สํคีตานิ อนฺตลิกฺเข วตฺตนฺติ ตถาคตสฺส ปูชาย.\csromanspacer{} น โข อานนฺท เอตฺตาวตา ตถาคโต สกฺกโต วา โหติ ครุกโต วา มานิโต วา ปูชิโต วา อปจิโต วา.\csromanspacer{} โย โข อานนฺท ภิกฺขุ วา ภิกฺขุนี วา อุปาสโก วา อุปาสิกา วา ธมฺมานุธมฺมปฺปฏิปนฺโน วิหรติ สามีจิปฺปฏิปนฺโน อนุธมฺมจารี, โส ตถาคตํ สกฺกโรติ ครุํ กโรติ มาเนติ ปูเชติ อปจิยติ\footnote{อิทํ ปทํ สี-สฺยา-อิ-โปตฺถเกสุ น ทิสฺสติ,} ปรมาย ปูชาย.\csromanspacer{} ตสฺมาติหานนฺท ธมฺมานุธมฺมปฺปฏิปนฺนา วิหริสฺสาม สามีจิปฺปฏิปนฺนา อนุธมฺมจาริโนติ.\csromanspacer{} เอวญฺหิ โว อานนฺท สิกฺขิตพฺพนฺติ.}

\csromanmatikaanchor{mtk.53}
\csromantocmarkref{section}{อุปวาณตฺเถร}{mtk.53}
\bookmark[dest={mtk.53},level=1]{อุปวาณตฺเถร}
\csromanheader{1}{\textbf{อุปวาณตฺเถร}}
\proseitem{200}{เตน โข ปน สมเยน อายสฺมา อุปวาโณ ภควโต ปุรโต ฐิโต โหติ ภควนฺตํ พีชยมาโน.\csromanspacer{} อถ โข ภควา อายสฺมนฺตํ อุปวาณํ อปสาเรสิ “อเปหิ ภิกฺขุ, มา เม ปุรโต อฏฺฐาสี”ติ.\csromanspacer{} อถ โข อายสฺมโต อานนฺทสฺส เอตทโหสิ “อยํ}

\vspace{\tipitakaparskip}
\csromancenter{มหาปรินิพฺพานสุตฺต โข อายสฺมา อุปวาโณ ทีฆรตฺตํ ภควโต อุปฏฺฐาโก สนฺติกาวจโร สมีปจารี, อถ จ ปน ภควา ปจฺฉิเม กาเล อายสฺมนฺตํ อุปวาณํ อปสาเรติ ‘อเปหิ ภิกฺขุ, มา เม ปุรโต อฏฺฐาสี’ติ.\csromanspacer{} โก นุ โข เหตุ โก ปจฺจโย, ยํ ภควา อายสฺมนฺตํ อุปวาณํ อปสาเรติ ‘อเปหิ ภิกฺขุ, มา เม ปุรโต อฏฺฐาสี’ติ”.\csromanspacer{} อถ โข อายสฺมา อานนฺโท ภควนฺตํ เอตทโวจ “อยํ ภนฺเต อายสฺมา อุปวาโณ ทีฆรตฺตํ ภควโต อุปฏฺฐาโก สนฺติกาวจโร สมีปจารี, อถ จ ปน ภควา ปจฺฉิเม กาเล อายสฺมนฺตํ อุปวาณํ อปสาเรติ ‘อเปหิ ภิกฺขุ, มา เม ปุรโต อฏฺฐาสี’ติ, โก นุ โข ภนฺเต เหตุ โก ปจฺจโย, ยํ ภควา อายสฺมนฺตํ อุปวาณํ อปสาเรติ ‘อเปหิ ภิกฺขุ, มา เม ปุรโต อฏฺฐาสี’ติ”.เยภุยฺเยน อานนฺท ทสสุ โลกธาตูสุ เทวตา สนฺนิปติตา ตถาคตํ ทสฺสนาย, ยาวตา อานนฺท กุสินารา อุปวตฺตนํ มลฺลานํ สาลวนํ สมนฺตโต ทฺวาทส โยชนานิ, นตฺถิ โส ปเทโส วาลคฺคโกฏินิตุทนมตฺโตปิ มเหสกฺขาหิ เทวตาหิ อปฺผุโฏ.\csromanspacer{} เทวตา อานนฺท อุชฺฌายนฺติ “ทูรา จ วตมฺห อาคตา ตถาคตํ ทสฺสนาย, กทาจิ กรหจิ ตถาคตา โลเก อุปฺปชฺชนฺติ อรหนฺโต สมฺมาสมฺพุทฺธา, อชฺเชว รตฺติยา ปจฺฉิเม ยาเม ตถาคตสฺส ปรินิพฺพานํ ภวิสฺสติ, อยํ จ มเหสกฺโข ภิกฺขุ ภควโต ปุรโต ฐิโต โอวาเรนฺโต, น มยํ ลภาม ปจฺฉิเม กาเล ตถาคตํ ทสฺสนายา”ติ.}
\proseitem{201}{\csromanfolio{115}กถํภูตา ปน ภนฺเต ภควา เทวตา มนสิกโรตีติ\footnote{มนสิ กโรนฺตีติ (สฺยา, ก)}.\csromanspacer{} สนฺตานนฺท เทวตา อากาเส ปถวีสญฺญินิโย เกเส ปกิริย กนฺทนฺติ, พาหา ปคฺคยฺห กนฺทนฺติ, ฉินฺนปาตํ ปปตนฺติ\footnote{ฉินฺนํปาทํวิย ปปตนฺติ (สฺยา)}, อาวฏฺฏนฺติ, วิวฏฺฏนฺติ, “อติขิปฺปํ ภควา ปรินิพฺพายิสฺสติ, อติขิปฺปํ สุคโต ปรินิพฺพายิสฺสติ, อติขิปฺปํ จกฺขุํ\footnote{จกฺขุมา (สฺยา, ก)} โลเก อนฺตรธายิสฺสตี”ติ.}

\prose{สนฺตานนฺท เทวตา ปถวิยํ ปถวีสญฺญินิโย เกเส ปกิริย กนฺทนฺติ, พาหา ปคฺคยฺห กนฺทนฺติ, ฉินฺนปาตํ ปปตนฺติ, อาวฏฺฏนฺติ, วิวฏฺฏนฺติ, “อติขิปฺปํ ภควา ปรินิพฺพายิสฺสติ, อติขิปฺปํ สุคโต ปรินิพฺพายิสฺสติ, อติขิปฺปํ จกฺขุํ โลเก อนฺตรธายิสฺสตี”ติ.}

\prose{\csromanfolio{116}ยา ปน ตา เทวตา วีตราคา, ตา สตา สมฺปชานา อธิวาเสนฺติ “อนิจฺจา สงฺขารา, ตํ กุเตตฺถ ลพฺภา”ติ.}

\csromanmatikaanchor{mtk.54}
\csromantocmarkref{section}{จตุสํเวชนียฐาน}{mtk.54}
\bookmark[dest={mtk.54},level=1]{จตุสํเวชนียฐาน}
\csromanheader{1}{\textbf{จตุสํเวชนียฐาน}}
\proseitem{202}{ปุพฺเพ ภนฺเต ทิสาสุ วสฺสํวุฏฺฐา\footnote{วสฺสํวุตฺถา (สี, สฺยา, กํ, อิ)} ภิกฺขู อาคจฺฉนฺติ ตถาคตํ ทสฺสนาย, เต มยํ ลภาม มโนภาวนีเย ภิกฺขู ทสฺสนาย, ลภาม ปยิรุปาสนาย.\csromanspacer{} ภควโต ปน มยํ ภนฺเต อจฺจเยน น ลภิสฺสาม มโนภาวนีเย ภิกฺขู ทสฺสนาย, น ลภิสฺสาม ปยิรุปาสนายาติ.}

\prose{จตฺตาริมานิ อานนฺท สทฺธสฺส กุลปุตฺตสฺส ทสฺสนียานิ สํเวชนียานิ ฐานานิ.\csromanspacer{} กตมานิ จตฺตาริ, “อิธ ตถาคโต ชาโต”ติ อานนฺท สทฺธสฺส กุลปุตฺตสฺส ทสฺสนียํ สํเวชนียํ ฐานํ. “อิธ ตถาคโต อนุตฺตรํ สมฺมาสมฺโพธิํ อภิสมฺพุทฺโธ”ติ อานนฺท สทฺธสฺส กุลปุตฺตสฺส ทสฺสนียํ สํเวชนียํ ฐานํ. “อิธ ตถาคเตน อนุตฺตรํ ธมฺมจกฺกํ ปวตฺติตนฺ”ติ อานนฺท สทฺธสฺส กุลปุตฺตสฺส ทสฺสนียํ สํเวชนียํ ฐานํ. “อิธ ตถาคโต อนุปาทิเสสาย นิพฺพานธาตุยา นิพฺพานํ”ติ อานนฺท สทฺธสฺส กุลปุตฺตสฺส ทสฺสนียํ สํเวชนียํ ฐานํ.\csromanspacer{} อิมานิ โข อานนฺท จตฺตาริ สทฺธสฺส กุลปุตฺตสฺส ทสฺสนียานิ สํเวชนียานิ ฐานานิ.}

\prose{อาคมิสฺสนฺติ โข อานนฺท สทฺธา ภิกฺขู ภิกฺขุนิโย อุปาสกา อุปาสิกาโย “อิธ ตถาคโต ชาโต”ติปิ, “อิธ ตถาคโต อนุตฺตรํ สมฺมาสมฺโพธิํ อภิสมฺพุทฺโธ”ติปิ, “อิธ ตถาคเตน อนุตฺตรํ ธมฺมจกฺกํ ปวตฺติตนฺ”ติปิ, “อิธ ตถาคโต อนุปาทิเสสาย นิพฺพานธาตุยา นิพฺพานํ”ติปิ.\csromanspacer{} เย หิ เกจิ อานนฺท เจติยจาริกํ อาหิณฺฑนฺตา ปสนฺนจิตฺตา กาลํ กริสฺสนฺติ, สพฺเพ เต กายสฺส เภทา ปรํ มรณา สุคติํ สคฺคํ โลกํ อุปปชฺชิสฺสนฺตีติ.}

\csromanmatikaanchor{mtk.55}
\csromantocmarkref{section}{อานนฺทปุจฺฉากถา}{mtk.55}
\bookmark[dest={mtk.55},level=1]{อานนฺทปุจฺฉากถา}
\csromanheader{1}{\textbf{อานนฺทปุจฺฉากถา}}
\proseitem{203}{กถํ มยํ ภนฺเต มาตุคาเม ปฏิปชฺชามาติ.\csromanspacer{} อทสฺสนํ อานนฺทาติ.\csromanspacer{} ทสฺสเน ภควา สติ กถํ ปฏิปชฺชิตพฺพนฺติ.\csromanspacer{} อนาลาโป}

\vspace{\tipitakaparskip}
\csromancenter{มหาปรินิพฺพานสุตฺต อานนฺทาติ.\csromanspacer{} อาลปนฺเตน ปน ภนฺเต กถํ ปฏิปชฺชิตพฺพนฺติ.\csromanspacer{} สติ อานนฺท อุปฏฺฐาเปตพฺพาติ.}
\proseitem{204}{\csromanfolio{117}กถํ มยํ ภนฺเต ตถาคตสฺส สรีเร ปฏิปชฺชามาติ.\csromanspacer{} อพฺยาวฏา ตุเมฺห อานนฺท โหถ ตถาคตสฺส สรีรปูชาย.\csromanspacer{} อิงฺฆ ตุเมฺห อานนฺท สารตฺเถ ฆฏถ อนุยุญฺชถ\footnote{สทตฺเถ อนุยุญฺชถ (สี, สฺยา), สทตฺถํ อนุยุญฺชถ (อิ), สารตฺเถ อนุยุญฺชถ (ก)} สารตฺเถ อปฺปมตฺตา อาตาปิโน ปหิตตฺตา วิหรถ.\csromanspacer{} สนฺตานนฺท ขตฺติยปณฺฑิตาปิ พฺราหฺมณปณฺฑิตาปิ คหปติปณฺฑิตาปิ ตถาคเต อภิปฺปสนฺนา, เต ตถาคตสฺส สรีรปูชํ กริสฺสนฺตีติ.}

\proseitem{205}{กถํ ปน ภนฺเต ตถาคตสฺส สรีเร ปฏิปชฺชิตพฺพนฺติ.\csromanspacer{} ยถา โข อานนฺท รญฺโญ จกฺกวตฺติสฺส สรีเร ปฏิปชฺชนฺติ, เอวํ ตถาคตสฺส สรีเร ปฏิปชฺชิตพฺพนฺติ.\csromanspacer{} กถํ ปน ภนฺเต รญฺโญ จกฺกวตฺติสฺส สรีเร ปฏิปชฺชนฺตีติ.\csromanspacer{} รญฺโญ อานนฺท จกฺกวตฺติสฺส สรีรํ อหเตน วตฺเถน เวเฐนฺติ, อหเตน วตฺเถน เวเฐตฺวา วิหเตน กปฺปาเสน เวเฐนฺติ, วิหเตน กปฺปาเสน เวเฐตฺวา อหเตน วตฺเถน เวเฐนฺติ, เอเตนุปาเยน ปญฺจหิ ยุคสเตหิ รญฺโญ จกฺกวตฺติสฺส สรีรํ\footnote{สรีเร (สฺยา, ก)} เวเฐตฺวา อายสาย เตลโทณิยา ปกฺขิปิตฺวา อญฺญิสฺสา อายสาย โทณิยา ปฏิกุชฺชิตฺวา สพฺพคนฺธานํ จิตกํ กริตฺวา รญฺโญ จกฺกวตฺติสฺส สรีรํ ฌาเปนฺติ, จาตุมหาปเถ\footnote{จาตุมฺมหาปเถ (สี, สฺยา, กํ, อิ)} รญฺโญ จกฺกวตฺติสฺส ถูปํ กโรนฺติ.\csromanspacer{} เอวํ โข อานนฺท รญฺโญ จกฺกวตฺติสฺส สรีเร ปฏิปชฺชนฺติ.\csromanspacer{} ยถา โข อานนฺท รญฺโญ จกฺกวตฺติสฺส สรีเร ปฏิปชฺชนฺติ, เอวํ ตถาคตสฺส สรีเร ปฏิปชฺชิตพฺพํ, จาตุมหาปเถ ตถาคตสฺส ถูโป กาตพฺโพ, ตตฺถ เย มาลํ วา คนฺธํ วา จุณฺณกํ\footnote{วณฺณกํ (สี, อิ)} วา อาโรเปสฺสนฺติ วา อภิวาเทสฺสนฺติ วา จิตฺตํ วา ปสาเทสฺสนฺติ, เตสํ ตํ ภวิสฺสติ ทีฆรตฺตํ หิตาย สุขาย.}

\csromanmatikaanchor{mtk.56}
\csromantocmarkref{section}{ถูปารหปุคฺคล}{mtk.56}
\bookmark[dest={mtk.56},level=1]{ถูปารหปุคฺคล}
\csromanheader{1}{\textbf{ถูปารหปุคฺคล}}
\proseitem{206}{จตฺตาโรเม อานนฺท ถูปารหา.\csromanspacer{} กตเม จตฺตาโร.\csromanspacer{} ตถาคโต อรหํ สมฺมาสมฺพุทฺโธ ถูปารโห, ปจฺเจกสมฺพุทฺโธ ถูปารโห, ตถาคตสฺส สาวโก ถูปารโห, ราชา จกฺกวตฺตี\footnote{จกฺกวตฺติ (สฺยา, ก)} ถูปารโหติ.}

\prose{\csromanfolio{118}กิญฺจานนฺท อตฺถวสํ ปฏิจฺจ ตถาคโต อรหํ สมฺมาสมฺพุทฺโธ ถูปารโห. “อยํ ตสฺส ภควโต อรหโต สมฺมาสมฺพุทฺธสฺส ถูโป”ติ อานนฺท พหุชนา จิตฺตํ ปสาเทนฺติ, เต ตตฺถ จิตฺตํ ปสาเทตฺวา กายสฺส เภทา ปรํ มรณา สุคติํ สคฺคํ โลกํ อุปปชฺชนฺติ.\csromanspacer{} อิทํ โข อานนฺท อตฺถวสํ ปฏิจฺจ ตถาคโต อรหํ สมฺมาสมฺพุทฺโธ ถูปารโห.}

\prose{กิญฺจานนฺท อตฺถวสํ ปฏิจฺจ ปจฺเจกสมฺพุทฺโธ ถูปารโห. “อยํ ตสฺส ภควโต ปจฺเจกสมฺพุทฺธสฺส ถูโป”ติ อานนฺท พหุชนา จิตฺตํ ปสาเทนฺติ, เต ตตฺถ จิตฺตํ ปสาเทตฺวา กายสฺส เภทา ปรํ มรณา สุคติํ สคฺคํ โลกํ อุปปชฺชนฺติ.\csromanspacer{} อิทํ โข อานนฺท อตฺถวสํ ปฏิจฺจ ปจฺเจกสมฺพุทฺโธ ถูปารโห.}

\prose{กิญฺจานนฺท อตฺถวสํ ปฏิจฺจ ตถาคตสฺส สาวโก ถูปารโห. “อยํ ตสฺส ภควโต อรหโต สมฺมาสมฺพุทฺธสฺส สาวกสฺส ถูโป”ติ อานนฺท พหุชนา จิตฺตํ ปสาเทนฺติ, เต ตตฺถ จิตฺตํ ปสาเทตฺวา กายสฺส เภทา ปรํ มรณา สุคติํ สคฺคํ โลกํ อุปปชฺชนฺติ.\csromanspacer{} อิทํ โข อานนฺท อตฺถวสํ ปฏิจฺจ ตถาคตสฺส สาวโก ถูปารโห.}

\prose{กิญฺจานนฺท อตฺถวสํ ปฏิจฺจ ราชา จกฺกวตฺตี ถูปารโห. “อยํ ตสฺส ธมฺมิกสฺส ธมฺมรญฺโญ ถูโป”ติ อานนฺท พหุชนา จิตฺตํ ปสาเทนฺติ, เต ตตฺถ จิตฺตํ ปสาเทตฺวา กายสฺส เภทา ปรํ มรณา สุคติํ สคฺคํ โลกํ อุปปชฺชนฺติ.\csromanspacer{} อิทํ โข อานนฺท อตฺถวสํ ปฏิจฺจ ราชา จกฺกวตฺตี ถูปารโห.\csromanspacer{} อิเม โข อานนฺท จตฺตาโร ถูปารหาติ.}

\csromanmatikaanchor{mtk.57}
\csromantocmarkref{section}{อานนฺท-อจฺฉริยธมฺม}{mtk.57}
\bookmark[dest={mtk.57},level=1]{อานนฺท-อจฺฉริยธมฺม}
\csromanheader{1}{\textbf{อานนฺท อจฺฉริยธมฺม}}
\proseitem{207}{อถ โข อายสฺมา อานนฺโท วิหารํ ปวิสิตฺวา กปิสีสํ อาลมฺพิตฺวา โรทมาโน อฏฺฐาสิ “อหญฺจ วตมฺหิ เสโข สกรณีโย, สตฺถุ จ เม ปรินิพฺพานํ ภวิสฺสติ, โย มม อนุกมฺปโก”ติ.\csromanspacer{} อถ โข ภควา ภิกฺขู อามนฺเตสิ “กหํ นุ โข ภิกฺขเว อานนฺโท”ติ.\csromanspacer{} เอโส ภนฺเต อายสฺมา อานนฺโท วิหารํ ปวิสิตฺวา กปิสีสํ อาลมฺพิตฺวา โรทมาโน ฐิโต “อหญฺจ วตมฺหิ เสโข สกรณีโย, สตฺถุ จ เม ปรินิพฺพานํ ภวิสฺสติ, โย มม อนุกมฺปโก”ติ.\csromanspacer{} อถ โข ภควา อญฺญตรํ ภิกฺขุํ อามนฺเตสิ “เอหิ ตฺวํ ภิกฺขุ มม วจเนน อานนฺทํ}

\vspace{\tipitakaparskip}
\csromancenter{มหาปรินิพฺพานสุตฺต อามนฺเตหิ ‘สตฺถา ตํ อาวุโส อานนฺท อามนฺเตตี’ติ”. “เอวํ ภนฺเต”ติ โข โส ภิกฺขุ ภควโต ปฏิสฺสุตฺวา เยนายสฺมา อานนฺโท เตนุปสงฺกมิ, อุปสงฺกมิตฺวา อายสฺมนฺตํ อานนฺทํ เอตทโวจ “สตฺถา ตํ อาวุโส อานนฺท อามนฺเตตี”ติ. “เอวมาวุโส”ติ โข อายสฺมา อานนฺโท ตสฺส ภิกฺขุโน ปฏิสฺสุตฺวา เยน ภควา เตนุปสงฺกมิ, อุปสงฺกมิตฺวา ภควนฺตํ อภิวาเทตฺวา เอกมนฺตํ นิสีทิ, เอกมนฺตํ นิสินฺนํ โข อายสฺมนฺตํ อานนฺทํ ภควา เอตทโวจ “อลํ อานนฺท มา โสจิ มา ปริเทวิ, นนุ เอตํ อานนฺท มยา ปฏิกจฺเจว อกฺขาตํ ‘สพฺเพเหว ปิเยหิ มนาเปหิ นานาภาโว วินาภาโว อญฺญถาภาโว’, ตํ กุเตตฺถ อานนฺท ลพฺภา, ‘ยํ ตํ ชาตํ ภูตํ สงฺขตํ ปโลกธมฺมํ, ตํ วต ตถาคตสฺสาปิ สรีรํ มา ปลุชฺชี’ติ, เนตํ ฐานํ วิชฺชติ.\csromanspacer{} ทีฆรตฺตํ โข เต อานนฺท ตถาคโต ปจฺจุปฏฺฐิโต เมตฺเตน กายกมฺเมน หิเตน สุเขน อทฺวเยน อปฺปมาเณน, เมตฺเตน วจีกมฺเมน หิเตน สุเขน อทฺวเยน อปฺปมาเณน, เมตฺเตน มโนกมฺเมน หิเตน สุเขน อทฺวเยน อปฺปมาเณน.\csromanspacer{} กตปุญฺโญสิ ตฺวํ อานนฺท, ปธานมนุยุญฺช ขิปฺปํ โหหิสิ อนาสโว”ติ.}
\proseitem{208}{\csromanfolio{119}อถ โข ภควา ภิกฺขู อามนฺเตสิ “เยปิ เต ภิกฺขเว อเหสุํ อตีตมทฺธานํ อรหนฺโต สมฺมาสมฺพุทฺธา, เตสมฺปิ ภควนฺตานํ เอตปฺปรมาเยว อุปฏฺฐากา อเหสุํ, เสยฺยถาปิ มยฺหํ อานนฺโท.\csromanspacer{} เยปิ เต ภิกฺขเว ภวิสฺสนฺติ อนาคตมทฺธานํ อรหนฺโต สมฺมาสมฺพุทฺธา, เตสมฺปิ ภควนฺตานํ เอตปฺปรมาเยว อุปฏฺฐากา ภวิสฺสนฺติ, เสยฺยถาปิ มยฺหํ อานนฺโท.\csromanspacer{} ปณฺฑิโต ภิกฺขเว อานนฺโท, เมธาวี ภิกฺขเว อานนฺโท, ชานาติ ‘อยํ กาโล ตถาคตํ ทสฺสนาย อุปสงฺกมิตุํ ภิกฺขูนํ, อยํ กาโล ภิกฺขุนีนํ, อยํ กาโล อุปาสกานํ, อยํ กาโล อุปาสิกานํ, อยํ กาโล รญฺโญ ราชมหามตฺตานํ ติตฺถิยานํ ติตฺถิยสาวกานนฺ”ติ.}

\proseitem{209}{จตฺตาโรเม ภิกฺขเว อจฺฉริยา อพฺภุตา ธมฺมา\footnote{อพฺภุตธมฺมา (สฺยา, ก)} อานนฺเท.\csromanspacer{} กตเม จตฺตาโร.\csromanspacer{} สเจ ภิกฺขเว ภิกฺขุปริสา อานนฺทํ ทสฺสนาย อุปสงฺกมติ, ทสฺสเนน สา อตฺตมนา โหติ.\csromanspacer{} ตตฺร เจ อานนฺโท \csromanfolio{120}ธมฺมํ ภาสติ, ภาสิเตนปิ สา อตฺตมนา โหติ.\csromanspacer{} อติตฺตาว ภิกฺขเว ภิกฺขุปริสา โหติ.\csromanspacer{} อถ โข อานนฺโท ตุณฺหี โหติ.\csromanspacer{} สเจ ภิกฺขเว ภิกฺขุนีปริสา อานนฺทํ ทสฺสนาย อุปสงฺกมติ, ทสฺสเนน สา อตฺตมนา โหติ.\csromanspacer{} ตตฺร เจ อานนฺโท ธมฺมํ ภาสติ, ภาสิเตนปิ สา อตฺตมนา โหติ.\csromanspacer{} อติตฺตาว ภิกฺขเว ภิกฺขุนีปริสา โหติ.\csromanspacer{} อถ โข อานนฺโท ตุณฺหี โหติ.\csromanspacer{} สเจ ภิกฺขเว อุปาสกปริสา อานนฺทํ ทสฺสนาย อุปสงฺกมติ, ทสฺสเนน สา อตฺตมนา โหติ.\csromanspacer{} ตตฺร เจ อานนฺโท ธมฺมํ ภาสติ, ภาสิเตนปิ สา อตฺตมนา โหติ.\csromanspacer{} อติตฺตาว ภิกฺขเว อุปาสกปริสา โหติ.\csromanspacer{} อถ โข อานนฺโท ตุณฺหี โหติ.\csromanspacer{} สเจ ภิกฺขเว อุปาสิกาปริสา อานนฺทํ ทสฺสนาย อุปสงฺกมติ, ทสฺสเนน สา อตฺตมนา โหติ.\csromanspacer{} ตตฺร เจ อานนฺโท ธมฺมํ ภาสติ, ภาสิเตนปิ สา อตฺตมนา โหติ.\csromanspacer{} อติตฺตาว ภิกฺขเว อุปาสิกาปริสา โหติ.\csromanspacer{} อถ โข อานนฺโท ตุณฺหี โหติ.\csromanspacer{} อิเม โข ภิกฺขเว จตฺตาโร อจฺฉริยา อพฺภุตา ธมฺมา อานนฺเท.}

\prose{จตฺตาโรเม ภิกฺขเว อจฺฉริยา อพฺภุตา ธมฺมา รญฺเญ จกฺกวตฺติมฺหิ.\csromanspacer{} กตเม จตฺตาโร.\csromanspacer{} สเจ ภิกฺขเว ขตฺติยปริสา ราชานํ จกฺกวตฺติํ ทสฺสนาย อุปสงฺกมติ, ทสฺสเนน สา อตฺตมนา โหติ.\csromanspacer{} ตตฺร เจ ราชา จกฺกวตฺตี ภาสติ, ภาสิเตนปิ สา อตฺตมนา โหติ.\csromanspacer{} อติตฺตาว ภิกฺขเว ขตฺติยปริสา โหติ.\csromanspacer{} อถ โข ราชา จกฺกวตฺตี ตุณฺหี โหติ.\csromanspacer{} สเจ ภิกฺขเว พฺราหฺมณปริสา ฯเปฯ คหปติปริสา ฯเปฯ สมณปริสา ราชานํ จกฺกวตฺติํ ทสฺสนาย อุปสงฺกมติ, ทสฺสเนน สา อตฺตมนา โหติ.\csromanspacer{} ตตฺร เจ ราชา จกฺกวตฺตี ภาสติ, ภาสิเตนปิ สา อตฺตมนา โหติ.\csromanspacer{} อติตฺตาว ภิกฺขเว สมณปริสา โหติ.\csromanspacer{} อถ โข ราชา จกฺกวตฺตี ตุณฺหี โหติ.\csromanspacer{} เอวเมว โข ภิกฺขเว จตฺตาโรเม อจฺฉริยา อพฺภุตา ธมฺมา อานนฺเท.\csromanspacer{} สเจ ภิกฺขเว ภิกฺขุปริสา อานนฺทํ ทสฺสนาย อุปสงฺกมติ, ทสฺสเนน สา อตฺตมนา โหติ.\csromanspacer{} ตตฺร เจ อานนฺโท ธมฺมํ ภาสติ, ภาสิเตนปิ สา อตฺตมนา โหติ.\csromanspacer{} อติตฺตาว ภิกฺขเว ภิกฺขุปริสา โหติ.\csromanspacer{} อถ โข อานนฺโท ตุณฺหี โหติ.\csromanspacer{} สเจ ภิกฺขเว ภิกฺขุนีปริสา ฯเปฯ อุปาสกปริสา ฯเปฯ อุปาสิกาปริสา อานนฺทํ ทสฺสนาย อุปสงฺกมติ, ทสฺสเนน สา}

\vspace{\tipitakaparskip}
\csromancenter{มหาปรินิพฺพานสุตฺต อตฺตมนา โหติ.\csromanspacer{} ตตฺร เจ อานนฺโท ธมฺมํ ภาสติ, ภาสิเตนปิ สา อตฺตมนา โหติ.\csromanspacer{} อติตฺตาว ภิกฺขเว อุปาสิกาปริสา โหติ.\csromanspacer{} อถ โข อานนฺโท ตุณฺหี โหติ.\csromanspacer{} อิเม โข ภิกฺขเว จตฺตาโร อจฺฉริยา อพฺภุตา ธมฺมา อานนฺเทติ.}
\csromanmatikaanchor{mtk.58}
\csromantocmarkref{section}{มหาสุทสฺสนสุตฺตเทสนา}{mtk.58}
\bookmark[dest={mtk.58},level=1]{มหาสุทสฺสนสุตฺตเทสนา}
\csromanheader{1}{\textbf{มหาสุทสฺสนสุตฺตเทสนา}}
\proseitem{210}{\csromanfolio{121}เอวํ วุตฺเต อายสฺมา อานนฺโท ภควนฺตํ เอตทโวจ “มา ภนฺเต ภควา อิมสฺมิํ ขุทฺทกนครเก อุชฺชงฺคลนครเก สาขานครเก ปรินิพฺพายิ.\csromanspacer{} สนฺติ ภนฺเต อญฺญานิ มหานครานิ.\csromanspacer{} เสยฺยถิทํ, จมฺปา ราชคหํ สาวตฺถี สาเกตํ โกสมฺพี พาราณสี, เอตฺถ ภควา ปรินิพฺพายตุ, เอตฺถ พหู ขตฺติยมหาสาลา พฺราหฺมณมหาสาลา คหปติมหาสาลา ตถาคเต อภิปฺปสนฺนา.\csromanspacer{} เต ตถาคตสฺส สรีรปูชํ กริสฺสนฺตี”ติ.\csromanspacer{} มาเหวํ อานนฺท อวจ, มาเหวํ อานนฺท อวจ “ขุทฺทกนครกํ อุชฺชงฺคลนครกํ สาขานครกนฺ”ติ.}

\prose{ภูตปุพฺพํ อานนฺท ราชา มหาสุทสฺสโน นาม อโหสิ จกฺกวตฺตี ธมฺมิโก ธมฺมราชา จาตุรนฺโต วิชิตาวี ชนปฺปทตฺถาวริยปฺปตฺโต สตฺตรตนสมนฺนาคโต.\csromanspacer{} รญฺโญ อานนฺท มหาสุทสฺสนสฺส อยํ กุสินารา กุสาวตี นาม ราชธานี อโหสิ, ปุรตฺถิเมน จ ปจฺฉิเมน จ ทฺวาทสโยชนานิ อายาเมน.\csromanspacer{} อุตฺตเรน จ ทกฺขิเณน จ สตฺตโยชนานิ วิตฺถาเรน.\csromanspacer{} กุสาวตี อานนฺท ราชธานี อิทฺธา เจว อโหสิ ผีตา จ พหุชนา จ อากิณฺณมนุสฺสา จ สุภิกฺขา จ.\csromanspacer{} เสยฺยถาปิ อานนฺท เทวานํ อาฬกมนฺทา นาม ราชธานี อิทฺธา เจว โหติ ผีตา จ พหุชนา จ อากิณฺณยกฺขา จ สุภิกฺขา จ.\csromanspacer{} เอวเมว โข อานนฺท กุสาวตี ราชธานี อิทฺธา เจว อโหสิ ผีตา จ พหุชนา จ อากิณฺณมนุสฺสา จ สุภิกฺขา จ.\csromanspacer{} กุสาวตี อานนฺท ราชธานี ทสหิ สทฺเทหิ อวิวิตฺตา อโหสิ ทิวา เจว รตฺติํ จ.\csromanspacer{} เสยฺยถิทํ, หตฺถิสทฺเทน อสฺสสทฺเทน รถสทฺเทน เภริสทฺเทน มุทิงฺคสทฺเทน วีณาสทฺเทน คีตสทฺเทน สงฺขสทฺเทน สมฺมสทฺเทน ปาณิตาฬสทฺเทน “อสฺนาถ ปิวถ ขาทถา”ติ ทสเมน สทฺเทน.}

\prose{คจฺฉ ตฺวํ อานนฺท กุสินารํ ปวิสิตฺวา โกสินารกานํ มลฺลานํ อาโรเจหิ “อชฺช โข วาเสฏฺฐา รตฺติยา ปจฺฉิเม ยาเม ตถาคตสฺส \csromanfolio{122}ปรินิพฺพานํ ภวิสฺสติ, อภิกฺกมถ วาเสฏฺฐา, อภิกฺกมถ วาเสฏฺฐา, มา ปจฺฉา วิปฺปฏิสาริโน อหุวตฺถ ‘อมฺหากํ จ โน คามกฺเขตฺเต ตถาคตสฺส ปรินิพฺพานํ อโหสิ, น มยํ ลภิมฺหา ปจฺฉิเม กาเล ตถาคตํ ทสฺสนายา’ติ”. “เอวํ ภนฺเต”ติ โข อายสฺมา อานนฺโท ภควโต ปฏิสฺสุตฺวา นิวาเสตฺวา ปตฺตจีวรมาทาย อตฺตทุติโย กุสินารํ ปาวิสิ.}

\csromanmatikaanchor{mtk.59}
\csromantocmarkref{section}{มลฺลานํ วนฺทนา}{mtk.59}
\bookmark[dest={mtk.59},level=1]{มลฺลานํ วนฺทนา}
\csromanheader{1}{\textbf{มลฺลานํ วนฺทนา}}
\proseitem{211}{เตน โข ปน สมเยน โกสินารกา มลฺลา สนฺธาคาเร\footnote{สนฺถาคาเร (สี, สฺยา, อิ)} สนฺนิปติตา โหนฺติ เกนจิเทว กรณีเยน.\csromanspacer{} อถ โข อายสฺมา อานนฺโท เยน โกสินารกานํ มลฺลานํ สนฺธาคารํ เตนุปสงฺกมิ, อุปสงฺกมิตฺวา โกสินารกานํ มลฺลานํ อาโรเจสิ “อชฺช โข วาเสฏฺฐา รตฺติยา ปจฺฉิเม ยาเม ตถาคตสฺส ปรินิพฺพานํ ภวิสฺสติ, อภิกฺกมถ วาเสฏฺฐา, อภิกฺกมถ วาเสฏฺฐา, มา ปจฺฉา วิปฺปฏิสาริโน อหุวตฺถ ‘อมฺหากํ จ โน คามกฺเขตฺเต ตถาคตสฺส ปรินิพฺพานํ อโหสิ, น มยํ ลภิมฺหา ปจฺฉิเม กาเล ตถาคตํ ทสฺสนายา’ติ”.\csromanspacer{} อิทมายสฺมโต อานนฺทสฺส วจนํ สุตฺวา มลฺลา จ มลฺลปุตฺตา จ มลฺลสุณิสา จ มลฺลปชาปติโย จ อฆาวิโน ทุมฺมนา เจโตทุกฺขสมปฺปิตา อปฺเปกจฺเจ เกเส ปกิริย กนฺทนฺติ, พาหา ปคฺคยฺห กนฺทนฺติ, ฉินฺนปาตํ ปปตนฺติ, อาวฏฺฏนฺติ, วิวฏฺฏนฺติ, “อติขิปฺปํ ภควา ปรินิพฺพายิสฺสติ, อติขิปฺปํ สุคโต ปรินิพฺพายิสฺสติ, อติขิปฺปํ จกฺขุํ โลเก อนฺตรธายิสฺสตี”ติ.\csromanspacer{} อถ โข มลฺลา จ มลฺลปุตฺตา จ มลฺลสุณิสา จ มลฺลปชาปติโย จ อฆาวิโน ทุมฺมนา เจโตทุกฺขสมปฺปิตา เยน อุปวตฺตนํ มลฺลานํ สาลวนํ, เยนายสฺมา อานนฺโท เตนุปสงฺกมิํสุ.\csromanspacer{} อถ โข อายสฺมโต อานนฺทสฺส เอตทโหสิ “สเจ โข อหํ โกสินารเก มลฺเล เอกเมกํ ภควนฺตํ วนฺทาเปสฺสามิ, อวนฺทิโต ภควา โกสินารเกหิ มลฺเลหิ ภวิสฺสติ, อถายํ รตฺติ วิภายิสฺสติ.\csromanspacer{} ยํนูนาหํ โกสินารเก มลฺเล กุลปริวตฺตโส กุลปริวตฺตโส ฐเปตฺวา ภควนฺตํ วนฺทาเปยฺยํ ‘อิตฺถนฺนาโม ภนฺเต มลฺโล สปุตฺโต สภริโย สปริโส สามจฺโจ ภควโต}

\vspace{\tipitakaparskip}
\csromancenter{มหาปรินิพฺพานสุตฺต ปาเท สิรสา วนฺทตี’ติ”.\csromanspacer{} อถ โข อายสฺมา อานนฺโท โกสินารเก มลฺเล กุลปริวตฺตโส กุลปริวตฺตโส ฐเปตฺวา ภควนฺตํ วนฺทาเปสิ “อิตฺถนฺนาโม ภนฺเต มลฺโล สปุตฺโต สภริโย สปริโส สามจฺโจ ภควโต ปาเท สิรสา วนฺทตี”ติ.\csromanspacer{} อถ โข อายสฺมา อานนฺโท เอเตน อุปาเยน ปฐเมเนว ยาเมน โกสินารเก มลฺเล ภควนฺตํ วนฺทาเปสิ.}
\csromanmatikaanchor{mtk.60}
\csromantocmarkref{section}{สุภทฺทปริพฺพาชกวตฺถุ}{mtk.60}
\bookmark[dest={mtk.60},level=1]{สุภทฺทปริพฺพาชกวตฺถุ}
\csromanheader{1}{\textbf{สุภทฺทปริพฺพาชกวตฺถุ}}
\proseitem{212}{\csromanfolio{123}เตน โข ปน สมเยน สุภทฺโท นาม ปริพฺพาชโก กุสินารายํ ปฏิวสติ.\csromanspacer{} อสฺโสสิ โข สุภทฺโท ปริพฺพาชโก “อชฺช กิร รตฺติยา ปจฺฉิเม ยาเม สมณสฺส โคตมสฺส ปรินิพฺพานํ ภวิสฺสตี”ติ.\csromanspacer{} อถ โข สุภทฺทสฺส ปริพฺพาชกสฺส เอตทโหสิ “สุตํ โข ปน เมตํ ปริพฺพาชกานํ วุฑฺฒานํ มหลฺลกานํ อาจริยปาจริยานํ ภาสมานานํ ‘กทาจิ กรหจิ ตถาคตา โลเก อุปฺปชฺชนฺติ อรหนฺโต สมฺมาสมฺพุทฺธา’ติ.\csromanspacer{} อชฺเชว รตฺติยา ปจฺฉิเม ยาเม สมณสฺส โคตมสฺส ปรินิพฺพานํ ภวิสฺสติ, อตฺถิ จ เม อยํ กงฺขาธมฺโม อุปฺปนฺโน, เอวํ ปสนฺโน อหํ สมเณ โคตเม ‘ปโหติ เม สมโณ โคตโม ตถา ธมฺมํ เทเสตุํ, ยถาหํ อิมํ กงฺขาธมฺมํ ปชเหยฺยนฺ’ติ”.\csromanspacer{} อถ โข สุภทฺโท ปริพฺพาชโก เยน อุปวตฺตนํ มลฺลานํ สาลวนํ, เยนายสฺมา อานนฺโท เตนุปสงฺกมิ, อุปสงฺกมิตฺวา อายสฺมนฺตํ อานนฺทํ เอตทโวจ “สุตํ เมตํ โภ อานนฺท ปริพฺพาชกานํ วุฑฺฒานํ มหลฺลกานํ อาจริยปาจริยานํ ภาสมานานํ ‘กทาจิ กรหจิ ตถาคตา โลเก อุปฺปชฺชนฺติ อรหนฺโต สมฺมาสมฺพุทฺธา’ติ.\csromanspacer{} อชฺเชว รตฺติยา ปจฺฉิเม ยาเม สมณสฺส โคตมสฺส ปรินิพฺพานํ ภวิสฺสติ, อตฺถิ จ เม อยํ กงฺขาธมฺโม อุปฺปนฺโน.\csromanspacer{} เอวํ ปสนฺโน อหํ สมเณ โคตเม ‘ปโหติ เม สมโณ โคตโม ตถา ธมฺมํ เทเสตุํ, ยถาหํ อิมํ กงฺขาธมฺมํ ปชเหยฺยนฺ’ติ.\csromanspacer{} สาธาหํ โภ อานนฺท ลเภยฺยํ สมณํ โคตมํ ทสฺสนายา”ติ.\csromanspacer{} เอวํ วุตฺเต อายสฺมา อานนฺโท สุภทฺทํ ปริพฺพาชกํ เอตทโวจ “อลํ อาวุโส สุภทฺท, มา ตถาคตํ วิเหเฐสิ, กิลนฺโต ภควา”ติ.\csromanspacer{} ทุติยมฺปิ โข สุภทฺโท ปริพฺพาชโก ฯเปฯ.\csromanspacer{} ตติยมฺปิ โข สุภทฺโท ปริพฺพาชโก อายสฺมนฺตํ อานนฺทํ เอตทโวจ “สุตํ เมตํ โภ อานนฺท ปริพฺพาชกานํ วุฑฺฒานํ มหลฺลกานํ อาจริยปาจริยานํ ภาสมานานํ ‘กทาจิ กรหจิ \csromanfolio{124}ตถาคตา โลเก อุปฺปชฺชนฺติ อรหนฺโต สมฺมาสมฺพุทฺธา’ติ.\csromanspacer{} อชฺเชว รตฺติยา ปจฺฉิเม ยาเม สมณสฺส โคตมสฺส ปรินิพฺพานํ ภวิสฺสติ, อตฺถิ จ เม อยํ กงฺขาธมฺโม อุปฺปนฺโน.\csromanspacer{} เอวํ ปสนฺโน อหํ สมเณ โคตเม ‘ปโหติ เม สมโณ โคตโม ตถา ธมฺมํ เทเสตุํ, ยถาหํ อิมํ กงฺขาธมฺมํ ปชเหยฺยนฺ’ติ.\csromanspacer{} สาธาหํ โภ อานนฺท ลเภยฺยํ สมณํ โคตมํ ทสฺสนายา” ติ.\csromanspacer{} ตติยมฺปิ โข อายสฺมา อานนฺโท สุภทฺทํ ปริพฺพาชกํ เอตทโวจ “อลํ อาวุโส สุภทฺท, มา ตถาคตํ วิเหเฐสิ, กิลนฺโต ภควา”ติ.}

\proseitem{213}{อสฺโสสิ โข ภควา อายสฺมโต อานนฺทสฺส สุภทฺเทน ปริพฺพาชเกน สทฺธิํ อิมํ กถาสลฺลาปํ.\csromanspacer{} อถ โข ภควา อายสฺมนฺตํ อานนฺทํ อามนฺเตสิ “อลํ อานนฺท มา สุภทฺทํ วาเรสิ, ลภตํ อานนฺท สุภทฺโท ตถาคตํ ทสฺสนาย.\csromanspacer{} ยํ กิญฺจิ มํ สุภทฺโท ปุจฺฉิสฺสติ, สพฺพํ ตํ อญฺญาเปกฺโขว ปุจฺฉิสฺสติ, โน วิเหสาเปกฺโข.\csromanspacer{} ยญฺจสฺสาหํ ปุฏฺโฐ พฺยากริสฺสามิ, ตํ ขิปฺปเมว อาชานิสฺสตี”ติ.\csromanspacer{} อถ โข อายสฺมา อานนฺโท สุภทฺทํ ปริพฺพาชกํ เอตทโวจ “คจฺฉาวุโส สุภทฺท, กโรติ เต ภควา โอกาสนฺ”ติ.\csromanspacer{} อถ โข สุภทฺโท ปริพฺพาชโก เยน ภควา เตนุปสงฺกมิ, อุปสงฺกมิตฺวา ภควตา สทฺธิํ สมฺโมทิ, สมฺโมทนียํ กถํ สารณียํ วีติสาเรตฺวา เอกมนฺตํ นิสีทิ, เอกมนฺตํ นิสินฺโน โข สุภทฺโท ปริพฺพาชโก ภควนฺตํ เอตทโวจ “เยเม โภ โคตม สมณพฺราหฺมณา สงฺฆิโน คณิโน คณาจริยา ญาตา ยสสฺสิโน ติตฺถกรา สาธุสมฺมตา พหุชนสฺส.\csromanspacer{} เสยฺยถิทํ, ปูรโณ กสฺสโป มกฺขลิ โคสาโล อชิโต เกสกมฺพโล ปกุโธ กจฺจายโน สญฺจโย เพลฏฺฐปุตฺโต นิคณฺโฐ นาฏปุตฺโต, สพฺเพเต สกาย ปฏิญฺญาย อพฺภญฺญิํสุ, สพฺเพว น อพฺภญฺญิํสุ, อุทาหุ เอกจฺเจ อพฺภญฺญิํสุ, เอกจฺเจ น อพฺภญฺญิํสู”ติ.\csromanspacer{} อลํ สุภทฺท ติฏฺฐเตตํ “สพฺเพเต สกาย ปฏิญฺญาย อพฺภญฺญิํสุ, สพฺเพว น อพฺภญฺญิํสุ, อุทาหุ เอกจฺเจ อพฺภญฺญิํสุ, เอกจฺเจ น อพฺภญฺญิํสู”ติ.\csromanspacer{} ธมฺมํ เต สุภทฺท เทเสสฺสามิ, ตํ สุณาหิ สาธุกํ มนสิกโรหิ, ภาสิสฺสามีติ. “เอวํ ภนฺเต”ติ โข สุภทฺโท ปริพฺพาชโก ภควโต ปจฺจสฺโสสิ.\csromanspacer{} ภควา เอตทโวจ\csromandash{}}

\proseitem{214}{“ยสฺมิํ โข สุภทฺท ธมฺมวินเย อริโย อฏฺฐงฺคิโก มคฺโค น อุปลพฺภติ, สมโณปิ ตตฺถ น อุปลพฺภติ, ทุติโยปิ ตตฺถ}

\vspace{\tipitakaparskip}
\csromancenter{มหาปรินิพฺพานสุตฺต สมโณ น อุปลพฺภติ, ตติโยปิ ตตฺถ สมโณ น อุปลพฺภติ, จตุตฺโถปิ ตตฺถ สมโณ น อุปลพฺภติ.\csromanspacer{} ยสฺมิํ จ โข สุภทฺท ธมฺมวินเย อริโย อฏฺฐงฺคิโก มคฺโค อุปลพฺภติ, สมโณปิ ตตฺถ อุปลพฺภติ, ทุติโยปิ ตตฺถ สมโณ อุปลพฺภติ, ตติโยปิ ตตฺถ สมโณ อุปลพฺภติ, จตุตฺโถปิ ตตฺถ สมโณ อุปลพฺภติ.\csromanspacer{} อิมสฺมิํ โข สุภทฺท ธมฺมวินเย อริโย อฏฺฐงฺคิโก มคฺโค อุปลพฺภติ, อิเธว สุภทฺท สมโณ, อิธ ทุติโย สมโณ, อิธ ตติโย สมโณ, อิธ จตุตฺโถ สมโณ, สุญฺญา ปรปฺปวาทา สมเณภิ อญฺเญหิ\footnote{อญฺเญ (อิ)}.\csromanspacer{} อิเม จ\footnote{อิเธว (ก)} สุภทฺท ภิกฺขู สมฺมา วิหเรยฺยุํ, อสุญฺโญ โลโก อรหนฺเตหิ อสฺสาติ.}
\vspace{0.5\baselineskip}
\csromangathameasure{เอกูนติํโส วยสา สุภทฺท, \\ ยํ ปพฺพชิํ กิํกุสลานุ-เอสี. \\ วสฺสานิ ปญฺญาส สมาธิกานิ, \\ ยโต อหํ ปพฺพชิโต สุภทฺท. \\ ญายสฺส ธมฺมสฺส ปเทสวตฺตี, \\ อิโต พหิทฺธา สมโณปิ นตฺถิ.}
\csromangathagroup{\csromangathastanza{\csromanfolio{125}เอกูนติํโส วยสา สุภทฺท, \\ ยํ ปพฺพชิํ กิํกุสลานุ-เอสี. \\ วสฺสานิ ปญฺญาส สมาธิกานิ, \\ ยโต อหํ ปพฺพชิโต สุภทฺท. \\ ญายสฺส ธมฺมสฺส ปเทสวตฺตี, \\ อิโต พหิทฺธา สมโณปิ นตฺถิ.}}

\prose{ทุติโยปิ สมโณ นตฺถิ.\csromanspacer{} ตติโยปิ สมโณ นตฺถิ.\csromanspacer{} จตุตฺโถปิ สมโณ นตฺถิ สุญฺญา ปรปฺปวาทา สมเณภิ อญฺเญหิ.\csromanspacer{} อิเม จ สุภทฺท ภิกฺขู สมฺมา วิหเรยฺยุํ, อสุญฺโญ โลโก อรหนฺเตหิ อสฺสา”ติ.}

\proseitem{215}{เอวํ วุตฺเต สุภทฺโท ปริพฺพาชโก ภควนฺตํ เอตทโวจ “อภิกฺกนฺตํ ภนฺเต, อภิกฺกนฺตํ ภนฺเต, เสยฺยถาปิ ภนฺเต นิกฺกุชฺชิตํ วา อุกฺกุชฺเชยฺย, ปฏิจฺฉนฺนํ วา วิวเรยฺย, มูฬฺหสฺส วา มคฺคํ อาจิกฺเขยฺย, อนฺธกาเร วา เตลปชฺโชตํ ธาเรยฺย ‘จกฺขุมนฺโต รูปานิ ทกฺขนฺตี’ติ.\csromanspacer{} เอวเมวํ ภควตา อเนกปริยาเยน ธมฺโม ปกาสิโต, เอสาหํ ภนฺเต ภควนฺตํ สรณํ คจฺฉามิ ธมฺมญฺจ ภิกฺขุสํฆญฺจ, ลเภยฺยาหํ ภนฺเต ภควโต สนฺติเก ปพฺพชฺชํ, ลเภยฺยํ อุปสมฺปทนฺ”ติ.\csromanspacer{} โย โข สุภทฺท อญฺญติตฺถิยปุพฺโพ อิมสฺมิํ ธมฺมวินเย อากงฺขติ ปพฺพชฺชํ, อากงฺขติ อุปสมฺปทํ, โส \csromanfolio{126}จตฺตาโร มาเส ปริวสติ, จตุนฺนํ มาสานํ อจฺจเยน อารทฺธจิตฺตา ภิกฺขู ปพฺพาเชนฺติ อุปสมฺปาเทนฺติ ภิกฺขุภาวาย.\csromanspacer{} อปิ จ เมตฺถ ปุคฺคลเวมตฺตตา วิทิตาติ.\csromanspacer{} สเจ ภนฺเต อญฺญติตฺถิยปุพฺพา อิมสฺมิํ ธมฺมวินเย อากงฺขนฺตา ปพฺพชฺชํ อากงฺขนฺตา อุปสมฺปทํ จตฺตาโร มาเส ปริวสนฺติ, จตุนฺนํ มาสานํ อจฺจเยน อารทฺธจิตฺตา ภิกฺขู ปพฺพาเชนฺติ อุปสมฺปาเทนฺติ ภิกฺขุภาวาย.\csromanspacer{} อหํ จตฺตาริ วสฺสานิ ปริวสิสฺสามิ, จตุนฺนํ วสฺสานํ อจฺจเยน อารทฺธจิตฺตา ภิกฺขู ปพฺพาเชนฺตุ อุปสมฺปาเทนฺตุ ภิกฺขุภาวายาติ.}

\prose{อถ โข ภควา อายสฺมนฺตํ อานนฺทํ อามนฺเตสิ “เตนหานนฺท สุภทฺทํ ปพฺพาเชหี”ติ. “เอวํ ภนฺเต”ติ โข อายสฺมา อานนฺโท ภควโต ปจฺจสฺโสสิ.\csromanspacer{} อถ โข สุภทฺโท ปริพฺพาชโก อายสฺมนฺตํ อานนฺทํ เอตทโวจ “ลาภา โว อาวุโส อานนฺท, สุลทฺธํ โว อาวุโส อานนฺท, เย เอตฺถ สตฺถุ\footnote{สตฺถารา (สฺยา)} สมฺมุขา อนฺเตวาสิกาภิเสเกน อภิสิตฺตา”ติ.\csromanspacer{} อลตฺถ โข สุภทฺโท ปริพฺพาชโก ภควโต สนฺติเก ปพฺพชฺชํ, อลตฺถ อุปสมฺปทํ.\csromanspacer{} อจิรูปสมฺปนฺโน โข ปนายสฺมา สุภทฺโท เอโก วูปกฏฺโฐ อปฺปมตฺโต อาตาปี ปหิตตฺโต วิหรนฺโต นจิรสฺเสว ‘ยสฺสตฺถาย กุลปุตฺตา สมฺมเทว อคารสฺมา อนคาริยํ ปพฺพชนฺติ’, ตทนุตฺตรํ พฺรหฺมจริยปริโยสานํ ทิฏฺเฐว ธมฺเม สยํ อภิญฺญา สจฺฉิกตฺวา อุปสมฺปชฺช วิหาสิ. “ขีณา ชาติ, วุสิตํ พฺรหฺมจริยํ, กตํ กรณียํ, นาปรํ อิตฺถตฺตายา”ติ อพฺภญฺญาสิ.\csromanspacer{} อญฺญตโร โข ปนายสฺมา สุภทฺโท อรหตํ อโหสิ.\csromanspacer{} โส ภควโต ปจฺฉิโม สกฺขิสาวโก อโหสีติ.}

\csromanheader{2}{ปญฺจโม ภาณวาโร.}
\csromansectionrule
\csromanmatikaanchor{mtk.61}
\csromantocmarkref{section}{ตถาคตปจฺฉิมวาจา}{mtk.61}
\bookmark[dest={mtk.61},level=1]{ตถาคตปจฺฉิมวาจา}
\csromanheader{1}{\textbf{ตถาคตปจฺฉิมวาจา}}
\proseitem{216}{อถ โข ภควา อายสฺมนฺตํ อานนฺทํ อามนฺเตสิ “สิยา โข ปนานนฺท ตุมฺหากํ เอวมสฺส ‘อตีตสตฺถุกํ ปาวจนํ, นตฺถิ โน สตฺถา’ติ.\csromanspacer{} น โข ปเนตํ อานนฺท เอวํ ทฏฺฐพฺพํ.\csromanspacer{} โย โว อานนฺท มยา ธมฺโม จ}

\vspace{\tipitakaparskip}
\csromancenter{มหาปรินิพฺพานสุตฺต วินโย จ เทสิโต ปญฺญตฺโต, โส โว มมจฺจเยน สตฺถา.\csromanspacer{} ยถา โข ปนานนฺท เอตรหิ ภิกฺขู อญฺญมญฺญํ อาวุโสวาเทน สมุทาจรนฺติ, น โข มมจฺจเยน เอวํ สมุทาจริตพฺพํ.\csromanspacer{} เถรตเรน อานนฺท ภิกฺขุนา นวกตโร ภิกฺขุ นาเมน วา โคตฺเตน วา อาวุโสวาเทน วา สมุทาจริตพฺโพ, นวกตเรน ภิกฺขุนา เถรตโร ภิกฺขุ ‘ภนฺเต’ติ วา ‘อายสฺมา’ติ วา สมุทาจริตพฺโพ.\csromanspacer{} อากงฺขมาโน อานนฺท สํโฆ มมจฺจเยน ขุทฺทานุขุทฺทกานิ สิกฺขาปทานิ สมูหนตุ.\csromanspacer{} ฉนฺนสฺส อานนฺท ภิกฺขุโน มมจฺจเยน พฺรหฺมทณฺโฑ ทาตพฺโพ”ติ.\csromanspacer{} กตโม ปน ภนฺเต พฺรหฺมทณฺโฑติ.\csromanspacer{} ฉนฺโน อานนฺท ภิกฺขุ ยํ อิจฺเฉยฺย, ตํ วเทยฺย.\csromanspacer{} โส ภิกฺขูหิ เนว วตฺตพฺโพ, น โอวทิตพฺโพ, น อนุสาสิตพฺโพติ.}
\proseitem{217}{\csromanfolio{127}อถ โข ภควา ภิกฺขู อามนฺเตสิ “สิยา โข ปน ภิกฺขเว เอกภิกฺขุสฺสาปิ กงฺขา วา วิมติ วา พุทฺเธ วา ธมฺเม วา สํเฆ วา มคฺเค วา ปฏิปทาย วา, ปุจฺฉถ ภิกฺขเว, มา ปจฺฉา วิปฺปฏิสาริโน อหุวตฺถ ‘สมฺมุขีภูโต โน สตฺถา อโหสิ.\csromanspacer{} น มยํ สกฺขิมฺหา ภควนฺตํ สมฺมุขา ปฏิปุจฺฉิตุนฺ’ติ”.\csromanspacer{} เอวํ วุตฺเต เต ภิกฺขู ตุณฺหี อเหสุํ.\csromanspacer{} ทุติยมฺปิ โข ภควา ฯเปฯ.\csromanspacer{} ตติยมฺปิ โข ภควา ภิกฺขู อามนฺเตสิ “สิยา โข ปน ภิกฺขเว เอกภิกฺขุสฺสาปิ กงฺขา วา วิมติ วา พุทฺเธ วา ธมฺเม วา สํเฆ วา มคฺเค วา ปฏิปทาย วา, ปุจฺฉถ ภิกฺขเว, มา ปจฺฉา วิปฺปฏิสาริโน อหุวตฺถ ‘สมฺมุขีภูโต โน สตฺถา อโหสิ.\csromanspacer{} น มยํ สกฺขิมฺหา ภควนฺตํ สมฺมุขา ปฏิปุจฺฉิตุนฺ’ติ”.\csromanspacer{} ตติยมฺปิ โข เต ภิกฺขู ตุณฺหี อเหสุํ.\csromanspacer{} อถ โข ภควา ภิกฺขู อามนฺเตสิ “สิยา โข ปน ภิกฺขเว สตฺถุคารเวนปิ น ปุจฺเฉยฺยาถ.\csromanspacer{} สหายโกปิ ภิกฺขเว สหายกสฺส อาโรเจตู”ติ.\csromanspacer{} เอวํ วุตฺเต เต ภิกฺขู ตุณฺหี อเหสุํ.\csromanspacer{} อถ โข อายสฺมา อานนฺโท ภควนฺตํ เอตทโวจ “อจฺฉริยํ ภนฺเต, อพฺภุตํ ภนฺเต, เอวํ ปสนฺโน อหํ ภนฺเต อิมสฺมิํ ภิกฺขุสํเฆ ‘นตฺถิ เอกภิกฺขุสฺสาปิ กงฺขา วา วิมติ วา พุทฺเธ วา ธมฺเม วา สํเฆ วา มคฺเค วา ปฏิปทาย วา’ติ”.\csromanspacer{} ปสาทา โข ตฺวํ อานนฺท วเทสิ, ญาณเมว เหตฺถ อานนฺท ตถาคตสฺส.\csromanspacer{} นตฺถิ อิมสฺมิํ ภิกฺขุสํเฆ เอกภิกฺขุสฺสาปิ กงฺขา วา วิมติ วา พุทฺเธ วา ธมฺเม วา สํเฆ วา มคฺเค วา ปฏิปทาย วา.\csromanspacer{} อิเมสํ หิ อานนฺท ปญฺจนฺนํ ภิกฺขุสตานํ โย ปจฺฉิมโก ภิกฺขุ, โส โสตาปนฺโน อวินิปาตธมฺโม นิยโต สมฺโพธิปรายโณติ.}

\proseitem{218}{\csromanfolio{128}อถ โข ภควา ภิกฺขู อามนฺเตสิ “หนฺท ทานิ ภิกฺขเว อามนฺตยามิ โว, วยธมฺมา สงฺขารา อปฺปมาเทน สมฺปาเทถา”ติ.\csromanspacer{} อยํ ตถาคตสฺส ปจฺฉิมา วาจา.}

\csromanmatikaanchor{mtk.62}
\csromantocmarkref{section}{ปรินิพฺพุตกถา}{mtk.62}
\bookmark[dest={mtk.62},level=1]{ปรินิพฺพุตกถา}
\csromanheader{1}{\textbf{ปรินิพฺพุตกถา}}
\proseitem{219}{อถ โข ภควา ปฐมํ ฌานํ สมาปชฺชิ, ปฐมชฺฌานา วุฏฺฐหิตฺวา ทุติยํ ฌานํ สมาปชฺชิ, ทุติยชฺฌานา วุฏฺฐหิตฺวา ตติยํ ฌานํ สมาปชฺชิ, ตติยชฺฌานา วุฏฺฐหิตฺวา จตุตฺถํ ฌานํ สมาปชฺชิ, จตุตฺถชฺฌานา วุฏฺฐหิตฺวา อากาสานญฺจายตนํ สมาปชฺชิ, อากาสานญฺจายตนสมาปตฺติยา วุฏฺฐหิตฺวา วิญฺญาณญฺจายตนํ สมาปชฺชิ, วิญฺญาณญฺจายตนสมาปตฺติยา วุฏฺฐหิตฺวา อากิญฺจญฺญายตนํ สมาปชฺชิ, อากิญฺจญฺญายตนสมาปตฺติยา วุฏฺฐหิตฺวา เนวสญฺญานาสญฺญายตนํ สมาปชฺชิ, เนวสญฺญานาสญฺญายตนสมาปตฺติยา วุฏฺฐหิตฺวา สญฺญาเวทยิตนิโรธํ สมาปชฺชิ.}

\prose{อถ โข อายสฺมา อานนฺโท อายสฺมนฺตํ อนุรุทฺธํ เอตทโวจ “นิพฺพานํ ภนฺเต อนุรุทฺธ ภควา”ติ.\csromanspacer{} นาวุโส อานนฺท ภควา นิพฺพานํ, สญฺญาเวทยิตนิโรธํ สมาปนฺโนติ.}

\prose{อถโข ภควา สญฺญาเวทยิตนิโรธสมาปตฺติยา วุฏฺฐหิตฺวา เนวสญฺญานาสญฺญายตนํ สมาปชฺชิ, เนวสญฺญานาสญฺญายตนสมาปตฺติยา วุฏฺฐหิตฺวา อากิญฺจญฺญายตนํ สมาปชฺชิ, อากิญฺจญฺญายตนสมาปตฺติยา วุฏฺฐหิตฺวา วิญฺญาณญฺจายตนํ สมาปชฺชิ, วิญฺญาณญฺจายตนสมาปตฺติยา วุฏฺฐหิตฺวา อากาสานญฺจายตนํ สมาปชฺชิ, อากาสานญฺจายตนสมาปตฺติยา วุฏฺฐหิตฺวา จตุตฺถํ ฌานํ สมาปชฺชิ, จตุตฺถชฺฌานา วุฏฺฐหิตฺวา ตติยํ ฌานํ สมาปชฺชิ, ตติยชฺฌานา วุฏฺฐหิตฺวา ทุติยํ ฌานํ สมาปชฺชิ, ทุติยชฺฌานา วุฏฺฐหิตฺวา ปฐมํ ฌานํ สมาปชฺชิ, ปฐมชฺฌานา วุฏฺฐหิตฺวา ทุติยํ ฌานํ สมาปชฺชิ, ทุติยชฺฌานา วุฏฺฐหิตฺวา ตติยํ ฌานํ สมาปชฺชิ, ตติยชฺฌานา วุฏฺฐหิตฺวา จตุตฺถํ ฌานํ สมาปชฺชิ, จตุตฺถชฺฌานา วุฏฺฐหิตฺวา สมนนฺตรา ภควา ปรินิพฺพายิ.}

\proseitem{220}{ปรินิพฺพุเต ภควติ สห ปรินิพฺพานา มหาภูมิจาโล อโหสิ ภิํสนโก สโลมหํโส.\csromanspacer{} เทวทุนฺทุภิโย จ ผลิํสุ.\csromanspacer{} ปรินิพฺพุเต ภควติ สห ปรินิพฺพานา พฺรหฺมาสหมฺปติ อิมํ คาถํ อภาสิ.}

\csromanheader{2}{3. มหาปรินิพฺพานสุตฺต}
\csromangathameasure{“สพฺเพว นิกฺขิปิสฺสนฺติ, ภูตา โลเก สมุสฺสยํ. \\ ยตฺถ เอตาทิโส สตฺถา, โลเก อปฺปฏิปุคฺคโล. \\ ตถาคโต พลปฺปตฺโต, สมฺพุทฺโธ นิพฺพานํ”ติ.}
\csromangathagroup{\csromangathastanza{\csromanfolio{129}“สพฺเพว นิกฺขิปิสฺสนฺติ, ภูตา โลเก สมุสฺสยํ. \\ ยตฺถ เอตาทิโส สตฺถา, โลเก อปฺปฏิปุคฺคโล. \\ ตถาคโต พลปฺปตฺโต, สมฺพุทฺโธ นิพฺพานํ”ติ.}}

\proseitem{221}{ปรินิพฺพุเต ภควติ สห ปรินิพฺพานา สกฺโก เทวานมินฺโท อิมํ คาถํ อภาสิ.}

\csromangathameasure{“อนิจฺจา วต สงฺขารา, อุปฺปาทวยธมฺมิโน. \\ อุปฺปชฺชิตฺวา นิรุชฺฌนฺติ, เตสํ วูปสโม สุโข”ติ.}
\csromangathagroup{\csromangathastanza{“อนิจฺจา วต สงฺขารา, อุปฺปาทวยธมฺมิโน. \\ อุปฺปชฺชิตฺวา นิรุชฺฌนฺติ, เตสํ วูปสโม สุโข”ติ.}}

\proseitem{222}{ปรินิพฺพุเต ภควติ สห ปรินิพฺพานา อายสฺมา อนุรุทฺโธ อิมา คาถาโย อภาสิ.}

\csromangathameasure{“นาหุ อสฺสาสปสฺสาโส, ฐิตจิตฺตสฺส ตาทิโน. \\ อเนโช สนฺติมารพฺภ, ยํ กาลมกรี มุนิ. \\ อสลฺลีเนน จิตฺเตน, เวทนํ อชฺฌวาสยิ. \\ ปชฺโชตสฺเสว นิพฺพานํ, วิโมกฺโข เจตโส อหู”ติ.}
\csromangathagroup{\csromangathastanza{“นาหุ อสฺสาสปสฺสาโส, ฐิตจิตฺตสฺส ตาทิโน. \\ อเนโช สนฺติมารพฺภ, ยํ กาลมกรี มุนิ.}\csromangathastanza{อสลฺลีเนน จิตฺเตน, เวทนํ อชฺฌวาสยิ. \\ ปชฺโชตสฺเสว นิพฺพานํ, วิโมกฺโข เจตโส อหู”ติ.}}

\proseitem{223}{ปรินิพฺพุเต ภควติ สห ปรินิพฺพานา อายสฺมา อานนฺโท อิมํ คาถํ อภาสิ.}

\csromangathameasure{“ตทาสิ ยํ ภิํสนกํ, ตทาสิ โลมหํสนํ. \\ สพฺพาการวรูเปเต, สมฺพุทฺเธ ปรินิพฺพุเต”ติ.}
\csromangathagroup{\csromangathastanza{“ตทาสิ ยํ ภิํสนกํ, ตทาสิ โลมหํสนํ. \\ สพฺพาการวรูเปเต, สมฺพุทฺเธ ปรินิพฺพุเต”ติ.}}

\proseitem{224}{ปรินิพฺพุเต ภควติ เย เต ตตฺถ ภิกฺขู อวีตราคา อปฺเปกจฺเจ พาหา ปคฺคยฺห กนฺทนฺติ, ฉินฺนปาตํ ปปตนฺติ, อาวฏฺฏนฺติ, วิวฏฺฏนฺติ, “อติขิปฺปํ ภควา ปรินิพฺพุโต, อติขิปฺปํ สุคโต ปรินิพฺพุโต, อติขิปฺปํ จกฺขุํ โลเก อนฺตรหิโต”ติ.\csromanspacer{} เย ปน เต ภิกฺขู วีตราคา, เต สตา สมฺปชานา อธิวาเสนฺติ “อนิจฺจา สงฺขารา, ตํ กุเตตฺถ ลพฺภา”ติ.}

\proseitem{225}{อถ โข อายสฺมา อนุรุทฺโธ ภิกฺขู อามนฺเตสิ “อลํ อาวุโส มา โสจิตฺถ มา ปริเทวิตฺถ.\csromanspacer{} นนุ เอตํ อาวุโส ภควตา ปฏิกจฺเจว อกฺขาตํ ‘สพฺเพเหว ปิเยหิ มนาเปหิ นานาภาโว วินาภาโว อญฺญถาภาโว’.\csromanspacer{} ตํ กุเตตฺถ อาวุโส ลพฺภา, ‘ยํ ตํ ชาตํ ภูตํ สงฺขตํ ปโลกธมฺมํ, ตํ วต มา ปลุชฺชี’ติ, เนตํ ฐานํ \csromanfolio{130}วิชฺชติ.\csromanspacer{} เทวตา อาวุโส อุชฺฌายนฺตี”ติ.\csromanspacer{} กถํภูตา ปน ภนฺเต อายสฺมา อนุรุทฺโธ เทวตา มนสิ กโรตีติ\footnote{ภนฺเต อนุรุทฺธ เทวตา มนสิ กโรนฺตีติ (สฺยา, ก)}.}

\prose{สนฺตาวุโส อานนฺท เทวตา อากาเส ปถวีสญฺญินิโย เกเส ปกิริย กนฺทนฺติ, พาหา ปคฺคยฺห กนฺทนฺติ, ฉินฺนปาตํ ปปตนฺติ, อาวฏฺฏนฺติ, วิวฏฺฏนฺติ, “อติขิปฺปํ ภควา ปรินิพฺพุโต, อติขิปฺปํ สุคโต ปรินิพฺพุโต, อติขิปฺปํ จกฺขุํ โลเก อนฺตรหิโต”ติ.\csromanspacer{} สนฺตาวุโส อานนฺท เทวตา ปถวิยา ปถวีสญฺญินิโย เกเส ปกิริย กนฺทนฺติ, พาหา ปคฺคยฺห กนฺทนฺติ, ฉินฺนปาตํ ปปตนฺติ, อาวฏฺฏนฺติ, วิวฏฺฏนฺติ, “อติขิปฺปํ ภควา ปรินิพฺพุโต, อติขิปฺปํ สุคโต ปรินิพฺพุโต, อติขิปฺปํ จกฺขุํ โลเก อนฺตรหิโต”ติ.\csromanspacer{} ยา ปน ตา เทวตา วีตราคา.\csromanspacer{} ตา สตา สมฺปชานา อธิวาเสนฺติ “อนิจฺจา สงฺขารา, ตํ กุเตตฺถ ลพฺภา”ติ.\csromanspacer{} อถ โข อายสฺมา จ อนุรุทฺโธ อายสฺมา จ อานนฺโท ตํ รตฺตาวเสสํ ธมฺมิยา กถาย วีตินาเมสุํ.}

\proseitem{226}{อถ โข อายสฺมา อนุรุทฺโธ อายสฺมนฺตํ อานนฺทํ อามนฺเตสิ “คจฺฉาวุโส อานนฺท กุสินารํ ปวิสิตฺวา โกสินารกานํ มลฺลานํ อาโรเจหิ ‘นิพฺพานํ วาเสฏฺฐา ภควา, ยสฺสทานิ กาลํ มญฺญถา’ติ”. “เอวํ ภนฺเต”ติ โข อายสฺมา อานนฺโท อายสฺมโต อนุรุทฺธสฺส ปฏิสฺสุตฺวา ปุพฺพณฺหสมยํ นิวาเสตฺวา ปตฺตจีวรมาทาย อตฺตทุติโย กุสินารํ ปาวิสิ.\csromanspacer{} เตน โข ปน สมเยน โกสินารกา มลฺลา สนฺธาคาเร สนฺนิปติตา โหนฺติ เตเนว กรณีเยน.\csromanspacer{} อถ โข อายสฺมา อานนฺโท เยน โกสินารกานํ มลฺลานํ สนฺธาคารํ เตนุปสงฺกมิ, อุปสงฺกมิตฺวา โกสินารกานํ มลฺลานํ อาโรเจสิ “นิพฺพานํ วาเสฏฺฐา ภควา, ยสฺสทานิ กาลํ มญฺญถา”ติ.\csromanspacer{} อิทมายสฺมโต อานนฺทสฺส วจนํ สุตฺวา มลฺลา จ มลฺลปุตฺตา จ มลฺลสุณิสา จ มลฺลปชาปติโย จ อฆาวิโน ทุมฺมนา เจโตทุกฺขสมปฺปิตา อปฺเปกจฺเจ เกเส ปกิริย กนฺทนฺติ, พาหา ปคฺคยฺห กนฺทนฺติ, ฉินฺนปาตํ ปปตนฺติ, อาวฏฺฏนฺติ, วิวฏฺฏนฺติ, “อติขิปฺปํ ภควาปรินิพฺพุโต, อติขิปฺปํ สุคโต ปรินิพฺพุโต, อติขิปฺปํ จกฺขุํ โลเก อนฺตรหิโต”ติ.}

\csromanmatikaanchor{mtk.63}
\csromantocmarkref{section}{พุทฺธสรีรปูชา}{mtk.63}
\bookmark[dest={mtk.63},level=1]{พุทฺธสรีรปูชา}
\csromanheader{1}{3. มหาปรินิพฺพานสุตฺต \textbf{พุทฺธสรีรปูชา}}
\proseitem{227}{\csromanfolio{131}อถ โข โกสินารกา มลฺลา ปุริเส อาณาเปสุํ “เตน หิ ภเณ กุสินารายํ คนฺธมาลญฺจ สพฺพญฺจ ตาฬาวจรํ สนฺนิปาเตถา”ติ.\csromanspacer{} อถ โข โกสินารกา มลฺลา คนฺธมาลญฺจ สพฺพญฺจ ตาฬาวจรํ ปญฺจ จ ทุสฺสยุคสตานิ อาทาย เยน อุปวตฺตนํ มลฺลานํ สาลวนํ, เยน ภควโต สรีรํ เตนุปสงฺกมิํสุ, อุปสงฺกมิตฺวา ภควโต สรีรํ นจฺเจหิ คีเตหิ วาทิเตหิ มาเลหิ คนฺเธหิ สกฺกโรนฺตา ครุํ กโรนฺตา มาเนนฺตา ปูเชนฺตา เจลวิตานานิ กโรนฺตา มณฺฑลมาเฬ ปฏิยาเทนฺตา เอกทิวสํ วีตินาเมสุํ.}

\prose{อถ โข โกสินารกานํ มลฺลานํ เอตทโหสิ “อติวิกาโล โข อชฺช ภควโต สรีรํ ฌาเปตุํ, เสฺว ทานิ มยํ ภควโต สรีรํ ฌาเปสฺสามา”ติ.\csromanspacer{} อถ โข โกสินารกา มลฺลา ภควโต สรีรํ นจฺเจหิ คีเตหิ วาทิเตหิ มาเลหิ คนฺเธหิ สกฺกโรนฺตา ครุํ กโรนฺตา มาเนนฺตา ปูเชนฺตา เจลวิตานานิ กโรนฺตา มณฺฑลมาเฬ ปฏิยาเทนฺตา ทุติยมฺปิ ทิวสํ วีตินาเมสุํ, ตติยมฺปิ ทิวสํ วีตินาเมสุํ, จตุตฺถมฺปิ ทิวสํ วีตินาเมสุํ, ปญฺจมมฺปิ ทิวสํ วีตินาเมสุํ, ฉฏฺฐมฺปิ ทิวสํ วีตินาเมสุํ.}

\prose{อถ โข สตฺตมํ ทิวสํ โกสินารกานํ มลฺลานํ เอตทโหสิ “มยํ ภควโต สรีรํ นจฺเจหิ คีเตหิ วาทิเตหิ มาเลหิ คนฺเธหิ สกฺกโรนฺตา ครุํ กโรนฺตา มาเนนฺตา ปูเชนฺตา ทกฺขิเณน ทกฺขิณํ นครสฺส หริตฺวา พาหิเรน พาหิรํ ทกฺขิณโต นครสฺส ภควโต สรีรํ ฌาเปสฺสามา”ติ.}

\proseitem{228}{เตน โข ปน สมเยน อฏฺฐ มลฺลปาโมกฺขา สีสํนฺหาตา อหตานิ วตฺถานิ นิวตฺถา “มยํ ภควโต สรีรํ อุจฺจาเรสฺสามา”ติ.\csromanspacer{} น สกฺโกนฺติ อุจฺจาเรตุํ.\csromanspacer{} อถ โข โกสินารกา มลฺลา อายสฺมนฺตํ อนุรุทฺธํ เอตทโวจุํ “โก นุ โข ภนฺเต อนุรุทฺธ เหตุ โก ปจฺจโย, เยนิเม อฏฺฐ มลฺลปาโมกฺขา สีสํนฺหาตา อหตานิ วตฺถานิ นิวตฺถา ‘มยํ ภควโต สรีรํ อุจฺจาเรสฺสามา’ติ น สกฺโกนฺติ อุจฺจาเรตุนฺ”ติ.\csromanspacer{} อญฺญถา โข วาเสฏฺฐา ตุมฺหากํ อธิปฺปาโย, อญฺญถา เทวตานํ อธิปฺปาโยติ.\csromanspacer{} กถํ ปน ภนฺเต เทวตานํ อธิปฺปาโยติ.\csromanspacer{} ตุมฺหากํ โข วาเสฏฺฐา อธิปฺปาโย “มยํ ภควโต สรีรํ นจฺเจหิ คีเตหิ วาทิเตหิ มาเลหิ \csromanfolio{132}คนฺเธหิ สกฺกโรนฺตา ครุํ กโรนฺตา มาเนนฺตา ปูเชนฺตา ทกฺขิเณน ทกฺขิณํ นครสฺส หริตฺวา พาหิเรน พาหิรํ ทกฺขิณโต นครสฺส ภควโต สรีรํ ฌาเปสฺสามา”ติ.\csromanspacer{} เทวตานํ โข วาเสฏฺฐา อธิปฺปาโย “มยํ ภควโต สรีรํ ทิพฺเพหิ นจฺเจหิ คีเตหิ วาทิเตหิ มาเลหิ คนฺเธหิ สกฺกโรนฺตา ครุํ กโรนฺตา มาเนนฺตา ปูเชนฺตา อุตฺตเรน อุตฺตรํ นครสฺส หริตฺวา อุตฺตเรน ทฺวาเรน นครํ ปเวเสตฺวา มชฺเฌน มชฺฌํ นครสฺส หริตฺวา ปุรตฺถิเมน ทฺวาเรน นิกฺขมิตฺวา ปุรตฺถิมโต นครสฺส มกุฏพนฺธนํ นาม มลฺลานํ เจติยํ, เอตฺถ ภควโต สรีรํ ฌาเปสฺสามา”ติ.\csromanspacer{} ยถา ภนฺเต เทวตานํ อธิปฺปาโย, ตถา โหตูติ.}

\proseitem{229}{เตน โข ปน สมเยน กุสินารา ยาว สนฺธิสมลสํกฏีรา ชณฺณุมตฺเตน โอธินา มนฺทารวปุปฺเผหิ สนฺถตา\footnote{สณฺฐิตา (สฺยา)} โหติ.\csromanspacer{} อถ โข เทวตา จ โกสินารกา จ มลฺลา ภควโต สรีรํ ทิพฺเพหิ จ มานุสเกหิ จ นจฺเจหิ คีเตหิ วาทิเตหิ มาเลหิ คนฺเธหิ สกฺกโรนฺตา ครุํ กโรนฺตา มาเนนฺตา ปูเชนฺตา อุตฺตเรน อุตฺตรํ นครสฺส หริตฺวา อุตฺตเรน ทฺวาเรน นครํ ปเวเสตฺวา มชฺเฌน มชฺฌํ นครสฺส หริตฺวา ปุรตฺถิเมน ทฺวาเรน นิกฺขมิตฺวา ปุรตฺถิมโต นครสฺส มกุฏพนฺธนํ นาม มลฺลานํ เจติยํ, เอตฺถ จ ภควโต สรีรํ นิกฺขิปิํสุ.}

\proseitem{230}{อถ โข โกสินารกา มลฺลา อายสฺมนฺตํ อานนฺทํ เอตทโวจุํ “กถํ มยํ ภนฺเต อานนฺท ตถาคตสฺส สรีเร ปฏิปชฺชามา”ติ.\csromanspacer{} ยถา โข วาเสฏฺฐา รญฺโญ จกฺกวตฺติสฺส สรีเร ปฏิปชฺชนฺติ, เอวํ ตถาคตสฺส สรีเร ปฏิปชฺชิตพฺพนฺติ.\csromanspacer{} กถํ ปน ภนฺเต อานนฺท รญฺโญ จกฺกวตฺติสฺส สรีเร ปฏิปชฺชนฺตีติ.\csromanspacer{} รญฺโญ วาเสฏฺฐา จกฺกวตฺติสฺส สรีรํ อหเตน วตฺเถน เวเฐนฺติ, อหเตน วตฺเถน เวเฐตฺวา วิหเตน กปฺปาเสน เวเฐนฺติ, วิหเตน กปฺปาเสน เวเฐตฺวา อหเตน วตฺเถน เวเฐนฺติ, เอเตน อุปาเยน ปญฺจหิ ยุคสเตหิ รญฺโญ จกฺกวตฺติสฺส สรีรํ เวเฐตฺวา อายสาย เตลโทณิยา ปกฺขิปิตฺวา อญฺญิสฺสา อายสาย โทณิยา ปฏิกุชฺชิตฺวา สพฺพคนฺธานํ จิตกํ กริตฺวา รญฺโญ จกฺกวตฺติสฺส สรีรํ ฌาเปนฺติ.\csromanspacer{} จาตุมหาปเถ รญฺโญ จกฺกวตฺติสฺส ถูปํ}

\vspace{\tipitakaparskip}
\csromancenter{มหาปรินิพฺพานสุตฺต กโรนฺติ.\csromanspacer{} เอวํ โข วาเสฏฺฐา รญฺโญ จกฺกวตฺติสฺส สรีเร ปฏิปชฺชนฺติ.\csromanspacer{} ยถา โข วาเสฏฺฐา รญฺโญ จกฺกวตฺติสฺส สรีเร ปฏิปชฺชนฺติ.\csromanspacer{} เอวํ ตถาคตสฺส สรีเร ปฏิปชฺชิตพฺพํ.\csromanspacer{} จาตุมหาปเถ ตถาคตสฺส ถูโป กาตพฺโพ, ตตฺถ เย มาลํ วา คนฺธํ วา จุณฺณกํ วา อาโรเปสฺสนฺติ วา อภิวาเทสฺสนฺติ วา จิตฺตํ วา ปสาเทสฺสนฺติ, เตสํ ตํ ภวิสฺสติ ทีฆรตฺตํ หิตาย สุขายาติ.\csromanspacer{} อถ โข โกสินารกา มลฺลา ปุริเส อาณาเปสุํ “เตน หิ ภเณ มลฺลานํ วิหตํ กปฺปาสํ สนฺนิปาเตถา”ติ.}
\prose{\csromanfolio{133}อถ โข โกสินารกา มลฺลา ภควโต สรีรํ อหเตน วตฺเถน เวเฐตฺวา วิหเตน กปฺปาเสน เวเฐสุํ, วิหเตน กปฺปาเสน เวเฐตฺวา อหเตน วตฺเถน เวเฐสุํ.\csromanspacer{} เอเตน อุปาเยน ปญฺจหิ ยุคสเตหิ ภควโต สรีรํ เวเฐตฺวา อายสาย เตลโทณิยา ปกฺขิปิตฺวา อญฺญิสฺสา อายสาย โทณิยา ปฏิกุชฺชิตฺวา สพฺพคนฺธานํ จิตกํ กริตฺวา ภควโต สรีรํ จิตกํ อาโรเปสุํ.}

\csromanmatikaanchor{mtk.64}
\csromantocmarkref{section}{มหากสฺสปตฺเถรวตฺถุ}{mtk.64}
\bookmark[dest={mtk.64},level=1]{มหากสฺสปตฺเถรวตฺถุ}
\csromanheader{1}{\textbf{มหากสฺสปตฺเถรวตฺถุ}}
\proseitem{231}{เตน โข ปน สมเยน อายสฺมา มหากสฺสโป ปาวาย กุสินารํ อทฺธานมคฺคปฺปฏิปนฺโน โหติ มหตา ภิกฺขุสํเฆน สทฺธิํ ปญฺจมตฺเตหิ ภิกฺขุสเตหิ.\csromanspacer{} อถ โข อายสฺมา มหากสฺสโป มคฺคา โอกฺกมฺม อญฺญตรสฺมิํ รุกฺขมูเล นิสีทิ.\csromanspacer{} เตน โข ปน สมเยน อญฺญตโร อาชีวโก กุสินาราย มนฺทารวปุปฺผํ คเหตฺวา ปาวํ อทฺธานมคฺคปฺปฏิปนฺโน โหติ.\csromanspacer{} อทฺทสา โข อายสฺมา มหากสฺสโป ตํ อาชีวกํ ทูรโตว อาคจฺฉนฺตํ, ทิสฺวา ตํ อาชีวกํ เอตทโวจ “อปาวุโส อมฺหากํ สตฺถารํ ชานาสี”ติ.\csromanspacer{} อามาวุโส ชานามิ, อชฺช สตฺตาหนิพฺพานํ สมโณ โคตโม, ตโต เม อิทํ มนฺทารวปุปฺผํ คหิตนฺติ.\csromanspacer{} ตตฺถ เย เต ภิกฺขู อวีตราคา อปฺเปกจฺเจ พาหา ปคฺคยฺห กนฺทนฺติ, ฉินฺนปาตํ ปปตนฺติ, อาวฏฺฏนฺติ, วิวฏฺฏนฺติ, “อติขิปฺปํ ภควา ปรินิพฺพุโต, อติขิปฺปํ สุคโตปรินิพฺพุโต, อติขิปฺปํ จกฺขุํ โลเก อนฺตรหิโต”ติ.\csromanspacer{} เย ปน เต ภิกฺขู วีตราคา, เต สตา สมฺปชานา อธิวาเสนฺติ “อนิจฺจา สงฺขารา, ตํ กุเตตฺถ ลพฺภา”ติ.}

\proseitem{232}{เตน โข ปน สมเยน สุภทฺโท นาม วุทฺธปพฺพชิโต ตสฺสํ ปริสายํ นิสินฺโน โหติ.\csromanspacer{} อถ โข สุภทฺโท วุทฺธปพฺพชิโต \csromanfolio{134}เต ภิกฺขู เอตทโวจ “อลํ อาวุโส มา โสจิตฺถ, มา ปริเทวิตฺถ, สุมุตฺตา มยํ เตน มหาสมเณน, อุปทฺทุตา จ โหม ‘อิทํ โว กปฺปติ, อิทํ โว น กปฺปตี’ติ.\csromanspacer{} อิทานิ ปน มยํ ยํ อิจฺฉิสฺสาม, ตํ กริสฺสาม, ยํ น อิจฺฉิสฺสาม, น ตํ กริสฺสามา”ติ.\csromanspacer{} อถ โข อายสฺมา มหากสฺสโป ภิกฺขู อามนฺเตสิ “อลํ อาวุโส มา โสจิตฺถ, มา ปริเทวิตฺถ.\csromanspacer{} นนุ เอตํ อาวุโส ภควตา ปฏิกจฺเจว อกฺขาตํ ‘สพฺเพเหว ปิเยหิ มนาเปหิ นานาภาโว วินาภาโว อญฺญถาภาโว’.\csromanspacer{} ตํ กุเตตฺถ อาวุโส ลพฺภา, ‘ยํ ตํ ชาตํ ภูตํ สงฺขตํ ปโลกธมฺมํ, ตํ ตถาคตสฺสาปิ สรีรํ มา ปลุชฺชี’ติ, เนตํ ฐานํ วิชฺชตี”ติ.}

\proseitem{233}{เตน โข ปน สมเยน จตฺตาโร มลฺลปาโมกฺขา สีสํนฺหาตา อหตานิ วตฺถานิ นิวตฺถา “มยํ ภควโต จิตกํ อาฬิมฺเปสฺสามา”ติ น สกฺโกนฺติ อาฬิมฺเปตุํ.\csromanspacer{} อถ โข โกสินารกา มลฺลา อายสฺมนฺตํ อนุรุทฺธํ เอตทโวจุํ “โก นุ โข ภนฺเต อนุรุทฺธ เหตุ โก ปจฺจโย, เยนิเม จตฺตาโร มลฺลปาโมกฺขา สีสํนฺหาตา อหตานิ วตฺถานิ นิวตฺถา ‘มยํ ภควโต จิตกํ อาฬิมฺเปสฺสามา’ติ น สกฺโกนฺติ อาฬิมฺเปตุนฺ”ติ.\csromanspacer{} อญฺญถา โข วาเสฏฺฐา เทวตานํ อธิปฺปาโยติ.\csromanspacer{} กถํ ปน ภนฺเต เทวตานํ อธิปฺปาโยติ.\csromanspacer{} เทวตานํ โข วาเสฏฺฐา อธิปฺปาโย “อยํ อายสฺมา มหากสฺสโป ปาวาย กุสินารํ อทฺธานมคฺคปฺปฏิปนฺโน มหตา ภิกฺขุสํเฆน สทฺธิํ ปญฺจมตฺเตหิ ภิกฺขุสเตหิ.\csromanspacer{} น ตาว ภควโต จิตโก ปชฺชลิสฺสติ, ยาวายสฺมา มหากสฺสโป ภควโต ปาเท สิรสา น วนฺทิสฺสตี”ติ.\csromanspacer{} ยถา ภนฺเต เทวตานํ อธิปฺปาโย, ตถา โหตูติ.}

\proseitem{234}{อถ โข อายสฺมา มหากสฺสโป เยน กุสินารา มกุฏพนฺธนํ นาม มลฺลานํ เจติยํ, เยน ภควโต จิตโก เตนุปสงฺกมิ, อุปสงฺกมิตฺวา เอกํสํ จีวรํ กตฺวา อญฺชลิํ ปณาเมตฺวา ติกฺขตฺตุํ จิตกํ ปทกฺขิณํ กตฺวา ภควโต ปาเท สิรสา วนฺทิ.\csromanspacer{} ตานิปิ โข ปญฺจภิกฺขุสตานิ เอกํสํ จีวรํ กตฺวา อญฺชลิํ ปณาเมตฺวา ติกฺขตฺตุํ จิตกํ ปทกฺขิณํ กตฺวา ภควโต ปาเท สิรสา วนฺทิํสุ.\csromanspacer{} วนฺทิเต จ ปนายสฺมตา มหากสฺสเปน เตหิ จ ปญฺจหิ ภิกฺขุสเตหิ สยเมว ภควโต จิตโก ปชฺชลิ.}

\csromanheader{2}{3. มหาปรินิพฺพานสุตฺต}
\proseitem{235}{\csromanfolio{135}ฌายมานสฺส โข ปน ภควโต สรีรสฺส ยํ อโหสิ ฉวีติ วา จมฺมนฺติ วา มํสนฺติ วา นฺหารูติ วา ลสิกาติ วา, ตสฺส เนว ฉาริกา ปญฺญายิตฺถ, น มสิ, สรีราเนว อวสิสฺสิํสุ.\csromanspacer{} เสยฺยถาปิ นาม สปฺปิสฺส วา เตลสฺส วา ฌายมานสฺส เนว ฉาริกา ปญฺญายติ, น มสิ.\csromanspacer{} เอวเมว ภควโต สรีรสฺส ฌายมานสฺส ยํ อโหสิ ฉวีติ วา จมฺมนฺติ วา มํสนฺติ วา นฺหารูติ วา ลสิกาติ วา, ตสฺส เนว ฉาริกา ปญฺญายิตฺถ, น มสิ, สรีราเนว อวสิสฺสิํสุ.\csromanspacer{} เตสญฺจ ปญฺจนฺนํ ทุสฺสยุคสตานํ เทฺวว ทุสฺสานิ น ฑยฺหิํสุ ยญฺจ สพฺพ-อพฺภนฺตริมํ ยญฺจ พาหิรํ.\csromanspacer{} ทฑฺเฒ จ โข ปน ภควโต สรีเร อนฺตลิกฺขา อุทกธารา ปาตุภวิตฺวา ภควโต จิตกํ นิพฺพาเปสิ.\csromanspacer{} อุทกสาลโตปิ\footnote{อุทกํ สาลโตปิ (สี, สฺยา, กํ)} อพฺภุนฺนมิตฺวา ภควโต จิตกํ นิพฺพาเปสิ.\csromanspacer{} โกสินารกาปิ มลฺลา สพฺพคนฺโธทเกน ภควโต จิตกํ นิพฺพาเปสุํ.\csromanspacer{} อถ โข โกสินารกา มลฺลา ภควโต สรีรานิ สตฺตาหํ สนฺธาคาเร สตฺติปญฺชรํ กริตฺวา ธนุปาการํ ปริกฺขิปาเปตฺวา\footnote{ปริกฺขิปิตฺวา (สฺยา)} นจฺเจหิ คีเตหิ วาทิเตหิ มาเลหิ คนฺเธหิ สกฺกริํสุ ครุํ กริํสุ มาเนสุํ ปูเชสุํ.}

\csromanmatikaanchor{mtk.65}
\csromantocmarkref{section}{สรีรธาตุวิภชน}{mtk.65}
\bookmark[dest={mtk.65},level=1]{สรีรธาตุวิภชน}
\csromanheader{1}{\textbf{สรีรธาตุวิภชน}}
\proseitem{236}{อสฺโสสิ โข ราชา มาคโธ อชาตสตฺตุ เวเทหิปุตฺโต “ภควา กิร กุสินารายํ นิพฺพานํ”ติ.\csromanspacer{} อถ โข ราชา มาคโธ อชาตสตฺตุ เวเทหิปุตฺโต โกสินารกานํ มลฺลานํ ทูตํ ปาเหสิ “ภควาปิ ขตฺติโย อหมฺปิ ขตฺติโย, อหมฺปิ อรหามิ ภควโต สรีรานํ ภาคํ, อหมฺปิ ภควโต สรีรานํ ถูปญฺจ มหญฺจ กริสฺสามี”ติ.}

\prose{อสฺโสสุํ โข เวสาลิกา ลิจฺฉวี “ภควา กิร กุสินารายํ นิพฺพานํ”ติ.\csromanspacer{} อถ โข เวสาลิกา ลิจฺฉวี โกสินารกานํ มลฺลานํ ทูตํ ปาเหสุํ “ภควาปิ ขตฺติโย มยมฺปิ ขตฺติยา, มยมฺปิ อรหาม ภควโต สรีรานํ ภาคํ, มยมฺปิ ภควโต สรีรานํ ถูปญฺจ มหญฺจ กริสฺสามา”ติ.}

\prose{อสฺโสสุํ โข กปิลวตฺถุวาสี สกฺยา “ภควา กิร กุสินารายํ นิพฺพานํ”ติ.\csromanspacer{} อถ โข กปิลวตฺถุวาสี สกฺยา โกสินารกานํ \csromanfolio{136}มลฺลานํ ทูตํ ปาเหสุํ “ภควา อมฺหากํ ญาติเสฏฺโฐ, มยมฺปิ อรหาม ภควโต สรีรานํ ภาคํ, มยมฺปิ ภควโต สรีรานํ ถูปญฺจ มหญฺจ กริสฺสามา”ติ.}

\prose{อสฺโสสุํ โข อลฺลกปฺปกา พุลโย\footnote{ถูลโย (สฺยา)} “ภควา กิร กุสินารายํ นิพฺพานํ”ติ.\csromanspacer{} อถ โข อลฺลกปฺปกา พุลโย โกสินารกานํ มลฺลานํ ทูตํ ปาเหสุํ “ภควาปิ ขตฺติโย มยมฺปิ ขตฺติยา, มยมฺปิ อรหาม ภควโต สรีรานํ ภาคํ, มยมฺปิ ภควโต สรีรานํ ถูปญฺจ มหญฺจ กริสฺสามา”ติ.}

\prose{อสฺโสสุํ โข รามคามกา โกฬิยา “ภควา กิร กุสินารายํ นิพฺพานํ”ติ.\csromanspacer{} อถ โข รามคามกา โกฬิยา โกสินารกานํ มลฺลานํ ทูตํ ปาเหสุํ “ภควาปิ ขตฺติโย มยมฺปิ ขตฺติยา, มยมฺปิ อรหาม ภควโต สรีรานํ ภาคํ, มยมฺปิ ภควโต สรีรานํ ถูปญฺจ มหญฺจ กริสฺสามา”ติ.}

\prose{อสฺโสสิ โข เวฏฺฐทีปโก พฺราหฺมโณ “ภควา กิร กุสินารายํ นิพฺพานํ”ติ.\csromanspacer{} อถ โข เวฏฺฐทีปโก พฺราหฺมโณ โกสินารกานํ มลฺลานํ ทูตํ ปาเหสิ “ภควาปิ ขตฺติโย อหมฺปิสฺมิ พฺราหฺมโณ, อหมฺปิ อรหามิ ภควโต สรีรานํ ภาคํ, อหมฺปิ ภควโต สรีรานํ ถูปญฺจ มหญฺจ กริสฺสามี”ติ.}

\prose{อสฺโสสุํ โข ปาเวยฺยกา มลฺลา “ภควา กิร กุสินารายํ นิพฺพานํ”ติ.\csromanspacer{} อถ โข ปาเวยฺยกา มลฺลา โกสินารกานํ มลฺลานํ ทูตํ ปาเหสุํ “ภควาปิ ขตฺติโย มยมฺปิ ขตฺติยา, มยมฺปิ อรหาม ภควโต สรีรานํ ภาคํ, มยมฺปิ ภควโต สรีรานํ ถูปญฺจ มหญฺจ กริสฺสามา”ติ.}

\prose{เอวํ วุตฺเต โกสินารกา มลฺลา เต สงฺเฆ คเณ เอตทโวจุํ “ภควา อมฺหากํ คามกฺเขตฺเต นิพฺพานํ, น มยํ ทสฺสาม ภควโต สรีรานํ ภาคนฺ”ติ.}

\csromanheader{2}{237. เอวํ วุตฺเต โทโณ พฺราหฺมโณ เต สงฺเฆ คเณ เอตทโวจ\csromandash{}}
\csromanheader{2}{3. มหาปรินิพฺพานสุตฺต}
\csromangathameasure{“สุณนฺตุ โภนฺโต มม เอกวาจํ, \\ อมฺหาก พุทฺโธ อหุ ขนฺติวาโท. \\ น หิ สาธุ ยํ อุตฺตมปุคฺคลสฺส, \\ สรีรภาเค สิยา สมฺปหาโร. \\ สพฺเพว โภนฺโต สหิตา สมคฺคา, \\ สมฺโมทมานา กโรมฏฺฐภาเค. \\ วิตฺถาริกา โหนฺตุ ทิสาสุ ถูปา, \\ พหู ชนา จกฺขุมโต ปสนฺนา”ติ.}
\csromangathagroup{\csromangathastanza{\csromanfolio{137}“สุณนฺตุ โภนฺโต มม เอกวาจํ, \\ อมฺหาก\footnote{ฉนฺทานุรกฺขณตฺถํ นิคฺคหีตโลโป.} พุทฺโธ อหุ ขนฺติวาโท. \\ น หิ สาธุ ยํ อุตฺตมปุคฺคลสฺส, \\ สรีรภาเค สิยา สมฺปหาโร.}\csromangathastanza{สพฺเพว โภนฺโต สหิตา สมคฺคา, \\ สมฺโมทมานา กโรมฏฺฐภาเค. \\ วิตฺถาริกา โหนฺตุ ทิสาสุ ถูปา, \\ พหู ชนา จกฺขุมโต ปสนฺนา”ติ.}}

\proseitem{238}{เตน หิ พฺราหฺมณ ตฺวญฺเญว ภควโต สรีรานิ อฏฺฐธา สมํ สุวิภตฺตํ วิภชาหีติ. “เอวํ โภ”ติ โข โทโณ พฺราหฺมโณ เตสํ สงฺฆานํ คณานํ ปฏิสฺสุตฺวา ภควโต สรีรานิ อฏฺฐธา สมํ สุวิภตฺตํ วิภชิตฺวา เต สงฺเฆ คเณ เอตทโวจ “อิมํ เม โภนฺโต ตุมฺพํ ททนฺตุ, อหมฺปิ ตุมฺพสฺส ถูปญฺจ มหญฺจ กริสฺสามี”ติ.\csromanspacer{} อทํสุ โข เต โทณสฺส พฺราหฺมณสฺส ตุมฺพํ.}

\prose{อสฺโสสุํ โข ปิปฺปลิวนิยา\footnote{ปิปฺผลิวนิยา (สฺยา)} โมริยา “ภควา กิร กุสินารายํ นิพฺพานํ”ติ.\csromanspacer{} อถ โข ปิปฺปลิวนิยา โมริยา โกสินารกานํ มลฺลานํ ทูตํ ปาเหสุํ “ภควาปิ ขตฺติโย มยมฺปิ ขตฺติยา, มยมฺปิ อรหาม ภควโต สรีรานํ ภาคํ, มยมฺปิ ภควโต สรีรานํ ถูปญฺจ มหญฺจ กริสฺสามา”ติ.\csromanspacer{} นตฺถิ ภควโต สรีรานํ ภาโค, วิภตฺตานิ ภควโต สรีรานิ.\csromanspacer{} อิโต องฺคารํ หรถาติ.\csromanspacer{} เต ตโต องฺคารํ หริํสุ\footnote{อาหริํสุ (สฺยา, ก)}.}

\csromanmatikaanchor{mtk.66}
\csromantocmarkref{section}{ธาตุถูปปูชา}{mtk.66}
\bookmark[dest={mtk.66},level=1]{ธาตุถูปปูชา}
\csromanheader{1}{\textbf{ธาตุถูปปูชา}}
\proseitem{239}{อถ โข ราชา มาคโธ อชาตสตฺตุ เวเทหิปุตฺโต ราชคเห ภควโต สรีรานํ ถูปญฺจ มหญฺจ อกาสิ.\csromanspacer{} เวสาลิกาปิ ลิจฺฉวี เวสาลิยํ ภควโต สรีรานํ ถูปญฺจ มหญฺจ อกํสุ.\csromanspacer{} กปิลวตฺถุวาสีปิ สกฺยา กปิลวตฺถุสฺมิํ ภควโต สรีรานํ ถูปญฺจ มหญฺจ อกํสุ.\csromanspacer{} อลฺลกปฺปกาปิ พุลโย อลฺลกปฺเป ภควโต สรีรานํ ถูปญฺจ มหญฺจ อกํสุ.\csromanspacer{} รามคามกาปิ โกฬิยา รามคาเม ภควโต สรีรานํ ถูปญฺจ มหญฺจ อกํสุ.\csromanspacer{} เวฏฺฐทีปโกปิ พฺราหฺมโณ เวฏฺฐทีเป \csromanfolio{138}ภควโต สรีรานํ ถูปญฺจ มหญฺจ อกาสิ.\csromanspacer{} ปาเวยฺยกาปิ มลฺลา ปาวายํ ภควโต สรีรานํ ถูปญฺจ มหญฺจ อกํสุ.\csromanspacer{} โกสินารกาปิ มลฺลา กุสินารายํ ภควโต สรีรานํ ถูปญฺจ มหญฺจ อกํสุ.\csromanspacer{} โทโณปิ พฺราหฺมโณ ตุมฺพสฺส ถูปญฺจ มหญฺจ อกาสิ.\csromanspacer{} ปิปฺปลิวนิยาปิ โมริยา ปิปฺปลิวเน องฺคารานํ ถูปญฺจ มหญฺจ อกํสุ.\csromanspacer{} อิติ อฏฺฐ สรีรถูปา นวโม ตุมฺพถูโป ทสโม องฺคารถูโป.\csromanspacer{} เอวเมตํ ภูตปุพฺพนฺติ.}

\csromanhangingbegin
\hangingheaditem{240}{อฏฺฐโทณํ จกฺขุมโต สรีรํ, สตฺตโทณํ ชมฺพุทีเป มเหนฺติ.}
\hangingbody{เอกญฺจ โทณํ ปุริสวรุตฺตมสฺส, รามคาเม นาคราชามเหติ.}
\hangingbody{เอกาหิ ทาฐา ติทิเวหิ ปูชิตา, เอกา ปน คนฺธารปุเร มหียติ.}
\hangingbody{กาลิงฺครญฺโญ วิชิเต ปุเนกํ, เอกํ ปน นาคราชา มเหติ.}
\hangingbody{ตสฺเสว เตเชน อยํ วสุนฺธรา,}
\hangingbody{อายาคเสฏฺเฐหิ มหี อลงฺกตา.}
\hangingbody{เอวํ อิมํ จกฺขุมโต สรีรํ,}
\hangingbody{สุสกฺกตํ สกฺกตสกฺกเตหิ.}
\hangingbody{เทวินฺท นาคินฺท นรินฺท ปูชิโต,}
\hangingbody{มนุสฺสินฺทเสฏฺเฐหิ ตเถว ปูชิโต.}
\hangingbody{ตํ วนฺทถ1 ปญฺชลิกา ลภิตฺวา,}
\hangingbody{พุทฺโธ หเว กปฺปสเตหิ ทุลฺลโภติ.}
\hangingbody{จตฺตาลีส สมา ทนฺตา, เกสา โลมา จ สพฺพโส.}
\hangingbody{เทวา หริํสุ เอเกกํ, จกฺกวาฬปรมฺปราติ.}
\csromanhangingend

\nitthitam{\textbf{มหาปรินิพฺพานสุตฺตํ นิฏฺฐิตํ ตติยํ.}}
\csromanmatikaanchor{mtk.67}
\csromantocmarknopage{chapter}{4. มหาสุทสฺสนสุตฺต}{mtk.67}
\bookmark[dest={mtk.67},level=0]{4. มหาสุทสฺสนสุตฺต}
\chapterheadpageread{4. มหาสุทสฺสนสุตฺต}
\proseitem{241}{\csromanfolio{139}เอวํ เม สุตํ\csromandash{}เอกํ สมยํ ภควา กุสินารายํ วิหรติ อุปวตฺตเน มลฺลานํ สาลวเน อนฺตเรน ยมกสาลานํ ปรินิพฺพานสมเย.\csromanspacer{} อถ โข อายสฺมา อานนฺโท เยน ภควา เตนุปสงฺกมิ, อุปสงฺกมิตฺวา ภควนฺตํ อภิวาเทตฺวา เอกมนฺตํ นิสีทิ.\csromanspacer{} เอกมนฺตํ นิสินฺโน โข อายสฺมา อานนฺโท ภควนฺตํ เอตทโวจ “มา ภนฺเต ภควา อิมสฺมิํ ขุทฺทกนครเก อุชฺชงฺคลนครเก สาขานครเก ปรินิพฺพายิ.\csromanspacer{} สนฺติ ภนฺเต อญฺญานิ มหานครานิ.\csromanspacer{} เสยฺยถิทํ, จมฺปา ราชคหํ สาวตฺถิ สาเกตํ โกสมฺพี พาราณสี, เอตฺถ ภควา ปรินิพฺพายตุ.\csromanspacer{} เอตฺถ พหู ขตฺติยมหาสาลา พฺราหฺมณมหาสาลา คหปติมหาสาลา ตถาคเต อภิปฺปสนฺนา, เต ตถาคตสฺส สรีรปูชํ กริสฺสนฺตี”ติ.}

\proseitem{242}{มา เหวํ อานนฺท อวจ, มา เหวํ อานนฺท อวจ “ขุทฺทกนครกํ อุชฺชงฺคลนครกํ สาขานครกนฺ”ติ.}

\csromanmatikaanchor{mtk.68}
\csromantocmarkref{section}{กุสาวตี ราชธานี}{mtk.68}
\bookmark[dest={mtk.68},level=1]{กุสาวตี ราชธานี}
\csromanheader{1}{\textbf{กุสาวตีราชธานี}}
\prose{ภูตปุพฺพํ อานนฺท ราชา มหาสุทสฺสโน นาม อโหสิ ขตฺติโย มุทฺธาวสิตฺโต\footnote{ขตฺติโย มุทฺธาภิสิตฺโต (ก), จกฺกวตฺตี ธมฺมิโก ธมฺมราชา (มหาปรินิพฺพานสุตฺต)} จาตุรนฺโต วิชิตาวี ชนปทตฺถาวริยปฺปตฺโต.\csromanspacer{} รญฺโญ อานนฺท มหาสุทสฺสนสฺส อยํ กุสินารา กุสาวตี นาม ราชธานี อโหสิ.\csromanspacer{} ปุรตฺถิเมน จ ปจฺฉิเมน จ ทฺวาทสโยชนานิ อายาเมน, อุตฺตเรน จ ทกฺขิเณน จ สตฺตโยชนานิ วิตฺถาเรน, กุสาวตี อานนฺท ราชธานี อิทฺธา เจว อโหสิ ผีตา จ พหุชนา จ อากิณฺณมนุสฺสา จ สุภิกฺขา จ.\csromanspacer{} เสยฺยถาปิ อานนฺท เทวานํ อาฬกมนฺทา นาม ราชธานี อิทฺธา เจว โหติ ผีตา จ\footnote{อิทฺธา เจว อโหสิ ผีตา จ (สฺยา)} พหุชนา จ อากิณฺณยกฺขา จ สุภิกฺขา จ.\csromanspacer{} เอวเมว โข อานนฺท \textbf{กุสาวตี ราชธานี} อิทฺธา เจว อโหสิ ผีตา จ พหุชนา จ อากิณฺณมนุสฺสา จ สุภิกฺขา จ.\csromanspacer{} กุสาวตี อานนฺท ราชธานี ทสหิ สทฺเทหิ อวิวิตฺตา อโหสิ ทิวา เจว รตฺติญฺจ.\csromanspacer{} เสยฺยถิทํ, หตฺถิสทฺเทน อสฺสสทฺเทน รถสทฺเทน เภริสทฺเทน มุทิงฺคสทฺเทน วีณาสทฺเทน คีตสทฺเทน}

\prose{\csromanfolio{140}สงฺขสทฺเทน สมฺมสทฺเทน ปาณิตาฬสทฺเทน “อสฺนาถ ปิวถ ขาทถา”ติ ทสเมน สทฺเทน.}

\prose{กุสาวตี อานนฺท ราชธานี สตฺตหิ ปากาเรหิ ปริกฺขิตฺตา อโหสิ.\csromanspacer{} เอโก ปากาโร โสวณฺณมโย, เอโก รูปิยมโย, เอโก เวฬุริยมโย, เอโก ผลิกมโย, เอโก โลหิตงฺกมโย\footnote{โลหิตงฺคมโย (ก), โลหิตกมโย (พฺยากรเณสุ)}, เอโก มสารคลฺลมโย, เอโก สพฺพรตนมโย.\csromanspacer{} กุสาวติยา อานนฺท ราชธานิยา จตุนฺนํ วณฺณานํ ทฺวารานิ อเหสุํ.\csromanspacer{} เอกํ ทฺวารํ โสวณฺณมยํ, เอกํ รูปิยมยํ, เอกํ เวฬุริยมยํ, เอกํ ผลิกมยํ, เอเกกสฺมิํ ทฺวาเร สตฺต สตฺต เอสิกา นิขาตา อเหสุํ ติโปริสงฺคา ติโปริสนิขาตา ทฺวาทสโปริสา อุพฺเพเธน.\csromanspacer{} เอกา เอสิกา โสวณฺณมยา, เอกา รูปิยมยา, เอกา เวฬุริยมยา, เอกา ผลิกมยา, เอกา โลหิตงฺกมยา, เอกา มสารคลฺลมยา, เอกา สพฺพรตนมยา.}

\prose{กุสาวตี อานนฺท ราชธานี สตฺตหิ ตาลปนฺตีหิ ปริกฺขิตฺตา อโหสิ.\csromanspacer{} เอกา ตาลปนฺติ โสวณฺณมยา, เอกา รูปิยมยา, เอกา เวฬุริยมยา, เอกา ผลิกมยา, เอกา โลหิตงฺกมยา, เอกา มสารคลฺลมยา, เอกา สพฺพรตนมยา.\csromanspacer{} โสวณฺณมยสฺส ตาลสฺส โสวณฺณมโย ขนฺโธ อโหสิ, รูปิยมยานิ ปตฺตานิ จ ผลานิ จ.\csromanspacer{} รูปิยมยสฺส ตาลสฺส รูปิยมโย ขนฺโธ อโหสิ, โสวณฺณมยานิ ปตฺตานิ จ ผลานิ จ.\csromanspacer{} เวฬุริยมยสฺส ตาลสฺส เวฬุริยมโย ขนฺโธ อโหสิ, ผลิกมยานิ ปตฺตานิ จ ผลานิ จ.\csromanspacer{} ผลิกมยสฺส ตาลสฺส ผลฺลิกมโย ขนฺโธ อโหสิ, เวฬุริยมยานิ ปตฺตานิ จ ผลานิ จ.\csromanspacer{} โลหิตงฺกมยสฺส ตาลสฺส โลหิตงฺกมโย ขนฺโธ อโหสิ, มสารคลฺลมยานิ ปตฺตานิ จ ผลานิ จ.\csromanspacer{} มสารคลฺลมยสฺส ตาลสฺส มสารคลฺลมโย ขนฺโธ อโหสิ, โลหิตงฺกมยานิ ปตฺตานิ จ ผลานิ จ.\csromanspacer{} สพฺพรตนมยสฺส ตาลสฺส สพฺพรตนมโย ขนฺโธ อโหสิ, สพฺพรตนมยานิ ปตฺตานิ จ ผลานิ จ.\csromanspacer{} ตาสํ โข ปนานนฺท ตาลปนฺตีนํ วาเตริตานํ สทฺโท อโหสิ วคฺคุ จ รชนีโย จ ขมนีโย\footnote{กมนีโย (สี, สฺยา, อิ)} จ มทนีโย จ.\csromanspacer{} เสยฺยถาปิ อานนฺท ปญฺจงฺคิกสฺส ตูริยสฺส สุวินีตสฺส สุปฺปฏิตาฬิตสฺส สุกุสเลหิ สมนฺนาหตสฺส สทฺโท โหติ วคฺคุ}

\vspace{\tipitakaparskip}
\csromancenter{มหาสุทสฺสนสุตฺต จ รชนีโย จ ขมนีโย จ มทนีโย จ.\csromanspacer{} เอวเมว โข อานนฺท ตาสํ ตาลปนฺตีนํ วาเตริตานํ สทฺโท อโหสิ วคฺคุ จ รชนีโย จ ขมนีโย จ มทนีโย จ.\csromanspacer{} เย โข ปนานนฺท เตน สมเยน กุสาวติยา ราชธานิยา ธุตฺตา อเหสุํ โสณฺฑา ปิปาสา, เต ตาสํ ตาลปนฺตีนํ วาเตริตานํ สทฺเทน ปริจาเรสุํ.}
\csromanmatikaanchor{mtk.69}
\csromantocmarkref{section}{จกฺกรตน}{mtk.69}
\bookmark[dest={mtk.69},level=1]{จกฺกรตน}
\csromanheader{1}{\textbf{จกฺกรตน}}
\proseitem{243}{\csromanfolio{141}ราชา อานนฺท มหาสุทสฺสโน สตฺตหิ รตเนหิ สมนฺนาคโต อโหสิ จตูหิ จ อิทฺธีหิ.\csromanspacer{} กตเมหิ สตฺตหิ.\csromanspacer{} อิธานนฺท รญฺโญ มหาสุทสฺสนสฺส ตทหุโปสเถ ปนฺนรเส สีสํนฺหาตสฺส อุโปสถิกสฺส อุปริปาสาทวรคตสฺส ทิพฺพํ จกฺกรตนํ ปาตุรโหสิ สหสฺสารํ สเนมิกํ สนาภิกํ สพฺพาการปริปูรํ, ทิสฺวา รญฺโญ มหาสุทสฺสนสฺส เอตทโหสิ “สุตํ โข ปน เมตํ ‘ยสฺส รญฺโญ ขตฺติยสฺส มุทฺธาวสิตฺตสฺส ตทหุโปสเถ ปนฺนรเส สีสํนฺหาตสฺส อุโปสถิกสฺส อุปริปาสาทวรคตสฺส ทิพฺพํ จกฺกรตนํ ปาตุภวติ สหสฺสารํ สเนมิกํ สนาภิกํ สพฺพาการปริปูรํ, โส โหติ ราชา จกฺกวตฺตี’ติ.\csromanspacer{} อสฺสํ นุ โข อหํ ราชา จกฺกวตฺตี”ติ.}

\proseitem{244}{อถ โข อานนฺท ราชา มหาสุทสฺสโน อุฏฺฐายาสนา เอกํสํ อุตฺตราสงฺคํ กริตฺวา วาเมน หตฺเถน สุวณฺณภิงฺการํ คเหตฺวา ทกฺขิเณน หตฺเถน จกฺกรตนํ อพฺภุกฺกิริ “ปวตฺตตุ ภวํ จกฺกรตนํ, อภิวิชินาตุ ภวํ จกฺกรตนนฺ”ติ, อถ โข ตํ อานนฺท จกฺกรตนํ ปุรตฺถิมํ ทิสํ ปวตฺติ\footnote{ปวตฺตติ (สฺยา, ก)}, อนฺวเทว\footnote{อนุเทว (สฺยา)} ราชา มหาสุทสฺสโน สทฺธิํ จตุรงฺคินิยา เสนาย, ยสฺมิํ โข ปนานนฺท ปเทเส จกฺกรตนํ ปติฏฺฐาสิ, ตตฺถ ราชา มหาสุทสฺสโน วาสํ อุปคจฺฉิ สทฺธิํ จตุรงฺคินิยา เสนาย.\csromanspacer{} เย โข ปนานนฺท ปุรตฺถิมาย ทิสาย ปฏิราชาโน, เต ราชานํ มหาสุทสฺสนํ อุปสงฺกมิตฺวา เอวมาหํสุ “เอหิ โข มหาราช, สฺวาคตํ เต มหาราช, สกํ เต มหาราช, อนุสาส มหาราชา”ติ.\csromanspacer{} ราชา มหาสุทสฺสโน เอวมาห “ปาโณ น หนฺตพฺโพ, อทินฺนํ น อาทาตพฺพํ, กาเมสุ มิจฺฉา น จริตพฺพา, มุสา น ภณิตพฺพา, มชฺชํ น ปาตพฺพํ, ยถาภุตฺตญฺจ \csromanfolio{142}ภุญฺชถา”ติ.\csromanspacer{} เย โข ปนานนฺท ปุรตฺถิมาย ทิสาย ปฏิราชาโน, เต รญฺโญ มหาสุทสฺสนสฺส อนุยนฺตา อเหสุํ.\csromanspacer{} อถ โข ตํ อานนฺท จกฺกรตนํ ปุรตฺถิมํ สมุทฺทํ อชฺโฌคาเหตฺวา ปจฺจุตฺตริตฺวา ทกฺขิณํ ทิสํ ปวตฺติ ฯเปฯ ทกฺขิณํ สมุทฺทํ อชฺโฌคาเหตฺวา ปจฺจุตฺตริตฺวา ปจฺฉิมํ ทิสํ ปวตฺติ ฯเปฯ ปจฺฉิมํ สมุทฺทํ อชฺโฌคาเหตฺวา ปจฺจุตฺตริตฺวา อุตฺตรํ ทิสํ ปวตฺติ, อนฺวเทว ราชา มหาสุทสฺสโน สทฺธิํ จตุรงฺคินิยา เสนาย, ยสฺมิํ โข ปนานนฺท ปเทเส จกฺกรตนํ ปติฏฺฐาสิ, ตตฺถ ราชา มหาสุทสฺสโน วาสํ อุปคจฺฉิ สทฺธิํ จตุรงฺคินิยา เสนาย.\csromanspacer{} เย โข ปนานนฺท อุตฺตราย ทิสาย ปฏิราชาโน, เต ราชานํ มหาสุทสฺสนํ อุปสงฺกมิตฺวา เอวมาหํสุ “เอหิ โข มหาราช, สฺวาคตํ เต มหาราช, สกํ เต มหาราช, อนุสาส มหาราชา”ติ.\csromanspacer{} ราชา มหาสุทสฺสโน เอวมาห “ปาโณ น หนฺตพฺโพ, อทินฺนํ น อาทาตพฺพํ, กาเมสุ มิจฺฉา น จริตพฺพา, มุสา น ภณิตพฺพา, มชฺชํ น ปาตพฺพํ, ยถาภุตฺตญฺจ ภุญฺชถา”ติ.\csromanspacer{} เย โข ปนานนฺท อุตฺตราย ทิสาย ปฏิราชาโน, เต รญฺโญ มหาสุทสฺสนสฺส อนุยนฺตา อเหสุํ.}

\proseitem{245}{อถ โข ตํ อานนฺท จกฺกรตนํ สมุทฺทปริยนฺตํ ปถวิํ อภิวิชินิตฺวา กุสาวติํ ราชธานิํ ปจฺจาคนฺตฺวา รญฺโญ มหาสุทสฺสนสฺส อนฺเตปุรทฺวาเร อตฺถกรณปมุเข อกฺขาหตํ มญฺเญ อฏฺฐาสิ รญฺโญ มหาสุทสฺสนสฺส อนฺเตปุรํ อุปโสภยมานํ.\csromanspacer{} รญฺโญ อานนฺท มหาสุทสฺสนสฺส เอวรูปํ จกฺกรตนํ ปาตุรโหสิ. (1)}

\csromanmatikaanchor{mtk.70}
\csromantocmarkref{section}{หตฺถิรตน}{mtk.70}
\bookmark[dest={mtk.70},level=1]{หตฺถิรตน}
\csromanheader{1}{\textbf{หตฺถิรตน}}
\proseitem{246}{ปุน จปรํ อานนฺท รญฺโญ มหาสุทสฺสนสฺส หตฺถิรตนํ ปาตุรโหสิ สพฺพเสโต สตฺตปฺปติฏฺโฐ อิทฺธิมา เวหาสงฺคโม อุโปสโถ นาม นาคราชา.\csromanspacer{} ตํ ทิสฺวา รญฺโญ มหาสุทสฺสนสฺส จิตฺตํ ปสีทิ “ภทฺทกํ วต โภ หตฺถิยานํ สเจ ทมถํ อุเปยฺยา”ติ.\csromanspacer{} อถ โข ตํ อานนฺท หตฺถิรตนํ เสยฺยถาปิ นาม คนฺธหตฺถาชานิโย ทีฆรตฺตํ สุปริทนฺโต, เอวเมว ทมถํ อุปคจฺฉิ.\csromanspacer{} ภูตปุพฺพํ อานนฺท ราชา มหาสุทสฺสโน ตเมว หตฺถิรตนํ วีมํสมาโน ปุพฺพณฺหสมยํ อภิรุหิตฺวา สมุทฺทปริยนฺตํ ปถวิํ อนุยายิตฺวา กุสาวติํ ราชธานิํ ปจฺจาคนฺตฺวา ปาตราสมกาสิ.\csromanspacer{} รญฺโญ อานนฺท มหาสุทสฺสนสฺส เอวรูปํ หตฺถิรตนํ ปาตุรโหสิ. (2)}

\csromanmatikaanchor{mtk.71}
\csromantocmarkref{section}{อสฺสรตน}{mtk.71}
\bookmark[dest={mtk.71},level=1]{อสฺสรตน}
\csromanheader{1}{4. มหาสุทสฺสนสุตฺต \textbf{อสฺสรตน}}
\proseitem{247}{\csromanfolio{143}ปุน จปรํ อานนฺท รญฺโญ มหาสุทสฺสนสฺส อสฺสรตนํ ปาตุรโหสิ สพฺพเสโต กาฬสีโส มุญฺชเกโส อิทฺธิมา เวหาสงฺคโม วลาหโก นาม อสฺสราชา.\csromanspacer{} ตํ ทิสฺวา รญฺโญ มหาสุทสฺสนสฺส จิตฺตํ ปสีทิ “ภทฺทกํ วต โภ อสฺสยานํ สเจ ทมถํ อุเปยฺยา”ติ.\csromanspacer{} อถ โข ตํ อานนฺท อสฺสรตนํ เสยฺยถาปิ นาม ภทฺโท อสฺสาชานิโย ทีฆรตฺตํ สุปริทนฺโต, เอวเมว ทมถํ อุปคจฺฉิ.\csromanspacer{} ภูตปุพฺพํ อานนฺท ราชา มหาสุทสฺสโน ตเมว อสฺสรตนํ วีมํสมาโน ปุพฺพณฺหสมยํ อภิรุหิตฺวา สมุทฺทปริยนฺตํ ปถวิํ อนุยายิตฺวา กุสาวติํ ราชธานิํ ปจฺจาคนฺตฺวา ปาตราสมกาสิ.\csromanspacer{} รญฺโญ อานนฺท มหาสุทสฺสนสฺส เอวรูปํ อสฺสรตนํ ปาตุรโหสิ. (3)}

\csromanmatikaanchor{mtk.72}
\csromantocmarkref{section}{มณิรตน}{mtk.72}
\bookmark[dest={mtk.72},level=1]{มณิรตน}
\csromanheader{1}{\textbf{มณิรตน}}
\proseitem{248}{ปุน จปรํ อานนฺท รญฺโญ มหาสุทสฺสนสฺส มณิรตนํ ปาตุรโหสิ.\csromanspacer{} โส อโหสิ มณิ เวฬุริโย สุโภ ชาติมา อฏฺฐํโส สุปริกมฺมกโต อจฺโฉ วิปฺปสนฺโน อนาวิโล สพฺพาการสมฺปนฺโน.\csromanspacer{} ตสฺส โข ปนานนฺท มณิรตนสฺส อาภา สมนฺตา โยชนํ ผุฏา อโหสิ.\csromanspacer{} ภูตปุพฺพํ อานนฺท ราชา มหาสุทสฺสโน ตเมว มณิรตนํ วีมํสมาโน จตุรงฺคินิํ เสนํ สนฺนยฺหิตฺวา มณิํ ธชคฺคํ อาโรเปตฺวา รตฺตนฺธการติมิสาย ปายาสิ.\csromanspacer{} เย โข ปนานนฺท สมนฺตา คามา อเหสุํ, เต เตโนภาเสน กมฺมนฺเต ปโยเชสุํ ทิวาติ มญฺญมานา.\csromanspacer{} รญฺโญ อานนฺท มหาสุทสฺสนสฺส เอวรูปํ มณิรตนํ ปาตุรโหสิ. (4)}

\csromanmatikaanchor{mtk.73}
\csromantocmarkref{section}{อิตฺถิรตน}{mtk.73}
\bookmark[dest={mtk.73},level=1]{อิตฺถิรตน}
\csromanheader{1}{\textbf{อิตฺถิรตน}}
\proseitem{249}{ปุน จปรํ อานนฺท รญฺโญ มหาสุทสฺสนสฺส อิตฺถิรตนํ ปาตุรโหสิ อภิรูปา ทสฺสนียา ปาสาทิกา ปรมาย วณฺณโปกฺขรตาย สมนฺนาคตา นาติทีฆา นาติรสฺสา นาติกิสา นาติถูลา นาติกาฬิกา นาจฺโจทาตา อติกฺกนฺตา มานุสิวณฺณํ\footnote{มานุสฺสิวณฺณํ (สฺยา)} อปฺปตฺตา ทิพฺพวณฺณํ.\csromanspacer{} ตสฺส โข ปนานนฺท อิตฺถิรตนสฺส เอวรูโป กายสมฺผสฺโส โหติ, \csromanfolio{144}เสยฺยถาปิ นาม ตูลปิจุโน วา กปฺปาสปิจุโน วา.\csromanspacer{} ตสฺส โข ปนานนฺท อิตฺถิรตนสฺส สีเต อุณฺหานิ คตฺตานิ โหนฺติ, อุเณฺห สีตานิ.\csromanspacer{} ตสฺส โข ปนานนฺท อิตฺถิรตนสฺส กายโต จนฺทนคนฺโธ วายติ, มุขโต อุปฺปลคนฺโธ.\csromanspacer{} ตํ โข ปนานนฺท อิตฺถิรตนํ รญฺโญ มหาสุทสฺสนสฺส ปุพฺพุฏฺฐายินี อโหสิ ปจฺฉานิปาตินี กิํการปฏิสฺสาวินี มนาปจารินี ปิยวาทินี.\csromanspacer{} ตํ โข ปนานนฺท อิตฺถิรตนํ ราชานํ มหาสุทสฺสนํ มนสาปิ โน อติจริ\footnote{อติจรี (ก), อติจารี (สี, สฺยา, อิ)}, กุโต ปน กาเยน.\csromanspacer{} รญฺโญ อานนฺท มหาสุทสฺสนสฺส เอวรูปํ อิตฺถิรตนํ ปาตุรโหสิ. (5)}

\csromanmatikaanchor{mtk.74}
\csromantocmarkref{section}{คหปติรตน}{mtk.74}
\bookmark[dest={mtk.74},level=1]{คหปติรตน}
\csromanheader{1}{\textbf{คหปติรตน}}
\proseitem{250}{ปุน จปรํ อานนฺท รญฺโญ มหาสุทสฺสนสฺส คหปติรตนํ ปาตุรโหสิ.\csromanspacer{} ตสฺส กมฺมวิปากชํ ทิพฺพจกฺขุ ปาตุรโหสิ, เยน นิธิํ ปสฺสติ สสฺสามิกํปิ อสฺสามิกํปิ.\csromanspacer{} โส ราชานํ มหาสุทสฺสนํ อุปสงฺกมิตฺวา เอวมาห “อปฺโปสฺสุกฺโก ตฺวํ เทว โหหิ, อหํ เต ธเนน ธนกรณียํ กริสฺสามี”ติ.\csromanspacer{} ภูตปุพฺพํ อานนฺท ราชา มหาสุทสฺสโน ตเมว คหปติรตนํ วีมํสมาโน นาวํ อภิรุหิตฺวา มชฺเฌ คงฺคาย นทิยา โสตํ โอคาหิตฺวา คหปติรตนํ เอตทโวจ “อตฺโถ เม คหปติ หิรญฺญสุวณฺเณนา”ติ.\csromanspacer{} เตน หิ มหาราช เอกํ ตีรํ นาวา อุเปตูติ.\csromanspacer{} อิเธว เม คหปติ อตฺโถ หิรญฺญสุวณฺเณนาติ.\csromanspacer{} อถ โข ตํ อานนฺท คหปติรตนํ อุโภหิ หตฺเถหิ อุทกํ โอมสิตฺวา ปูรํ หิรญฺญสุวณฺณสฺส กุมฺภิํ อุทฺธริตฺวา ราชานํ มหาสุทสฺสนํ เอตทโวจ “อลเมตฺตาวตา มหาราช, กตเมตฺตาวตา มหาราช, ปูชิตเมตฺตาวตา มหาราชา”ติ.\csromanspacer{} ราชา มหาสุทสฺสโน เอวมาห “อลเมตฺตาวตา คหปติ, กตเมตฺตาวตา คหปติ, ปูชิตเมตฺตาวตา คหปตี”ติ.\csromanspacer{} รญฺโญ อานนฺท มหาสุทสฺสนสฺส เอวรูปํ คหปติรตนํ ปาตุรโหสิ. (6)}

\csromanmatikaanchor{mtk.75}
\csromantocmarkref{section}{ปริณายกรตน}{mtk.75}
\bookmark[dest={mtk.75},level=1]{ปริณายกรตน}
\csromanheader{1}{\textbf{ปริณายกรตน}}
\proseitem{251}{ปุน จปรํ อานนฺท รญฺโญ มหาสุทสฺสนสฺส ปริณายกรตนํ ปาตุรโหสิ ปณฺฑิโต วิยตฺโต เมธาวี ปฏิพโล ราชานํ}

\vspace{\tipitakaparskip}
\csromancenter{มหาสุทสฺสนสุตฺต มหาสุทสฺสนํ อุปยาเปตพฺพํ อุปยาเปตุํ อปยาเปตพฺพํ อปยาเปตุํ ฐเปตพฺพํ ฐเปตุํ.\csromanspacer{} โส ราชานํ มหาสุทสฺสนํ อุปสงฺกมิตฺวา เอวมาห “อปฺโปสฺสุกฺโก ตฺวํ เทว โหหิ, อหมนุสาสิสฺสามี”ติ.\csromanspacer{} รญฺโญ อานนฺท มหาสุทสฺสนสฺส เอวรูปํ ปริณายกรตนํ ปาตุรโหสิ. (7)}
\csromancenter{ราชา อานนฺท มหาสุทสฺสโน อิเมหิ สตฺตหิ รตเนหิ สมนฺนาคโต อโหสิ.}
\csromanmatikaanchor{mtk.76}
\csromantocmarkref{section}{จตุ-อิทฺธิสมนฺนาคต}{mtk.76}
\bookmark[dest={mtk.76},level=1]{จตุ-อิทฺธิสมนฺนาคต}
\csromanheader{1}{\textbf{จตุ-อิทฺธิสมนฺนาคต}}
\proseitem{252}{\csromanfolio{145}(ราชา อานนฺท มหาสุทสฺสโน จตูหิ อิทฺธีหิ สมนฺนาคโต อโหสิ.) กตมาหิ จตูหิ อิทฺธีหิ.\csromanspacer{} อิธานนฺท ราชา มหาสุทสฺสโน อภิรูโป อโหสิ ทสฺสนีโย ปาสาทิโก ปรมาย วณฺณโปกฺขรตาย สมนฺนาคโต อติวิย อญฺเญหิ มนุสฺเสหิ.\csromanspacer{} ราชา อานนฺท มหาสุทสฺสโน อิมาย ปฐมาย อิทฺธิยา สมนฺนาคโต อโหสิ.}

\prose{ปุน จปรํ อานนฺท ราชา มหาสุทสฺสโน ทีฆายุโก อโหสิ จิรฏฺฐิติโก อติวิย อญฺเญหิ มนุสฺเสหิ.\csromanspacer{} ราชา อานนฺท มหาสุทสฺสโน อิมาย ทุติยาย อิทฺธิยา สมนฺนาคโต อโหสิ.}

\prose{ปุน จปรํ อานนฺท ราชา มหาสุทสฺสโน อปฺปาพาโธ อโหสิ อปฺปาตงฺโก สมเวปากินิยา คหณิยา สมนฺนาคโต นาติสีตาย นาจฺจุณฺหาย อติวิย อญฺเญหิ มนุสฺเสหิ.\csromanspacer{} ราชา อานนฺท มหาสุทสฺสโน อิมาย ตติยาย อิทฺธิยา สมนฺนาคโต อโหสิ.}

\prose{ปุน จปรํ อานนฺท ราชา มหาสุทสฺสโน พฺราหฺมณคหปติกานํ ปิโย อโหสิ มนาโป.\csromanspacer{} เสยฺยถาปิ อานนฺท ปิตา ปุตฺตานํ ปิโย โหติ มนาโป, เอวเมว โข อานนฺท ราชา มหาสุทสฺสโน พฺราหฺมณคหปติกานํ ปิโย อโหสิ มนาโป, รญฺโญปิ อานนฺท มหาสุทสฺสนสฺส พฺราหฺมณคหปติกา ปิยา อเหสุํ มนาปา.\csromanspacer{} เสยฺยถาปิ อานนฺท ปิตุ ปุตฺตา ปิยา โหนฺติ มนาปา, เอวเมว โข อานนฺท รญฺโญปิ มหาสุทสฺสนสฺส พฺราหฺมณคหปติกา ปิยา อเหสุํ มนาปา.}

\prose{ภูตปุพฺพํ อานนฺท ราชา มหาสุทสฺสโน จตุรงฺคินิยา เสนาย อุยฺยานภูมิํ นิยฺยาสิ.\csromanspacer{} อถ โข อานนฺท พฺราหฺมณคหปติกา ราชานํ มหาสุทสฺสนํ อุปสงฺกมิตฺวา เอวมาหํสุ “อตรมาโน เทว ยาหิ, \csromanfolio{146}ยถา ตํ มยํ จิรตรํ ปสฺเสยฺยามา”ติ.\csromanspacer{} ราชาปิ อานนฺท มหาสุทสฺสโน สารถิํ อามนฺเตสิ “อตรมาโน สารถิ รถํ เปเสหิ, ยถา อหํ พฺราหฺมณคหปติเก จิรตรํ ปสฺเสยฺยนฺ”ติ.\csromanspacer{} ราชา อานนฺท มหาสุทสฺสโน อิมาย จตุตฺถิยา1 อิทฺธิยา สมนฺนาคโต อโหสิ.\csromanspacer{} ราชา อานนฺท มหาสุทสฺสโน อิมาหิ จตูหิ อิทฺธีหิ สมนฺนาคโต อโหสิ.}

\csromanmatikaanchor{mtk.77}
\csromantocmarkref{section}{ธมฺมปาสาทโปกฺขรณี}{mtk.77}
\bookmark[dest={mtk.77},level=1]{ธมฺมปาสาทโปกฺขรณี}
\csromanheader{1}{\textbf{ธมฺมปาสาทโปกฺขรณี}}
\proseitem{253}{อถ โข อานนฺท รญฺโญ มหาสุทสฺสนสฺส เอตทโหสิ “ยํนูนาหํ อิมาสุ ตาลนฺตริกาสุ ธนุสเต ธนุสเต โปกฺขรณิโย มาเปยฺยนฺ”ติ.}

\prose{มาเปสิ โข อานนฺท ราชา มหาสุทสฺสโน ตาสุ ตาลนฺตริกาสุ ธนุสเต ธนุสเต โปกฺขรณิโย.\csromanspacer{} ตา โข ปนานนฺท โปกฺขรณิโย จตุนฺนํ วณฺณานํ อิฏฺฐกาหิ จิตา อเหสุํ, เอกา อิฏฺฐกา โสวณฺณมยา เอกา รูปิยมยา เอกา เวฬุริยมยา เอกา ผลิกมยา.}

\prose{ตาสุ โข ปนานนฺท โปกฺขรณีสุ จตฺตาริ จตฺตาริ โสปานานิ อเหสุํ จตุนฺนํ วณฺณานํ, เอกํ โสปานํ โสวณฺณมยํ เอกํ รูปิยมยํ เอกํ เวฬุริยมยํ เอกํ ผลิกมยํ.\csromanspacer{} โสวณฺณมยสฺส โสปานสฺส โสวณฺณมยา ถมฺภา อเหสุํ รูปิยมยา สูจิโย จ อุณฺหีสญฺจ.\csromanspacer{} รูปิยมยสฺส โสปานสฺส รูปิยมยา ถมฺภา อเหสุํ โสวณฺณมยา สูจิโย จ อุณฺหีสญฺจ.\csromanspacer{} เวฬุริยมยสฺส โสปานสฺส เวฬุริยมยา ถมฺภา อเหสุํ ผลิกมยา สูจิโย จ อุณฺหีสญฺจ.\csromanspacer{} ผลิกมยสฺส โสปานสฺส ผลิกมยา ถมฺภา อเหสุํ เวฬุริยมยา สูจิโย จ อุณฺหีสญฺจ.\csromanspacer{} ตา โข ปนานนฺท โปกฺขรณิโย ทฺวีหิ เวทิกาหิ ปริกฺขิตฺตา อเหสุํ เอกา เวทิกา โสวณฺณมยา เอกา รูปิยมยา.\csromanspacer{} โสวณฺณมยาย เวทิกาย โสวณฺณมยา ถมฺภา อเหสุํ รูปิยมยา สูจิโย จ อุณฺหีสญฺจ.\csromanspacer{} รูปิยมยาย เวทิกาย รูปิยมยา ถมฺภา อเหสุํ โสวณฺณมยา สูจิโย จ อุณฺหีสญฺจ.\csromanspacer{} อถ โข อานนฺท รญฺโญ มหาสุทสฺสนสฺส เอตทโหสิ “ยํนูนาหํ อิมาสุ โปกฺขรณีสุ เอวรูปํ มาลํ โรปาเปยฺยํ อุปฺปลํ ปทุมํ กุมุทํ ปุณฺฑรีกํ สพฺโพตุกํ สพฺพชนสฺส อนาวฏนฺ”ติ.\csromanspacer{} โรปาเปสิ}

\vspace{\tipitakaparskip}
\csromancenter{จตุตฺถาย (สฺยา)}
\csromancenter{มหาสุทสฺสนสุตฺต โข อานนฺท ราชา มหาสุทสฺสโน ตาสุ โปกฺขรณีสุ เอวรูปํ มาลํ อุปฺปลํ ปทุมํ กุมุทํ ปุณฺฑรีกํ สพฺโพตุกํ สพฺพชนสฺส อนาวฏํ.}
\proseitem{254}{\csromanfolio{147}อถ โข อานนฺท รญฺโญ มหาสุทสฺสนสฺส เอตทโหสิ “ยํนูนาหํ อิมาสํ โปกฺขรณีนํ ตีเร นฺหาปเก ปุริเส ฐเปยฺยํ, เย อาคตาคตํ ชนํ นฺหาเปสฺสนฺตี”ติ.\csromanspacer{} ฐเปสิ โข อานนฺท ราชา มหาสุทสฺสโน ตาสํ โปกฺขรณีนํ ตีเร นฺหาปเก ปุริเส, เย อาคตาคตํ ชนํ นฺหาเปสุํ.}

\prose{อถ โข อานนฺท รญฺโญ มหาสุทสฺสนสฺส เอตทโหสิ “ยํนูนาหํ อิมาสํ โปกฺขรณีนํ ตีเร เอวรูปํ ทานํ ปฏฺฐเปยฺยํ อนฺนํ อนฺนฏฺฐิกสฺส1 ปานํ ปานฏฺฐิกสฺส วตฺถํ วตฺถฏฺฐิกสฺส ยานํ ยานฏฺฐิกสฺส สยนํ สยนฏฺฐิกสฺส อิตฺถิํ อิตฺถิฏฺฐิกสฺส หิรญฺญํ หิรญฺญฏฺฐิกสฺส สุวณฺณํ สุวณฺณฏฺฐิกสฺสา”ติ.\csromanspacer{} ปฏฺฐเปสิ โข อานนฺท ราชา มหาสุทสฺสโน ตาสํ โปกฺขรณีนํ ตีเร เอวรูปํ ทานํ อนฺนํ อนฺนฏฺฐิกสฺส ปานํ ปานฏฺฐิกสฺส วตฺถํ วตฺถฏฺฐิกสฺส ยานํ ยานฏฺฐิกสฺส สยนํ สยนฏฺฐิกสฺส อิตฺถิํ อิตฺถิฏฺฐิกสฺส หิรญฺญํ หิรญฺญฏฺฐิกสฺส สุวณฺณํ สุวณฺณฏฺฐิกสฺส.}

\proseitem{255}{อถ โข อานนฺท พฺราหฺมณคหปติกา ปหูตํ สาปเตยฺยํ อาทาย ราชานํ มหาสุทสฺสนํ อุปสงฺกมิตฺวา เอวมาหํสุ “อิทํ เทว ปหูตํ สาปเตยฺยํ เทวญฺเญว อุทฺทิสฺส อาภตํ, ตํ เทโว ปฏิคฺคณฺหตู”ติ.\csromanspacer{} อลํ โภ มมปิทํ ปหูตํ สาปเตยฺยํ ธมฺมิเกน พลินา อภิสงฺขตํ, ตญฺจ โว โหตุ, อิโต จ ภิยฺโย หรถาติ.\csromanspacer{} เต รญฺญา ปฏิกฺขิตฺตา เอกมนฺตํ อปกฺกมฺม เอวํ สมจินฺเตสุํ “น โข เอตํ อมฺหากํ ปติรูปํ, ยํ มยํ อิมานิ สาปเตยฺยานิ ปุนเทว สกานิ ฆรานิ ปฏิหเรยฺยาม.\csromanspacer{} ยํนูน มยํ รญฺโญ มหาสุทสฺสนสฺส นิเวสนํ มาเปยฺยามา”ติ.\csromanspacer{} เต ราชานํ มหาสุทสฺสนํ อุปสงฺกมิตฺวา เอวมาหํสุ “นิเวสนํ เต เทว มาเปสฺสามา”ติ.\csromanspacer{} อธิวาเสสิ โข อานนฺท ราชา มหาสุทสฺสโน ตุณฺหีภาเวน.}

\proseitem{256}{อถ โข อานนฺท สกฺโก เทวานมินฺโท รญฺโญ มหาสุทสฺสนสฺส เจตสา เจโตปริวิตกฺกมญฺญาย วิสฺสกมฺมํ2 เทวปุตฺตํ อามนฺเตสิ}

\proseitem{1}{อนฺนตฺถิตสฺส (สี, สฺยา, กํ, อิ), เอวํ สพฺพตฺถ ปกภิรูเปเนว ทิสฺสติ.}

\proseitem{2}{วิสุกมฺมํ (ก) \csromanfolio{148}“เอหิ ตฺวํ สมฺม วิสฺสกมฺม, รญฺโญ มหาสุทสฺสนสฺส นิเวสนํ มาเปหิ ธมฺมํ นาม ปาสาทนฺ”ติ. “เอวํ ภทฺทนฺตวา”ติ โข อานนฺท วิสฺสกมฺโม เทวปุตฺโต สกฺกสฺส เทวานมินฺทสฺส ปฏิสฺสุตฺวา เสยฺยถาปิ นาม พลวา ปุริโส สมิญฺชิตํ วา พาหํ ปสาเรยฺย ปสาริตํ วา พาหํ สมิญฺเชยฺย, เอวเมว เทเวสุ ตาวติํเสสุ อนฺตรหิโต รญฺโญ มหาสุทสฺสนสฺส ปุรโต ปาตุรโหสิ.\csromanspacer{} อถ โข อานนฺท วิสฺสกมฺโม เทวปุตฺโต ราชานํ มหาสุทสฺสนํ เอตทโวจ “นิเวสนํ เต เทว มาเปสฺสามิ ธมฺมํ นาม ปาสาทนฺ”ติ.\csromanspacer{} อธิวาเสสิ โข อานนฺท ราชา มหาสุทสฺสโน ตุณฺหีภาเวน.}

\prose{มาเปสิ โข อานนฺท วิสฺสกมฺโม เทวปุตฺโต รญฺโญ มหาสุทสฺสนสฺส นิเวสนํ ธมฺมํ นาม ปาสาทํ.\csromanspacer{} ธมฺโม อานนฺท ปาสาโท ปุรตฺถิเมน ปจฺฉิเมน จ โยชนํ อายาเมน อโหสิ.\csromanspacer{} อุตฺตเรน ทกฺขิเณน จ อฑฺฒโยชนํ วิตฺถาเรน.\csromanspacer{} ธมฺมสฺส อานนฺท ปาสาทสฺส ติโปริสํ อุจฺจตเรน วตฺถุ จิตํ อโหสิ จตุนฺนํ วณฺณานํ อิฏฺฐกาหิ, เอกา อิฏฺฐกา โสวณฺณมยา เอกา รูปิยมยา เอกา เวฬุริยมยา เอกา ผลิกมยา.}

\prose{ธมฺมสฺส อานนฺท ปาสาทสฺส จตุราสีติ ถมฺภสหสฺสานิ อเหสุํ จตุนฺนํ วณฺณานํ, เอโก ถมฺโภ โสวณฺณมโย เอโก รูปิยมโย เอโก เวฬุริยมโย เอโก ผลฺลิกมโย.\csromanspacer{} ธมฺโม อานนฺท ปาสาโท จตุนฺนํ วณฺณานํ ผลเกหิ สนฺถโต อโหสิ, เอกํ ผลกํ โสวณฺณมยํ เอกํ รูปิยมยํ เอกํ เวฬุริยมยํ เอกํ ผลิกมยํ.}

\prose{ธมฺมสฺส อานนฺท ปาสาทสฺส จตุวีสติ โสปานานิ อเหสุํ จตุนฺนํ วณฺณานํ, เอกํ โสปานํ โสวณฺณมยํ เอกํ รูปิยมยํ เอกํ เวฬุริยมยํ เอกํ ผลฺลิกมยํ.\csromanspacer{} โสวณฺณมยสฺส โสปานสฺส โสวณฺณมยา ถมฺภา อเหสุํ รูปิยมยา สูจิโย จ อุณฺหีสญฺจ.\csromanspacer{} รูปิยมยสฺส โสปานสฺส รูปิยมยา ถมฺภา อเหสุํ โสวณฺณมยา สูจิโย จ อุณฺหีสญฺจ.\csromanspacer{} เวฬุริยมยสฺส โสปานสฺส เวฬุริยมยา ถมฺภา อเหสุํ ผลฺลิกมยา สูจิโย จ อุณฺหีสญฺจ.\csromanspacer{} ผลิกมยสฺส โสปานสฺส ผลิกมยา ถมฺภา อเหสุํ เวฬุริยมยา สูจิโย จ อุณฺหีสญฺจ.}

\prose{ธมฺเม อานนฺท ปาสาเท จตุราสีติ กูฏาคารสหสฺสานิ อเหสุํ จตุนฺนํ วณฺณานํ, เอกํ กูฏาคารํ โสวณฺณมยํ เอกํ รูปิยมยํ เอกํ}

\vspace{\tipitakaparskip}
\csromancenter{มหาสุทสฺสนสุตฺต เวฬุริยมยํ เอกํ ผลิกมยํ.\csromanspacer{} โสวณฺณมเย กูฏาคาเร รูปิยมโย ปลฺลงฺโก ปญฺญตฺโต อโหสิ, รูปิยมเย กูฏาคาเร โสวณฺณมโย ปลฺลงฺโก ปญฺญตฺโต อโหสิ, เวฬุริย-มเย กูฏาคาเร ทนฺตมโย ปลฺลงฺโก ปญฺญตฺโต อโหสิ, ผลิกมเย กูฏาคาเร สารมโย ปลฺลงฺโก ปญฺญตฺโต อโหสิ.\csromanspacer{} โสวณฺณมยสฺส กูฏาคารสฺส ทฺวาเร รูปิยมโย ตาโล ฐิโต อโหสิ, ตสฺส รูปิยมโย ขนฺโธ โสวณฺณมยานิ ปตฺตานิ จ ผลานิ จ.\csromanspacer{} รูปิยมยสฺส กูฏาคารสฺส ทฺวาเร โสวณฺณมโย ตาโล ฐิโต อโหสิ, ตสฺส โสวณฺณมโย ขนฺโธ, รูปิยมยานิ ปตฺตานิ จ ผลานิ จ.\csromanspacer{} เวฬุริยมยสฺส กูฏาคารสฺส ทฺวาเร ผลิกมโย ตาโล ฐิโต อโหสิ, ตสฺส ผลิกมโย ขนฺโธ, เวฬุริยมยานิ ปตฺตานิ จ ผลานิ จ.\csromanspacer{} ผลิกมยสฺส กูฏาคารสฺส ทฺวาเร เวฬุริยมโย ตาโล ฐิโต อโหสิ, ตสฺส เวฬุริยมโย ขนฺโธ, ผลิกมยานิ ปตฺตานิ จ ผลานิ จ.}
\proseitem{257}{\csromanfolio{149}อถ โข อานนฺท รญฺโญ มหาสุทสฺสนสฺส เอตทโหสิ “ยํนูนาหํ มหาวิยูหสฺส กูฏาคารสฺส ทฺวาเร สพฺพโสวณฺณมยํ ตาลวนํ มาเปยฺยํ, ยตฺถ ทิวาวิหารํ นิสีทิสฺสามี”ติ.\csromanspacer{} มาเปสิ โข อานนฺท ราชา มหาสุทสฺสโน มหาวิยูหสฺส กูฏาคารสฺส ทฺวาเร สพฺพโสวณฺณมยํ ตาลวนํ, ยตฺถ ทิวาวิหารํ นิสีทิ.\csromanspacer{} ธมฺโม อานนฺท ปาสาโท ทฺวีหิ เวทิกาหิ ปริกฺขิตฺโต อโหสิ, เอกา เวทิกา โสวณฺณมยา เอกา รูปิยมยา.\csromanspacer{} โสวณฺณมยาย เวทิกาย โสวณฺณมยา ถมฺภา อเหสุํ รูปิยมยา สูจิโย จ อุณฺหีสญฺจ.\csromanspacer{} รูปิยมยาย เวทิกาย รูปิยมยา ถมฺภา อเหสุํ โสวณฺณมยา สูจิโย จ อุณฺหีสญฺจ.}

\proseitem{258}{ธมฺโม อานนฺท ปาสาโท ทฺวีหิ กิงฺกิณิกชาเลหิ1 ปริกฺขิตฺโต อโหสิ, เอกํ ชาลํ โสวณฺณมยํ เอกํ รูปิยมยํ.\csromanspacer{} โสวณฺณมยสฺส ชาลสฺส รูปิยมยา กิงฺกิณิกา อเหสุํ, รูปิยมยสฺส ชาลสฺส โสวณฺณมยา กิงฺกิณิกา อเหสุํ.\csromanspacer{} เตสํ โข ปนานนฺท กิงฺกิณิกชาลานํ วาเตริตานํ สทฺโท อโหสิ วคฺคุ จ รชนีโย จ ขมนีโย จ มทนีโย จ.\csromanspacer{} เสยฺยถาปิ อานนฺท ปญฺจงฺคิกสฺส ตูริยสฺส สุวินีตสฺส}

\vspace{\tipitakaparskip}
\csromancenter{กิงฺกณิกชาเลหิ (สฺยา, ก)}
\prosecont{\csromanfolio{150}สุปฺปฏิตาฬิตสฺส สุกุสเลหิ\footnote{กุสเลหิ (สี, สฺยา, กํ, อิ)} สมนฺนาหตสฺส สทฺโท โหติ วคฺคุ จ รชนีโย จ ขมนีโย จ มทนีโย จ.\csromanspacer{} เอวเมว โข อานนฺท เตสํ กิงฺกิณิกชาลานํ วาเตริตานํ สทฺโท อโหสิ วคฺคุ จ รชนีโย จ กมนีโย จ มทนีโย จ.\csromanspacer{} เย โข ปนานนฺท เตน สมเยน กุสาวติยา ราชธานิยา ธุตฺตา อเหสุํ โสณฺฑา ปิปาสา, เต เตสํ กิงฺกิณิกชาลานํ วาเตริตานํ สทฺเทน ปริจาเรสุํ.\csromanspacer{} นิฏฺฐิโต โข ปนานนฺท ธมฺโม ปาสาโท ทุทฺทิกฺโข อโหสิ มุสติ จกฺขูนิ.\csromanspacer{} เสยฺยถาปิ อานนฺท วสฺสานํ ปจฺฉิเม มาเส สรทสมเย วิทฺเธ วิคตวลาหเก เทเว อาทิจฺโจ นภํ อพฺภุสฺสกฺกมาโน\footnote{อพฺภุคฺคมมาโน (สี, อิ, ก)} ทุทฺทิกฺโข\footnote{ทุทิกฺโข (อิ)} โหติ มุสติ จกฺขูนิ, เอวเมว โข อานนฺท ธมฺโม ปาสาโท ทุทฺทิกฺโข อโหสิ มุสติ จกฺขูนิ.}

\proseitem{259}{อถ โข อานนฺท รญฺโญ มหาสุทสฺสนสฺส เอตทโหสิ. “ยํนูนาหํ ธมฺมสฺส ปาสาทสฺส ปุรโต ธมฺมํ นาม โปกฺขรณิํ มาเปยฺยนฺ”ติ.\csromanspacer{} มาเปสิ โข อานนฺท ราชา มหาสุทสฺสโน ธมฺมสฺส ปาสาทสฺส ปุรโต ธมฺมํ นาม โปกฺขรณิํ.\csromanspacer{} ธมฺมา อานนฺท โปกฺขรณี ปุรตฺถิเมน ปจฺฉิเมน จ โยชนํ อายาเมน อโหสิ, อุตฺตเรน ทกฺขิเณน จ อฑฺฒโยชนํ วิตฺถาเรน.\csromanspacer{} ธมฺมา อานนฺท โปกฺขรณี จตุนฺนํ วณฺณานํ อิฏฺฐกาหิ จิตา อโหสิ, เอกา อิฏฺฐกา โสวณฺณมยา เอกา รูปิยมยา เอกา เวฬุริยมยา เอกา ผลิกมยา.}

\prose{ธมฺมาย อานนฺท โปกฺขรณิยา จตุวีสติ โสปานานิ อเหสุํ จตุนฺนํ วณฺณานํ, เอกํ โสปานํ โสวณฺณมยํ เอกํ รูปิยมยํ เอกํ เวฬุริยมยํ เอกํ ผลิกมยํ.\csromanspacer{} โสวณฺณมยสฺส โสปานสฺส โสวณฺณมยา ถมฺภา อเหสุํ รูปิยมยา สูจิโย จ อุณฺหีสญฺจ.\csromanspacer{} รูปิยมยสฺส โสปานสฺส รูปิยมยา ถมฺภา อเหสุํ โสวณฺณมยา สูจิโย จ อุณฺหีสญฺจ.\csromanspacer{} เวฬุริยมยสฺส โสปานสฺส เวฬุริยมยา ถมฺภา อเหสุํ ผลฺลิกมยา สูจิโย จ อุณฺหีสญฺจ.\csromanspacer{} ผลิกมยสฺส โสปานสฺส ผลิกมยา ถมฺภา อเหสุํ เวฬุริยมยา สุจิโย จ อุณฺหีสญฺจ.}

\prose{ธมฺมา อานนฺท โปกฺขรณี ทฺวีหิ เวทิกาหิ ปริกฺขิตฺตา อโหสิ, เอกา เวทิกา โสวณฺณมยา เอกา รูปิยมยา.\csromanspacer{} โสวณฺณมยาย เวทิกาย}

\vspace{\tipitakaparskip}
\csromancenter{มหาสุทสฺสนสุตฺต โสวณฺณมยา ถมฺภา อเหสุํ รูปิยมยา สูจิโย จ อุณฺหีสญฺจ.\csromanspacer{} รูปิยมยาย เวทิกาย รูปิยมยา ถมฺภา อเหสุํ โสวณฺณมยา สูจิโย จ อุณฺหีสญฺจ.}
\prose{\csromanfolio{151}ธมฺมา อานนฺท โปกฺขรณี สตฺตหิ ตาลปนฺตีหิ ปริกฺขิตฺตา อโหสิ, เอกา ตาลปนฺติ โสวณฺณมยา เอกา รูปิยมยา เอกา เวฬุริยมยา เอกา ผลิกมยา เอกา โลหิตงฺกมยา เอกา มสารคลฺลมยา เอกา สพฺพรตนมยา.\csromanspacer{} โสวณฺณมยสฺส ตาลสฺส โสวณฺณมโย ขนฺโธ อโหสิ, รูปิยมยานิ ปตฺตานิ จ ผลานิ จ.\csromanspacer{} รูปิยมยสฺส ตาลสฺส รูปิยมโย ขนฺโธ อโหสิ, โสวณฺณมยานิ ปตฺตานิ จ ผลานิ จ.\csromanspacer{} เวฬุริยมยสฺส ตาลสฺส เวฬุริยมโย ขนฺโธ อโหสิ, ผลิกมยานิ ปตฺตานิ จ ผลานิ จ.\csromanspacer{} ผลิกมยสฺส ตาลสฺส ผลิกมโย ขนฺโธ อโหสิ, เวฬุริยมยานิ ปตฺตานิ จ ผลานิ จ.\csromanspacer{} โลหิตงฺกมยสฺส ตาลสฺส โลหิตงฺกมโย ขนฺโธ อโหสิ, มสารคลฺลมยานิ ปตฺตานิ จ ผลานิ จ.\csromanspacer{} มสารคลฺลมยสฺส ตาลสฺส มสารคลฺลมโย ขนฺโธ อโหสิ, โลหิตงฺกมยานิ ปตฺตานิ จ ผลานิ จ.\csromanspacer{} สพฺพรตนมยสฺส ตาลสฺส สพฺพรตนมโย ขนฺโธ อโหสิ, สพฺพรตนมยานิ ปตฺตานิ จ ผลานิ จ.\csromanspacer{} ตาสํ โข ปนานนฺท ตาลปนฺตีนํ วาเตริตานํ สทฺโท อโหสิ วคฺคุ จ รชนีโย จ ขมนีโย จ มทนีโย จ.\csromanspacer{} เสยฺยถาปิ อานนฺท ปญฺจงฺคิกสฺส ตูริยสฺส สุวินีตสฺส สุปฺปฏิตาฬิตสฺส สุกุสเลหิ สมนฺนาหตสฺส สทฺโท โหติ วคฺคุ จ รชนีโย จ กมนีโย จ มทนีโย จ, เอวเมว โข อานนฺท ตาสํ ตาลปนฺตีนํ วาเตริตานํ สทฺโท อโหสิ วคฺคุ จ รชนีโย จ ขมนีโย จ มทนีโย จ.\csromanspacer{} เย โข ปนานนฺท เตน สมเยน กุสาวติยา ราชธานิยา ธุตฺตา อเหสุํ โสณฺฑา ปิปาสา, เต ตาสํ ตาลปนฺตีนํ วาเตริตานํ สทฺเทน ปริจาเรสุํ.}

\prose{นิฏฺฐิเต โข ปนานนฺท ธมฺเม ปาสาเท นิฏฺฐิตาย ธมฺมาย จ โปกฺขรณิยา ราชา มหาสุทสฺสโน ‘เย\footnote{เย โข ปนานนฺท (สฺยา, ก)} เตน สมเยน สมเณสุ วา สมณสมฺมตา พฺราหฺมเณสุ วา พฺราหฺมณสมฺมตา’, เต สพฺพกาเมหิ สนฺตปฺเปตฺวา ธมฺมํ ปาสาทํ อภิรุหิ.}

\vspace{\tipitakaparskip}
\csromancenter{ปฐมภาณวาโร.}
\csromanmatikaanchor{mtk.78}
\csromantocmarkref{section}{ฌานสมฺปตฺติ}{mtk.78}
\bookmark[dest={mtk.78},level=1]{ฌานสมฺปตฺติ}
\csromanheader{1}{\textbf{ฌานสมฺปตฺติ}}
\proseitem{260}{\csromanfolio{152}อถ โข อานนฺท รญฺโญ มหาสุทสฺสนสฺส เอตทโหสิ “กิสฺส นุ โข เม อิทํ กมฺมสฺส ผลํ กิสฺส กมฺมสฺส วิปาโก, เยนาหํ เอตรหิ เอวํมหิทฺธิโก เอวํมหานุภาโว”ติ.\csromanspacer{} อถ โข อานนฺท รญฺโญ มหาสุทสฺสนสฺส เอตทโหสิ “ติณฺณํ โข เม อิทํ กมฺมานํ ผลํ ติณฺณํ กมฺมานํ วิปาโก, เยนาหํ เอตรหิ เอวํมหิทฺธิโก เอวํมหานุภาโว.\csromanspacer{} เสยฺยถิทํ, ทานสฺส ทมสฺส สํยมสฺสา”ติ.}

\prose{อถ โข อานนฺท ราชา มหาสุทสฺสโน เยน มหาวิยูหํ กูฏาคารํ เตนุปสงฺกมิ, อุปสงฺกมิตฺวา มหาวิยูหสฺส กูฏาคารสฺส ทฺวาเร ฐิโต อุทานํ อุทาเนสิ “ติฏฺฐ กามวิตกฺก, ติฏฺฐ พฺยาปาทวิตกฺก, ติฏฺฐ วิหิํสาวิตกฺก, เอตฺตาวตา กามวิตกฺก, เอตฺตาวตา พฺยาปาทวิตกฺก, เอตฺตาวตา วิหิํสาวิตกฺกา”ติ.}

\proseitem{261}{อถ โข อานนฺท ราชา มหาสุทสฺสโน มหาวิยูหํ กูฏาคารํ ปวิสิตฺวา โสวณฺณมเย ปลฺลงฺเก นิสฺสินฺโน วิวิจฺเจว กาเมหิ วิวิจฺจ อกุสเลหิ ธมฺเมหิ สวิตกฺกํ สวิจารํ วิเวกชํ ปีติสุขํ ปฐมํ ฌานํ อุปสมฺปชฺช วิหาสิ, วิตกฺกวิจารานํ วูปสมา อชฺฌตฺตํ สมฺปสาทนํ เจตโส เอโกทิภาวํ อวิตกฺกํ อวิจารํ สมาธิชํ ปีติสุขํ ทุติยํ ฌานํ อุปสมฺปชฺช วิหาสิ, ปีติยา จ วิราคา อุเปกฺขโก จ วิหาสิ, สโต จ สมฺปชาโน สุขญฺจ กาเยน ปฏิสํเวเทสิ, ยํ ตํ อริยา อาจิกฺขนฺติ “อุเปกฺขโก สติมา สุขวิหารี”ติ ตติยํ ฌานํ อุปสมฺปชฺช วิหาสิ, สุขสฺส จ ปหานา ทุกฺขสฺส จ ปหานา ปุพฺเพว โสมนสฺสโทมนสฺสานํ อตฺถงฺคมา อทุกฺขมสุขํ อุเปกฺขาสติปาริสุทฺธิํ จตุตฺถํ ฌานํ อุปสมฺปชฺช วิหาสิ.}

\proseitem{262}{อถ โข อานนฺท ราชา มหาสุทสฺสโน มหาวิยูหา กูฏาคารา นิกฺขมิตฺวา โสวณฺณมยํ กูฏาคารํ ปวิสิตฺวา รูปิยมเย ปลฺลงฺเก นิสินฺโน เมตฺตาสหคเตน เจตสา เอกํ ทิสํ ผริตฺวา วิหาสิ.\csromanspacer{} ตถา ทุติยํ ตถา ตติยํ ตถา จตุตฺถํ.\csromanspacer{} อิติ อุทฺธมโธ ติริยํ สพฺพธิ สพฺพตฺตตาย สพฺพาวนฺตํ โลกํ เมตฺตาสหคเตน เจตสา วิปุเลน มหคฺคเตน อปฺปมาเณน อเวเรน อพฺยาปชฺเชน ผริตฺวา วิหาสิ.}

\vspace{\tipitakaparskip}
\csromancenter{มหาสุทสฺสนสุตฺต กรุณาสหคเตน เจตสา ฯเปฯ.\csromanspacer{} มุทิตาสหคเตน เจตสา ฯเปฯ.\csromanspacer{} อุเปกฺขาสหคเตน เจตสา เอกํ ทิสํ ผริตฺวา วิหาสิ.\csromanspacer{} ตถา ทุติยํ ตถา ตติยํ ตถา จตุตฺถํ.\csromanspacer{} อิติ อุทฺธมโธ ติริยํ สพฺพธิ สพฺพตฺตตาย สพฺพาวนฺตํ โลกํ อุเปกฺขาสหคเตน เจตสา วิปุเลน มหคฺคเตน อปฺปมาเณน อเวเรน อพฺยาปชฺเชน ผริตฺวา วิหาสิ.}
\csromanmatikaanchor{mtk.79}
\csromantocmarkref{section}{จตุราสีตินครสหสฺสาทิ}{mtk.79}
\bookmark[dest={mtk.79},level=1]{จตุราสีตินครสหสฺสาทิ}
\csromanheader{1}{\textbf{จตุราสีตินครสหสฺสาทิ}}
\proseitem{263}{\csromanfolio{153}รญฺโญ อานนฺท มหาสุทสฺสนสฺส จตุราสีติ นครสหสฺสานิ อเหสุํ กุสาวตีราชธานิปฺปมุขานิ, จตุราสีติ ปาสาทสหสฺสานิ อเหสุํ ธมฺมปาสาทปฺปมุขานิ, จตุราสีติ กูฏาคารสหสฺสานิ อเหสุํ มหาวิยูหกูฏาคารปฺปมุขานิ, จตุราสีติ ปลฺลงฺกสหสฺสานิ อเหสุํ โสวณฺณมยานิ รูปิยมยานิ ทนฺตมยานิ สารมยานิ โคนกตฺถตานิ ปฏิกตฺถตานิ ปฏลิกตฺถตานิ กทลิมิคปวรปจฺจตฺถรณานิ ส- อุตฺตรจฺฉทานิ อุภโตโลหิตกูปธานานิ, จตุราสีติ นาคสหสฺสานิ อเหสุํ โสวณฺณาลงฺการานิ โสวณฺณธชานิ เหมชาลปฏิจฺฉนฺนานิ โอโปสถนาคราชปฺปมุขานิ, จตุราสีติ อสฺสสหสฺสานิ อเหสุํ โสวณฺณาลงฺการานิ โสวณฺณธชานิ เหมชาลปฏิจฺฉนฺนานิ วลาหก- อสฺสราชปฺปมุขานิ, จตุราสีติ รถสหสฺสานิ อเหสุํ สีหจมฺมปริวารานิ พฺยคฺฆจมฺมปริวารานิ ทีปิจมฺมปริวารานิ ปณฺฑุกมฺพลปริวารานิ โสวณฺณาลงฺการานิ โสวณฺณธชานิ เหมชาลปฏิจฺฉนฺนานิ เวชยนฺตรถปฺปมุขานิ, จตุราสีติ มณิสหสฺสานิ อเหสุํ มณิรตนปฺปมุขานิ, จตุราสีติ อิตฺถิสหสฺสานิ อเหสุํ สุภทฺทาเทวิปฺปมุขานิ, จตุราสีติ คหปติสหสฺสานิ อเหสุํ คหปติรตนปฺปมุขานิ, จตุราสีติ ขตฺติยสหสฺสานิ อเหสุํ อนุยนฺตานิ ปริณายกรตนปฺปมุขานิ, จตุราสีติ เธนุสหสฺสานิ อเหสุํ ทุหสนฺทนานิ\footnote{ทุกูลสนฺทนานิ (อิ) ทุกูลสนฺทานานิ (สํ 2. 118)} กํสูปธารณานิ, จตุราสีติ วตฺถโกฏิสหสฺสานิ อเหสุํ โขมสุขุมานํ กปฺปาสิกสุขุมานํ โกเสยฺยสุขุมานํ กมฺพลสุขุมานํ, (รญฺโญ อานนฺท มหาสุทสฺสนสฺส)\footnote{( ) สี-อิ-โปตฺถเกสุ นตฺถิ.} จตุราสีติ ถาลิปากสหสฺสานิ อเหสุํ สายํ ปาตํ ภตฺตาภิหาโร อภิหริยิตฺถ.}

\proseitem{264}{\csromanfolio{154}เตน โข ปนานนฺท สมเยน รญฺโญ มหาสุทสฺสนสฺส จตุราสีติ นาคสหสฺสานิ สายํ ปาตํ อุปฏฺฐานํ อาคจฺฉนฺติ.\csromanspacer{} อถ โข อานนฺท รญฺโญ มหาสุทสฺสนสฺส เอตทโหสิ “อิมานิ โข เม จตุราสีติ นาคสหสฺสานิ สายํ ปาตํ อุปฏฺฐานํ อาคจฺฉนฺติ, ยํนูน วสฺสสตสฺส วสฺสสตสฺส อจฺจเยน เทฺวจตฺตาลีสํ เทฺวจตฺตาลีสํ นาคสหสฺสานิ สกิํ สกิํ อุปฏฺฐานํ อาคจฺเฉยฺยุนฺ”ติ.\csromanspacer{} อถ โข อานนฺท ราชา มหาสุทสฺสโน ปริณายกรตนํ อามนฺเตสิ “อิมานิ โข เม สมฺม ปริณายกรตน จตุราสีติ นาคสหสฺสานิ สายํ ปาตํ อุปฏฺฐานํ อาคจฺฉนฺติ, เตน หิ สมฺม ปริณายกรตน วสฺสสตสฺส วสฺสสตสฺส อจฺจเยน เทฺวจตฺตาลีสํ เทฺวจตฺตาลีสํ นาคสหสฺสานิ สกิํ สกิํ อุปฏฺฐานํ อาคจฺฉนฺตู”ติ. “เอวํ เทวา”ติ โข อานนฺท ปริณายกรตนํ รญฺโญ มหาสุทสฺสนสฺส ปจฺจสฺโสสิ.\csromanspacer{} อถ โข อานนฺท รญฺโญ มหาสุทสฺสนสฺส อปเรน สมเยน วสฺสสตสฺส วสฺสสตสฺส อจฺจเยน เทฺวจตฺตาลีสํ เทฺวจตฺตาลีสํ นาคสหสฺสานิ สกิํ สกิํ อุปฏฺฐานํ อาคมํสุ.}

\csromanmatikaanchor{mtk.80}
\csromantocmarkref{section}{สุภทฺทาเทวิ-อุปสงฺกมน}{mtk.80}
\bookmark[dest={mtk.80},level=1]{สุภทฺทาเทวิ-อุปสงฺกมน}
\csromanheader{1}{\textbf{สุภทฺทาเทวิ-อุปสงฺกมน}}
\proseitem{265}{อถ โข อานนฺท สุภทฺทาย เทวิยา พหุนฺนํ วสฺสานํ พหุนฺนํ วสฺสสตานํ พหุนฺนํ วสฺสสหสฺสานํ อจฺจเยน เอตทโหสิ “จิรํ ทิฏฺโฐ โข เม ราชา มหาสุทสฺสโน, ยํนูนาหํ ราชานํ มหาสุทสฺสนํ ทสฺสนาย อุปสงฺกเมยฺยนฺ”ติ.\csromanspacer{} อถ โข อานนฺท สุภทฺทา เทวี อิตฺถาคารํ อามนฺเตสิ “เอถ ตุเมฺห สีสานิ นฺหายถ ปีตานิ วตฺถานิ ปารุปถ, จิรํ ทิฏฺโฐ โน ราชา มหาสุทสฺสโน, ราชานํ มหาสุทสฺสนํ ทสฺสนาย อุปสงฺกมิสฺสามา”ติ. “เอวํ อยฺเย”ติ โข อานนฺท อิตฺถาคารํ สุภทฺทาย เทวิยา ปฏิสฺสุตฺวา สีสานิ นฺหายิตฺวา ปีตานิ วตฺถานิ ปารุปิตฺวา เยน สุภทฺทา เทวี เตนุปสงฺกมิ.\csromanspacer{} อถ โข อานนฺท สุภทฺทา เทวี ปริณายกรตนํ อามนฺเตสิ “กปฺเปหิ สมฺม ปริณายกรตน จตุรงฺคินิํ เสนํ, จิรํ ทิฏฺโฐ โน ราชา มหาสุทสฺสโน, ราชานํ มหาสุทสฺสนํ ทสฺสนาย อุปสงฺกมิสฺสามา”ติ. “เอวํ เทวี”ติ โข อานนฺท ปริณายกรตนํ สุภทฺทาย เทวิยา ปฏิสฺสุตฺวา จตุรงฺคินิํ เสนํ กปฺปาเปตฺวา สุภทฺทาย เทวิยา ปฏิเวเทสิ “กปฺปิตา โข เทวิ จตุรงฺคินี เสนา, ยสฺสทานิ กาลํ มญฺญสี”ติ.\csromanspacer{} อถ โข อานนฺท}

\vspace{\tipitakaparskip}
\csromancenter{มหาสุทสฺสนสุตฺต สุภทฺทา เทวี จตุรงฺคินิยา เสนาย สทฺธิํ อิตฺถาคาเรน เยน ธมฺโม ปาสาโท เตนุปสงฺกมิ, อุปสงฺกมิตฺวา ธมฺมํ ปาสาทํ อภิรุหิตฺวา เยน มหาวิยูหํ กูฏาคารํ เตนุปสงฺกมิ, อุปสงฺกมิตฺวา มหาวิยูหสฺส กูฏาคารสฺส ทฺวารพาหํ อาลมฺพิตฺวา อฏฺฐาสิ.\csromanspacer{} อถ โข อานนฺท ราชา มหาสุทสฺสโน สทฺทํ สุตฺวา “กิํ นุ โข มหโต วิย ชนกายสฺส สทฺโท”ติ มหาวิยูหา กูฏาคารา นิกฺขมนฺโต อทฺทส สุภทฺทํ เทวิํ ทฺวารพาหํ อาลมฺพิตฺวา ฐิตํ, ทิสฺวาน สุภทฺทํ เทวิํ เอตทโวจ “เอตฺเถว เทวิ ติฏฺฐ มา ปาวิสี”ติ.\csromanspacer{} อถ โข อานนฺท ราชา มหาสุทสฺสโน อญฺญตรํ ปุริสํ อามนฺเตสิ “เอหิ ตฺวํ อมฺโภ ปุริส มหาวิยูหา กูฏาคารา โสวณฺณมยํ ปลฺลงฺกํ นีหริตฺวา สพฺพโสวณฺณมเย ตาลวเน ปญฺญเปหี”ติ. “เอวํ เทวา”ติ โข อานนฺท โส ปุริโส รญฺโญ มหาสุทสฺสนสฺส ปฏิสฺสุตฺวา มหาวิยูหา กูฏาคารา โสวณฺณมยํ ปลฺลงฺกํ นีหริตฺวา สพฺพโสวณฺณมเย ตาลวเน ปญฺญเปสิ.\csromanspacer{} อถ โข อานนฺท ราชา มหาสุทสฺสโน ทกฺขิเณน ปสฺเสน สีหเสยฺยํ กปฺเปสิ ปาเท ปาทํ อจฺจาธาย สโต สมฺปชาโน.}
\proseitem{266}{\csromanfolio{155}อถ โข อานนฺท สุภทฺทาย เทวิยา เอตทโหสิ “วิปฺปสนฺนานิ โข รญฺโญ มหาสุทสฺสนสฺส อินฺทฺริยานิ, ปริสุทฺโธ ฉวิวณฺโณ ปริโยทาโต, มา เหว โข ราชา มหาสุทสฺสโน กาลมกาสี”ติ ราชานํ มหาสุทสฺสนํ เอตทโวจ\csromandash{}}

\prose{“อิมานิ เต เทว จตุราสีติ นครสหสฺสานิ กุสาวตีราชธานิปฺปมุขานิ, เอตฺถ เทว ฉนฺทํ ชเนหิ ชีวิเต อเปกฺขํ กโรหิ.\csromanspacer{} อิมานิ เต เทว จตุราสีติ ปาสาทสหสฺสานิ ธมฺมปาสาทปฺปมุขานิ, เอตฺถ เทว ฉนฺทํ ชเนหิ ชีวิเต อเปกฺขํ กโรหิ.\csromanspacer{} อิมานิ เต เทว จตุราสีติ กูฏาคารสหสฺสานิ มหาวิยูหกูฏาคารปฺปมุขานิ, เอตฺถ เทว ฉนฺทํ ชเนหิ ชีวิเต อเปกฺขํ กโรหิ.\csromanspacer{} อิมานิ เต เทว จตุราสีติ ปลฺลงฺกสหสฺสานิ โสวณฺณมยานิ รูปิยมยานิ ทนฺตมยานิ สารมยานิ โคนกตฺถตานิ ปฏิกตฺถตานิ ปฏลิกตฺถตานิ กทลิมิคปวรปจฺจตฺถรณานิ ส-อุตฺตรจฺฉทานิ อุภโตโลหิตกูปธานานิ, เอตฺถ เทว ฉนฺทํ ชเนหิ ชีวิเต อเปกฺขํ กโรหิ.\csromanspacer{} อิมานิ เต เทว จตุราสีติ นาคสหสฺสานิ โสวณฺณาลงฺการานิ โสวณฺณธชานิ เหมชาลปฏิจฺฉนฺนานิ อุโปสถนาคราชปฺปมุขานิ, เอตฺถ \csromanfolio{156}เทว ฉนฺทํ ชเนหิ ชีวิเต อเปกฺขํ กโรหิ.\csromanspacer{} อิมานิ เต เทว จตุราสีติ อสฺสสหสฺสานิ โสวณฺณาลงฺการานิ โสวณฺณธชานิ เหมชาลปฏิจฺฉนฺนานิ วลาหก-อสฺสราชปฺปมุขานิ, เอตฺถ เทว ฉนฺทํ ชเนหิ ชีวิเต อเปกฺขํ กโรหิ.\csromanspacer{} อิมานิ เต เทว จตุราสีติ รถสหสฺสานิ สีหจมฺมปริวารานิ พฺยคฺฆจมฺมปริวารานิ ทีปิจมฺมปริวารานิ ปณฺฑุกมฺพลปริวารานิ โสวณฺณาลงฺการานิ โสวณฺณธชานิ เหมชาลปฏิจฺฉนฺนานิ เวชยนฺตรถปฺปมุขานิ, เอตฺถ เทว ฉนฺทํ ชเนหิ ชีวิเต อเปกฺขํ กโรหิ.\csromanspacer{} อิมานิ เต เทว จตุราสีติ มณิสหสฺสานิ มณิรตนปฺปมุขานิ, เอตฺถ เทว ฉนฺทํ ชเนหิ ชีวิเต อเปกฺขํ กโรหิ.\csromanspacer{} อิมานิ เต เทว จตุราสีติ อิตฺถิสหสฺสานิ อิตฺถิรตนปฺปมุขานิ, เอตฺถ เทว ฉนฺทํ ชเนหิ ชีวิเต อเปกฺขํ กโรหิ.\csromanspacer{} อิมานิ เต เทว จตุราสีติ คหปติสหสฺสานิ คหปติรตนปฺปมุขานิ, เอตฺถ เทว ฉนฺทํ ชเนหิ ชีวิเต อเปกฺขํ กโรหิ.\csromanspacer{} อิมานิ เต เทว จตุราสีติ ขตฺติยสหสฺสานิ อนุยนฺตานิ ปริณายกรตนปฺปมุขานิ, เอตฺถ เทว ฉนฺทํ ชเนหิ ชีวิเต อเปกฺขํ กโรหิ.\csromanspacer{} อิมานิ เต เทว จตุราสีติ เธนุสหสฺสานิ ทุหสนฺทนานิ กํสูปธารณานิ, เอตฺถ เทว ฉนฺทํ ชเนหิ ชีวิเต อเปกฺขํ กโรหิ.\csromanspacer{} อิมานิ เต เทว จตุราสีติ วตฺถโกฏิสหสฺสานิ โขมสุขุมานํ กปฺปาสิกสุขุมานํ โกเสยฺยสุขุมานํ กมฺพลสุขุมานํ, เอตฺถ เทว ฉนฺทํ ชเนหิ ชีวิเต อเปกฺขํ กโรหิ.\csromanspacer{} อิมานิ เต เทว จตุราสีติ ถาลิปากสหสฺสานิ สายํ ปาตํ ภตฺตาภิหาโร อภิหริยติ, เอตฺถ เทว ฉนฺทํ ชเนหิ ชีวิเต อเปกฺขํ กโรหี”ติ.}

\proseitem{267}{เอวํ วุตฺเต อานนฺท ราชา มหาสุทสฺสโน สุภทฺทํ เทวิํ เอตทโวจ\csromandash{}“ทีฆรตฺตํ โข มํ ตฺวํ เทวิ อิฏฺเฐหิ กนฺเตหิ ปิเยหิ มนาเปหิ สมุทาจริตฺถ, อถ จ ปน มํ ตฺวํ ปจฺฉิเม กาเล อนิฏฺเฐหิ อกนฺเตหิ อปฺปิเยหิ อมนาเปหิ สมุทาจรสี”ติ.\csromanspacer{} กถํ จรหิ ตํ เทว สมุทาจรามีติ.\csromanspacer{} เอวํ โข มํ ตฺวํ เทวิ สมุทาจร “สพฺเพเหว เทว ปิเยหิ มนาเปหิ นานาภาโว วินาภาโว อญฺญถาภาโว, มา โข ตฺวํ เทว สาเปกฺโข กาลมกาสิ, ทุกฺขา สาเปกฺขสฺส กาลํกิริยา, ครหิตา จ สาเปกฺขสฺส กาลํกิริยา.\csromanspacer{} อิมานิ เต เทว จตุราสีติ}

\vspace{\tipitakaparskip}
\csromancenter{มหาสุทสฺสนสุตฺต นครสหสฺสานิ กุสาวตีราชธานิปฺปมุขานิ, เอตฺถ เทว ฉนฺทํ ปชห ชีวิเต อเปกฺขํ มากาสิ.\csromanspacer{} อิมานิ เต เทว จตุราสีติ ปาสาทสหสฺสานิ ธมฺมปาสาทปฺปมุขานิ, เอตฺถ เทว ฉนฺทํ ปชห ชีวิเต อเปกฺขํ มากาสิ.\csromanspacer{} อิมานิ เต เทว จตุราสีติ กูฏาคารสหสฺสานิ มหาวิยูหกูฏาคารปฺปมุขานิ, เอตฺถ เทว ฉนฺทํ ปชห ชีวิเต อเปกฺขํ มากาสิ.\csromanspacer{} อิมานิ เต เทว จตุราสีติ ปลฺลงฺกสหสฺสานิ โสวณฺณมยานิ รูปิยมยานิ ทนฺตมยานิ สารมยานิ โคนกตฺถตานิ ปฏิกตฺถตานิ ปฏลิกตฺถตานิ กทลิมิคปวรปจฺจตฺถรณานิ ส- อุตฺตรจฺฉทานิ อุภโตโลหิตกูปธานานิ, เอตฺถ เทว ฉนฺทํ ปชห ชีวิเต อเปกฺขํ มากาสิ.\csromanspacer{} อิมานิ เต เทว จตุราสีติ นาคสหสฺสานิ โสวณฺณาลงฺการานิ โสวณฺณธชานิ เหมชาลปฏิจฺฉนฺนานิ อุโปสถนาคราชปฺปมุขานิ, เอตฺถ เทว ฉนฺทํ ปชห ชีวิเต อเปกฺขํ มากาสิ.\csromanspacer{} อิมานิ เต เทว จตุราสีติ อสฺสสหสฺสานิ โสวณฺณาลงฺการานิ โสวณฺณธชานิ เหมชาลปฏิจฺฉนฺนานิ วลาหก-อสฺสราชปฺปมุขานิ, เอตฺถ เทว ฉนฺทํ ปชห ชีวิเต อเปกฺขํ มากาสิ.\csromanspacer{} อิมานิ เต เทว จตุราสีติ รถสหสฺสานิ สีหจมฺมปริวารานิ พฺยคฺฆจมฺมปริวารานิ ทีปิจมฺมปริวารานิ ปณฺฑุกมฺพลปริวารานิ โสวณฺณาลงฺการานิ โสวณฺณธชานิ เหมชาลปฏิจฺฉนฺนานิ เวชยนฺตรถปฺปมุขานิ, เอตฺถ เทว ฉนฺทํ ปชห ชีวิเต อเปกฺขํ มากาสิ.\csromanspacer{} อิมานิ เต เทว จตุราสีติ มณิสหสฺสานิ มณิรตนปฺปมุขานิ, เอตฺถ เทว ฉนฺทํ ปชห ชีวิเต อเปกฺขํ มากาสิ.\csromanspacer{} อิมานิ เต เทว จตุราสีติ อิตฺถิสหสฺสานิ สุภทฺทาเทวิปฺปมุขานิ, เอตฺถ เทว ฉนฺทํ ปชห ชีวิเต อเปกฺขํ มากาสิ.\csromanspacer{} อิมานิ เต เทว จตุราสีติ คหปติสหสฺสานิ คหปติรตนปฺปมุขานิ, เอตฺถ เทว ฉนฺทํ ปชห ชีวิเต อเปกฺขํ มากาสิ.\csromanspacer{} อิมานิ เต เทว จตุราสีติ ขตฺติยสหสฺสานิ อนุยนฺตานิ ปริณายกรตนปฺปมุขานิ, เอตฺถ เทว ฉนฺทํ ปชห ชีวิเต อเปกฺขํ มากาสิ.\csromanspacer{} อิมานิ เต เทว จตุราสีติ เธนุสหสฺสานิ ทุหสนฺทนานิ กํสูปธารณานิ, เอตฺถ เทว ฉนฺทํ ปชห ชีวิเต อเปกฺขํ มากาสิ.\csromanspacer{} อิมานิ เต เทว จตุราสีติ วตฺถโกฏิสหสฺสานิ โขมสุขุมานํ กปฺปาสิกสุขุมานํ โกเสยฺยสุขุมานํ กมฺพลสุขุมานํ, เอตฺถ เทว ฉนฺทํ ปชห ชีวิเต อเปกฺขํ มากาสิ.\csromanspacer{} อิมานิ เต เทว จตุราสีติ ถาลิปากสหสฺสานิ สายํ ปาตํ ภตฺตาภิหาโร อภิหริยติ, เอตฺถ เทว ฉนฺทํ ปชห ชีวิเต อเปกฺขํ มากาสี”ติ.}
\proseitem{268}{\csromanfolio{158}เอวํ วุตฺเต อานนฺท สุภทฺทา เทวี ปโรทิ อสฺสูนิ ปวตฺเตสิ.\csromanspacer{} อถ โข อานนฺท สุภทฺทา เทวี อสฺสูนิ ปุญฺฉิตฺวา\footnote{ปมชฺชิตฺวา (สี, สฺยา, อิ), ปุญฺชิตฺวา (ก)} ราชานํ มหาสุทสฺสนํ เอตทโวจ\csromandash{}}

\prose{“สพฺเพเหว เทว ปิเยหิ มนาเปหิ นานาภาโว วินาภาโว อญฺญถาภาโว, มา โข ตฺวํ เทว สาเปกฺโข กาลมกาสิ, ทุกฺขา สาเปกฺขสฺส กาลํกิริยา, ครหิตา จ สาเปกฺขสฺส กาลํกิริยา.\csromanspacer{} อิมานิ เต เทว จตุราสีติ นครสหสฺสานิ กุสาวตีราชธานิปฺปมุขานิ, เอตฺถ เทว ฉนฺทํ ปชห ชีวิเต อเปกฺขํ มากาสิ.\csromanspacer{} อิมานิ เต เทว จตุราสีติ ปาสาทสหสฺสานิ ธมฺมปาสาทปฺปมุขานิ, เอตฺถ เทว ฉนฺทํ ปชห ชีวิเต อเปกฺขํ มากาสิ.\csromanspacer{} อิมานิ เต เทว จตุราสีติ กูฏาคารสหสฺสานิ มหาวิยูหกูฏาคารปฺปมุขานิ, เอตฺถ เทว ฉนฺทํ ปชห ชีวิเต อเปกฺขํ มากาสิ.\csromanspacer{} อิมานิ เต เทว จตุราสีติ ปลฺลงฺกสหสฺสานิ โสวณฺณมยานิ รูปิยมยานิ ทนฺตมยานิ สารมยานิ โคนกตฺถตานิ ปฏิกตฺถตานิ ปฏลิกตฺถตานิ กทลิมิคปวรปจฺจตฺถรณานิ ส- อุตฺตรจฺฉทานิ อุภโตโลหิตกูปธานานิ, เอตฺถ เทว ฉนฺทํ ปชห ชีวิเต อเปกฺขํ มากาสิ.\csromanspacer{} อิมานิ เต เทว จตุราสีติ นาคสหสฺสานิ โสวณฺณาลงฺการานิ โสวณฺณธชานิ เหมชาลปฏิจฺฉนฺนานิ อุโปสถนาคราชปฺปมุขานิ, เอตฺถ เทว ฉนฺทํ ปชห ชีวิเต อเปกฺขํ มากาสิ.\csromanspacer{} อิมานิ เต เทว จตุราสีติ อสฺสสหสฺสานิ โสวณฺณาลงฺการานิ โสวณฺณธชานิ เหมชาลปฏิจฺฉนฺนานิ วลาหก-อสฺสราชปฺปมุขานิ, เอตฺถ เทว ฉนฺทํ ปชห ชีวิเต อเปกฺขํ มากาสิ.\csromanspacer{} อิมานิ เต เทว จตุราสีติ รถสหสฺสานิ สีหจมฺมปริวารานิ พฺยคฺฆจมฺมปริวารานิ ทีปิจมฺมปริวารานิ ปณฺฑุกมฺพลปริวารานิ โสวณฺณาลงฺการานิ โสวณฺณธชานิ เหมชาลปฏิจฺฉนฺนานิ เวชยนฺตรถปฺปมุขานิ, เอตฺถ เทว ฉนฺทํ ปชห ชีวิเต อเปกฺขํ มากาสิ.\csromanspacer{} อิมานิ เต เทว จตุราสีติ มณิสหสฺสานิ มณิรตนปฺปมุขานิ, เอตฺถ เทว ฉนฺทํ ปชห ชีวิเต อเปกฺขํ มากาสิ.\csromanspacer{} อิมานิ เต เทว จตุราสีติ อิตฺถิสหสฺสานิ อิตฺถิรตนปฺปมุขานิ, เอตฺถ เทว ฉนฺทํ ปชห ชีวิเต อเปกฺขํ มากาสิ.\csromanspacer{} อิมานิ เต เทว จตุราสีติ คหปติสหสฺสานิ คหปติรตนปฺปมุขานิ, เอตฺถ เทว ฉนฺทํ ปชห ชีวิเต อเปกฺขํ มากาสิ.\csromanspacer{} อิมานิ}

\vspace{\tipitakaparskip}
\csromancenter{มหาสุทสฺสนสุตฺต เต เทว จตุราสีติ ขตฺติยสหสฺสานิ อนุยนฺตานิ ปริณายกรตนปฺปมุขานิ, เอตฺถ เทว ฉนฺทํ ปชห ชีวิเต อเปกฺขํ มากาสิ.\csromanspacer{} อิมานิ เต เทว จตุราสีติ เธนุสหสฺสานิ ทุหสนฺทนานิ กํสูปธารณานิ, เอตฺถ เทว ฉนฺทํ ปชห ชีวิเต อเปกฺขํ มากาสิ.\csromanspacer{} อิมานิ เต เทว จตุราสีติ วตฺถโกฏิสหสฺสานิ โขมสุขุมานํ กปฺปาสิกสุขุมานํ โกเสยฺยสุขุมานํ กมฺพลสุขุมานํ, เอตฺถ เทว ฉนฺทํ ปชห ชีวิเต อเปกฺขํ มากาสิ.\csromanspacer{} อิมานิ เต เทว จตุราสีติ ถาลิปากสหสฺสานิ สายํ ปาตํ ภตฺตาภิหาโร อภิหริยติ, เอตฺถ เทว ฉนฺทํ ปชห ชีวิเต อเปกฺขํ มากาสี”ติ.}
\csromanmatikaanchor{mtk.81}
\csromantocmarkref{section}{พฺรหฺมโลกูปคม}{mtk.81}
\bookmark[dest={mtk.81},level=1]{พฺรหฺมโลกูปคม}
\csromanheader{1}{\textbf{พฺรหฺมโลกูปคม}}
\proseitem{269}{\csromanfolio{159}อถ โข อานนฺท ราชา มหาสุทสฺสโน น จิรสฺเสว กาลมกาสิ.\csromanspacer{} เสยฺยถาปิ อานนฺท คหปติสฺส วา คหปติปุตฺตสฺส วา มนุญฺญํ โภชนํ ภุตฺตาวิสฺส ภตฺตสมฺมโท โหติ.\csromanspacer{} เอวเมว โข อานนฺท รญฺโญ มหาสุทสฺสนสฺส มารณนฺติกา เวทนา อโหสิ.\csromanspacer{} กาลงฺกโต จ อานนฺท ราชา มหาสุทสฺสโน สุคติํ พฺรหฺมโลกํ อุปปชฺชิ.\csromanspacer{} ราชา อานนฺท มหาสุทสฺสโน จตุราสีติ วสฺสสหสฺสานิ กุมารกีฬํ\footnote{กีฬิตํ (ก), กีฬิกํ (สี, อิ)} กีฬิ.\csromanspacer{} จตุราสีติ วสฺสสหสฺสานิ โอปรชฺชํ กาเรสิ.\csromanspacer{} จตุราสีติ วสฺสสหสฺสานิ รชฺชํ กาเรสิ.\csromanspacer{} จตุราสีติ วสฺสสหสฺสานิ คิหิภูโต\footnote{คิหีภูโต (สี, อิ)} ธมฺเม ปาสาเท พฺรหฺมจริยํ จริ\footnote{พฺรหฺมจริยมจริ (ก)}.\csromanspacer{} โส จตฺตาโร พฺรหฺมวิหาเร ภาเวตฺวา กายสฺส เภทา ปรํ มรณา พฺรหฺมโลกูปโค อโหสิ.}

\proseitem{270}{สิยา โข ปนานนฺท เอวมสฺส “อญฺโญ นูน เตน สมเยน ราชา มหาสุทสฺสโน อโหสี”ติ.\csromanspacer{} น โข ปเนตํ อานนฺท เอวํ ทฏฺฐพฺพํ.\csromanspacer{} อหํ เตน สมเยน ราชา มหาสุทสฺสโน อโหสิํ.\csromanspacer{} มม ตานิ จตุราสีติ นครสหสฺสานิ กุสาวตีราชธานิปฺปมุขานิ, มม ตานิ จตุราสีติ ปาสาทสหสฺสานิ ธมฺมปาสาทปฺปมุขานิ, มม ตานิ จตุราสีติ กูฏาคารสหสฺสานิ มหาวิยูหกูฏาคารปฺปมุขานิ, มม ตานิ จตุราสีติ ปลฺลงฺกสหสฺสานิ โสวณฺณมยานิ รูปิยมยานิ ทนฺตมยานิ สารมยานิ โคนกตฺถตานิ ปฏิกตฺถตานิ ปฏิลิกตฺถตานิ กทลิมิคปวรปจฺจตฺถรณานิ \csromanfolio{160}ส-อุตฺตรจฺฉทานิ อุภโตโลหิตกูปธานานิ, มม ตานิ จตุราสีติ นาคสหสฺสานิ โสวณฺณาลงฺการานิ โสวณฺณธชานิ เหมชาลปฏิจฺฉนฺนานิ อุโปสถนาคราชปฺปมุขานิ, มม ตานิ จตุราสีติ อสฺสสหสฺสานิ โสวณฺณาลงฺการานิ โสวณฺณธชานิ เหมชาลปฏิจฺฉนฺนานิ วลาหก- อสฺสราชปฺปมุขานิ, มม ตานิ จตุราสีติ รถสหสฺสานิ สีหจมฺมปริวารานิ พฺยคฺฆจมฺมปริวารานิ ทีปิจมฺมปริวารานิ ปณฺฑุกมฺพลปริวารานิ โสวณฺณาลงฺการานิ โสวณฺณธชานิ เหมชาลปฏิจฺฉนฺนานิ เวชยนฺตรถปฺปมุขานิ, มม ตานิ จตุราสีติ มณิสหสฺสานิ มณิรตนปฺปมุขานิ, มม ตานิ จตุราสีติ อิตฺถิสหสฺสานิ สุภทฺทาเทวิปฺปมุขานิ, มม ตานิ จตุราสีติ คหปติสหสฺสานิ คหปติรตนปฺปมุขานิ, มม ตานิ จตุราสีติ ขตฺติยสหสฺสานิ อนุยนฺตานิ ปริณายกรตนปฺปมุขานิ, มม ตานิ จตุราสีติ เธนุสหสฺสานิ ทุหสนฺทนานิ กํสูปธารณานิ, มม ตานิ จตุราสีติ วตฺถโกฏิสหสฺสานิ โขมสุขุมานํ กปฺปาสิกสุขุมานํ โกเสยฺยสุขุมานํ กมฺพลสุขุมานํ, มม ตานิ จตุราสีติ ถาลิปากสหสฺสานิ สายํ ปาตํ ภตฺตาภิหาโร อภิหริยิตฺถ.}

\proseitem{271}{เตสํ โข ปนานนฺท จตุราสีติ นครสหสฺสานํ เอกญฺเญว ตํ นครํ โหติ, ยํ เตน สมเยน อชฺฌาวสามิ ยทิทํ กุสาวตี ราชธานี.\csromanspacer{} เตสํ โข ปนานนฺท จตุราสีติ ปาสาทสหสฺสานํ เอโกเยว โส ปาสาโท โหติ, ยํ เตน สมเยน อชฺฌาวสามิ ยทิทํ ธมฺโม ปาสาโท.\csromanspacer{} เตสํ โข ปนานนฺท จตุราสีติ กูฏาคารสหสฺสานํ เอกญฺเญว ตํ กูฏาคารํ โหติ, ยํ เตน สมเยน อชฺฌาวสามิ ยทิทํ มหาวิยูหํ กูฏาคารํ.\csromanspacer{} เตสํ โข ปนานนฺท จตุราสีติ ปลฺลงฺกสหสฺสานํ เอโกเยว โส ปลฺลงฺโก โหติ, ยํ เตน สมเยน ปริภุญฺชามิ ยทิทํ โสวณฺณมโย วา รูปิยมโย วา ทนฺตมโย วา สารมโย วา.\csromanspacer{} เตสํ โข ปนานนฺท จตุราสีติ นาคสหสฺสานํ เอโกเยว โส นาโค โหติ, ยํ เตน สมเยน อภิรุหามิ ยทิทํ อุโปสโถ นาคราชา.\csromanspacer{} เตสํ โข ปนานนฺท จตุราสีติ อสฺสสหสฺสานํ เอโกเยว โส อสฺโส โหติ, ยํ เตน สมเยน อภิรุหามิ ยทิทํ วลาหโก อสฺสราชา.\csromanspacer{} เตสํ โข ปนานนฺท จตุราสีติ รถสหสฺสานํ เอโกเยว โส รโถ โหติ, ยํ เตน สมเยน อภิรุหามิ ยทิทํ เวชยนฺตรโถ.\csromanspacer{} เตสํ โข ปนานนฺท จตุราสีติ อิตฺถิสหสฺสานํ}

\vspace{\tipitakaparskip}
\csromancenter{มหาสุทสฺสนสุตฺต เอกาเยว สา อิตฺถี โหติ, ยา เตน สมเยน ปจฺจุปฏฺฐาติ ขตฺติยานี วา เวสฺสินี\footnote{เวสฺสายินี (สฺยา) เวลามิกานี (ก-สี, อิ) เวลามิกา (สํ 2. 119) สํยุตฺตฏฺฐกถา โอโลเกตพฺพา.} วา.\csromanspacer{} เตสํ โข ปนานนฺท จตุราสีติ วตฺถโกฏิสหสฺสานํ เอกํเยว ตํ ทุสฺสยุคํ โหติ, ยํ เตน สมเยน ปริทหามิ โขมสุขุมํ วา กปฺปาสิกสุขุมํ วา โกเสยฺยสุขุมํ วา กมฺพลสุขุมํ วา.\csromanspacer{} เตสํ โข ปนานนฺท จตุราสีติ ถาลิปากสหสฺสานํ เอโกเยว โส ถาลิปาโก โหติ, ยโต นาฬิโกทนปรมํ ภุญฺชามิ ตทุปิยญฺจ สูเปยฺยํ.}
\proseitem{272}{\csromanfolio{161}ปสฺสานนฺท สพฺเพเต สงฺขารา อตีตา นิรุทฺธา วิปริณตา.\csromanspacer{} เอวํ อนิจฺจา โข อานนฺท สงฺขารา, เอวํ อทฺธุวา โข อานนฺท สงฺขารา, เอวํ อนสฺสาสิกา โข อานนฺท สงฺขารา.\csromanspacer{} ยาวญฺจิทํ อานนฺท อลเมว สพฺพสงฺขาเรสุ นิพฺพินฺทิตุํ, อลํ วิรชฺชิตุํ, อลํ วิมุจฺจิตุํ.}

\prose{ฉกฺขตฺตุํ โข ปนาหํ อานนฺท อภิชานามิ อิมสฺมิํ ปเทเส สรีรํ นิกฺขิปิตํ, ตญฺจ โข ราชาว สมาโน จกฺกวตฺตี ธมฺมิโก ธมฺมราชา จาตุรนฺโต วิชิตาวี ชนปทตฺถาวริยปตฺโต สตฺตรตนสมนฺนาคโต, อยํ สตฺตโม สรีรนิกฺเขโป.\csromanspacer{} น โข ปนาหํ อานนฺท ตํ ปเทสํ สมนุปสฺสามิ สเทวเก โลเก สมารเก สพฺรหฺมเก สสฺสมณพฺราหฺมณิยา ปชาย สเทวมนุสฺสาย, ยตฺถ ตถาคโต อฏฺฐมํ สรีรํ นิกฺขิเปยฺยาติ.\csromanspacer{} อิทมโวจ ภควา, อิทํ วตฺวาน สุคโต อถาปรํ เอตทโวจ สตฺถา\csromandash{}}

\csromangathameasure{อนิจฺจา วต สงฺขารา, อุปฺปาทวยธมฺมิโน. \\ อุปฺปชฺชิตฺวา นิรุชฺฌนฺติ, เตสํ วูปสโม สุโขติ.}
\csromangathagroup{\csromangathastanza{อนิจฺจา วต สงฺขารา, อุปฺปาทวยธมฺมิโน. \\ อุปฺปชฺชิตฺวา นิรุชฺฌนฺติ, เตสํ วูปสโม สุโขติ.}}

\nitthitam{\textbf{มหาสุทสฺสนสุตฺตํ นิฏฺฐิตํ จตุตฺถํ.}}
\csromanmatikaanchor{mtk.82}
\csromantocmarknopage{chapter}{5. ชนวสภสุตฺต}{mtk.82}
\bookmark[dest={mtk.82},level=0]{5. ชนวสภสุตฺต}
\chapterheadpageread{5. ชนวสภสุตฺต}
\csromanmatikaanchor{mtk.83}
\csromantocmarkref{section}{นาติกิยาทิพฺยากรณ}{mtk.83}
\bookmark[dest={mtk.83},level=1]{นาติกิยาทิพฺยากรณ}
\csromanheader{1}{\textbf{นาติกิยาทิพฺยากรณ}}
\proseitem{273}{\csromanfolio{162}เอวํ เม สุตํ\csromandash{}เอกํ สมยํ ภควา นาติเก\footnote{นาทิเก (สี, สฺยา, อิ)} วิหรติ คิญฺชกาวสเถ.\csromanspacer{} เตน โข ปน สมเยน ภควา ปริโต ปริโต ชนปเทสุ ปริจารเก อพฺภตีเต กาลงฺกเต อุปปตฺตีสุ พฺยากโรติ กาสิโกสเลสุ วชฺชิมลฺเลสุ เจติวํเสสุ\footnote{เจติยวเสสุ (ก)} กุรุปญฺจาเลสุ มชฺฌสูรเสเนสุ\footnote{มจฺฉสุรเสเนสุ (สฺยา), มจฺฉสูรเสเนสุ (สี, อิ)} “อสุ อมุตฺร อุปปนฺโน อสุ อมุตฺร อุปปนฺโน\footnote{อุปปนฺโนติ (ก)}, ปโรปญฺญาส นาติกิยา ปริจารกา อพฺภตีตา กาลงฺกตา ปญฺจนฺนํ โอรมฺภาคิยานํ สํโยชนานํ ปริกฺขยา โอปปาติกา ตตฺถ ปรินิพฺพายิโน อนาวตฺติธมฺมา ตสฺมา โลกา.\csromanspacer{} สาธิกา นวุติ นาติกิยา ปริจารกา อพฺภตีตา กาลงฺกตา ติณฺณํ สํโยชนานํ ปริกฺขยา ราคโทสโมหานํ ตนุตฺตา สกทาคามิโน สกิเทว\footnote{สกิํเทว (ก)} อิมํ โลกํ อาคนฺตฺวา ทุกฺขสฺสนฺตํ กริสฺสนฺติ.\csromanspacer{} สาติเรกานิ ปญฺจสตานิ นาติกิยา ปริจารกา อพฺภตีตา กาลงฺกตา ติณฺณํ สํโยชนานํ ปริกฺขยา โสตาปนฺนา อวินิปาตธมฺมา นิยตา สมฺโพธิปรายณา”ติ.}

\proseitem{274}{อสฺโสสุํ โข นาติกิยา ปริจารกา “ภควา กิร ปริโต ปริโต ชนปเทสุ ปริจารเก อพฺภตีเต กาลงฺกเต อุปปตฺตีสุ พฺยากโรติ กาสิโกสเลสุ วชฺชิมลฺเลสุ เจติวํเสสุ กุรุปญฺจาเลสุ มชฺฌสูรเสเนสุ ‘อสุ อมุตฺร อุปปนฺโน อสุ อมุตฺร อุปปนฺโน, ปโรปญฺญาส นาติกิยา ปริจารกา อพฺภตีตา กาลงฺกตา ปญฺจนฺนํ โอรมฺภาคิยานํ สํโยชนานํ ปริกฺขยา โอปปาติกา ตตฺถ ปรินิพฺพายิโน อนาวตฺติธมฺมา ตสฺมา โลกา.\csromanspacer{} สาธิกา นวุติ นาติกิยา ปริจารกา อพฺภตีตา กาลงฺกตา ติณฺณํ สํโยชนานํ ปริกฺขยา ราคโทสโมหานํ ตนุตฺตา สกทาคามิโน สกิเทว อิมํ โลกํ อาคนฺตฺวา ทุกฺขสฺสนฺตํ กริสฺสนฺติ.\csromanspacer{} สาติเรกานิ ปญฺจสตานิ นาติกิยา ปริจารกา อพฺภตีตา กาลงฺกตา ติณฺณํ สํโยชนานํ}

\vspace{\tipitakaparskip}
\csromancenter{ชนวสภสุตฺต ปริกฺขยา โสตาปนฺนา อวินิปาตธมฺมา นิยตา สมฺโพธิปรายณา’ติ”.\csromanspacer{} เตน จ นาติกิยา ปริจารกา อตฺตมนา อเหสุํ ปมุทิตา ปีติโสมนสฺสชาตา ภควโต ปญฺหเวยฺยากรณํ\footnote{ปญฺหาเวยฺยากรณํ (สฺยา, ก)} สุตฺวา.}
\proseitem{275}{\csromanfolio{163}อสฺโสสิ โข อายสฺมา อานนฺโท “ภควา กิร ปริโต ปริโต ชนปเทสุ ปริจารเก อพฺภตีเต กาลงฺกเต อุปปตฺตีสุ พฺยากโรติ กาสิโกสเลสุ วชฺชิมลฺเลสุ เจติวํเสสุ กุรุปญฺจาเลสุ มชฺฌสูรเสเนสุ ‘อสุ อมุตฺร อุปปนฺโน อสุ อมุตฺร อุปปนฺโน, ปโรปญฺญาส นาติกิยา ปริจารกา อพฺภตีตา กาลงฺกตา ปญฺจนฺนํ โอรมฺภาคิยานํ สํโยชนานํ ปริกฺขยา โอปปาติกา ตตฺถ ปรินิพฺพายิโน อนาวตฺติธมฺมา ตสฺมา โลกา.\csromanspacer{} สาธิกา นวุติ นาติกิยา ปริจารกา อพฺภตีตา กาลงฺกตา ติณฺณํ สํโยชนานํ ปริกฺขยา ราคโทสโมหานํ ตนุตฺตา สกทาคามิโน สกิเทว อิมํ โลกํ อาคนฺตฺวา ทุกฺขสฺสนฺตํ กริสฺสนฺติ.\csromanspacer{} สาติเรกานิ ปญฺจสตานิ นาติกิยา ปริจารกา อพฺภตีตา กาลงฺกตา ติณฺณํ สํโยชนานํ ปริกฺขยา โสตาปนฺนา อวินิปาตธมฺมา นิยตา สมฺโพธิปรายณา’ติ.\csromanspacer{} เตน จ นาติกิยา ปริจารกา อตฺตมนา อเหสุํ ปมุทิตา ปีติโสมนสฺสชาตา ภควโต ปญฺหเวยฺยากรณํ สุตฺวา”ติ.}

\csromanmatikaanchor{mtk.84}
\csromantocmarkref{section}{อานนฺทปริกถา}{mtk.84}
\bookmark[dest={mtk.84},level=1]{อานนฺทปริกถา}
\csromanheader{1}{\textbf{อานนฺทปริกถา}}
\proseitem{276}{อถ โข อายสฺมโต อานนฺทสฺส เอตทโหสิ “อิเม โข ปนาปิ อเหสุํ มาคธกา ปริจารกา พหู เจว รตฺตญฺญู จ อพฺภตีตา กาลงฺกตา.\csromanspacer{} สุญฺญา มญฺเญ องฺคมคธา องฺคมาคธเกหิ\footnote{องฺคมาคธิเกหิ (สฺยา)} ปริจารเกหิ อพฺภตีเตหิ กาลงฺกเตหิ.\csromanspacer{} เต โข ปนาปิ\footnote{เตน โข ปนาปิ (สฺยา)} อเหสุํ พุทฺเธ ปสนฺนา ธมฺเม ปสนฺนา สํเฆ ปสนฺนา สีเลสุ ปริปูรการิโน, เต อพฺภตีตา กาลงฺกตา ภควตา อพฺยากตา.\csromanspacer{} เตสมฺปิสฺส สาธุ เวยฺยากรณํ, พหุชโน ปสีเทยฺย, ตโต คจฺเฉยฺย สุคติํ.\csromanspacer{} อยํ โข ปนาปิ อโหสิ ราชา มาคโธ เสนิโย พิมฺพิสาโร ธมฺมิโก ธมฺมราชา หิโต พฺราหฺมณคหปติกานํ เนคมานญฺเจว ชานปทานญฺจ.\csromanspacer{} อปิสฺสุทํ มนุสฺสา กิตฺตยมานรูปา วิหรนฺติ ‘เอวํ โน โส ธมฺมิโก ธมฺมราชา \csromanfolio{164}สุขาเปตฺวา กาลงฺกโต, เอวํ มยํ ตสฺส ธมฺมิกสฺส ธมฺมรญฺโญ วิชิเต ผาสุ\footnote{ผาสุกํ (สฺยา) 2-3. นินฺนมนา (สฺยา), ทีนมานา (สี, อิ)} วิหริมฺหา’ติ.\csromanspacer{} โส โข ปนาปิ อโหสิ พุทฺเธ ปสนฺโน ธมฺเม ปสนฺโน สํเฆ ปสนฺโน สีเลสุ ปริปูรการี.\csromanspacer{} อปิสฺสุทํ มนุสฺสา เอวมาหํสุ ‘ยาว มรณกาลาปิ ราชา มาคโธ เสนิโย พิมฺพิสาโร ภควนฺตํ กิตฺตยมานรูโป กาลงฺกโต’ติ.\csromanspacer{} โส อพฺภตีโต กาลงฺกโต ภควตา อพฺยากโต.\csromanspacer{} ตสฺสปิสฺส สาธุ เวยฺยากรณํ พหุชโน ปสีเทยฺย, ตโต คจฺเฉยฺย สุคติํ.\csromanspacer{} ภควโต โข ปน สมฺโพธิ มคเธสุ.\csromanspacer{} ยตฺถ โข ปน ภควโต สมฺโพธิ มคเธสุ, กถํ ตตฺร ภควา มาคธเก ปริจารเก อพฺภตีเต กาลงฺกเต อุปปตฺตีสุ น พฺยากเรยฺย.\csromanspacer{} ภควา เจ โข ปน มาคธเก ปริจารเก อพฺภตีเต กาลงฺกเต อุปปตฺตีสุ น พฺยากเรยฺย.\csromanspacer{} ทีนมนา2 เตนสฺสุ มาคธกา ปริจารกา, เยน โข ปนสฺสุ ทีนมนา3 มาคธกา ปริจารกา.\csromanspacer{} กถํ เต ภควา น พฺยากเรยฺยา”ติ.}

\proseitem{277}{อิทมายสฺมา อานนฺโท มาคธเก ปริจารเก อารพฺภ เอโก รโห อนุวิจินฺเตตฺวา รตฺติยา ปจฺจูสสมยํ ปจฺจุฏฺฐาย เยน ภควา เตนุปสงฺกมิ, อุปสงฺกมิตฺวา ภควนฺตํ อภิวาเทตฺวา เอกมนฺตํ นิสีทิ, เอกมนฺตํ นิสินฺโน โข อายสฺมา อานนฺโท ภควนฺตํ เอตทโวจ “สุตํ เมตํ ภนฺเต ‘ภควา กิร ปริโต ปริโต ชนปเทสุ ปริจารเก อพฺภตีเต กาลงฺกเต อุปปตฺตีสุ พฺยากโรติ กาสิโกสเลสุ วชฺชิมลฺเลสุ เจติวํเสสุ กุรุปญฺจาเลสุ มชฺฌสูรเสเนสุ ‘อสุ อมุตฺร อุปปนฺโน อสุ อมุตฺร อุปปนฺโน, ปโรปญฺญาส นาติกิยา ปริจารกา อพฺภตีตา กาลงฺกตา ปญฺจนฺนํ โอรมฺภาคิยานํ สํโยชนานํ ปริกฺขยา โอปปาติกา ตตฺถ ปรินิพฺพายิโน อนาวตฺติธมฺมา ตสฺมา โลกา.\csromanspacer{} สาธิกา นวุติ นาติกิยา ปริจารกา อพฺภตีตา กาลงฺกตา ติณฺณํ สํโยชนานํ ปริกฺขยา ราคโทสโมหานํ ตนุตฺตา สกทาคามิโน สกิเทว อิมํ โลกํ อาคนฺตฺวา ทุกฺขสฺสนฺตํ กริสฺสนฺติ.\csromanspacer{} สาติเรกานิ ปญฺจสตานิ นาติกิยา ปริจารกา อพฺภตีตา กาลงฺกตา ติณฺณํ สํโยชนานํ ปริกฺขยา โสตาปนฺนา อวินิปาตธมฺมา นิยตา สมฺโพธิปรายณา’ติ.\csromanspacer{} เตน จ นาติกิยา ปริจารกา อตฺตมนา อเหสุํ ปมุทิตา ปีติโสมนสฺสชาตา ภควโต ปญฺหเวยฺยากรณํ}

\vspace{\tipitakaparskip}
\csromancenter{ชนวสภสุตฺต สุตฺวา’ติ.\csromanspacer{} อิเม โข ปนาปิ ภนฺเต อเหสุํ มาคธกา ปริจารกา พหู เจว รตฺตญฺญู จ อพฺภตีตา กาลงฺกตา.\csromanspacer{} สุญฺญา มญฺเญ องฺคมคธา องฺคมาคธเกหิ ปริจารเกหิ อพฺภตีเตหิ กาลงฺกเตหิ.\csromanspacer{} เต โข ปนาปิ ภนฺเต อเหสุํ พุทฺเธ ปสนฺนา ธมฺเม ปสนฺนา สํเฆ ปสนฺนา สีเลสุ ปริปูรการิโน, เต อพฺภตีตา กาลงฺกตา ภควตา อพฺยากตา.\csromanspacer{} เตสมฺปิสฺส สาธุ เวยฺยากรณํ, พหุชโน ปสีเทยฺย, ตโต คจฺเฉยฺย สุคติํ.\csromanspacer{} อยํ โข ปนาปิ ภนฺเต อโหสิ ราชา มาคโธ เสนิโย พิมฺพิสาโร ธมฺมิโก ธมฺมราชา หิโต พฺราหฺมณคหปติกานํ เนคมานญฺเจว ชนปทานญฺจ.\csromanspacer{} อปิสฺสุทํ มนุสฺสา กิตฺตยมานรูปา วิหรนฺติ ‘เอวํ โน โส ธมฺมิโก ธมฺมราชา สุขาเปตฺวา กาลงฺกโต.\csromanspacer{} เอวํ มยํ ตสฺส ธมฺมิกสฺส ธมฺมรญฺโญ วิชิเต ผาสุ วิหริมฺหา’ติ.\csromanspacer{} โส โข ปนาปิ ภนฺเต อโหสิ พุทฺเธ ปสนฺโน ธมฺเม ปสนฺโน สํเฆ ปสนฺโน สีเลสุ ปริปูรการี.\csromanspacer{} อปิสฺสุทํ มนุสฺสา เอวมาหํสุ ‘ยาว มรณกาลาปิ ราชา มาคโธ เสนิโย พิมฺพิสาโร ภควนฺตํ กิตฺตยมานรูโป กาลงฺกโต’ติ.\csromanspacer{} โส อพฺภตีโต กาลงฺกโต ภควตา อพฺยากโต, ตสฺสปิสฺส สาธุ เวยฺยากรณํ, พหุชโน ปสีเทยฺย, ตโต คจฺเฉยฺย สุคติํ.\csromanspacer{} ภควโต โข ปน ภนฺเต สมฺโพธิ มคเธสุ.\csromanspacer{} ยตฺถ โข ปน ภนฺเต ภควโต สมฺโพธิ มคเธสุ, กถํ ตตฺร ภควา มาคธเก ปริจารเก อพฺภตีเต กาลงฺกเต อุปปตฺตีสุ น พฺยากเรยฺย.\csromanspacer{} ภควา เจ โข ปน ภนฺเต มาคธเก ปริจารเก อพฺภตีเต กาลงฺกเต อุปปตฺตีสุ น พฺยากเรยฺย.\csromanspacer{} ทินมนา เตนสฺสุ มาคธกา ปริจารกา, เยน โข ปนสฺสุ ทีนมนา มาคธกา ปริจารกา.\csromanspacer{} กถํ เต ภควา น พฺยากเรยฺยา”ติ.\csromanspacer{} อิทมายสฺมา อานนฺโท มาคธเก ปริจารเก อารพฺภ ภควโต สมฺมุขา ปริกถํ กตฺวา อุฏฺฐายาสนา ภควนฺตํ อภิวาเทตฺวา ปทกฺขิณํ กตฺวา ปกฺกามิ.}
\proseitem{278}{\csromanfolio{165}อถ โข ภควา อจิรปกฺกนฺเต อายสฺมนฺเต อานนฺเท ปุพฺพณฺหสมยํ นิวาเสตฺวา ปตฺตจีวรมาทาย นาติกํ ปิณฺฑาย ปาวิสิ, นาติเก ปิณฺฑาย จริตฺวา ปจฺฉาภตฺตํ ปิณฺฑปาตปฏิกฺกนฺโต ปาเท ปกฺขาเลตฺวา คิญฺชกาวสถํ ปวิสิตฺวา มาคธเก ปริจารเก อารพฺภ อฏฺฐิํกตฺวา\footnote{อฏฺฐิกตฺวา (สี, สฺยา, อิ)} มนสิกตฺวา สพฺพํ เจตสา\footnote{สพฺพเจตสา (อิ)} สมนฺนาหริตฺวา ปญฺญตฺเต \csromanfolio{166}อาสเน นิสีทิ “คติํ เนสํ ชานิสฺสามิ อภิสมฺปรายํ, ยํคติกา เต ภวนฺโต ยํอภิสมฺปรายา”ติ.\csromanspacer{} อทฺทสา โข ภควา มาคธเก ปริจารเก “ยํคติกา เต ภวนฺโต ยํอภิสมฺปรายา”ติ.\csromanspacer{} อถ โข ภควา สายนฺหสมยํ ปฏิสลฺลานา วุฏฺฐิโต คิญฺชกาวสถา นิกฺขมิตฺวา วิหารปจฺฉายายํ ปญฺญตฺเต อาสเน นิสีทิ.}

\proseitem{279}{อถ โข อายสฺมา อานนฺโท เยน ภควา เตนุปสงฺกมิ, อุปสงฺกมิตฺวา ภควนฺตํ อภิวาเทตฺวา เอกมนฺตํ นิสีทิ, เอกมนฺตํ นิสินฺโน โข อายสฺมา อานนฺโท ภควนฺตํ เอตทโวจ “อุปสนฺตปทิสฺโส\footnote{อุปสนฺตปติโส (ก)} ภนฺเต ภควา, ภาติริว ภควโต มุขวณฺโณ วิปฺปสนฺนตฺตา อินฺทฺริยานํ.\csromanspacer{} สนฺเตน นูนชฺช ภนฺเต ภควา วิหาเรน วิหาสี”ติ.\csromanspacer{} ยเทว โข เม ตฺวํ อานนฺท มาคธเก ปริจารเก อารพฺภ สมฺมุขา ปริกถํ กตฺวา อุฏฺฐายาสนา ปกฺกนฺโต, ตเทวาหํ นาติเก ปิณฺฑาย จริตฺวา ปจฺฉาภตฺตํ ปิณฺฑปาตปฏิกฺกนฺโต ปาเท ปกฺขาเลตฺวา คิญฺชกาวสถํ ปวิสิตฺวา มาคธเก ปริจารเก อารพฺภ อฏฺฐิํกตฺวา มนสิกตฺวา สพฺพํ เจตสา สมนฺนาหริตฺวา ปญฺญตฺเต อาสเน นิสีทิํ “คติํ เนสํ ชานิสฺสามิ อภิสมฺปรายํ, ยํคติกา เต ภวนฺโต ยํอภิสมฺปรายา”ติ.\csromanspacer{} อทฺทสํ โข อหํ อานนฺท มาคธเก ปริจารเก “ยํคติกา เต ภวนฺโต ยํอภิสมฺปรายา”ติ.}

\csromanmatikaanchor{mtk.85}
\csromantocmarkref{section}{ชนวสภยกฺข}{mtk.85}
\bookmark[dest={mtk.85},level=1]{ชนวสภยกฺข}
\csromanheader{1}{\textbf{ชนวสภยกฺข}}
\proseitem{280}{อถ โข อานนฺท อนฺตรหิโต ยกฺโข สทฺทมนุสฺสาเวสิ “ชนวสโภ อหํ ภควา ชนวสโภ อหํ สุคตา”ติ.\csromanspacer{} อภิชานาสิ โน ตฺวํ อานนฺท อิโต ปุพฺเพ เอวรูปํ นามเธยฺยํ สุตํ\footnote{สุตฺวา (อิ)} ยทิทํ ชนวสโภติ.}

\prose{น โข อหํ ภนฺเต อภิชานามิ อิโต ปุพฺเพ เอวรูปํ นามเธยฺยํ สุตํ ยทิทํ ชนวสโภติ, อปิ จ เม ภนฺเต โลมานิ หฏฺฐานิ “ชนวสโภ”ติ นามเธยฺยํ สุตฺวา.\csromanspacer{} ตสฺส มยฺหํ ภนฺเต เอตทโหสิ “น หิ นูน โส โอรโก ยกฺโข ภวิสฺสติ ยทิทํ เอวรูปํ นามเธยฺยํ สุปญฺญตฺตํ ยทิทํ ชนวสโภ”ติ.\csromanspacer{} อนนฺตรา โข อานนฺท สทฺทปาตุภาวา}

\vspace{\tipitakaparskip}
\csromancenter{ชนวสภสุตฺต อุฬารวณฺโณ เม ยกฺโข สมฺมุเข ปาตุรโหสิ.\csromanspacer{} ทุติยกมฺปิ สทฺทมนุสฺสาเวสิ “พิมฺพิสาโร อหํ ภควา, พิมฺพิสาโร อหํ สุคตา”ติ.\csromanspacer{} อิทํ สตฺตมํ โข อหํ ภนฺเต เวสฺสวณสฺส มหาราชสฺส สหพฺยตํ อุปปชฺชามิ, โส ตโต จุโต มนุสฺสราชา ภวิตุํ ปโหมิ\footnote{โส ตโต จุโต มนุสฺสราชา, อมนุสฺสราชา ทิวิ โหมิ (สี, อิ)}.}
\vspace{0.5\baselineskip}
\csromangathameasure{อิโต สตฺต ตโต สตฺต, สํสารานิ จตุทฺทส. \\ นิวาสมภิชานามิ, ยตฺถ เม วุสิตํ ปุเร.}
\csromangathagroup{\csromangathastanza{\csromanfolio{167}อิโต สตฺต ตโต สตฺต, สํสารานิ จตุทฺทส. \\ นิวาสมภิชานามิ, ยตฺถ เม วุสิตํ ปุเร.}}

\proseitem{281}{ทีฆรตฺตํ โข อหํ ภนฺเต อวินิปาโต อวินิปาตํ สญฺชานามิ, อาสา จ ปน เม สนฺติฏฺฐติ สกทาคามิตายาติ.\csromanspacer{} อจฺฉริยมิทํ อายสฺมโต ชนวสภสฺส ยกฺขสฺส อพฺภุตมิทํ อายสฺมโต ชนวสภสฺส ยกฺขสฺส. “ทีฆรตฺตํ โข อหํ ภนฺเต อวินิปาโต อวินิปาตํ สญฺชานามี”ติ จ วเทสิ, “อาสา จ ปน เม สนฺติฏฺฐติ สกทาคามิตายา”ติ จ วเทสิ, กุโตนิทานํ ปนายสฺมา ชนวสโภ ยกฺโข เอวรูปํ อุฬารํ วิเสสาธิคมํ สญฺชานาตีติ.\csromanspacer{} น อญฺญตฺร ภควา ตว สาสนา น อญฺญตฺร\footnote{อญฺญตฺถ (สี, อิ)} สุคต ตว สาสนา, ยทคฺเค อหํ ภนฺเต ภควติ เอกนฺติกโต\footnote{เอกนฺตโต (สฺยา), เอกนฺตคโต (อิ)} อภิปฺปสนฺโน, ตทคฺเค อหํ ภนฺเต ทีฆรตฺตํ อวินิปาโต อวินิปาตํ สญฺชานามิ, อาสา จ ปน เม สนฺติฏฺฐติ สกทาคามิตาย.\csromanspacer{} อิธาหํ ภนฺเต เวสฺสวเณน มหาราเชน เปสิโต วิรูฬฺหกสฺส มหาราชสฺส สนฺติเก เกนจิเทว กรณีเยน อทฺทสํ ภควนฺตํ อนฺตรามคฺเค คิญฺชกาวสถํ ปวิสิตฺวา มาคธเก ปริจารเก อารพฺภ อฏฺฐิํกตฺวา มนสิกตฺวา สพฺพํ เจตสา สมนฺนาหริตฺวา นิสินฺนํ “คติํ เนสํ ชานิสฺสามิ อภิสมฺปรายํ, ยํคติกา เต ภวนฺโต ยํอภิสมฺปรายา”ติ.\csromanspacer{} อนจฺฉริยํ โข ปเนตํ ภนฺเต, ยํ เวสฺสวณสฺส มหาราชสฺส ตสฺสํ ปริสายํ ภาสโต สมฺมุขา สุตํ สมฺมุขา ปฏิคฺคหิตํ “ยํคติกา เต ภวนฺโต ยํอภิสมฺปรายา”ติ.\csromanspacer{} ตสฺส มยฺหํ ภนฺเต เอตทโหสิ “ภควนฺตญฺจ ทกฺขามิ อิทญฺจ ภควโต อาโรเจสฺสามี”ติ.\csromanspacer{} อิเม โข เม ภนฺเต เทฺวปจฺจยา ภควนฺตํ ทสฺสนาย อุปสงฺกมิตุํ.}

\csromanmatikaanchor{mtk.86}
\csromantocmarkref{section}{เทวสภา}{mtk.86}
\bookmark[dest={mtk.86},level=1]{เทวสภา}
\csromanheader{1}{\textbf{เทวสภา}}
\proseitem{282}{\csromanfolio{168}ปุริมานิ ภนฺเต ทิวสานิ ปุริมตรานิ ตทหุโปสเถ ปนฺนรเส วสฺสูปนายิกาย ปุณฺณาย ปุณฺณมาย รตฺติยา เกวลกปฺปา จ เทวา ตาวติํสา สุธมฺมายํ สภายํ สนฺนิสินฺนา โหนฺติ สนฺนิปติตา, มหตี จ ทิพฺพปริสา\footnote{ทิพฺพา ปริสา (สี, อิ)} สมนฺตโต นิสินฺนา โหนฺติ\footnote{นิสินฺนา โหติ (สี), สนฺนิสินฺนา โหนฺติ สนฺนิปติตา (ก) 3. ปจฺฉาภิมุโข (ก)}, จตฺตาโร จ มหาราชาโน จตุทฺทิสา นิสินฺนา โหนฺติ.\csromanspacer{} ปุรตฺถิมาย ทิสาย ธตรฏฺโฐ มหาราชา ปจฺฉิมาภิมุโข3 นิสินฺโน โหติ เทเว ปุรกฺขตฺวา, ทกฺขิณาย ทิสาย วิรูฬฺหโก มหาราชา อุตฺตราภิมุโข นิสินฺโน โหติ เทเว ปุรกฺขตฺวา, ปจฺฉิมาย ทิสาย วิรูปกฺโข มหาราชา ปุรตฺถาภิมุโข นิสินฺโน โหติ เทเว ปุรกฺขตฺวา, อุตฺตราย ทิสาย เวสฺสวโณ มหาราชา ทกฺขิณาภิมุโข นิสินฺโน โหติ เทเว ปุรกฺขตฺวา.\csromanspacer{} ยทา ภนฺเต เกวลกปฺปา จ เทวา ตวติํสา สุธมฺมายํ สภายํ สนฺนิสินฺนา โหนฺติ สนฺนิปติตา, มหตี จ ทิพฺพปริสา สมนฺตโต นิสินฺนา โหนฺติ, จตฺตาโร จ มหาราชาโน จตุทฺทิสา นิสินฺนา โหนฺติ, อิทํ เนสํ โหติ อาสนสฺมิํ, อถ ปจฺฉา อมฺหากํ อาสนํ โหติ.\csromanspacer{} เย เต ภนฺเต เทวา ภควติ พฺรหฺมจริยํ จริตฺวา อธุนูปปนฺนา ตาวติํสกายํ, เต อญฺเญ เทเว อติโรจนฺติ วณฺเณน เจว ยสสา จ, เตน สุทํ ภนฺเต เทวา ตาวติํสา อตฺตมนา โหนฺติ ปมุทิตา ปีติโสมนสฺสชาตา “ทิพฺพา วต โภ กายา ปริปูเรนฺติ, หายนฺติ อสุรกายา”ติ.\csromanspacer{} อถ โข ภนฺเต สกฺโก เทวานมินฺโท เทวานํ ตาวติํสานํ สมฺปสาทํ วิทิตฺวา อิมาหิ คาถาหิ อนุโมทิ\csromandash{}}

\csromangathameasure{“โมทนฺติ วต โภ เทวา, ตาวติํสา สหินฺทกา. \\ ตถาคตํ นมสฺสนฺตา, ธมฺมสฺส จ สุธมฺมตํ. \\ นเว เทเว จ ปสฺสนฺตา, วณฺณวนฺเต ยสสฺสิเน. \\ สุคตสฺมิํ พฺรหฺมจริยํ, จริตฺวาน อิธาคเต. \\ เต อญฺเญ อติโรจนฺติ, วณฺเณน ยสสายุนา. \\ สาวกา ภูริปญฺญสฺส, วิเสสูปคตา อิธ.}
\csromangathagroup{\csromangathastanza{“โมทนฺติ วต โภ เทวา, ตาวติํสา สหินฺทกา\footnote{ส-อินฺทกา (สี)}. \\ ตถาคตํ นมสฺสนฺตา, ธมฺมสฺส จ สุธมฺมตํ.}\csromangathastanza{นเว เทเว จ ปสฺสนฺตา, วณฺณวนฺเต ยสสฺสิเน\footnote{ยสสฺสิโน (สฺยา)}. \\ สุคตสฺมิํ พฺรหฺมจริยํ, จริตฺวาน อิธาคเต.}\csromangathastanza{เต อญฺเญ อติโรจนฺติ, วณฺเณน ยสสายุนา. \\ สาวกา ภูริปญฺญสฺส, วิเสสูปคตา อิธ.}}

\csromanheader{2}{5. ชนวสภสุตฺต}
\csromangathameasure{อิทํ ทิสฺวาน นนฺทนฺติ, ตาวติํสา สหินฺทกา. \\ ตถาคตํ นมสฺสนฺตา, ธมฺมสฺส จ สุธมฺมตนฺ”ติ.}
\csromangathagroup{\csromangathastanza{\csromanfolio{169}อิทํ ทิสฺวาน นนฺทนฺติ, ตาวติํสา สหินฺทกา. \\ ตถาคตํ นมสฺสนฺตา, ธมฺมสฺส จ สุธมฺมตนฺ”ติ.}}

\prose{เตน สุทํ ภนฺเต เทวา ตาวติํสา ภิยฺโยโส มตฺตาย อตฺตมนา โหนฺติ ปมุทิตา ปีติโสมนสฺสชาตา “ทิพฺพา วต โภ กายา ปริปูเรนฺติ, หายนฺติ อสุรกายา”ติ.\csromanspacer{} อถ โข ภนฺเต เยนตฺเถน เทวา ตาวติํสา สุธมฺมายํ สภายํ สนฺนิสินฺนา โหนฺติ สนฺนิปติตา, ตํ อตฺถํ จินฺตยิตฺวา ตํ อตฺถํ มนฺตยิตฺวา วุตฺตวจนาปิ ตํ\footnote{วุตฺตวจนา นามิทํ (ก)} จตฺตาโร มหาราชาโน ตสฺมิํ อตฺเถ โหนฺติ, ปจฺจานุสิฏฺฐวจนาปิ ตํ\footnote{ปจฺจานุสิฏฺฐวจนา นามิทํ (ก)} จตฺตาโร มหาราชาโน ตสฺมิํ อตฺเถ โหนฺติ, สเกสุ สเกสุ อาสเนสุ ฐิตา อวิปกฺกนฺตา\footnote{อธิปกฺกนฺตา (ก)}.}

\csromangathameasure{เต วุตฺตวากฺยา ราชาโน, ปฏิคฺคยฺหานุสาสนิํ. \\ วิปฺปสนฺนมนา สนฺตา, อฏฺฐํสุ สมฺหิ อาสเนติ.}
\csromangathagroup{\csromangathastanza{เต วุตฺตวากฺยา ราชาโน, ปฏิคฺคยฺหานุสาสนิํ. \\ วิปฺปสนฺนมนา สนฺตา, อฏฺฐํสุ สมฺหิ อาสเนติ.}}

\proseitem{283}{อถ โข ภนฺเต อุตฺตราย ทิสาย อุฬาโร อาโลโก สญฺชายิ, โอภาโส ปาตุรโหสิ อติกฺกมฺเมว เทวานํ เทวานุภาวํ.\csromanspacer{} อถ โข ภนฺเต สกฺโก เทวานมินฺโท เทเว ตาวติํเส อามนฺเตสิ “ยถา โข มาริสา นิมิตฺตานิ ทิสฺสนฺติ, อุฬาโร อาโลโก สญฺชายติ, โอภาโส ปาตุ ภวติ, พฺรหฺมา ปาตุ ภวิสฺสติ.\csromanspacer{} พฺรหฺมุโน เหตํ ปุพฺพนิมิตฺตํ ปาตุภาวาย, ยทิทํ อาโลโก สญฺชายติ โอภาโส ปาตุ ภวตี”ติ.}

\csromangathameasure{ยถา นิมิตฺตา ทิสฺสนฺติ, พฺรหฺมา ปาตุ ภวิสฺสติ. \\ พฺรหฺมุโน เหตํ นิมิตฺตํ, โอภาโส วิปุโล มหาติ.}
\csromangathagroup{\csromangathastanza{ยถา นิมิตฺตา ทิสฺสนฺติ, พฺรหฺมา ปาตุ ภวิสฺสติ. \\ พฺรหฺมุโน เหตํ นิมิตฺตํ, โอภาโส วิปุโล มหาติ.}}

\csromanmatikaanchor{mtk.87}
\csromantocmarkref{section}{สนงฺกุมารกถา}{mtk.87}
\bookmark[dest={mtk.87},level=1]{สนงฺกุมารกถา}
\csromanheader{1}{\textbf{สนงฺกุมารกถา}}
\proseitem{284}{อถ โข ภนฺเต เทวา ตาวติํสา ยถาสเกสุ อาสเนสุ นิสีทิํสุ “โอภาสเมตํ ญสฺสาม, ยํวิปาโก ภวิสฺสติ, สจฺฉิกตฺวาว นํ คมิสฺสามา”ติ.\csromanspacer{} จตฺตาโรปิ มหาราชาโน ยถาสเกสุ อาสเนสุ นิสีทิํสุ “โอภาสเมตํ ญสฺสาม, ยํวิปาโก ภวิสฺสติ, สจฺฉิกตฺวาว นํ คมิสฺสามา”ติ.\csromanspacer{} อิทํ สุตฺวา เทวา ตาวติํสา \csromanfolio{170}เอกคฺคา สมาปชฺชิํสุ “โอภาสเมตํ ญสฺสาม, ยํวิปาโก ภวิสฺสติ, สจฺฉิกตฺวาว นํ คมิสฺสามา”ติ.}

\prose{ยทา ภนฺเต พฺรหฺมา สนงฺกุมาโร เทวานํ ตาวติํสานํ ปาตุ ภวติ, โอฬาริกํ อตฺตภาวํ อภินิมฺมินิตฺวา ปาตุ ภวติ.\csromanspacer{} โย โข ปน ภนฺเต พฺรหฺมุโน ปกติวณฺโณ, อนภิสมฺภวนีโย โส เทวานํ ตาวติํสานํ จกฺขุปถสฺมิํ.\csromanspacer{} ยทา ภนฺเต พฺรหฺมา สนงฺกุมาโร เทวานํ ตาวติํสานํ ปาตุ ภวติ, โส อญฺเญ เทเว อติโรจติ วณฺเณน เจว ยสสา จ.\csromanspacer{} เสยฺยถาปิ ภนฺเต โสวณฺโณ วิคฺคโห มานุสํ วิคฺคหํ อติโรจติ, เอวเมว โข ภนฺเต ยทา พฺรหฺมา สนงฺกุมาโร เทวานํ ตาวติํสานํ ปาตุ ภวติ, โส อญฺเญ เทเว อติโรจติ วณฺเณน เจว ยสสา จ.\csromanspacer{} ยทา ภนฺเต พฺรหฺมา สนงฺกุมาโร เทวานํ ตาวติํสานํ ปาตุ ภวติ, น ตสฺสํ ปริสายํ โกจิ เทโว อภิวาเทติ วา ปจฺจุฏฺเฐติ วา อาสเนน วา นิมนฺเตติ.\csromanspacer{} สพฺเพว ตุณฺหีภูตา ปญฺชลิกา ปลฺลงฺเกน นิสีทนฺติ “ยสฺสทานิ เทวสฺส ปลฺลงฺกํ อิจฺฉิสฺสติ พฺรหฺมา สนงฺกุมาโร, ตสฺส เทวสฺส ปลฺลงฺเก นิสีทิสฺสตี”ติ.}

\prose{ยสฺส โข ปน ภนฺเต เทวสฺส พฺรหฺมา สนงฺกุมาโร ปลฺลงฺเก นิสีทติ, อุฬารํ โส ลภติ เทโว เวทปฏิลาภํ, อุฬารํ โส ลภติ เทโว โสมนสฺสปฏิลาภํ.\csromanspacer{} เสยฺยถาปิ ภนฺเต ราชา ขตฺติโย มุทฺธาวสิตฺโต อธุนาภิสิตฺโต รชฺเชน, อุฬารํ โส ลภติ เวทปฏิลาภํ, อุฬารํ โส ลภติ โสมนสฺสปฏิลาภํ.\csromanspacer{} เอวเมว โข ภนฺเต ยสฺส เทวสฺส พฺรหฺมา สนงฺกุมาโร ปลฺลงฺเก นิสีทติ, อุฬารํ โส ลภติ เทโว เวทปฏิลาภํ, อุฬารํ โส ลภติ เทโว โสมนสฺสปฏิลาภํ.\csromanspacer{} อถ ภนฺเต พฺรหฺมา สนงฺกุมาโร โอฬาริกํ อตฺตภาวํ อภินิมฺมินิตฺวา กุมารวณฺณี\footnote{กุมารวณฺโณ (สฺยา, ก)} หุตฺวา ปญฺจสิโข เทวานํ ตาวติํสานํ ปาตุรโหสิ, โส เวหาสํ อพฺภุคฺคนฺตฺวา อากาเส อนฺตลิกฺเข ปลฺลงฺเกน นิสีทิ.\csromanspacer{} เสยฺยถาปิ ภนฺเต พลวา ปุริโส สุปจฺจตฺถเต วา ปลฺลงฺเก สเม วา ภูมิภาเค ปลฺลงฺเกน นิสีเทยฺย, เอวเมว โข ภนฺเต พฺรหฺมา สนงฺกุมาโร เวหาสํ อพฺภุคฺคนฺตฺวา อากาเส อนฺตลิกฺเข ปลฺลงฺเกน นิสีทิตฺวา เทวานํ ตาวติํสานํ สมฺปสาทํ วิทิตฺวา อิมาหิ คาถาหิ อนุโมทิ\csromandash{}}

\csromanheader{2}{5. ชนวสภสุตฺต}
\csromangathameasure{“โมทนฺติ วต โภ เทวา, ตาวติํสา สหินฺทกา. \\ ตถาคตํ นมสฺสนฺตา, ธมฺมสฺส จ สุธมฺมตํ. \\ นเว เทเว จ ปสฺสนฺตา, วณฺณวนฺเต ยสสฺสิเน. \\ สุคตสฺมิํ พฺรหฺมจริยํ, จริตฺวาน อิธาคเต. \\ เต อญฺเญ อติโรจนฺติ, วณฺเณน ยสสายุนา. \\ สาวกา ภูริปญฺญสฺส, วิเสสูปคตา อิธ. \\ อิทํ ทิสฺวาน นนฺทนฺติ, ตาวติํสา สหินฺทกา. \\ ตถาคตํ นมสฺสนฺตา, ธมฺมสฺส จ สุธมฺมตนฺ”ติ.}
\csromangathagroup{\csromangathastanza{\csromanfolio{171}“โมทนฺติ วต โภ เทวา, ตาวติํสา สหินฺทกา. \\ ตถาคตํ นมสฺสนฺตา, ธมฺมสฺส จ สุธมฺมตํ.}\csromangathastanza{นเว เทเว จ ปสฺสนฺตา, วณฺณวนฺเต ยสสฺสิเน. \\ สุคตสฺมิํ พฺรหฺมจริยํ, จริตฺวาน อิธาคเต.}\csromangathastanza{เต อญฺเญ อติโรจนฺติ, วณฺเณน ยสสายุนา. \\ สาวกา ภูริปญฺญสฺส, วิเสสูปคตา อิธ.}\csromangathastanza{อิทํ ทิสฺวาน นนฺทนฺติ, ตาวติํสา สหินฺทกา. \\ ตถาคตํ นมสฺสนฺตา, ธมฺมสฺส จ สุธมฺมตนฺ”ติ.}}

\proseitem{285}{อิมมตฺถํ ภนฺเต พฺรหฺมา สนงฺกุมาโร ภาสิตฺถ, อิมมตฺถํ ภนฺเต พฺรหฺมุโน สนงฺกุมารสฺส ภาสโต อฏฺฐงฺคสมนฺนาคโต สโร โหติ วิสฺสฏฺโฐ จ วิญฺเญยฺโย จ มญฺชุ จ สวนีโย จ พินฺทุ จ อวิสารี จ คมฺภีโร จ นินฺนาที จ.\csromanspacer{} ยถาปริสํ โข ปน ภนฺเต พฺรหฺมา สนงฺกุมาโร สเรน วิญฺญาเปติ, น จสฺส พหิทฺธา ปริสาย โฆโส นิจฺฉรติ.\csromanspacer{} ยสฺส โข ปน ภนฺเต เอวํ อฏฺฐงฺคสมนฺนาคโต สโร โหติ, โส วุจฺจติ “พฺรหฺมสฺสโร”ติ.}

\prose{อถ โข ภนฺเต พฺรหฺมา สนงฺกุมาโร เตตฺติํเส อตฺตภาเว อภินิมฺมินิตฺวา เทวานํ ตาวติํสานํ ปจฺเจกปลฺลงฺเกสุ ปลฺลงฺเกน\footnote{ปจฺเจกปลฺลงฺเกน (ก)} นิสีทิตฺวา เทเว ตาวติํเส อามนฺเตสิ “ตํ กิํ มญฺญนฺติ โภนฺโต เทวา ตาวติํสา, ยาวญฺจ โส ภควา พหุชนหิตาย ปฏิปนฺโน พหุชนสุขาย โลกานุกมฺปาย อตฺถาย หิตาย สุขาย เทวมนุสฺสานํ.\csromanspacer{} เย หิ เกจิ โภ พุทฺธํ สรณํ คตา ธมฺมํ สรณํ คตา สํฆํ สรณํ คตา สีเลสุ ปริปูรการิโน.\csromanspacer{} เต กายสฺส เภทา ปรํ มรณา อปฺเปกจฺเจ ปรนิมฺมิตวสวตฺตีนํ เทวานํ สหพฺยตํ อุปปชฺชนฺติ, อปฺเปกจฺเจ นิมฺมานรตีนํ เทวานํ สหพฺยตํ อุปปชฺชนฺติ, อปฺเปกจฺเจ ตุสิตานํ เทวานํ สหพฺยตํ อุปปชฺชนฺติ, อปฺเปกจฺเจ ยามานํ เทวานํ สหพฺยตํ อุปปชฺชนฺติ, อปฺเปกจฺเจ ตาวติํสานํ เทวานํ สหพฺยตํ อุปปชฺชนฺติ, อปฺเปกจฺเจ จาตุมหาราชิกานํ เทวานํ สหพฺยตํ อุปปชฺชนฺติ.\csromanspacer{} เย สพฺพนิหีนํ กายํ ปริปูเรนฺติ, เต คนฺธพฺพกายํ ปริปูเรนฺตี”ติ.}

\proseitem{286}{\csromanfolio{172}อิมมตฺถํ ภนฺเต พฺรหฺมา สนงฺกุมาโร ภาสิตฺถ, อิมมตฺถํ ภนฺเต พฺรหฺมุโน สนงฺกุมารสฺส ภาสโต โฆโสเยว เทวา มญฺญนฺติ “ยฺวายํ มม ปลฺลงฺเก, สฺวายํ เอโกว ภาสตี”ติ.}

\csromangathameasure{เอกสฺมิํ ภาสมานสฺมิํ, สพฺเพ ภาสนฺติ นิมฺมิตา. \\ เอกสฺมิํ ตุณฺหิมาสีเน, สพฺเพ ตุณฺหี ภวนฺติ เต. \\ ตทาสุ เทวา มญฺญนฺติ, ตาวติํสา สหินฺทกา. \\ ยฺวายํ มม ปลฺลงฺกสฺมิํ, สฺวายํ เอโกว ภาสตีติ.}
\csromangathagroup{\csromangathastanza{เอกสฺมิํ ภาสมานสฺมิํ, สพฺเพ ภาสนฺติ นิมฺมิตา. \\ เอกสฺมิํ ตุณฺหิมาสีเน, สพฺเพ ตุณฺหี ภวนฺติ เต.}\csromangathastanza{ตทาสุ เทวา มญฺญนฺติ, ตาวติํสา สหินฺทกา. \\ ยฺวายํ มม ปลฺลงฺกสฺมิํ, สฺวายํ เอโกว ภาสตีติ.}}

\prose{อถ โข ภนฺเต พฺรหฺมา สนงฺกุมาโร เอกตฺเตน อตฺตานํ อุปสํหรติ, เอกตฺเตน อเตนํ อุปสํหริตฺวา สกฺกสฺส เทวานมินฺทสฺส ปลฺลงฺเก ปลฺลงฺเกน นิสีทิตฺวา เทเว ตาวติํเส อามนฺเตสิ\csromandash{}}

\csromanmatikaanchor{mtk.88}
\csromantocmarkref{section}{ภาวิต-อิทฺธิปาท}{mtk.88}
\bookmark[dest={mtk.88},level=1]{ภาวิต-อิทฺธิปาท}
\csromanheader{1}{\textbf{ภาวิต-อิทฺธิปาท}}
\proseitem{287}{ตํ กิํ มญฺญนฺติ โภนฺโต เทวา ตาวติํสา, ยาว สุปญฺญตฺตา จิเม เตน ภควตา ชานตา ปสฺสตา อรหตา สมฺมาสมฺพุทฺเธน จตฺตาโร อิทฺธิปาทา ปญฺญตฺตา อิทฺธิปหุตาย\footnote{อิทฺธิพหุลีกตาย (สฺยา)} อิทฺธิวิสวิตาย\footnote{อิทฺธิวิเสวิตาย (สฺยา)} อิทฺธิวิกุพฺพนตาย.\csromanspacer{} กตเม จตฺตาโร.\csromanspacer{} อิธ โภ ภิกฺขุ ฉนฺทสมาธิปฺปธานสงฺขารสมนฺนาคตํ อิทฺธิปาทํ ภาเวติ.\csromanspacer{} วีริยสมาธิปฺปธานสงฺขารสมนฺนาคตํ อิทฺธิปาทํ ภาเวติ.\csromanspacer{} จิตฺตสมาธิปฺปธานสงฺขารสมนฺนาคตํ อิทฺธิปาทํ ภาเวติ.\csromanspacer{} วีมํสาสมาธิปฺปธานสงฺขารสมนฺนาคตํ อิทฺธิปาทํ ภาเวติ.\csromanspacer{} อิเม โข โภ เตน ภควตา ชานตา ปสฺสตา อรหตา สมฺมาสมฺพุทฺเธน จตฺตาโร อิทฺธิปาทา ปญฺญตฺตา อิทฺธิปหุตาย อิทฺธิวิสวิตาย อิทฺธิวิกุพฺพนตาย.}

\prose{เย หิ เกจิ โภ อตีตมทฺธานํ สมณา วา พฺราหฺมณา วา อเนกวิหิตํ อิทฺธิวิธํ ปจฺจนุโภสุํ, สพฺเพ เต อิเมสํเยว จตุนฺนํ อิทฺธิปาทานํ ภาวิตตฺตา พหุลีกตตฺตา.\csromanspacer{} เยปิ หิ เกจิ โภ อนาคตมทฺธานํ สมณา วา พฺราหฺมณา วา อเนกวิหิตํ อิทฺธิวิธํ ปจฺจนุโภสฺสนฺติ, สพฺเพ เต อิเมสํเยว จตุนฺนํ อิทฺธิปาทานํ ภาวิตตฺตา พหุลีกตตฺตา.\csromanspacer{} เยปิ หิ เกจิ โภ เอตรหิ สมณา วา พฺราหฺมณา วา อเนกวิหิตํ อิทฺธิวิธํ ปจฺจนุโภนฺติ, สพฺเพ เต อิเมสํเยว จตุนฺนํ อิทฺธิปาทานํ ภาวิตตฺตา พหุลีกตตฺตา.\csromanspacer{} ปสฺสนฺติ โน โภนฺโต เทวา ตาวติํสา}

\vspace{\tipitakaparskip}
\csromancenter{ชนวสภสุตฺต มมปิมํ เอวรูปํ อิทฺธานุภาวนฺติ.\csromanspacer{} เอวํ มหาพฺรหฺเมติ.\csromanspacer{} อหมฺปิ โข โภ อิเมสํเยว จตุนฺนญฺจ อิทฺธิปาทานํ ภาวิตตฺตา พหุลีกตตฺตา เอวํมหิทฺธิโก เอวํมหานุภาโวติ.\csromanspacer{} อิมมตฺถํ ภนฺเต พฺรหฺมา สนงฺกุมาโร ภาสิตฺถ.\csromanspacer{} อิมมตฺถํ ภนฺเต พฺรหฺมา สนงฺกุมาโร ภาสิตฺวา เทเว ตาวติํเส อามนฺเตสิ.}
\csromanmatikaanchor{mtk.89}
\csromantocmarkref{section}{ติวิธ-โอกาสาธิคม}{mtk.89}
\bookmark[dest={mtk.89},level=1]{ติวิธ-โอกาสาธิคม}
\csromanheader{1}{\textbf{ติวิธ โอกาสาธิคม}}
\proseitem{288}{\csromanfolio{173}“ตํ กิํ มญฺญนฺติ โภนฺโต เทวา ตาวติํสา, ยาวญฺจิทํ เตน ภควตา ชานตา ปสฺสตา อรหตา สมฺมาสมฺพุทฺเธน ตโย โอกาสาธิคมา อนุพุทฺธา สุขสฺสาธิคมาย.\csromanspacer{} กตเม ตโย.\csromanspacer{} อิธ โภ เอกจฺโจ สํสฏฺโฐ วิหรติ กาเมหิ สํสฏฺโฐ อกุสเลหิ ธมฺเมหิ.\csromanspacer{} โส อปเรน สมเยน อริยธมฺมํ สุณาติ, โยนิโส มนสิ กโรติ, ธมฺมานุธมฺมํ ปฏิปชฺชติ.\csromanspacer{} โส อริยธมฺมสฺสวนํ อาคมฺม โยนิโสมนสิการํ ธมฺมานุธมฺมปฺปฏิปตฺติํ อสํสฏฺโฐ วิหรติ กาเมหิ อสํสฏฺโฐ อกุสเลหิ ธมฺเมหิ.\csromanspacer{} ตสฺส อสํสฏฺฐสฺส กาเมหิ อสํสฏฺฐสฺส อกุสเลหิ ธมฺเมหิ อุปฺปชฺชติ สุขํ, สุขา ภิยฺโย โสมนสฺสํ.\csromanspacer{} เสยฺยถาปิ โภ ปมุทา ปาโมชฺชํ\footnote{ปามุชฺชํ (อิ, ก)} ชาเยถ, เอวเมว โข โภ อสํสฏฺฐสฺส กาเมหิ อสํสฏฺฐสฺส อกุสเลหิ ธมฺเมหิ อุปฺปชฺชติ สุขํ, สุขา ภิยฺโย โสมนสฺสํ.\csromanspacer{} อยํ โข โภ เตน ภควตา ชานตา ปสฺสตา อรหตา สมฺมาสมฺพุทฺเธน ปฐโม โอกาสาธิคโม อนุพุทฺโธ สุขสฺสาธิคมาย.}

\prose{ปุน จปรํ โภ อิเธกจฺจสฺส โอฬาริกา กายสงฺขารา อปฺปฏิปฺปสฺสทฺธา โหนฺติ, โอฬาริกา วจีสงฺขารา อปฺปฏิปฺปสฺสทฺธา โหนฺติ, โอฬาริกา จิตฺตสงฺขารา อปฺปฏิปฺปสฺสทฺธา โหนฺติ.\csromanspacer{} โส อปเรน สมเยน อริยธมฺมํ สุณาติ, โยนิโส มนสิ กโรติ, ธมฺมานุธมฺมํ ปฏิปชฺชติ.\csromanspacer{} ตสฺส อริยธมฺมสฺสวนํ อาคมฺม โยนิโสมนสิการํ ธมฺมานุธมฺมปฺปฏิปตฺติํ โอฬาริกา กายสงฺขารา ปฏิปฺปสฺสมฺภนฺติ, โอฬาริกา วจีสงฺขารา ปฏิปฺปสฺสมฺภนฺติ, โอฬาริกา จิตฺตสงฺขารา ปฏิปฺปสฺสมฺภนฺติ.\csromanspacer{} ตสฺส โอฬาริกานํ กายสงฺขารานํ ปฏิปฺปสฺสทฺธิยา โอฬาริกานํ วจีสงฺขารานํ ปฏิปฺปสฺสทฺธิยา โอฬาริกานํ \csromanfolio{174}จิตฺตสงฺขารานํ ปฏิปฺปสฺสทฺธิยา อุปฺปชฺชติ สุขํ, สุขา ภิยฺโย โสมนสฺสํ.\csromanspacer{} เสยฺยถาปิ โภ ปมุทา ปาโมชฺชํ ชาเยถ, เอวเมว โข โภ โอฬาริกานํ กายสงฺขารานํ ปฏิปฺปสฺสทฺธิยา โอฬาริกานํ วจีสงฺขารานํ ปฏิปฺปสฺสทฺธิยา โอฬาริกานํ จิตฺตสงฺขารานํ ปฏิปฺปสฺสทฺธิยา อุปฺปชฺชติ สุขํ, สุขา ภิยฺโย โสมนสฺสํ.\csromanspacer{} อยํ โข โภ เตน ภควตา ชานตา ปสฺสตา อรหตา สมฺมาสมฺพุทฺเธน ทุติโย โอกาสาธิคโม อนุพุทฺโธ สุขสฺสาธิคมาย.}

\prose{ปุน จปรํ โภ อิเธกจฺโจ “อิทํ กุสลนฺ”ติ ยถาภูตํ นปฺปชานาติ, “อิทํ อกุสลนฺ”ติ ยถาภูตํ นปฺปชานาติ. “อิทํ สาวชฺชํ อิทํ อนวชฺชํ, อิทํ เสวิตพฺพํ อิทํ นเสวิตพฺพํ, อิทํ หีนํ อิทํ ปณีตํ, อิทํ กณฺหสุกฺกสปฺปฏิภาคนฺ”ติ ยถาภูตํ นปฺปชานาติ.\csromanspacer{} โส อปเรน สมเยน อริยธมฺมํ สุณาติ, โยนิโสมนสิกโรติ, ธมฺมานุธมฺมํ ปฏิปชฺชติ.\csromanspacer{} โส อริยธมฺมสฺสวนํ อาคมฺม โยนิโสมนสิการํ ธมฺมานุธมฺมปฺปฏิปตฺติํ “อิทํ กุสลนฺ”ติ ยถาภูตํ ปชานาติ, “อิทํ อกุสลนฺ”ติ ยถาภูตํ ปชานาติ. “อิทํ สาวชฺชํ อิทํ อนวชฺชํ, อิทํ เสวิตพฺพํ อิทํ น เสวิตพฺพํ, อิทํ หีนํ อิทํ ปณีตํ, อิทํ กณฺหสุกฺกสปฺปฏิภาคนฺ”ติ ยถาภูตํ ปชานาติ.\csromanspacer{} ตสฺส เอวํ ชานโต เอวํ ปสฺสโต อวิชฺชา ปหียติ, วิชฺชา อุปฺปชฺชติ.\csromanspacer{} ตสฺส อวิชฺชาวิราคา วิชฺชุปฺปาทา อุปฺปชฺชติ สุขํ, สุขา ภิยฺโย โสมนสฺสํ.\csromanspacer{} เสยฺยถาปิ โภ ปมุทา ปาโมชฺชํ ชาเยถ, เอวเมว โข โภ อวิชฺชาวิราคา วิชฺชุปฺปาทา อุปฺปชฺชติ สุขํ, สุขา ภิยฺโย โสมนสฺสํ.\csromanspacer{} อยํ โข โภ เตน ภควตา ชานตา ปสฺสตา อรหตา สมฺมาสมฺพุทฺเธน ตติโย โอกาสาธิคโม อนุพุทฺโธ สุขสฺสาธิคมาย.\csromanspacer{} อิเม โข โภ เตน ภควตา ชานตา ปสฺสตา อรหตา สมฺมาสมฺพุทฺเธน ตโย โอกาสาธิคมา อนุพุทฺธา สุขสฺสาธิคมายา”ติ.\csromanspacer{} อิมมตฺถํ ภนฺเต พฺรหฺมา สนงฺกุมาโร ภาสิตฺถ, อิมมตฺถํ ภนฺเต พฺรหฺมา สนงฺกุมาโร ภาสิตฺวา เทเว ตาวติํเส อามนฺเตสิ\csromandash{}}

\csromanmatikaanchor{mtk.90}
\csromantocmarkref{section}{จตุสติปฏฺฐาน}{mtk.90}
\bookmark[dest={mtk.90},level=1]{จตุสติปฏฺฐาน}
\csromanheader{1}{\textbf{จตุสติปฏฺฐาน}}
\proseitem{289}{“ตํ กิํ มญฺญนฺติ โภนฺโต เทวา ตาวติํสา, ยาว สุปญฺญตฺตา จิเม เตน ภควตา ชานตา ปสฺสตา อรหตา สมฺมาสมฺพุทฺเธน จตฺตาโร สติปฏฺฐานา ปญฺญตฺตา กุสลสฺสาธิคมาย.\csromanspacer{} กตเม จตฺตาโร.}

\vspace{\tipitakaparskip}
\csromancenter{ชนวสภสุตฺต อิธ โภ ภิกฺขุ อชฺฌตฺตํ กาเย กายานุปสฺสี วิหรติ อาตาปี สมฺปชาโน สติมา วิเนยฺย โลเก อภิชฺฌาโทมนสฺสํ.\csromanspacer{} อชฺฌตฺตํ กาเย กายานุปสฺสี วิหรนฺโต ตตฺถ สมฺมา สเมธิยติ, สมฺมา วิปฺปสีทติ.\csromanspacer{} โส ตตฺถ สมฺมา สมาหิโต สมฺมา วิปฺปสนฺโน พหิทฺธา ปรกาเย ญาณทสฺสนํ อภินิพฺพตฺเตติ.\csromanspacer{} อชฺฌตฺตํ เวทนาสุ เวทนานุปสฺสี วิหรติ ฯเปฯ พหิทฺธา ปรเวทนาสุ ญาณทสฺสนํ อภินิพฺพตฺเตติ.\csromanspacer{} อชฺฌตฺตํ จิตฺเต จิตฺตานุปสฺสี วิหรติ ฯเปฯ พหิทฺธา ปรจิตฺเต ญาณทสฺสนํ อภินิพฺพตฺเตติ.\csromanspacer{} อชฺฌตฺตํ ธมฺเมสุ ธมฺมานุปสฺสี วิหรติ อาตาปี สมฺปชาโน สติมา วิเนยฺย โลเก อภิชฺฌาโทมนสฺสํ.\csromanspacer{} อชฺฌตฺตํ ธมฺเมสุ ธมฺมานุปสฺสี วิหรนฺโต ตตฺถ สมฺมา สมาธิยติ, สมฺมา วิปฺปสีทติ.\csromanspacer{} โส ตตฺถ สมฺมา สมาหิโต สมฺมา วิปฺปสนฺโน พหิทฺธา ปรธมฺเมสุ ญาณทสฺสนํ อภินิพฺพตฺเตติ.\csromanspacer{} อิเม โข โภ เตน ภควตา ชานตา ปสฺสตา อรหตา สมฺมาสมฺพุทฺเธน จตฺตาโร สติปฏฺฐานา ปญฺญตฺตา กุสลสฺสาธิคมายา”ติ.\csromanspacer{} อิมมตฺถํ ภนฺเต พฺรหฺมา สนงฺกุมาโร ภาสิตฺถ.\csromanspacer{} อิมมตฺถํ ภนฺเต พฺรหฺมา สนงฺกุมาโร ภาสิตฺวา เทเว ตาวติํเส อามนฺเตสิ\csromandash{}}
\csromanmatikaanchor{mtk.91}
\csromantocmarkref{section}{สตฺตสมาธิปริกฺขาร}{mtk.91}
\bookmark[dest={mtk.91},level=1]{สตฺตสมาธิปริกฺขาร}
\csromanheader{1}{\textbf{สตฺตสมาธิปริกฺขาร}}
\proseitem{290}{\csromanfolio{175}“ตํ กิํ มญฺญนฺติ โภนฺโต เทวา ตาวติํสา, ยาว สุปญฺญตฺตา จิเม เตน ภควตา ชานตา ปสฺสตา อรหตา สมฺมาสมฺพุทฺเธน สตฺต สมาธิปริกฺขารา สเมสมาธิสฺส ปริภาวนาย สมฺมาสมาธิสฺส ปาริปูริยา.\csromanspacer{} กตเม สตฺต.\csromanspacer{} สมฺมาทิฏฺฐิ สมฺมาสงฺกปฺโป สมฺมาวาจา สมฺมากมฺมนฺโต สมฺมา-อาชีโว สมฺมาวายาโม สมฺมาสติ.\csromanspacer{} ยา โข โภ อิเมหิ สตฺตหงฺเคหิ จิตฺตสฺส เอกคฺคตา ปริกฺขตา.\csromanspacer{} อยํ วุจฺจติ โภ อริโย สมฺมาสมาธิ ส-อุปนิโส อิติปิ สปริกฺขาโร อิติปิ.\csromanspacer{} สมฺมาทิฏฺฐิสฺส โภ สมฺมาสงฺกปฺโป ปโหติ, สมฺมาสงฺกปฺปสฺส สมฺมาวาจา ปโหติ, สมฺมาวาจสฺส สมฺมากมฺมนฺโต ปโหติ.\csromanspacer{} สมฺมากมฺมนฺตสฺส สมฺมา-อาชีโว ปโหติ, สมฺมา-อาชีวสฺส สมฺมาวายาโม ปโหติ, สเมวายามสฺส สมฺมาสติ ปโหติ, สมฺมาสติสฺส สมฺมาสมาธิ ปโหติ, สเมสมาธิสฺส สมฺมาญาณํ ปโหติ, สมฺมาญาณสฺส สมฺมาวิมุตฺติ ปโหติ.\csromanspacer{} ยญฺหิ ตํ โภ สมฺมา วทมาโน วเทยฺย “สฺวากฺขาโต ภควตา ธมฺโม สนฺทิฏฺฐิโก อกาลิโก เอหิปสฺสิโก โอปเนยฺยิโก ปจฺจตฺตํ เวทิตพฺโพ วิญฺญูหิ อปารุตา \csromanfolio{176}อมตสฺส ทฺวารา”ติ, อิทเมว ตํ สมฺมา วทมาโน วเทยฺย.\csromanspacer{} สฺวากฺขาโต หิ โภ ภควตา ธมฺโม สนฺทิฏฺฐิโก อกาลิโก เอหิปสฺสิโก โอปเนยฺยิโก ปจฺจตฺตํ เวทิตพฺโพ วิญฺญูหิ อปารุตา อมตสฺส ทฺวารา\footnote{ทฺวาราติ (สฺยา, ก)}.}

\prose{เย หิ เกจิ โภ พุทฺเธ อเวจฺจปฺปสาเทน สมนฺนาคตา, ธมฺเม อเวจฺจปฺปสาเทน สมนฺนาคตา, สํเฆ อเวจฺจปฺปสาเทน สมนฺนาคตา, อริยกนฺเตหิ สีเลหิ สมนฺนาคตา.\csromanspacer{} เย จิเม โอปปาติกา ธมฺมวินีตา สาติเรกานิ จตุวีสติสตสหสฺสานิ มาคธกา ปริจารกา อพฺภตีตา กาลงฺกตา ติณฺณํ สํโยชนานํ ปริกฺขยา โสตาปนฺนา อวินิปาตธมฺมา นิยตา สมฺโพธิปรายณา.\csromanspacer{} อตฺถิ เจเวตฺถ สกทาคามิโน.}

\csromangathameasure{อตฺถายํ อิตรา ปชา, ปุญฺญภาคาติ เม มโน. \\ สงฺขาตุํ โนปิ สกฺโกมิ, มุสาวาทสฺส โอตฺตปฺปนฺ”ติ.}
\csromangathagroup{\csromangathastanza{อตฺถายํ\footnote{อถายํ (สี, สฺยา)} อิตรา ปชา, ปุญฺญภาคาติ เม มโน. \\ สงฺขาตุํ โนปิ สกฺโกมิ, มุสาวาทสฺส โอตฺตปฺปนฺ”ติ.}}

\proseitem{291}{อิมมตฺถํ ภนฺเต พฺรหฺมา สนงฺกุมาโร ภาสิตฺถ, อิมมตฺถํ ภนฺเต พฺรหฺมุโน สนงฺกุมารสฺส ภาสโต เวสฺสวณสฺส มหาราชสฺส เอวํ เจตโส ปริวิตกฺโก อุทปาทิ “อจฺฉริยํ วต โภ อพฺภุตํ วต โภ, เอวรูโปปิ นาม อุฬาโร สตฺถา ภวิสฺสติ, เอวรูปํ อุฬารํ ธมฺมกฺขานํ, เอวรูปา อุฬารา วิเสสาธิคมา ปญฺญายิสฺสนฺตี”ติ.\csromanspacer{} อถ ภนฺเต พฺรหฺมา สนงฺกุมาโร เวสฺสวณสฺส มหาราชสฺส เจตสา เจโตปริวิตกฺกมญฺญาย เวสฺสวณํ มหาราชานํ เอตทโวจ “ตํ กิํ มญฺญติ ภวํ เวสฺสวโณ มหาราชา อตีตมฺปิ อทฺธานํ เอวรูโป อุฬาโร สตฺถา อโหสิ, เอวรูปํ อุฬารํ ธมฺมกฺขานํ, เอวรูปา อุฬารา วิเสสาธิคมา ปญฺญายิํสุ.\csromanspacer{} อนาคตมฺปิ อทฺธานํ เอวรูโป อุฬาโร สตฺถา ภวิสฺสติ, เอวรูปํ อุฬารํ ธมฺมกฺขานํ, เอวรูปา อุฬารา วิเสสาธิคมา ปญฺญายิสฺสนฺตี”ติ.}

\proseitem{292}{อิมมตฺถํ ภนฺเต พฺรหฺมา สนงฺกุมาโร เทวานํ ตาวติํสานํ อภาสิ, อิมมตฺถํ เวสฺสวโณ มหาราชา พฺรหฺมุโน สนงฺกุมารสฺส}

\vspace{\tipitakaparskip}
\csromancenter{ชนวสภสุตฺต เทวานํ ตาวติํสานํ ภาสโต สมฺมุขาสุตํ\footnote{สุตฺวา (สี, อิ)} สมฺมุขา ปฏิคฺคหิตํ\footnote{ปฏิคฺคเหตฺวา (สี, อิ)} สยํ ปริสายํ อาโรเจสิ, อิมมตฺถํ ชนวสโภ ยกฺโข เวสฺสวณสฺส มหาราชสฺส สยํ ปริสายํ ภาสโต สมฺมุขา สุตํ1 สมฺมุขา ปฏิคฺคหิตํ2 ภควโต อาโรเจสิ.\csromanspacer{} อิมมตฺถํ ภควา ชนวสภสฺส ยกฺขสฺส สมฺมุขา สุตฺวา สมฺมุขา ปฏิคฺคเหตฺวา สามญฺจ อภิญฺญาย อายสฺมโต อานนฺทสฺส อาโรเจสิ, อิมมตฺถมายสฺมา อานนฺโท ภควโต สมฺมุขา สุตฺวา สมฺมุขา ปฏิคฺคเหตฺวา อาโรเจสิ ภิกฺขูนํ ภิกฺขุนีนํ อุปาสกานํ อุปาสิกานํ.\csromanspacer{} ตยิทํ พฺรหฺมจริยํ อิทฺธญฺเจว ผีตญฺจ วิตฺถาริกํ พาหุชญฺญํ ปุถุภูตํ ยาว เทวมนุสฺเสหิ สุปฺปกาสิตนฺติ.}
\nitthitam{\textbf{ชนวสภสุตฺตํ นิฏฺฐิตํ ปญฺจมํ.}}
\csromanmatikaanchor{mtk.92}
\csromantocmarknopage{chapter}{6. มหาโควินฺทสุตฺต}{mtk.92}
\bookmark[dest={mtk.92},level=0]{6. มหาโควินฺทสุตฺต}
\chapterheadpageread{6. มหาโควินฺทสุตฺต}
\proseitem{293}{\csromanfolio{178}เอวํ เม สุตํ\csromandash{}เอกํ สมยํ ภควา ราชคเห วิหรติ คิชฺฌกูเฏ ปพฺพเต.\csromanspacer{} อถ โข ปญฺจสิโข คนฺธพฺพปุตฺโต อภิกฺกนฺตาย รตฺติยา อภิกฺกนฺตวณฺโณ เกวลกปฺปํ คิชฺฌกูฏํ ปพฺพตํ โอภาเสตฺวา เยน ภควา เตนุปสงฺกมิ, อุปสงฺกมิตฺวา ภควนฺตํ อภิวาเทตฺวา เอกมนฺตํ อฏฺฐาสิ, เอกมนฺตํ ฐิโต โข ปญฺจสิโข คนฺธพฺพปุตฺโต ภควนฺตํ เอตทโวจ “ยํ โข เม ภนฺเต เทวานํ ตาวติํสานํ สมฺมุขา สุตํ สมฺมุขา ปฏิคฺคหิตํ, อาโรเจมิ ตํ ภควโต”ติ. “อาโรเจหิ เม ตฺวํ ปญฺจสิขา”ติ ภควา อโวจ.}

\csromanmatikaanchor{mtk.93}
\csromantocmarkref{section}{เทวสภา}{mtk.93}
\bookmark[dest={mtk.93},level=1]{เทวสภา}
\csromanheader{1}{\textbf{เทวสภา}}
\proseitem{294}{ปุริมานิ ภนฺเต ทิวสานิ ปุริมตรานิ ตทหุโปสเถ ปนฺนรเส ปวารณาย ปุณฺณาย ปุณฺณมาย รตฺติยา เกวลกปฺปา จ เทวา ตาวติํสา สุธมฺมายํ สภายํ สนฺนิสินฺนา โหนฺติ สนฺนิปติตา, มหตี จ ทิพฺพปริสา สมนฺตโต นิสินฺนา โหนฺติ, จตฺตาโร จ มหาราชาโน จตุทฺทิสา นิสินฺนา โหนฺติ, ปุรตฺถิมาย ทิสาย ธตรฏฺโฐ มหาราชา ปจฺฉิมาภิมุโข นิสินฺโน โหติ เทเว ปุรกฺขตฺวา, ทกฺขิณาย ทิสาย วิรูฬฺหโก มหาราชา อุตฺตราภิมุโข นิสินฺโน โหติ เทเว ปุรกฺขตฺวา, ปจฺฉิมาย ทิสาย วิรูปกฺโข มหาราชา ปุรตฺถาภิมุโข นิสินฺโน โหติ เทเว ปุรกฺขตฺวา, อุตฺตราย ทิสาย เวสฺสวโณ มหาราชา ทกฺขิณาภิมุโข นิสินฺโน โหติ เทเว ปุรกฺขตฺวา.\csromanspacer{} ยทา ภนฺเต เกวลกปฺปา จ เทวา ตาวติํสา สุธมฺมายํ สภายํ สนฺนิสินฺนา โหนฺติ สนฺนิปติตา, มหตี จ ทิพฺพปริสา สมนฺตโต นิสินฺนา โหนฺติ, จตฺตาโร จ มหาราชาโน จตุทฺทิสา นิสินฺนา โหนฺติ, อิทํ เนสํ โหติ อาสนสฺมิํ, อถ ปจฺฉา อมฺหากํ อาสนํ โหติ.}

\prose{เย เต ภนฺเต เทวา ภควติ พฺรหฺมจริยํ จริตฺวา อธุนูปปนฺนา ตาวติํสกายํ.\csromanspacer{} เต อญฺเญ เทเว อติโรจนฺติ วณฺเณน เจว ยสสา จ.\csromanspacer{} เตน สุทํ ภนฺเต เทวา ตาวติํสา อตฺตมนา โหนฺติ}

\vspace{\tipitakaparskip}
\csromancenter{มหาโควินฺทสุตฺต ปมุทิตา ปีติโสมนสฺสชาตา “ทิพฺพา วต โภ กายา ปริปูเรนฺติ, หายนฺติ อสุรกายา”ติ.}
\proseitem{295}{\csromanfolio{179}อถ โข ภนฺเต สกฺโก เทวานมินฺโท เทวานํ ตาวติํสานํ สมฺปสาทํ วิทิตฺวา อิมาหิ คาถาหิ อนุโมทิ\csromandash{}}

\csromangathameasure{“โมทนฺติ วต โภ เทวา, ตาวติํสา สหินฺทกา. \\ ตถาคตํ นมสฺสนฺตา, ธมฺมสฺส จ สุธมฺมตํ. \\ นเว เทเว จ ปสฺสนฺตา, วณฺณวนฺเต ยสสฺสิเน. \\ สุคตสฺมิํ พฺรหฺมจริยํ, จริตฺวาน อิธาคเต. \\ เต อญฺเญ อติโรจนฺติ, วณฺเณน ยสสายุนา. \\ สาวกา ภูริปญฺญสฺส, วิเสสูปคตา อิธ. \\ อิทํ ทิสฺวาน นนฺทนฺติ, ตาวติํสา สหินฺทกา. \\ ตถาคตํ นมสฺสนฺตา, ธมฺมสฺส จ สุธมฺมตนฺ”ติ.}
\csromangathagroup{\csromangathastanza{“โมทนฺติ วต โภ เทวา, ตาวติํสา สหินฺทกา. \\ ตถาคตํ นมสฺสนฺตา, ธมฺมสฺส จ สุธมฺมตํ.}\csromangathastanza{นเว เทเว จ ปสฺสนฺตา, วณฺณวนฺเต ยสสฺสิเน. \\ สุคตสฺมิํ พฺรหฺมจริยํ, จริตฺวาน อิธาคเต.}\csromangathastanza{เต อญฺเญ อติโรจนฺติ, วณฺเณน ยสสายุนา. \\ สาวกา ภูริปญฺญสฺส, วิเสสูปคตา อิธ.}\csromangathastanza{อิทํ ทิสฺวาน นนฺทนฺติ, ตาวติํสา สหินฺทกา. \\ ตถาคตํ นมสฺสนฺตา, ธมฺมสฺส จ สุธมฺมตนฺ”ติ.}}

\prose{เตน สุทํ ภนฺเต เทวา ตวติํสา ภิยฺโยโส มตฺตาย อตฺตมนา โหนฺติ ปมุทิตา ปีติโสมนสฺสชาตา “ทิพฺพา วต โภ กายา ปริปูเรนฺติ, หายนฺติ อสุรกายา”ติ.}

\csromanmatikaanchor{mtk.94}
\csromantocmarkref{section}{อฏฺฐยถาภุจฺจวณฺณ}{mtk.94}
\bookmark[dest={mtk.94},level=1]{อฏฺฐยถาภุจฺจวณฺณ}
\csromanheader{1}{\textbf{อฏฺฐยถาภุจฺจวณฺณ}}
\proseitem{296}{อถ โข ภนฺเต สกฺโก เทวานมินฺโท เทวานํ ตาวติํสานํ สมฺปสาทํ วิทิตฺวา เทเว ตาวติํเส อามนฺเตสิ “อิจฺเฉยฺยาถ โน ตุเมฺห มาริสา ตสฺส ภควโต อฏฺฐ ยถาภุจฺเจ วณฺเณ โสตุนฺ”ติ.\csromanspacer{} อิจฺฉาม มยํ มาริส ตสฺส ภควโต อฏฺฐ ยถาภุจฺเจ วณฺเณ โสตุนฺติ.\csromanspacer{} อถ โข ภนฺเต สกฺโก เทวานมินฺโท เทวานํ ตาวติํสานํ ภควโต อฏฺฐ ยถาภุจฺเจ วณฺเณ ปยิรุทาหาสิ “ตํ กิํมญฺญนฺติ โภนฺโต เทวา ตาวติํสา, ยาวญฺจ โส ภควา พหุชนหิตาย ปฏิปนฺโน พหุชนสุขาย โลกานุกมฺปาย อตฺถาย หิตาย สุขาย เทวมนุสฺสานํ, เอวํ พหุชนหิตาย ปฏิปนฺนํ พหุชนสุขาย โลกานุกมฺปาย อตฺถาย หิตาย สุขาย เทวมนุสฺสานํ อิมินาปงฺเคน สมนฺนาคตํ สตฺถารํ เนว อตีตํเส สมนุปสฺสาม, น ปเนตรหิ อญฺญตฺร เตน ภควตา. (1)}

\prose{\csromanfolio{180}สฺวากฺขาโต โข ปน เตน ภควตา ธมฺโม สนฺทิฏฺฐิโก อกาลิโก เอหิปสฺสิโก โอปเนยฺยิโก ปจฺจตฺตํ เวทิตพฺโพ วิญฺญูหิ, เอวํ โอปเนยฺยิกสฺส ธมฺมสฺส เทเสตารํ อิมินาปงฺเคน สมนฺนาคตํ สตฺถารํ เนว อตีตํเส สมนุปสฺสาม, น ปเนตรหิ อญฺญตฺร เตน ภควตา. (2)}

\prose{“อิทํ กุสลนฺ”ติ โข ปน เตน ภควตา สุปญฺญตฺตํ, “อิทํ อกุสลนฺ”ติ สุปญฺญตฺตํ. “อิทํ สาวชฺชํ อิทํ อนวชฺชํ, อิทํ เสวิตพฺพํ อิทํ น เสวิตพฺพํ, อิทํ หีนํ อิทํ ปณีตํ, อิทํ กณฺหสุกฺกสปฺปฏิภาคนฺ”ติ สุปญฺญตฺตํ.\csromanspacer{} เอวํ กุสลากุสล สาวชฺชานวชฺช เสวิตพฺพาเสวิตพฺพ หีนปณีต กณฺหสุกฺกสปฺปฏิภาคานํ ธมฺมานํ ปญฺญาเปตารํ อิมินาปงฺเคน สมนฺนาคตํ สตฺถารํ เนว อตีตํเส สมนุปสฺสาม, น ปเนตรหิ อญฺญตฺร เตน ภควตา. (3)}

\prose{สุปญฺญตฺตา โข ปน เตน ภควตา สาวกานํ นิพฺพานคามินี ปฏิปทา สํสนฺทติ นิพฺพานญฺจ ปฏิปทา จ.\csromanspacer{} เสยฺยถาปิ นาม คงฺโคทกํ ยมุโนทเกน สํสนฺทติ สเมติ, เอวเมว สุปญฺญตฺตา เตน ภควตา สาวกานํ นิพฺพานคามินี ปฏิปทา สํสนฺทติ นิพฺพานญฺจ ปฏิปทา จ.\csromanspacer{} เอวํ นิพฺพานคามินิยา ปฏิปทาย ปญฺญาเปตารํ อิมินาปงฺเคน สมนฺนาคตํ สตฺถารํ เนว อตีตํเส สมนุปสฺสาม, น ปเนตรหิ อญฺญตฺร เตน ภควตา. (4)}

\prose{อภินิปฺผนฺโน\footnote{อภินิปฺปนฺโน (อิ, ก)} โข ปน ตสฺส ภควโต ลาโภ อภินิปฺผนฺโน สิโลโก, ยาว มญฺเญ ขตฺติยา สํปิยายมานรูปา วิหรนฺติ, วิคตมโท โข ปน โส ภควา อาหารํ อาหาเรติ.\csromanspacer{} เอวํ วิคตมทํ อาหารํ อาหรยมานํ อิมินาปงฺเคน สมนฺนาคตํ สตฺถารํ เนว อตีตํเส สมนุปสฺสาม, น ปเนตรหิ อญฺญตฺร เตน ภควตา. (5)}

\prose{ลทฺธสหาโย โข ปน โส ภควา เสขานญฺเจว ปฏิปนฺนานํ ขีณาสวานญฺจ วุสิตวตํ, เต ภควา อปนุชฺช เอการามตํ อนุยุตฺโต วิหรติ.\csromanspacer{} เอวํ เอการามตํ อนุยุตฺตํ อิมินาปงฺเคน สมนฺนาคตํ สตฺถารํ เนว อตีตํเส สมนุปสฺสาม, น ปเนตรหิ อญฺญตฺร เตน ภควตา. (6)}

\prose{ยถาวาที โข ปน โส ภควา ตถาการี, ยถาการี ตถาวาที, อิติ ยถาวาที ตถาการี, ยถาการี ตถาวาที.\csromanspacer{} เอวํ ธมฺมานุธมฺมปฺปฏิปนฺนํ อิมินาปงฺเคน สมนฺนาคตํ สตฺถารํ เนว อตีตํเส สมนุปสฺสาม, น ปเนตรหิ อญฺญตฺร เตน ภควตา. (7)}

\csromanheader{2}{6. มหาโควินฺทสุตฺต}
\prose{\csromanfolio{181}ติณฺณวิจิกิจฺโฉ โข ปน โส ภควา วิคตกถํกโถ ปริโยสิตสงฺกปฺโป อชฺฌาสยํ อาทิพฺรหฺมจริยํ.\csromanspacer{} เอวํ ติณฺณวิจิกิจฺฉํ วิคตกถํกถํ ปริโยสิตสงฺกปฺปํ อชฺฌาสยํ อาทิพฺรหฺมจริยํ อิมินาปงฺเคน สมนฺนาคตํ สตฺถารํ เนว อตีตํเส สมนุปสฺสาม, น ปเนตรหิ อญฺญตฺร เตน ภควตา”ติ. (8)}

\proseitem{297}{อิเม โข ภนฺเต สกฺโก เทวานมินฺโท เทวานํ ตาวติํสานํ ภควโต อฏฺฐ ยถาภุจฺเจ วณฺเณ ปยิรุทาหาสิ.\csromanspacer{} เตน สุทํ ภนฺเต เทวา ตาวติํสา ภิยฺโยโส มตฺตาย อตฺตมนา โหนฺติ ปมุทิตา ปีติโสมนสฺสชาตา ภควโต อฏฺฐ ยถาภุจฺเจ วณฺเณ สุตฺวา.\csromanspacer{} ตตฺร ภนฺเต เอกจฺเจ เทวา เอวมาหํสุ “อโห วต มาริสา จตฺตาโร สมฺมาสมฺพุทฺธา โลเก อุปฺปชฺเชยฺยุํ ธมฺมญฺจ เทเสยฺยุํ ยถริว ภควา.\csromanspacer{} ตทสฺส พหุชนหิตาย พหุชนสุขาย โลกานุกมฺปาย อตฺถาย หิตาย สุขาย เทวมนุสฺสานนฺ”ติ.\csromanspacer{} เอกจฺเจ เทวา เอวมาหํสุ “ติฏฺฐนฺตุมาริสา จตฺตาโร สมฺมาสมฺพุทฺธา, อโห วต มาริสา ตโย สมฺมาสมฺพุทฺธา โลเก อุปฺปชฺเชยฺยุํ ธมฺมญฺจ เทเสยฺยุํ ยถริว ภควา.\csromanspacer{} ตทสฺส พหุชนหิตาย พหุชนสุขาย โลกานุกมฺปาย อตฺถาย หิตาย สุขาย เทวมนุสฺสานนฺ”ติ.\csromanspacer{} เอกจฺเจ เทวา เอวมาหํสุ “ติฏฺฐนฺตุ มาริสา ตโย สมฺมาสมฺพุทฺธา, อโห วต มาริสา เทฺว สมฺมาสมฺพุทฺธา โลเก อุปฺปชฺเชยฺยุํ ธมฺมญฺจ เทเสยฺยุํ ยถริว ภควา.\csromanspacer{} ตทสฺส พหุชนหิตาย พหุชนสุขาย โลกานุกมฺปาย อตฺถาย หิตาย สุขาย เทวมนุสฺสานนฺ”ติ.}

\proseitem{298}{เอวํ วุตฺเต ภนฺเต สกฺโก เทวานมินฺโท เทเว ตาวติํเส เอตทโวจ “อฏฺฐานํ โข เอตํ มาริสา อนวกาโส, ยํ เอกิสฺสา โลกธาตุยา เทฺว อรหนฺโต สมฺมาสมฺพุทฺธา อปุพฺพํ อจริมํ อุปฺปชฺเชยฺยุํ, เนตํ ฐานํ วิชฺชติ.\csromanspacer{} อโห วต มาริสา โส ภควา อปฺปาพาโธ อปฺปาตงฺโก จิรํ ทีฆมทฺธานํ ติฏฺเฐยฺย.\csromanspacer{} ตทสฺส พหุชนหิตาย พหุชนสุขาย โลกานุกมฺปาย อตฺถาย หิตาย สุขาย เทวมนุสฺสานนฺ”ติ.\csromanspacer{} อถ โข ภนฺเต เยนตฺเถน เทวา ตาวติํสา สุธมฺมายํ สภายํ สนฺนิสินฺนา โหนฺติ สนฺนิปติตา, ตํ อตฺถํ จินฺตยิตฺวา ตํ อตฺถํ มนฺตยิตฺวา วุตฺตวจนาปิ ตํ จตฺตาโร มหาราชาโน ตสฺมิํ อตฺเถ โหนฺติ, \csromanfolio{182}ปจฺจานุสิฏฺฐวจนาปิ ตํ จตฺตาโร มหาราชาโน ตสฺมิํ อตฺเถ โหนฺติ, สเกสุ สเกสุ อาสเนสุ ฐิตา อวิปกฺกนฺตา.}

\csromangathameasure{เต วุตฺตวากฺยา ราชาโน, ปฏิคฺคยฺหานุสาสนิํ. \\ วิปฺปสนฺนมนา สนฺตา, อฏฺฐํสุ สมฺหิ อาสเนติ.}
\csromangathagroup{\csromangathastanza{เต วุตฺตวากฺยา ราชาโน, ปฏิคฺคยฺหานุสาสนิํ. \\ วิปฺปสนฺนมนา สนฺตา, อฏฺฐํสุ สมฺหิ อาสเนติ.}}

\proseitem{299}{อถ โข ภนฺเต อุตฺตราย ทิสาย อุฬาโร อาโลโก สญฺชายิ, โอภาโส ปาตุรโหสิ อติกฺกมฺเมว เทวานํ เทวานุภาวํ.\csromanspacer{} อถ โข ภนฺเต สกฺโก เทวานมินฺโท เทเว ตาวติํเส อามนฺเตสิ “ยถา โข มาริสา นิมิตฺตานิ ทิสฺสนฺติ, อุฬาโร อาโลโก สญฺชายติ, โอภาโส ปาตุ ภวติ, พฺรหฺมา ปาตุ ภวิสฺสติ, พฺรหฺมุโน เหตํ ปุพฺพนิมิตฺตํ ปาตุภาวาย, ยทิทํ อาโลโก สญฺชายติ โอภาโส ปาตุภวตี”ติ.}

\csromangathameasure{ยถา นิมิตฺตา ทิสฺสนฺติ, พฺรหฺมา ปาตุ ภวิสฺสติ. \\ พฺรหฺมุโน เหตํ นิมิตฺตํ, โอภาโส วิปุโล มหาติ.}
\csromangathagroup{\csromangathastanza{ยถา นิมิตฺตา ทิสฺสนฺติ, พฺรหฺมา ปาตุ ภวิสฺสติ. \\ พฺรหฺมุโน เหตํ นิมิตฺตํ, โอภาโส วิปุโล มหาติ.}}

\csromanmatikaanchor{mtk.95}
\csromantocmarkref{section}{สนงฺกุมารกถา}{mtk.95}
\bookmark[dest={mtk.95},level=1]{สนงฺกุมารกถา}
\csromanheader{1}{\textbf{สนงฺกุมารกถา}}
\proseitem{300}{อถ โข ภนฺเต เทวา ตาวติํสา ยถาสเกสุ อาสเนสุ นิสีทิํสุ “โอภาสเมตํ ญสฺสาม, ยํวิปาโก ภวิสฺสติ, สจฺฉิกตฺวาว นํ คมิสฺสามา”ติ.\csromanspacer{} จตฺตาโรปิ มหาราชาโน ยถาสเกสุ อาสเนสุ นิสีทิํสุ “โอภาสเมตํ ญสฺสาม, ยํวิปาโก ภวิสฺสติ, สจฺฉิกตฺวาว นํ คมิสฺสามา”ติ.\csromanspacer{} อิทํ สุตฺวา เทวา ตาวติํสา เอกคฺคา สมาปชฺชิํสุ “โอภาสเมตํ ญสฺสาม, ยํวิปาโก ภวิสฺสติ, สจฺฉิกตฺวาว นํ คมิสฺสามา”ติ.}

\prose{ยทา ภนฺเต พฺรหฺมา สนงฺกุมาโร เทวานํ ตาวติํสานํ ปาตุ ภวติ.\csromanspacer{} โอฬาริกํ อตฺตภาวํ อภินิมฺมินิตฺวา ปาตุ ภวติ.\csromanspacer{} โย โข ปน ภนฺเต พฺรหฺมุโน ปกติวณฺโณ, อนภิสมฺภวนีโย โส เทวานํ ตาวติํสานํ จกฺขุปถสฺมิํ.\csromanspacer{} ยทา ภนฺเต พฺรหฺมา สนงฺกุมาโร เทวานํ ตาวติํสานํ ปาตุ ภวติ, โส อญฺเญ เทเว อติโรจติ วณฺเณน เจว ยสสา จ.\csromanspacer{} เสยฺยถาปิ ภนฺเต โส วณฺโณ วิคฺคโห มานุสํ วิคฺคหํ อติโรจติ, เอวเมว โข ภนฺเต ยทา พฺรหฺมา สนงฺกุมาโร}

\vspace{\tipitakaparskip}
\csromancenter{มหาโควินฺทสุตฺต เทวานํ ตาวติํสานํ ปาตุ ภวติ, โส อญฺเญ เทเว อติโรจติ วณฺเณน เจว ยสสา จ.\csromanspacer{} ยทา ภนฺเต พฺรหฺมา สนงฺกุมาโร เทวานํ ตาวติํสานํ ปาตุ ภวติ, น ตสฺสํ ปริสายํ โกจิ เทโว อภิวาเทติ วา ปจฺจุฏฺเฐติ วา อาสเนน วา นิมนฺเตติ, สพฺเพว ตุณฺหีภูตา ปญฺชลิกา ปลฺลงฺเกน นิสีทนฺติ “ยสฺสทานิ เทวสฺส ปลฺลงฺกํ อิจฺฉิสฺสติ พฺรหฺมา สนงฺกุมาโร, ตสฺส เทวสฺส ปลฺลงฺเก นิสีทิสฺสตี”ติ.\csromanspacer{} ยสฺส โข ปน ภนฺเต เทวสฺส พฺรหฺมา สนงฺกุมาโร ปลฺลงฺเก นิสีทติ, อุฬารํ โส ลภติ เทโว เวทปฏิลาภํ, อุฬารํ โส ลภติ เทโว โสมนสฺสปฏิลาภํ.\csromanspacer{} เสยฺยถาปิ ภนฺเต ราชา ขตฺติโย มุทฺธาวสิตฺโต อธุนาภิสิตฺโต รชฺเชน, อุฬารํ โส ลภติ เวทปฏิลาภํ, อุฬารํ โส ลภติ โสมนสฺสปฏิลาภํ.\csromanspacer{} เอวเมว โข ภนฺเต ยสฺส เทวสฺส พฺรหฺมา สนงฺกุมาโร ปลฺลงฺเก นิสีทติ, อุฬารํ โส ลภติ เทโว เวทปฏิลาภํ, อุฬารํ โส ลภติ เทโว โสมนสฺสปฏิลาภํ.\csromanspacer{} อถ ภนฺเต พฺรหฺมา สนงฺกุมาโร เทวานํ ตาวติํสานํ สมฺปสาทํ วิทิตฺวา อนฺตรหิโต อิมาหิ คาถาหิ อนุโมทิ\csromandash{}}
\vspace{0.5\baselineskip}
\csromangathameasure{“โมทนฺติ วต โภ เทวา, ตาวติํสา สหินฺทกา. \\ ตถาคตํ นมสฺสนฺตา, ธมฺมสฺส จ สุธมฺมตํ. \\ นเว เทเว จ ปสฺสนฺตา, วณฺณวนฺเต ยสสฺสิเน. \\ สุคตสฺมิํ พฺรหฺมจริยํ, จริตฺวาน อิธาคเต. \\ เต อญฺเญ อติโรจนฺติ, วณฺเณน ยสสายุนา. \\ สาวกา ภูริปญฺญสฺส, วิเสสูปคตา อิธ. \\ อิทํ ทิสฺวาน นนฺทนฺติ, ตาวติํสา สหินฺทกา. \\ ตถาคตํ นมสฺสนฺตา, ธมฺมสฺส จ สุธมฺมตนฺ”ติ.}
\csromangathagroup{\csromangathastanza{\csromanfolio{183}“โมทนฺติ วต โภ เทวา, ตาวติํสา สหินฺทกา. \\ ตถาคตํ นมสฺสนฺตา, ธมฺมสฺส จ สุธมฺมตํ.}\csromangathastanza{นเว เทเว จ ปสฺสนฺตา, วณฺณวนฺเต ยสสฺสิเน. \\ สุคตสฺมิํ พฺรหฺมจริยํ, จริตฺวาน อิธาคเต.}\csromangathastanza{เต อญฺเญ อติโรจนฺติ, วณฺเณน ยสสายุนา. \\ สาวกา ภูริปญฺญสฺส, วิเสสูปคตา อิธ.}\csromangathastanza{อิทํ ทิสฺวาน นนฺทนฺติ, ตาวติํสา สหินฺทกา. \\ ตถาคตํ นมสฺสนฺตา, ธมฺมสฺส จ สุธมฺมตนฺ”ติ.}}

\proseitem{301}{อิมมตฺถํ ภนฺเต พฺรหฺมา สนงฺกุมาโร อภาสิตฺถ.\csromanspacer{} อิมมตฺถํ ภนฺเต พฺรหฺมุโน สนงฺกุมารสฺส ภาสโต อฏฺฐงฺคสมนฺนาคโต สโร โหติ วิสฺสฏฺโฐ จ วิญฺเญยฺโย จ มญฺชุ จ สวนีโย จ พินฺทุ จ อวิสารี จ คมฺภีโร จ นินฺนาที จ.\csromanspacer{} ยถาปริสํ โข ปน ภนฺเต พฺรหฺมา สนงฺกุมาโร สเรน วิญฺญาเปติ, น จสฺส พหิทฺธา ปริสาย โฆโส นิจฺฉรติ, ยสฺส โข ปน ภนฺเต เอวํ อฏฺฐงฺคสมนฺนาคโต สโร โหติ, โส วุจฺจติ พฺรหฺมสฺสโรติ.\csromanspacer{} อถ โข ภนฺเต เทวา ตาวติํสา พฺรหฺมานํ สนงฺกุมารํ \csromanfolio{184}เอตทโวจุํ “สาธุ มหาพฺรหฺเม เอตเทว มยํ สงฺขาย โมทาม, อตฺถิ จ สกฺเกน เทวานมินฺเทน ตสฺส ภควโต อฏฺฐ ยถาภุจฺจา วณฺณา ภาสิตา, เต จ มยํ สงฺขาย โมทามา”ติ.}

\csromanmatikaanchor{mtk.96}
\csromantocmarkref{section}{อฏฺฐยถาภุจฺจวณฺณ}{mtk.96}
\bookmark[dest={mtk.96},level=1]{อฏฺฐยถาภุจฺจวณฺณ}
\csromanheader{1}{\textbf{อฏฺฐยถาภุจฺจวณฺณ}}
\proseitem{302}{อถ ภนฺเต พฺรหฺมา สนงฺกุมาโร สกฺกํ เทวานมินฺทํ เอตทโวจ “สาธุ เทวานมินฺท มยมฺปิ ตสฺส ภควโต อฏฺฐ ยถาภุจฺเจ วณฺเณ สุเณยฺยามา”ติ. “เอวํ มหาพฺรหฺเม”ติ โข ภนฺเต สกฺโก เทวานมินฺโท พฺรหฺมุโน สนงฺกุมารสฺส ภควโต อฏฺฐ ยถาภุจฺเจ วณฺเณ ปยิรุทาหาสิ.}

\prose{ตํ กิํมญฺญติ ภวํ มหาพฺรหฺมา, ยาวญฺจ โส ภควา พหุชนหิตาย ปฏิปนฺโน พหุชนสุขาย โลกานุกมฺปาย อตฺถาย หิตาย สุขาย เทวมนุสฺสานํ, เอวํ พหุชนหิตาย ปฏิปนฺนํ พหุชนสุขาย โลกานุกมฺปาย อตฺถาย หิตาย สุขาย เทวมนุสฺสานํ อิมินาปงฺเคน สมนฺนาคตํ สตฺถารํ เนว อตีตํเส สมนุปสฺสาม, น ปเนตรหิ อญฺญตฺร เตน ภควตา. (1)}

\prose{สฺวากฺขาโต โข ปน เตน ภควตา ธมฺโม สนฺทิฏฺฐิโก อกาลิโก เอหิปสฺสิโก โอปเนยฺยิโก ปจฺจตฺตํ เวทิตพฺโพ วิญฺญูหิ.\csromanspacer{} เอวํ โอปเนยฺยิกสฺส ธมฺมสฺส เทเสตารํ อิมินาปงฺเคน สมนฺนาคตํ สตฺถารํ เนว อตีตํเส สมนุปสฺสาม, น ปเนตรหิ อญฺญตฺร เตน ภควตา. (2)}

\prose{“อิทํ กุสลนฺ”ติ โข ปน เตน ภควตา สุปญฺญตฺตํ, “อิทํ อกุสลนฺ”ติ สุปญฺญตฺตํ, “อิทํ สาวชฺชํ อิทํ อนวชฺชํ, อิทํ เสวิตพฺพํ อิทํ น เสวิตพฺพํ, อิทํ หีนํ อิทํ ปณีตํ, อิทํ กณฺหสุกฺกสปฺปฏิภาคนฺ”ติ สุปญฺญตฺตํ.\csromanspacer{} เอวํ กุสลากุสล สาวชฺชานวชฺช เสวิตพฺพาเสวิตพฺพ หีนปณีต กณฺหสุกฺกสปฺปฏิภาคานํ ธมฺมานํ ปญฺญาเปตารํ อิมินาปงฺเคน สมนฺนาคตํ สตฺถารํ เนว อตีตํเส สมนุปสฺสาม, น ปเนตรหิ อญฺญตฺร เตน ภควตา. (3)}

\prose{สุปญฺญตฺตา โข ปน เตน ภควตา สาวกานํ นิพฺพานคามินี ปฏิปทา สํสนฺทติ นิพฺพานญฺจ ปฏิปทา จ.\csromanspacer{} เสยฺยถาปิ นาม คงฺโคทกํ ยมุโนทเกน}

\vspace{\tipitakaparskip}
\csromancenter{มหาโควินฺทสุตฺต สํสนฺทติ สเมติ, เอวเมว สุปญฺญตฺตา เตน ภควตา สาวกานํ นิพฺพานคามินี ปฏิปทา สํสนฺทติ นิพฺพานญฺจ ปฏิปทา จ.\csromanspacer{} เอวํ นิพฺพานคามินิยา ปฏิปทาย ปญฺญาเปตารํ อิมินาปงฺเคน สมนฺนาคตํ สตฺถารํ เนว อตีตํเส สมนุปสฺสาม, น ปเนตรหิ อญฺญตฺร เตน ภควตา. (4)}
\prose{\csromanfolio{185}อภินิปฺผนฺโน โข ปน ตสฺส ภควโต ลาโภ อภินิปฺผนฺโน สิโลโก, ยาว มญฺเญ ขตฺติยา สํปิยายมานรูปา วิหรนฺติ.\csromanspacer{} วิคตมโท โข ปน โส ภควา อาหารํ อาหาเรติ.\csromanspacer{} เอวํ วิคตมทํ อาหารํ อาหรยมานํ อิมินาปงฺเคน สมนฺนาคตํ สตฺถารํ เนว อตีตํเส สมนุปสฺสาม, น ปเนตรหิ อญฺญตฺร เตน ภควตา. (5)}

\prose{ลทฺธสหาโย โข ปน โส ภควา เสขานญฺเจว ปฏิปนฺนานํ ขีณาสวานญฺจ วุสิตวตํ, เต ภควา อปนุชฺช เอการามตํ อนุยุตฺโต วิหรติ.\csromanspacer{} เอวํ เอการามตํ อนุยุตฺตํ อิมินาปงฺเคน สมนฺนาคตํ สตฺถารํ เนว อตีตํเส สมนุปสฺสาม, น ปเนตรหิ อญฺญตฺร เตน ภควตา. (6)}

\prose{ยถาวาที โข ปน โส ภควา ตถาการี, ยถาการี ตถาวาที, อิติ ยถาวาที ตถาการี, ยถาการี ตถาวาที, เอวํ ธมฺมานุธมฺมปฺปฏิปฺปนฺนํ อิมินาปงฺเคน สมนฺนาคตํ สตฺถารํ เนว อตีตํเส สมนุปสฺสาม, น ปเนตรหิ อญฺญตฺร เตน ภควตา. (7)}

\prose{ติณฺณวิจิกิจฺโฉ โข ปน โส ภควา วิคตกถํกโถ ปริโยสิตสงฺกปฺโป อชฺฌาสยํ อาทิพฺรหฺมจริยํ.\csromanspacer{} เอวํ ติณฺณวิจิกิจฺฉํ วิคตกถํกถํ ปริโยสิตสงฺกปฺปํ อชฺฌาสยํ อาทิพฺรหฺมจริยํ อิมินาปงฺเคน สมนฺนาคตํ สตฺถารํ เนว อตีตํเส สมนุปสฺสาม, น ปเนตรหิ อญฺญตฺร เตน ภควตาติ. (8)}

\proseitem{303}{อิเม โข ภนฺเต สกฺโก เทวานมินฺโท พฺรหฺมุโน สนงฺกุมารสฺส ภควโต อฏฺฐ ยถาภุจฺเจ วณฺเณ ปยิรุทาหาสิ.\csromanspacer{} เตน สุทํ ภนฺเต พฺรหฺมา สนงฺกุมาโร อตฺตมโน โหติ ปมุทิโต ปีติโสมนสฺสชาโต ภควโต อฏฺฐ ยถาภุจฺเจ วณฺเณ สุตฺวา.\csromanspacer{} อถ ภนฺเต พฺรหฺมา สนงฺกุมาโร โอฬาริกํ อตฺตภาวํ อภินิมฺมินิตฺวา กุมารวณฺณี หุตฺวา ปญฺจสิโข เทวานํ ตาวติํสานํ ปาตุรโหสิ.\csromanspacer{} โส เวหาสํ อพฺภุคฺคนฺตฺวา อากาเส อนฺตลิกฺเข ปลฺลงฺเกน นิสีทิ.\csromanspacer{} เสยฺยถาปิ ภนฺเต \csromanfolio{186}พลวา ปุริโส สุปจฺจตฺถเต วา ปลฺลงฺเก สเม วา ภูมิภาเค ปลฺลงฺเกน นิสีเทยฺย, เอวเมว โข ภนฺเต พฺรหฺมา สนงฺกุมาโร เวหาสํ อพฺภุคฺคนฺตฺวา อากาเส อนฺตลิกฺเข ปลฺลงฺเกน นิสีทิตฺวา เทเว ตาวติํเส อามนฺเตสิ\csromandash{}}

\csromanmatikaanchor{mtk.97}
\csromantocmarkref{section}{โควินฺทพฺราหฺมณวตฺถุ}{mtk.97}
\bookmark[dest={mtk.97},level=1]{โควินฺทพฺราหฺมณวตฺถุ}
\csromanheader{1}{\textbf{โควินฺทพฺราหฺมณวตฺถุ}}
\proseitem{304}{ตํ กิํมญฺญนฺติ โภนฺโต เทวา ตาวติํสา, ยาว ทีฆรตฺตํ มหาปญฺโญว โส ภควา อโหสิ.}

\prose{ภูตปุพฺพํ โภ ราชา ทิสมฺปติ นาม อโหสิ, ทิสมฺปติสฺส รญฺโญ โควินฺโท นาม พฺราหฺมโณ ปุโรหิโต อโหสิ, ทิสมฺปติสฺส รญฺโญ เรณุ นาม กุมาโร ปุตฺโต อโหสิ, โควินฺทสฺส พฺราหฺมณสฺส โชติปาโล นาม มาณโว ปุตฺโต อโหสิ, อิติ เรณุ จ ราชปุตฺโต โชติปาโล จ มาณโว อญฺเญ จ ฉ ขตฺติยา อิจฺเจเต อฏฺฐ สหายา อเหสุํ.\csromanspacer{} อถ โข โภ อโหรตฺตานํ อจฺจเยน โควินฺโท พฺราหฺมโณ กาลมกาสิ.\csromanspacer{} โควินฺเท พฺราหฺมเณ กาลงฺกเต ราชา ทิสมฺปติ ปริเทเวสิ “ยสฺมิํ วต โภ มยํ สมเย โควินฺเท พฺราหฺมเณ สพฺพกิจฺจานิ สมฺมา โวสฺสชฺชิตฺวา ปญฺจหิ กามคุเณหิ สมปฺปิตา สมงฺคีภูตา ปริจาเรม, ตสฺมิํ โน สมเย โควินฺโท พฺราหฺมโณ กาลงฺกโต”ติ.\csromanspacer{} เอวํ วุตฺเต โภ เรณุ ราชปุตฺโต ราชานํ ทิสมฺปติํ เอตทโวจ “มา โข ตฺวํ เทว โควินฺเท พฺราหฺมเณ กาลงฺกเต อติพาฬฺหํ ปริเทเวสิ, อตฺถิ เทว โควินฺทสฺส พฺราหฺมณสฺส โชติปาโล นาม มาณโว ปุตฺโต ปณฺฑิตตโร เจว ปิตรา, อลมตฺถทสตโร เจว ปิตรา, เยปิสฺส ปิตา อตฺเถ อนุสาสิ, เตปิ โชติปาลสฺเสว มาณวสฺส อนุสาสนิยา”ติ.\csromanspacer{} เอวํ กุมาราติ.\csromanspacer{} เอวํ เทวาติ.}

\csromanmatikaanchor{mtk.98}
\csromantocmarkref{section}{มหาโควินฺทวตฺถุ}{mtk.98}
\bookmark[dest={mtk.98},level=1]{มหาโควินฺทวตฺถุ}
\csromanheader{1}{\textbf{มหาโควินฺทวตฺถุ}}
\proseitem{305}{อถ โข โภ ราชา ทิสมฺปติ อญฺญตรํ ปุริสํ อามนฺเตสิ “เอหิ ตฺวํ อมฺโภ ปุริส เยน โชติปาโล มาณโว เตนุปสงฺกม, อุปสงฺกมิตฺวา โชติปาลํ มาณวํ เอวํ วเทหิ ‘ภวมตฺถุ ภวนฺตํ โชติปาลํ, ราชา ทิสมฺปติ ภวนฺตํ โชติปาลํ มาณวํ อามนฺตยติ, ราชา ทิสมฺปติ โภโต โชติปาลสฺส มาณวสฺส ทสฺสนกาโม’ติ”. “เอวํ เทวา” ติ โข โภ โส ปุริโส ทิสมฺปติสฺส รญฺโญ ปฏิสฺสุตฺวา}

\vspace{\tipitakaparskip}
\csromancenter{มหาโควินฺทสุตฺต เยน โชติปาโล มาณโว เตนุปสงฺกมิ, อุปสงฺกมิตฺวา โชติปาลํ มาณวํ เอตทโวจ “ภวมตฺถุ ภวนฺตํ โชติปาลํ, ราชา ทิสมฺปติ ภวนฺตํ โชติปาลํ มาณวํ อามนฺตยติ, ราชา ทิสมฺปติ โภโต โชติปาลสฺส มาณวสฺส ทสฺสนกาโม”ติ. “เอวํ โภ”ติ โข โภ โชติปาโล มาณโว ตสฺส ปุริสสฺส ปฏิสฺสุตฺวา เยน ราชา ทิสมฺปติ เตนุปสงฺกมิ, อุปสงฺกมิตฺวา ทิสมฺปตินา รญฺญา สทฺธิํ สมฺโมทิ, สมฺโมทนียํ กถํ สารณียํ วีติสาเรตฺวา เอกมนฺตํ นิสีทิ.\csromanspacer{} เอกมนฺตํ นิสินฺนํ โข โภ โชติปาลํ มาณวํ ราชา ทิสมฺปติ เอตทโวจ “อนุสาสตุ โน ภวํ โชติปาโล, มา โน ภวํ โชติปาโล อนุสาสนิยา ปจฺจพฺยาหาสิ.\csromanspacer{} เปตฺติเก ตํ ฐาเน ฐเปสฺสามิ โควินฺทิเย อภิสิญฺจิสฺสามี”ติ. “เอวํ โภ”ติ โข โภ โส โชติปาโล มาณโว ทิสมฺปติสฺส รญฺโญ ปจฺจสฺโสสิ.\csromanspacer{} อถ โข โภ ราชา ทิสมฺปติ โชติปาลํ มาณวํ โควินฺทิเย อภิสิญฺจิ, ตํ เปตฺติเก ฐาเน ฐเปสิ.\csromanspacer{} อภิสิตฺโต โชติปาโล มาณโว โควินฺทิเย เปตฺติเก ฐาเน ฐปิโต เยปิสฺส ปิตา อตฺเถ อนุสาสิ, เตปิ อตฺเถ อนุสาสติ, เยปิสฺส ปิตา อตฺเถ นานุสาสิ, เตปิ อตฺเถ อนุสาสติ, เยปิสฺส ปิตา กมฺมนฺเต อภิสมฺโภสิ, เตปิ กมฺมนฺเต อภิสมฺโภติ, เยปิสฺส ปิตา กมฺมนฺเต นาภิสมฺโภสิ, เตปิ กมฺมนฺเต อภิสมฺโภติ.\csromanspacer{} ตเมนํ มนุสฺสา เอวมาหํสุ “โควินฺโท วต โภ พฺราหฺมโณ, มหาโควินฺโท วต โภ พฺราหฺมโณ”ติ.\csromanspacer{} อิมินา โข เอวํ โภ ปริยาเยน โชติปาลสฺส มาณวสฺส โควินฺโท มหาโควินฺโทเตฺวว สมญฺญา อุทปาทิ.}
\csromanmatikaanchor{mtk.99}
\csromantocmarkref{section}{รชฺชสํวิภชน}{mtk.99}
\bookmark[dest={mtk.99},level=1]{รชฺชสํวิภชน}
\csromanheader{1}{\textbf{รชฺชสํวิภชน}}
\proseitem{306}{\csromanfolio{187}อถ โข โภ มหาโควินฺโท พฺราหฺมโณ เยน เต ฉ ขตฺติยา เตนุปสงฺกมิ, อุปสงฺกมิตฺวา เต ฉ ขตฺติเย เอตทโวจ “ทิสมฺปติ โข โภ ราชา ชิณฺโณ วุทฺโธ มหลฺลโก อทฺธคโต วโย-อนุปฺปตฺโต, โก นุ โข ปน โภ ชานาติ ชีวิตํ.\csromanspacer{} ฐานํ โข ปเนตํ วิชฺชติ, ยํ ทิสมฺปติมฺหิ รญฺเญ กาลงฺกเต ราชกตฺตาโร เรณุํ ราชปุตฺตํ รชฺเช อภิสิญฺเจยฺยุํ.\csromanspacer{} อายนฺตุ โภนฺโต, เยน เรณุ ราชปุตฺโต เตนุปสงฺกมถ, อุปสงฺกมิตฺวา เรณุํ ราชปุตฺตํ เอวํ วเทถ ‘มยํ โข โภโต เรณุสฺส สหายา ปิยา มนาปา อปฺปฏิกูลา, ยํสุโข \csromanfolio{188}ภวํ ตํสุขา มยํ, ยํทุกฺโข ภวํ ตํทุกฺขา มยํ.\csromanspacer{} ทิสมฺปติ โข โภ ราชา ชิณฺโณ วุทฺโธ มหลฺลโก อทฺธคโต วโย-อนุปฺปตฺโต, โก นุ โข ปน โภ ชานาติ ชีวิตํ.\csromanspacer{} ฐานํ โข ปเนตํ วิชฺชติ, ยํ ทิสมฺปติมฺหิ รญฺเญ กาลงฺกเต ราชกตฺตาโร ภวนฺตํ เรณุํ รชฺเช อภิสิญฺเจยฺยุํ.\csromanspacer{} สเจ ภวํ เรณุ รชฺชํ ลเภถ, สํวิภเชถ โน รชฺเชนา’ติ”. “เอวํ โภ”ติ โข โภ เต ฉ ขตฺติยา มหาโควินฺทสฺส พฺราหฺมณสฺส ปฏิสฺสุตฺวา เยน เรณุ ราชปุตฺโต เตนุปสงฺกมิํสุ, อุปสงฺกมิตฺวา เรณุํ ราชปุตฺตํ เอตทโวจุํ “มยํ โข โภโต เรณุสฺส สหายา ปิยา มนาปา อปฺปฏิกูลา, ยํสุโข ภวํ ตํสุขา มยํ, ยํทุกฺโข ภวํ ตํทุกฺขา มยํ.\csromanspacer{} ทิสมฺปติ โข โภ ราชา ชิณฺโณ วุทฺโธ มหลฺลโก อทฺธคโต วโย-อนุปฺปตฺโต, โก นุ โข ปน โภ ชานาติ ชีวิตํ.\csromanspacer{} ฐานํ โข ปเนตํ วิชฺชติ, ยํ ทิสมฺปติมฺหิ รญฺเญ กาลงฺกเต ราชกตฺตาโร ภวนฺตํ เรณุํ รชฺเช อภิสิญฺเจยฺยุํ.\csromanspacer{} สเจ ภวํ เรณุ รชฺชํ ลเภถ, สํวิภเชถ โน รชฺเชนา”ติ.\csromanspacer{} โก นุ โข โภ อญฺโญ มม วิชิเต สุโข ภเวถ\footnote{สุขา ภเวยฺยาถ (ก), สุขํ ภเวยฺยาถ (สฺยา), สุขเมเธยฺยาถ (สี, อิ), สุข เมเธถ (?)} อญฺญตฺร ภวนฺเตภิ.\csromanspacer{} สจาหํ โภ รชฺชํ ลภิสฺสามิ, สํวิภชิสฺสามิ โว รชฺเชนาติ.}

\proseitem{307}{อถ โข โภ อโหรตฺตานํ อจฺจเยน ราชา ทิสมฺปติ กาลมกาสิ.\csromanspacer{} ทิสมฺปติมฺหิ รญฺเญ กาลงฺกเต ราชกตฺตาโร เรณุํ ราชปุตฺตํ รชฺเช อภิสิญฺจิํสุ.\csromanspacer{} อภิสิตฺโต เรณุ รชฺเชน ปญฺจหิ กามคุเณหิ สมปฺปิโต สมงฺคีภูโต ปริจาเรติ.\csromanspacer{} อถ โข โภ มหาโควินฺโท พฺราหฺมโณ เยน เต ฉ ขตฺติยา เตนุปสงฺกมิ, อุปสงฺกมิตฺวา เต ฉ ขตฺติเย เอตทโวจ ทิสมฺปติ โข โภ ราชา กาลงฺกโต, อภิสิตฺโต เรณุ รชฺเชน ปญฺจหิ กามคุเณหิ สมปฺปิโต สมงฺคีภูโต ปริจาเรติ.\csromanspacer{} โก นุ โข ปน โภ ชานาติ มทนียา กามา, อายนฺตุ โภนฺโต, เยน เรณุ ราชา เตนุปสงฺกมถ, อุปสงฺกมิตฺวา เรณุํ ราชานํ เอวํ วเทถ ทิสมฺปติ โข โภ ราชา กาลงฺกโต, อภิสิตฺโต ภวํ เรณุ รชฺเชน, สรติ ภวํ ตํ วจนนฺติ.}

\csromanheader{2}{6. มหาโควินฺทสุตฺต}
\proseitem{308}{\csromanfolio{189}“เอวํ โภ”ติ โข โภ เต ฉ ขตฺติยา มหาโควินฺทสฺส พฺราหฺมณสฺส ปฏิสฺสุตฺวา เยน เรณุ ราชา เตนุปสงฺกมิํสุ, อุปสงฺกมิตฺวา เรณุํ ราชานํ เอตทโวจุํ “ทิสมฺปติ โข โภ ราชา กาลงฺกโต, อภิสิตฺโต ภวํ เรณุ รชฺเชน, สรติ ภวํ ตํ วจนนฺ”ติ.\csromanspacer{} สรามหํ โภ ตํ วจนํ\footnote{วจนนฺติ (สฺยา, ก)}.\csromanspacer{} โก นุ โข โภ ปโหติ อิมํ มหาปถวิํ อุตฺตเรน อายตํ ทกฺขิเณน สกฏมุขํ สตฺตธา สมํ สุวิภตฺตํ วิภชิตุนฺติ.\csromanspacer{} โก นุ โข โภ อญฺโญ ปโหติ อญฺญตฺร มหาโควินฺเทน พฺราหฺมเณนาติ.\csromanspacer{} อถ โข โภ เรณุ ราชา อญฺญตรํ ปุริสํ อามนฺเตสิ “เอหิ ตฺวํ อมฺโภ ปุริส เยน มหาโควินฺโท พฺราหฺมโณ เตนุปสงฺกม, อุปสงฺกมิตฺวา มหาโควินฺทํ พฺราหฺมณํ เอวํ วเทหิ “ราชา ตํ ภนฺเต เรณุ อามนฺเตตี”ติ. “เอวํ เทวา”ติ โข โภ โส ปุริโส เรณุสฺส รญฺโญ ปฏิสฺสุตฺวา เยน มหาโควินฺโท พฺราหฺมโณ เตนุปสงฺกมิ, อุปสงฺกมิตฺวา มหาโควินฺทํ พฺราหฺมณํ เอตทโวจ “ราชา ตํ ภนฺเต เรณุ อามนฺเตตี”ติ. “เอวํ โภ”ติ โข โภ มหาโควินฺโท พฺราหฺมโณ ตสฺส ปุริสสฺส ปฏิสฺสุตฺวา เยน เรณุ ราชา เตนุปสงฺกมิ, อุปสงฺกมิตฺวา เรณุนา รญฺญา สทฺธิํ สมฺโมทิ, สมฺโมทนียํ กถํ สารณียํ วีติสาเรตฺวา เอกมนฺตํ นิสีทิ.\csromanspacer{} เอกมนฺตํ นิสินฺนํ โข โภ มหาโควินฺทํ พฺราหฺมณํ เรณุ ราชา เอตทโวจ “เอตุ ภวํ โควินฺโท, อิมํ มหาปถวิํ อุตฺตเรน อายตํ ทกฺขิเณน สกฏมุขํ สตฺตธา สมํ สุวิภตฺตํ วิภชตู”ติ. “เอวํ โภ”ติ โข โภ มหาโควินฺโท พฺราหฺมโณ เรณุสฺส รญฺโญ ปฏิสฺสุตฺวา อิมํ มหาปถวิํ อุตฺตเรน อายตํ ทกฺขิเณน สกฏมุขํ สตฺตธา สมํ สุวิภตฺตํ วิภชิ.\csromanspacer{} สพฺพานิ สกฏมุขานิ ปฏฺฐเปสิ\footnote{อฏฺฐเปสิ (สี, อิ)}.\csromanspacer{} ตตฺร สุทํ มชฺเฌ เรณุสฺส รญฺโญ ชนปโท โหติ.}

\csromanhangingbegin
\hangingheaditem{309}{ทนฺตปุรํ กลิงฺคานํ\footnote{กาลิงฺคานํ (สฺยา, อิ, ก)}, อสฺสกานญฺจ โปตนํ.}
\hangingbody{มเหสยํ4 อวนฺตีนํ, โสวีรานญฺจ โรรุกํ.}
\hangingbody{มิถิลา จ วิเทหานํ, จมฺปา องฺเคสุ มาปิตา.}
\hangingbody{พาราณสี จ กาสีนํ, เอเต โควินฺทมาปิตาติ.}
\csromanhangingend

\proseitem{310}{\csromanfolio{190}อถ โข โภ เต ฉ ขตฺติยา ยถาสเกน ลาเภน อตฺตมนา อเหสุํ ปริปุณฺณสงฺกปฺปา “ยํ วต โน อโหสิ อิจฺฉิตํ, ยํ อากงฺขิตํ, ยํ อธิปฺเปตํ, ยํ อภิปตฺถิตํ, ตํ โน ลทฺธนฺ”ติ.}

\csromangathameasure{สตฺตภู พฺรหฺมทตฺโต จ, เวสฺสภู ภรโต สห. \\ เรณุ เทฺว ธตรฏฺฐา จ, ตทาสุํ สตฺต ภารธาติ.}
\csromangathagroup{\csromangathastanza{สตฺตภู พฺรหฺมทตฺโต จ, เวสฺสภู ภรโต สห. \\ เรณุ เทฺว ธตรฏฺฐา จ, ตทาสุํ สตฺต ภารธาติ.}}

\nitthitamruled{ปฐมภาณวาโร นิฏฺฐิโต.}
\csromanmatikaanchor{mtk.100}
\csromantocmarkref{section}{กิตฺติสทฺท-อพฺภุคฺคมน}{mtk.100}
\bookmark[dest={mtk.100},level=1]{กิตฺติสทฺท-อพฺภุคฺคมน}
\csromanheader{1}{\textbf{กิตฺติสทฺท-อพฺภุคฺคมน}}
\proseitem{311}{อถ โข โภ เต ฉ ขตฺติยา เยน มหาโควินฺโท พฺราหฺมโณ เตนุปสงฺกมิํสุ, อุปสงฺกมิตฺวา มหาโควินฺทํ พฺราหฺมณํ เอตทโวจุํ “ยถา โข ภวํ โควินฺโท เรณุสฺส รญฺโญ สหาโย ปิโย มนาโป อปฺปฏิกูโล, เอวเมว โข ภวํ โควินฺโท อมฺหากํปิ สหาโย ปิโย มนาโป อปฺปฏิกูโล, อนุสาสตุ โน ภวํ โควินฺโท, มา โน ภวํ โควินฺโท อนุสาสนิยา ปจฺจพฺยาหาสี”ติ. “เอวํ โภ”ติ โข โภ มหาโควินฺโท พฺราหฺมโณ เตสํ ฉนฺนํ ขตฺติยานํ ปจฺจสฺโสสิ.\csromanspacer{} อถ โข โภ มหาโควินฺโท พฺราหฺมโณ สตฺต จ ราชาโน ขตฺติเย มุทฺธาวสิตฺเต รชฺเช\footnote{มุทฺธาภิสิตฺเต รชฺเชน (สฺยา)} อนุสาสิ, สตฺต จ พฺราหฺมณมหาสาเล สตฺต จ นฺหาตกสตานิ มนฺเต วาเจสิ.}

\proseitem{312}{อถ โข โภ มหาโควินฺทสฺส พฺราหฺมณสฺส อปเรน สมเยน เอวํ กลฺยาโณ กิตฺติสทฺโท อพฺภุคฺคจฺฉิ\footnote{อพฺภุคฺคญฺฉิ (สี, อิ)} “สกฺขิ มหาโควินฺโท พฺราหฺมโณ พฺรหฺมานํ ปสฺสติ, สกฺขิ มหาโควินฺโท พฺราหฺมโณ พฺรหฺมุนา สากจฺเฉติ สลฺลปติ มนฺเตตี”ติ.\csromanspacer{} อถ โข โภ มหาโควินฺทสฺส พฺราหฺมณสฺส เอตทโหสิ “มยฺหํ โข เอวํ กลฺยาโณ กิตฺติสทฺโท อพฺภุคฺคโต ‘สกฺขิ มหาโควินฺโท พฺราหฺมโณ พฺรหฺมานํ ปสฺสติ, สกฺขิ มหาโควินฺโท พฺราหฺมโณ พฺรหฺมุนา สากจฺเฉติ สลฺลปติ มนฺเตตี’ติ, น โข ปนาหํ พฺรหฺมานํ ปสฺสามิ, น พฺรหฺมุนา สากจฺเฉมิ, น พฺรหฺมุนา}

\vspace{\tipitakaparskip}
\csromancenter{มหาโควินฺทสุตฺต สลฺลปามิ, น พฺรหฺมุนา มนฺเตมิ, สุตํ โข ปน เมตํ พฺราหฺมณานํ วุทฺธานํ มหลฺลกานํ อาจริยปาจริยานํ ภาสมานานํ ‘โย วสฺสิเก จตฺตาโร มาเส ปฏิสลฺลียติ, กรุณํ ฌานํ ฌายติ, โส พฺรหฺมานํ ปสฺสติ พฺรหฺมุนา สากจฺเฉติ พฺรหฺมุนา สลฺลปติ พฺรหฺมุนา มนฺเตตี’ติ, ยํนูนาหํ วสฺสิเก จตฺตาโร มาเส ปฏิสลฺลีเยยฺยํ, กรุณํ ฌานํ ฌาเยยฺยนฺ”ติ.}
\proseitem{313}{\csromanfolio{191}อถ โข โภ มหาโควินฺโท พฺราหฺมโณ เยน เรณุ ราชา เตนุปสงฺกมิ, อุปสงฺกมิตฺวา เรณุํ ราชานํ เอตทโวจ “มยฺหํ โข โภ เอวํ กลฺยาโณ กิตฺติสทฺโท อพฺภุคฺคโต ‘สกฺขิ มหาโควินฺโท พฺราหฺมโณ พฺรหฺมานํ ปสฺสติ, สกฺขิ มหาโควินฺโท พฺราหฺมโณ พฺรหฺมุนา สากจฺเฉติ สลฺลปติ มนฺเตตี’ติ, น โข ปนาหํ โภ พฺรหฺมานํ ปสฺสฺéมิ, น พฺรหฺมุนา สากจฺเฉมิ, น พฺรหฺมุนา สลฺลปามิ, น พฺรหฺมุนา มนฺเตมิ, สุตํ โข ปน เมตํ พฺราหฺมณานํ วุทฺธานํ มหลฺลกานํ อาจริยปาจริยานํ ภาสมานานํ ‘โย วสฺสิเก จตฺตาโร มาเส ปฏิสลฺลียติ, กรุณํ ฌานํ ฌายติ, โส พฺรหฺมานํ ปสฺสติ, พฺรหฺมุนา สากจฺเฉติ พฺรหฺมุนา สลฺลปติ พฺรหฺมุนา มนฺเตตี’ติ, อิจฺฉามหํ โภ วสฺสิเก จตฺตาโร มาเส ปฏิสลฺลียิตุํ, กรุณํ ฌานํ ฌายิตุํ, นมฺหิ เกนจิ อุปสงฺกมิตพฺโพ อญฺญตฺร เอเกน ภตฺตาภิหาเรนา”ติ.\csromanspacer{} ยสฺสทานิ ภวํ โควินฺโท กาลํ มญฺญตีติ.}

\proseitem{314}{อถ โข โภ มหาโควินฺโท พฺราหฺมโณ เยน เต ฉ ขตฺติยา เตนุปสงฺกมิ, อุปสงฺกมิตฺวา เต ฉ ขตฺติเย เอตทโวจ “มยฺหํ โข โภ เอวํ กลฺยาโณ กิตฺติสทฺโท อพฺภุคฺคโต ‘สกฺขิ มหาโควินฺโท พฺราหฺมโณ พฺรหฺมานํ ปสฺสติ, สกฺขิ มหาโควินฺโท พฺราหฺมโณ พฺรหฺมุนา สากจฺเฉติ สลฺลปติ มนฺเตตี’ติ, น โข ปนาหํ โภ พฺรหฺมานํ ปสฺสามิ, น พฺรหฺมุนา สากจฺเฉมิ, น พฺรหฺมุนา สลฺลปามิ, น พฺรหฺมุนา มนฺเตมิ, สุตํ โข ปน เมตํ พฺราหฺมณานํ วุทฺธานํ มหลฺลกานํ อาจริยปาจริยานํ ภาสมานานํ ‘โย วสฺสิเก จตฺตาโร มาเส ปฏิสลฺลียติ, กรุณํ ฌานํ ฌายติ, โส พฺรหฺมานํ ปสฺสติ พฺรหฺมุนา สากจฺเฉติ พฺรหฺมุนา สลฺลปติ พฺรหฺมุนา มนฺเตตี’ติ, อิจฺฉามหํ โภ วสฺสิเก จตฺตาโร มาเส ปฏิสลฺลียิตุํ, กรุณํ ฌานํ ฌายิตุํ, นมฺหิ เกนจิ อุปสงฺกมิตพฺโพ อญฺญตฺร เอเกน ภตฺตาภิหาเรนา”ติ.\csromanspacer{} ยสฺสทานิ ภวํ โควินฺโท กาลํ มญฺญตีติ.}

\proseitem{315}{\csromanfolio{192}อถ โข โภ มหาโควินฺโท พฺราหฺมโณ เยน เต สตฺต จ พฺราหฺมณมหาสาลา สตฺต จ นฺหาตกสตานิ เตนุปสงฺกมิ, อุปสงฺกมิตฺวา เต สตฺต จ พฺราหฺมณมหาสาเล สตฺต จ นฺหาตกสตานิ เอตทโวจ “มยฺหํ โข โภ เอวํ กลฺยาโณ กิตฺติสทฺโท อพฺภุคฺคโต ‘สกฺขิ มหาโควินฺโท พฺราหฺมโณ พฺรหฺมานํ ปสฺสติ, สกฺขิ มหาโควินฺโท พฺราหฺมโณ พฺรหฺมุนา สากจฺเฉติ สลฺลปติ มนฺเตตี’ติ, น โข ปนาหํ โภ พฺรหฺมานํ ปสฺสามิ, น พฺรหฺมุนา สากจฺเฉมิ, น พฺรหฺมุนา สลฺลปามิ, น พฺรหฺมุนา มนฺเตมิ, สุตํ โข ปน เมตํ พฺราหฺมณานํ วุทฺธานํ มหลฺลกานํ อาจริยปาจริยานํ ภาสมานานํ ‘โย วสฺสิเก จตฺตาโร มาเส ปฏิสลฺลียติ, กรุณํ ฌานํ ฌายติ, โส พฺรหฺมานํ ปสฺสติ, พฺรหฺมุนา สากจฺเฉติ, พฺรหฺมุนา สลฺลปติ พฺรหฺมุนา มนฺเตตี’ติ, เตน หิ โภ ยถาสุเต ยถาปริยตฺเต มนฺเต วิตฺถาเรน สชฺฌายํ กโรถ, อญฺญมญฺญํ จ มนฺเต วาเจถ, อิจฺฉามหํ โภ วสฺสิเก จตฺตาโร มาเส ปฏิสลฺลียิตุํ, กรุณํ ฌานํ ฌายิตุํ, นมฺหิ เกนจิ อุปสงฺกมิตพฺโพ อญฺญตฺร เอเกน ภตฺตาภิหาเรนา”ติ.\csromanspacer{} ยสฺสทานิ ภวํ โควินฺโท กาลํ มญฺญตีติ.}

\proseitem{316}{อถ โข โภ มหาโควินฺโท พฺราหฺมโณ เยน จตฺตารีสา ภริยา สาทิสิโย เตนุปสงฺกมิ, อุปสงฺกมิตฺวา จตฺตารีสา ภริยา สาทิสิโย เอตทโวจ “มยฺหํ โข โภตี เอวํ กลฺยาโณ กิตฺติสทฺโท อพฺภุคฺคโต ‘สกฺขิ มหาโควินฺโท พฺราหฺมโณ พฺรหฺมานํ ปสฺสติ, สกฺขิ มหาโควินฺโท พฺราหฺมโณ พฺรหฺมุนา สากจฺเฉติ สลฺลปติ มนฺเตตี’ติ, น โข ปเนหํ โภตี พฺรหฺมานํ ปสฺสามิ, น พฺรหฺมุนา สากจฺเฉมิ, น พฺรหฺมุนา สลฺลปามิ, น พฺรหฺมุนา มนฺเตมิ, สุตํ โข ปน เมตํ พฺราหฺมณานํ วุทฺธานํ มหลฺลกานํ อาจริยปาจริยานํ ภาสมานานํ ‘โย วสฺสิเก จตฺตาโร มาเส ปฏิสลฺลียติ, กรุณํ ฌานํ ฌายติ, โส พฺรหฺมานํ ปสฺสติ, พฺรหฺมุนา สากจฺเฉติ, พฺรหฺมุนา สลฺลปติ, พฺรหฺมุนา มนฺเตตี’ติ, อิจฺฉามหํ โภตี วสฺสิเก จตฺตาโร มาเส ปฏิสลฺลียิตุํ, กรุณํ ฌานํ ฌายิตุํ, นมฺหิ เกนจิ อุปสงฺกมิตพฺโพ อญฺญตฺร เอเกน ภตฺตาภิหาเรนา”ติ.\csromanspacer{} ยสฺสทานิ ภวํ โควินฺโท กาลํ มญฺญตีติ.}

\proseitem{317}{อถ โข โภ มหาโควินฺโท พฺราหฺมโณ ปุรตฺถิเมน นครสฺส นวํ สนฺธาคารํ การาเปตฺวา วสฺสิเก จตฺตาโร มาเส ปฏิสลฺลียิ,}

\vspace{\tipitakaparskip}
\csromancenter{มหาโควินฺทสุตฺต กรุณํ ฌานํ ฌายิ, นาสฺสุธ โกจิ อุปสงฺกมติ\footnote{อุปสงฺกมิ (อิ)} อญฺญตฺร เอเกน ภตฺตาภิหาเรน.\csromanspacer{} อถ โข โภ มหาโควินฺทสฺส พฺราหฺมณสฺส จตุนฺนํ มาสานํ อจฺจเยน อหุเทว อุกฺกณฺฐนา อหุ ปริตสฺสนา “สุตํ โข ปน เมตํ พฺราหฺมณานํ วุทฺธานํ มหลฺลกานํ อาจริยปาจริยานํ ภาสมานานํ ‘โย วสฺสิเก จตฺตาโร มาเส ปฏิสลฺลียติ, กรุณํ ฌานํ ฌายติ, โส พฺรหฺมานํ ปสฺสติ, พฺรหฺมุนา สากจฺเฉติ พฺรหฺมุนา สลฺลปติ พฺรหฺมุนา มนฺเตตี’ติ, น โข ปนาหํ พฺรหฺมานํ ปสฺสามิ, น พฺรหฺมุนา สากจฺเฉมิ น พฺรหฺมุนา สลฺลปามิ น พฺรหฺมุนา มนฺเตมี”ติ.}
\csromanmatikaanchor{mtk.101}
\csromantocmarkref{section}{พฺรหฺมุนาสากจฺฉา}{mtk.101}
\bookmark[dest={mtk.101},level=1]{พฺรหฺมุนาสากจฺฉา}
\csromanheader{1}{\textbf{พฺรหฺมุนาสากจฺฉา}}
\proseitem{318}{\csromanfolio{193}อถ โข โภ พฺรหฺมา สนงฺกุมาโร มหาโควินฺทสฺส พฺราหฺมณสฺส เจตสา เจโตปริวิตกฺกมญฺญาย เสยฺยถาปิ นาม พลวา ปุริโส สมิญฺชิตํ วา พาหํ ปสาเรยฺย, ปสาริตํ วา พาหํ สมิญฺเชยฺย, เอวเมว พฺรหฺมโลเก อนฺตรหิโต มหาโควินฺทสฺส พฺราหฺมณสฺส สมฺมุเข ปาตุรโหสิ.\csromanspacer{} อถ โข โภ มหาโควินฺทสฺส พฺราหฺมณสฺส อหุเทว ภยํ อหุ ฉมฺภิตตฺตํ อหุ โลมหํโส ยถา ตํ อทิฏฺฐปุพฺพํ รูปํ ทิสฺวา.\csromanspacer{} อถ โข โภ มหาโควินฺโท พฺราหฺมโณ ภีโต สํวิคฺโค โลมหฏฺฐชาโต พฺรหฺมานํ สนงฺกุมารํ คาถาย อชฺฌภาสิ\csromandash{}}

\csromangathameasure{“วณฺณวา ยสวา สิริมา, โก นุ ตฺวมสิ มาริส. \\ อชานนฺตา ตํ ปุจฺฉาม, กถํ ชาเนมุ ตํ มยนฺ”ติ. \\ มํ เว กุมารํ ชานนฺติ, พฺรหฺมโลเก สนนฺตนํ. \\ สพฺเพ ชานนฺติ มํ เทวา, เอวํ โควินฺท ชานหิ. \\ “อาสนํ อุทกํ ปชฺชํ, มธุสากญฺจ พฺรหฺมุโน. \\ อคฺเฆ ภวนฺตํ ปุจฺฉาม, อคฺฆํ กุรุตุ โน ภวํ”. \\ ปฏิคฺคณฺหาม เต อคฺฆํ, ยํ ตฺวํ โควินฺท ภาสสิ. \\ ทิฏฺฐธมฺมหิตตฺถาย, สมฺปราย สุขาย จ. \\ กตาวกาโส ปุจฺฉสฺสุ, ยํ กิญฺจิ อภิปตฺถิตนฺติ.}
\csromangathagroup{\csromangathastanza{“วณฺณวา ยสวา สิริมา, โก นุ ตฺวมสิ มาริส. \\ อชานนฺตา ตํ ปุจฺฉาม, กถํ ชาเนมุ ตํ มยนฺ”ติ.}\csromangathastanza{มํ เว กุมารํ ชานนฺติ, พฺรหฺมโลเก สนนฺตนํ\footnote{สนนฺติจ (ก)}. \\ สพฺเพ ชานนฺติ มํ เทวา, เอวํ โควินฺท ชานหิ.}\csromangathastanza{“อาสนํ อุทกํ ปชฺชํ, มธุสากญฺจ\footnote{มธุปากญฺจ (สี, สฺยา, อิ)} พฺรหฺมุโน. \\ อคฺเฆ ภวนฺตํ ปุจฺฉาม, อคฺฆํ กุรุตุ โน ภวํ”.}\csromangathastanza{ปฏิคฺคณฺหาม เต อคฺฆํ, ยํ ตฺวํ โควินฺท ภาสสิ. \\ ทิฏฺฐธมฺมหิตตฺถาย, สมฺปราย สุขาย จ. \\ กตาวกาโส ปุจฺฉสฺสุ, ยํ กิญฺจิ อภิปตฺถิตนฺติ.}}

\proseitem{319}{อถ โข โภ มหาโควินฺทสฺส พฺราหฺมณสฺส เอตทโหสิ “กตาวกาโส โขมฺหิ พฺรหฺมุนา สนงฺกุมาเรน, กิํ นุ โข อหํ พฺรหฺมานํ \csromanfolio{194}สนงฺกุมารํ ปุจฺเฉยฺยํ ทิฏฺฐธมฺมิกํ วา อตฺถํ สมฺปรายิกํ วา”ติ.\csromanspacer{} อถ โข โภ มหาโควินฺทสฺส พฺราหฺมณสฺส เอตทโหสิ “กุสโล โข อหํ ทิฏฺฐธมฺมิกานํ อตฺถานํ, อญฺเญปิ มํ ทิฏฺฐธมฺมิกํ อตฺถํ ปุจฺฉนฺติ, ยํนูนาหํ พฺรหฺมานํ สนงฺกุมารํ สมฺปรายิกญฺเญว อตฺถํ ปุจฺเฉยฺยนฺ”ติ.\csromanspacer{} อถ โข โภ มหาโควินฺโท พฺราหฺมโณ พฺรหฺมานํ สนงฺกุมารํ คาถาย อชฺฌภาสิ\csromandash{}}

\csromangathameasure{“ปุจฺฉามิ พฺรหฺมานํ สนงฺกุมารํ, \\ กงฺขี อกงฺขิํ ปรเวทิเยสุ. \\ กตฺถฏฺฐิโต กิมฺหิ จ สิกฺขมาโน, \\ ปปฺโปติ มจฺโจ อมตํ พฺรหฺมโลกนฺ”ติ. \\ หิตฺวา มมตฺตํ มนุเชสุ พฺรหฺเม, \\ เอโกทิภูโต กรุเณธิมุตฺโต. \\ นิรามคนฺโธ วิรโต เมถุนสฺมา, \\ เอตฺถฏฺฐิโต เอตฺถ จ สิกฺขมาโน. \\ ปปฺโปติ มจฺโจ อมตํ พฺรหฺมโลกนฺติ.}
\csromangathagroup{\csromangathastanza{“ปุจฺฉามิ พฺรหฺมานํ สนงฺกุมารํ, \\ กงฺขี อกงฺขิํ ปรเวทิเยสุ. \\ กตฺถฏฺฐิโต กิมฺหิ จ สิกฺขมาโน, \\ ปปฺโปติ มจฺโจ อมตํ พฺรหฺมโลกนฺ”ติ.}\csromangathastanza{หิตฺวา มมตฺตํ มนุเชสุ พฺรหฺเม, \\ เอโกทิภูโต กรุเณธิมุตฺโต\footnote{กรุณาธิมุตฺโต (สี, สฺยา, อิ)}. \\ นิรามคนฺโธ วิรโต เมถุนสฺมา, \\ เอตฺถฏฺฐิโต เอตฺถ จ สิกฺขมาโน.}\csromangathastanza{ปปฺโปติ มจฺโจ อมตํ พฺรหฺมโลกนฺติ.}}

\proseitem{320}{“หิตฺวา มมตฺตนฺ”ติ อหํ โภโต อาชานามิ, อิเธกจฺโจ อปฺปํ วา โภคกฺขนฺธํ ปหาย มหนฺตํ วา โภคกฺขนฺธํ ปหาย อปฺปํ วา ญาติปริวฏฺฏํ ปหาย มหนฺตํ วา ญาติปริวฏฺฏํ ปหาย เกสมสฺสุํ โอหาเรตฺวา กาสายานิ วตฺถานิ อจฺฉาเทตฺวา อคารสฺมา อนคาริยํ ปพฺพชติ, อิติ “หิตฺวา มมตฺตนฺ”ติ อหํ โภโต อาชานามิ. “เอโกทิภูโต”ติ อหํ โภโต อาชานามิ, อิเธกจฺโจ วิวิตฺตํ เสนาสนํ ภชติ อรญฺญํ รุกฺขมูลํ ปพฺพตํ กนฺทรํ คิริคุหํ สุสานํ วนปตฺถํ อพฺโภกาสํ ปลาลปุญฺชํ, อิติ “เอโกทิภูโต”ติ อหํ โภโต อาชานามิ. “กรุเณธิมุตฺโต”ติ อหํ โภโต อาชานามิ, อิเธกจฺโจ กรุณาสหคเตน เจตสา เอกํ ทิสํ ผริตฺวา วิหรติ.\csromanspacer{} ตถา ทุติยํ.\csromanspacer{} ตถา ตติยํ.\csromanspacer{} ตถา จตุตฺถํ.\csromanspacer{} อิติ อุทฺธมโธติริยํ สพฺพธิ สพฺพตฺตตาย สพฺพาวนฺตํ โลกํ กรุณาสหคเตน เจตสา วิปุเลน มหคฺคเตน อปฺปมาเณน อเวเรน อพฺยาปชฺเชน ผริตฺวา วิหรติ, อิติ “กรุเณธิมุตฺโต”ติ อหํ โภโต อาชานามิ.\csromanspacer{} อามคนฺเธ จ โข อหํ โภโต ภาสมานสฺส น อาชานามิ.}

\csromanheader{2}{6. มหาโควินฺทสุตฺต}
\csromangathameasure{เก อามคนฺธา มนุเชสุ พฺรหฺเม, \\ เอเต อวิทฺวา อิธ พฺรูหิ ธีร. \\ เกนาวฏา วาติ ปชา กุรุตุ, \\ อาปายิกา นิวุตพฺรหฺมโลกาติ. \\ โกโธ โมสวชฺชํ นิกติ จ ทุพฺโภ, \\ กทริยตา อติมาโน อุสูยา. \\ อิจฺฉา วิวิจฺฉา ปรเหฐนา จ, \\ โลโภ จ โทโส จ มโท จ โมโห. \\ เอเตสุ ยุตฺตา อนิรามคนฺธา, \\ อาปายิกา นิวุตพฺรหฺมโลกาติ.}
\csromangathagroup{\csromangathastanza{\csromanfolio{195}เก อามคนฺธา มนุเชสุ พฺรหฺเม, \\ เอเต อวิทฺวา อิธ พฺรูหิ ธีร. \\ เกนาวฏา\footnote{เกนาวุฏา (สฺยา)} วาติ ปชา กุรุตุ\footnote{กุรุรู (สฺยา), กุรุฏฺฐรู (อิ), กุรูรุ (?)}, \\ อาปายิกา นิวุตพฺรหฺมโลกาติ.}\csromangathastanza{โกโธ โมสวชฺชํ นิกติ จ ทุพฺโภ, \\ กทริยตา อติมาโน อุสูยา. \\ อิจฺฉา วิวิจฺฉา ปรเหฐนา จ, \\ โลโภ จ โทโส จ มโท จ โมโห. \\ เอเตสุ ยุตฺตา อนิรามคนฺธา, \\ อาปายิกา นิวุตพฺรหฺมโลกาติ.}}

\prose{ยถา โข อหํ โภโต อามคนฺเธ ภาสมานสฺส อาชานามิ.\csromanspacer{} เต น สุนิมฺมทยา อคารํ อชฺฌาวสตา.\csromanspacer{} ปพฺพชิสฺสามหํ โภ อคารสฺมา อนคาริยนฺติ.\csromanspacer{} ยสฺสทานิ ภวํ โควินฺโท กาลํ มญฺญตีติ.}

\csromanmatikaanchor{mtk.102}
\csromantocmarkref{section}{เรณุราช-อามนฺตนา}{mtk.102}
\bookmark[dest={mtk.102},level=1]{เรณุราช-อามนฺตนา}
\csromanheader{1}{\textbf{เรณุราช อามนฺตนา}}
\proseitem{321}{อถ โข โภ มหาโควินฺโท พฺราหฺมโณ เยน เรณุ ราชา เตนุปสงฺกมิ, อุปสงฺกมิตฺวา เรณุํ ราชานํ เอตทโวจ “อญฺญํ ทานิ ภวํ ปุโรหิตํ ปริเยสตุ, โย โภโต รชฺชํ อนุสาสิสฺสติ, อิจฺฉามหํ โภ อคารสฺมา อนคาริยํ ปพฺพชิตุํ.\csromanspacer{} ยถา โข ปน เม สุตํ พฺรหฺมุโน อามคนฺเธ ภาสมานสฺส, เต น สุนิมฺมทยา อคารํ อชฺฌาวสตา, ปพฺพชิสฺสามหํ โภ อคารสฺมา อนคาริยนฺ”ติ.}

\csromangathameasure{อามนฺตยามิ ราชานํ, เรณุํ ภูมิปติํ อหํ. \\ ตฺวํ ปชานสฺสุ รชฺเชน, นาหํ โปโรหิจฺเจ รเม. \\ สเจ เต อูนํ กาเมหิ, อหํ ปริปูรยามิ เต. \\ โย ตํ หิํสติ วาเรมิ, ภูมิเสนาปติ อหํ. \\ ตุวํ ปิตา อหํ ปุตฺโต, มา โน โควินฺท ปาชหิ. \\ นมตฺถิ อูนํ กาเมหิ, หิํสิตา เม น วิชฺชติ. \\ อมนุสฺสวโจ สุตฺวา, ตสฺมาหํ น คเห รเม. \\ อมนุสฺโส กถํวณฺโณ, กิํ เต อตฺถํ อภาสถ. \\ ยญฺจ สุตฺวา ชหาสิ โน, เคเห อเมฺห จ เกวลี. \\ อุปวุตฺถสฺส เม ปุพฺเพ, ยิฏฺฐุกามสฺส เม สโต. \\ อคฺคิ ปชฺชลิโต อาสิ, กุสปตฺตปริตฺถโต. \\ ตโต เม พฺรหฺมา ปาตุรหุ, พฺรหฺมโลกา สนนฺตโน. \\ โส เม ปญฺหํ วิยากาสิ, ตํ สุตฺวา น คเห รเม. \\ สทฺทหามิ อหํ โภโต, ยํ ตฺวํ โควินฺท ภาสสิ. \\ อมนุสฺสวโจ สุตฺวา, กถํ วตฺเตถ อญฺญถา. \\ เต ตํ อนุวตฺติสฺสาม, สตฺถา โควินฺท โน ภวํ. \\ มณิ ยถา เวฬุริโย, อกาโจ วิมโล สุโภ. \\ เอวํ สุทฺธา จริสฺสาม, โควินฺทสฺสนุสาสเนติ.}
\csromangathagroup{\csromangathastanza{อามนฺตยามิ ราชานํ, เรณุํ ภูมิปติํ อหํ. \\ ตฺวํ ปชานสฺสุ รชฺเชน, นาหํ โปโรหิจฺเจ รเม.}\csromangathastanza{สเจ เต อูนํ กาเมหิ, อหํ ปริปูรยามิ เต. \\ โย ตํ หิํสติ วาเรมิ, ภูมิเสนาปติ อหํ.}\csromangathastanza{ตุวํ ปิตา อหํ ปุตฺโต, มา โน โควินฺท ปาชหิ\footnote{ปาเชหิ (อฏฺฐกถายํ สํวณฺณิตปาฐนฺตรํ)}. \\ นมตฺถิ อูนํ กาเมหิ, หิํสิตา เม น วิชฺชติ.}\csromangathastanza{อมนุสฺสวโจ สุตฺวา, ตสฺมาหํ น คเห รเม. \\ อมนุสฺโส กถํวณฺโณ, กิํ เต อตฺถํ อภาสถ.}\csromangathastanza{\csromanfolio{196}ยญฺจ สุตฺวา ชหาสิ โน, เคเห อเมฺห จ เกวลี. \\ อุปวุตฺถสฺส เม ปุพฺเพ, ยิฏฺฐุกามสฺส เม สโต.}\csromangathastanza{อคฺคิ ปชฺชลิโต อาสิ, กุสปตฺตปริตฺถโต. \\ ตโต เม พฺรหฺมา ปาตุรหุ, พฺรหฺมโลกา สนนฺตโน.}\csromangathastanza{โส เม ปญฺหํ วิยากาสิ, ตํ สุตฺวา น คเห รเม. \\ สทฺทหามิ อหํ โภโต, ยํ ตฺวํ โควินฺท ภาสสิ.}\csromangathastanza{อมนุสฺสวโจ สุตฺวา, กถํ วตฺเตถ อญฺญถา. \\ เต ตํ อนุวตฺติสฺสาม, สตฺถา โควินฺท โน ภวํ.}\csromangathastanza{มณิ ยถา เวฬุริโย, อกาโจ วิมโล สุโภ. \\ เอวํ สุทฺธา จริสฺสาม, โควินฺทสฺสนุสาสเนติ.}}

\prose{สเจ ภวํ โควินฺโท อคารสฺมา อนคาริยํ ปพฺพชิสฺสติ, มยมฺปิ อคารสฺมา อนคาริยํ ปพฺพชิสฺสาม.\csromanspacer{} อถ ยา เต คติ, สา โน คติ ภวิสฺสตีติ.}

\csromanmatikaanchor{mtk.103}
\csromantocmarkref{section}{ฉขตฺติย-อามนฺตนา}{mtk.103}
\bookmark[dest={mtk.103},level=1]{ฉขตฺติย-อามนฺตนา}
\csromanheader{1}{\textbf{ฉขตฺติย อามนฺตนา}}
\proseitem{322}{อถ โข โภ มหาโควินฺโท พฺราหฺมโณ เยน เต ฉ ขตฺติยา เตนุปสงฺกมิ, อุปสงฺกมิตฺวา เต ฉ ขตฺติเย เอตทโวจ “อญฺญํ ทานิ ภวนฺโต ปุโรหิตํ ปริเยสนฺตุ, โย ภวนฺตานํ รชฺเช อนุสาสิสฺสติ, อิจฺฉามหํ โภ อคารสฺมา อนคาริยํ ปพฺพชิตุํ, ยถา โข ปน เม สุตํ พฺรหฺมุโน อามคนฺเธ ภาสมานสฺส, เต น สุนิมฺมทยา อคารํ อชฺฌาวสตา, ปพฺพชิสฺสามหํ โภ อคารสฺมา อนคาริยนฺ”ติ.\csromanspacer{} อถ โข โภ เต ฉ ขตฺติยา เอกมนฺตํ อปกฺกมฺม เอวํ สมจินฺเตสุํ “อิเม โข พฺราหฺมณา นาม ธนลุทฺธา, ยํนูน มยํ มหาโควินฺทํ พฺราหฺมณํ ธเนน สิกฺเขยฺยามา”ติ.\csromanspacer{} เต มหาโควินฺทํ พฺราหฺมณํ อุปสงฺกมิตฺวา เอวมาหํสุ “สํวิชฺชติ โข โภ อิเมสุ สตฺตสุ รชฺเชสุ ปหูตํ สาปเตยฺยํ, ตโต โภโต ยาวตเกน อตฺโถ, ตาวตกํ อาหรียตนฺ”ติ.\csromanspacer{} อลํ โภ มมปิทํ ปหูตํ สาปเตยฺยํ ภวนฺตานํเยว วาหสา, ตมหํ สพฺพํ ปหาย อคารสฺมา อนคาริยํ ปพฺพชิสฺสามิ.\csromanspacer{} ยถา โข ปน เม สุตํ พฺรหฺมุโน อามคนฺเธ ภาสมานสฺส, เต น สุนิมฺมทยา}

\vspace{\tipitakaparskip}
\csromancenter{มหาโควินฺทสุตฺต อคารํ อชฺฌาวสตา, ปพฺพชิสฺสามหํ โภ อคารสฺมา อนคาริยนฺติ.\csromanspacer{} อถ โข โภ เต ฉ ขตฺติยา เอกมนฺตํ อปกฺกมฺม เอวํ สมจินฺเตสุํ “อิเม โข พฺราหฺมณา นาม อิตฺถิลุทฺธา, ยํนูน มยํ มหาโควินฺทํ พฺราหฺมณํ อิตฺถีหิ สิกฺเขยฺยามา”ติ.\csromanspacer{} เต มหาโควินฺทํ พฺราหฺมณํ อุปสงฺกมิตฺวา เอวมาหํสุ “สํวิชฺชนฺติ โข โภ อิเมสุ สตฺตสุ รชฺเชสุ ปหูตา อิตฺถิโย, ตโต โภโต ยาวติกาหิ อตฺโถ, ตาวติกา อานียตนฺ”ติ.\csromanspacer{} อลํ โภ มมปิมา\footnote{มมปิตา(ก), มมปิ (สี)} จตฺตารีสา ภริยา สาทิสิโย, ตาปาหํ สพฺพา ปหาย อคารสฺมา อนคาริยํ ปพฺพชิสฺสามิ.\csromanspacer{} ยถา โข ปน เม สุตํ พฺรหฺมุโน อามคนฺเธ ภาสมานสฺส, เต น สุนิมฺมทยา อคารํ อชฺฌาวสตา, ปพฺพชิสฺสามหํ โภ อคารสฺมา อนคาริยนฺติ.}
\proseitem{323}{\csromanfolio{197}สเจ ภวํ โควินฺโท อคารสฺมา อนคาริยํ ปพฺพชิสฺสติ, มยมฺปิ อคารสฺมา อนคาริยํ ปพฺพชิสฺสาม, อถ ยา เต คติ, สา โน คติ ภวิสฺสตีติ.}

\csromangathameasure{สเจ ชหถ กามานิ, ยตฺถ สตฺโต ปุถุชฺชโน. \\ อารมฺภโวฺห ทฬฺหา โหถ, ขนฺตีพลสมาหิตา. \\ เอส มคฺโค อุชุมคฺโค, เอส มคฺโค อนุตฺตโร. \\ สทฺธมฺโม สพฺภิ รกฺขิโต, พฺรหฺมโลกูปปตฺติยาติ.}
\csromangathagroup{\csromangathastanza{สเจ ชหถ กามานิ, ยตฺถ สตฺโต ปุถุชฺชโน. \\ อารมฺภโวฺห ทฬฺหา โหถ, ขนฺตีพลสมาหิตา.}\csromangathastanza{เอส มคฺโค อุชุมคฺโค, เอส มคฺโค อนุตฺตโร. \\ สทฺธมฺโม สพฺภิ รกฺขิโต, พฺรหฺมโลกูปปตฺติยาติ.}}

\prose{เตน หิ ภวํ โควินฺโท สตฺต วสฺสานิ อาคเมตุ, สตฺตนฺนํ วสฺสานํ อจฺจเยน มยมฺปิ อคารสฺมา อนคาริยํ ปพฺพชิสฺสาม, อถ ยา เต คติ, สา โน คติ ภวิสฺสตีติ.}

\prose{อติจิรํ โข โภ สตฺต วสฺสานิ, นาหํ สกฺโกมิ ภวนฺเต สตฺต วสฺสานิ อาคเมตุํ, โก นุ โข ปน โภ ชานาติ ชีวิตานํ.\csromanspacer{} คมนีโย สมฺปราโย, มนฺตายํ\footnote{มนฺตาย (พหูสุ)} โพทฺธพฺพํ, กตฺตพฺพํ กุสลํ, จริตพฺพํ พฺรหฺมจริยํ, นตฺถิ ชาตสฺส อมรณํ.\csromanspacer{} ยถา โข ปน เม สุตํ พฺรหฺมุโน อามคนฺเธ ภาสมานสฺส, เต น สุนิมฺมทยา อคารํ อชฺฌาวสตา, ปพฺพชิสฺสามหํ โภ อคารสฺมา อนคาริยนฺติ.\csromanspacer{} เตน หิ ภวํ โควินฺโท ฉพฺพสฺสานิ อาคเมตุ ฯเปฯ ปญฺจ วสฺสานิ อาคเมตุ.\csromanspacer{} จตฺตาริ วสฺสานิ อาคเมตุ.\csromanspacer{} ตีณิ วสฺสานิ อาคเมตุ.\csromanspacer{} เทฺว วสฺสานิ อาคเมตุ.\csromanspacer{} เอกํ \csromanfolio{198}วสฺสํ อาคเมตุ, เอกสฺส วสฺสสฺส อจฺจเยน มยมฺปิ อคารสฺมา อนคาริยํ ปพฺพชิสฺสาม, อถ ยา เต คติ, สา โน คติ ภวิสฺสตีติ.}

\prose{อติจิรํ โข โภ เอกํ วสฺสํ, นาหํ สกฺโกมิ ภวนฺเต เอกํ วสฺสํ อาคเมตุํ, โก นุ โข ปน โภ ชานาติ ชีวิตานํ.\csromanspacer{} คมนีโย สมฺปราโย, มนฺตายํ โพทฺธพฺพํ, กตฺตพฺพํ กุสลํ, จริตพฺพํ พฺรหฺมจริยํ, นตฺถิ ชาตสฺส อมรณํ.\csromanspacer{} ยถา โข ปน เม สุตํ พฺรหฺมุโน อามคนฺเธ ภาสมานสฺส, เต น สุนิมฺมทยา อคารํ อชฺฌาวสตา, ปพฺพชิสฺสามหํ โภ อคารสฺมา อนคาริยนฺติ.\csromanspacer{} เตน หิ ภวํ โควินฺโท สตฺต มาสานิ อาคเมตุ, สตฺตนฺนํ มาสานํ อจฺจเยน มยมฺปิ อคารสฺมา อนคาริยํ ปพฺพชิสฺสาม, อถ ยา เต คติ, สา โน คติ ภวิสฺสตีติ.}

\prose{อติจิรํ โข โภ สตฺต มาสานิ, นาหํ สกฺโกมิ ภวนฺเต สตฺต มาสานิ อาคเมตุํ, โก นุ โข ปน โภ ชานาติ ชีวิตานํ.\csromanspacer{} คมนีโย สมฺปราโย, มนฺตายํ โพทฺธพฺพํ, กตฺตพฺพํ กุสลํ, จริตพฺพํ พฺรหฺมจริยํ, นตฺถิ ชาตสฺส อมรณํ.\csromanspacer{} ยถา โข ปน เม สุตํ พฺรหฺมุโน อามคนฺเธ ภาสมานสฺส, เต น สุนิมฺมทยา อคารํ อชฺฌาวสตา, ปพฺพชิสฺสามหํ โภ อคารสฺมา อนคาริยนฺติ.\csromanspacer{} เตน หิ ภวํ โควินฺโท ฉ มาสานิ อาคเมตุ ฯเปฯ ปญฺจ มาสานิ อาคเมตุ.\csromanspacer{} จตฺตาริ มาสานิ อาคเมตุ.\csromanspacer{} ตีณิ มาสานิ อาคเมตุ.\csromanspacer{} เทฺว มาสานิ อาคเมตุ.\csromanspacer{} เอกํ มาสํ อาคเมตุ.\csromanspacer{} อทฺธมาสํ อาคเมตุ, อทฺธมาสสฺส อจฺจเยน มยมฺปิ อคารสฺมา อนคาริยํ ปพฺพชิสฺสาม, อถ ยา เต คติ, สา โน คติ ภวิสฺสตีติ.}

\prose{อติจิรํ โข โภ อทฺธมาโส, นาหํ สกฺโกมิ ภวนฺเต อทฺธมาสํ อาคเมตุํ, โก นุ โข ปน โภ ชานาติ ชีวิตานํ.\csromanspacer{} คมนีโย สมฺปราโย, มนฺตายํ โพทฺธพฺพํ, กตฺตพฺพํ กุสลํ, จริตพฺพํ พฺรหฺมจริยํ, นตฺถิ ชาตสฺส อมรณํ.\csromanspacer{} ยถา โข ปน เม สุตํ พฺรหฺมุโน อามคนฺเธ ภาสมานสฺส, เต น สุนิมฺมทยา อคารํ อชฺฌาวสตา, ปพฺพชิสฺสามหํ โภ อคารสฺมา อนคาริยนฺติ.\csromanspacer{} เตน หิ ภวํ โควินฺโท สตฺตาหํ อาคเมตุ ยาว มยํ สเก ปุตฺตภาตโร รชฺเชน\footnote{รชฺเช (สฺยา)} อนุสาสิสฺสาม,}

\vspace{\tipitakaparskip}
\csromancenter{มหาโควินฺทสุตฺต สตฺตาหสฺส อจฺจเยน มยมฺปิ อคารสฺมา อนคาริยํ ปพฺพชิสฺสาม, อถ ยา เต คติ, สา โน คติ ภวิสฺสตีติ.\csromanspacer{} น จิรํ โข โภ สตฺตาหํ, อาคเมสฺสามหํ ภวนฺเต สตฺตาหนฺติ.}
\csromanmatikaanchor{mtk.104}
\csromantocmarkref{section}{พฺราหฺมณมหาสาลาทีนํอามนฺตนา}{mtk.104}
\bookmark[dest={mtk.104},level=1]{พฺราหฺมณมหาสาลาทีนํอามนฺตนา}
\csromanheader{1}{\textbf{พฺราหฺมณมหาสาลาทีนํอามนฺตนา}}
\proseitem{324}{\csromanfolio{199}อถ โข โภ มหาโควินฺโท พฺราหฺมโณ เยน เต สตฺต จ พฺราหฺมณมหาสาลา สตฺต จ นฺหาตกสตานิ เตนุปสงฺกมิ, อุปสงฺกมิตฺวา เต สตฺต จ พฺราหฺมณมหาสาเล สตฺต จ นฺหาตกสตานิ เอตทโวจ “อญฺญํ ทานิ ภวนฺโต อาจริยํ ปริเยสนฺตุ, โย ภวนฺตานํ มนฺเต วาเจสฺสติ.\csromanspacer{} อิจฺฉามหํ โภ อคารสฺมา อนคาริยํ ปพฺพชิตุํ.\csromanspacer{} ยถา โข ปน เม สุตํ พฺรหฺมุโน อามคนฺเธ ภาสมานสฺส, เต น สุนิมฺมทยา อคารํ อชฺฌาวสตา, ปพฺพชิสฺสามหํ โภ อคารสฺมา อนคาริยนฺ”ติ.\csromanspacer{} มา ภวํ โควินฺโท อคารสฺมา อนคาริยํ ปพฺพชิ.\csromanspacer{} ปพฺพชฺชา โภ อปฺเปสกฺขา จ อปฺปลาภา จ, พฺรหฺมญฺญํ มเหสกฺขญฺจ มหาลาภญฺจาติ.\csromanspacer{} มา ภวนฺโต เอวํ อวจุตฺถ “ปพฺพชฺชา อปฺเปสกฺขา จ อปฺปลาภา จ, พฺรหฺมญฺญํ มเหสกฺขญฺจ มหาลาภญฺจา”ติ.\csromanspacer{} โก นุ โข โภ อญฺญตฺร มยา มเหสกฺขตโร วา มหาลาภตโร วา.\csromanspacer{} อหํ หิ โภ เอตรหิ ราชาว รญฺญํ พฺรหฺมาว พฺราหฺมณานํ\footnote{พฺรหฺมานํ (สี, อิ, ก)} เทวตาว คหปติกานํ.\csromanspacer{} ตมหํ สพฺพํ ปหาย อคารสฺมา อนคาริยํ ปพฺพชิสฺสามิ.\csromanspacer{} ยถา โข ปน เม สุตํ พฺรหฺมุโน อามคนฺเธ ภาสมานสฺส, เต น สุนิมฺมทยา อคารํ อชฺฌาวสตา.\csromanspacer{} ปพฺพชิสฺสามหํ โภ อคารสฺมา อนคาริยนฺติ.\csromanspacer{} สเจ ภวํ โควินฺโท อคารสฺมา อนคาริยํ ปพฺพชิสฺสติ, มยมฺปิ อคารสฺมา อนคาริยํ ปพฺพชิสฺสาม, อถ ยา เต คติ, สา โน คติ ภวิสฺสตีติ.}

\csromanmatikaanchor{mtk.105}
\csromantocmarkref{section}{ภริยานํอามนฺตนา}{mtk.105}
\bookmark[dest={mtk.105},level=1]{ภริยานํอามนฺตนา}
\csromanheader{1}{\textbf{ภริยานํอามนฺตนา}}
\proseitem{325}{อถ โข โภ มหาโควินฺโท พฺราหฺมโณ เยน จตฺตารีสา ภริยา สาทิสิโย เตนุปสงฺกมิ, อุปสงฺกมิตฺวา จตฺตารีสา ภริยา สาทิสิโย เอตทโวจ “ยา โภตีนํ อิจฺฉติ, สกานิ วา ญาติกุลานิ คจฺฉตุ อญฺญํ วา ภตฺตารํ ปริเยสตุ.\csromanspacer{} อิจฺฉามหํ โภตี \csromanfolio{200}อคารสฺมา อนคาริยํ ปพฺพชิตุํ.\csromanspacer{} ยถา โข ปน เม สุตํ พฺรหฺมุโน อามคนฺเธ ภาสมานสฺส, เต น สุนิมฺมทยา อคารํ อชฺฌาวสตา.\csromanspacer{} ปพฺพชิสฺสามหํ โภตี อคารสฺมา อนคาริยนฺ”ติ.\csromanspacer{} ตฺวญฺเญว โน ญาติ ญาติกามานํ, ตฺวํ ปน ภตฺตา ภตฺตุกามานํ.\csromanspacer{} สเจ ภวํ โควินฺโท อคารสฺมา อนคาริยํ ปพฺพชิสฺสติ, มยมฺปิ อคารสฺมา อนคาริยํ ปพฺพชิสฺสาม, อถ ยา เต คติ, สา โน คติ ภวิสฺสตีติ.}

\csromanmatikaanchor{mtk.106}
\csromantocmarkref{section}{มหาโควินฺทปพฺพชฺชา}{mtk.106}
\bookmark[dest={mtk.106},level=1]{มหาโควินฺทปพฺพชฺชา}
\csromanheader{1}{\textbf{มหาโควินฺทปพฺพชฺชา}}
\proseitem{326}{อถ โข โภ มหาโควินฺโท พฺราหฺมโณ ตสฺส สตฺตาหสฺส อจฺจเยน เกสมสฺสุํ โอหาเรตฺวา กาสายานิ วตฺถานิ อจฺฉาเทตฺวา อคารสฺมา อนคาริยํ ปพฺพชิ.\csromanspacer{} ปพฺพชิตํ ปน มหาโควินฺทํ พฺราหฺมณํ สตฺต จ ราชาโน ขตฺติยา มุทฺธาวสิตฺตา สตฺต จ พฺราหฺมณมหาสาลา สตฺต จ นฺหาตกสตานิ จตฺตารีสา จ ภริยา สาทิสิโย อเนกานิ จ ขตฺติยสหสฺสานิ อเนกานิ จ พฺราหฺมณสหสฺสานิ อเนกานิ จ คหปติสหสฺสานิ อเนเกหิ จ อิตฺถาคาเรหิ อิตฺถิโย เกสมสฺสุํ โอหาเรตฺวา กาสายานิ วตฺถานิ อจฺฉาเทตฺวา มหาโควินฺทํ พฺราหฺมณํ อคารสฺมา อนคาริยํ ปพฺพชิตํ อนุปพฺพชิํสุ.\csromanspacer{} ตาย สุทํ โภ ปริสาย ปริวุโต มหาโควินฺโท พฺราหฺมโณ คามนิคมราชธานีสุ จาริกํ จรติ.\csromanspacer{} ยํ โข ปน โภ เตน สมเยน มหาโควินฺโท พฺราหฺมโณ คามํ วา นิคมํ วา อุปสงฺกมติ.\csromanspacer{} ตตฺถ ราชาว โหติ รญฺญํ พฺรหฺมาว พฺราหฺมณานํ เทวตาว คหปติกานํ.\csromanspacer{} เตน โข ปน สมเยน มนุสฺสา ขิปนฺติ วา อุปกฺขลนฺติ วา.\csromanspacer{} เต เอวมาหํสุ “นมตฺถุ มหาโควินฺทสฺส พฺราหฺมณสฺส, นมตฺถุ สตฺต ปุโรหิตสฺสา”ติ.}

\proseitem{327}{มหาโควินฺโท โภ พฺราหฺมโณ เมตฺตาสหคเตน เจตสา เอกํ ทิสํ ผริตฺวา วิหาสิ.\csromanspacer{} ตถา ทุติยํ.\csromanspacer{} ตถา ตติยํ.\csromanspacer{} ตถา จตุตฺถํ.\csromanspacer{} อิติ อุทฺธมโธ ติริยํ สพฺพธิ สพฺพตฺตตาย สพฺพาวนฺตํ โลกํ เมตฺตาสหคเตน เจตสา วิปุเลน มหคฺคเตน อปฺปมาเณน อเวเรน อพฺยาปชฺเชน ผริตฺวา วิหาสิ.\csromanspacer{} กรุณาสหคเตน เจตสา ฯเปฯ.\csromanspacer{} มุทิตาสหคเตน เจตสา ฯเปฯ.\csromanspacer{} อุเปกฺขาสหคเตน เจตสา ฯเปฯ อพฺยาปชฺเชน ผริตฺวา วิหาสิ.\csromanspacer{} สาวกานํ จ พฺรหฺมโลกสหพฺยตาย มคฺคํ เทเสสิ.}

\csromanheader{2}{6. มหาโควินฺทสุตฺต}
\proseitem{328}{\csromanfolio{201}เย โข ปน โภ เตน สมเยน มหาโควินฺทสฺส พฺราหฺมณสฺส สาวกา สพฺเพนสพฺพํ สาสนํ อาชานิํสุ.\csromanspacer{} เต กายสฺส เภทา ปรํ มรณา สุคติํ พฺรหฺมโลกํ อุปปชฺชิํสุ.\csromanspacer{} เย น สพฺเพนสพฺพํ สาสนํ อาชานิํสุ.\csromanspacer{} เต กายสฺส เภทา ปรํ มรณา อปฺเปกจฺเจ ปรนิมฺมิตวสวตฺตีนํ เทวานํ สหพฺยตํ อุปปชฺชิํสุ.\csromanspacer{} อปฺเปกจฺเจ นิมฺมานรตีนํ เทวานํ สหพฺยตํ อุปปชฺชิํสุ.\csromanspacer{} อปฺเปกจฺเจ ตุสิตานํ เทวานํ สหพฺยตํ อุปปชฺชิํสุ.\csromanspacer{} อปฺเปกจฺเจ ยามานํ เทวานํ สหพฺยตํ อุปปชฺชิํสุ.\csromanspacer{} อปฺเปกจฺเจ ตาวติํสานํ เทวานํ สหพฺยตํ อุปปชฺชิํสุ.\csromanspacer{} อปฺเปกจฺเจ จาตุมหาราชิกานํ เทวานํ สหพฺยตํ อุปปชฺชิํสุ.\csromanspacer{} เย สพฺพนิหีนํ กายํ ปริปูเรสุํ, เต คนฺธพฺพกายํ ปริปูเรสุํ.\csromanspacer{} อิติ โข โภ\footnote{ปน (สฺยา, ก)} สพฺเพสํเยว เตสํ กุลปุตฺตานํ อโมฆา ปพฺพชฺชา อโหสิ อวญฺฌา สผลา ส-อุทฺรยาติ.}

\proseitem{329}{สรติ ตํ ภควาติ.\csromanspacer{} สรามหํ ปญฺจสิข.\csromanspacer{} อหํ เตน สมเยน มหาโควินฺโท พฺราหฺมโณ อโหสิํ.\csromanspacer{} อหํ เตสํ สาวกานํ พฺรหฺมโลกสหพฺยตาย มคฺคํ เทเสสิํ.\csromanspacer{} ตํ โข ปน เม ปญฺจสิข พฺรหฺมจริยํ น นิพฺพิทาย น วิราคาย น นิโรธาย น อุปสมาย น อภิญฺญาย น สมฺโพธาย น นิพฺพานาย สํวตฺตติ ยาวเทว พฺรหฺมโลกูปปตฺติยา.}

\prose{อิทํ โข ปน เม ปญฺจสิข พฺรหฺมจริยํ เอกนฺตนิพฺพิทาย วิราคาย นิโรธาย อุปสมาย อภิญฺญาย สมฺโพธาย นิพฺพานาย สํวตฺตติ.\csromanspacer{} กตมญฺจ ตํ ปญฺจสิข พฺรหฺมจริยํ เอกนฺตนิพฺพิทาย วิราคาย นิโรธาย อุปสมาย อภิญฺญาย สมฺโพธาย นิพฺพานาย สํวตฺตติ.\csromanspacer{} อยเมว อริโย อฏฺฐงฺคิโก มคฺโค.\csromanspacer{} เสยฺยถิทํ, สมฺมาทิฏฺฐิ สมฺมาสงฺกปฺโป สมฺมาวาจา สมฺมากมฺมนฺโต สมฺมา-อาชีโว สมฺมาวายาโม สมฺมาสติ สมฺมาสมาธิ.\csromanspacer{} อิทํ โข ตํ ปญฺจสิข พฺรหฺมจริยํ เอกนฺตนิพฺพิทาย วิราคาย นิโรธาย อุปสมาย อภิญฺญาย สมฺโพธาย นิพฺพานาย สํวตฺตติ.}

\proseitem{330}{เย โข ปน เม ปญฺจสิข สาวกา สพฺเพนสพฺพํ สาสนํ อาชานนฺติ.\csromanspacer{} เต อาสวานํ ขยา อนาสวํ เจโตวิมุตฺติํ ปญฺญาวิมุตฺติํ ทิฏฺเฐวธมฺเม สยํ อภิญฺญา สจฺฉิกตฺวา อุปสมฺปชฺช วิหรนฺติ.\csromanspacer{} เย น สพฺเพนสพฺพํ สาสนํ อาชานนฺติ.\csromanspacer{} เต ปญฺจนฺนํ โอรมฺภาคิยานํ สํโยชนานํ \csromanfolio{202}ปริกฺขยา โอปปาติกา โหนฺติ ตตฺถ ปรินิพฺพายิโน อนาวตฺติธมฺมา ตสฺมา โลกา.\csromanspacer{} เย น สพฺเพนสพฺพํ สาสนํ อาชานนฺติ.\csromanspacer{} อปฺเปกจฺเจ ติณฺณํ สํโยชนานํ ปริกฺขยา ราคโทสโมหานํ ตนุตฺตา สกทาคามิโน โหนฺติ สกิเทว อิมํ โลกํ อาคนฺตฺวา ทุกฺขสฺสนฺตํ กริสฺสนฺติ\footnote{กโรนฺติ (สี, อิ)}.\csromanspacer{} เย น สพฺเพนสพฺพํ สาสนํ อาชานนฺติ.\csromanspacer{} อปฺเปกจฺเจ ติณฺณํ สํโยชนานํ ปริกฺขยา โสตาปนฺนา โหนฺติ อวินิปาตธมฺมา นิยตา สมฺโพธิปรายณา.\csromanspacer{} อิติ โข ปญฺจสิข สพฺเพสํเยว อิเมสํ กุลปุตฺตานํ อโมฆา ปพฺพชฺชา\footnote{ปพฺพชฺชา อโหสิ (ก)} อวญฺฌา สผลา ส-อุทฺรยาติ.}

\prose{อิทมโวจ ภควา.\csromanspacer{} อตฺตมโน ปญฺจสิโข คนฺธพฺพปุตฺโต ภควโต ภาสิตํ อภินนฺทิตฺวา อนุโมทิตฺวา ภควนฺตํ อภิวาเทตฺวา ปทกฺขิณํ กตฺวา ตตฺเถวนฺตรธายีติ.}

\nitthitam{\textbf{มหาโควินฺทสุตฺตํ นิฏฺฐิตํ ฉฏฺฐํ.}}
\csromanmatikaanchor{mtk.107}
\csromantocmarknopage{chapter}{7. มหาสมยสุตฺต}{mtk.107}
\bookmark[dest={mtk.107},level=0]{7. มหาสมยสุตฺต}
\chapterheadpageread{7. มหาสมยสุตฺต}
\proseitem{331}{\csromanfolio{203}เอวํ เม สุตํ\csromandash{}เอกํ สมยํ ภควา สกฺเกสุ วิหรติ กปิลวตฺถุสฺมิํ มหาวเน มหตา ภิกฺขุสํเฆน สทฺธิํ ปญฺจมตฺเตหิ ภิกฺขุสเตหิ สพฺเพเหว อรหนฺเตหิ, ทสหิ จ โลกธาตูหิ เทวตา เยภุยฺเยน สนฺนิปติตา โหนฺติ ภควนฺตํ ทสฺสนาย ภิกฺขุสํฆญฺจ.\csromanspacer{} อถ โข จตุนฺนํ สุทฺธาวาสกายิกานํ เทวตานํ\footnote{เทวานํ (สี, สฺยา, อิ)} เอตทโหสิ “อยํ โข ภควา สกฺเกสุ วิหรติ กปิลวตฺถุสฺมิํ มหาวเน มหตา ภิกฺขุสํเฆน สทฺธิํ ปญฺจมตฺเตหิ ภิกฺขุสเตหิ สพฺเพเหว อรหนฺเตหิ, ทสหิ จ โลกธาตูหิ เทวตา เยภุยฺเยน สนฺนิปติตา โหนฺติ ภควนฺตํ ทสฺสนาย ภิกฺขุสํฆญฺจ.\csromanspacer{} ยํนูน มยมฺปิ เยน ภควา เตนุปสงฺกเมยฺยาม, อุปสงฺกมิตฺวา ภควโต สนฺติเก ปจฺเจกํ คาถํ\footnote{ปจฺเจกคาถํ (สี, สฺยา, อิ), ปจฺเจกคาถา (ก-สี) 3. อุชุก มกํสุ (สี, สฺยา, อิ)} ภาเสยฺยามา”ติ.}

\proseitem{332}{อถ โข ตา เทวตา เสยฺยถาปิ นาม พลวา ปุริโส สมิญฺชิตํ วา พาหํ ปสาเรยฺย ปสาริตํ วา พาหํ สมิญฺเชยฺย.\csromanspacer{} เอวเมว สุทฺธาวาเสสุ เทเวสุ อนฺตรหิตา ภควโต ปุรโต ปาตุรเหสุํ.\csromanspacer{} อถ โข ตา เทวตา ภควนฺตํ อภิวาเทตฺวา เอกมนฺตํ อฏฺฐํสุ.\csromanspacer{} เอกมนฺตํ ฐิตา โข เอกา เทวตา ภควโต สนฺติเก อิมํ คาถํ อภาสิ.}

\prose{“มหาสมโย ปวนสฺมิํ, เทวกายา สมาคตา.}

\prose{อาคตมฺห อิมํ ธมฺมสมยํ, ทกฺขิตาเย อปราชิตสํฆนฺ”ติ.}

\prose{อถ โข อปรา เทวตา ภควโต สนฺติเก อิมํ คาถํ อภาสิ.}

\prose{“ตตฺร ภิกฺขโว สมาทหํสุ, จิตฺตมตฺตโน อุชุกํ อกํสุ3.}

\prose{สารถีว เนตฺตานิ คเหตฺวา, อินฺทฺริยานิ รกฺขนฺติ ปณฺฑิตา”ติ.}

\prose{อถ โข อปรา เทวตา ภควโต สนฺติเก อิมํ คาถํ อภาสิ.}

\prose{“เฉตฺวา ขีลํ เฉตฺวา ปลิฆํ, อินฺทขีลํ อูหจฺจ\footnote{อุหจฺจ (ก)} มเนชา.}

\prose{เต จรนฺติ สุทฺธา วิมลา, จกฺขุมตา สุทนฺตา สุสุนาคา”ติ.}

\prose{\csromanfolio{204}อถ โข อปรา เทวตา ภควโต สนฺติเก อิมํ คาถํ อภาสิ.}

\prose{“เยเกจิ พุทฺธํ สรณํ คตาเส, น เต คมิสฺสนฺติ อปายภูมิํ.}

\prose{ปหาย มานุสํ เทหํ, เทวกายํ ปริปูเรสฺสนฺตี”ติ.}

\csromanmatikaanchor{mtk.108}
\csromantocmarkref{section}{เทวตาสนฺนิปาต}{mtk.108}
\bookmark[dest={mtk.108},level=1]{เทวตาสนฺนิปาต}
\csromanheader{1}{\textbf{เทวตาสนฺนิปาต}}
\proseitem{333}{อถ โข ภควา ภิกฺขู อามนฺเตสิ “เยภุยฺเยน ภิกฺขเว ทสสุ โลกธาตูสุ เทวตา สนฺนิปติตา (โหนฺติ)\footnote{( ) สี-อิ-โปตฺถเกสุ นตฺถิ.} ตถาคตํ ทสฺสนาย ภิกฺขุสํฆญฺจ.\csromanspacer{} เยปิ เต ภิกฺขเว อเหสุํ อตีตมทฺธานํ อรหนฺโต สมฺมาสมฺพุทฺธา.\csromanspacer{} เตสมฺปิ ภควนฺตานํ เอตํปรมาเยว\footnote{เอตปรมาเยว (สี, สฺยา, อิ)} เทวตา สนฺนิปติตา อเหสุํ เสยฺยถาปิ มยฺหํ เอตรหิ.\csromanspacer{} เยปิ เต ภิกฺขเว ภวิสฺสนฺติ อนาคตมทฺธานํ อรหนฺโต สมฺมาสมฺพุทฺธา, เตสมฺปิ ภควนฺตานํ เอตํปรมาเยว เทวตา สนฺนิปติตา ภวิสฺสนฺติ เสยฺยถาปิ มยฺหํ เอตรหิ.\csromanspacer{} อาจิกฺขิสฺสามิ ภิกฺขเว เทวกายานํ นามานิ, กิตฺตยิสฺสามิ ภิกฺขเว เทวกายานํ นามานิ, เทเสสฺสามิ ภิกฺขเว เทวกายานํ นามานิ, ตํ สุณาถ สาธุกํ มนสิกโรถ ภาสิสฺสามี”ติ. “เอวํ ภนฺเต”ติ โข เต ภิกฺขู ภควโต ปจฺจสฺโสสุํ.}

\csromanheader{2}{334. ภควา เอตทโวจ\csromandash{}}
\csromangathameasure{สิโลกมนุกสฺสามิ, ยตฺถ ภุมฺมา ตทสฺสิตา. \\ เย สิตา คิริคพฺภรํ, ปหิตตฺตา สมาหิตา. \\ ปุถูสีหาว สลฺลีนา, โลมหํสาภิสมฺภุโน. \\ โอทาตมนสา สุทฺธา, วิปฺปสนฺนมนาวิลา. \\ ภิยฺโย ปญฺจสเต ญตฺวา, วเน กาปิลวตฺถเว. \\ ตโต อามนฺตยี สตฺถา, สาวเก สาสเน รเต. \\ เทวกายา อภิกฺกนฺตา, เต วิชานาถ ภิกฺขโว. \\ เต จ อาตปฺปมกรุํ, สุตฺวา พุทฺธสฺส สาสนํ. \\ เตสํ ปาตุรหุ ญาณํ, อมนุสฺสานทสฺสนํ. \\ อปฺเปเก สตมทฺทกฺขุํ, สหสฺสํ อถ สตฺตริํ.}
\csromangathagroup{\csromangathastanza{สิโลกมนุกสฺสามิ, ยตฺถ ภุมฺมา ตทสฺสิตา. \\ เย สิตา คิริคพฺภรํ, ปหิตตฺตา สมาหิตา.}\csromangathastanza{ปุถูสีหาว สลฺลีนา, โลมหํสาภิสมฺภุโน. \\ โอทาตมนสา สุทฺธา, วิปฺปสนฺนมนาวิลา\footnote{วิปฺปสนฺนามนาวิลา (อิ, ก)}.}\csromangathastanza{ภิยฺโย ปญฺจสเต ญตฺวา, วเน กาปิลวตฺถเว. \\ ตโต อามนฺตยี สตฺถา, สาวเก สาสเน รเต.}\csromangathastanza{เทวกายา อภิกฺกนฺตา, เต วิชานาถ ภิกฺขโว. \\ เต จ อาตปฺปมกรุํ, สุตฺวา พุทฺธสฺส สาสนํ.}\csromangathastanza{เตสํ ปาตุรหุ ญาณํ, อมนุสฺสานทสฺสนํ. \\ อปฺเปเก สตมทฺทกฺขุํ, สหสฺสํ อถ สตฺตริํ.}}

\csromanheader{2}{7. มหาสมยสุตฺต}
\csromangathameasure{สตํ เอเก สหสฺสานํ, อมนุสฺสานมทฺทสุํ. \\ อปฺเปเกนนฺตมทฺทกฺขุํ, ทิสา สพฺพา ผุฏา อหุํ. \\ ตญฺจ สพฺพํ อภิญฺญาย, ววตฺถิตฺวาน จกฺขุมา. \\ ตโต อามนฺตยี สตฺถา, สาวเก สาสเน รเต. \\ เทวกายา อภิกฺกนฺตา, เต วิชานาถ ภิกฺขโว. \\ เย โวหํ กิตฺตยิสฺสามิ, คิราหิ อนุปุพฺพโส.}
\csromangathagroup{\csromangathastanza{\csromanfolio{205}สตํ เอเก สหสฺสานํ, อมนุสฺสานมทฺทสุํ. \\ อปฺเปเกนนฺตมทฺทกฺขุํ, ทิสา สพฺพา ผุฏา อหุํ.}\csromangathastanza{ตญฺจ สพฺพํ อภิญฺญาย, ววตฺถิตฺวาน\footnote{ววกฺขิตฺวาน (สี, สฺยา, อิ), อเวกฺขิตฺวาน (ฏีกา)} จกฺขุมา. \\ ตโต อามนฺตยี สตฺถา, สาวเก สาสเน รเต.}\csromangathastanza{เทวกายา อภิกฺกนฺตา, เต วิชานาถ ภิกฺขโว. \\ เย โวหํ กิตฺตยิสฺสามิ, คิราหิ อนุปุพฺพโส.}}

\csromanhangingbegin
\hangingheaditem{335}{สตฺตสหสฺสา เต ยกฺขา, ภุมฺมา กาปิลวตฺถวา.}
\hangingbody{อิทฺธิมนฺโต ชุติมนฺโต, วณฺณวนฺโต ยสสฺสิโน.}
\hangingbody{โมทมานา อภิกฺกามุํ, ภิกฺขูนํ สมิติํ วนํ.}
\hangingbody{ฉสหสฺสา เหมวตา, ยกฺขา นานตฺตวณฺณิโน.}
\hangingbody{อิทฺธิมนฺโต ชุติมนฺโต2, วณฺณวนฺโต ยสสฺสิโน.}
\hangingbody{โมทมานา อภิกฺกามุํ, ภิกฺขูนํ สมิติํ วนํ.}
\hangingbody{สาตาคิรา ติสหสฺสา, ยกฺขา นานตฺตวณฺณิโน.}
\hangingbody{อิทฺธิมนฺโต ชุติมนฺโต, วณฺณวนฺโต ยสสฺสิโน.}
\hangingbody{โมทมานา อภิกฺกามุํ, ภิกฺขูนํ สมิติํ วนํ.}
\hangingbody{อิจฺเจเต โสฬสสหสฺสา, ยกฺขา นานตฺตวณฺณิโน.}
\hangingbody{อิทฺธิมนฺโต ชุติมนฺโต, วณฺณวนฺโต ยสสฺสิโน.}
\hangingbody{โมทมานา อภิกฺกามุํ, ภิกฺขูนํ สมิติํ วนํ.}
\hangingbody{เวสฺสามิตฺตา ปญฺจสตา, ยกฺขา นานตฺตวณฺณิโน.}
\hangingbody{อิทฺธิมนฺโต ชุติมนฺโต, วณฺณวนฺโต ยสสฺสิโน.}
\hangingbody{โมทมานา อภิกฺกามุํ, ภิกฺขูนํ สมิติํ วนํ.}
\hangingbody{กุมฺภีโร ราชคหิโก, เวปุลฺลสฺส นิเวสนํ.}
\hangingbody{ภิยฺโย นํ สตสหสฺสํ, ยกฺขานํ ปยิรุปาสติ.}
\hangingbody{กุมฺภีโร ราชคหิโก, โสปาคา สมิติํ วนํ.}
\csromanhangingend

\csromanhangingbegin
\hangingheaditem{336}{ปุริมญฺจ ทิสํ ราชา, ธตรฏฺโฐ ปสาสติ.}
\hangingbody{คนฺธพฺพานํ อธิปติ, มหาราชา ยสสฺสิโส.}
\csromanhangingend

\csromangathameasure{ปุตฺตาปิ ตสฺส พหโว, อินฺทนามา มหพฺพลา. \\ อิทฺธิมนฺโต ชุติมนฺโต, วณฺณวนฺโต ยสสฺสิโน. \\ โมทมานา อภิกฺกามุํ, ภิกฺขูนํ สมิติํ วนํ. \\ ทกฺขิณญฺจ ทิสํ ราชา, วิรูโฬฺห ตํ ปสาสติ. \\ กุมฺภณฺฑานํ อธิปติ, มหาราชา ยสสฺสิโส. \\ ปุตฺตาปิ ตสฺส พหโว, อินฺทนามา มหพฺพลา. \\ อิทฺธิมนฺโต ชุติมนฺโต, วณฺณวนฺโต ยสสฺสิโน. \\ โมทมานา อภิกฺกามุํ, ภิกฺขูนํ สมิติํ วนํ. \\ ปจฺฉิมญฺจ ทิสํ ราชา, วิรูปกฺโข ปสาสติ. \\ นาคานญฺจ อธิปติ, มหาราชา ยสสฺสิโส. \\ ปุตฺตาปิ ตสฺส พหโว, อินฺทนามา มหพฺพลา. \\ อิทฺธิมนฺโต ชุติมนฺโต, วณฺณวนฺโต ยสสฺสิโน. \\ โมทมานา อภิกฺกามุํ, ภิกฺขูนํ สมิติํ วนํ. \\ อุตฺตรญฺจ ทิสํ ราชา, กุเวโร ตํ ปสาสติ. \\ ยกฺขานญฺจ อธิปติ, มหาราชา ยสสฺสิโส. \\ ปุตฺตาปิ ตสฺส พหโว, อินฺทนามา มหพฺพลา. \\ อิทฺธิมนฺโต ชุติมนฺโต, วณฺณวนฺโต ยสสฺสิโน. \\ โมทมานา อภิกฺกามุํ, ภิกฺขูนํ สมิติํ วนํ. \\ ปุริมํ ทิสํ ธตรฏฺโฐ, ทกฺขิเณน วิรูฬฺหโก. \\ ปจฺฉิเมน วิรูปกฺโข, กุเวโร อุตฺตรํ ทิสํ. \\ จตฺตาโร เต มหาราชา, สมนฺตา จตุโร ทิสา. \\ ททฺทลฺลมานา อฏฺฐํสุ, วเน กาปิลวตฺถเว.}
\csromangathagroup{\csromangathastanza{\csromanfolio{206}ปุตฺตาปิ ตสฺส พหโว, อินฺทนามา มหพฺพลา. \\ อิทฺธิมนฺโต ชุติมนฺโต, วณฺณวนฺโต ยสสฺสิโน.}\csromangathastanza{โมทมานา อภิกฺกามุํ, ภิกฺขูนํ สมิติํ วนํ. \\ ทกฺขิณญฺจ ทิสํ ราชา, วิรูโฬฺห ตํ ปสาสติ\footnote{ตปฺปสาสติ (สฺยา)}.}\csromangathastanza{กุมฺภณฺฑานํ อธิปติ, มหาราชา ยสสฺสิโส. \\ ปุตฺตาปิ ตสฺส พหโว, อินฺทนามา มหพฺพลา.}\csromangathastanza{อิทฺธิมนฺโต ชุติมนฺโต, วณฺณวนฺโต ยสสฺสิโน. \\ โมทมานา อภิกฺกามุํ, ภิกฺขูนํ สมิติํ วนํ.}\csromangathastanza{ปจฺฉิมญฺจ ทิสํ ราชา, วิรูปกฺโข ปสาสติ. \\ นาคานญฺจ อธิปติ, มหาราชา ยสสฺสิโส.}\csromangathastanza{ปุตฺตาปิ ตสฺส พหโว, อินฺทนามา มหพฺพลา. \\ อิทฺธิมนฺโต ชุติมนฺโต, วณฺณวนฺโต ยสสฺสิโน.}\csromangathastanza{โมทมานา อภิกฺกามุํ, ภิกฺขูนํ สมิติํ วนํ. \\ อุตฺตรญฺจ ทิสํ ราชา, กุเวโร ตํ ปสาสติ.}\csromangathastanza{ยกฺขานญฺจ อธิปติ, มหาราชา ยสสฺสิโส. \\ ปุตฺตาปิ ตสฺส พหโว, อินฺทนามา มหพฺพลา.}\csromangathastanza{อิทฺธิมนฺโต ชุติมนฺโต, วณฺณวนฺโต ยสสฺสิโน. \\ โมทมานา อภิกฺกามุํ, ภิกฺขูนํ สมิติํ วนํ.}\csromangathastanza{ปุริมํ ทิสํ ธตรฏฺโฐ, ทกฺขิเณน วิรูฬฺหโก. \\ ปจฺฉิเมน วิรูปกฺโข, กุเวโร อุตฺตรํ ทิสํ.}\csromangathastanza{จตฺตาโร เต มหาราชา, สมนฺตา จตุโร ทิสา. \\ ททฺทลฺลมานา\footnote{ททฺทฬฺหมานา (ก)} อฏฺฐํสุ, วเน กาปิลวตฺถเว.}}

\csromanhangingbegin
\hangingheaditem{337}{เตสํ มายาวิโน ทาสา, อาคุํ\footnote{อาคู (สฺยา), อาคุ (สี, อิ) เอวมุปริปิ.} วญฺจนิกา สฐา.}
\hangingbody{มายา กุเฏณฺฑุ วิเฏณฺฑุ4, วิฏุจฺจ5 วิฏุโฏ สห.}
\hangingbody{จนฺทโน กามเสฏฺโฐ จ, กินฺนิฆณฺฑุ6 นิฆณฺฑุ จ.}
\hangingbody{ปนาโท โอปมญฺโญ จ, เทวสูโต จ มาตลิ.}
\csromanhangingend

\csromanheader{2}{7. มหาสมยสุตฺต}
\csromangathameasure{จิตฺตเสโน จ คนฺธพฺโพ, นโฬราชา ชเนสโภ. \\ อาคา ปญฺจสิโข เจว, ติมฺพรู สูริยวจฺฉสา. \\ เอเต จญฺเญ จ ราชาโน, คนฺธพฺพา สห ราชุภิ. \\ โมทมานา อภิกฺกามุํ, ภิกฺขูนํ สมิติํ วนํ.}
\csromangathagroup{\csromangathastanza{\csromanfolio{207}จิตฺตเสโน จ คนฺธพฺโพ, นโฬราชา ชเนสโภ\footnote{ชโนสโภ (สฺยา)}. \\ อาคา ปญฺจสิโข เจว, ติมฺพรู สูริยวจฺฉสา\footnote{ชโนสโภ (สฺยา)}.}\csromangathastanza{เอเต จญฺเญ จ ราชาโน, คนฺธพฺพา สห ราชุภิ. \\ โมทมานา อภิกฺกามุํ, ภิกฺขูนํ สมิติํ วนํ.}}

\csromanhangingbegin
\hangingheaditem{338}{อถาคุํ นาคสา นาคา, เวสาลา สหตจฺฉกา.}
\hangingbody{กมฺปลสฺสตรา อาคุํ, ปายาคา สห ญาติภิ.}
\hangingbody{ยามุนา ธตรฏฺฐา จ, อาคู นาคา ยสสฺสิโน.}
\hangingbody{เอราวโณ มหานาโค, โสปาคา สมิติํ วนํ.}
\hangingbody{เย นาคราเช สหสา หรนฺติ, ทิพฺพา ทิชา ปกฺขิ วิสุทฺธจกฺขู.}
\hangingbody{เวหายสา3 เต วนมชฺฌปตฺตา, จิตฺรา สุปณฺณา อิติ เตส นามํ.}
\hangingbody{อภยํ ตทา นาคราชานมาสิ,}
\hangingbody{สุปณฺณโต เขมมกาสิ พุทฺโธ.}
\hangingbody{สณฺหาหิ วาจาหิ อุปวฺหยนฺตา,}
\hangingbody{นาคา สุปณฺณา สรณมกํสุ พุทฺธํ.}
\csromanhangingend

\csromanhangingbegin
\hangingheaditem{339}{ชิตา วชิรหตฺเถน, สมุทฺทํ อสุราสิตา.}
\hangingbody{ภาตโร วาสวสฺเสเต, อิทฺธิมนฺโต ยสสฺสิโน.}
\hangingbody{กาลกญฺจา มหาภิสฺมา4, อสุรา ทานเวฆสา.}
\hangingbody{เวปจิตฺติ สุจิตฺติ จ, ปหาราโท นมุจี สห.}
\hangingbody{สตญฺจ พลิปุตฺตานํ, สพฺเพ เวโรจนามกา.}
\hangingbody{สนฺนยฺหิตฺวา พลิเสนํ5, ราหุภทฺทมุปาคมุํ.}
\hangingbody{สมโย ทานิ ภทฺทนฺเต, ภิกฺขูนํ สมิติํ วนํ.}
\csromanhangingend

\csromanhangingbegin
\hangingheaditem{340}{อาโป จ เทวา ปถวี, เตโช วาโย ตทาคมุํ.}
\hangingbody{วรุณา วารณา6 เทวา, โสโม จ ยสสา สห.}
\hangingbody{เมตฺตา กรุณา กายิกา, อาคุํ เทวา ยสสฺสิโน.}
\hangingbody{ทเสเต ทสธา กายา, สพฺเพ นานตฺตวณฺณิโน.}
\csromanhangingend

\csromangathameasure{อิทฺธิมนฺโต ชุติมนฺโต, วณฺณวนฺโต ยสสฺสิโน. \\ โมทมานา อภิกฺกามุํ, ภิกฺขูนํ สมิติํ วนํ. \\ เวณฺฑุเทวา สหลิ จ, อสมา จ ทุเว ยมา. \\ จนฺทสฺสูปนิสา เทวา, จนฺทมาคุํ ปุรกฺขตฺวา. \\ สูริยสฺสูปนิสา เทวา, สูริยมาคุํ ปุรกฺขตฺวา. \\ นกฺขตฺตานิ ปุรกฺขตฺวา, อาคุํ มนฺทวลาหกา. \\ วสูนํ วาสโว เสฏฺโฐ, สกฺโกปาคา ปุรินฺทโท. \\ ทเสเต ทสธา กายา, สพฺเพ นานตฺตวณฺณิโน. \\ อิทฺธิมนฺโต ชุติมนฺโต, วณฺณวนฺโต ยสสฺสิโน. \\ โมทมานา อภิกฺกามุํ, ภิกฺขูนํ สมิติํ วนํ. \\ อถาคุํ สหภู เทวา, ชลมคฺคิ สิขาริว. \\ อริฏฺฐกา จ โรชา จ, อุมาปุปฺผนิภาสิโน. \\ วรุณา สหธมฺมา จ, อจฺจุตา จ อเนชกา. \\ สูเลยฺยรุจิรา อาคุํ, อาคุํ วาสวเนสิโน. \\ ทเสเต ทสธา กายา, สพฺเพ นานตฺตวณฺณิโน. \\ อิทฺธิมนฺโต ชุติมนฺโต, วณฺณวนฺโต ยสสฺสิโน. \\ โมทมานา อภิกฺกามุํ, ภิกฺขูนํ สมิติํ วนํ. \\ สมานา มหาสมนา, มานุสา มานุสุตฺตมา. \\ ขิฑฺฑาปโทสิกา อาคุํ, อาคุํ มโนปโทสิกา. \\ อถาคุํ หรโย เทวา, เย จ โลหิตวาสิโน. \\ ปารคา มหาปารคา, อาคุํ เทวา ยสสฺสิโน. \\ ทเสเต ทสธา กายา, สพฺเพ นานตฺตวณฺณิโน. \\ อิทฺธิมนฺโต ชุติมนฺโต, วณฺณวนฺโต ยสสฺสิโน. \\ โมทมานา อภิกฺกามุํ, ภิกฺขูนํ สมิติํ วนํ. \\ สุกฺกา กรมฺภา อรุณา, อาคุํ เวฆนสา สห. \\ โอทาตคยฺหา ปาโมกฺขา, อาคุํ เทวา วิจกฺขณา.}
\csromangathagroup{\csromangathastanza{\csromanfolio{208}อิทฺธิมนฺโต ชุติมนฺโต, วณฺณวนฺโต ยสสฺสิโน. \\ โมทมานา อภิกฺกามุํ, ภิกฺขูนํ สมิติํ วนํ.}\csromangathastanza{เวณฺฑุเทวา สหลิ จ\footnote{เวณฺหูจ เทวา สหลีจ (สี, อิ)}, อสมา จ ทุเว ยมา. \\ จนฺทสฺสูปนิสา เทวา, จนฺทมาคุํ ปุรกฺขตฺวา.}\csromangathastanza{สูริยสฺสูปนิสา\footnote{สุริยสฺสูปนิสา (สี, สฺยา, อิ)} เทวา, สูริยมาคุํ ปุรกฺขตฺวา. \\ นกฺขตฺตานิ ปุรกฺขตฺวา, อาคุํ มนฺทวลาหกา.}\csromangathastanza{วสูนํ วาสโว เสฏฺโฐ, สกฺโกปาคา ปุรินฺทโท. \\ ทเสเต ทสธา กายา, สพฺเพ นานตฺตวณฺณิโน.}\csromangathastanza{อิทฺธิมนฺโต ชุติมนฺโต, วณฺณวนฺโต ยสสฺสิโน. \\ โมทมานา อภิกฺกามุํ, ภิกฺขูนํ สมิติํ วนํ.}\csromangathastanza{อถาคุํ สหภู เทวา, ชลมคฺคิ สิขาริว. \\ อริฏฺฐกา จ โรชา จ, อุมาปุปฺผนิภาสิโน.}\csromangathastanza{วรุณา สหธมฺมา จ, อจฺจุตา จ อเนชกา. \\ สูเลยฺยรุจิรา อาคุํ, อาคุํ วาสวเนสิโน.}\csromangathastanza{ทเสเต ทสธา กายา, สพฺเพ นานตฺตวณฺณิโน. \\ อิทฺธิมนฺโต ชุติมนฺโต, วณฺณวนฺโต ยสสฺสิโน.}\csromangathastanza{โมทมานา อภิกฺกามุํ, ภิกฺขูนํ สมิติํ วนํ. \\ สมานา มหาสมนา, มานุสา มานุสุตฺตมา.}\csromangathastanza{ขิฑฺฑาปโทสิกา อาคุํ, อาคุํ มโนปโทสิกา. \\ อถาคุํ หรโย เทวา, เย จ โลหิตวาสิโน.}\csromangathastanza{ปารคา มหาปารคา, อาคุํ เทวา ยสสฺสิโน. \\ ทเสเต ทสธา กายา, สพฺเพ นานตฺตวณฺณิโน.}\csromangathastanza{อิทฺธิมนฺโต ชุติมนฺโต, วณฺณวนฺโต ยสสฺสิโน. \\ โมทมานา อภิกฺกามุํ, ภิกฺขูนํ สมิติํ วนํ.}\csromangathastanza{สุกฺกา กรมฺภา\footnote{กรุมฺหา (สี, สฺยา, อิ)} อรุณา, อาคุํ เวฆนสา สห. \\ โอทาตคยฺหา ปาโมกฺขา, อาคุํ เทวา วิจกฺขณา.}}

\csromanheader{2}{7. มหาสมยสุตฺต}
\csromangathameasure{สทามตฺตา หารคชา, มิสฺสกา จ ยสสฺสิโน. \\ ถนยํ อาค ปชฺชุนฺโน, โย ทิสา อภิวสฺสติ. \\ ทเสเต ทสธา กายา, สพฺเพ นานตฺตวณฺณิโน. \\ อิทฺธิมนฺโต ชุติมนฺโต, วณฺณวนฺโต ยสสฺสิโน. \\ โมทมานา อภิกฺกามุํ, ภิกฺขูนํ สมิติํ วนํ. \\ เขมิยา ตุสิตา ยามา, กฏฺฐกา จ ยสสฺสิโน. \\ ลมฺพีตกา ลามเสฏฺฐา, โชตินามา จ อาสวา. \\ นิมฺมานรติโน อาคุํ, อถาคุํ ปรนิมฺมิตา. \\ ทเสเต ทสธา กายา, สพฺเพ นานตฺตวณฺณิโน. \\ อิทฺธิมนฺโต ชุติมนฺโต, วณฺณวนฺโต ยสสฺสิโน. \\ โมทมานา อภิกฺกามุํ, ภิกฺขูนํ สมิติํ วนํ. \\ สฏฺเฐเต เทวนิกายา, สพฺเพ นานตฺตวณฺณิโน. \\ นามนฺวเยน อาคจฺฉุํ, เย จญฺเญ สทิสา สห. \\ ปวุฏฺฐชาติมขิลํ, โอฆติณฺณมนาสวํ. \\ ทกฺเขโมฆตรํ นาคํ, จนฺทํว อสิตาติคํ.}
\csromangathagroup{\csromangathastanza{\csromanfolio{209}สทามตฺตา หารคชา, มิสฺสกา จ ยสสฺสิโน. \\ ถนยํ อาค ปชฺชุนฺโน, โย ทิสา อภิวสฺสติ.}\csromangathastanza{ทเสเต ทสธา กายา, สพฺเพ นานตฺตวณฺณิโน. \\ อิทฺธิมนฺโต ชุติมนฺโต, วณฺณวนฺโต ยสสฺสิโน.}\csromangathastanza{โมทมานา อภิกฺกามุํ, ภิกฺขูนํ สมิติํ วนํ. \\ เขมิยา ตุสิตา ยามา, กฏฺฐกา จ ยสสฺสิโน.}\csromangathastanza{ลมฺพีตกา ลามเสฏฺฐา, โชตินามา จ อาสวา. \\ นิมฺมานรติโน อาคุํ, อถาคุํ ปรนิมฺมิตา.}\csromangathastanza{ทเสเต ทสธา กายา, สพฺเพ นานตฺตวณฺณิโน. \\ อิทฺธิมนฺโต ชุติมนฺโต, วณฺณวนฺโต ยสสฺสิโน.}\csromangathastanza{โมทมานา อภิกฺกามุํ, ภิกฺขูนํ สมิติํ วนํ. \\ สฏฺเฐเต เทวนิกายา, สพฺเพ นานตฺตวณฺณิโน.}\csromangathastanza{นามนฺวเยน อาคจฺฉุํ\footnote{อาคญฺฉุํ (สี, สฺยา, อิ)}, เย จญฺเญ สทิสา สห. \\ ปวุฏฺฐชาติมขิลํ\footnote{อาคญฺฉุํ (สี, สฺยา, อิ)}, โอฆติณฺณมนาสวํ. \\ ทกฺเขโมฆตรํ นาคํ, จนฺทํว อสิตาติคํ.}}

\csromanhangingbegin
\hangingheaditem{341}{สุพฺรหฺมา ปรมตฺโต จ\footnote{ปรมตฺโถ จ (ก)}, ปุตฺตา อิทฺธิมโต สห.}
\hangingbody{สนงฺกุมาโร ติสฺโส จ, โสปาค สมิติํ วนํ.}
\hangingbody{สหสฺสํ พฺรหฺมโลกานํ, มหาพฺรหฺมาภิติฏฺฐติ.}
\hangingbody{อุปปนฺโน ชุติมนฺโต, ภิสฺมากาโย ยสสฺสิโส.}
\hangingbody{ทเสตฺถ อิสฺสรา อาคุํ, ปจฺเจกวสวตฺติโน.}
\hangingbody{เตสญฺจ มชฺฌโต อาค, หาริโต ปริวาริโต.}
\csromanhangingend

\csromanhangingbegin
\hangingheaditem{342}{เต จ สพฺเพ อภิกฺกนฺเต, ส-อินฺเท\footnote{สินฺเท (สฺยา)} เทเว สพฺรหฺมเก.}
\hangingbody{มารเสนา อภิกฺกามิ, ปสฺส กณฺหสฺส มนฺทิยํ.}
\hangingbody{เอถ คณฺหถ พนฺธถ, ราเคน พทฺธมตฺถุ โว.}
\hangingbody{สมนฺตา ปริวาเรถ, มา โว มุญฺจิตฺถ โกจิ นํ.}
\csromanhangingend

\csromangathameasure{อิติ ตตฺถ มหาเสโน, กโณฺห เสนํ อเปสยิ. \\ ปาณินา ตลมาหจฺจ, สรํ กตฺวาน เภรวํ. \\ ยถา ปาวุสฺสโก เมโฆ, ถนยนฺโต สวิชฺชุโก. \\ ตทา โส ปจฺจุทาวตฺติ, สงฺกุทฺโธ อสยํวเส.}
\csromangathagroup{\csromangathastanza{\csromanfolio{210}อิติ ตตฺถ มหาเสโน, กโณฺห เสนํ อเปสยิ. \\ ปาณินา ตลมาหจฺจ, สรํ กตฺวาน เภรวํ.}\csromangathastanza{ยถา ปาวุสฺสโก เมโฆ, ถนยนฺโต สวิชฺชุโก. \\ ตทา โส ปจฺจุทาวตฺติ, สงฺกุทฺโธ อสยํวเส\footnote{อสยํวสี (สี, อิ)}.}}

\csromanhangingbegin
\hangingheaditem{343}{อญฺจ สพฺพํ อภิญฺญาย, ววตฺถิตฺวาน จกฺขุมา.}
\hangingbody{ตโต อามนฺตยี สตฺถา, สาวเก สาสเน รเต.}
\hangingbody{มารเสนา อภิกฺกนฺตา, เต วิชานาถ ภิกฺขโว.}
\hangingbody{เต จ อาตปฺปมกรุํ, สุตฺวา พุทฺธสฺส สาสนํ.}
\hangingbody{วีตราเคหิ ปกฺกามุํ, เนสํ โลมาปิ อิญฺชยุํ.}
\hangingbody{สพฺเพ วิชิตสงฺคามา, ภยาตีตา ยสสฺสิโน.}
\hangingbody{โมทนฺติ สห ภูเตหิ, สาวกา เต ชเนสุตาติ.}
\csromanhangingend

\nitthitam{\textbf{มหาสมยสุตฺตํ นิฏฺฐิตํ สตฺตมํ.}}
\csromanmatikaanchor{mtk.109}
\csromantocmarknopage{chapter}{8. สกฺกปญฺหสุตฺต}{mtk.109}
\bookmark[dest={mtk.109},level=0]{8. สกฺกปญฺหสุตฺต}
\chapterheadpageread{8. สกฺกปญฺหสุตฺต}
\proseitem{344}{\csromanfolio{211}เอวํ เม สุตํ\csromandash{}เอกํ สมยํ ภควา มคเธสุ วิหรติ ปาจีนโต ราชคหสฺส อมฺพสณฺฑา นาม พฺราหฺมณคาโม, ตสฺสุตฺตรโต เวทิยเก ปพฺพเต อินฺทสาลคุหายํ.\csromanspacer{} เตน โข ปน สมเยน สกฺกสฺส เทวานมินฺทสฺส อุสฺสุกฺกํ อุทปาทิ ภควนฺตํ ทสฺสนาย.\csromanspacer{} อถ โข สกฺกสฺส เทวานมินฺทสฺส เอตทโหสิ “กหํ นุ โข ภควา เอตรหิ วิหรติ อรหํ สมฺมาสมฺพุทฺโธ”ติ.\csromanspacer{} อทฺทสา โข สกฺโก เทวานมินฺโท ภควนฺตํ มคเธสุ วิหรนฺตํ ปาจีนโต ราชคหสฺส อมฺพสณฺฑา นาม พฺราหฺมณคาโม, ตสฺสุตฺตรโต เวทิยเก ปพฺพเต อินฺทสาลคุหายํ, ทิสฺวาน เทเว ตาวติํเส อามนฺเตสิ “อยํ มาริสา ภควา มคเธสุ วิหรติ ปาจีนโต ราชคหสฺส อมฺพสณฺฑา นาม พฺราหฺมณคาโม, ตสฺสุตฺตรโต เวทิยเก ปพฺพเต อินฺทสาลคุหายํ.\csromanspacer{} ยทิ ปน มาริสา มยํ ตํ ภควนฺตํ ทสฺสนาย อุปสงฺกเมยฺยาม อรหนฺตํ สมฺมาสมฺพุทฺธนฺ”ติ. “เอวํ ภทฺทนฺตวา”ติ โข เทวา ตาวติํสา สกฺกสฺส เทวานมินฺทสฺส ปจฺจสฺโสสุํ.}

\proseitem{345}{อถ โข สกฺโก เทวานมินฺโท ปญฺจสิขํ คนฺธพฺพเทวปุตฺตํ\footnote{คนฺธพฺพปุตฺตํ (สฺยา)} อามนฺเตสิ “อยํ ตาต ปญฺจสิข ภควา มคเธสุ วิหรติ ปาจีนโต ราชคหสฺส อมฺพสณฺฑา นาม พฺราหฺมณคาโม, ตสฺสุตฺตรโต เวทิยเก ปพฺพเต อินฺทสาลคุหายํ.\csromanspacer{} ยทิ ปน ตาต ปญฺจสิข มยํ ตํ ภควนฺตํ ทสฺสนาย อุปสงฺกเมยฺยาม อรหนฺตํ สมฺมาสมฺพุทฺธนฺ”ติ. “เอวํ ภทฺทนฺตวา”ติ โข ปญฺจสิโข คนฺธพฺพเทวปุตฺโต สกฺกสฺส เทวานมินฺทสฺส ปฏิสฺสุตฺวา เพลุวปณฺฑุวีณํ อาทาย สกฺกสฺส เทวานมินฺทสฺส อนุจริยํ อุปาคมิ.}

\proseitem{346}{อถ โข สกฺโก เทวานมินฺโท เทเวหิ ตาวติํเสหิ ปริวุโต ปญฺจสิเขน คนฺธพฺพเทวปุตฺเตน ปุรกฺขโต เสยฺยถาปิ นาม พลวา ปุริโส สมิญฺชิตํ วา พาหํ ปสาเรยฺย ปสาริตํ วา พาหํ สมิญฺเชยฺย, เอวเมว เทเวสุ ตาวติํเสสุ อนฺตรหิโต มคเธสุ ปาจีนโต ราชคหสฺส อมฺพสณฺฑา นาม พฺราหฺมณคาโม, ตสฺสุตฺตรโต เวทิยเก ปพฺพเต ปจฺจุฏฺฐาสิ.\csromanspacer{} เตน โข ปน สมเยน เวทิยโก \csromanfolio{212}ปพฺพโต อติริว โอภาสชาโต โหติ อมฺพสณฺฑา จ พฺราหฺมณคาโม ยถา ตํ เทวานํ เทวานุภาเวน.\csromanspacer{} อปิสฺสุทํ ปริโต คาเมสุ มนุสฺสา เอวมาหํสุ “อาทิตฺตสฺสุ นามชฺช เวทิยโก ปพฺพโต.\csromanspacer{} ฌายติสุ\footnote{ฌายตสฺสุ (สฺยา), ปชฺฌาริตสฺสุ (สี, อิ)} นามชฺช เวทิยโก ปพฺพโต.\csromanspacer{} ชลติสุ\footnote{ชลตสฺสุ (สฺยา), ชลิตสฺสุ (สี, อิ)} นามชฺช เวทิยโก ปพฺพโต.\csromanspacer{} กิํสุนามชฺช เวทิยโก ปพฺพโต อติริว โอภาสชาโต อมฺพสณฺฑา จ พฺราหฺมณคาโม”ติ.\csromanspacer{} สํวิคฺคา โลมหฏฺฐชาตา อเหสุํ.}

\proseitem{347}{อถ โข สกฺโก เทวานมินฺโท ปญฺจสิขํ คนฺธพฺพเทวปุตฺตํ อามนฺเตสิ “ทุรุปสงฺกมา โข ตาต ปญฺจสิข ตถาคตา มาทิเสน, ฌายี ฌานรตา, ตทนฺตรํ\footnote{ตทนนฺตรํ (สี, สฺยา, อิ, ก)} ปฏิสลฺลีนา.\csromanspacer{} ยทิ ปน ตฺวํ ตาต ปญฺจสิข ภควนฺตํ ปฐมํ ปสาเทยฺยาสิ, ตยา ตาต ปฐมํ ปสาทิตํ ปจฺฉา มยํ ตํ ภควนฺตํ ทสฺสนาย อุปสงฺกเมยฺยาม อรหนฺตํ สมฺมาสมฺพุทฺธนฺ”ติ. “เอวํ ภทฺทนฺตวา”ติ โข ปญฺจสิโข คนฺธพฺพเทวปุตฺโต สกฺกสฺส เทวานมินฺทสฺส ปฏิสฺสุตฺวา เพลุวปณฺฑุวีณํ อาทาย เยน อินฺทสาลคุหา เตนุปสงฺกมิ, อุปสงฺกมิตฺวา “เอตฺตาวตา เม ภควา เนว อติทูเร ภวิสฺสติ นาจฺจาสนฺเน, สทฺทญฺจ เม โสสฺสตี”ติ เอกมนฺตํ อฏฺฐาสิ.}

\csromanmatikaanchor{mtk.110}
\csromantocmarkref{section}{ปญฺจสิขคีตคาถา}{mtk.110}
\bookmark[dest={mtk.110},level=1]{ปญฺจสิขคีตคาถา}
\csromanheader{1}{\textbf{ปญฺจสิขคีตคาถา}}
\proseitem{348}{เอกมนฺตํ ฐิโต โข ปญฺจสิโข คนฺธพฺพเทวปุตฺโต เพลุวปณฺฑุวีณํ\footnote{เวฬุวปณฺฑุวีณํ อาทาย (สฺยา)} อสฺสาเวสิ, อิมา จ คาถา อภาสิ พุทฺธูปสญฺหิตา ธมฺมูปสญฺหิตา สํฆูปสญฺหิตา อรหนฺตูปสญฺหิตา กามูปสญฺหิตา.}

\csromangathameasure{“วนฺเท เต ปิตรํ ภทฺเท, ติมฺพรุํ สูริยวจฺฉเส. \\ เยน ชาตาสิ กลฺยาณี, อานนฺทชนนี มม. \\ วาโตว เสทตํ กนฺโต, ปานียํว ปิปาสโต. \\ องฺคีรสิ ปิยาเมสิ, ธมฺโม อรหตามิว. \\ อาตุรสฺเสว เภสชฺชํ, โภชนํว ชิฆจฺฉโต. \\ ปรินิพฺพาปย มํ ภทฺเท, ชลนฺตมิว วารินา.}
\csromangathagroup{\csromangathastanza{“วนฺเท เต ปิตรํ ภทฺเท, ติมฺพรุํ สูริยวจฺฉเส. \\ เยน ชาตาสิ กลฺยาณี, อานนฺทชนนี มม.}\csromangathastanza{วาโตว เสทตํ กนฺโต, ปานียํว ปิปาสโต. \\ องฺคีรสิ ปิยาเมสิ, ธมฺโม อรหตามิว.}\csromangathastanza{อาตุรสฺเสว เภสชฺชํ, โภชนํว ชิฆจฺฉโต. \\ ปรินิพฺพาปย มํ ภทฺเท, ชลนฺตมิว วารินา.}}

\csromanheader{2}{8. สกฺกปญฺหสุตฺต}
\csromangathameasure{สีโตทกํ โปกฺขรณิํ, ยุตฺตํ กิญฺชกฺขเรณุนา. \\ นาโค ฆมฺมาภิตตฺโตว, โอคาเห เต ถนูทรํ. \\ อจฺจงฺกุโสว นาโคว, ชิตํ เม ตุตฺตโตมรํ. \\ การณํ นปฺปชานามิ, สมฺมตฺโต ลกฺขณูรุยา. \\ ตยิ เคธิตจิตฺโตสฺมิ, จิตฺตํ วิปริณามิตํ. \\ ปฏิคนฺตุํ น สกฺโกมิ, วงฺกฆสฺโตว อมฺพุโช. \\ วามูรุ สช มํ ภทฺเท, สช มํ มนฺทโลจเน. \\ ปลิสฺสช มํ กลฺยาณิ, เอตํ เม อภิปตฺถิตํ. \\ อปฺปโก วต เม สนฺโต, กาโม เวลฺลิตเกสิยา. \\ อเนกภาโว สมุปฺปาทิ, อรหนฺเตว ทกฺขิณา. \\ ยํ เม อตฺถิ กตํ ปุญฺญํ, อรหนฺเตสุ ตาทิสุ. \\ ตํ เม สพฺพงฺคกลฺยาณิ, ตยา สทฺธิํ วิปจฺจตํ. \\ ยํ เม อตฺถิ กตํ ปุญฺญํ, อสฺมิํ ปถวิมณฺฑเล. \\ ตํ เม สพฺพงฺคกลฺยาณิ, ตยา สทฺธิํ วิปจฺจตํ. \\ สกฺยปุตฺโตว ฌาเนน, เอโกทิ นิปโก สโต. \\ อมตํ มุนิ ชิคีสาโน, ตมหํ สูริยวจฺฉเส. \\ ยถาปิ มุนิ นนฺเทยฺย, ปตฺวา สมฺโพธิมุตฺตมํ. \\ เอวํ นนฺเทยฺยํ กลฺยาณิ, มิสฺสีภาวํ คโต ตยา. \\ สกฺโก เจ เม วรํ ทชฺชา, ตาวติํสานมิสฺสโร. \\ ตาหํ ภทฺเท วเรยฺยาเห, เอวํ กาโม ทโฬฺห มม. \\ สาลํว น จิรํ ผุลฺลํ, ปิตรํ เต สุเมธเส. \\ วนฺทมาโน นมสฺสามิ, ยสฺสาเสตาทิสี ปชา”ติ.}
\csromangathagroup{\csromangathastanza{\csromanfolio{213}สีโตทกํ โปกฺขรณิํ, ยุตฺตํ กิญฺชกฺขเรณุนา. \\ นาโค ฆมฺมาภิตตฺโตว, โอคาเห เต ถนูทรํ.}\csromangathastanza{อจฺจงฺกุโสว นาโคว, ชิตํ เม ตุตฺตโตมรํ. \\ การณํ นปฺปชานามิ, สมฺมตฺโต ลกฺขณูรุยา.}\csromangathastanza{ตยิ เคธิตจิตฺโตสฺมิ, จิตฺตํ วิปริณามิตํ. \\ ปฏิคนฺตุํ น สกฺโกมิ, วงฺกฆสฺโตว อมฺพุโช.}\csromangathastanza{วามูรุ สช มํ ภทฺเท, สช มํ มนฺทโลจเน. \\ ปลิสฺสช มํ กลฺยาณิ, เอตํ เม อภิปตฺถิตํ.}\csromangathastanza{อปฺปโก วต เม สนฺโต, กาโม เวลฺลิตเกสิยา. \\ อเนกภาโว สมุปฺปาทิ, อรหนฺเตว ทกฺขิณา.}\csromangathastanza{ยํ เม อตฺถิ กตํ ปุญฺญํ, อรหนฺเตสุ ตาทิสุ. \\ ตํ เม สพฺพงฺคกลฺยาณิ, ตยา สทฺธิํ วิปจฺจตํ.}\csromangathastanza{ยํ เม อตฺถิ กตํ ปุญฺญํ, อสฺมิํ ปถวิมณฺฑเล. \\ ตํ เม สพฺพงฺคกลฺยาณิ, ตยา สทฺธิํ วิปจฺจตํ.}\csromangathastanza{สกฺยปุตฺโตว ฌาเนน, เอโกทิ นิปโก สโต. \\ อมตํ มุนิ ชิคีสาโน\footnote{ชิคิํสาโน (สี, สฺยา, อิ)}, ตมหํ สูริยวจฺฉเส.}\csromangathastanza{ยถาปิ มุนิ นนฺเทยฺย, ปตฺวา สมฺโพธิมุตฺตมํ. \\ เอวํ นนฺเทยฺยํ กลฺยาณิ, มิสฺสีภาวํ คโต ตยา.}\csromangathastanza{สกฺโก เจ เม วรํ ทชฺชา, ตาวติํสานมิสฺสโร. \\ ตาหํ ภทฺเท วเรยฺยาเห, เอวํ กาโม ทโฬฺห มม.}\csromangathastanza{สาลํว น จิรํ ผุลฺลํ, ปิตรํ เต สุเมธเส. \\ วนฺทมาโน นมสฺสามิ, ยสฺสาเสตาทิสี ปชา”ติ.}}

\proseitem{349}{เอวํ วุตฺเต ภควา ปญฺจสิขํ คนฺธพฺพเทวปุตฺตํ เอตทโวจ “สํสนฺทติ โข เต ปญฺจสิข ตนฺติสฺสโร คีตสฺสเรน, คีตสฺสโร จ ตนฺติสฺสเรน, น จ ปน\footnote{เนว ปน (สฺยา)} เต ปญฺจสิข ตนฺติสฺสโร คีตสฺสรํ อติวตฺตติ, คีตสฺสโร จ ตนฺติสฺสรํ.\csromanspacer{} กทา สํยูฬฺหา ปน เต ปญฺจสิข อิมา คาถา พุทฺธูปสญฺหิตา \csromanfolio{214}ธมฺมูปสญฺหิตา สํฆูปสญฺหิตา อรหนฺตูปสญฺหิตา กามูปสญฺหิตา”ติ.\csromanspacer{} เอกมิทํ ภนฺเต สมยํ ภควา อุรุเวลายํ วิหรติ นชฺชา เนรญฺชราย ตีเร อชปาลนิโคฺรเธ ปฐมาภิสมฺพุทฺโธ.\csromanspacer{} เตน โข ปนาหํ ภนฺเต สมเยน ภทฺทา นาม สูริยวจฺฉสา ติมฺพรุโน คนฺธพฺพรญฺโญ ธีตา, ตมภิกงฺขามิ.\csromanspacer{} สา โข ปน ภนฺเต ภคินี ปรกามินี โหติ, สิขณฺฑี นาม มาตลิสฺส สงฺคาหกสฺส ปุตฺโต, ตมภิกงฺขติ.\csromanspacer{} ยโต โข อหํ ภนฺเต ตํ ภคินิํ นาลตฺถํ เกนจิ ปริยาเยน.\csromanspacer{} อถาหํ เพลุวปณฺฑุวีณํ อาทาย เยน ติมฺพรุโน คนฺธพฺพรญฺโญ นิเวสนํ เตนุปสงฺกมิํ, อุปสงฺกมิตฺวา เพลุวปณฺฑุวีณํ อสฺสาเวสิํ, อิมา จ คาถา อภาสิํ พุทฺธูปสญฺหิตา ธมฺมูปสญฺหิตา สํฆูปสญฺหิตา อรหนฺตูปสญฺหิตา กามูปสญฺหิตา.}

\csromangathameasure{“วนฺเท เต ปิตรํ ภทฺเท, ติมฺพรุํ สูริยวจฺฉเส. \\ เยน ชาตาสิ กลฺยาณี, อานนฺทชนนี มม ฯเปฯ. \\ สาลํว น จิรํ ผุลฺลํ, ปิตรํ เต สุเมธเส. \\ วนฺทมาโน นมสฺสามิ, ยสฺสาเสตาทิสี ปชา”ติ.}
\csromangathagroup{\csromangathastanza{“วนฺเท เต ปิตรํ ภทฺเท, ติมฺพรุํ สูริยวจฺฉเส. \\ เยน ชาตาสิ กลฺยาณี, อานนฺทชนนี มม ฯเปฯ.}\csromangathastanza{สาลํว น จิรํ ผุลฺลํ, ปิตรํ เต สุเมธเส. \\ วนฺทมาโน นมสฺสามิ, ยสฺสาเสตาทิสี ปชา”ติ.}}

\prose{เอวํ วุตฺเต ภนฺเต ภทฺทา สูริยวจฺฉสา มํ เอตทโวจ “น โข เม มาริส โส ภควา สมฺมุขา ทิฏฺโฐ.\csromanspacer{} อปิ จ สุโตเยว เม โส ภควา เทวานํ ตาวติํสานํ สุธมฺมายํ สภายํ อุปนจฺจนฺติยา.\csromanspacer{} ยโต โข ตฺวํ มาริส ตํ ภควนฺตํ กิตฺเตสิ, โหตุ โน อชฺช สมาคโม”ติ.\csromanspacer{} โสเยว โน ภนฺเต ตสฺสา ภคินิยา สทฺธิํ สมาคโม อโหสิ.\csromanspacer{} น จ ทานิ ตโต ปจฺฉาติ.}

\csromanmatikaanchor{mtk.111}
\csromantocmarkref{section}{สกฺกูปสงฺกม}{mtk.111}
\bookmark[dest={mtk.111},level=1]{สกฺกูปสงฺกม}
\csromanheader{1}{\textbf{สกฺกูปสงฺกม}}
\proseitem{350}{อถ โข สกฺกสฺส เทวานมินฺทสฺส เอตทโหสิ “ปฏิสมฺโมทติ ปญฺจสิโข คนฺธพฺพเทวปุตฺโต ภควตา, ภควา จ ปญฺจสิเขนา”ติ.\csromanspacer{} อถ โข สกฺโก เทวานมินฺโท ปญฺจสิขํ คนฺธพฺพเทวปุตฺตํ อามนฺเตสิ “อภิวาเทหิ เม ตฺวํ ตาต ปญฺจสิข ภควนฺตํ ‘สกฺโก ภนฺเต เทวานมินฺโท สามจฺโจ สปริชโน ภควโต ปาเท สิรสา วนฺทตี’ติ”. “เอวํ ภทฺทนฺตวา”ติ โข ปญฺจสิโข คนฺธพฺพเทวปุตฺโต สกฺกสฺส เทวานมินฺทสฺส ปฏิสฺสุตฺวา ภควนฺตํ อภิวาเทติ “สกฺโก ภนฺเต เทวานฺมินฺโท สามจฺโจ สปริชโน}

\vspace{\tipitakaparskip}
\csromancenter{สกฺกปญฺหสุตฺต ภควโต ปาเท สิรสา วนฺทตี”ติ.\csromanspacer{} เอวํ สุขี โหตุ ปญฺจสิข สกฺโก เทวานมินฺโท สามจฺโจ สปริชโน, สุขกามา หิ เทวา มนุสฺสา อสุรา นาคา คนฺธพฺพา เย จญฺเญ สนฺติ ปุถุกายาติ.}
\proseitem{351}{\csromanfolio{215}เอวญฺจ ปน ตถาคตา เอวรูเป มเหสกฺเข ยกฺเข อภิวทนฺติ.\csromanspacer{} อภิวทิโต สกฺโก เทวานมินฺโท ภควโต อินฺทสาลคุหํ ปวิสิตฺวา ภควนฺตํ อภิวาเทตฺวา เอกมนฺตํ อฏฺฐาสิ.\csromanspacer{} เทวาปิ ตาวติํสา อินฺทสาลคุหํ ปวิสิตฺวา ภควนฺตํ อภิวาเทตฺวา เอกมนฺตํ อฏฺฐํสุ.\csromanspacer{} ปญฺจสิโขปิ คนฺธพฺพเทวปุตฺโต อินฺทสาลคุหํ ปวิสิตฺวา ภควนฺตํ อภิวาเทตฺวา เอกมนฺตํ อฏฺฐาสิ.}

\prose{เตน โข ปน สมเยน อินฺทสาลคุหา วิสมา สนฺตี สมา สมปาทิ, สมฺพาธา สนฺตี อุรุนฺทา\footnote{อุรุทฺทา (ก)} สมปาทิ, อนฺธกาโร คุหายํ อนฺตรธายิ, อาโลโก อุทปาทิ.\csromanspacer{} ยถา ตํ เทวานํ เทวานุภาเวน.}

\proseitem{352}{อถ โข ภควา สกฺกํ เทวานมินฺทํ เอตทโวจ “อจฺฉริยมิทํ อายสฺมโต โกสิยสฺส, อพฺภุตมิทํ อายสฺมโต โกสิยสฺส ตาว พหุกิจฺจสฺส พหุกรณียสฺส ยทิทํ อิธาคมนนฺ”ติ.\csromanspacer{} จิรปฏิกาหํ ภนฺเต ภควนฺตํ ทสฺสนาย อุปสงฺกมิตุกาโม, อปิ จ เทวานํ ตาวติํสานํ เกหิจิ เกหิจิ\footnote{เกหิจิ (สฺยา)} กิจฺจกรณีเยหิ พฺยาวโฏ, เอวาหํ นาสกฺขิํ ภควนฺตํ ทสฺสนาย อุปสงฺกมิตุํ.\csromanspacer{} เอกมิทํ ภนฺเต สมยํ ภควา สาวตฺถิยํ วิหรติ สลฬาคารเก.\csromanspacer{} อถ ขฺวาหํ ภนฺเต สาวตฺถิํ อคมาสิํ ภควนฺตํ ทสฺสนาย.\csromanspacer{} เตน โข ปน ภนฺเต สมเยน ภควา อญฺญตเรน สมาธินา นิสินฺโน โหติ.\csromanspacer{} ภูชติ จ\footnote{ภุญฺชตี จ (สี, อิ), ภุชคี (สฺยา)} นาม เวสฺสวณสฺส มหาราชสฺส ปริจาริกา ภควนฺตํ ปจฺจุปฏฺฐิตา โหติ, ปญฺชลิกา นมสฺสมานา ติฏฺฐติ.\csromanspacer{} อถ ขฺวาหํ ภนฺเต ภูชติํ เอตทโวจํ “อภิวาเทหิ เม ตฺวํ ภคินิ ภควนฺตํ ‘สกฺโก ภนฺเต เทวานมินฺโท สามจฺโจ สปริชโน ภควโต ปาเท สิรสา วนฺทตี’ติ”.\csromanspacer{} เอวํ วุตฺเต ภนฺเต สา ภูชติ มํ เอตทโวจ “อกาโล โข มาริส ภควนฺตํ ทสฺสนาย, ปฏิสลฺลีโน ภควา”ติ.\csromanspacer{} เตน หิ ภคินิ ยทา ภควา ตมฺหา สมาธิมฺหา วุฏฺฐิโต โหติ.\csromanspacer{} อถ มม วจเนน ภควนฺตํ อภิวาเทหิ “สกฺโก ภนฺเต \csromanfolio{216}เทวานมินฺโท สามจฺโจ สปริชโน ภควโต ปาเท สิรสา วนฺทตี”ติ.\csromanspacer{} กจฺจิ เม สา ภนฺเต ภคินี ภควนฺตํ อภิวาเทสิ.\csromanspacer{} สรติ ภควา ตสฺสา ภคินิยา วจนนฺติ.\csromanspacer{} อภิวาเทสิ มํ สา เทวานมินฺท ภคินี, สรามหํ ตสฺสา ภคินิยา วจนํ.\csromanspacer{} อปิจาหํ อายสฺมโต เนมิสทฺเทน\footnote{จกฺกเนมิสทฺเทน (สฺยา)} ตมฺหา สมาธิมฺหา วุฏฺฐิโตติ.\csromanspacer{} เย เต ภนฺเต เทวา อเมฺหหิ ปฐมตรํ ตาวติํสกายํ อุปปนฺนา, เตสํ เม สมฺมุขา สุตํ สมฺมุขา ปฏิคฺคหิตํ “ยทา ตถาคตา โลเก อุปฺปชฺชนฺติ อรหนฺโต สมฺมาสมฺพุทฺธา, ทิพฺพา กายา ปริปูเรนฺติ, หายนฺติ อสุรกายา”ติ.\csromanspacer{} ตํ เม อิทํ ภนฺเต สกฺขิทิฏฺฐํ, ยโต ตถาคโต โลเก อุปฺปนฺโน อรหํ สมฺมาสมฺพุทฺโธ, ทิพฺพา กายา ปริปูเรนฺติ, หายนฺติ อสุรกายาติ.}

\csromanmatikaanchor{mtk.112}
\csromantocmarkref{section}{โคปกวตฺถุ}{mtk.112}
\bookmark[dest={mtk.112},level=1]{โคปกวตฺถุ}
\csromanheader{1}{\textbf{โคปกวตฺถุ}}
\proseitem{353}{อิเธว ภนฺเต กปิลวตฺถุสฺมิํ โคปิกา นาม สกฺยธีตา อโหสิ พุทฺเธ ปสนฺนา ธมฺเม ปสนฺนา สํเฆ ปสนฺนา สีเลสุ ปริปูรการินี, สา อิตฺถิตฺตํ\footnote{อิตฺถิจิตฺตํ (สฺยา)} วิราเชตฺวา ปุริสตฺตํ\footnote{ปุริสจิตฺตํ (สฺยา)} ภาเวตฺวา กายสฺส เภทา ปรํ มรณา สุคติํ สคฺคํ โลกํ อุปปนฺนา, เทวานํ ตาวติํสานํ สหพฺยตํ อมฺหากํ ปุตฺตตฺตํ อชฺฌุปคตา.\csromanspacer{} ตตฺรปิ นํ เอวํ ชานนฺติ “โคปโก เทวปุตฺโต โคปโก เทวปุตฺโต”ติ.\csromanspacer{} อญฺเญปิ ภนฺเต ตโย ภิกฺขู ภควติ พฺรหฺมจริยํ จริตฺวา หีนํ คนฺธพฺพกายํ อุปปนฺนา.\csromanspacer{} เต ปญฺจหิ กามคุเณหิ สมปฺปิตา สมงฺคีภูตา ปริจารยมานา อมฺหากํ อุปฏฺฐานํ อาคจฺฉนฺติ อมฺหากํ ปาริจริยํ.\csromanspacer{} เต อมฺหากํ อุปฏฺฐานํ อาคเต อมฺหากํ ปาริจริยํ โคปโก เทวปุตฺโต ปฏิโจเทสิ “กุโตมุขา นาม ตุเมฺห มาริสา ตสฺส ภควโต ธมฺมํ อสฺสุตฺถ\footnote{อายุหิตฺถ (สฺยา)}.\csromanspacer{} อหํ หิ นาม อิตฺถิกา สมานา พุทฺเธ ปสนฺนา ธมฺเม ปสนฺนา สํเฆ ปสนฺนา สีเลสุ ปริปูรการินี อิตฺถิตฺตํ วิราเชตฺวา ปุริสตฺตํ ภาเวตฺวา กายสฺส เภทา ปรํ มรณา สุคติํ สคฺคํ โลกํ อุปปนฺนา, เทวานํ ตาวติํสานํ สหพฺยตํ สกฺกสฺส เทวานมินฺทสฺส ปุตฺตตฺตํ อชฺฌุปคตา.\csromanspacer{} อิธาปิ มํ เอวํ ชานนฺติ ‘โคปโก เทวปุตฺโต โคปโก เทวปุตฺโต’ติ.\csromanspacer{} ตุเมฺห ปน มาริสา ภควติ พฺรหฺมจริยํ จริตฺวา หีนํ คนฺธพฺพกายํ อุปปนฺนา.\csromanspacer{} ทุทฺทิฏฺฐรูปํ วต โภ อทฺทสาม, เย มยํ}

\vspace{\tipitakaparskip}
\csromancenter{สกฺกปญฺหสุตฺต อทฺทสาม สหธมฺมิเก หีนํ คนฺธพฺพกายํ อุปปนฺเน”ติ.\csromanspacer{} เตสํ ภนฺเต โคปเกน เทวปุตฺเตน ปฏิโจทิตานํ เทฺว เทวา ทิฏฺเฐวธมฺเม สติํ ปฏิลภิํสุ กายํ พฺรหฺมปุโรหิตํ, เอโก ปน เทโว กาเม อชฺฌาวสิ.}
\csromanhangingbegin
\hangingheaditem{354}{\csromanfolio{217}อุปาสิกา จกฺขุมโต อโหสิํ,}
\hangingbody{นามมฺปิ มยฺหํ อหุ “โคปิกา”ติ.}
\hangingbody{พุทฺเธ จ ธมฺเม จ อภิปฺปสนฺนา,}
\hangingbody{สํฆญฺจุปฏฺฐาสิํ ปสนฺนจิตฺตา.}
\hangingbody{ตสฺเสว พุทฺธสฺส สุธมฺมตาย,}
\hangingbody{สกฺกสฺส ปุตฺโตมฺหิ มหานุภาโว.}
\hangingbody{มหาชุตีโก ติทิวูปปนฺโน,}
\hangingbody{ชานนฺติ มํ อิธาปิ “โคปโก”ติ.}
\hangingbody{อถทฺทสํ ภิกฺขโว ทิฏฺฐปุพฺเพ,}
\hangingbody{คนฺธพฺพกายูปคเต วสีเน.}
\hangingbody{อิเมหิ เต โคตมสาวกาเส,}
\hangingbody{เย จ มยํ ปุพฺเพ มนุสฺสภูตา.}
\hangingbody{อนฺเนน ปาเนน อุปฏฺฐหิมฺหา,}
\hangingbody{ปาทูปสงฺคยฺห สเก นิเวสเน.}
\hangingbody{กุโตมุขา นาม อิเม ภวนฺโต,}
\hangingbody{พุทฺธสฺส ธมฺมานิ ปฏิคฺคเหสุํ1.}
\hangingbody{ปจฺจตฺตํ เวทิตพฺโพ หิ ธมฺโม,}
\hangingbody{สุเทสิโต จกฺขุมตานุพุทฺโธ.}
\hangingbody{อหํ หิ ตุเมฺหว อุปาสมาโน,}
\hangingbody{สุตฺวาน อริยาน สุภาสิตานิ.}
\hangingbody{สกฺกสฺส ปุตฺโตมฺหิ มหานุภาโว,}
\hangingbody{มหาชุตีโก ติทิวูปปนฺโน.}
\hangingbody{ตุเมฺห ปน เสฏฺฐมุปาสมานา,}
\hangingbody{อนุตฺตรํ พฺรหฺมจริยํ จริตฺวา.}
\csromanhangingend

\csromangathameasure{หีนํ กายํ อุปปนฺนา ภวนฺโต, \\ อนานุโลมา ภวตูปปตฺติ. \\ ทุทฺทิฏฺฐรูปํ วต อทฺทสาม, \\ สหธมฺมิเก หีนกายูปปนฺเน. \\ คนฺธพฺพกายูปคตา ภวนฺโต, \\ เทวานมาคจฺฉถ ปาริจริยํ. \\ อคาเร วสโต มยฺหํ, \\ อิมํ ปสฺส วิเสสตํ. \\ อิตฺถี หุตฺวา สฺวชฺช ปุโมมฺหิ เทโว, \\ ทิพฺเพหิ กาเมหิ สมงฺคิภูโต. \\ เต โจทิตา โคตมสาวเกน, \\ สํเวคมาปาทุ สเมจฺจ โคปกํ. \\ หนฺท วิยายาม พฺยายาม, \\ มา โน มยํ ปรเปสฺสา อหุมฺหา. \\ เตสํ ทุเว วีริยมารภิํสุ, \\ อนุสฺสรํ โคตมสาสนานิ. \\ อิเธว จิตฺตานิ วิราชยิตฺวา, \\ กาเมสุ อาทีนวมทฺทสํสุ. \\ เต กามสํโยชนพนฺธนานิ, \\ ปาปิมโยคานิ ทุรจฺจยานิ. \\ นาโคว สนฺนานิ คุณานิ เฉตฺวา, \\ เทเว ตาวติํเส อติกฺกมิํสุ. \\ ส-อินฺทา เทวา สปชาปติกา, \\ สพฺเพ สุธมฺมาย สภายุปวิฏฺฐา. \\ เตสํ นิสินฺนานํ อภิกฺกมิํสุ, \\ วีรา วิราคา วิรชํ กโรนฺตา. \\ เต ทิสฺวา สํเวคมกาสิ วาสโว, \\ เทวาภิภู เทวคณสฺส มชฺเฌ.}
\csromangathagroup{\csromangathastanza{\csromanfolio{218}หีนํ กายํ อุปปนฺนา ภวนฺโต, \\ อนานุโลมา ภวตูปปตฺติ. \\ ทุทฺทิฏฺฐรูปํ วต อทฺทสาม, \\ สหธมฺมิเก หีนกายูปปนฺเน.}\csromangathastanza{คนฺธพฺพกายูปคตา ภวนฺโต, \\ เทวานมาคจฺฉถ ปาริจริยํ. \\ อคาเร วสโต มยฺหํ, \\ อิมํ ปสฺส วิเสสตํ.}\csromangathastanza{อิตฺถี หุตฺวา สฺวชฺช ปุโมมฺหิ เทโว, \\ ทิพฺเพหิ กาเมหิ สมงฺคิภูโต. \\ เต โจทิตา โคตมสาวเกน, \\ สํเวคมาปาทุ สเมจฺจ โคปกํ.}\csromangathastanza{หนฺท วิยายาม\footnote{วิคายาม (สฺยา), วิตายาม (อิ)} พฺยายาม\footnote{วิยายมาม (สี, อิ)}, \\ มา โน มยํ ปรเปสฺสา อหุมฺหา. \\ เตสํ ทุเว วีริยมารภิํสุ, \\ อนุสฺสรํ โคตมสาสนานิ.}\csromangathastanza{อิเธว จิตฺตานิ วิราชยิตฺวา, \\ กาเมสุ อาทีนวมทฺทสํสุ. \\ เต กามสํโยชนพนฺธนานิ, \\ ปาปิมโยคานิ ทุรจฺจยานิ.}\csromangathastanza{นาโคว สนฺนานิ คุณานิ\footnote{สนฺทานคุณานิ (สี, อิ), สนฺตานิ คุณานิ (สฺยา)} เฉตฺวา, \\ เทเว ตาวติํเส อติกฺกมิํสุ. \\ ส-อินฺทา เทวา สปชาปติกา, \\ สพฺเพ สุธมฺมาย สภายุปวิฏฺฐา.}\csromangathastanza{เตสํ นิสินฺนานํ อภิกฺกมิํสุ, \\ วีรา วิราคา วิรชํ กโรนฺตา. \\ เต ทิสฺวา สํเวคมกาสิ วาสโว, \\ เทวาภิภู เทวคณสฺส มชฺเฌ.}}

\csromanheader{2}{8. สกฺกปญฺหสุตฺต}
\csromangathameasure{อิเมหิ เต หีนกายูปปนฺนา, \\ เทเว ตาวติํเส อภิกฺกมนฺติ. \\ สํเวคชาตสฺส วโจ นิสมฺม, \\ โส โคปโก วาสวมชฺฌภาสิ. \\ พุทฺโธ ชนินฺทตฺถิ มนุสฺสโลเก, \\ กามาภิภู สกฺยมุนีติ ญายติ. \\ ตสฺเสว เต ปุตฺตา สติยา วิหีนา, \\ โจทิตา มยา เต สติมชฺฌลตฺถุํ. \\ ติณฺณํ เตสํ อาวสิเนตฺถ เอโก, \\ คนฺธพฺพกายูปคโต วสีโน. \\ เทฺว จ สมฺโพธิปถานุสาริโน, \\ เทเวปิ หีเฬนฺติ สมาหิตตฺตา. \\ เอตาทิสี ธมฺมปฺปกาสเนตฺถ, \\ น ตตฺถ กิํกงฺขติ โกจิ สาวโก. \\ นิติณฺณ-โอฆํ วิจิกิจฺฉฉินฺนํ, \\ พุทฺธํ นมสฺสาม ชินํ ชนินฺทํ. \\ ยํ เต ธมฺมํ อิธญฺญาย, \\ วิเสสํ อชฺฌคํสุ เต. \\ กายํ พฺรหฺมปุโรหิตํ, \\ ทุเว เตสํ วิเสสคู. \\ ตสฺส ธมฺมสฺส ปตฺติยา, \\ อาคตมฺหาสิ มาริส. \\ กตาวกาสา ภควตา, \\ ปญฺหํ ปุจฺเฉมุ มาริสาติ.}
\csromangathagroup{\csromangathastanza{\csromanfolio{219}อิเมหิ เต หีนกายูปปนฺนา, \\ เทเว ตาวติํเส อภิกฺกมนฺติ. \\ สํเวคชาตสฺส วโจ นิสมฺม, \\ โส โคปโก วาสวมชฺฌภาสิ.}\csromangathastanza{พุทฺโธ ชนินฺทตฺถิ มนุสฺสโลเก, \\ กามาภิภู สกฺยมุนีติ ญายติ. \\ ตสฺเสว เต ปุตฺตา สติยา วิหีนา, \\ โจทิตา มยา เต สติมชฺฌลตฺถุํ.}\csromangathastanza{ติณฺณํ เตสํ อาวสิเนตฺถ\footnote{อวสีเนตฺถ (อิ)} เอโก, \\ คนฺธพฺพกายูปคโต วสีโน. \\ เทฺว จ สมฺโพธิปถานุสาริโน, \\ เทเวปิ หีเฬนฺติ สมาหิตตฺตา.}\csromangathastanza{เอตาทิสี ธมฺมปฺปกาสเนตฺถ, \\ น ตตฺถ กิํกงฺขติ โกจิ สาวโก. \\ นิติณฺณ-โอฆํ วิจิกิจฺฉฉินฺนํ, \\ พุทฺธํ นมสฺสาม ชินํ ชนินฺทํ.}\csromangathastanza{ยํ เต ธมฺมํ อิธญฺญาย, \\ วิเสสํ อชฺฌคํสุ\footnote{อชฺฌคมํสุ (สฺยา)} เต. \\ กายํ พฺรหฺมปุโรหิตํ, \\ ทุเว เตสํ วิเสสคู.}\csromangathastanza{ตสฺส ธมฺมสฺส ปตฺติยา, \\ อาคตมฺหาสิ มาริส. \\ กตาวกาสา ภควตา, \\ ปญฺหํ ปุจฺเฉมุ มาริสาติ.}}

\proseitem{355}{อถ โข ภควโต เอตทโหสิ “ทีฆรตฺตํ วิสุทฺโธ โข อยํ ยกฺโข\footnote{สกฺโก (สี, สฺยา, อิ)}, ยํ กิญฺจิ มํ ปญฺหํ ปุจฺฉิสฺสติ, สพฺพํ ตํ อตฺถสญฺหิตํเยว ปุจฺฉิสฺสติ, โน อนตฺถสญฺหิตํ.\csromanspacer{} ยญฺจสฺสาหํ ปุฏฺโฐ พฺยากริสฺสามิ, ตํ ขิปฺปเมว อาชานิสฺสตี”ติ.}

\csromanheader{2}{356. อถ โข ภควา สกฺกํ เทวานมินฺทํ คาถาย อชฺฌภาสิ\csromandash{}}
\csromangathameasure{“ปุจฺฉ วาสว มํ ปญฺหํ, ยํ กิญฺจิ มนสิจฺฉสิ. \\ ตสฺส ตสฺเสว ปญฺหสฺส, อหํ อนฺตํ กโรมิ เต”ติ.}
\csromangathagroup{\csromangathastanza{\csromanfolio{220}“ปุจฺฉ วาสว มํ ปญฺหํ, ยํ กิญฺจิ มนสิจฺฉสิ. \\ ตสฺส ตสฺเสว ปญฺหสฺส, อหํ อนฺตํ กโรมิ เต”ติ.}}

\nitthitam{ปฐมภาณวาโร นิฏฺฐิโต.}
\proseitem{357}{กตาวกาโส สกฺโก เทวานมินฺโท ภควตา อิมํ ภควนฺตํ\footnote{เทวานมินฺโท ภควนฺตํ อิมํ (สี, อิ)} ปฐมํ ปญฺหํ อปุจฺฉิ\csromandash{}}

\prose{“กิํสํโยชนา นุ โข มาริส เทวา มนุสฺสา อสุรา นาคา คนฺธพฺพา เย จญฺเญ สนฺติ ปุถุกายา, เต ‘อเวรา อทณฺฑา อสปตฺตา อพฺยาปชฺชา วิหเรมุ อเวริโน’ติ อิติ จ เนสํ โหติ, อถ จ ปน สเวรา สทณฺฑา สสปตฺตา สพฺยาปชฺชา วิหรนฺติ สเวริโน”ติ.\csromanspacer{} อิตฺถํ สกฺโก เทวานมินฺโท ภควนฺตํ ปญฺหํ\footnote{อิมํ ปฐมํ ปญฺหํ (สี, อิ)} อปุจฺฉิ.\csromanspacer{} ตสฺส ภควา ปญฺหํ ปุฏฺโฐ พฺยากาสิ\csromandash{}}

\prose{“อิสฺสามจฺฉริยสํโยชนา โข เทวานมินฺท เทวา มนุสฺสา อสุรา นาคา คนฺธพฺพา เย จญฺเญ สนฺติ ปุถุกายา, เต ‘อเวรา อทณฺฑา อสปตฺตา อพฺยาปชฺชา วิหเรมุ อเวริโน’ติ อิติ จ เนสํ โหติ, อถ จ ปน สเวรา สทณฺฑา สสปตฺตา สพฺยาปชฺชา วิหรนฺติ สเวริโน”ติ.\csromanspacer{} อิตฺถํ ภควา สกฺกสฺส เทวานมินฺทสฺส ปญฺหํ ปุฏฺโฐ พฺยากาสิ.\csromanspacer{} อตฺตมโน สกฺโก เทวานมินฺโท ภควโต ภาสิตํ อภินนฺทิ อนุโมทิ “เอวเมตํ ภควา เอวเมตํ สุคต, ติณฺณา เมตฺถ กงฺขา วิคตา กถํกถา ภควโต ปญฺหเวยฺยากรณํ สุตฺวา”ติ.}

\proseitem{358}{อิติห สกฺโก เทวานมินฺโท ภควโต ภาสิตํ อภินนฺทิตฺวา อนุโมทิตฺวา ภควนฺตํ อุตฺตริ\footnote{อุตฺตริํ (สี, สฺยา, อิ)} ปญฺหํ อปุจฺฉิ\csromandash{}}

\prose{“อิสฺสามจฺฉริยํ ปน มาริส กิํนิทานํ กิํสมุทยํ กิํชาติกํ กิํปภวํ, กิสฺมิํ สติ อิสฺสามจฺฉริยํ โหติ, กิสฺมิํ อสติ อิสฺสามจฺฉริยํ น โหตี”ติ.\csromanspacer{} อิสฺสามจฺฉริยํ โข เทวานมินฺท ปิยาปฺปิยนิทานํ ปิยาปฺปิยสมุทยํ ปิยาปฺปิยชาติกํ ปิยาปฺปิยปภวํ, ปิยาปฺปิเย สติ อิสฺสามจฺฉริยํ โหติ, ปิยาปฺปิเย อสติ อิสฺสามจฺฉริยํ น โหตีติ.}

\csromanheader{2}{8. สกฺกปญฺหสุตฺต}
\prose{\csromanfolio{221}“ปิยาปฺปิยํ ปน มาริส กิํนิทานํ กิํสมุทยํ กิํชาติกํ กิํปภวํ, กิสฺมิํ สติ ปิยาปฺปิยํ โหติ, กิสฺมิํ อสติ ปิยาปฺปิยํ น โหตี”ติ.\csromanspacer{} ปิยาปฺปิยํ โข เทวานมินฺท ฉนฺทนิทานํ ฉนฺทสมุทยํ ฉนฺทชาติกํ ฉนฺทปภวํ, ฉนฺเท สติ ปิยาปฺปิยํ โหติ, ฉนฺเท อสติ ปิยาปฺปิยํ น โหตีติ.}

\prose{“ฉนฺโท ปน มาริส กิํนิทาโน กิํสมุทโย กิํชาติโก กิํปภโว, กิสฺมิํ สติ ฉนฺโท โหติ, กิสฺมิํ อสติ ฉนฺโท น โหตี”ติ.\csromanspacer{} ฉนฺโท โข เทวานมินฺท วิตกฺกนิทาโน วิตกฺกสมุทโย วิตกฺกชาติโก วิตกฺกปภโว, วิตกฺเก สติ ฉนฺโท โหติ, วิตกฺเก อสติ ฉนฺโท น โหตีติ.}

\prose{“วิตกฺโก ปน มาริส กิํนิทาโน กิํสมุทโย กิํชาติโก กิํปภโว, กิสฺมิํ สติ วิตกฺโก โหติ, กิสฺมิํ อสติ วิตกฺโก น โหตี”ติ.\csromanspacer{} วิตกฺโก โข เทวานมินฺท ปปญฺจสญฺญาสงฺขานิทาโน ปปญฺจสญฺญาสงฺขาสมุทโย ปปญฺจสญฺญาสงฺขาชาติโก ปปญฺจสญฺญาสงฺขาปภโว, ปปญฺจสญฺญาสงฺขาย สติ วิตกฺโก โหติ, ปปญฺจสญฺญาสงฺขาย อสติ วิตกฺโก น โหตีติ.}

\prose{“กถํ ปฏิปนฺโน ปน มาริส ภิกฺขุ ปปญฺจสญฺญาสงฺขานิโรธสารุปฺปคามินิํ ปฏิปทํ ปฏิปนฺโน โหตี”ติ.}

\csromanmatikaanchor{mtk.113}
\csromantocmarkref{section}{เวทนากมฺมฏฺฐาน}{mtk.113}
\bookmark[dest={mtk.113},level=1]{เวทนากมฺมฏฺฐาน}
\csromanheader{1}{\textbf{เวทนากมฺมฏฺฐาน}}
\proseitem{359}{โสมนสฺสํปาหํ\footnote{ปหํ (สี, อิ), จาหํ (สฺยา, กํ)} เทวานมินฺท ทุวิเธน วทามิ เสวิตพฺพํปิ อเสวิตพฺพํปิ.\csromanspacer{} โทมนสฺสํปาหํ เทวานมินฺท ทุวิเธน วทามิ เสวิตพฺพํปิ อเสวิตพฺพํปิ.\csromanspacer{} อุเปกฺขํปาหํ เทวานมินฺท ทุวิเธน วทามิ เสวิตพฺพํปิ อเสวิตพฺพํปิ.}

\proseitem{360}{“โสมนสฺสํปาหํ เทวานมินฺท ทุวิเธน วทามิ เสวิตพฺพํปิ อเสวิตพฺพํปี”ติ อิติ โข ปเนตํ วุตฺตํ, กิญฺเจตํ ปฏิจฺจ วุตฺตํ.\csromanspacer{} ตตฺถ ยํ ชญฺญา โสมนสฺสํ “อิมํ โข เม โสมนสฺสํ เสวโต อกุสลา ธมฺมา อภิวฑฺฒนฺติ, กุสลา ธมฺมา ปริหายนฺตี”ติ, เอวรูปํ โสมนสฺสํ น เสวิตพฺพํ.\csromanspacer{} ตตฺถ ยํ ชญฺญา โสมนสฺสํ “อิมํ โข เม โสมนสฺสํ \csromanfolio{222}เสวโต อกุสลา ธมฺมา ปริหายนฺติ, กุสลา ธมฺมา อภิวฑฺฒนฺตี”ติ, เอวรูปํ โสมนสฺสํ เสวิตพฺพํ.\csromanspacer{} ตตฺถ ยํ เจ สวิตกฺกํ สวิจารํ, ยํ เจ อวิตกฺกํ อวิจารํ, เย อวิตกฺเก อวิจาเร, เต\footnote{เส (สี, อิ)} ปณีตตเร. “โสมนสฺสํปาหํ เทวานมินฺท ทุวิเธน วทามิ เสวิตพฺพํปิ อเสวิตพฺพํปี”ติ อิติ ยํ ตํ วุตฺตํ, อิทเมตํ ปฏิจฺจ วุตฺตํ.}

\proseitem{361}{“โทมนสฺสํปาหํ เทวานมินฺท ทุวิเธน วทามิ เสวิตพฺพํปิ อเสวิตพฺพํปี”ติ อิติ โข ปเนตํ วุตฺตํ, กิญฺเจตํ ปฏิจฺจ วุตฺตํ.\csromanspacer{} ตตฺถ ยํ ชญฺญา โทมนสฺสํ “อิมํ โข เม โทมนสฺสํ เสวโต อกุสลา ธมฺมา อภิวฑฺฒนฺติ, กุสลา ธมฺมา ปริหายนฺตี”ติ, เอวรูปํ โทมนสฺสํ น เสวิตพฺพํ.\csromanspacer{} ตตฺถ ยํ ชญฺญา โทมนสฺสํ “อิมํ โข เม โทมนสฺสํ เสวโต อกุสลา ธมฺมา ปริหายนฺติ, กุสลา ธมฺมา อภิวฑฺฒนฺตี”ติ, เอวรูปํ โทมนสฺสํ เสวิตพฺพํ.\csromanspacer{} ตตฺถ ยํ เจ สวิตกฺกํ สวิจารํ, ยํ เจ อวิตกฺกํ อวิจารํ, เย อวิตกฺเก อวิจาเร, เต ปณีตตเร. “โทมนสฺสํปาหํ เทวานมินฺท ทุวิเธน วทามิ เสวิตพฺพํปิ อเสวิตพฺพํปี”ติ อิติ ยํ ตํ วุตฺตํ, อิทเมตํ ปฏิจฺจ วุตฺตํ.}

\proseitem{362}{“อุเปกฺขํปาหํ เทวานมินฺท ทุวิเธน วทามิ เสวิตพฺพํปิ อเสวิตพฺพํปี”ติ อิติ โข ปเนตํ วุตฺตํ, กิญฺเจตํ ปฏิจฺจ วุตฺตํ.\csromanspacer{} ตตฺถ ยํ ชญฺญา อุเปกฺขํ “อิมํ โข เม อุเปกฺขํ เสวโต อกุสลา ธมฺมา อภิวฑฺฒนฺติ, กุสลา ธมฺมา ปริหายนฺตี”ติ, เอวรูปา อุเปกฺขา น เสวิตพฺพา.\csromanspacer{} ตตฺถ ยํ ชญฺญา อุเปกฺขํ “อิมํ โข เม อุเปกฺขํ เสวโต อกุสลา ธมฺมา ปริหายนฺติ, กุสลา ธมฺมา อภิวฑฺฒนฺตี”ติ, เอวรูปา อุเปกฺขา เสวิตพฺพา.\csromanspacer{} ตตฺถ ยํ เจ สวิตกฺกํ สวิจารํ, ยํ เจ อวิตกฺกํ อวิจารํ, เย อวิตกฺเก อวิจาเร, เต ปณีตตเร. “อุเปกฺขํปาหํ เทวานมินฺท ทุวิเธน วทามิ เสวิตพฺพํปิ อเสวิตพฺพํปี”ติ อิติ ยํ ตํ วุตฺตํ, อิทเมตํ ปฏิจฺจ วุตฺตํ.}

\proseitem{363}{เอวํ ปฏิปนฺโน โข เทวานมินฺท ภิกฺขุ ปปญฺจสญฺญาสงฺขานิโรธสารุปฺปคามินิํ ปฏิปทํ ปฏิปนฺโน โหตีติ.\csromanspacer{} อิตฺถํ ภควา สกฺกสฺส เทวานมินฺทสฺส ปญฺหํ ปุฏฺโฐ พฺยากาสิ.\csromanspacer{} อตฺตมโน สกฺโก เทวานมินฺโท ภควโต}

\vspace{\tipitakaparskip}
\csromancenter{สกฺกปญฺหสุตฺต ภาสิตํ อภินนฺทิ อนุโมทิ “เอวเมตํ ภควา เอวเมตํ สุคต, ติณฺณา เมตฺถ กงฺขา วิคตา กถํกถา ภควโต ปญฺหเวยฺยากรณํ สุตฺวา”ติ.}
\csromanmatikaanchor{mtk.114}
\csromantocmarkref{section}{ปาติโมกฺขสํวร}{mtk.114}
\bookmark[dest={mtk.114},level=1]{ปาติโมกฺขสํวร}
\csromanheader{1}{\textbf{ปาติโมกฺขสํวร}}
\proseitem{364}{\csromanfolio{223}อิติห สกฺโก เทวานมินฺโท ภควโต ภาสิตํ อภินนฺทิตฺวา อนุโมทิตฺวา ภควนฺตํ อุตฺตริ ปญฺหํ อปุจฺฉิ\csromandash{}}

\prose{“กถํ ปฏิปนฺโน ปน มาริส ภิกฺขุ ปาติโมกฺขสํวราย ปฏิปนฺโน โหตี”ติ.\csromanspacer{} กายสมาจารํปาหํ เทวานมินฺท ทุวิเธน วทามิ เสวิตพฺพํปิ อเสวิตพฺพํปิ.\csromanspacer{} วจีสมาจารํปาหํ เทวานมินฺท ทุวิเธน วทามิ เสวิตพฺพํปิ อเสวิตพฺพํปิ.\csromanspacer{} ปริเยสนํปาหํ เทวานมินฺท ทุวิเธน วทามิ เสวิตพฺพํปิ อเสวิตพฺพํปิ.}

\prose{“กายสมาจารํปาหํ เทวานมินฺท ทุวิเธน วทามิ เสวิตพฺพํปิ อเสวิตพฺพํปี”ติ อิติ โข ปเนตํ วุตฺตํ, กิญฺเจตํ ปฏิจฺจ วุตฺตํ.\csromanspacer{} ตตฺถ ยํ ชญฺญา กายสมาจารํ “อิมํ โข เม กายสมาจารํ เสวโต อกุสลา ธมฺมา อภิวฑฺฒนฺติ, กุสลา ธมฺมา ปริหายนฺตี”ติ, เอวรูโป กายสมาจาโร น เสวิตพฺโพ.\csromanspacer{} ตตฺถ ยํ ชญฺญา กายสมาจารํ “อิมํ โข เม กายสมาจารํ เสวโต อกุสลา ธมฺมา ปริหายนฺติ, กุสลา ธมฺมา อภิวฑฺฒนฺตี”ติ, เอวรูโป กายสมาจาโร เสวิตพฺโพ. “กายสมาจารํปาหํ เทวานมินฺท ทุวิเธน วทามิ เสวิตพฺพํปิ อเสวิตพฺพํปี”ติ อิติ ยํ ตํ วุตฺตํ, อิทเมตํ ปฏิจฺจ วุตฺตํ.}

\prose{“วจีสมาจารํปาหํ เทวานมินฺท ทุวิเธน วทามิ เสวิตพฺพํปิ อเสวิตพฺพํปี”ติ อิติ โข ปเนตํ วุตฺตํ, กิญฺเจตํ ปฏิจฺจ วุตฺตํ.\csromanspacer{} ตตฺถ ยํ ชญฺญา วจีสมาจารํ “อิมํ โข เม วจีสมาจารํ เสวโต อกุสลา ธมฺมา อภิวฑฺฒนฺติ, กุสลา ธมฺมา ปริหายนฺตี”ติ, เอวรูโป วจีสมาจาโร น เสวิตพฺโพ.\csromanspacer{} ตตฺถ ยํ ชญฺญา วจีสมาจารํ “อิมํ โข เม วจีสมาจารํ เสวโต อกุสลา ธมฺมา ปริหายนฺติ, กุสลา ธมฺมา อภิวฑฺฒนฺตี”ติ, เอวรูโป วจีสมาจาโร เสวิตพฺโพ. “วจีสมาจารํปาหํ เทวานมินฺท ทุวิเธน วทามิ เสวิตพฺพํปิ อเสวิตพฺพํปี”ติ อิติ ยํ ตํ วุตฺตํ, อิทเมตํ ปฏิจฺจ วุตฺตํ.}

\csromanmatikaanchor{mtk.115}
\csromantocmarkref{section}{อินฺทฺริยสํวร}{mtk.115}
\bookmark[dest={mtk.115},level=1]{อินฺทฺริยสํวร}
\prose{\csromanfolio{224}“ปริเยสนํปาหํ เทวานมินฺท ทุวิเธน วทามิ เสวิตพฺพํปิ อเสวิตพฺพํปี”ติ อิติ โข ปเนตํ วุตฺตํ, กิญฺเจตํ ปฏิจฺจ วุตฺตํ.\csromanspacer{} ตตฺถ ยํ ชญฺญา ปริเยสนํ “อิมํ โข เม ปริเยสนํ เสวโต อกุสลา ธมฺมา อภิวฑฺฒนฺติ, กุสลา ธมฺมา ปริหายนฺตี”ติ, เอวรูปา ปริเยสนา น เสวิตพฺพา.\csromanspacer{} ตตฺถ ยํ ชญฺญา ปริเยสนํ “อิมํ โข เม ปริเยสนํ เสวโต อกุสลา ธมฺมา ปริหายนฺติ, กุสลา ธมฺมา อภิวฑฺฒนฺตี”ติ.\csromanspacer{} เอวรูปา ปริเยสนา เสวิตพฺพา. “ปริเยสนํปาหํ เทวานมินฺท ทุวิเธน วทามิ เสวิตพฺพํปิ อเสวิตพฺพํปี”ติ อิติ ยํ ตํ วุตฺตํ, อิทเมตํ ปฏิจฺจ วุตฺตํ.}

\prose{เอวํ ปฏิปนฺโน โข เทวานมินฺท ภิกฺขุ ปาติโมกฺขสํวราย ปฏิปนฺโน โหตีติ.\csromanspacer{} อิตฺถํ ภควา สกฺกสฺส เทวานมินฺทสฺส ปญฺหํ ปุฏฺโฐ พฺยากาสิ.\csromanspacer{} อตฺตมโน สกฺโก เทวานมินฺโท ภควโต ภาสิตํ อภินนฺทิ อนุโมทิ “เอวเมตํ ภควา เอวเมตํ สุคต, ติณฺณา เมตฺถ กงฺขา วิคตา กถํกถา ภควโต ปญฺหเวยฺยากรณํ สุตฺวา”ติ.}

\csromanheader{2}{\textbf{อินฺทฺรียสํวร}}
\proseitem{365}{อิติห สกฺโก เทวานมินฺโท ภควโต ภาสิตํ อภินนฺทิตฺวา อนุโมทิตฺวา ภควนฺตํ อุตฺตริ ปญฺหํ อปุจฺฉิ\csromandash{}}

\prose{“กถํ ปฏิปนฺโน ปน มาริส ภิกฺขุ อินฺทฺริยสํวราย ปฏิปนฺโน โหตี”ติ.\csromanspacer{} จกฺขุวิญฺเญยฺยํ รูปํปาหํ เทวานมินฺท ทุวิเธน วทามิ เสวิตพฺพํปิ อเสวิตพฺพํปิ.\csromanspacer{} โสตวิญฺเญยฺยํ สทฺทํปาหํ เทวานมินฺท ทุวิเธน วทามิ เสวิตพฺพํปิ อเสวิตพฺพํปิ.\csromanspacer{} ฆานวิญฺเญยฺยํ คนฺธํปาหํ เทวานมินฺท ทุวิเธน วทามิ เสวิตพฺพํปิ อเสวิตพฺพํปิ.\csromanspacer{} ชิวฺหาวิญฺเญยฺยํ รสํปาหํ เทวานมินฺท ทุวิเธน วทามิ เสวิตพฺพํปิ อเสวิตพฺพํปิ.\csromanspacer{} กายวิญฺเญยฺยํ โผฏฺฐพฺพํปาหํ เทวานมินฺท ทุวิเธน วทามิ เสวิตพฺพํปิ อเสวิตพฺพํปิ.\csromanspacer{} มโนวิญฺเญยฺยํ ธมฺมํปาหํ เทวานมินฺท ทุวิเธน วทามิ เสวิตพฺพํปิ อเสวิตพฺพํปีติ.}

\csromanheader{2}{เอวํ วุตฺเต สกฺโก เทวานมินฺโท ภควนฺตํ เอตทโวจ\csromandash{}}
\prose{“อิมสฺส โข อหํ ภนฺเต ภควตา สํขิตฺเตน ภาสิตสฺส เอวํ วิตฺถาเรน อตฺถํ อาชานามิ.\csromanspacer{} ยถารูปํ ภนฺเต จกฺขุวิญฺเญยฺยํ รูปํ เสวโต อกุสลา ธมฺมา อภิวฑฺฒนฺติ, กุสลา ธมฺมา ปริหายนฺติ, เอวรูปํ จกฺขุวิญฺเญยฺยํ}

\vspace{\tipitakaparskip}
\csromancenter{สกฺกปญฺหสุตฺต รูปํ น เสวิตพฺพํ.\csromanspacer{} ยถารูปญฺจ โข ภนฺเต จกฺขุวิญฺเญยฺยํ รูปํ เสวโต อกุสลา ธมฺมา ปริหายนฺติ, กุสลา ธมฺมา อภิวฑฺฒนฺติ, เอวรูปํ จกฺขุวิญฺเญยฺยํ รูปํ เสวิตพฺพํ.\csromanspacer{} ยถารูปญฺจ โข ภนฺเต โสตวิญฺเญยฺยํ สทฺทํ เสวโต ฯเปฯ ฆานวิญฺเญยฺยํ คนฺธํ เสวโต.\csromanspacer{} ชิวฺหาวิญฺเญยฺยํ รสํ เสวโต.\csromanspacer{} กายวิญฺเญยฺยํ โผฏฺฐพฺพํ เสวโต.\csromanspacer{} มโนวิญฺเญยฺยํ ธมฺมํ เสวโต อกุสลา ธมฺมา อภิวฑฺฒนฺติ, กุสลา ธมฺมา ปริหายนฺติ, เอวรูโป มโนวิญฺเญยฺโย ธมฺโม น เสวิตพฺโพ.\csromanspacer{} ยถารูปญฺจ โข ภนฺเต มโนวิญฺเญยฺยํ ธมฺมํ เสวโต อกุสลา ธมฺมา ปริหายนฺติ, กุสลา ธมฺมา อภิวฑฺฒนฺติ, เอวรูโป มโนวิญฺเญยฺโย ธมฺโม เสวิตพฺโพ.}
\prose{\csromanfolio{225}อิมสฺส โข เม ภนฺเต ภควตา สํขิตฺเตน ภาสิตสฺส เอวํ วิตฺถาเรน อตฺถํ อาชานโต ติณฺณา เมตฺถ กงฺขา วิคตา กถํกถา ภควโต ปญฺหเวยฺยากรณํ สุตฺวา”ติ.}

\proseitem{366}{อิติห สกฺโก เทวานมินฺโท ภควโต ภาสิตํ อภินนฺทิตฺวา อนุโมทิตฺวา ภควนฺตํ อุตฺตริ ปญฺหํ อปุจฺฉิ\csromandash{}}

\prose{“สพฺเพว นุ โข มาริส สมณพฺราหฺมณา เอกนฺตวาทา เอกนฺตสีลา เอกนฺตฉนฺทา เอกนฺต-อชฺโฌสานา”ติ.\csromanspacer{} น โข เทวานมินฺท สพฺเพ สมณพฺราหฺมณา เอกนฺตวาทา เอกนฺตสีลา เอกนฺตฉนฺทา เอกนฺต-อชฺโฌสานาติ.}

\prose{“กสฺมา ปน มาริส น สพฺเพ สมณพฺราหฺมณา เอกนฺตวาทา เอกนฺตสีลา เอกนฺตฉนฺทา เอกนฺต-อชฺโฌสานา”ติ.\csromanspacer{} อเนกธาตุ นานาธาตุ โข เทวานมินฺท โลโก.\csromanspacer{} ตสฺมิํ อเนกธาตุ นานาธาตุสฺมิํ โลเก ยํ ยเทว สตฺตา ธาตุํ อภินิวิสนฺติ, ตํ ตเทว ถามสา ปรามาสา อภินิวิสฺส โวหรนฺติ “อิทเมว สจฺจํ โมฆมญฺญนฺ”ติ.\csromanspacer{} ตสฺมา น สพฺเพ สมณพฺราหฺมณา เอกนฺตวาทา เอกนฺตสีลา เอกนฺตฉนฺทา เอกนฺต-อชฺโฌสานาติ.}

\prose{“สพฺเพว นุ โข มาริส สมณพฺราหฺมณา อจฺจนฺตนิฏฺฐา อจฺจนฺตโยคกฺเขมี อจฺจนฺตพฺรหฺมจารี อจฺจนฺตปริโยสานา”ติ.\csromanspacer{} น โข เทวานมินฺท สพฺเพ สมณพฺราหฺมณา อจฺจนฺตนิฏฺฐา อจฺจนฺตโยคกฺเขมี อจฺจนฺตพฺรหฺมจารี อจฺจนฺตปริโยสานาติ.}

\prose{\csromanfolio{226}“กสฺมา ปน มาริส น สพฺเพ สมณพฺราหฺมณา อจฺจนฺตนิฏฺฐา อจฺจนฺตโยคกฺเขมี อจฺจนฺตพฺรหฺมจารี อจฺจนฺตปริโยสานา”ติ.\csromanspacer{} เย โข เทวานมินฺท ภิกฺขู ตณฺหาสงฺขยวิมุตฺตา, เต อจฺจนฺตนิฏฺฐา อจฺจนฺตโยคกฺเขมี อจฺจนฺตพฺรหฺมจารี อจฺจนฺตปริโยสานา.\csromanspacer{} ตสฺมา น สพฺเพ สมณพฺราหฺมณา อจฺจนฺตนิฏฺฐา อจฺจนฺตโยคกฺเขมี อจฺจนฺตพฺรหฺมจารี อจฺจนฺตปริโยสานาติ.}

\prose{อิตฺถํ ภควา สกฺกสฺส เทวานมินฺทสฺส ปญฺหํ ปุฏฺโฐ พฺยากาสิ.\csromanspacer{} อตฺตมโน สกฺโก เทวานมินฺโท ภควโต ภาสิตํ อภินนฺทิ อนุโมทิ “เอวเมตํ ภควา เอวเมตํ สุคต, ติณฺณา เมตฺถ กงฺขา วิคตา กถํกถา ภควโต ปญฺหเวยฺยากรณํ สุตฺวา”ติ.}

\proseitem{367}{อิติห สกฺโก เทวานมินฺโท ภควโต ภาสิตํ อภินนฺทิตฺวา อนุโมทิตฺวา ภควนฺตํ เอตทโวจ\csromandash{}}

\prose{“เอชา ภนฺเต โรโค, เอชา คณฺโฑ, เอชา สลฺลํ, เอชา อิมํ ปุริสํ ปริกฑฺฒติ ตสฺส ตสฺเสว ภวสฺส อภินิพฺพตฺติยา.\csromanspacer{} ตสฺมา อยํ ปุริโส อุจฺจาวจมาปชฺชติ.\csromanspacer{} เยสาหํ ภนฺเต ปญฺหานํ อิโต พหิทฺธา อญฺเญสุ สมณพฺราหฺมเณสุ โอกาสกมฺมํปิ นาลตฺถํ, เต เม ภควตา พฺยากตา.\csromanspacer{} ทีฆรตฺตานุสยิหญฺจ ปน\footnote{ทีฆรตฺตานุปสฺสตา, ยญฺจ ปน (สฺยา), ทีฆรตฺตานุสยิโน, ยญฺจ ปน (สี, อิ)} เม วิจิกิจฺฉากถํกถาสลฺลํ, ตญฺจ ภควตา อพฺพุฬฺหนฺ”ติ.}

\prose{อภิชานาสิ โน ตฺวํ เทวานมินฺท อิเม ปเญฺห อญฺเญ สมณพฺราหฺมเณ ปุจฺฉิตาติ.\csromanspacer{} อภิชานามหํ ภนฺเต อิเม ปเญฺห อญฺเญ สมณพฺราหฺมเณ ปุจฺฉิตาติ.\csromanspacer{} ยถา กถํ ปน เต เทวานมินฺท พฺยากํสุ, สเจ เต อครุ ภาสสฺสูติ.\csromanspacer{} น โข เม ภนฺเต ครุ ยตฺถสฺส ภควา นิสินฺโน ภควนฺตรูโป วาติ.\csromanspacer{} เตน หิ เทวานมินฺท ภาสสฺสูติ.\csromanspacer{} เยสฺวาหํ\footnote{เยสาหํ (สี, สฺยา, อิ)} ภนฺเต มญฺญามิ สมณพฺราหฺมณา อารญฺญิกา ปนฺตเสนาสนาติ.\csromanspacer{} ตฺยาหํ อุปสงฺกมิตฺวา อิเม ปเญฺห ปุจฺฉามิ, เต มยา ปุฏฺฐา น สมฺปายนฺติ, อสมฺปายนฺตา มมํเยว ปฏิปุจฺฉนฺติ “โก นาโม อายสฺมา”ติ.\csromanspacer{} เตสาหํ ปุฏฺโฐ พฺยากโรมิ “อหํ โข มาริส สกฺโก เทวานมินฺโท”ติ.\csromanspacer{} เต มมํเยว อุตฺตริ ปฏิปุจฺฉนฺติ “กิํ ปนายสฺมา เทวานมินฺท\footnote{เทวานมินฺโท (สี, อิ)} กมฺมํ กตฺวา อิมํ ฐานํ ปตฺโต”ติ.\csromanspacer{} เตสาหํ ยถาสุตํ ยถาปริยตฺตํ ธมฺมํ เทเสมิ.\csromanspacer{} เต ตาวตเกเนว}

\vspace{\tipitakaparskip}
\csromancenter{สกฺกปญฺหสุตฺต อตฺตมนา โหนฺติ “สกฺโก จ โน เทวานมินฺโท ทิฏฺโฐ, ยญฺจ โน อปุจฺฉิมฺหา, ตญฺจ โน พฺยากาสี”ติ.\csromanspacer{} เต อญฺญทตฺถุ มมํเยว สาวกา สมฺปชฺชนฺติ, น จาหํ เตสํ.\csromanspacer{} อหํ โข ปน ภนฺเต ภควโต สาวโก โสตาปนฺโน อวินิปาตธมฺโม นิยโต สมฺโพธิปรายโณติ.}
\csromanmatikaanchor{mtk.116}
\csromantocmarkref{section}{โสมนสฺสปฏิลาภกถา}{mtk.116}
\bookmark[dest={mtk.116},level=1]{โสมนสฺสปฏิลาภกถา}
\csromanheader{1}{\textbf{โสมนสฺสปฏิลาภกถา}}
\proseitem{368}{\csromanfolio{227}อภิชานาสิ โน ตฺวํ เทวานมินฺท อิโต ปุพฺเพ เอวรูปํ เวทปฏิลาภํ โสมนสฺสปฏิลาภนฺติ.\csromanspacer{} อภิชานามหํ ภนฺเต อิโต ปุพฺเพ เอวรูปํ เวทปฏิลาภํ โสมนสฺสปฏิลาภนฺติ.\csromanspacer{} ยถา กถํ ปน ตฺวํ เทวานมินฺท อภิชานาสิ อิโต ปุพฺเพ เอวรูปํ เวทปฏิลาภํ โสมนสฺสปฏิลาภนฺติ.}

\prose{ภูตปุพฺพํ ภนฺเต เทวาสุรสงฺคาโม สมุปพฺยูโฬฺห\footnote{สมูปพฺพูโฬฺห (สี, อิ)} อโหสิ.\csromanspacer{} ตสฺมิํ โข ปน ภนฺเต สงฺคาเม เทวา ชินิํสุ, อสุรา ปราชยิํสุ\footnote{ปราชิํสุ (สี, อิ)}.\csromanspacer{} ตสฺส มยฺหํ ภนฺเต ตํ สงฺคามํ อภิวิชินิตฺวา วิชิตสงฺคามสฺส เอตทโหสิ “ยา เจว ทานิ ทิพฺพา โอชา ยา จ อสุรา โอชา, อุภยเมตํ\footnote{อุภยเมตฺถ (สฺยา)} เทวา ปริภุญฺชิสฺสนฺตี”ติ.\csromanspacer{} โส โข ปน เม ภนฺเต เวทปฏิลาโภ โสมนสฺสปฏิลาโภ สทณฺฑาวจโร สสตฺถาวจโร น นิพฺพิทาย น วิราคาย น นิโรธาย น อุปสมาย น อภิญฺญาย น สมฺโพธาย น นิพฺพานาย สํวตฺตติ.\csromanspacer{} โย โข ปน เม อยํ ภนฺเต ภควโต ธมฺมํ สุตฺวา เวทปฏิลาโภ โสมนสฺสปฏิลาโภ, โส อทณฺฑาวจโร อสตฺถาวจโร เอกนฺตนิพฺพิทาย วิราคาย นิโรธาย อุปสมาย อภิญฺญาย สมฺโพธาย นิพฺพานาย สํวตฺตตีติ.}

\proseitem{369}{กิํ ปน ตฺวํ เทวานมินฺท อตฺถวสํ สมฺปสฺสมาโน เอวรูปํ เวทปฏิลาภํ โสมนสฺสปฏิลาภํ ปเวเทสีติ.\csromanspacer{} ฉ โข อหํ ภนฺเต อตฺถวเส สมฺปสฺสมาโน เอวรูปํ เวทปฏิลาภํ โสมนสฺสปฏิลาภํ ปเวเทมิ.}

\csromangathameasure{อิเธว ติฏฺฐมานสฺส, เทวภูตสฺส เม สโต. \\ ปุนรายุ จ เม ลทฺโธ, เอวํ ชานาหิ มาริส.}
\csromangathagroup{\csromangathastanza{อิเธว ติฏฺฐมานสฺส, เทวภูตสฺส เม สโต. \\ ปุนรายุ จ เม ลทฺโธ, เอวํ ชานาหิ มาริส.}}

\prose{\csromanfolio{228}อิมํ โข อหํ ภนฺเต ปฐมํ อตฺถวสํ สมฺปสฺสมาโน เอวรูปํ เวทปฏิลาภํ โสมนสฺสปฏิลาภํ ปเวเทมิ.}

\csromangathameasure{จุตาหํ ทิวิยา กายา, อายุํ หิตฺวา อมานุสํ. \\ อมูโฬฺห คพฺภเมสฺสามิ, ยตฺถ เม รมตี มโน.}
\csromangathagroup{\csromangathastanza{จุตาหํ ทิวิยา กายา, อายุํ หิตฺวา อมานุสํ. \\ อมูโฬฺห คพฺภเมสฺสามิ, ยตฺถ เม รมตี มโน.}}

\prose{อิมํ โข อหํ ภนฺเต ทุติยํ อตฺถวสํ สมฺปสฺสมาโน เอวรูปํ เวทปฏิลาภํ โสมนสฺสปฏิลาภํ ปเวเทมิ.}

\csromangathameasure{สฺวาหํ อมูฬฺหปญฺญสฺส, วิหรํ สาสเน รโต. \\ ญาเยน วิหริสฺสามิ, สมฺปชาโน ปฏิสฺสโต.}
\csromangathagroup{\csromangathastanza{สฺวาหํ อมูฬฺหปญฺญสฺส\footnote{อมูฬฺหปญฺหสฺส (?)}, วิหรํ สาสเน รโต. \\ ญาเยน วิหริสฺสามิ, สมฺปชาโน ปฏิสฺสโต.}}

\prose{อิมํ โข อหํ ภนฺเต ตติยํ อตฺถวสํ สมฺปสฺสมาโน เอวรูปํ เวทปฏิลาภํ โสมนสฺสปฏิลาภํ ปเวเทมิ.}

\csromangathameasure{ญาเยน เม จรโต จ, สมฺโพธิ เจ ภวิสฺสติ. \\ อญฺญาตา วิหริสฺสามิ, เสฺวว อนฺโต ภวิสฺสติ.}
\csromangathagroup{\csromangathastanza{ญาเยน เม จรโต จ, สมฺโพธิ เจ ภวิสฺสติ. \\ อญฺญาตา วิหริสฺสามิ, เสฺวว อนฺโต ภวิสฺสติ.}}

\prose{อิมํ โข อหํ ภนฺเต จตุตฺถํ อตฺถวสํ สมฺปสฺสมาโน เอวรูปํ เวทปฏิลาภํ โสมนสฺสปฏิลาภํ ปเวเทมิ.}

\csromangathameasure{จุตาหํ มานุสา กายา, อายุํ หิตฺวาน มานุสํ. \\ ปุน เทโว ภวิสฺสามิ, เทวโลกมฺหิ อุตฺตโม.}
\csromangathagroup{\csromangathastanza{จุตาหํ มานุสา กายา, อายุํ หิตฺวาน มานุสํ. \\ ปุน เทโว ภวิสฺสามิ, เทวโลกมฺหิ อุตฺตโม.}}

\prose{อิมํ โข อหํ ภนฺเต ปญฺจมํ อตฺถวสํ สมฺปสฺสมาโน เอวรูปํ เวทปฏิลาภํ โสมนสฺสปฏิลาภํ ปเวเทมิ.}

\csromangathameasure{เต ปณีตตรา เทวา, อกนิฏฺฐา ยสสฺสิโน. \\ อนฺติเม วตฺตมานมฺหิ, โส นิวาโส ภวิสฺสติ.}
\csromangathagroup{\csromangathastanza{เต\footnote{เย (?)} ปณีตตรา เทวา, อกนิฏฺฐา ยสสฺสิโน. \\ อนฺติเม วตฺตมานมฺหิ, โส นิวาโส ภวิสฺสติ.}}

\prose{อิมํ โข อหํ ภนฺเต ฉฏฺฐํ อตฺถวสํ สมฺปสฺสมาโน เอวรูปํ เวทปฏิลาภํ โสมนสฺสปฏิลาภํ ปเวเทมิ.}

\prose{อิเม โข อหํ ภนฺเต ฉ อตฺถวเส สมฺปสฺสมาโน เอวรูปํ เวทปฏิลาภํ โสมนสฺสปฏิลาภํ ปเวเทมิ.}

\csromanheader{2}{8. สกฺกปญฺหสุตฺต}
\csromanhangingbegin
\hangingheaditem{370}{\csromanfolio{229}อปริโยสิตสงฺกปฺโป, วิจิกิจฺโฉ กถํกถี.}
\hangingbody{วิจริํ ทีฆมทฺธานํ, อเนฺวสนฺโต ตถาคตํ.}
\hangingbody{ยสฺสุ มญฺญามิ สมเณ, ปวิวิตฺตวิหาริโน.}
\hangingbody{สมฺพุทฺธา อิติ มญฺญาโน, คจฺฉามิ เต อุปาสิตุํ.}
\hangingbody{กถํ อาราธนา โหติ, กถํ โหติ วิราธนา.}
\hangingbody{อิติ ปุฏฺฐา น สมฺปายนฺติ1, มคฺเค ปฏิปทาสุ จ.}
\hangingbody{ตฺยสฺสุ ยทา มํ ชานนฺติ, สกฺโก เทวานมาคโต.}
\hangingbody{ตฺยสฺสุ มเมว ปุจฺฉนฺติ, กิํ กตฺวา ปาปุณี อิทํ.}
\hangingbody{เตสํ ยถาสุตํ ธมฺมํ, เทสยามิ ชเน สุตํ2.}
\hangingbody{เตน อตฺตมนา โหนฺติ, ทิฏฺโฐ โน วาสโวติ จ.}
\hangingbody{ยทา จ พุทฺธมทฺทกฺขิํ, วิจิกิจฺฉาวิตารณํ.}
\hangingbody{โสมฺหิ วีตภโย อชฺช, สมฺพุทฺธํ ปยิรุปาสิย3.}
\hangingbody{ตณฺหาสลฺลสฺส หนฺตารํ, พุทฺธํ อปฺปฏิปุคฺคลํ.}
\hangingbody{อหํ วนฺเท มหาวีรํ, พุทฺธมาทิจฺจพนฺธุนํ.}
\hangingbody{ยํ กโรมสิ พฺรหฺมุโน, สมํ เทเวหิ มาริส.}
\hangingbody{ตทชฺช ตุยฺหํ กสฺสาม4, หนฺท สามํ กโรม เต.}
\hangingbody{ตฺวเมว อสิ5 สมฺพุทฺโธ, ตุวํ สตฺถา อนุตฺตโร.}
\hangingbody{สเทวกสฺมิํ โลกสฺมิํ, นตฺถิ เต ปฏิปุคฺคโลติ.}
\csromanhangingend

\proseitem{371}{อถ โข สกฺโก เทวานมินฺโท ปญฺจสิขํ คนฺธพฺพปุตฺตํ อามนฺเตสิ “พหูปกาโร โข เมสิ ตฺวํ ตาต ปญฺจสิข, ยํ ตฺวํ ภควนฺตํ ปฐมํ ปสาเทสิ, ตยา ตาต ปฐมํ ปสาทิตํ ปจฺฉา มยํ ตํ ภควนฺตํ ทสฺสนาย อุปสงฺกมิมฺหา อรหนฺตํ สมฺมาสมฺพุทฺธํ.\csromanspacer{} เปตฺติเก วา ฐาเน \csromanfolio{230}ฐปยิสฺสามิ, คนฺธพฺพราชา ภวิสฺสสิ, ภทฺทญฺจ เต สูริยวจฺฉสํ ทมฺมิ, สา หิ เต อภิปตฺถิตา”ติ.}

\prose{อถ โข สกฺโก เทวานมินฺโท ปาณินา ปถวิํ ปรามสิตฺวา ติกฺขตฺตุํ อุทานํ อุทาเนสิ “นโม ตสฺส ภควโต อรหโต สมฺมาสมฺพุทฺธสฺสา”ติ.}

\prose{อิมสฺมิํ จ ปน เวยฺยากรณสฺมิํ ภญฺญมาเน สกฺกสฺส เทวานมินฺทสฺส วิรชํ วีตมลํ ธมฺมจกฺขุํ อุทปาทิ “ยํ กิญฺจิ สมุทยธมฺมํ, สพฺพํ ตํ นิโรธธมฺมนฺ”ติ, อญฺเญสญฺจ อสีติยา เทวตาสหสฺสานํ.\csromanspacer{} อิติ เย สกฺเกน เทวานมินฺเทน อชฺฌิฏฺฐปญฺหา ปุฏฺฐา, เต ภควตา พฺยากตา.\csromanspacer{} ตสฺมา อิมสฺส เวยฺยากรณสฺส สกฺกปญฺหาเตฺวว อธิวจนนฺติ.}

\nitthitam{\textbf{สกฺกปญฺหสุตฺตํ นิฏฺฐิตํ อฏฺฐมํ.}}
\csromanmatikaanchor{mtk.117}
\csromantocmarknopage{chapter}{9. มหาสติปฏฺฐานสุตฺต}{mtk.117}
\bookmark[dest={mtk.117},level=0]{9. มหาสติปฏฺฐานสุตฺต}
\chapterheadpageread{9. มหาสติปฏฺฐานสุตฺต}
\proseitem{372}{\csromanfolio{231}เอวํ เม สุตํ\csromandash{}เอกํ สมยํ ภควา กุรูสุ วิหรติ กมฺมาสธมฺมํ นาม กุรูนํ นิคโม.\csromanspacer{} ตตฺร โข ภควา ภิกฺขู อามนฺเตสิ “ภิกฺขโว”ติ. “ภทฺทนฺเต”ติ\footnote{ภทนฺเตติ (สี, สฺยา, อิ)} เต ภิกฺขู ภควโต ปจฺจสฺโสสุํ.\csromanspacer{} ภควา เอตทโวจ\csromandash{}}

\csromanmatikaanchor{mtk.118}
\csromantocmarkref{section}{อุทฺเทส}{mtk.118}
\bookmark[dest={mtk.118},level=1]{อุทฺเทส}
\csromanheader{1}{\textbf{อุทฺเทส}}
\proseitem{373}{เอกายโน อยํ ภิกฺขเว มคฺโค สตฺตานํ วิสุทฺธิยา โสกปริเทวานํ สมติกฺกมาย ทุกฺขโทมนสฺสานํ อตฺถงฺคมาย ญายสฺส อธิคมาย นิพฺพานสฺส สจฺฉิกิริยาย, ยทิทํ จตฺตาโร สติปฏฺฐานา.}

\prose{กตเม จตฺตาโร.\csromanspacer{} อิธ ภิกฺขเว ภิกฺขุ กาเย กายานุปสฺสี วิหรติ อาตาปี สมฺปชาโน สติมา วิเนยฺย โลเก อภิชฺฌาโทมนสฺสํ, เวทนาสุ เวทนานุปสฺสี วิหรติ อาตาปี สมฺปชาโน สติมา วิเนยฺย โลเก อภิชฺฌาโทมนสฺสํ, จิตฺเต จิตฺตานุปสฺสี วิหรติ อาตาปี สมฺปชาโน สติมา วิเนยฺย โลเก อภิชฺฌาโทมนสฺสํ, ธมฺเมสุ ธมฺมานุปสฺสี วิหรติ อาตาปี สมฺปชาโน สติมา วิเนยฺย โลเก อภิชฺฌาโทมนสฺสํ.}

\csromanheader{2}{อุทฺเทโส นิฐิโต.}
\csromansectionrule
\csromanmatikaanchor{mtk.119}
\csromantocmarkref{section}{กายานุปสฺสนา อานาปานปพฺพ}{mtk.119}
\bookmark[dest={mtk.119},level=1]{กายานุปสฺสนา อานาปานปพฺพ}
\csromanheader{1}{\textbf{กายานุปสฺสนา อานาปานปพฺพ}}
\proseitem{374}{กถญฺจ ปน ภิกฺขเว ภิกฺขุ กาเย กายานุปสฺสี วิหรติ.\csromanspacer{} อิธ ภิกฺขเว ภิกฺขุ อรญฺญคโต วา รุกฺขมูลคโต วา สุญฺญาคารคโตวา นิสีทติ ปลฺลงฺกํ อาภุชิตฺวา อุชุํ กายํ ปณิธาย ปริมุขํ สติํ อุปฏฺฐเปตฺวา.\csromanspacer{} โส สโตว อสฺสสติ, สโตว ปสฺสสติ.\csromanspacer{} ทีฆํ วา อสฺสสนฺโต “ทีฆํ อสฺสสามี”ติ ปชานาติ, ทีฆํ วา ปสฺสสนฺโต “ทีฆํ ปสฺสสามี”ติ ปชานาติ.\csromanspacer{} รสฺสํ วา อสฺสสนฺโต “รสฺสํ อสฺสสามี”ติ ปชานาติ, รสฺสํ วา ปสฺสสนฺโต “รสฺสํ ปสฺสสามี”ติ ปชานาติ. “สพฺพกายปฏิสํเวที อสฺสสิสฺสามี”ติ \csromanfolio{232}สิกฺขติ, “สพฺพกายปฏิสํเวที ปสฺสสิสฺสามี”ติ สิกฺขติ. “ปสฺสมฺภยํ กายสงฺขารํ อสฺสสิสฺสามี”ติ สิกฺขติ, “ปสฺสมฺภยํ กายสงฺขารํ ปสฺสสิสฺสามี”ติ สิกฺขติ.}

\prose{เสยฺยถาปิ ภิกฺขเว ทกฺโข ภมกาโร วา ภมการนฺเตวาสี วา ทีฆํ วา อญฺฉนฺโต “ทีฆํ อญฺฉามี”ติ ปชานาติ, รสฺสํ วา อญฺฉนฺโต “รสฺสํ อญฺฉามี”ติ ปชานาติ.\csromanspacer{} เอวเมว โข ภิกฺขเว ภิกฺขุ ทีฆํ วา อสฺสสนฺโต “ทีฆํ อสฺสสามี”ติ ปชานาติ, ทีฆํ วา ปสฺสสนฺโต “ทีฆํ ปสฺสสามี”ติ ปชานาติ.\csromanspacer{} รสฺสํ วา อสฺสสนฺโต “รสฺสํ อสฺสสามี”ติ ปชานาติ, รสฺสํ วา ปสฺสสนฺโต “รสฺสํ ปสฺสสามี”ติ ปชานาติ. “สพฺพกายปฏิสํเวที อสฺสสิสฺสามี”ติ สิกฺขติ, “สพฺพกายปฏิสํเวที ปสฺสสิสฺสามี”ติ สิกฺขติ. “ปสฺสมฺภยํ กายสงฺขารํ อสฺสสิสฺสามี”ติ สิกฺขติ, “ปสฺสมฺภยํ กายสงฺขารํ ปสฺสสิสฺสามี”ติ สิกฺขติ.\csromanspacer{} อิติ อชฺฌตฺตํ วา กาเย กายานุปสฺสี วิหรติ, พหิทฺธา วา กาเย กายานุปสฺสี วิหรติ, อชฺฌตฺตพหิทฺธา วา กาเย กายานุปสฺสี วิหรติ.\csromanspacer{} สมุทยธมฺมานุปสฺสี วา กายสฺมิํ วิหรติ, วยธมฺมานุปสฺสี วา กายสฺมิํ วิหรติ, สมุทยวยธมฺมานุปสฺสี วา กายสฺมิํ วิหรติ. “อตฺถิ กาโย”ติ วา ปนสฺส สติ ปจฺจุปฏฺฐิตา โหติ ยาวเทว ญาณมตฺตาย ปฏิสฺสติมตฺตาย.\csromanspacer{} อนิสฺสิโต จ วิหรติ, น จ กิญฺจิ โลเก อุปาทิยติ.\csromanspacer{} เอวมฺปิ โข\footnote{เอวมฺปิ (สี, สฺยา, อิ)} ภิกฺขเว ภิกฺขุ กาเย กายนุปสฺสี วิหรติ.}

\nitthitamruled{อานาปานปพฺพํ นิฏฺฐิตํ}
\csromanmatikaanchor{mtk.120}
\csromantocmarkref{section}{กายานุปสฺสนา อิริยาปถปพฺพ}{mtk.120}
\bookmark[dest={mtk.120},level=1]{กายานุปสฺสนา อิริยาปถปพฺพ}
\csromanheader{1}{\textbf{กายานุปสฺสนา อิริยาปถปพฺพ}}
\proseitem{375}{ปุน จปรํ ภิกฺขเว ภิกฺขุ คจฺฉนฺโต วา “คจฺฉามี”ติ ปชานาติ, ฐิโต วา “ฐิโตมฺหี”ติ ปชานาติ, นิสินฺโน วา “นิสินฺโนมฺหี”ติ ปชานาติ, สยาโน วา “สยาโนมฺหี”ติ ปชานาติ, ยถา ยถา วา ปนสฺส กาโย ปณิหิโต โหติ, ตถา ตถา นํ ปชานาติ.\csromanspacer{} อิติ อชฺฌตฺตํ วา กาเย กายานุปสฺสี วิหรติ, พหิทฺธา วา กาเย}

\vspace{\tipitakaparskip}
\csromancenter{มหาสติปฏฺฐานสุตฺต กายานุปสฺสี วิหรติ, อชฺฌตฺตพหิทฺธา วา กาเย กายานุปสฺสี วิหรติ.\csromanspacer{} สมุทยธมฺมานุปสฺสี วา กายสฺมิํ วิหรติ, วยธมฺมานุปสฺสี วา กายสฺมิํ วิหรติ, สมุทยวยธมฺมานุปสฺสี วา กายสฺมิํ วิหรติ. “อตฺถิ กาโย”ติ วา ปนสฺส สติ ปจฺจุปฏฺฐิตา โหติ ยาวเทว ญณมตฺตาย ปฏิสฺสติมตฺตาย.\csromanspacer{} อนิสฺสิโต จ วิหรติ, น จ กิญฺจิ โลเก อุปาทิยติ.\csromanspacer{} เอวมฺปิ โข ภิกฺขเว ภิกฺขุ กาเย กายนุปสฺสี วิหรติ.}
\nitthitamruled{อิริยาปถปพฺพํ นิฏฺฐิตํ.}
\csromanmatikaanchor{mtk.121}
\csromantocmarkref{section}{กายานุปสฺสนา สมฺปชานปพฺพ}{mtk.121}
\bookmark[dest={mtk.121},level=1]{กายานุปสฺสนา สมฺปชานปพฺพ}
\csromanheader{1}{\textbf{กายานุปสฺสนา สมฺปชานปพฺพ}}
\proseitem{376}{\csromanfolio{233}ปุน จปรํ ภิกฺขเว ภิกฺขุ อภิกฺกนฺเต ปฏิกฺกนฺเต สมฺปชานการี โหติ, อาโลกิเต วิโลกิเต สมฺปชานการี โหติ, สมิญฺชิเต ปสาริเต สมฺปชานการี โหติ, สงฺฆาฏิปตฺตจีวรธารเณ สมฺปชานการี โหติ, อสิเต ปีเต ขายิเต สายิเต สมฺปชานการี โหติ, อุจฺจารปสฺสาวกมฺเม สมฺปชานการี โหติ, คเต ฐิเต นิสินฺเน สุตฺเต ชาคริเต ภาสิเต ตุณฺหิภาเว สมฺปชานการี โหติ.\csromanspacer{} อิติ อชฺฌตฺตํ วา ฯเปฯ.\csromanspacer{} เอวมฺปิ โข ภิกฺขเว ภิกฺขุ กาเย กายานุปสฺสี วิหรติ.}

\nitthitamruled{สมฺปชานปพฺพํ นิฏฺฐิตํ.}
\csromanmatikaanchor{mtk.122}
\csromantocmarkref{section}{กายานุปสฺสนา ปฏิกูลมนสิการปพฺพ}{mtk.122}
\bookmark[dest={mtk.122},level=1]{กายานุปสฺสนา ปฏิกูลมนสิการปพฺพ}
\csromanheader{1}{\textbf{กายานุปสฺสนา ปฏิกูลมนสิการปพฺพ}}
\proseitem{377}{ปุน จปรํ ภิกฺขเว ภิกฺขุ อิมเมว กายํ อุทฺธํ ปาทตลา อโธ เกสมตฺถกา ตจปริยนฺตํ ปูรํ นานปฺปการสฺส อสุจิโน ปจฺจเวกฺขติ “อตฺถิ อิมสฺมิํ กาเย เกสา โลมา นขา ทนฺตา ตโจ, มํสํ นฺหารุ อฏฺฐิ อฏฺฐิมิญฺชํ วกฺกํ, หทยํ ยกนํ กิโลมกํ ปิหกํ ปปฺผาสํ, อนฺตํ \csromanfolio{234}อนฺตคุณํ อุทริยํ กรีสํ\footnote{กรีสํ มตฺถลุงฺคํ (ก)}, ปิตฺตํ เสมฺหํ ปุพฺโพ โลหิตํ เสโท เมโท, อสฺสุ วสา เขโฬ สิงฺฆาณิกา ลสิกา, มุตฺตนฺ”ติ.}

\prose{เสยฺยถาปิ ภิกฺขเว อุภโตมุขา ปุโตฬิ\footnote{มูโตฬี (สฺยา), มุโตลิ (อิ)} ปูรา นานาวิหิตสฺส ธญฺญสฺส.\csromanspacer{} เสยฺยถิทํ, สาลีนํ วีหีนํ มุคฺคานํ มาสานํ ติลานํ ตณฺฑุลานํ.\csromanspacer{} ตเมนํ จกฺขุมา ปุริโส มุญฺจิตฺวา ปจฺจเวกฺเขยฺย “อิเม สาลี อิเม วีหี อิเม มุคฺคา อิเม มาสา อิเม ติลา อิเม ตณฺฑุลา”ติ.\csromanspacer{} เอวเมว โข ภิกฺขเว ภิกฺขุ อิมเมว กายํ อุทฺธํ ปาทตลา อโธ เกสมตฺถกา ตจปริยนฺตํ ปูรํ นานปฺปการสฺส อสุจิโน ปจฺจเวกฺขติ “อตฺถิ อิมสฺมิํ กาเย เกสา โลมา ฯเปฯ มุตฺตนฺ”ติ.\csromanspacer{} อิติ อชฺฌตฺตํ วา ฯเปฯ.\csromanspacer{} เอวมฺปิ โข ภิกฺขเว ภิกฺขุ กาเย กายานุปสฺสี วิหรติ.}

\nitthitamruled{ปฏิกูลมนสิการปพฺพํ นิฏฺฐิตํ.}
\csromanmatikaanchor{mtk.123}
\csromantocmarkref{section}{กายานุปสฺสนา ธาตุมนสิการปพฺพ}{mtk.123}
\bookmark[dest={mtk.123},level=1]{กายานุปสฺสนา ธาตุมนสิการปพฺพ}
\csromanheader{1}{\textbf{กายานุปสฺสนา ธาตุมนสิการปพฺพ}}
\proseitem{378}{ปุน จปรํ ภิกฺขเว ภิกฺขุ อิมเมว กายํ ยถาฐิตํ ยถาปณิหิตํ ธาตุโส ปจฺจเวกฺขติ “อตฺถิ อิมสฺมิํ กาเย ปถวีธาตุ อาโปธาตุ เตโชธาตุ วาโยธาตู”ติ.}

\prose{เสยฺยถาปิ ภิกฺขเว ทกฺโข โคฆาตโก วา โคฆาตกนฺเตวาสี วา คาวิํ วธิตฺวา จตุมหาปเถ พิลโส วิภชิตฺวา นิสินฺโน อสฺส.\csromanspacer{} เอวเมว โข ภิกฺขเว ภิกฺขุ อิมเมว กายํ ยถาฐิตํ ยถาปณิหิตํ ธาตุโส ปจฺจเวกฺขติ “อตฺถิ อิมสฺมิํ กาเย ปถวีธาตุ อาโปธาตุ เตโชธาตุ วาโยธาตู”ติ.\csromanspacer{} อิติ อชฺฌตฺตํ วา กาเย กายานุปสฺสี วิหรติ ฯเปฯ.\csromanspacer{} เอวมฺปิ โข ภิกฺขเว ภิกฺขุ กาเย กายานุปสฺสี วิหรติ.}

\nitthitam{ธาตุมนสิการปพฺพํ นิฏฺฐิตํ.}
\csromanmatikaanchor{mtk.124}
\csromantocmarkref{section}{กายานุปสฺสนา นวสิวถิกปพฺพ}{mtk.124}
\bookmark[dest={mtk.124},level=1]{กายานุปสฺสนา นวสิวถิกปพฺพ}
\csromanheader{1}{9. มหาสติปฏฺฐานสุตฺต \textbf{กายานุปสฺสนา นวสิวถิกปพฺพ}}
\proseitem{379}{\csromanfolio{235}ปุน จปรํ ภิกฺขเว ภิกฺขุ เสยฺยถาปิ ปสฺเสยฺย สรีรํ สิวถิกาย ฉฑฺฑิตํ เอกาหมตํ วา ทฺวีหมตํ วา ตีหมตํ วา อุทฺธุมาตกํ วินีลกํ วิปุพฺพกชาตํ.\csromanspacer{} โส อิมเมว กายํ อุปสํหรติ “อยมฺปิ โข กาโย เอวํธมฺโม เอวํภาวี เอวํ-อนตีโต”ติ.\csromanspacer{} อิติ อชฺฌตฺตํ วา ฯเปฯ.\csromanspacer{} เอวมฺปิ โข ภิกฺขเว ภิกฺขุ กาเย กายานุปสฺสี วิหรติ.}

\prose{ปุน จปรํ ภิกฺขเว ภิกฺขุ เสยฺยถาปิ ปสฺเสยฺย สรีรํ สิวถิกาย ฉฑฺฑิตํ กาเกหิ วา ขชฺชมานํ กุลเลหิ วา ขชฺชมานํ คิชฺเฌหิ วา ขชฺชมานํ กงฺเกหิ วา ขชฺชมานํ สุนเขหิ วา ขชฺชมานํ พฺยคฺเฆหิ วา ขชฺชมานํ ทีปีหิ วา ขชฺชมานํ สิงฺคาเลหิ วา\footnote{คิชฺเฌหิ วา ขชฺชมานํ, สุวาเนหิ วา ขชฺชมานํ, สิคาเลหิ วา ขชฺชมานํ (สฺยา, อิ) 2. อปคตนฺหารุสมฺพนฺธานิ (สฺยา) 3. “อญฺเญน โคปฺผกฏฺฐิกนฺ”ติ อิทํ สี-สฺยา-อิ-โปตฺถเกสุ นตฺถิ. 4-4. อญฺเญน กฏฏฺฐินํ อญฺเญน ปิฏฺฐฏฺฐิกํ อญฺเญน กณฺฏกฏฺฐิกํ อญฺเญน ผาสุกฏฺฐิกํ อญฺเญน อุรฏฺฐิกํ อญฺเญน อํสฏฺฐิกํ อญฺเญน พาหุฏฺฐิกํ (สฺยา)} ขชฺชมานํ วิวิเธหิ วา ปาณกชาเตหิ ขชฺชมานํ.\csromanspacer{} โส อิมเมว กายํ อุปสํหรติ “อยมฺปิ โข กาโย เอวํธมฺโม เอวํภาวี เอวํอนตีโต”ติ.\csromanspacer{} อิติ อชฺฌตฺตํ วา ฯเปฯ.\csromanspacer{} เอวมฺปิ โข ภิกฺขเว ภิกฺขุ กาเย กายานุปสฺสี วิหรติ.}

\prose{ปุน จปรํ ภิกฺขเว ภิกฺขุ เสยฺยถาปิ ปสฺเสยฺย สรีรํ สิวถิกาย ฉฑฺฑิตํ อฏฺฐิกสงฺขลิกํ สมํสโลหิตํ นฺหารุสมฺพนฺธํ ฯเปฯ อฏฺฐิกสงฺขลิกํ นิมํสโลหิตมกฺขิตํ นฺหารุสมฺพนฺธํ ฯเปฯ อฏฺฐิกสงฺขลิกํ อปคตมํสโลหิตํ นฺหารุสมฺพนฺธํ ฯเปฯ อฏฺฐิกานิ อปคตสมฺพนฺธานิ2 ทิสา วิทิสา วิกฺขิตฺตานิ, อญฺเญน หตฺถฏฺฐิกํ อญฺเญน ปาทฏฺฐิกํ อญฺเญน โคปฺผกฏฺฐิกํ3 อญฺเญน ชงฺฆฏฺฐิกํ อญฺเญน อูรุฏฺฐิกํ อญฺเญน กฏิฏฺฐิกํ4 อญฺเญน ผาสุกฏฺฐิกํ อญฺเญน ปิฏฺฐิฏฺฐิกํ อญฺเญน ขนฺธฏฺฐิกํ4 อญฺเญน คีวฏฺฐิกํ อญฺเญน หนุกฏฺฐิกํ อญฺเญน ทนฺตฏฺฐิกํ อญฺเญน สีสกฏาหํ.\csromanspacer{} โส อิมเมว กายํ อุปสํหรติ “อยมฺปิ โข กาโย เอวํธมฺโม เอวํภาวี เอวํอนตีโต”ติ.\csromanspacer{} อิติ อชฺฌตฺตํ วา ฯเปฯ วิหรติ.}

\prose{\csromanfolio{236}ปุน จปรํ ภิกฺขเว ภิกฺขุ เสยฺยถาปิ ปสฺเสยฺย สรีรํ สิวถิกาย ฉฑฺฑิตํ อฏฺฐิกานิ เสตานิ สงฺขวณฺณปฏิภาคานิ ฯเปฯ อฏฺฐิกานิ ปุญฺชกิตานิ เตโรวสฺสิกานิ ฯเปฯ อฏฺฐิกานิ ปูตีนิ จุณฺณกชาตานิ.\csromanspacer{} โส อิมเมว กายํ อุปสํหรติ “อยมฺปิ โข กาโย เอวํธมฺโม เอวํภาวี เอวํอนตีโต”ติ.\csromanspacer{} อิติ อชฺฌตฺตํ วา กาเย กายานุปสฺสี วิหรติ, พหิทฺธา วา กาเย กายานุปสฺสี วิหรติ, อชฺฌตฺตพหิทฺธา วา กาเย กายานุปสฺสี วิหรติ.\csromanspacer{} สมุทยธมฺมานุปสฺสี วา กายสฺมิํ วิหรติ, วยธมฺมานุปสฺสี วา กายสฺมิํ วิหรติ, สมุทยวยธมฺมานุปสฺสี วา กายสฺมิํ วิหรติ. “อตฺถิ กาโย”ติ วา ปนสฺส สติ ปจฺจุปฏฺฐิตา โหติ ยาวเทว ญาณมตฺตาย ปฏิสฺสติมตฺตาย.\csromanspacer{} อนิสฺสิโต จ วิหรติ, น จ กิญฺจิ โลเก อุปาทิยติ.\csromanspacer{} เอวมฺปิ โข ภิกฺขเว ภิกฺขุ กาเย กายานุปสฺสี วิหรติ.}

\nitthitam{นวสิวถิกปพฺพํ นิฏฺฐิตํ.}
\nitthitamruled{จุทฺทส กายานุปสฺสนา นิฏฺฐิตา.}
\csromanmatikaanchor{mtk.125}
\csromantocmarkref{section}{เวทนานุปสฺสนา}{mtk.125}
\bookmark[dest={mtk.125},level=1]{เวทนานุปสฺสนา}
\csromanheader{1}{\textbf{เวทนานุปสฺสนา}}
\proseitem{380}{กถญฺจ ปน ภิกฺขเว ภิกฺขุ เวทนาสุ เวทนานุปสฺสี วิหรติ.\csromanspacer{} อิธ ภิกฺขเว ภิกฺขุ สุขํ วา เวทนํ เวทยมาโน “สุขํ เวทนํ เวทยามี”ติ ปชานาติ.\csromanspacer{} ทุกฺขํ วา เวทนํ เวทยมาโน “ทุกฺขํ เวทนํ เวทยามี”ติ ปชานาติ.\csromanspacer{} อทุกฺขมสุขํ วา เวทนํ เวทยมาโน “อทุกฺขมสุขํ เวทนํ เวทยามี”ติ ปชานาติ.\csromanspacer{} สามิสํ วา สุขํ เวทนํ เวทยมาโน “สามิสํ สุขํ เวทนํ เวทยามี”ติ ปชานาติ, นิรามิสํ วา สุขํ เวทนํ เวทยมาโน “นิรามิสํ สุขํ เวทนํ เวทยามี”ติ ปชานาติ.\csromanspacer{} สามิสํ วา ทุกฺขํ เวทนํ เวทยมาโน “สามิสํ ทุกฺขํ เวทนํ เวทยามี”ติ ปชานาติ, นิรามิสํ วา ทุกฺขํ เวทนํ เวทยมาโน “นิรามิสํ ทุกฺขํ เวทนํ เวทยามี”ติ ปชานาติ.\csromanspacer{} สามิสํ วา อทุกฺขมสุขํ เวทนํ เวทยมาโน “สามิสํ อทุกฺขมสุขํ เวทนํ เวทยามี”ติ ปชานาติ, นิรามิสํ วา อทุกฺขมสุขํ เวทนํ เวทยมาโน “นิรามิสํ อทุกฺขมสุขํ เวทนํ เวทยามี”ติ ปชานาติ.\csromanspacer{} อิติ อชฺฌตฺตํ วา เวทนาสุ เวทนานุปสฺสี วิหรติ, พหิทฺธา วา}

\vspace{\tipitakaparskip}
\csromancenter{มหาสติปฏฺฐานสุตฺต เวทนาสุ เวทนานุปสฺสี วิหรติ, อชฺฌตฺตพหิทฺธา วา เวทนาสุ เวทนานุปสฺสี วิหรติ.\csromanspacer{} สมุทยธมฺมานุปสฺสี วา เวทนาสุ วิหรติ, วยธมฺมานุปสฺสี วา เวทนาสุ วิหรติ, สมุทยวยธมฺมานุปสฺสี วา เวทนาสุ วิหรติ. “อตฺถิ เวทนา”ติ วา ปนสฺส สติ ปจฺจุปฏฺฐิตา โหติ ยาวเทว ญาณมตฺตาย ปฏิสฺสติมตฺตาย.\csromanspacer{} อนิสฺสิโต จ วิหรติ, น จ กิญฺจิ โลเก อุปาทิยติ.\csromanspacer{} เอวมฺปิ โข ภิกฺขเว ภิกฺขุ เวทนาสุ เวทนานุปสฺสี วิหรติ.}
\nitthitamruled{เวทนานุปสฺสนา นิฏฺฐิตา.}
\csromanmatikaanchor{mtk.126}
\csromantocmarkref{section}{จิตฺตานุปสฺสนา}{mtk.126}
\bookmark[dest={mtk.126},level=1]{จิตฺตานุปสฺสนา}
\csromanheader{1}{\textbf{จิตฺตานุปสฺสนา}}
\proseitem{381}{\csromanfolio{237}กถญฺจ ปน ภิกฺขเว ภิกฺขุ จิตฺเต จิตฺตานุปสฺสี วิหรติ.\csromanspacer{} อิธ ภิกฺขเว ภิกฺขุ สราคํ วา จิตฺตํ “สราคํ จิตฺตนฺ”ติ ปชานาติ, วีตราคํ วา จิตฺตํ “วีตราคํ จิตฺตนฺ”ติ ปชานาติ.\csromanspacer{} สโทสํ วา จิตฺตํ “สโทสํ จิตฺตนฺ”ติ ปชานาติ, วีตโทสํ วา จิตฺตํ “วีตโทสํ จิตฺตนฺ”ติ ปชานาติ.\csromanspacer{} สโมหํ วา จิตฺตํ “สโมหํ จิตฺตนฺ”ติ ปชานาติ, วีตโมหํ วา จิตฺตํ “วีตโมหํ จิตฺตนฺ”ติ ปชานาติ.\csromanspacer{} สํขิตฺตํ วา จิตฺตํ “สํขิตฺตํ จิตฺตนฺ”ติ ปชานาติ, วิกฺขิตฺตํ วา จิตฺตํ “วิกฺขิตฺตํ จิตฺตนฺ”ติ ปชานาติ.\csromanspacer{} มหคฺคตํ วา จิตฺตํ “มหคฺคตํ จิตฺตนฺ”ติ ปชานาติ, อมหคฺคตํ วา จิตฺตํ “อมหคฺคตํ จิตฺตนฺ”ติ ปชานาติ.\csromanspacer{} ส-อุตฺตรํ วา จิตฺตํ “ส-อุตฺตรํ จิตฺตนฺ”ติ ปชานาติ, อนุตฺตรํ วา จิตฺตํ “อนุตฺตรํ จิตฺตนฺ”ติ ปชานาติ.\csromanspacer{} สมาหิตํ วา จิตฺตํ “สมาหิตํ จิตฺตนฺ”ติ ปชานาติ, อสมาหิตํ วา จิตฺตํ “อสมาหิตํ จิตฺตนฺ”ติ ปชานาติ.\csromanspacer{} วิมุตฺตํ วา จิตฺตํ “วิมุตฺตํ จิตฺตนฺ”ติ ปชานาติ, อวิมุตฺตํ วา จิตฺตํ “อวิมุตฺตํ จิตฺตนฺ”ติ ปชานาติ.\csromanspacer{} อิติ อชฺฌตฺตํ วา จิตฺเต จิตฺตานุปสฺสี วิหรติ, พหิทฺธา วา จิตฺเต จิตฺตานุปสฺสี วิหรติ, อชฺฌตฺตพหิทฺธา วา จิตฺเต จิตฺตานุปสฺสี วิหรติ.\csromanspacer{} สมุทยธมฺมานุปสฺสี วา จิตฺตสฺมิํ วิหรติ, วยธมฺมานุปสฺสี วา จิตฺตสฺมิํ วิหรติ, สมุทยวยธมฺมานุปสฺสี วา จิตฺตสฺมิํ วิหรติ. “อตฺถิ จิตฺตนฺ”ติ วา ปนสฺส สติ ปจฺจุปฏฺฐิตา โหติ ยาวเทว ญาณมตฺตาย ปฏิสฺสติมตฺตาย.\csromanspacer{} อนิสฺสิโต จ วิหรติ, น จ กิญฺจิ โลเก อุปาทิยติ.\csromanspacer{} เอวมฺปิ โข ภิกฺขเว ภิกฺขุ จิตฺเต จิตฺตานุปสฺสี วิหรติ.}

\nitthitam{\textbf{จิตฺตานุปสฺสนา} นิฏฺฐิตา.}
\csromanmatikaanchor{mtk.127}
\csromantocmarkref{section}{ธมฺมานุปสฺสนา นีวรณปพฺพ}{mtk.127}
\bookmark[dest={mtk.127},level=1]{ธมฺมานุปสฺสนา นีวรณปพฺพ}
\csromanheader{1}{\textbf{ธมฺมานุปสฺสนา นีวรณปพฺพ}}
\proseitem{382}{\csromanfolio{238}กถญฺจ ปน ภิกฺขเว ภิกฺขุ ธมฺเมสุ ธมฺมานุปสฺสี วิหรติ.\csromanspacer{} อิธ ภิกฺขเว ภิกฺขุ ธมฺเมสุ ธมฺมานุปสฺสี วิหรติ ปญฺจสุ นีวรเณสุ.\csromanspacer{} กถญฺจ ปน ภิกฺขเว ภิกฺขุ ธมฺเมสุ ธมฺมานุปสฺสี วิหรติ ปญฺจสุ นีวรเณสุ.}

\prose{อิธ ภิกฺขเว ภิกฺขุ สนฺตํ วา อชฺฌตฺตํ กามจฺฉนฺทํ “อตฺถิ เม อชฺฌตฺตํ กามจฺฉนฺโท”ติ ปชานาติ, อสนฺตํ วา อชฺฌตฺตํ กามจฺฉนฺทํ “นตฺถิ เม อชฺฌตฺตํ กามจฺฉนฺโท”ติ ปชานาติ, ยถา จ อนุปฺปนฺนสฺส กามจฺฉนฺทสฺส อุปฺปาโท โหติ ตญฺจ ปชานาติ, ยถา จ อุปฺปนฺนสฺส กามจฺฉนฺทสฺส ปหานํ โหติ ตญฺจ ปชานาติ, ยถา จ ปหีนสฺส กามจฺฉนฺทสฺส อายติํ อนุปฺปาโท โหติ ตญฺจ ปชานาติ.}

\prose{สนฺตํ วา อชฺฌตฺตํ พฺยาปาทํ “อตฺถิ เม อชฺฌตฺตํ พฺยาปาโท”ติ ปชานาติ, อสนฺตํ วา อชฺฌตฺตํ พฺยาปาทํ “นตฺถิ เม อชฺฌตฺตํ พฺยาปาโท”ติ ปชานาติ, ยถา จ อนุปฺปนฺนสฺส พฺยาปาทสฺส อุปฺปาโท โหติ ตญฺจ ปชานาติ, ยถา จ อุปฺปนฺนสฺส พฺยาปาทสฺส ปหานํ โหติ ตญฺจ ปชานาติ, ยถา จ ปหีนสฺส พฺยาปาทสฺส อายติํ อนุปฺปาโท โหติ ตญฺจ ปชานาติ.}

\prose{สนฺตํ วา อชฺฌตฺตํ ถินมิทฺธํ “อตฺถิ เม อชฺฌตฺตํ ถินมิทฺธนฺ”ติ ปชานาติ, อสนฺตํ วา อชฺฌตฺตํ ถินมิทฺธํ “นตฺถิ เม อชฺฌตฺตํ ถินมิทฺธนฺ”ติ ปชานาติ, ยถา จ อนุปฺปนฺนสฺส ถินมิทฺธสฺส อุปฺปาโท โหติ ตญฺจ ปชานาติ, ยถา จ อุปฺปนฺนสฺส ถินมิทฺธสฺส ปหานํ โหติ ตญฺจ ปชานาติ, ยถา จ ปหีนสฺส ถินมิทฺธสฺส อายติํ อนุปฺปาโท โหติ ตญฺจ ปชานาติ.}

\prose{สนฺตํ วา อชฺฌตฺตํ อุทฺธจฺจกุกฺกุจฺจํ “อตฺถิ เม อชฺฌตฺตํ อุทฺธจฺจกุกฺกุจฺจนฺ”ติ ปชานาติ, อสนฺตํ วา อชฺฌตฺตํ อุทฺธจฺจกุกฺกุจฺจํ “นตฺถิ เม อชฺฌตฺตํ อุทฺธจฺจกุกฺกุจฺจนฺ”ติ ปชานาติ, ยถา จ อนุปฺปนฺนสฺส อุทฺธจฺจกุกฺกุจฺจสฺส อุปฺปาโท โหติ ตญฺจ ปชานาติ, ยถา จ อุปฺปนฺนสฺส อุทฺธจฺจกุกฺกุจฺจสฺส ปหานํ โหติ ตญฺจ ปชานาติ, ยถา จ ปหีนสฺส อุทฺธจฺจกุกฺกุจฺจสฺส อายติํ อนุปฺปาโท โหติ ตญฺจ ปชานาติ.}

\prose{สนฺตํ วา อชฺฌตฺตํ วิจิกิจฺฉํ “อตฺถิ เม อชฺฌตฺตํ วิจิกิจฺฉา”ติ ปชานาติ, อสนฺตํ วา อชฺฌตฺตํ วิจิกิจฺฉํ “นตฺถิ เม อชฺฌตฺตํ วิจิกิจฺฉา”ติ ปชานาติ, ยถา จ อนุปฺปนฺนาย วิจิกิจฺฉาย อุปฺปาโท โหติ ตญฺจ ปชานาติ,}

\vspace{\tipitakaparskip}
\csromancenter{มหาสติปฏฺฐานสุตฺต ยถา จ อุปฺปนฺนาย วิจิกิจฺฉาย ปหานํ โหติ ตญฺจ ปชานาติ, ยถา จ ปหีนาย วิจิกิจฺฉาย อายติํ อนุปฺปาโท โหติ ตญฺจ ปชานาติ.}
\prose{\csromanfolio{239}อิติ อชฺฌตฺตํ วา ธมฺเมสุ ธมฺมานุปสฺสี วิหรติ, พหิทฺธา วา ธมฺเมสุ ธมฺมานุปสฺสี วิหรติ, อชฺฌตฺตพหิทฺธา วา ธมฺเมสุ ธมฺมานุปสฺสี วิหรติ.\csromanspacer{} สมุทยธมฺมานุปสฺสี วา ธมฺเมสุ วิหรติ, วยธมฺมานุปสฺสี วา ธมฺเมสุ วิหรติ, สมุทยวยธมฺมานุปสฺสี วา ธมฺเมสุ วิหรติ. “อตฺถิ ธมฺมา”ติ วา ปนสฺส สติ ปจฺจุปฏฺฐิตา โหติ ยาวเทว ญาณมตฺตาย ปฏิสฺสติมตฺตาย.\csromanspacer{} อนิสฺสิโต จ วิหรติ, น จ กิญฺจิ โลเก อุปาทิยติ.\csromanspacer{} เอวมฺปิ โข ภิกฺขเว ภิกฺขุ ธมฺเมสุ ธมฺมานุปสฺสี วิหรติ ปญฺจสุ นีวรเณสุ.}

\nitthitamruled{นีวรณปพฺพํ นิฏฺฐิตํ.}
\csromanmatikaanchor{mtk.128}
\csromantocmarkref{section}{ธมฺมานุปสฺสนา ขนฺธปพฺพ}{mtk.128}
\bookmark[dest={mtk.128},level=1]{ธมฺมานุปสฺสนา ขนฺธปพฺพ}
\csromanheader{1}{\textbf{ธมฺมานุปสฺสนา ขนฺธปพฺพ}}
\proseitem{383}{ปุน จปรํ ภิกฺขเว ภิกฺขุ ธมฺเมสุ ธมฺมานุปสฺสี วิหรติ ปญฺจสุ อุปาทานกฺขนฺเธสุ.\csromanspacer{} กถญฺจ ปน ภิกฺขเว ภิกฺขุ ธมฺเมสุ ธมฺมานุปสฺสี วิหรติ ปญฺจสุ อุปาทานกฺขนฺเธสุ.\csromanspacer{} อิธ ภิกฺขเว ภิกฺขุ อิติ รูปํ, อิติ รูปสฺส สมุทโย, อิติ รูปสฺส อตฺถงฺคโม.\csromanspacer{} อิติ เวทนา, อิติ เวทนาย สมุทโย, อิติ เวทนาย อตฺถงฺคโม.\csromanspacer{} อิติ สญฺญา, อิติ สญฺญาย สมุทโย, อิติ สญฺญาย อตฺถงฺคโม.\csromanspacer{} อิติ สงฺขารา, อิติ สงฺขารานํ สมุทโย, อิติ สงฺขารานํ อตฺถงฺคโม.\csromanspacer{} อิติ วิญฺญาณํ, อิติ วิญฺญาณสฺส สมุทโย, อิติ วิญฺญาณสฺส อตฺถงฺคโมติ.\csromanspacer{} อิติ อชฺฌตฺตํ วา ธมฺเมสุ ธมฺมานุปสฺสี วิหรติ, พหิทฺธา วา ธมฺเมสุ ธมฺมานุปสฺสี วิหรติ, อชฺฌตฺตพหิทฺธา วา ธมฺเมสุ ธมฺมานุปสฺสี วิหรติ.\csromanspacer{} สมุทยธมฺมานุปสฺสี วา ธมฺเมสุ วิหรติ, วยธมฺมานุปสฺสี วา ธมฺเมสุ วิหรติ, สมุทยวยธมฺมานุปสฺสี วา ธมฺเมสุ วิหรติ. “อตฺถิ ธมฺมา”ติ วา ปนสฺส สติ ปจฺจุปฏฺฐิตา โหติ ยาวเทว ญาณมตฺตาย ปฏิสฺสติมตฺตาย.\csromanspacer{} อนิสฺสิโต จ วิหรติ, น จ กิญฺจิ โลเก อุปาทิยติ.\csromanspacer{} เอวมฺปิ โข ภิกฺขเว ภิกฺขุ ธมฺเมสุ ธมฺมานุปสฺสี วิหรติ ปญฺจสุ อุปาทานกฺขนฺเธสุ.}

\nitthitam{ขนฺธปพฺพํ นิฏฺฐิตํ.}
\csromanmatikaanchor{mtk.129}
\csromantocmarkref{section}{ธมฺมานุปสฺสนา อายตนปพฺพ}{mtk.129}
\bookmark[dest={mtk.129},level=1]{ธมฺมานุปสฺสนา อายตนปพฺพ}
\csromanheader{1}{\textbf{ธมฺมานุปสฺสนา อายตนปพฺพ}}
\proseitem{384}{\csromanfolio{240}ปุน จปรํ ภิกฺขเว ภิกฺขุ ธมฺเมสุ ธมฺมานุปสฺสี วิหรติ ฉสุ อชฺฌตฺติกพาหิเรสุ อายตเนสุ.\csromanspacer{} กถญฺจ ปน ภิกฺขเว ภิกฺขุ ธมฺเมสุ ธมฺมานุปสฺสี วิหรติ ฉสุ อชฺฌตฺติกพาหิเรสุ อายตเนสุ.}

\prose{อิธ ภิกฺขเว ภิกฺขุ จกฺขุญฺจ ปชานาติ, รูเป จ ปชานาติ, ยญฺจ ตทุภยํ ปฏิจฺจ อุปฺปชฺชติ สํโยชนํ ตญฺจ ปชานาติ, ยถา จ อนุปฺปนฺนสฺส สํโยชนสฺส อุปฺปาโท โหติ ตญฺจ ปชานาติ, ยถา จ อุปฺปนฺนสฺส สํโยชนสฺส ปหานํ โหติ ตญฺจ ปชานาติ, ยถา จ ปหีนสฺส สํโยชนสฺส อายติํ อนุปฺปาโท โหติ ตญฺจ ปชานาติ.}

\prose{โสตญฺจ ปชานาติ, สทฺเท จ ปชานาติ, ยญฺจ ตทุภยํ ปฏิจฺจ อุปฺปชฺชติ สํโยชนํ ตญฺจ ปชานาติ, ยถา จ อนุปฺปนฺนสฺส สํโยชนสฺส อุปฺปาโท โหติ ตญฺจ ปชานาติ, ยถา จ อุปฺปนฺนสฺส สํโยชนสฺส ปหานํ โหติ ตญฺจ ปชานาติ, ยถา จ ปหีนสฺส สํโยชนสฺส อายติํ อนุปฺปาโท โหติ ตญฺจ ปชานาติ.}

\prose{ฆานญฺจ ปชานาติ, คนฺเธ จ ปชานาติ, ยญฺจ ตทุภยํ ปฏิจฺจ อุปฺปชฺชติ สํโยชนํ ตญฺจ ปชานาติ, ยถา จ อนุปฺปนฺนสฺส สํโยชนสฺส อุปฺปาโท โหติ ตญฺจ ปชานาติ, ยถา จ อุปฺปนฺนสฺส สํโยชนสฺส ปหานํ โหติ ตญฺจ ปชานาติ, ยถา จ ปหีนสฺส สํโยชนสฺส อายติํ อนุปฺปาโท โหติ ตญฺจ ปชานาติ.}

\prose{ชิวฺหญฺจ ปชานาติ, รเส จ ปชานาติ, ยญฺจ ตทุภยํ ปฏิจฺจ อุปฺปชฺชติ สํโยชนํ ตญฺจ ปชานาติ, ยถา จ อนุปฺปนฺนสฺส สํโยชนสฺส อุปฺปาโท โหติ ตญฺจ ปชานาติ, ยถา จ อุปฺปนฺนสฺส สํโยชนสฺส ปหานํ โหติ ตญฺจ ปชานาติ, ยถา จ ปหีนสฺส สํโยชนสฺส อายติํ อนุปฺปาโท โหติ ตญฺจ ปชานาติ.}

\prose{กายญฺจ ปชานาติ, โผฏฺฐพฺเพ จ ปชานาติ, ยญฺจ ตทุภยํ ปฏิจฺจ อุปฺปชฺชติ สํโยชนํ ตญฺจ ปชานาติ, ยถา จ อนุปฺปนฺนสฺส สํโยชนสฺส อุปฺปาโท โหติ ตญฺจ ปชานาติ, ยถา จ อุปฺปนฺนสฺส สํโยชนสฺส ปหานํ โหติ ตญฺจ ปชานาติ, ยถา จ ปหีนสฺส สํโยชนสฺส อายติํ อนุปฺปาโท โหติ, ตญฺจ ปชานาติ.}

\csromanheader{2}{9. มหาสติปฏฺฐานสุตฺต}
\prose{\csromanfolio{241}มนญฺจ ปชานาติ, ธมฺเม จ ปชานาติ, ยญฺจ ตทุภยํ ปฏิจฺจ อุปฺปชฺชติ สํโยชนํ ตญฺจ ปชานาติ, ยถา จ อนุปฺปนฺนสฺส สํโยชนสฺส อุปฺปาโท โหติ ตญฺจ ปชานาติ, ยถา จ อุปฺปนฺนสฺส สํโยชนสฺส ปหานํ โหติ ตญฺจ ปชานาติ, ยถา จ ปหีนสฺส สํโยชนสฺส อายติํ อนุปฺปาโท โหติ ตญฺจ ปชานาติ.}

\prose{อิติ อชฺฌตฺตํ วา ธมฺเมสุ ธมฺมานุปสฺสี วิหรติ, พหิทฺธา วา ธมฺเมสุ ธมฺมานุปสฺสี วิหรติ, อชฺฌตฺตพหิทฺธา วา ธมฺเมสุ ธมฺมานุปสฺสี วิหรติ.\csromanspacer{} สมุทยธมฺมานุปสฺสี วา ธมฺเมสุ วิหรติ, วยธมฺมานุปสฺสี วา ธมฺเมสุ วิหรติ, สมุทยวยธมฺมานุปสฺสี วา ธมฺเมสุ วิหรติ. “อตฺถิ ธมฺมา”ติ วา ปนสฺส สติ ปจฺจุปฏฺฐิตา โหติ ยาวเทว ญาณมตฺตาย ปฏิสฺสติมตฺตาย.\csromanspacer{} อนิสฺสิโต จ วิหรติ, น จ กิญฺจิ โลเก อุปาทิยติ.\csromanspacer{} เอวมฺปิ โข ภิกฺขเว ภิกฺขุ ธมฺเมสุ ธมฺมานุปสฺสี วิหรติ ฉสุ อชฺฌตฺติกพาหิเรสุ อายตเนสุ.}

\nitthitamruled{อายตนปพฺพํ นิฏฺฐิตํ.}
\csromanmatikaanchor{mtk.130}
\csromantocmarkref{section}{ธมฺมานุปสฺสนา โพชฺฌงฺคปพฺพ}{mtk.130}
\bookmark[dest={mtk.130},level=1]{ธมฺมานุปสฺสนา โพชฺฌงฺคปพฺพ}
\csromanheader{1}{\textbf{ธมฺมานุปสฺสนา โพชฺฌงฺคปพฺพ}}
\proseitem{385}{ปุน จปรํ ภิกฺขเว ภิกฺขุ ธมฺเมสุ ธมฺมานุปสฺสี วิหรติ สตฺตสุ โพชฺฌงฺเคสุ.\csromanspacer{} กถญฺจ ปน ภิกฺขเว ภิกฺขุ ธมฺเมสุ ธมฺมานุปสฺสี วิหรติ สตฺตสุ โพชฺฌงฺเคสุ.\csromanspacer{} อิธ ภิกฺขเว ภิกฺขุ สนฺตํ วา อชฺฌตฺตํ สติสมฺโพชฺฌงฺคํ “อตฺถิ เม อชฺฌตฺตํ สติสมฺโพชฺฌงฺโค”ติ ปชานาติ, อสนฺตํ วา อชฺฌตฺตํ สติสมฺโพชฺฌงฺคํ “นตฺถิ เม อชฺฌตฺตํ สติสมฺโพชฺฌงฺโค”ติ ปชานาติ, ยถา จ อนุปฺปนฺนสฺส สติสมฺโพชฺฌงฺคสฺส อุปฺปาโท โหติ ตญฺจ ปชานาติ, ยถา จ อุปฺปนฺนสฺส สติสมฺโพชฺฌงฺคสฺส ภาวนาย ปาริปูรี โหติ ตญฺจ ปชานาติ.}

\prose{สนฺตํ วา อชฺฌตฺตํ ธมฺมวิจยสมฺโพชฺฌงฺคํ “อตฺถิ เม อชฺฌตฺตํ ธมฺมวิจยสมฺโพชฺฌงฺโค”ติ ปชานาติ, อสนฺตํ วา อชฺฌตฺตํ ธมฺมวิจยสมฺโพชฺฌงฺคํ “นตฺถิ เม อชฺฌตฺตํ ธมฺมวิจยสมฺโพชฺฌงฺโค”ติ ปชานาติ, ยถา จ อนุปฺปนฺนสฺส ธมฺมวิจยสมฺโพชฺฌงฺคสฺส อุปฺปาโท โหติ ตญฺจ ปชานาติ, ยถา จ \csromanfolio{242}อุปฺปนฺนสฺส ธมฺมวิจยสมฺโพชฺฌงฺคสฺส ภาวนาย ปาริปูรี โหติ ตญฺจ ปชานาติ.}

\prose{สนฺตํ วา อชฺฌตฺตํ วีริยสมฺโพชฺฌงฺคํ “อตฺถิ เม อชฺฌตฺตํ วีริยสมฺโพชฺฌงฺโค”ติ ปชานาติ, อสนฺตํ วา อชฺฌตฺตํ วีริยสมฺโพชฺฌงฺคํ “นตฺถิ เม อชฺฌตฺตํ วีริยสมฺโพชฺฌงฺโค”ติ ปชานาติ, ยถา จ อนุปฺปนฺนสฺส วีริยสมฺโพชฺฌงฺคสฺส อุปฺปาโท โหติ ตญฺจ ปชานาติ, ยถา จ อุปฺปนฺนสฺส วีริยสมฺโพชฺฌงฺคสฺส ภาวนาย ปาริปูรี โหติ ตญฺจ ปชานาติ.}

\prose{สนฺตํ วา อชฺฌตฺตํ ปีติสมฺโพชฺฌงฺคํ “อตฺถิ เม อชฺฌตฺตํ ปีติสมฺโพชฺฌงฺโค”ติ ปชานาติ, อสนฺตํ วา อชฺฌตฺตํ ปีติสมฺโพชฺฌงฺคํ “นตฺถิ เม อชฺฌตฺตํ ปีติสมฺโพชฺฌงฺโค”ติ ปชานาติ, ยถา จ อนุปฺปนฺนสฺส ปีติสมฺโพชฺฌงฺคสฺส อุปฺปาโท โหติ ตญฺจ ปชานาติ, ยถา จ อุปฺปนฺนสฺส ปีติสมฺโพชฺฌงฺคสฺส ภาวนาย ปาริปูรี โหติ ตญฺจ ปชานาติ.}

\prose{สนฺตํ วา อชฺฌตฺตํ ปสฺสทฺธิสมฺโพชฺฌงฺคํ “อตฺถิ เม อชฺฌตฺตํ ปสฺสทฺธิสมฺโพชฺฌงฺโค”ติ ปชานาติ, อสนฺตํ วา อชฺฌตฺตํ ปสฺสทฺธิสมฺโพชฺฌงฺคํ “นตฺถิ เม อชฺฌตฺตํ ปสฺสทฺธิสมฺโพชฺฌงฺโค”ติ ปชานาติ, ยถา จ อนุปฺปนฺนสฺส ปสฺสทฺธิสมฺโพชฺฌงฺคสฺส อุปฺปาโท โหติ ตญฺจ ปชานาติ, ยถา จ อุปฺปนฺนสฺส ปสฺสทฺธิสมฺโพชฺฌงฺคสฺส ภาวนาย ปาริปูรี โหติ ตญฺจ ปชานาติ.}

\prose{สนฺตํ วา อชฺฌตฺตํ สมาธิสมฺโพชฺฌงฺคํ “อตฺถิ เม อชฺฌตฺตํ สมาธิสมฺโพชฺฌงฺโค”ติ ปชานาติ, อสนฺตํ วา อชฺฌตฺตํ สมาธิสมฺโพชฺฌงฺคํ “นตฺถิ เม อชฺฌตฺตํ สมาธิสมฺโพชฺฌงฺโค”ติ ปชานาติ, ยถา จ อนุปฺปนฺนสฺส สมาธิสมฺโพชฺฌงฺคสฺส อุปฺปาโท โหติ ตญฺจ ปชานาติ, ยถา จ อุปฺปนฺนสฺส สมาธิสมฺโพชฺฌงฺคสฺส ภาวนาย ปาริปูรี โหติ ตญฺจ ปชานาติ.}

\prose{สนฺตํ วา อชฺฌตฺตํ อุเปกฺขาสมฺโพชฺฌงฺคํ “อตฺถิ เม อชฺฌตฺตํ อุเปกฺขาสมฺโพชฺฌงฺโค”ติ ปชานาติ, อสนฺตํ วา อชฺฌตฺตํ อุเปกฺขาสมฺโพชฺฌงฺคํ “นตฺถิ เม อชฺฌตฺตํ อุเปกฺขาสมฺโพชฺฌงฺโค”ติ ปชานาติ, ยถา จ อนุปฺปนฺนสฺส อุเปกฺขาสมฺโพชฺฌงฺคสฺส อุปฺปาโท โหติ ตญฺจ ปชานาติ, ยถา จ อุปฺปนฺนสฺส อุเปกฺขาสมฺโพชฺฌงฺคสฺส ภาวนาย ปาริปูรี โหติ ตญฺจ ปชานาติ.}

\prose{อิติ อชฺฌตฺตํ วา ธมฺเมสุ ธมฺมานุปสฺสี วิหรติ, พหิทฺธา วา ธมฺเมสุ ธมฺมานุปสฺสี วิหรติ, อชฺฌตฺตพหิทฺธา วา ธมฺเมสุ ธมฺมานุปสฺสี วิหรติ.}

\vspace{\tipitakaparskip}
\csromancenter{มหาสติปฏฺฐานสุตฺต สมุทยธมฺมานุปสฺสี วา ธมฺเมสุ วิหรติ, วยธมฺมานุปสฺสี วา ธมฺเมสุ วิหรติ, สมุทยวยธมฺมานุปสฺสี วา ธมฺเมสุ วิหรติ. “อตฺถิ ธมฺมา”ติ วา ปนสฺส สติ ปจฺจุปฏฺฐิตา โหติ ยาวเทว ญาณมตฺตาย ปฏิสฺสติมตฺตาย.\csromanspacer{} อนิสฺสิโต จ วิหรติ, น จ กิญฺจิ โลเก อุปาทิยติ.\csromanspacer{} เอวมฺปิ โข ภิกฺขเว ภิกฺขุ ธมฺเมสุ ธมฺมานุปสฺสี วิหรติ สตฺตสุ โพชฺฌงฺเคสุ.}
\nitthitamruled{โพชฺฌงฺคปพฺพํ นิฏฺฐิตํ\footnote{โพชฺฌงฺคปพฺพํ นิฏฺฐิตํ. ปฐมภาณวารํ (สฺยา)}.}
\csromanmatikaanchor{mtk.131}
\csromantocmarkref{section}{ธมฺมานุปสฺสนา สจฺจปพฺพ}{mtk.131}
\bookmark[dest={mtk.131},level=1]{ธมฺมานุปสฺสนา สจฺจปพฺพ}
\csromanheader{1}{\textbf{ธมฺมานุปสฺสนา สจฺจปพฺพ}}
\proseitem{386}{\csromanfolio{243}ปุน จปรํ ภิกฺขเว ภิกฺขุ ธมฺเมสุ ธมฺมานุปสฺสี วิหรติ จตูสุ อริยสจฺเจสุ.\csromanspacer{} กถญฺจ ปน ภิกฺขเว ภิกฺขุ ธมฺเมสุ ธมฺมานุปสฺสี วิหรติ จตูสุ อริยสจฺเจสุ.\csromanspacer{} อิธ ภิกฺขเว ภิกฺขุ “อิทํ ทุกฺขนฺ”ติ ยถาภูตํ ปชานาติ, “อยํ ทุกฺขสมุทโย”ติ ยถาภูตํ ปชานาติ, “อยํ ทุกฺขนิโรโธ”ติ ยถาภูตํ ปชานาติ, “อยํ ทุกฺขนิโรธคามินี ปฏิปทา”ติ ยถาภูตํ ปชานาติ.}

\nitthitamruled{ปฐมภาณวาโร นิฏฺฐิโต.}
\csromanmatikaanchor{mtk.132}
\csromantocmarkref{section}{ทุกฺขสจฺจนิทฺเทส}{mtk.132}
\bookmark[dest={mtk.132},level=1]{ทุกฺขสจฺจนิทฺเทส}
\csromanheader{1}{\textbf{ทุกฺขสจฺจนิทฺเทส}}
\proseitem{387}{กตมญฺจ ภิกฺขเว ทุกฺขํ อริยสจฺจํ.\csromanspacer{} ชาติปิ ทุกฺขา, ชราปิ ทุกฺขา, มรณมฺปิ ทุกฺขํ, โสกปริเทวทุกฺขโทมนสฺสุปายาสาปิ ทุกฺขา, อปฺปิเยหิ สมฺปโยโคปิ ทุกฺโข, ปิเยหิ วิปฺปโยโคปิ ทุกฺโข\footnote{อปฺปิเยหิ ฯเปฯ วิปฺปโยโค ทุกฺโขติปาโฐ เจว ตํนิทฺเทโส จ กตฺถจิ น ทิสฺสติ, อฏฺฐกถายํปิ ตํสํวณฺณนา นตฺถิ.}, ยมฺปิจฺฉํ น ลภติ ตมฺปิ ทุกฺขํ, สํขิตฺเตน ปญฺจุปาทานกฺขนฺธา\footnote{ปญฺจุปาทานกฺขนฺธาปิ (ก)}, ทุกฺขา.}

\proseitem{388}{กตมา จ ภิกฺขเว ชาติ.\csromanspacer{} ยา เตสํ เตสํ สตฺตานํ ตมฺหิ ตมฺหิ สตฺตนิกาเย ชาติ สญฺชาติ โอกฺกนฺติ อภินิพฺพตฺติ ขนฺธานํ ปาตุภาโว อายตนานํ ปฏิลาโภ.\csromanspacer{} อยํ วุจฺจติ ภิกฺขเว ชาติ.}

\proseitem{389}{\csromanfolio{244}กตมา จ ภิกฺขเว ชรา.\csromanspacer{} ยา เตสํ เตสํ สตฺตานํ ตมฺหิ ตมฺหิ สตฺตนิกาเย ชรา ชีรณตา ขณฺฑิจฺจํ ปาลิจฺจํ วลิตฺตจตา อายุโน สํหานิ อินฺทฺริยานํ ปริปาโก.\csromanspacer{} อยํ วุจฺจติ ภิกฺขเว ชรา.}

\proseitem{390}{กตมญฺจ ภิกฺขเว มรณํ.\csromanspacer{} ยํ\footnote{อฏฺฐกถา โอโลเกตพฺพา.} เตสํ เตสํ สตฺตานํ ตมฺหา ตมฺหา สตฺตนิกายา จุติ จวนตา เภโท อนฺตรธานํ มจฺจุ มรณํ กาลกิริยา ขนฺธานํ เภโท กเฬวรสฺส นิกฺเขโป ชีวิตินฺทฺริยสฺสุปจฺเฉโท.\csromanspacer{} อิทํ วุจฺจติ ภิกฺขเว มรณํ.}

\proseitem{391}{กตโม จ ภิกฺขเว โสโก.\csromanspacer{} โย โข ภิกฺขเว อญฺญตรญฺญตเรน พฺยสเนน สมนฺนาคตสฺส อญฺญตรญฺญตเรน ทุกฺขธมฺเมน ผุฏฺฐสฺส โสโก โสจนา โสจิตตฺตํ อนฺโตโสโก อนฺโตปริโสโก.\csromanspacer{} อยํ วุจฺจติ ภิกฺขเว โสโก.}

\proseitem{392}{กตโม จ ภิกฺขเว ปริเทโว.\csromanspacer{} โย โข ภิกฺขเว อญฺญตรญฺญตเรน พฺยสเนน สมนฺนาคตสฺส อญฺญตรญฺญตเรน ทุกฺขธมฺเมน ผุฏฺฐสฺส อาเทโว ปริเทโว อาเทวนา ปริเทวนา อาเทวิตตฺตํ ปริเทวิตตฺตํ.\csromanspacer{} อยํ วุจฺจติ ภิกฺขเว ปริเทโว.}

\proseitem{393}{กตมญฺจ ภิกฺขเว ทุกฺขํ.\csromanspacer{} ยํ โข ภิกฺขเว กายิกํ ทุกฺขํ กายิกํ อสาตํ กายสมฺผสฺสชํ ทุกฺขํ อสาตํ เวทยิตํ.\csromanspacer{} อิทํ วุจฺจติ ภิกฺขเว ทุกฺขํ.}

\proseitem{394}{กตมญฺจ ภิกฺขเว โทมนสฺสํ.\csromanspacer{} ยํ โข ภิกฺขเว เจตสิกํ ทุกฺขํ เจตสิกํ อสาตํ มโนสมฺผสฺสชํ ทุกฺขํ อสาตํ เวทยิตํ.\csromanspacer{} อิทํ วุจฺจติ ภิกฺขเว โทมนสฺสํ.}

\proseitem{395}{กตโม จ ภิกฺขเว อุปายาโส.\csromanspacer{} โย โข ภิกฺขเว อญฺญตรญฺญตเรน พฺยสเนน สมนฺนาคตสฺส อญฺญตรญฺญตเรน ทุกฺขธมฺเมน ผุฏฺฐสฺส อายาโส อุปายาโส อายาสิตตฺตํ อุปายาสิตตฺตํ.\csromanspacer{} อยํ วุจฺจติ ภิกฺขเว อุปายาโส.}

\csromanheader{2}{9. มหาสติปฏฺฐานสุตฺต}
\proseitem{396}{\csromanfolio{245}กตโม จ ภิกฺขเว อปฺปิเยหิ สมฺปโยโค ทุกฺโข.\csromanspacer{} อิธ ยสฺส เต โหนฺติ อนิฏฺฐา อกนฺตา อมนาปา รูปา สทฺทา คนฺธา รสา โผฏฺฐพฺพา ธมฺมา, เย วา ปนสฺส เต โหนฺติ อนตฺถกามา อหิตกามา อผาสุกกามา อโยคกฺเขมกามา, ยา เตหิ สทฺธิํ สงฺคติ สมาคโม สโมธานํ มิสฺสีภาโว.\csromanspacer{} อยํ วุจฺจติ ภิกฺขเว อปฺปิเยหิ สมฺปโยโค ทุกฺโข.}

\proseitem{397}{กตโม จ ภิกฺขเว ปิเยหิ วิปฺปโยโค ทุกฺโข.\csromanspacer{} อิธ ยสฺส เต โหนฺติ อิฏฺฐา กนฺตา มนาปา รูปา สทฺทา คนฺธา รสา โผฏฺฐพฺพา ธมฺมา, เย วา ปนสฺส เต โหนฺติ อตฺถกามา หิตกามา ผาสุกกามา โยคกฺเขมกามา มาตา วา ปิตา วา ภาตา วา ภคินี วา มิตฺตา วา อมจฺจา วา ญาติสาโลหิตา วา, ยา เตหิ สทฺธิํ อสงฺคติ อสมาคโม อสโมธานํ อมิสฺสีภาโว.\csromanspacer{} อยํ วุจฺจติ ภิกฺขเว ปิเยหิ วิปฺปโยโค ทุกฺโข.}

\proseitem{398}{กตมญฺจ ภิกฺขเว ยมฺปิจฺฉํ น ลภติ ตมฺปิ ทุกฺขํ.\csromanspacer{} ชาติธมฺมานํ ภิกฺขเว สตฺตานํ เอวํ อิจฺฉา อุปฺปชฺชติ “อโห วต มยํ น ชาติธมฺมา อสฺสาม, น จ วต โน ชาติ อาคจฺเฉยฺยา”ติ.\csromanspacer{} น โข ปเนตํ อิจฺฉาย ปตฺตพฺพํ, อิทมฺปิ ยมฺปิจฺฉํ น ลภติ ตมฺปิ ทุกฺขํ.\csromanspacer{} ชราธมฺมานํ ภิกฺขเว สตฺตานํ เอวํ อิจฺฉา อุปฺปชฺชติ “อโห วต มยํ น ชราธมฺมา อสฺสาม, น จ วต โน ชรา อาคจฺเฉยฺยา”ติ.\csromanspacer{} น โข ปเนตํ อิจฺฉาย ปตฺตพฺพํ, อิทมฺปิ ยมฺปิจฺฉํ น ลภติ ตมฺปิ ทุกฺขํ.\csromanspacer{} พฺยาธิธมฺมานํ ภิกฺขเว สตฺตานํ เอวํ อิจฺฉา อุปฺปชฺชติ “อโห วต มยํ น พฺยาธิธมฺมา อสฺสาม, น จ วต โน พฺยาธิ อาคจฺเฉยฺยา”ติ.\csromanspacer{} น โข ปเนตํ อิจฺฉาย ปตฺตพฺพํ, อิทมฺปิ ยมฺปิจฺฉํ น ลภติ ตมฺปิ ทุกฺขํ.\csromanspacer{} มรณธมฺมานํ ภิกฺขเว สตฺตานํ เอวํ อิจฺฉา อุปฺปชฺชติ “อโห วต มยํ น มรณธมฺมา อสฺสาม, น จ วต โน มรณํ อาคจฺเฉยฺยา”ติ.\csromanspacer{} น โข ปเนตํ อิจฺฉาย ปตฺตพฺพํ, อิทมฺปิ ยมฺปิจฺฉํ น ลภติ ตมฺปิ ทุกฺขํ.\csromanspacer{} โสกปริเทวทุกฺขโทมนสฺสุปายาสธมฺมานํ ภิกฺขเว สตฺตานํ เอวํ อิจฺฉา อุปฺปชฺชติ “อโห วต มยํ น โสกปริเทวทุกฺขโทมนสฺสุปายาสา อสฺสาม, น จ วต โน โสกปริเทวทุกฺขโทมนสฺสุปายาสธมฺมา อาคจฺเฉยฺยุนฺ”ติ.\csromanspacer{} น โข ปเนตํ อิจฺฉาย ปตฺตพฺพํ, อิทมฺปิ ยมฺปิจฺฉํ น ลภติ ตมฺปิ ทุกฺขํ.}

\proseitem{399}{\csromanfolio{246}กตเม จ ภิกฺขเว สํขิตฺเตน ปญฺจุปาทานกฺขนฺธา ทุกฺขา.\csromanspacer{} เสยฺยถิทํ, รูปุปาทานกฺขนฺโธ, เวทนุปาทานกฺขนฺโธ, สญฺญุปาทานกฺขนฺโธ, สงฺขารุปาทานกฺขนฺโธ, วิญฺญาณุปาทานกฺขนฺโธ.\csromanspacer{} อิเม วุจฺจนฺติ ภิกฺขเว สํขิตฺเตน ปญฺจุปาทานกฺขนฺธา ทุกฺขา.\csromanspacer{} อิทํ วุจฺจติ ภิกฺขเว ทุกฺขํ อริยสจฺจํ.}

\csromanmatikaanchor{mtk.133}
\csromantocmarkref{section}{สมุทยสจฺจนิทฺเทส}{mtk.133}
\bookmark[dest={mtk.133},level=1]{สมุทยสจฺจนิทฺเทส}
\csromanheader{1}{\textbf{สมุทยสจฺจนิทฺเทส}}
\proseitem{400}{กตมญฺจ ภิกฺขเว ทุกฺขสมุทยํ\footnote{ทุกฺขสมุทโย (สฺยา)} อริยสจฺจํ.\csromanspacer{} ยายํ ตณฺหา โปโนพฺภวิกา\footnote{โปโนภวิกา (สี, อิ)} นนฺทีราคสหคตา\footnote{นนฺทิราคสหคตา (สี, สฺยา, อิ)} ตตฺรตตฺราภินนฺทินี.\csromanspacer{} เสยฺยถิทํ, กามตณฺหา ภวตณฺหา วิภวตณฺหา.}

\prose{สา โข ปเนสา ภิกฺขเว ตณฺหา กตฺถ อุปฺปชฺชมานา อุปฺปชฺชติ, กตฺถ นิวิสมานา นิวิสติ.\csromanspacer{} ยํ โลเก ปิยรูปํ สาตรูปํ, เอตฺเถสา ตณฺหา อุปฺปชฺชมานา อุปฺปชฺชติ, เอตฺถ นิวิสมานา นิวิสติ.}

\prose{กิญฺจ โลเก ปิยรูปํ สาตรูปํ.\csromanspacer{} จกฺขุ โลเก ปิยรูปํ สาตรูปํ, เอตฺเถสา ตณฺหา อุปฺปชฺชมานา อุปฺปชฺชติ, เอตฺถ นิวิสมานา นิวิสติ.\csromanspacer{} โสตํ โลเก ฯเปฯ.\csromanspacer{} ฆานํ โลเก.\csromanspacer{} ชิวฺหา โลเก.\csromanspacer{} กาโย โลเก.\csromanspacer{} มโน โลเก ปิยรูปํ สาตรูปํ, เอตฺเถสา ตณฺหา อุปฺปชฺชมานา อุปฺปชฺชติ, เอตฺถ นิวิสมานา นิวิสติ.}

\prose{รูปา โลเก.\csromanspacer{} สทฺทา โลเก.\csromanspacer{} คนฺธา โลเก.\csromanspacer{} รสา โลเก.\csromanspacer{} โผฏฺฐพฺพา โลเก.\csromanspacer{} ธมฺมา โลเก ปิยรูปํ สาตรูปํ, เอตฺเถสา ตณฺหา อุปฺปชฺชมานา อุปฺปชฺชติ, เอตฺถ นิวิสมานา นิวิสติ.}

\prose{จกฺขุวิญฺญาณํ โลเก.\csromanspacer{} โสตวิญฺญาณํ โลเก.\csromanspacer{} ฆานวิญฺญาณํ โลเก.\csromanspacer{} ชิวฺหาวิญฺญาณํ โลเก.\csromanspacer{} กายวิญฺญาณํ โลเก.\csromanspacer{} มโนวิญฺญาณํ โลเก ปิยรูปํ สาตรูปํ, เอตฺเถสา ตณฺหา อุปฺปชฺชมานา อุปฺปชฺชติ, เอตฺถ นิวิสมานา นิวิสติ.}

\prose{จกฺขุสมฺผสฺโส โลเก.\csromanspacer{} โสตสมฺผสฺโส โลเก.\csromanspacer{} ฆานสมฺผสฺโส โลเก.\csromanspacer{} ชิวฺหาสมฺผสฺโส โลเก.\csromanspacer{} กายสมฺผสฺโส โลเก.\csromanspacer{} มโนสมฺผสฺโส โลเก ปิยรูปํ สาตรูปํ, เอตฺเถสา ตณฺหา อุปฺปชฺชมานา อุปฺปชฺชติ, เอตฺถ นิวิสมานา นิวิสติ.}

\csromanheader{2}{9. มหาสติปฏฺฐานสุตฺต}
\prose{\csromanfolio{247}จกฺขุสมฺผสฺสชา เวทนา โลเก.\csromanspacer{} โสตสมฺผสฺสชา เอทนาโลเก.ฆานสมฺผสฺสชา เวทนา โลเก.\csromanspacer{} ชิวฺหาสมฺผสฺสชา เวทนา โลเก.\csromanspacer{} กายสมฺผสฺสชา เวทนา โลเก.\csromanspacer{} มโนสมฺผสฺสชา เวทนา โลเก ปิยรูปํ สาตรูปํ, เอตฺเถสา ตณฺหา อุปฺปชฺชมานา อุปฺปชฺชติ, เอตฺถ นิวิสมานา นิวิสติ.}

\prose{รูปสญฺญา โลเก.\csromanspacer{} สทฺทสญฺญา โลเก.\csromanspacer{} คนฺธสญฺญา โลเก.\csromanspacer{} รสสญฺญา โลเก.\csromanspacer{} โผฏฺฐพฺพสญฺญา โลเก.\csromanspacer{} ธมฺมสญฺญา โลเก ปิยรูปํ สาตรูปํ, เอตฺเถสา ตณฺหา อุปฺปชฺชมานา อุปฺปชฺชติ, เอตฺถ นิวิสมานา นิวิสติ.}

\prose{รูปสญฺเจตนา โลเก.\csromanspacer{} สทฺทสญฺเจตนา โลเก.\csromanspacer{} คนฺธสญฺเจตนา โลเก.\csromanspacer{} รสสญฺเจตนา โลเก.\csromanspacer{} โผฏฺฐพฺพสญฺเจตนา โลเก.\csromanspacer{} ธมฺมสญฺเจตนา โลเก ปิยรูปํ สาตรูปํ, เอตฺเถสา ตณฺหา อุปฺปชฺชมานา อุปฺปชฺชติ, เอตฺถ นิวิสมานา นิวิสติ.}

\prose{รูปตณฺหา โลเก.\csromanspacer{} สทฺทตณฺหา โลเก.\csromanspacer{} คนฺธตณฺหา โลเก.\csromanspacer{} รสตณฺหา โลเก.\csromanspacer{} โผฏฺฐพฺพตณฺหา โลเก.\csromanspacer{} ธมฺมตณฺหา โลเก ปิยรูปํ สาตรูปํ, เอตฺเถสา ตณฺหา อุปฺปชฺชมานา อุปฺปชฺชติ, เอตฺถ นิวิสมานา นิวิสติ.}

\prose{รูปวิตกฺโก โลเก.\csromanspacer{} สทฺทวิตกฺโก โลเก.\csromanspacer{} คนฺธวิตกฺโก โลเก.\csromanspacer{} รสวิตกฺโก โลเก.\csromanspacer{} โผฏฺฐพฺพวิตกฺโก โลเก.\csromanspacer{} ธมฺมวิตกฺโก โลเก ปิยรูปํ สาตรูปํ, เอตฺเถสา ตณฺหา อุปฺปชฺชมานา อุปฺปชฺชติ, เอตฺถ นิวิสมานา นิวิสติ.}

\prose{รูปวิจาโร โลเก.\csromanspacer{} สทฺทวิจาโร โลเก.\csromanspacer{} คนฺธวิจาโร โลเก.\csromanspacer{} รสวิจาโร โลเก.\csromanspacer{} โผฏฺฐพฺพวิจาโร โลเก.\csromanspacer{} ธมฺมวิจาโร โลเก ปิยรูปํ สาตรูปํ, เอตฺเถสา ตณฺหา อุปฺปชฺชมานา อุปฺปชฺชติ, เอตฺถ นิวิสมานา นิวิสติ.\csromanspacer{} อิทํ วุจฺจติ ภิกฺขเว ทุกฺขสมุทยํ อริยสจฺจํ.}

\csromanmatikaanchor{mtk.134}
\csromantocmarkref{section}{นิโรธสจฺจนิทฺเทส}{mtk.134}
\bookmark[dest={mtk.134},level=1]{นิโรธสจฺจนิทฺเทส}
\csromanheader{1}{\textbf{นิโรธสจฺจนิทฺเทส}}
\proseitem{401}{กตมญฺจ ภิกฺขเว ทุกฺขนิโรธํ\footnote{ทุกฺขนิโรโธ (สฺยา)} อริยสจฺจํ.\csromanspacer{} โย ตสฺสาเยว ตณฺหาย อเสสวิราคนิโรโธ จาโค ปฏินิสฺสคฺโค มุตฺติ อนาลโย.}

\prose{\csromanfolio{248}สา โข ปเนสา ภิกฺขเว ตณฺหา กตฺถ ปหียมานา ปหียติ, กตฺถ นิรุชฺฌมานา นิรุชฺฌติ.\csromanspacer{} ยํ โลเก ปิยรูปํ สาตรูปํ, เอตฺเถสา ตณฺหา ปหียมานา ปหียติ, เอตฺถ นิรุชฺฌมานา นิรุชฺฌติ.}

\prose{กิญฺจ โลเก ปิยรูปํ สาตรูปํ.\csromanspacer{} จกฺขุ โลเก ปิยรูปํ สาตรูปํ, เอตฺเถสา ตณฺหา ปหียมานา ปหียติ, เอตฺถ นิรุชฺฌมานา นิรุชฺฌติ.\csromanspacer{} โสตํ โลเก ฯเปฯ.\csromanspacer{} ฆานํ โลเก.\csromanspacer{} ชิวฺหา โลเก.\csromanspacer{} กาโย โลเก.\csromanspacer{} มโน โลเก ปิยรูปํ สาตรูปํ, เอตฺเถสา ตณฺหา ปหียมานา ปหียติ, เอตฺถ นิรุชฺฌมานา นิรุชฺฌติ.}

\prose{รูปา โลเก.\csromanspacer{} สทฺทา โลเก.\csromanspacer{} คนฺธา โลเก.\csromanspacer{} รสา โลเก.\csromanspacer{} โผฏฺฐพฺพา โลเก.\csromanspacer{} ธมฺมา โลเก ปิยรูปํ สาตรูปํ, เอตฺเถสา ตณฺหา ปหิยมานา ปหียติ, เอตฺถ นิรุชฺฌมานา นิรุชฺฌติ.}

\prose{จกฺขุวิญฺญาณํ โลเก.\csromanspacer{} โสตวิญฺญาณํ โลเก.\csromanspacer{} ฆานวิญฺญาณํ โลเก.\csromanspacer{} ชิวฺหาวิญฺญาณํ โลเก.\csromanspacer{} กายวิญฺญาณํ โลเก.\csromanspacer{} มโนวิญฺญาณํ โลเก ปิยรูปํ สาตรูปํ, เอตฺเถสา ตณฺหา ปหียมานา ปหียติ, เอตฺถ นิรุชฺฌมานา นิรุชฺฌติ.}

\prose{จกฺขุสมฺผสฺโส โลเก.\csromanspacer{} โสตสมฺผสฺโส โลเก.\csromanspacer{} ฆานสมฺผสฺโส โลเก.\csromanspacer{} ชิวฺหาสมฺผสฺโส โลเก.\csromanspacer{} กายสมฺผสฺโส โลเก.\csromanspacer{} มโนสมฺผสฺโส โลเก ปิยรูปํ สาตรูปํ, เอตฺเถสา ตณฺหา ปหียมานา ปหียติ, เอตฺถ นิรุชฺฌมานา นิรุชฺฌติ.}

\prose{จกฺขุสมฺผสฺสชา เวทนา โลเก.\csromanspacer{} โสตสมฺผสฺสชา เวทนา โลเก.\csromanspacer{} ฆานสมฺผสฺสชา เวทนา โลเก.\csromanspacer{} ชิวฺหาสมฺผสฺสชา เวทนา โลเก.\csromanspacer{} กายสมฺผสฺสชา เวทนา โลเก.\csromanspacer{} มโนสมฺผสฺสชา เวทนา โลเก ปิยรูปํ สาตรูปํ, เอตฺเถสา ตณฺหา ปหียมานา ปหียติ, เอตฺถ นิรุชฺฌมานา นิรุชฺฌติ.}

\prose{รูปสญฺญา โลเก.\csromanspacer{} สทฺทสญฺญา โลเก.\csromanspacer{} คนฺธสญฺญา โลเก.\csromanspacer{} รสสญฺญา โลเก.\csromanspacer{} โผฏฺฐพฺพสญฺญา โลเก.\csromanspacer{} ธมฺมสญฺญา โลเก ปิยรูปํ สาตรูปํ, เอตฺเถสา ตณฺหา ปหียมานา ปหียติ, เอตฺถ นิรุชฺฌมานา นิรุชฺฌติ.}

\csromanheader{2}{9. มหาสติปฏฺฐานสุตฺต}
\prose{\csromanfolio{249}รูปสญฺเจตนา โลเก.\csromanspacer{} สทฺทสญฺเจตนา โลเก.\csromanspacer{} คนฺธสญฺเจตนา โลเก.\csromanspacer{} รสสญฺเจตนา โลเก.\csromanspacer{} โผฏฺฐพฺพสญฺเจตนา โลเก.\csromanspacer{} ธมฺมสญฺเจตนา โลเก ปิยรูปํ สาตรูปํ, เอตฺเถสา ตณฺหา ปหียมานา ปหียติ, เอตฺถ นิรุชฺฌมานา นิรุชฺฌติ.}

\prose{รูปตณฺหา โลเก.\csromanspacer{} สทฺทตณฺหา โลเก.คนฺธตณฺหา โลเก.\csromanspacer{} รสตณฺหา โลเก.\csromanspacer{} โผฏฺฐพฺพตณฺหา โลเก.\csromanspacer{} ธมฺมตณฺหา โลเก ปิยรูปํ สาตรูปํ, เอตฺเถสา ตณฺหา ปหียมานา ปหียติ, เอตฺถ นิรุชฺฌมานา นิรุชฺฌติ.}

\prose{รูปวิตกฺโก โลเก.\csromanspacer{} สทฺทวิตกฺโก โลเก.\csromanspacer{} คนฺธวิตกฺโก โลเก.\csromanspacer{} รสวิตกฺโก โลเก.\csromanspacer{} โผฏฺฐพฺพวิตกฺโก โลเก.\csromanspacer{} ธมฺมวิตกฺโก โลเก ปิยรูปํ สาตรูปํ, เอตฺเถสา ตณฺหา ปหียมานา ปหียติ, เอตฺถ นิรุชฺฌมานา นิรุชฺฌติ.}

\prose{รูปวิจาโร โลเก.\csromanspacer{} สทฺทวิจาโร โลเก.\csromanspacer{} คนฺธวิจาโร โลเก.\csromanspacer{} รสวิจาโร โลเก.\csromanspacer{} โผฏฺฐพฺพวิจาโร โลเก.\csromanspacer{} ธมฺมวิจาโร โลเก ปิยรูปํ สาตรูปํ, เอตฺเถสา ตณฺหา ปหียมานา ปหียติ, เอตฺถ นิรุชฺฌมานา นิรุชฺฌติ.\csromanspacer{} อิทํ วุจฺจติ ภิกฺขเว ทุกฺขนิโรธํ อริยสจฺจํ.}

\csromanmatikaanchor{mtk.135}
\csromantocmarkref{section}{มคฺคสจฺจนิทฺเทส}{mtk.135}
\bookmark[dest={mtk.135},level=1]{มคฺคสจฺจนิทฺเทส}
\csromanheader{1}{\textbf{มคฺคสจฺจนิทฺเทส}}
\proseitem{402}{กตมญฺจ ภิกฺขเว ทุกฺขนิโรธคามินี ปฏิปทา อริยสจฺจํ.\csromanspacer{} อยเมว อริโย อฏฺฐงฺคิโก มคฺโค.\csromanspacer{} เสยฺยถิทํ, สมฺมาทิฏฺฐิ สมฺมาสงฺกปฺโป สมฺมาวาจา สเมกมฺมนฺโต สมฺมา-อาชีโว สมฺมาวายาโม สมฺมาสติ สมฺมาสมาธิ.}

\prose{กตมา จ ภิกฺขเว สมฺมาทิฏฺฐิ.\csromanspacer{} ยํ โข ภิกฺขเว ทุกฺเข ญาณํ ทุกฺขสมุทเย ญาณํ ทุกฺขนิโรเธ ญาณํ ทุกฺขนิโรธคามินิยา ปฏิปทาย ญาณํ.\csromanspacer{} อยํ วุจฺจติ ภิกฺขเว สมฺมาทิฏฺฐิ.}

\prose{กตโม จ ภิกฺขเว สมฺมาสงฺกปฺโป.\csromanspacer{} เนกฺขมฺมสงฺกปฺโป อพฺยาปาทสงฺกปฺโป อวิหิํสาสงฺกปฺโป.\csromanspacer{} อยํ วุจฺจติ ภิกฺขเว สมฺมาสงฺกปฺโป.}

\prose{\csromanfolio{250}กตมา จ ภิกฺขเว สมฺมาวาจา.\csromanspacer{} มุสาวาทา เวรมณี\footnote{เวรมณิ (ก)} ปิสุณาย วาจาย เวรมณี ผรุสาย วาจาย เวรมณี สมฺผปฺปลาปา เวรมณี.\csromanspacer{} อยํ วุจฺจติ ภิกฺขเว สเมวาจา.}

\prose{กตโม จ ภิกฺขเว สมฺมากมฺมนฺโต.\csromanspacer{} ปาณาติปาตา เวรมณี อทินฺนาทานา เวรมณี กาเมสุมิจฺฉาจารา เวรมณี.\csromanspacer{} อยํ วุจฺจติ ภิกฺขเว สมฺมากมฺมนฺโต.}

\prose{กตโม จ ภิกฺขเว สมฺมา-อาชีโว.\csromanspacer{} อิธ ภิกฺขเว อริยสาวโก มิจฺฉา-อาชีวํ ปหาย สมฺมา-อาชีเวน ชีวิตํ กปฺเปติ.\csromanspacer{} อยํ วุจฺจติ ภิกฺขเว สมฺมา-อาชีโว.}

\prose{กตโม จ ภิกฺขเว สมฺมาวายาโม.\csromanspacer{} อิธ ภิกฺขเว ภิกฺขุ อนุปฺปนฺนานํ ปาปกานํ อกุสลานํ ธมฺมานํ อนุปฺปาทาย ฉนฺทํ ชเนติ วายมติ วีริยํ อารภติ จิตฺตํ ปคฺคณฺหาติ ปทหติ, อุปฺปนฺนานํ ปาปกานํ อกุสลานํ ธมฺมานํ ปหานาย ฉนฺทํ ชเนติ วายมติ วีริยํ อารภติ จิตฺตํ ปคฺคณฺหาติ ปทหติ, อนุปฺปนฺนานํ กุสลานํ ธมฺมานํ อุปฺปาทาย ฉนฺทํ ชเนติ วายมติ วีริยํ อารภติ จิตฺตํ ปคฺคณฺหาติ ปทหติ, อุปฺปนฺนานํ กุสลานํ ธมฺมานํ ฐิติยา อสมฺโมสาย ภิยฺโยภาวาย เวปุลฺลาย ภาวนาย ปาริปูริยา ฉนฺทํ ชเนติ วายมติ วีริยํ อารภติ จิตฺตํ ปคฺคณฺหาติ ปทหติ.\csromanspacer{} อยํ วุจฺจติ ภิกฺขเว สมฺมาวายาโม.}

\prose{กตมา จ ภิกฺขเว สมฺมาสติ.\csromanspacer{} อิธ ภิกฺขเว ภิกฺขุ กาเย กายานุปสฺสี วิหรติ อาตาปี สมฺปชาโน สติมา วิเนยฺย โลเก อภิชฺฌาโทมนสฺสํ, เวทนาสุ เวทนานุปสฺสี วิหรติ อาตาปี สมฺปชาโน สติมา วิเนยฺย โลเก อภิชฺฌาโทมนสฺสํ, จิตฺเต จิตฺตานุปสฺสี วิหรติ อาตาปี สมฺปชาโน สติมา วิเนยฺย โลเก อภิชฺฌาโทมนสฺสํ, ธมฺเมสุ ธมฺมานุปสฺสี วิหรติ อาตาปี สมฺปชาโน สติมา วิเนยฺย โลเก อภิชฺฌาโทมนสฺสํ.\csromanspacer{} อยํ วุจฺจติ ภิกฺขเว สมฺมาสติ.}

\prose{กตโม จ ภิกฺขเว สมฺมาสมาธิ.\csromanspacer{} อิธ ภิกฺขเว ภิกฺขุ วิวิจฺเจว กาเมหิ วิวิจฺจ อกุสเลหิ ธมฺเมหิ สวิตกฺกํ สวิจารํ วิเวกชํ ปีติสุขํ ปฐมํ ฌานํ อุปสมฺปชฺช วิหรติ.\csromanspacer{} วิตกฺกวิจารานํ วูปสมา}

\vspace{\tipitakaparskip}
\csromancenter{มหาสติปฏฺฐานสุตฺต อชฺฌตฺตํ สมฺปสาทนํ เจตโส เอโกทิภาวํ อวิตกฺกํ อวิจารํ สมาธิชํ ปีติสุขํ ทุติยํ ฌานํ อุปสมฺปชฺช วิหรติ.\csromanspacer{} ปีติยา จ วิราคา อุเปกฺขโก จ วิหรติ, สโต จ สมฺปชาโน, สุขญฺจ กาเยน ปฏิสํเวเทติ, ยํ ตํ อริยา อาจิกฺขนฺติ “อุเปกฺขโก สติมา สุขวิหารี”ติ, ตติยํ ฌานํ อุปสมฺปชฺช วิหรติ.\csromanspacer{} สุขสฺส จ ปหานา ทุกฺขสฺส จ ปหานา ปุพฺเพว โสมนสฺสโทมนสฺสานํ อตฺถงฺคมา อทุกฺขมสุขํ อุเปกฺขาสติปาริสุทฺธิํ จตุตฺตํ ฌานํ อุปสมฺปชฺช วิหรติ.\csromanspacer{} อยํ วุจฺจติ ภิกฺขเว สมฺมาสมาธิ.\csromanspacer{} อิทํ วุจฺจติ ภิกฺขเว ทุกฺขนิโรธคามินี ปฏิปทา อริยสจฺจํ.}
\proseitem{403}{\csromanfolio{251}อิติ อชฺฌตฺตํ วา ธมฺเมสุ ธมฺมานุปสฺสี วิหรติ, พหิทฺธา วา ธมฺเมสุ ธมฺมานุปสฺสี วิหรติ, อชฺฌตฺตพหิทฺธา วา ธมฺเมสุ ธมฺมานุปสฺสี วิหรติ.\csromanspacer{} สมุทยธมฺมานุปสฺสี วา ธมฺเมสุ วิหรติ, วยธมฺมานุปสฺสี วา ธมฺเมสุ วิหรติ, สมุทยวยธมฺมานุปสฺสี วา ธมฺเมสุ วิหรติ. “อตฺถิ ธมฺมา”ติ วา ปนสฺส สติ ปจฺจุปฏฺฐิตา โหติ ยาวเทว ญาณมตฺตาย ปฏิสฺสติมตฺตาย.\csromanspacer{} อนิสฺสิโต จ วิหรติ, น จ กิญฺจิ โลเก อุปาทิยติ.\csromanspacer{} เอวมฺปิ โข ภิกฺขเว ภิกฺขุ ธมฺเมสุ ธมฺมานุปสฺสี วิหรติ จตูสุ อริยสจฺเจสุ.}

\nitthitam{สจฺจปพฺพํ นิฏฺฐิตํ.}
\nitthitam{ธมฺมานุปสฺสนา นิฏฺฐิตา.}
\proseitem{404}{โย หิ โกจิ ภิกฺขเว อิเม จตฺตาโย สติปฏฺฐาเน เอวํ ภาเวยฺย สตฺตวสฺสานิ, ตสฺส ทฺวินฺนํ ผลานํ อญฺญตรํ ผลํ ปาฏิกงฺขํ ทิฏฺเฐว ธมฺเม อญฺญา สติ วา อุปาทิเสเส อนาคามิตา.}

\prose{ติฏฺฐนฺตุ ภิกฺขเว สตฺตวสฺสานิ.\csromanspacer{} โย หิ โกจิ ภิกฺขเว อิเม จตฺตาโย สติปฏฺฐาเน เอวํ ภาเวยฺย ฉ วสฺสานิ ฯเปฯ ปญฺจ วสฺสานิ.\csromanspacer{} จตฺตาริ วสฺสานิ.\csromanspacer{} ตีณิ วสฺสานิ.\csromanspacer{} เทฺว วสฺสานิ.\csromanspacer{} เอกํ วสฺสํ.\csromanspacer{} ติฏฺฐตุ ภิกฺขเว เอกํ วสฺสํ.\csromanspacer{} โย หิ โกจิ ภิกฺขเว อิเม จตฺตาโร สติปฏฺฐาเน เอวํ ภาเวยฺย สตฺตมาสานิ, ตสฺส ทฺวินฺนํ ผลานํ อญฺญตรํ ผลํ ปาฏิกงฺขํ ทิฏฺเฐว ธมฺเม อญฺญา สติ วา อุปาทิเสเส อนาคามิตา.}

\prose{\csromanfolio{252}ติฏฺฐนฺตุ ภิกฺขเว สตฺต มาสานิ.\csromanspacer{} โย หิ โกจิ ภิกฺขเว อิเม จตฺตาโร สติปฏฺฐาเน เอวํ ภาเวยฺย ฉ มาสานิ ฯเปฯ ปญฺจ มาสานิ.\csromanspacer{} จตฺตาริ มาสานิ.\csromanspacer{} ตีณิ มาสานิ.\csromanspacer{} เทฺว มาสานิ.\csromanspacer{} เอกํ มาสํ.\csromanspacer{} อฑฺฒมาสํ.\csromanspacer{} ติฏฺฐตุ ภิกฺขเว อฑฺฒมาโส.\csromanspacer{} โย หิ โกจิ ภิกฺขเว อิเม จตฺตาโร สติปฏฺฐาเน เอวํ ภาเวยฺย สตฺตาหํ, ตสฺส ทฺวินฺนํ ผลานํ อญฺญตรํ ผลํ ปาฏิกงฺขํ ทิฏฺเฐว ธมฺเม อญฺญา สติ วา อุปาทิเสเส อนาคามิตาติ.}

\proseitem{405}{“เอกายโน อยํ ภิกฺขเว มคฺโค สตฺตานํ วิสุทฺธิยา โสกปริเทวานํ สมติกฺกมาย ทุกฺขโทมนสฺสานํ อตฺถงฺคมาย ญายสฺส อธิคมาย นิพฺพานสฺส สจฺฉิกิริยาย, ยทิทํ จตฺตาโร สติปฏฺฐานา”ติ อิติ ยํ ตํ วุตฺตํ, อิทเมตํ ปฏิจฺจ วุตฺตนฺติ.\csromanspacer{} อิทมโวจ ภควา.\csromanspacer{} อตฺตมนา เต ภิกฺขู ภควโต ภาสิตํ อภินนฺทุนฺติ.}

\nitthitam{\textbf{มหาสติปฏฺฐานสุตฺตํ นิฏฺฐิตํ นวมํ.}}
\csromanmatikaanchor{mtk.136}
\csromantocmarknopage{chapter}{10. ปายาสิสุตฺต}{mtk.136}
\bookmark[dest={mtk.136},level=0]{10. ปายาสิสุตฺต}
\chapterheadpageread{10. ปายาสิสุตฺต}
\proseitem{406}{\csromanfolio{253}เอวํ เม สุตํ\csromandash{}เอกํ สมยํ อายสฺมา กุมารกสฺสโป โกสเลสุ จาริกํ จรมาโน มหตา ภิกฺขุสํเฆน สทฺธิํ ปญฺจมตฺเตหิ ภิกฺขุสเตหิ เยน เสตพฺยา นาม โกสลานํ นครํ ตทวสริ, ตตฺร สุทํ อายสฺมา กุมารกสฺสโป เสตพฺยายํ วิหรติ อุตฺตเรน เสตพฺยํ สิํสปาวเน\footnote{สีสปาวเน (สฺยา)}.\csromanspacer{} เตน โข ปน สมเยน ปายาสิ ราชญฺโญ เสตพฺยํ อชฺฌาวสติ สตฺตุสฺสทํ สติณกฏฺโฐทกํ สธญฺญํ ราชโภคฺคํ รญฺญา ปเสนทินา โกสเลน ทินฺนํ ราชทายํ พฺรหฺมเทยฺยํ.}

\csromanmatikaanchor{mtk.137}
\csromantocmarkref{section}{ปายาสิราชญฺญวตฺถุ}{mtk.137}
\bookmark[dest={mtk.137},level=1]{ปายาสิราชญฺญวตฺถุ}
\csromanheader{1}{\textbf{ปายาสิราชญฺญวตฺถุ}}
\proseitem{407}{เตน โข ปน สมเยน ปายาสิสฺส ราชญฺญสฺส เอวรูปํ ปาปกํ ทิฏฺฐิคตํ อุปฺปนฺนํ โหติ “อิติปิ นตฺถิ ปโร โลโก, นตฺถิ สตฺตา โอปปาติกา, นตฺถิ สุกตทุกฺกฏานํ\footnote{สุกฏทุกฺกฏานํ (สี, อิ)} กมฺมานํ ผลํ วิปาโก”ติ.\csromanspacer{} อสฺโสสุํ โข เสตพฺยกา พฺราหฺมณคหปติกา “สมโณ ขลุ โภ กุมารกสฺสโป สมณสฺส โคตมสฺส สาวโก โกสเลสุ จาริกํ จรมาโน มหตา ภิกฺขุสํเฆน สทฺธิํ ปญฺจมตฺเตหิ ภิกฺขุสเตหิ เสตพฺยํ อนุปฺปตฺโต เสตพฺยายํ วิหรติ อุตฺตเรน เสตพฺยํ สิํสปาวเน.\csromanspacer{} ตํ โข ปน ภวนฺตํ กุมารกสฺสปํ เอวํ กลฺยาโณ กิตฺติสทฺโท อพฺภุคฺคโต ‘ปณฺฑิโต พฺยตฺโต เมธาวี พหุสฺสุโต จิตฺตกถี กลฺยาณปฏิภาโน วุทฺโธ\footnote{พุทฺโธ (สฺยา, ก)} เจว อรหา จ’.\csromanspacer{} สาธุ โข ปน ตถารูปานํ อรหตํ ทสฺสนํ โหตี”ติ.\csromanspacer{} อถ โข เสตพฺยกา พฺราหฺมณคหปติกา เสตพฺยาย นิกฺขมิตฺวา สงฺฆสงฺฆี คณีภูตา อุตฺตเรนมุขา คจฺฉนฺติ เยน สิํสปาวเน\footnote{เยน สิํสปาวเน, เตนุปสงฺกมนฺติ (สี, อิ)}.}

\proseitem{408}{เตน โข ปน สมเยน ปายาสิ ราชญฺโญ อุปริปาสาเท ทิวาเสยฺยํ อุปคโต โหติ.\csromanspacer{} อทฺทสา โข ปายาสิ ราชญฺโญ เสตพฺยเก พฺราหฺมณคหปติเก เสตพฺยาย นิกฺขมิตฺวา สงฺฆสงฺฆี คณีภูเต อุตฺตเรนมุเข คจฺฉนฺเต เยน สิํสปาวเน\footnote{เยน สิํสปาวเน, เตนุปสงฺกมนฺเต (สี, อิ)}, ทิสฺวา ขตฺตํ \csromanfolio{254}อามนฺเตสิ “กิํ นุ โข โภ ขตฺเต เสตพฺยกา พฺราหฺมณคหปติกา เสตพฺยาย นิกฺขมิตฺวา สงฺฆสงฺฆี คณีภูตา อุตฺตเรนมุขา คจฺฉนฺติ เยน สิํสปาวนนฺ”ติ\footnote{เอตฺถ ปน สพฺพตฺถปิ เอวเมว ทิสฺสติ, นตฺถิ ปาฐนฺตรํ.}.}

\prose{อตฺถิ โข โภ สมโณ กุมารกสฺสโป สมณสฺส โคตมสฺส สาวโก โกสเลสุ จาริกํ จรมาโน มหตา ภิกฺขุสํเฆน สทฺธิํ ปญฺจมตฺเตหิ ภิกฺขุสเตหิ เสตพฺยํ อนุปฺปตฺโต เสตพฺยายํ วิหรติ อุตฺตเรน เสตพฺยํ สิํสปาวเน.\csromanspacer{} ตํ โข ปน ภวนฺตํ กุมารกสฺสปํ เอวํ กลฺยาโณ กิตฺติสทฺโท อพฺภุคฺคโต “ปณฺฑิโต พฺยตฺโต เมธาวี พหุสฺสุโต จิตฺตกถี กลฺยาณปฏิภาโน วุทฺโธ เจว อรหา จา”ติ\footnote{อรหา จ (สฺยา, ก)}.\csromanspacer{} ตเมเต\footnote{ตเมนํ เต (สี, ก), ตเมนํ (อิ)} ภวนฺตํ กุมารกสฺสปํ ทสฺสนาย อุปสงฺกมนฺตีติ.\csromanspacer{} เตน หิ โภ ขตฺเต เยน เสตพฺยกา พฺราหฺมณคหปติกา เตนุปสงฺกม, อุปสงฺกมิตฺวา เสตพฺยเก พฺราหฺมณคหปติเก เอวํ วเทหิ “ปายาสิ โภ ราชญฺโญ เอวมาห ‘อาคเมนฺตุ กิร ภวนฺโต ปายาสิปิ ราชญฺโญ สมณํ กุมารกสฺสปํ ทสฺสนาย อุปสงฺกมิสฺสติ, ปุรา สมโณ กุมารกสฺสโป เสตพฺยเก พฺราหฺมณคหปติเก พาเล อพฺยตฺเต สญฺญาเปติ ‘อิติปิ อตฺถิ ปโร โลโก, อตฺถิ สตฺตา โอปปาติกา, อตฺถิ สุกตทุกฺกฏานํ กมฺมานํ ผลํ วิปาโก’ติ.\csromanspacer{} นตฺถิ หิ โภ ขตฺเต ปโร โลโก, นตฺถิ สตฺตา โอปปาติกา, นตฺถิ สุกตทุกฺกฏานํ กมฺมานํ ผลํ วิปาโก”ติ.\csromanspacer{} เอวํ โภติ โข โส ขตฺตา ปายาสิสฺส ราชญฺญสฺส ปฏิสฺสุตฺวา เยน เสตพฺยกา พฺราหฺมณคหปติกา เตนุปสงฺกมิ, อุปสงฺกมิตฺวา เสตพฺยเก พฺราหฺมณคหปติเก เอตทโวจ “ปายาสิ โภ ราชญฺโญ เอวมาห ‘อาคเมนฺตุ กิร ภวนฺโต ปายาสิปิ ราชญฺโญ สมณํ กุมารกสฺสปํ ทสฺสนาย อุปสงฺกมิสฺสตี”ติ.}

\proseitem{409}{อถ โข ปายาสิ ราชญฺโญ เสตพฺยเกหิ พฺราหฺมณคหปติเกหิ ปริวุโต เยน สิํสปาวเน เยนายสฺมา กุมารกสฺสโป เตนุปสงฺกมิ, อุปสงฺกมิตฺวา อายสฺมตา กุมารกสฺสเปน}

\vspace{\tipitakaparskip}
\csromancenter{ปายาสิสุตฺต สทฺธิํ สมฺโมทิ, สมฺโมทนียํ กถํ สารณียํ วีติสาเรตฺวา เอกมนฺตํ นิสีทิ.\csromanspacer{} เสตพฺยกาปิ โข พฺราหฺมณคหปติกา อปฺเปกจฺเจ อายสฺมนฺตํ กุมารกสฺสปํ อภิวาเทตฺวา เอกมนฺตํ นิสีทิํสุ, อปฺเปกจฺเจ อายสฺมตา กุมารกสฺสเปน สทฺธิํ สมฺโมทิํสุ, สมฺโมทนียํ กถํ สารณียํ วีติสาเรตฺวา เอกมนฺตํ นิสีทิํสุ, อปฺเปกจฺเจ เยนายสฺมา กุมารกสฺสโป เตนญฺชลิํ ปณาเมตฺวา เอกมนฺตํ นิสีทิํสุ, อปฺเปกจฺเจ นามโคตฺตํ สาเวตฺวา เอกมนฺตํ นิสีทิํสุ, อปฺเปกจฺเจ ตุณฺหีภูตา เอกมนฺตํ นิสีทิํสุ.}
\csromanmatikaanchor{mtk.138}
\csromantocmarkref{section}{นตฺถิกวาท}{mtk.138}
\bookmark[dest={mtk.138},level=1]{นตฺถิกวาท}
\csromanheader{1}{\textbf{นตฺถิกวาท}}
\csromanmatikaanchor{mtk.139}
\csromantocmarkref{section}{จนฺทิมสูริย-อุปมา}{mtk.139}
\bookmark[dest={mtk.139},level=1]{จนฺทิมสูริย-อุปมา}
\proseitem{410}{\csromanfolio{255}เอกมนฺตํ นิสินฺโน โข ปายาสิ ราชญฺโญ อายสฺมนฺตํ กุมารกสฺสปํ เอตทโวจ “อหญฺหิ โภ กสฺสป เอวํวาที เอวํทิฏฺฐี ‘อิติปิ นตฺถิ ปโร โลโก, นตฺถิ สตฺตา โอปปาติกา, นตฺถิ สุกตทุกฺกฏานํ กมฺมานํ ผลํ วิปาโก’ติ”.\csromanspacer{} นาหํ ราชญฺญ เอวํวาทิํ เอวํทิฏฺฐิํ อทฺทสํ วา อสฺโสสิํ วา.\csromanspacer{} กถญฺหิ นาม เอวํ วเทยฺย “อิติปิ นตฺถิ ปโร โลโก, นตฺถิ สตฺตา โอปปาติกา, นตฺถิ สุกตทุกฺกฏานํ กมฺมานํ ผลํ วิปาโก”ติ.}

\vspace{\tipitakaparskip}
\csromancenter{\textbf{(1) จนฺทิมสูริย-อุปมา}}
\proseitem{411}{เตน หิ ราชญฺญ ตญฺเญเวตฺถ ปฏิปุจฺฉิสฺสามิ, ยถา เต ขเมยฺย, ตถา นํ พฺยากเรยฺยาสิ.\csromanspacer{} ตํ กิํ มญฺญสิ ราชญฺญ, อิเม จนฺทิมสูริยา อิมสฺมิํ วา โลเก ปรสฺมิํ วา, เทวา วา เต มนุสฺสา วา ติ.\csromanspacer{} อิเม โภ กสฺสป จนฺทิมสูริยา ปรสฺมิํ โลเก น อิมสฺมิํ, เทวา เต น มนุสฺสาติ.\csromanspacer{} อิมินาปิ โข เต ราชญฺญ ปริยาเยน เอวํ โหตุ “อิติปิ อตฺถิ ปโร โลโก, อตฺถิ สตฺตา โอปปาติกา, อตฺถิ สุกตทุกฺกฏานํ กมฺมานํ ผลํ วิปาโก”ติ.}

\csromanmatikaanchor{mtk.140}
\csromantocmarkref{section}{โจร-อุปมา}{mtk.140}
\bookmark[dest={mtk.140},level=1]{โจร-อุปมา}
\proseitem{412}{กิญฺจาปิ ภวํ กสฺสโป เอวมาห.\csromanspacer{} อถ โข เอวํ เม เอตฺถ โหติ “อิติปิ นตฺถิ ปโร โลโก, นตฺถิ สตฺตา โอปปาติกา, นตฺถิ สุกตทุกฺกฏานํ กมฺมานํ ผลํ วิปาโก”ติ.\csromanspacer{} อตฺถิ ปน ราชญฺญ ปริยาโย, เยน เต ปริยาเยน เอวํ โหติ “อิติปิ นตฺถิ ปโร โลโก, นตฺถิ สตฺตา โอปปาติกา, นตฺถิ สุกตทุกฺกฏานํ กมฺมานํ ผลํ วิปาโก”ติ. \csromanfolio{256}อตฺถิ โภ กสฺสป ปริยาโย, เยน เม ปริยาเยน เอวํ โหติ “อิติปิ นตฺถิ ปโร โลโก, นตฺถิ สตฺตา โอปปาติกา, นตฺถิ สุกตทุกฺกฏานํ กมฺมานํ ผลํ วิปาโก”ติ.\csromanspacer{} ยถา กถํ วิย ราชญฺญาติ.\csromanspacer{} อิธ เม โภ กสฺสป มิตฺตามจฺจา ญาติสาโลหิตา ปาณาติปาตี อทินฺนาทายี กาเมสุมิจฺฉาจารี มุสาวาที ปิสุณวาจา ผรุสวาจา สมฺผปฺปลาปี อภิชฺฌาลู พฺยาปนฺนจิตฺตา มิจฺฉาทิฏฺฐี, เต อปเรน สมเยน อาพาธิกา โหนฺติ ทุกฺขิตา พาฬฺหคิลานา.\csromanspacer{} ยทาหํ ชานามิ “น ทานิเม อิมมฺหา อาพาธา วุฏฺฐหิสฺสนฺตี”ติ, ตฺยาหํ อุปสงฺกมิตฺวา เอวํ วทามิ “สนฺติ โข โภ เอเก สมณพฺราหฺมณา เอวํวาทิโน เอวํทิฏฺฐิโน ‘เย เต ปาณาติปาตี อทินฺนาทายี กาเมสุมิจฺฉาจารี มุสาวาที ปิสุณวาจา ผรุสวาจา สมฺผปฺปลาปี อภิชฺฌาลู พฺยาปนฺนจิตฺตา มิจฺฉาทิฏฺฐี, เต กายสฺส เภทา ปรํ มรณา อปายํ ทุคฺคติํ วินิปาตํ นิรยํ อุปปชฺชนฺตี’ติ.\csromanspacer{} ภวนฺโต โข ปาณาติปาตี, อทินฺนาทายี, กาเมสุมิจฺฉาจารี มุสาวาที ปิสุณวาจา ผรุสวาจา สมฺผปฺปลาปี อภิชฺฌาลู พฺยาปนฺนจิตฺตา มิจฺฉาทิฏฺฐี, สเจ เตสํ ภวตํ สมณพฺราหฺมณานํ สจฺจํ วจนํ, ภวนฺโต กายสฺส เภทา ปรํ มรณา อปายํ ทุคฺคติํ วินิปาตํ นิรยํ อุปปชฺชิสฺสนฺติ, สเจ โภ กายสฺส เภทา ปรํ มรณา อปายํ ทุคฺคติํ วินิปาตํ นิรยํ อุปปชฺเชยฺยาถ, เยน เม อาคนฺตฺวา อาโรเจยฺยาถ ‘อิติปิ อตฺถิ ปโร โลโก, อตฺถิ สตฺตา โอปปาติกา, อตฺถิ สุกตทุกฺกฏานํ กมฺมานํ ผลํ วิปาโก’ติ, ภวนฺโต โข ปน เม สทฺธายิกา ปจฺจยิกา, ยํ ภวนฺเตหิ ทิฏฺฐํ, ยถา สามํ ทิฏฺฐํ เอวเมตํ ภวิสฺสตี”ติ.\csromanspacer{} เต เม “สาธู”ติ ปฏิสฺสุตฺวา เนว อาคนฺตฺวา อาโรเจนฺติ, น ปน ทูตํ ปหิณนฺติ.\csromanspacer{} อยมฺปิ โข โภ กสฺสป ปริยาโย, เยน เม ปริยาเยน เอวํ โหติ “อิติปิ นตฺถิ ปโร โลโก, นตฺถิ สตฺตา โอปปาติกา, นตฺถิ สุกตทุกฺกฏานํ กมฺมานํ ผลํ วิปาโก”ติ.}

\vspace{\tipitakaparskip}
\csromancenter{\textbf{(2) โจร-อุปมา}}
\proseitem{413}{เตน หิ ราชญฺญ ตญฺเญเวตฺถ ปฏิปุจฺฉิสฺสามิ, ยถา เต ขเมยฺย, ตถา นํ พฺยากเรยฺยาสิ, ตํ กิํ มญฺญสิ ราชญฺญ, อิธ เต ปุริสา โจรํ อาคุจาริํ คเหตฺวา ทสฺเสยฺยุํ “อยํ เต ภนฺเต โจโร อาคุจารี, อิมสฺส ยํ อิจฺฉสิ, ตํ ทณฺฑํ ปเณหี”ติ.\csromanspacer{} เต ตฺวํ เอวํ วเทยฺยาสิ “เตน หิ โภ อิมํ ปุริสํ ทฬฺหาย รชฺชุยา ปจฺฉาพาหํ คาฬฺหพนฺธนํ}

\vspace{\tipitakaparskip}
\csromancenter{ปายาสิสุตฺต พนฺธิตฺวา ขุรมุณฺฑํ กริตฺวา\footnote{กาเรตฺวา (สฺยา, ก)} ขรสฺสเรน ปณเวน รถิกาย รถิกํ\footnote{รถิยาย รถิยํ (พหูสุ)} สิงฺฆาฏเกน สิงฺฆาฏกํ ปริเนตฺวา ทกฺขิเณน ทฺวาเรน นิกฺขมิตฺวา ทกฺขิณโต นครสฺส อาฆาตเน สีสํ ฉินฺทถา”ติ.\csromanspacer{} เต “สาธู”ติ ปฏิสฺสุตฺวา ตํ ปุริสํ ทฬฺหาย รชฺชุยา ปจฺฉาพาหํ คาฬฺหพนฺธนํ พนฺธิตฺวา ขุรมุณฺฑํ กริตฺวา ขรสฺสเรน ปณเวน รถิกาย รถิกํ สิงฺฆาฏเกน สิงฺฆาฏกํ ปริเนตฺวา ทกฺขิเณน ทฺวาเรน นิกฺขมิตฺวา ทกฺขิณโต นครสฺส อาฆาตเน นิสีทาเปยฺยุํ.\csromanspacer{} ลเภยฺย นุ โข โส โจโร โจรฆาเตสุ “อาคเมนฺตุ ตาว ภวนฺโต โจรฆาตา อมุกสฺมิํ เม คาเม วา นิคเม วา มิตฺตามจฺจา ญาติสาโลหิตา, ยาวาหํ เตสํ อุทฺทิสิตฺวา อาคจฺฉามี”ติ, อุทาหุ วิปฺปลปนฺตสฺเสว โจรฆาตา สีสํ ฉินฺเทยฺยุนฺติ.\csromanspacer{} น หิ โส โภ กสฺสป โจโร ลเภยฺย โจรฆาเตสุ “อาคเมนฺตุ ตาว ภวนฺโต โจรฆาตา อมุกสฺมิํ เม คาเม วา นิคเม วา มิตฺตามจฺจา ญาติสาโลหิตา, ยาวาหํ เตสํ อุทฺทิสิตฺวา อาคจฺฉามี”ติ.\csromanspacer{} อถ โข นํ วิปฺปลปนฺตสฺเสว โจรฆาตา สีสํ ฉินฺเทยฺยุนฺติ.\csromanspacer{} โส หิ นาม ราชญฺญ โจโร มนุสฺโส มนุสฺสภูเตสุ โจรฆาเตสุ น ลภิสฺสติ “อาคเมนฺตุ ตาว ภวนฺโต โจรฆาตา อมุกสฺมิํ เม คาเม วา นิคเม วา มิตฺตามจฺจา ญาติสาโลหิตา, ยาวาหํ เตสํ อุทฺทิสิตฺวา อาคจฺฉามี”ติ.\csromanspacer{} กิํ ปน เต มิตฺตามจฺจา ญาติสาโลหิตา ปาณาติปาตี อทินฺนาทายี กาเมสุมิจฺฉาจารี มุสาวาที ปิสุณวาจา ผรุสวาจา สมฺผปฺปลาปี อภิชฺฌาลู พฺยาปนฺนจิตฺตา มิจฺฉาทิฏฺฐี, เต กายสฺส เภทา ปรํ มรณา อปายํ ทุคฺคติํ วินิปาตํ นิรยํ อุปปนฺนา ลภิสฺสนฺติ นิรยปาเลสุ “อาคเมนฺตุ ตาว ภวนฺโต นิรยปาลา, ยาว มยํ ปายาสิสฺส ราชญฺญสฺส คนฺตฺวา อาโรเจม ‘อิติปิ อตฺถิ ปโร โลโก, อตฺถิ สตฺตา โอปปาติกา, อตฺถิ สุกตทุกฺกฏานํ กมฺมานํ ผลํ วิปาโก’ติ”.\csromanspacer{} อิมินาปิ โข เต ราชญฺญ ปริยาเยน เอวํ โหตุ “อิติปิ อตฺถิ ปโร โลโก, อตฺถิ สตฺตา โอปปาติกา, อตฺถิ สุกตทุกฺกฏานํ กมฺมานํ ผลํ วิปาโก”ติ.}
\proseitem{414}{\csromanfolio{257}กิญฺจาปิ ภวํ กสฺสโป เอวมาห.\csromanspacer{} อถ โข เอวํ เม เอตฺถ โหติ “อิติปิ นตฺถิ ปโร โลโก, นตฺถิ สตฺตา โอปปาติกา, นตฺถิ \csromanfolio{258}สุกตทุกฺกฏานํ กมฺมานํ ผลํ วิปาโก”ติ.\csromanspacer{} อตฺถิ ปน ราชญฺญ ปริยาโย, เยน เต ปริยาเยน เอวํ โหติ “อิติปิ นตฺถิ ปโร โลโก, นตฺถิ สตฺตา โอปปาติกา, นตฺถิ สุกตทุกฺกฏานํ กมฺมานํ ผลํ วิปาโก”ติ.\csromanspacer{} อตฺถิ โภ กสฺสป ปริยาโย, เยน เม ปริยาเยน เอวํ โหติ “อิติปิ นตฺถิ ปโร โลโก, นตฺถิ สตฺตา โอปปาติกา, นตฺถิ สุกตทุกฺกฏานํ กมฺมานํ ผลํ วิปาโก”ติ.\csromanspacer{} ยถา กถํ วิย ราชญฺญาติ.\csromanspacer{} อิธ เม โภ กสฺสป มิตฺตามจฺจา ญาติสาโลหิตา ปาณาติปาตา ปฏิวิรตา อทินฺนาทานา ปฏิวิรตา กาเมสุมิจฺฉาจารา ปฏิวิรตา มุสาวาทา ปฏิวิรตา ปิสุณาย วาจาย ปฏิวิรตา ผรุสาย วาจาย ปฏิวิรตา สมฺผปฺปลาปา ปฏิวิรตา อนภิชฺฌาลู อพฺยาปนฺนจิตฺตา สมฺมาทิฏฺฐี, เต อปเรน สมเยน อาพาธิกา โหนฺติ ทุกฺขิตา พาฬฺหคิลานา.\csromanspacer{} ยทาหํ ชานามิ “น ทานิเม อิมมฺหา อาพาธา วุฏฺฐหิสฺสนฺตี”ติ, ตฺยาหํ อุปสงฺกมิตฺวา เอวํ วทามิ “สนฺติ โข โภ เอเก สมณพฺราหฺมณา เอวํวาทิโน เอวํทิฏฺฐิโน ‘เย เต ปาณาติปาตา ปฏิวิรตา อทินฺนาทานา ปฏิวิรตา กาเมสุมิจฺฉาจารา ปฏิวิรตา มุสาวาทา ปฏิวิรตา ปิสุณาย วาจาย ปฏิวิรตา ผรุสาย วาจาย ปฏิวิรตา สมฺผปฺปลาปา ปฏิวิรตา อนภิชฺฌาลู อพฺยาปนฺนจิตฺตา สมฺมาทิฏฺฐี, เต กายสฺส เภทา ปรํ มรณา สุคติํ สคฺคํ โลกํ อุปปชฺชนฺตี’ติ, ภวนฺโต โข ปาณาติปาตา ปฏิวิรตา อทินฺนาทานา ปฏิวิรตา กาเมสุมิจฺฉาจารา ปฏิวิรตา มุสาวาทา ปฏิวิรตา ปิสุณาย วาจาย ปฏิวิรตา ผรุสาย วาจาย ปฏิวิรตา สมฺผปฺปลาปา ปฏิวิรตา อนภิชฺฌาลู อพฺยาปนฺนจิตฺตา สมฺมาทิฏฺฐี, สเจ เตสํ ภวตํ สมณพฺราหฺมณานํ สจฺจํ วจนํ, ภวนฺโต กายสฺส เภทา ปรํ มรณา สุคติํ สคฺคํ โลกํ อุปปชฺชิสฺสนฺติ.\csromanspacer{} สเจ โภ กายสฺส เภทา ปรํ มรณา สุคติํ สคฺคํ โลกํ อุปปชฺเชยฺยาถ, เยน เม อาคนฺตฺวา อาโรเจยฺยาถ ‘อิติปิ อตฺถิ ปโร โลโก, อตฺถิ สตฺตา โอปปาติกา, อตฺถิ สุกตทุกฺกฏานํ กมฺมานํ ผลํ วิปาโก’ติ.\csromanspacer{} ภวนฺโต โข ปน เม สทฺธายิกา ปจฺจยิกา, ยํ ภวนฺเตหิ ทิฏฺฐํ, ยถา สามํ ทิฏฺฐํ เอวเมตํ ภวิสฺสตี”ติ.\csromanspacer{} เต เม “สาธู”ติ ปฏิสฺสุตฺวา เนว อาคนฺตฺวา อาโรเจนฺติ, น ปน ทูตํ ปหิณนฺติ.\csromanspacer{} อยมฺปิ โข โภ กสฺสป ปริยาโย, เยน เม ปริยาเยน เอวํ โหติ “อิติปิ นตฺถิ ปโร โลโก, นตฺถิ สตฺตา โอปปาติกา, นตฺถิ สุกตทุกฺกฏานํ กมฺมานํ ผลํ วิปาโก”ติ.}

\csromanmatikaanchor{mtk.141}
\csromantocmarkref{section}{คูถกูปปุริส-อุปมา}{mtk.141}
\bookmark[dest={mtk.141},level=1]{คูถกูปปุริส-อุปมา}
\csromanheader{1}{10. ปายาสิสุตฺต \textbf{(3) คูถกูปปุริส-อุปมา}}
\proseitem{415}{\csromanfolio{259}เตน หิ ราชญฺญ อุปมํ เต กริสฺสามิ.\csromanspacer{} อุปมาย มิเธกจฺเจ\footnote{อุปมายปิเธกจฺเจ (สี, สฺยา), อุปมายปิ อิเธกจฺเจ (อิ)} วิญฺญู ปุริสา ภาสิตสฺส อตฺถํ อาชานนฺติ.\csromanspacer{} เสยฺยถาปิ ราชญฺญ ปุริโส คูถกูเป สสีสกํ\footnote{สสีสโก (สฺยา)} นิมุคฺโค อสฺส.\csromanspacer{} อถ ตฺวํ ปุริเส อาณาเปยฺยาสิ “เตน หิ โภ ตํ ปุริสํ ตมฺหา คูถกูปา อุทฺธรถา”ติ.\csromanspacer{} เต “สาธู”ติ ปฏิสฺสุตฺวา ตํ ปุริสํ ตมฺหา คูถกูปา อุทฺธเรยฺยุํ.\csromanspacer{} เต ตฺวํ เอวํ วเทยฺยาสิ “เตน หิ โภ ตสฺส ปุริสสฺส กายา เวฬุเปสิกาหิ คูถํ สุนิมฺมชฺชิตํ นิมฺมชฺชถา”ติ.\csromanspacer{} เต “สาธู”ติ ปฏิสฺสุตฺวา ตสฺส ปุริสสฺส กายา เวฬุเปสิกาหิ คูถํ สุนิมฺมชฺชิตํ นิมฺมชฺเชยฺยุํ.\csromanspacer{} เต ตฺวํ เอวํ วเทยฺยาสิ “เตน หิ โภ ตสฺส ปุริสสฺส กายํ ปณฺฑุมตฺติกาย ติกฺขตฺตุํ สุพฺพฏฺฏิตํ อุพฺพฏฺเฏถา”ติ\footnote{สุปฺปฏฺฏิตํ อุปฺปฏฺเฏถาติ (ก)}.\csromanspacer{} เต ตสฺส ปุริสสฺส กายํ ปณฺฑุมตฺติกาย ติกฺขตฺตุํ สุพฺพฏฺฏิตํ อุพฺพฏฺเฏยฺยุํ.\csromanspacer{} เต ตฺวํ เอวํ วเทยฺยาสิ “เตน หิ โภ ตํ ปุริสํ เตเลน อพฺภญฺชิตฺวา สุขุเมน จุณฺเณน ติกฺขตฺตุํ สุปฺปโธตํ กโรถา”ติ.\csromanspacer{} เต ตํ ปุริสํ เตเลน อพฺภญฺชิตฺวา สุขุเมน จุณฺเณน ติกฺขตฺตุํ สุปฺปโธตํ กเรยฺยุํ.\csromanspacer{} เต ตฺวํ เอวํ วเทยฺยาสิ “เตน หิ โภ ตสฺส ปุริสสฺส เกสมสฺสุํ กปฺเปถา”ติ.\csromanspacer{} เต ตสฺส ปุริสสฺส เกสมสฺสุํ กปฺเปยฺยุํ.\csromanspacer{} เต ตฺวํ เอวํ วเทยฺยาสิ “เตน หิ โภ ตสฺส ปุริสสฺส มหคฺฆญฺจ มาลํ มหคฺฆญฺจ วิเลปนํ มหคฺฆานิ จ วตฺถานิ อุปหรถา”ติ.\csromanspacer{} เต ตสฺส ปุริสสฺส มหคฺฆญฺจ มาลํ มหคฺฆญฺจ วิเลปนํ มหคฺฆานิ จ วตฺถานิ อุปหเรยฺยุํ.\csromanspacer{} เต ตฺวํ เอวํ วเทยฺยาสิ “เตน หิ โภ ตํ ปุริสํ ปาสาทํ อาโรเปตฺวา ปญฺจกามคุณานิ อุปฏฺฐาเปถา”ติ.\csromanspacer{} เต ตํ ปุริสํ ปาสาทํ อาโรเปตฺวา ปญฺจกามคุณานิ อุปฏฺฐาเปยฺยุํ.}

\prose{ตํ กิํ มญฺญสิ ราชญฺญ, อปิ นุ ตสฺส ปุริสสฺส สุนฺหาตสฺส สุวิลิตฺตสฺส สุกปฺปิตเกสมสฺสุสฺส อามุกฺกมาลาภรณสฺส โอทาตวตฺถวสนสฺส อุปริปาสาทวรคตสฺส ปญฺจหิ กามคุเณหิ สมปฺปิตสฺส สมงฺคีภูตสฺส ปริจารยมานสฺส ปุนเทว ตสฺมิํ คูถกูเป นิมุชฺชิตุกามตา\footnote{นิมุชฺชิตุกามฺยตา (สฺยา, ก)} อสฺสาติ.\csromanspacer{} โน หิทํ โภ กสฺสป.\csromanspacer{} ตํ กิสฺส เหตุ.\csromanspacer{} อสุจิ โภ กสฺสป คูถกูโป อสุจิ เจว อสุจิสงฺขาโต จ ทุคฺคนฺโธ จ ทุคฺคนฺธสงฺขาโต จ เชคุจฺโฉ จ \csromanfolio{260}เชคุจฺฉสงฺขาโต จ ปฏิกูโล จ ปฏิกูลสงฺขาโต จาติ.\csromanspacer{} เอวเมว โข ราชญฺญ มนุสฺสา เทวานํ อสุจี เจว อสุจิสงฺขาตา จ ทุคฺคนฺธา จ ทุคฺคนฺธสงฺขาตา จ เชคุจฺฉา จ เชคุจฺฉสงฺขาตา จ ปฏิกูลา จ ปฏิกูลสงฺขาตา จ.\csromanspacer{} โยชนสตํ โข ราชญฺญ มนุสฺสคนฺโธ เทเว อุพฺพาธติ.\csromanspacer{} กิํ ปน เต มิตฺตามจฺจา ญาติสาโลหิตา ปาณาติปาตา ปฏิวิรตา อทินฺนาทานา ปฏิวิรตา กาเมสุมิจฺฉาจารา ปฏิวิรตา มุสาวาทา ปฏิวิรตา ปิสุณาย วาจาย ปฏิวิรตา ผรุสาย วาจาย ปฏิวิรตา สมฺผปฺปลาปา ปฏิวิรตา อนภิชฺฌาลู อพฺยาปนฺนจิตฺตา สมฺมาทิฏฺฐี, กายสฺส เภทา ปรํ มรณา สุคติํ สคฺคํ โลกํ อุปปนฺนา, เต อาคนฺตฺวา อาโรเจสฺสนฺติ “อิติปิ อตฺถิ ปโร โลโก, อตฺถิ สตฺตา โอปปาติกา, อตฺถิ สุกตทุกฺกฏานํ กมฺมานํ ผลํ วิปาโก”ติ.\csromanspacer{} อิมินาปิ โข เต ราชญฺญ ปริยาเยน เอวํ โหตุ “อิติปิ อตฺถิ ปโร โลโก, อตฺถิ สตฺตา โอปปาติกา, อตฺถิ สุกตทุกฺกฏานํ กมฺมานํ ผลํ วิปาโก”ติ.}

\proseitem{416}{กิญฺจาปิ ภวํ กสฺสโป เอวมาห.\csromanspacer{} อถ โข เอวํ เม เอตฺถ โหติ “อิติปิ นตฺถิ ปโร โลโก, นตฺถิ สตฺตา โอปปาติกา, นตฺถิ สุกตทุกฺกฏานํ กมฺมานํ ผลํ วิปาโก”ติ.\csromanspacer{} อตฺถิ ปน ราชญฺญ ปริยาโย ฯเปฯ.\csromanspacer{} อตฺถิ โภ กสฺสป ปริยาโย ฯเปฯ.\csromanspacer{} ยถา กถํ วิย ราชญฺญาติ.\csromanspacer{} อิธ เม โภ กสฺสป มิตฺตามจฺจา ญาติสาโลหิตา ปาณาติปาตา ปฏิวิรตา อทินฺนาทานา ปฏิวิรตา กาเมสุมิจฺฉาจารา ปฏิวิรตา มุสาวาทา ปฏิวิรตา สุราเมรยมชฺชปมาทฏฺฐานา ปฏิวิรตา, เต อปเรน สมเยน อาพาธิกา โหนฺติ ทุกฺขิตา พาฬฺหคิลานา.\csromanspacer{} ยทาหํ ชานามิ “น ทานิเม อิมมฺหา อาพาธา วุฏฺฐหิสฺสนฺตี”ติ, ตฺยาหํ อุปสงฺกมิตฺวา เอวํ วทามิ “สนฺติ โข โภ เอเก สมณพฺราหฺมณา เอวํวาทิโน เอวํทิฏฺฐิโน ‘เย เต ปาณาติปาตา ปฏิวิรตา อทินฺนาทานา ปฏิวิรตา กาเมสุมิจฺฉาจารา ปฏิวิรตา มุสาวาทา ปฏิวิรตา สุราเมรยมชฺชปมาทฏฺฐานา ปฏิวิรตา, เต กายสฺส เภทา ปรํ มรณา สุคติํ สคฺคํ โลกํ อุปปชฺชนฺติ เทวานํ ตาวติํสานํ สหพฺยตนฺ’ติ.\csromanspacer{} ภวนฺโต โข ปาณาติปาตา ปฏิวิรตา อทินฺนาทานา ปฏิวิรตา กาเมสุมิจฺฉาจารา ปฏิวิรตา มุสาวาทา ปฏิวิรตา สุราเมรยมชฺชปมาทฏฺฐานา ปฏิวิรตา, สเจ เตสํ ภวตํ สมณพฺราหฺมณานํ สจฺจํ วจนํ, ภวนฺโต กายสฺส เภทา}

\vspace{\tipitakaparskip}
\csromancenter{ปายาสิสุตฺต ปรํ มรณา สุคติํ สคฺคํ โลกํ อุปปชฺชิสฺสนฺติ เทวานํ ตาวติํสานํ สหพฺยตํ, สเจ โภ กายสฺส เภทา ปรํ มรณา สุคติํ สคฺคํ โลกํ อุปปชฺเชยฺยาถ เทวานํ ตาวติํสานํ สหพฺยตํ, เยน เม อาคนฺตฺวา อาโรเจยฺยาถ ‘อิติปิ อตฺถิ ปโร โลโก, อตฺถิ สตฺตา โอปปาติกา, อตฺถิ สุกตทุกฺกฏานํ กมฺมานํ ผลํ วิปาโก’ติ.\csromanspacer{} ภวนฺโต โข ปน เม สทฺธายิกา ปจฺจยิกา, ยํ ภวนฺเตหิ ทิฏฺฐํ, ยถา สามํ ทิฏฺฐํ เอวเมตํ ภวิสฺสตี”ติ.\csromanspacer{} เต เม “สาธู”ติ ปฏิสฺสุตฺวา เนว อาคนฺตฺวา อาโรเจนฺติ, น ปน ทูตํ ปหิณนฺติ.\csromanspacer{} อยมฺปิ โข โภ กสฺสป ปริยาโย, เยน เม ปริยาเยน เอวํ โหติ “อิติปิ นตฺถิ ปโร โลโก, นตฺถิ สตฺตา โอปปาติกา, นตฺถิ สุกตทุกฺกฏานํ กมฺมานํ ผลํ วิปาโก”ติ.}
\csromancenter{\textbf{(4) ตาวติํสเทว-อุปมา}}
\csromanmatikaanchor{mtk.142}
\csromantocmarkref{section}{ตาวติํสเทว-อุปมา}{mtk.142}
\bookmark[dest={mtk.142},level=1]{ตาวติํสเทว-อุปมา}
\csromanmatikaanchor{mtk.143}
\csromantocmarkref{section}{ชจฺจนฺธ-อุปมา}{mtk.143}
\bookmark[dest={mtk.143},level=1]{ชจฺจนฺธ-อุปมา}
\proseitem{417}{\csromanfolio{261}เตน หิ ราชญฺญ ตญฺเญเวตฺถ ปฏิปุจฺฉิสฺสามิ, ยถา เต ขเมยฺย, ตถา นํ พฺยากเรยฺยาสิ.\csromanspacer{} ยํ โข ปน ราชญฺญ มานุสฺสกํ วสฺสสตํ เทวานํ ตาวติํสานํ เอโส เอโก รตฺตินฺทิโว\footnote{รตฺติทิโว (ก)}, ตาย รตฺติยา ติํสรตฺติโย มาโส, เตน มาเสน ทฺวาทสมาสิโย สํวจฺฉโร, เตน สํวจฺฉเรน ทิพฺพํ วสฺสสหสฺสํ เทวานํ ตาวติํสานํ อายุปฺปมาณํ.\csromanspacer{} เย เต มิตฺตามจฺจา ญาติสาโลหิตา ปาณาติปาตา ปฏิวิรตา อทินฺนาทานา ปฏิวิรตา กาเมสุมิจฺฉาจารา ปฏิวิรตา มุสาวาทา ปฏิวิรตา สุราเมรยมชฺชปมาทฏฺฐานา ปฏิวิรตา, เต กายสฺส เภทา ปรํ มรณา สุคติํ สคฺคํ โลกํ อุปปนฺนา เทวานํ ตาวติํสานํ สหพฺยตํ.\csromanspacer{} สเจ ปน เตสํ เอวํ ภวิสฺสติ “ยาว มยํ เทฺว วา ตีณิ วา รตฺตินฺทิวา ทิพฺเพหิ ปญฺจหิ กามคุเณหิ สมปฺปิตา สมงฺคีภูตา ปริจาเรม, อถ มยํ ปายาสิสฺส ราชญฺญสฺส คนฺตฺวา อาโรเจยฺยาม ‘อิติปิ อตฺถิ ปโร โลโก, อตฺถิ สตฺตา โอปปาติกา, อตฺถิ สุกตทุกฺกฏานํ กมฺมานํ ผลํ วิปาโก’ติ”.\csromanspacer{} อปิ นุ เต อาคนฺตฺวา อาโรเจยฺยุํ “อิติปิ อตฺถิ ปโร โลโก, อตฺถิ สตฺตา โอปปาติกา, อตฺถิ สุกตทุกฺกฏานํ กมฺมานํ ผลํ วิปาโก”ติ.\csromanspacer{} โน หิทํ โภ กสฺสป.\csromanspacer{} อปิ หิ มยํ โภ กสฺสป จิรํ กาลงฺกตาปิ ภเวยฺยาม.\csromanspacer{} โก ปเนตํ โภโต กสฺสปสฺส \csromanfolio{262}อาโรเจติ “อตฺถิ เทวา ตาวติํสา”ติ วา “เอวํทีฆายุกา เทวา ตาวติํสา”ติ วา, น มยํ โภโต กสฺสปสฺส สทฺทหาม “อตฺถิ เทวา ตาวติํสา”ติ วา “เอวํทีฆายุกา เทวา ตาวติํสา”ติ วาติ.}

\vspace{\tipitakaparskip}
\csromancenter{\textbf{(5) ชจฺจนฺธ อุปมา}}
\proseitem{418}{เสยฺยถาปิ ราชญฺญ ชจฺจนฺโธ ปุริโส น ปสฺเสยฺย กณฺหสุกฺกานิ รูปานิ, น ปสฺเสยฺย นีลกานิ รูปานิ, น ปสฺเสยฺย ปีตกานิ รูปานิ, น ปสฺเสยฺย โลหิตกานิ รูปานิ, น ปสฺเสยฺย มญฺชิฏฺฐกานิ\footnote{มญฺเชฏฺฐกานิ (สฺยา)} รูปานิ, น ปสฺเสยฺย สมวิสมํ, น ปสฺเสยฺย ตารกานิ รูปานิ, น ปสฺเสยฺย จนฺทิมสูริเย, โส เอวํ วเทยฺย “นตฺถิ กณฺหสุกฺกานิ รูปานิ, นตฺถิ กณฺหสุกฺกานํ รูปานํ ทสฺสาวี.\csromanspacer{} นตฺถิ นีลกานิ รูปานิ, นตฺถิ นีลกานํ รูปานํ ทสฺสาวี.\csromanspacer{} นตฺถิ ปีตกานิ รูปานิ, นตฺถิ ปีตกานํ รูปานํ ทสฺสาวี.\csromanspacer{} นตฺถิ โลหิตกานิ รูปานิ, นตฺถิ โลหิตกานํ รูปานํ ทสฺสาวี.\csromanspacer{} นตฺถิ มญฺชิฏฺฐกานิ รูปานิ, นตฺถิ มญฺชิฏฺฐกานํ รูปานํ ทสฺสาวี.\csromanspacer{} นตฺถิ สมวิสมํ, นตฺถิ สมวิสมสฺส ทสฺสาวี.\csromanspacer{} นตฺถิ ตารกานิ รูปานิ, นตฺถิ ตารกานํ รูปานํ ทสฺสาวี.\csromanspacer{} นตฺถิ จนฺทิมสูริยา, นตฺถิ จนฺทิมสูริยานํ ทสฺสาวี.\csromanspacer{} อหเมตํ น ชานามิ, อหเมตํ น ปสฺสามิ, ตสฺมา ตํ นตฺถี”ติ, สมฺมา นุ โข โส ราชญฺญ วทมาโน วเทยฺยาติ.\csromanspacer{} โน หิทํ โภ กสฺสป, อตฺถิ กณฺหสุกฺกานิ รูปานิ, อตฺถิ กณฺหสุกฺกานํ รูปานํ ทสฺสาวี.\csromanspacer{} อตฺถิ นีลกานิ รูปานิ, อตฺถิ นีลกานํ รูปานํ ทสฺสาวี ฯเปฯ.\csromanspacer{} อตฺถิ สมวิสมํ, อตฺถิ สมวิสมสฺส ทสฺสาวี.\csromanspacer{} อตฺถิ ตารกานิ รูปานิ, อตฺถิ ตารกานํ รูปานํ ทสฺสาวี.\csromanspacer{} อตฺถิ จนฺทิมสูริยา, อตฺถิ จนฺทิมสูริยานํ ทสฺสาวี. “อหเมตํ น ชานามิ, อหเมตํ น ปสฺสามิ, ตสฺมา ตํ นตฺถี”ติ, น หิ โส โภ กสฺสป สมฺมา วทมาโน วเทยฺยาติ.\csromanspacer{} เอวเมว โข ตฺวํ ราชญฺญ ชจฺจนฺธูปโม มญฺเญ ปฏิภาสิ ยํ มํ ตฺวํ เอวํ วเทสิ.}

\prose{โก ปเนตํ โภโต กสฺสปสฺส อาโรเจติ “อตฺถิ เทวา ตาวติํสา”ติ วา “เอวํทีฆายุกา เทวา ตาวติํสา”ติ วา.\csromanspacer{} น มยํ โภโต กสฺสปสฺส สทฺทหาม “อตฺถิ เทวา ตาวติํสา”ติ วา “เอวํทีฆายุกา เทวา ตาวติํสา”ติ วาติ.\csromanspacer{} น โข ราชญฺญ เอวํ ปโร โลโก ทฏฺฐพฺโพ, ยถา ตฺวํ มญฺญสิ อิมินา มํสจกฺขุนา.\csromanspacer{} เย โข เต ราชญฺญ สมณพฺราหฺมณา อรญฺญวนปตฺถานิ ปนฺตานิ เสนาสนานิ}

\vspace{\tipitakaparskip}
\csromancenter{ปายาสิสุตฺต ปฏิเสวนฺติ, เต ตตฺถ อปฺปมตฺตา อาตาปิโน ปหิตตฺตา วิหรนฺตา ทิพฺพจกฺขุํ วิโสเธนฺติ, เต ทิพฺเพน จกฺขุนา วิสุทฺเธน อติกฺกนฺตมานุสเกน อิมญฺเจว โลกํ ปสฺสนฺติ ปรญฺจ สตฺเต จ โอปปาติเก.\csromanspacer{} เอวญฺจ โข ราชญฺญ ปโร โลโก ทฏฺฐพฺโพ, นเตฺวว ยถา ตฺวํ มญฺญสิ อิมินา มํสจกฺขุนา.\csromanspacer{} อิมินาปิ โข เต ราชญฺญ ปริยาเยน เอวํ โหตุ “อิติปิ อตฺถิ ปโร โลโก, อตฺถิ สตฺตา โอปปาติกา, อตฺถิ สุกตทุกฺกฏานํ กมฺมานํ ผลํ วิปาโก”ติ.}
\csromanmatikaanchor{mtk.144}
\csromantocmarkref{section}{คพฺภินี-อุปมา}{mtk.144}
\bookmark[dest={mtk.144},level=1]{คพฺภินี-อุปมา}
\proseitem{419}{\csromanfolio{263}กิญฺจาปิ ภวํ กสฺสโป เอวมาห.\csromanspacer{} อถ โข เอวํ เม เอตฺถ โหติ “อิติปิ นตฺถิ ปโร โลโก, นตฺถิ สตฺตา โอปปาติกา, นตฺถิ สุกตทุกฺกฏานํ กมฺมานํ ผลํ วิปาโก”ติ.\csromanspacer{} อตฺถิ ปน ราชญฺญ ปริยาโย ฯเปฯ.\csromanspacer{} อตฺถิ โภ กสฺสป ปริยาโย ฯเปฯ.\csromanspacer{} ยถา กถํ วิย ราชญฺญาติ.\csromanspacer{} อิธาหํ โภ กสฺสป ปสฺสามิ สมณพฺราหฺมเณ สีลวนฺเต กลฺยาณธมฺเม ชีวิตุกาเม อมริตุกาเม สุขกาเม ทุกฺขปฏิกูเล.\csromanspacer{} ตสฺส มยฺหํ โภ กสฺสป เอวํ โหติ, สเจ โข อิเม โภนฺโต สมณพฺราหฺมณา สีลวนฺโต กลฺยาณธมฺมา เอวํ ชาเนยฺยุํ “อิโต โน มตานํ เสยฺโย ภวิสฺสตี”ติ.\csromanspacer{} อิทานิเม โภนฺโต สมณพฺราหฺมณา สีลวนฺโต กลฺยาณธมฺมา วิสํ วา ขาเทยฺยุํ, สตฺถํ วา อาหเรยฺยุํ, อุพฺพนฺธิตฺวา วา กาลงฺกเรยฺยุํ, ปปาเต วา ปปเตยฺยุํ.\csromanspacer{} ยสฺมา จ โข อิเม โภนฺโต สมณพฺราหฺมณา สีลวนฺโต กลฺยาณธมฺมา น เอวํ ชานนฺติ “อิโต โน มตานํ เสยฺโย ภวิสฺสตี”ติ.\csromanspacer{} ตสฺมา อิเม โภนฺโต สมณพฺราหฺมณา สีลวนฺโต กลฺยาณธมฺมา ชีวิตุกามา อมริตุกามา สุขกามา ทุกฺขปฏิกูลา (อตฺตานํ น มาเรนฺติ.)\footnote{( ) นตฺถิ (สฺยา, อิ)} อยํปิ โข โภ กสฺสป ปริยาโย, เยน เม ปริยาเยน เอวํ โหติ “อิติปิ นตฺถิ ปโร โลโก, นตฺถิ สตฺตา โอปปาติกา, นตฺถิ สุกตทุกฺกฏานํ กมฺมานํ ผลํ วิปาโก”ติ.}

\vspace{\tipitakaparskip}
\csromancenter{\textbf{(6) คพฺภินี-อุปมา}}
\proseitem{420}{เตน หิ ราชญฺญ อุปมํ เต กริสฺสามิ.\csromanspacer{} อุปมาย มิเธกจฺเจ วิญฺญู ปุริสา ภาสิตสฺส อตฺถํ อาชานนฺติ.\csromanspacer{} ภูตปุพฺพํ ราชญฺญ อญฺญตรสฺส \csromanfolio{264}พฺราหฺมณสฺส เทฺว ปชาปติโย อเหสุํ, เอกิสฺสา ปุตฺโต อโหสิ ทสวสฺสุทฺเทสิโก วา ทฺวาทสวสฺสุทฺเทสิโก วา, เอกา คพฺภินี อุปวิชญฺญา.\csromanspacer{} อถ โข โส พฺราหฺมโณ กาลมกาสิ.\csromanspacer{} อถ โข โส มาณวโก มาตุสปตฺติํ\footnote{มาตุสปติํ (สฺยา)} เอตทโวจ “ยมิทํ โภติ ธนํ วา ธญฺญํ วา รชตํ วา ชาตรูปํ วา, สพฺพํ ตํ มยฺหํ, นตฺถิ ตุเยฺหตฺถ กิญฺจิ, ปิตุ เม\footnote{ปิตุ เม สนฺตโก (สฺยา)} โภติ ทายชฺชํ นิยฺยาเทหี”ติ\footnote{นียฺยาเตหีติ (สี, อิ)}.\csromanspacer{} เอวํ วุตฺเต สา พฺราหฺมณี ตํ มาณวกํ เอตทโวจ “อาคเมหิ ตาว ตาต, ยาว วิชายามิ, สเจ กุมารโก ภวิสฺสติ, ตสฺสปิ เอกเทโส ภวิสฺสติ, สเจ กุมาริกา ภวิสฺสติ, สาปิ เต โอปโภคฺคา\footnote{อุปโภคฺคา (สฺยา)} ภวิสฺสตี”ติ.\csromanspacer{} ทุติยมฺปิ โข โส มาณวโก มาตุสปตฺติํ เอตทโวจ “ยมิทํ โภติ ธนํ วา ธญฺญํ วา รชตํ วา ชาตรูปํ วา, สพฺพํ ตํ มยฺหํ, นตฺถิ ตุเยฺหตฺถ กิญฺจิ, ปิตุ เม โภติ ทายชฺชํ นิยฺยาเทหี”ติ.\csromanspacer{} ทุติยมฺปิ โข สา พฺราหฺมณี ตํ มาณวกํ เอตทโวจ “อาคเมหิ ตาว ตาต, ยาว วิชายามิ, สเจ กุมารโก ภวิสฺสติ, ตสฺสปิ เอกเทโส ภวิสฺสติ, สเจ กุมาริกา ภวิสฺสติ, สาปิ เต โอปโภคฺคา4 ภวิสฺสตี”ติ.\csromanspacer{} ตติยมฺปิ โข โส มาณวโก มาตุสปตฺติํ เอตทโวจ “ยมิทํ โภติ ธนํ วา ธญฺญํ วา รชตํ วา ชาตรูปํ วา, สพฺพํ ตํ มยฺหํ, นตฺถิ ตุเยฺหตฺถ กิญฺจิ, ปิตุ เม โภติ ทายชฺชํ นิยฺยาเทหี”ติ.}

\prose{อถ โข สา พฺราหฺมณี สตฺถํ คเหตฺวา โอวรกํ ปวิสิตฺวา อุทรํ โอปาเทสิ\footnote{อุปฺปาเตสิ (สฺยา)} ยาว วิชายามิ ยทิ วา กุมารโก ยทิ วา กุมาริกาติ.\csromanspacer{} สา อตฺตานํ เจว ชีวิตญฺจ คพฺภญฺจ สาปเตยฺยญฺจ วินาเสสิ, ยถา ตํ พาลา อพฺยตฺตา อนยพฺยสนํ อาปนฺนา อโยนิโส ทายชฺชํ คเวสนฺตี.\csromanspacer{} เอวเมว โข ตฺวํ ราชญฺญ พาโล อพฺยตฺโต อนยพฺยสนํ อาปชฺชิสฺสสิ อโยนิโส ปรโลกํ คเวสนฺโต, เสยฺยถาปิ สา พฺราหฺมณี พาลา อพฺยตฺตา อนยพฺยสนํ อาปนฺนา อโยนิโส ทายชฺชํ คเวสนฺตี.\csromanspacer{} น โข ราชญฺญ สมณพฺราหฺมณา สีลวนฺโต กลฺยาณธมฺมา อปกฺกํ ปริปาเจนฺติ, อปิ จ ปริปากํ อาคเมนฺติ.\csromanspacer{} ปณฺฑิตานํ อตฺโถ หิ ราชญฺญ สมณพฺราหฺมณานํ สีลวนฺตานํ กลฺยาณธมฺมานํ ชีวิเตน.\csromanspacer{} ยถา ยถา โข ราชญฺญ สมณพฺราหฺมณา สีลวนฺโต กลฺยาณธมฺมา จิรํ ทีฆมทฺธานํ ติฏฺฐนฺติ,}

\vspace{\tipitakaparskip}
\csromancenter{ปายาสิสุตฺต ตถา ตถา พหุํ ปุญฺญํ ปสวนฺติ, พหุชนหิตาย จ ปฏิปชฺชนฺติ พหุชนสุขาย โลกานุกมฺปาย อตฺถาย หิตาย สุขาย เทวมนุสฺสานํ.\csromanspacer{} อิมินาปิ โข เต ราชญฺญ ปริยาเยน เอวํ โหตุ “อิติปิ อตฺถิ ปโร โลโก, อตฺถิ สตฺตา โอปปาติกา, อตฺถิ สุกตทุกฺกฏานํ กมฺมานํ ผลํ วิปาโก”ติ.}
\csromanmatikaanchor{mtk.145}
\csromantocmarkref{section}{สุปินก-อุปมา}{mtk.145}
\bookmark[dest={mtk.145},level=1]{สุปินก-อุปมา}
\proseitem{421}{\csromanfolio{265}กิญฺจาปิ ภวํ กสฺสโป เอวมาห.\csromanspacer{} อถ โข เอวํ เม เอตฺถ โหติ “อิติปิ นตฺถิ ปโร โลโก, นตฺถิ สตฺตา โอปปาติกา, นตฺถิ สุกตทุกฺกฏานํ กมฺมานํ ผลํ วิปาโก”ติ.\csromanspacer{} อตฺถิ ปน ราชญฺญ ปริยาโย ฯเปฯ.\csromanspacer{} อตฺถิ โภ กสฺสป ปริยาโย ฯเปฯ.\csromanspacer{} ยถา กถํ วิย ราชญฺญาติ.\csromanspacer{} อิธ เม โภ กสฺสป ปุริสา โจรํ อาคุจาริํ คเหตฺวา ทสฺเสนฺติ “อยํ เต ภนฺเต โจโร อาคุจารี, อิมสฺส ยํ อิจฺฉสิ, ตํ ทณฺฑํ ปเณหี”ติ.\csromanspacer{} ตฺยาหํ เอวํ วทามิ “เตน หิ โภ อิมํ ปุริสํ ชีวนฺตํเยว กุมฺภิยา ปกฺขิปิตฺวา มุขํ ปิทหิตฺวา อลฺเลน จมฺเมน โอนนฺธิตฺวา อลฺลาย มตฺติกาย พหลาวเลปนํ\footnote{พหลวิเลปนํ (สฺยา, ก)} กริตฺวา อุทฺธนํ อาโรเปตฺวา อคฺคิํ เทถา”ติ.\csromanspacer{} เต เม “สาธู”ติ ปฏิสฺสุตฺวา ตํ ปุริสํ ชีวนฺตํเยว กุมฺภิยา ปกฺขิปิตฺวา มุขํ ปิทหิตฺวา อลฺเลน จมฺเมน โอนนฺธิตฺวา อลฺลาย มตฺติกาย พหลาวเลปนํ กริตฺวา อุทฺธนํ อาโรเปตฺวา อคฺคิํ เทนฺติ.\csromanspacer{} ยทา มยํ ชานาม “กาลงฺกโต โส ปุริโส”ติ.\csromanspacer{} อถ นํ กุมฺภิํ โอโรเปตฺวา อุพฺภินฺทิตฺวา มุขํ วิวริตฺวา สณิกํ นิลฺโลเกม\footnote{วิโลเกม (สฺยา)} “อปฺเปว นามสฺส ชีวํ นิกฺขมนฺตํ ปสฺเสยฺยามา”ติ.\csromanspacer{} เนวสฺส มยํ ชีวํ นิกฺขมนฺตํ ปสฺสาม.\csromanspacer{} อยมฺปิ โข โภ กสฺสป ปริยาโย, เยน เม ปริยาเยน เอวํ โหติ “อิติปิ นตฺถิ ปโร โลโก, นตฺถิ สตฺตา โอปปาติกา, นตฺถิ สุกตทุกฺกฏานํ กมฺมานํ ผลํ วิปาโก”ติ.}

\vspace{\tipitakaparskip}
\csromancenter{\textbf{(7) สุปินก-อุปมา}}
\csromanmatikaanchor{mtk.146}
\csromantocmarkref{section}{สนฺตตฺต-อโยคุฬ-อุปมา}{mtk.146}
\bookmark[dest={mtk.146},level=1]{สนฺตตฺต-อโยคุฬ-อุปมา}
\proseitem{422}{เตน หิ ราชญฺญ ตญฺเญเวตฺถ ปฏิปุจฺฉิสฺสามิ, ยถา เต ขเมยฺย, ตถา นํ พฺยากเรยฺยาสิ.\csromanspacer{} อภิชานาสิ โน ตฺวํ ราชญฺญ ทิวา เสยฺยํ อุปคโต สุปินกํ ปสฺสิตา อารามรามเณยฺยกํ วนรามเณยฺยกํ ภูมิรามเณยฺยกํ โปกฺขรณีรามเณยฺยกนฺติ.\csromanspacer{} อภิชานามหํ โภ กสฺสป ทิวาเสยฺยํ อุปคโต สุปินกํ ปสฺสิตา อารามรามเณยฺยกํ วนรามเณยฺยกํ ภูมิรามเณยฺยกํ โปกฺขรณีรามเณยฺยกนฺติ.\csromanspacer{} รกฺขนฺติ ตํ ตมฺหิ สมเย \csromanfolio{266}ขุชฺชาปิ วามนกาปิ เวลาสิกาปิ\footnote{เจลาวิกาปิ (สฺยา), เกฬายิกาปิ (สี)} โกมาริกาปีติ.\csromanspacer{} เอวํ โภ กสฺสป รกฺขนฺติ มํ ตมฺหิ สมเย ขุชฺชาปิ วามนกาปิ เวลาสิกาปิ1 โกมาริกาปีติ.\csromanspacer{} อปินุ ตา ตุยฺหํ ชีวํ ปสฺสนฺติ ปวิสนฺตํ วา นิกฺขมนฺตํ วาติ.\csromanspacer{} โน หิทํ โภ กสฺสป.\csromanspacer{} ตาหิ นาม ราชญฺญ ตุยฺหํ ชีวนฺตสฺส ชีวนฺติโย ชีวํ น ปสฺสิสฺสนฺติ ปวิสนฺตํ วา นิกฺขมนฺตํ วา.\csromanspacer{} กิํ ปน ตฺวํ กาลงฺกตสฺส ชีวํ ปสฺสิสฺสสิ ปวิสนฺตํ วา นิกฺขมนฺตํ วา.\csromanspacer{} อิมินาปิ โข เต ราชญฺญ ปริยาเยน เอวํ โหตุ “อิติปิ อตฺถิ ปโร โลโก, อตฺถิ สตฺตา โอปปาติกา, อตฺถิ สุกตทุกฺกฏานํ กมฺมานํ ผลํ วิปาโก”ติ.}

\proseitem{423}{กิญฺจาปิ ภวํ กสฺสโป เอวมาห.\csromanspacer{} อถ โข เอวํ เม เอตฺถ โหติ “อิติปิ นตฺถิ ปโร โลโก, นตฺถิ สตฺตา โอปปาติกา, นตฺถิ สุกตทุกฺกฏานํ กมฺมานํ ผลํ วิปาโก”ติ.\csromanspacer{} อตฺถิ ปน ราชญฺญ ปริยาโย ฯเปฯ.\csromanspacer{} อตฺถิ โภ กสฺสป ปริยาโย ฯเปฯ.\csromanspacer{} ยถา กถํ วิย ราชญฺญาติ.\csromanspacer{} อิธ เม โภ กสฺสป ปุริสา โจรํ อาคุจาริํ คเหตฺวา ทสฺเสนฺติ “อยํ เต ภนฺเต โจโร อาคุจารี, อิมสฺส ยํ อิจฺฉสิ, ตํ ทณฺฑํ ปเณหี”ติ.\csromanspacer{} ตฺยาหํ เอวํ วทามิ “เตน หิ โภ อิมํ ปุริสํ ชีวนฺตํเยว ตุลาย ตุเลตฺวา ชิยาย อนสฺสาสกํ มาเรตฺวา ปุนเทว ตุลาย ตุเลถา”ติ.\csromanspacer{} เต เม “สาธู”ติ ปฏิสฺสุตฺวา ตํ ปุริสํ ชีวนฺตํเยว ตุลาย ตุเลตฺวา ชิยาย อนสฺสาสกํ มาเรตฺวา ปุนเทว ตุลาย ตุเลนฺติ.\csromanspacer{} ยทา โส ชีวติ, ตทา ลหุตโร จ โหติ มุทุตโร จ กมฺมญฺญตโร จ.\csromanspacer{} ยทา ปน โส กาลงฺกโต โหติ, ตทา ครุตโร จ โหติ ปตฺถินฺนตโร จ อกมฺมญฺญตโร จ.\csromanspacer{} อยมฺปิ โข โภ กสฺสป ปริยาโย, เยน เม ปริยาเยน เอวํ โหติ “อิติปิ นตฺถิ ปโร โลโก, นตฺถิ สตฺตา โอปปาติกา, นตฺถิ สุกตทุกฺกฏานํ กมฺมานํ ผลํ วิปาโก”ติ.}

\vspace{\tipitakaparskip}
\csromancenter{\textbf{(8) สนฺตตฺต-อโยคุฬ-อุปมา}}
\proseitem{424}{เตน หิ ราชญฺญ อุปมํ เต กริสฺสามิ.\csromanspacer{} อุปมาย มิเธกจฺเจ วิญฺญู ปุริสา ภาสิตสฺส อตฺถํ อาชานนฺติ.\csromanspacer{} เสยฺยถาปิ ราชญฺญ ปุริโส ทิวสํ}

\vspace{\tipitakaparskip}
\csromancenter{ปายาสิสุตฺต สนฺตตฺตํ อโยคุฬํ อาทิตฺตํ สมฺปชฺชลิตํ สโชติภูตํ ตุลาย ตุเลยฺย.\csromanspacer{} ตเมนํ อปเรน สมเยน สีตํ นิพฺพุตํ ตุลาย ตุเลยฺย.\csromanspacer{} กทา นุ โข โส อโยคุโฬ ลหุตโร วา โหติ มุทุตโร วา กมฺมญฺญตโร วา, ยทา วา อาทิตฺโต สมฺปชฺชลิโต สโชติภูโต, ยทา วา สีโต นิพฺพุโตติ.\csromanspacer{} ยทา โส โภ กสฺสป อโยคุโฬ เตโชสหคโต จ โหติ วาโยสหคโต จ อาทิตฺโต สมฺปชฺชลิโต สโชติภูโต, ตทา ลหุตโร จ โหติ มุทุตโร จ กมฺมญฺญตโร จ.\csromanspacer{} ยทา ปน โส อโยคุโฬ เนว เตโชสหคโต โหติ น วาโยสหคโต สีโต นิพฺพุโต, ตทา ครุตโร จ โหติ ปตฺถินฺนตโร จ อกมฺมญฺญตโร จาติ.\csromanspacer{} เอวเมว โข ราชญฺญ ยทายํ กาโย อายุสหคโต จ โหติ อุสฺมาสหคโต จ วิญฺญาณสหคโต จ, ตทา ลหุตโร จ โหติ มุทุตโร จ กมฺมญฺญตโร จ.\csromanspacer{} ยทา ปนายํ กาโย เนว อายุสหคโต โหติ น อุสฺมาสหคโต น วิญฺญาณสหคโต, ตทา ครุตโร จ โหติ ปตฺถินฺนตโร จ อกมฺมญฺญตโร จ.\csromanspacer{} อิมินาปิ โข เต ราชญฺญ ปริยาเยน เอวํ โหตุ “อิติปิ อตฺถิ ปโร โลโก, อตฺถิ สตฺตา โอปปาติกา, อตฺถิ สุกตทุกฺกฏานํ กมฺมานํ ผลํ วิปาโก”ติ.}
\csromanmatikaanchor{mtk.147}
\csromantocmarkref{section}{สงฺขธม-อุปมา}{mtk.147}
\bookmark[dest={mtk.147},level=1]{สงฺขธม-อุปมา}
\proseitem{425}{\csromanfolio{267}กิญฺจาปิ ภวํ กสฺสโป เอวมาห.\csromanspacer{} อถ โข เอวํ เม เอตฺถ โหติ.\csromanspacer{} อิติปิ นตฺถิ ปโร โลโก, นตฺถิ สตฺตา โอปปาติกา, นตฺถิ สุกตทุกฺกฏานํ กมฺมานํ ผลํ วิปาโกติ.\csromanspacer{} อตฺถิ ปน ราชญฺญ ปริยาโย ฯเปฯ.\csromanspacer{} อตฺถิ โภ กสฺสป ปริยาโย ฯเปฯ.\csromanspacer{} ยถา กถํ วิย ราชญฺญาติ.\csromanspacer{} อิธ เม โภ กสฺสป ปุริสา โจรํ อาคุจาริํ คเหตฺวา ทสฺเสนฺติ “อยํ เต ภนฺเต โจโร อาคุจารี, อิมสฺส ยํ อิจฺฉสิ, ตํ ทณฺฑํ ปเณหี”ติ.\csromanspacer{} ตฺยาหํ เอวํ วทามิ “เตน หิ โภ อิมํ ปุริสํ อนุปหจฺจ ฉวิญฺจ จมฺมญฺจ มํสญฺจ นฺหารุํ จ อฏฺฐิญฺจ อฏฺฐิมิญฺชญฺจ ชีวิตา โวโรเปถ, อปฺเปว นามสฺส ชีวํ นิกฺขมนฺตํ ปสฺเสยฺยามา”ติ.\csromanspacer{} เต เม “สาธู”ติ ปฏิสฺสุตฺวา ตํ ปุริสํ อนุปหจฺจ ฉวิญฺจ ฯเปฯ ชีวิตา โวโรเปนฺติ.\csromanspacer{} ยทา โส อามโต โหติ.\csromanspacer{} ตฺยาหํ เอวํ วทามิ “เตน หิ โภ อิมํ ปุริสํ อุตฺตานํ นิปาเตถ, อปฺเปว นามสฺส ชีวํ นิกฺขมนฺตํ ปสฺเสยฺยามา”ติ.\csromanspacer{} เต ตํ ปุริสํ อุตฺตานํ นิปาเตนฺติ.\csromanspacer{} เนวสฺส มยํ ชีวํ นิกฺขมนฺตํ ปสฺสาม.\csromanspacer{} ตฺยาหํ เอวํ วทามิ “เตน หิ โภ อิมํ ปุริสํ อวกุชฺชํ นิปาเตถ, ปสฺเสน นิปาเตถ, ทุติเยน \csromanfolio{268}ปสฺเสน นิปาเตถ, อุทฺธํ ฐเปถ, โอมุทฺธกํ ฐเปถ, ปาณินา อาโกเฏถ, เลฑฺฑุนา อาโกเฏถ, ทณฺเฑน อาโกเฏถ, สตฺเถน อาโกเฏถ, โอธุนาถ สนฺธุนาถ นิทฺธุนาถ, อปฺเปว นามสฺส ชีวํ นิกฺขมนฺตํ ปสฺเสยฺยามา”ติ.\csromanspacer{} เต ตํ ปุริสํ โอธุนนฺติ สนฺธุนนฺติ นิทฺธุนนฺติ.\csromanspacer{} เนวสฺส มยํ ชีวํ นิกฺขมนฺตํ ปสฺสาม.\csromanspacer{} ตสฺส ตเทว จกฺขุ โหติ เต รูปา, ตญฺจายตนํ นปฺปฏิสํเวเทติ.\csromanspacer{} ตเทว โสตํ โหติ เต สทฺทา, ตญฺจายตนํ นปฺปฏิสํเวเทติ.\csromanspacer{} ตเทว ฆานํ โหติ เต คนฺธา, ตญฺจายตนํ นปฺปฏิสํเวเทติ.\csromanspacer{} สาว ชิวฺหา โหติ เต รสา, ตญฺจายตนํ นปฺปฏิสํเวเทติ.\csromanspacer{} เสฺวว กาโย โหติ เต โผฏฺฐพฺพา, ตญฺจายตนํ นปฺปฏิสํเวเทติ.\csromanspacer{} อยมฺปิ โข โภ กสฺสป ปริยาโย, เยน เม ปริยาเยน เอวํ โหติ “อิติปิ นตฺถิ ปโร โลโก, นตฺถิ สตฺตา โอปปาติกา, นตฺถิ สุกตทุกฺกฏานํ กมฺมานํ ผลํ วิปาโก”ติ.}

\vspace{\tipitakaparskip}
\csromancenter{\textbf{(9) สงฺขธม-อุปมา}}
\proseitem{426}{เตน หิ ราชญฺญ อุปมํ เต กริสฺสามิ.\csromanspacer{} อุปมาย มิเธกจฺเจ วิญฺญู ปุริสา ภาสิตสฺส อตฺถํ อาชานนฺติ.\csromanspacer{} ภูตปุพฺพํ ราชญฺญ อญฺญตโร สงฺขธโม สงฺขํ อาทาย ปจฺจนฺติมํ ชนปทํ อคมาสิ.\csromanspacer{} โส เยน อญฺญตโร คาโม เตนุปสงฺกมิ, อุปสงฺกมิตฺวา มชฺเฌ คามสฺส ฐิโต ติกฺขตฺตุํ สงฺขํ อุปลาเปตฺวา สงฺขํ ภูมิยํ นิกฺขิปิตฺวา เอกมนฺตํ นิสีทิ.\csromanspacer{} อถ โข ราชญฺญ เตสํ ปจฺจนฺตชนปทานํ\footnote{ปจฺจนฺตชานํ (สี)} มนุสฺสานํ เอตทโหสิ “อมฺโภ กสฺส นุ โข\footnote{เอตทโหสิ “กิสฺส นุโข (อิ)} เอโส สทฺโท เอวํรชนีโย เอวํกมนีโย เอวํมทนีโย เอวํพนฺธนีโย เอวํมุจฺฉนีโย”ติ สนฺนิปติตฺวา ตํ สงฺขธมํ เอตทโวจุํ “อมฺโภ กสฺส นุ โข เอโส สทฺโท เอวํรชนีโย เอวํกมนีโย เอวํมทนีโย เอวํพนฺธนีโย เอวํมุจฺฉนีโยติ.\csromanspacer{} เอโส โข โภ สงฺโข นาม, ยสฺเสโส สทฺโท เอวํรชนีโย เอวํกมนีโย เอวํมทนีโย เอวํพนฺธนีโย เอวํมุจฺฉนีโยติ.\csromanspacer{} เต ตํ สงฺขํ อุตฺตานํ นิปาเตสุํ “วเทหิ โภ สงฺข วเทหิ โภ สงฺขา”ติ.\csromanspacer{} เนว โส สงฺโข สทฺทมกาสิ.\csromanspacer{} เต ตํ สงฺขํ อวกุชฺชํ นิปาเตสุํ, ปสฺเสน นิปาเตสุํ, ทุติเยน ปสฺเสน นิปาเตสุํ, อุทฺธํ ฐเปสุํ, โอมุทฺธกํ ฐเปสุํ, ปาณินา อาโกเฏสุํ, ลฑฺฑุนา อาโกเฏสุํ, ทณฺเฑน อาโกเฏสุํ, สตฺเถน อาโกเฏสุํ,}

\vspace{\tipitakaparskip}
\csromancenter{ปายาสิสุตฺต โอธุนิํสุ สนฺธุนิํสุ นิทฺธุนิํสุ “วเทหิ โภ สงฺข วเทหิ โภ สงฺขา”ติ.\csromanspacer{} เนว โส สงฺโข สทฺทมกาสิ.}
\prose{\csromanfolio{269}อถ โข ราชญฺญ ตสฺส สงฺขธมสฺส เอตทโหสิ “ยาว พาลา อิเม ปจฺจนฺตชนปทามนุสฺสา, กถญฺหิ นาม อโยนิโส สงฺขสทฺทํ คเวสิสฺสนฺตี”ติ.\csromanspacer{} เตสํ เปกฺขมานานํ สงฺขํ คเหตฺวา ติกฺขตฺตุํ สงฺขํ อุปลาเปตฺวา สงฺขํ อาทาย ปกฺกามิ.\csromanspacer{} อถ โข ราชญฺญ เตสํ ปจฺจนฺตชนปทานํ มนุสฺสานํ เอตทโหสิ “ยทา กิร โภ อยํ สงฺโข นาม ปุริสสหคโต จ โหติ วายามสหคโต\footnote{วาโยสหคโต (สฺยา)} จ วายุสหคโต จ, ตทายํ สงฺโข สทฺทํ กโรติ.\csromanspacer{} ยทา ปนายํ สงฺโข เนว ปุริสสหคโต โหติ น วายามสหคโต น วายุสหคโต, นายํ สงฺโข สทฺทํ กโรตี”ติ.\csromanspacer{} เอวเมว โข ราชญฺญ ยทายํ กาโย อายุสหคโต จ โหติ อุสฺมาสหคโต จ วิญฺญาณสหคโต จ, ตทา อภิกฺกมติปิ ปฏิกฺกมติปิ ติฏฺฐติปิ นิสีทติปิ เสยฺยํปิ กปฺเปติ, จกฺขุนาปิ รูปํ ปสฺสติ, โสเตนปิ สทฺทํ สุณาติ, ฆาเนนปิ คนฺธํ ฆายติ, ชิวฺหายปิ รสํ สายติ, กาเยนปิ โผฏฺฐพฺพํ ผุสติ, มนสาปิ ธมฺมํ วิชานาติ.\csromanspacer{} ยทา ปนายํ กาโย เนว อายุสหคโต โหติ, น อุสฺมาสหคโต, น วิญฺญาณสหคโต, ตทา เนว อภิกฺกมติ น ปฏิกฺกมติ น ติฏฺฐติ น นิสีทติ น เสยฺยํ กปฺเปติ, จกฺขุนาปิ รูปํ น ปสฺสติ, โสเตนปิ สทฺทํ น สุณาติ, ฆาเนนปิ คนฺธํ น ฆายติ, ชิวฺหายปิ รสํ น สายติ, กาเยนปิ โผฏฺฐพฺพํ น ผุสติ, มนสาปิ ธมฺมํ น วิชานาติ.\csromanspacer{} อิมินาปิ โข เต ราชญฺญ ปริยาเยน เอวํ โหตุ “อิติปิ อตฺถิ ปโร โลโก, อตฺถิ สตฺตา โอปปาติกา, อตฺถิ สุกตทุกฺกฏานํ กมฺมานํ ผลํ วิปาโก”ติ\footnote{วิปาโกติ. ปฐมภาณวารํ (สฺยา)}.}

\csromanmatikaanchor{mtk.148}
\csromantocmarkref{section}{อคฺคิกชฏิล-อุปมา}{mtk.148}
\bookmark[dest={mtk.148},level=1]{อคฺคิกชฏิล-อุปมา}
\proseitem{427}{กิญฺจาปิ ภวํ กสฺสโป เอวมาห.\csromanspacer{} อถ โข เอวํ เม เอตฺถ โหติ “อิติปิ นตฺถิ ปโร โลโก, นตฺถิ สตฺตา โอปปาติกา, นตฺถิ สุกตทุกฺกฏานํ กมฺมานํ ผลํ วิปาโก”ติ.\csromanspacer{} อตฺถิ ปน ราชญฺญ ปริยาโย ฯเปฯ.\csromanspacer{} อตฺถิ โภ กสฺสป ปริยาโย ฯเปฯ.\csromanspacer{} ยถา กถํ วิย ราชญฺญาติ.\csromanspacer{} อิธ เม โภ กสฺสป ปุริสา โจรํ อาคุจาริํ คเหตฺวา ทสฺเสนฺติ “อยํ เต ภนฺเต โจโร อาคุจารี, อิมสฺส ยํ อิจฺฉสิ, ตํ ทณฺฑํ ปเณหี”ติ.\csromanspacer{} ตฺยาหํ เอวํ วทามิ “เตน หิ โภ อิมสฺส ปุริสสฺส ฉวิํ \csromanfolio{270}ฉินฺทถ, อปฺเปว นามสฺส ชีวํ ปสฺเสยฺยามา”ติ.\csromanspacer{} เต ตสฺส ปุริสสฺส ฉวิํ ฉินฺทนฺติ, เนวสฺส มยํ ชีวํ ปสฺสาม.\csromanspacer{} ตฺยาหํ เอวํ วทามิ “เตน หิ โภ อิมสฺส ปุริสสฺส จมฺมํ ฉินฺทถ, มํสํ ฉินฺทถ, นฺหารุํ ฉินฺทถ, อฏฺฐิํ ฉินฺทถ, อฏฺฐิมิญฺชํ ฉินฺทถ, อปฺเปว นามสฺส ชีวํ ปสฺเสยฺยามา”ติ.\csromanspacer{} เต ตสฺส ปุริสสฺส อฏฺฐิมิญฺชํ ฉินฺทนฺติ, เนวสฺส มยํ ชีวํ ปสฺเสยฺยาม.\csromanspacer{} อยํปิ โข โภ กสฺสป ปริยาโย, เยน เม ปริยาเยน เอวํ โหติ “อิติปิ นตฺถิ ปโร โลโก, นตฺถิ สตฺตา โอปปาติกา, นตฺถิ สุกตทุกฺกฏานํ กมฺมานํ ผลํ วิปาโก”ติ.}

\vspace{\tipitakaparskip}
\csromancenter{\textbf{(10) อคฺคิกชฏิล-อุปมา}}
\proseitem{428}{เตน หิ ราชญฺญ อุปมํ เต กริสฺสามิ.\csromanspacer{} อุปมาย มิเธกจฺเจ วิญฺญู ปุริสา ภาสิตสฺส อตฺถํ อาชานนฺติ.\csromanspacer{} ภูตปุพฺพํ ราชญฺญ อญฺญตโร อคฺคิโก ชฏิโล อรญฺญายตเน ปณฺณกุฏิยา สมฺมติ\footnote{วสติ (สี, อิ)}.\csromanspacer{} อถ โข ราชญฺญ อญฺญตโร ชนปเท สตฺโถ\footnote{สตฺโถ ชนปทปเทสา (สี), ชนปโท สตฺถวาโส (สฺยา), ชนปทปเทโส (อิ)} วุฏฺฐาสิ.\csromanspacer{} อถ โข โส สตฺโถ\footnote{สตฺถวาโส (สฺยา)} ตสฺส อคฺคิกสฺส ชฏิลสฺส อสฺสมสฺส สามนฺตา เอกรตฺติํ วสิตฺวา ปกฺกามิ.\csromanspacer{} อถ โข ราชญฺญ ตสฺส อคฺคิกสฺส ชฏิลสฺส เอตทโหสิ “ยํนูนาหํ เยน โส สตฺถวาโส เตนุปสงฺกเมยฺยํ, อปฺเปว นาเมตฺถ กิญฺจิ อุปกรณํ อธิคจฺเฉยฺยนฺ”ติ.\csromanspacer{} อถ โข โส อคฺคิโก ชฏิโล กาลสฺเสว วุฏฺฐาย เยน โส สตฺถวาโส เตนุปสงฺกมิ, อุปสงฺกมิตฺวา อทฺทส ตสฺมิํ สตฺถวาเส ทหรํ กุมารํ มนฺทํ อุตฺตานเสยฺยกํ ฉฑฺฑิตํ, ทิสฺวานสฺส เอตทโหสิ “น โข เม ตํ ปติรูปํ ยํ เม เปกฺขมานสฺส มนุสฺสภูโต กาลงฺกเรยฺย, ยํนูนาหํ อิมํ ทารกํ อสฺสมํ เนตฺวา อาปาเทยฺยํ โปเสยฺยํ วฑฺเฒยฺยนฺ”ติ.\csromanspacer{} อถ โข โส อคฺคิโก ชฏิโล ตํ ทารกํ อสฺสมํ เนตฺวา อาปาเทสิ โปเสสิ วฑฺเฒสิ.\csromanspacer{} ยทา โส ทารโก ทสวสฺสุทฺเทสิโก วา โหติ\footnote{อโหสิ (?)} ทฺวาทสวสฺสุทฺเทสิโก วา.\csromanspacer{} อถ โข ตสฺส อคฺคิกสฺส ชฏิลสฺส ชนปเท กิญฺจิเทว กรณียํ อุปฺปชฺชิ.\csromanspacer{} อถ โข โส อคฺคิโก ชฏิโล ตํ ทารกํ เอตทโวจ “อิจฺฉามหํ ตาต ชนปทํ\footnote{นครํ (ก)} คนฺตุํ, อคฺคิํ ตาต ปริจเรยฺยาสิ, มา จ เต อคฺคิ นิพฺพายิ, สเจ จ เต อคฺคิ นิพฺพาเยยฺย, อยํ วาสี อิมานิ กฏฺฐานิ อิทํ อรณิสหิตํ, อคฺคิํ นิพฺพตฺเตตฺวา}

\vspace{\tipitakaparskip}
\csromancenter{ปายาสิสุตฺต อคฺคิํ ปริจเรยฺยาสี”ติ.\csromanspacer{} อถ โข โส อคฺคิโก ชฏิโล ตํ ทารกํ เอวํ อนุสาสิตฺวา ชนปทํ อคมาสิ.\csromanspacer{} ตสฺส ขิฑฺฑาปสุตสฺส อคฺคิ นิพฺพายิ.}
\prose{\csromanfolio{271}อถ โข ตสฺส ทารกสฺส เอตทโหสิ “ปิตา โข มํ เอวํ อวจ ‘อคฺคิํ ตาต ปริจเรยฺยาสิ, มา จ เต อคฺคิ นิพฺพายิ, สเจ จ เต อคฺคิ นิพฺพาเยยฺย, อยํ วาสี อิมานิ กฏฺฐานิ อิทํ อรณิสหิตํ, อคฺคิํ นิพฺพตฺเตตฺวา อคฺคิํ ปริจเรยฺยาสี’ติ, ยํนูนาหํ อคฺคิํ นิพฺพตฺเตตฺวา อคฺคิํ ปริจเรยฺยนฺ”ติ.\csromanspacer{} อถ โข โส ทารโก อรณิสหิตํ วาสิยา ตจฺฉิ “อปฺเปว นาม อคฺคิํ อธิคจฺเฉยฺยนฺ”ติ.\csromanspacer{} เนว โส อคฺคิํ อธิคจฺฉิ.\csromanspacer{} อรณิสหิตํ ทฺวิธา ผาเลสิ, ติธา ผาเลสิ, จตุธา ผาเลสิ, ปญฺจธา ผาเลสิ, ทสธา ผาเลสิ, สตธา\footnote{วีสติธา (สฺยา)} ผาเลสิ, สกลิกํ สกลิกํ อกาสิ, สกลิกํ สกลิกํ กริตฺวา อุทุกฺขเล โกฏฺเฏสิ, อุทุกฺขเล โกฏฺเฏตฺวา มหาวาเต โอปุนิ\footnote{โอผุนิ (สฺยา, ก)} “อปฺเปว นาม อคฺคิํ อธิคจฺเฉยฺยนฺ”ติ.\csromanspacer{} เนว โส อคฺคิํ อธิคจฺฉิ.}

\csromanmatikaanchor{mtk.149}
\csromantocmarkref{section}{เทฺวสตฺถวาห-อุปมา}{mtk.149}
\bookmark[dest={mtk.149},level=1]{เทฺวสตฺถวาห-อุปมา}
\prose{อถ โข โส อคฺคิโก ชฏิโล ชนปเท ตํ กรณียํ ตีเรตฺวา เยน สโก อสฺสโม เตนุปสงฺกมิ, อุปสงฺกมิตฺวา ตํ ทารกํ เอตทโวจ “กจฺจิ เต ตาต อคฺคิ น นิพฺพุโต”ติ.\csromanspacer{} อิธ เม ตาต ขิฑฺฑาปสุตสฺส อคฺคิ นิพฺพายิ, ตสฺส เม เอตทโหสิ “ปิตา โข มํ เอวํ อวจ ‘อคฺคิํ ตาต ปริจเรยฺยาสิ, มา จ เต ตาต อคฺคิ นิพฺพายิ, สเจ จ เต อคฺคิ นิพฺพาเยยฺย, อยํ วาสี อิมานิ กฏฺฐานิ อิทํ อรณิสหิตํ, อคฺคิํ นิพฺพตฺเตตฺวา อคฺคิํ ปริจเรยฺยาสี’ติ.\csromanspacer{} ยํนูนาหํ อคฺคิํ นิพฺพตฺเตตฺวา อคฺคิํ ปริจเรยฺยนฺ”ติ.\csromanspacer{} อถ ขฺวาหํ ตาต อรณิสหิตํ วาสิยา ตจฺฉิํ “อปฺเปว นาม อคฺคิํ อธิคจฺเฉยฺยนฺ”ติ.\csromanspacer{} เนวาหํ อคฺคิํ อธิคจฺฉิํ.\csromanspacer{} อรณิสหิตํ ทฺวิธา ผาเลสิํ, ติธา ผาเลสิํ, จตุธา ผาเลสิํ, ปญฺจธา ผาเลสิํ, ทสธา ผาเลสิํ, สตธา ผาเลสิํ, สกลิกํ สกลิกํ อกาสิํ, สกลิกํ สกลิกํ กริตฺวา อุทุกฺขเล โกฏฺเฏสิํ, อุทุกฺขเล โกฏฺเฏตฺวา มหาวาเต โอปุนิํ “อปฺเปว นาม อคฺคิํ อธิคจฺเฉยฺยนฺ”ติ.\csromanspacer{} เนวาหํ อคฺคิํ อธิคจฺฉินฺติ.\csromanspacer{} อถ โข ตสฺส อคฺคิกสฺส ชฏิลสฺส เอตทโหสิ “ยาว พาโล อยํ ทารโก อพฺยตฺโต, กถํ หิ นาม อโยนิโส อคฺคิํ คเวสิสฺสตี”ติ.\csromanspacer{} ตสฺส เปกฺขมานสฺส อรณิสหิตํ คเหตฺวา อคฺคิํ นิพฺพตฺเตตฺวา ตํ ทารกํ \csromanfolio{272}เอตทโวจ “เอวํ โข ตาต อคฺคิ นิพฺพตฺเตตพฺโพ, น เตฺวว ยถา ตฺวํ พาโล อพฺยตฺโต อโยนิโส อคฺคิํ คเวสี”ติ.\csromanspacer{} เอวเมว โข ตฺวํ ราชญฺญ พาโล อพฺยตฺโต อโยนิโส ปรโลกํ คเวสิสฺสสิ.\csromanspacer{} ปฏินิสฺสชฺเชตํ ราชญฺญ ปาปกํ ทิฏฺฐิคตํ, ปฏินิสฺสชฺเชตํ ราชญฺญ ปาปกํ ทิฏฺฐิคตํ, มา เต อโหสิ ทีฆรตฺตํ อหิตาย ทุกฺขายาติ.}

\proseitem{429}{กิญฺจาปิ ภวํ กสฺสโป เอวมาห, อถ โข เนวาหํ สกฺโกมิ อิทํ ปาปกํ ทิฏฺฐิคตํ ปฏินิสฺสชฺชิตุํ.\csromanspacer{} ราชาปิ มํ ปเสนทิ โกสโล ชานาติ ติโรราชาโนปิ “ปายาสิ ราชญฺโญ เอวํวาที เอวํทิฏฺฐี ‘อิติปิ นตฺถิ ปโร โลโก, นตฺถิ สตฺตา โอปปาติกา, นตฺถิ สุกตทุกฺกฏานํ กมฺมานํ ผลํ วิปาโก’ติ”.\csromanspacer{} สจาหํ โภ กสฺสป อิทํ ปาปกํ ทิฏฺฐิคตํ ปฏินิสฺสชฺชิสฺสามิ, ภวิสฺสนฺติ เม วตฺตาโร “ยาว พาโล ปายาสิ ราชญฺโญ อพฺยตฺโต ทุคฺคหิตคาหี”ติ.\csromanspacer{} โกเปนปิ นํ หริสฺสามิ, มกฺเขนปิ นํ หริสฺสามิ, ปลาเสนปิ นํ หริสฺสามีติ.}

\vspace{\tipitakaparskip}
\csromancenter{\textbf{(11) เทฺวสตฺถวาห-อุปมา}}
\proseitem{430}{เตน หิ ราชญฺญ อุปมํ เต กริสฺสามิ, อุปมาย มิเธกจฺเจ วิญฺญู ปุริสา ภาสิตสฺส อตฺถํ อาชานนฺติ.\csromanspacer{} ภูตปุพฺพํ ราชญฺญ มหาสกฏสตฺโถ สกฏสหสฺสํ ปุรตฺถิมา ชนปทา ปจฺฉิมํ ชนปทํ อคมาสิ, โส เยน เยน คจฺฉิ, ขิปฺปํเยว ปริยาทิยติ ติณกฏฺโฐทกํ หริตกปณฺณํ.\csromanspacer{} ตสฺมิํ โข ปน สตฺเถ เทฺว สตฺถวาหา อเหสุํ เอโก ปญฺจนฺนํ สกฏสตานํ, เอโก ปญฺจนฺนํ สกฏสตานํ.\csromanspacer{} อถ โข เตสํ สตฺถวาหานํ เอตทโหสิ “อยํ โข มหาสกฏสตฺโถ สกฏสหสฺสํ, เต มยํ เยน เยน คจฺฉาม, ขิปฺปเมว ปริยาทิยติ ติณกฏฺโฐทกํ หริตกปณฺณํ.\csromanspacer{} ยํนูน มยํ อิมํ สตฺถํ ทฺวิธา วิภเชยฺยาม เอกโต ปญฺจ สกฏสตานิ เอกโต ปญฺจ สกฏสตานี”ติ.\csromanspacer{} เต ตํ สตฺถํ ทฺวิธา วิภชิํสุ\footnote{วิภเชสุํ (ก)} เอกโต ปญฺจ สกฏสตานิ เอกโต ปญฺจ สกฏสตานิ.\csromanspacer{} เอโก สตฺถวาโห พหุํ ติณญฺจ กฏฺฐญฺจ อุทกญฺจ อาโรเปตฺวา สตฺถํ ปยาเปสิ\footnote{ปายาเปสิ (สี, อิ)}.\csromanspacer{} ทฺวีหตีหปยาโต โข ปน โส สตฺโถ อทฺทส ปุริสํ กาฬํ โลหิตกฺขํ\footnote{โลหิตกฺขิํ (สฺยา)} สนฺนทฺธกลาปํ\footnote{อาสนฺนทฺธกลาปํ (สฺยา)} กุมุทมาลิํ อลฺลวตฺถํ อลฺลเกสํ กทฺทมมกฺขิเตหิ}

\vspace{\tipitakaparskip}
\csromancenter{ปายาสิสุตฺต จกฺเกหิ ภเทฺรน รเถน ปฏิปถํ อาคจฺฉนฺตํ, ทิสฺวา เอตทโวจ “กุโต โภ อาคจฺฉสี”ติ.\csromanspacer{} อมุกมฺหา ชนปทาติ. “กุหิํ คมิสฺสสี”ติ.\csromanspacer{} อมุกํ นาม ชนปทนฺติ. “กจฺจิ โภ ปุรโต กนฺตาเร มหาเมโฆ อภิปฺปวุฏฺโฐ”ติ.\csromanspacer{} เอวํ โภ ปุรโต กนฺตาเร มหาเมโฆ อภิปฺปวุฏฺโฐ, อาสิตฺโตทกานิ วฏุมานิ, พหุ ติณญฺจ กฏฺฐญฺจ อุทกญฺจ, ฉฑฺเฑถ โภ ปุราณานิ ติณานิ กฏฺฐานิ อุทกานิ, ลหุภาเรหิ สกเฏหิ สีฆํ สีฆํ คจฺฉถ, มา โยคฺคานิ กิลมิตฺถาติ.}
\prose{\csromanfolio{273}อถ โข โส สตฺถวาโห สตฺถิเก อามนฺเตสิ “อยํ โภ ปุริโส เอวมาห ‘ปุรโต กนฺตาเร มหาเมโฆ อภิปฺปวุฏฺโฐ, อาสิตฺโตทกานิ วฏุมานิ, พหุ ติณญฺจ กฏฺฐญฺจ อุทกญฺจ, ฉฑฺเฑถ โภ ปุราณานิ ติณานิ กฏฺฐานิ อุทกานิ, ลหุภาเรหิ สกเฏหิ สีฆํ สีฆํ คจฺฉถ, มา โยคฺคานิ กิลมิตฺถา’ติ”.\csromanspacer{} ฉฑฺเฑถ โภ ปุราณานิ ติณานิ กฏฺฐานิ อุทกานิ, ลหุภาเรหิ สกเฏหิ สตฺถํ ปยาเปถาติ. “เอวํ โภ”ติ โข เต สตฺถิกา ตสฺส สตฺถวาหสฺส ปฏิสฺสุตฺวา ฉฑฺเฑตฺวา ปุราณานิ ติณานิ กฏฺฐานิ อุทกานิ ลหุภาเรหิ สกเฏหิ สตฺถํ ปยาเปสุํ.\csromanspacer{} เต ปฐเมปิ สตฺถวาเส น อทฺทสํสุ ติณํ วา กฏฺฐํ วา อุทกํ วา.\csromanspacer{} ทุติเยปิ สตฺถวาเส.\csromanspacer{} ตติเยปิ สตฺถวาเส.\csromanspacer{} จตุตฺเถปิ สตฺถวาเส.\csromanspacer{} ปญฺจเมปิ สตฺถวาเส.\csromanspacer{} ฉฏฺเฐปิ สตฺถวาเส.\csromanspacer{} สตฺตเมปิ สตฺถวาเส น อทฺทสํสุ ติณํ วา กฏฺฐํ วา อุทกํ วา, สพฺเพว อนยพฺยสนํ อาปชฺชิํสุ.\csromanspacer{} เย จ ตสฺมิํ สตฺเถ อเหสุํ มนุสฺสา วา ปสู วา, สพฺเพ โส ยกฺโข อมนุสฺโส ภกฺเขสิ.\csromanspacer{} อฏฺฐิกาเนว เสสานิ.}

\prose{ยทา อญฺญาสิ ทุติโย สตฺถวาโห “พหุนิกฺขนฺโต โข โภ ทานิ โส สตฺโถ”ติ.\csromanspacer{} พหุํ ติณญฺจ กฏฺฐญฺจ อุทกญฺจ อาโรเปตฺวา สตฺถํ ปยาเปสิ.\csromanspacer{} ทฺวีหตีหปยาโต โข ปน โส สตฺโถ อทฺทส ปุริสํ กาฬํ โลหิตกฺขํ สนฺนทฺธกลาปํ กุมุทมาลิํ อลฺลวตฺถํ อลฺลเกสํ กทฺทมมกฺขิเตหิ จกฺเกหิ ภเทฺรน รเถน ปฏิปถํ อาคจฺฉนฺตํ, ทิสฺวา เอตทโวจ “กุโต โภ อาคจฺฉสี”ติ.\csromanspacer{} อมุกมฺหา ชนปทาติ. “กุหิํ คมิสฺสสี”ติ.\csromanspacer{} อมุกํ นาม ชนปทนฺติ. “กจฺจิ โภ ปุรโต กนฺตาเร มหาเมโฆ อภิปฺปวุฏฺโฐ”ติ.\csromanspacer{} เอวํ โภ ปุรโต กนฺตาเร มหาเมโฆ อภิปฺปวุฏฺโฐ, อาสิตฺโตทกานิ วฏุมานิ, พหุ ติณญฺจ กฏฺฐญฺจ อุทกญฺจ, \csromanfolio{274}ฉฑฺเฑถ โภ ปุราณานิ ติณานิ กฏฺฐานิ อุทกานิ, ลหุภาเรหิ สกเฏหิ สีฆํ สีฆํ คจฺฉถ, มา โยคฺคานิ กิลมิตฺถาติ.}

\prose{อถ โข โส สตฺถวาโห สตฺถิเก อามนฺเตสิ “อยํ โภ ปุริโส เอวมาห ‘ปุรโต กนฺตาเร มหาเมโฆ อภิปฺปวุฏฺโฐ, อาสิตฺโตทกานิ วฏุมานิ, พหุ ติณญฺจ กฏฺฐญฺจ อุทกญฺจ, ฉฑฺเฑถ โภ ปุราณานิ ติณานิ กฏฺฐานิ อุทกานิ, ลหุภาเรหิ สกเฏหิ สีฆํ สีฆํ คจฺฉถ, มา โยคฺคานิ กิลมิตฺถา’ติ.\csromanspacer{} อยํ โภ ปุริโส เนว อมฺหากํ มิตฺโต, น ญาติสาโลหิโต, กถํ มยํ อิมสฺส สทฺธาย คมิสฺสาม.\csromanspacer{} น โว ฉฑฺเฑตพฺพานิ ปุราณานิ ติณานิ กฏฺฐานิ อุทกานิ, ยถาภเตน ภณฺเฑน สตฺถํ ปยาเปถ.\csromanspacer{} น โน ปุราณํ ฉฑฺเฑสฺสามา”ติ. “เอวํ โภ”ติ โข เต สตฺถิกา ตสฺส สตฺถวาหสฺส ปฏิสฺสุตฺวา ยถาภเตน ภณฺเฑน สตฺถํ ปยาเปสุํ.\csromanspacer{} เต ปฐเมปิ สตฺถวาเส น อทฺทสํสุ ติณํ วา กฏฺฐํ วา อุทกํ วา.\csromanspacer{} ทุติเยปิ สตฺถวาเส.\csromanspacer{} ตติเยปิ สตฺถวาเส.\csromanspacer{} จตุตฺเถปิ สตฺถวาเส.\csromanspacer{} ปญฺจเมปิ สตฺถวาเส.\csromanspacer{} ฉฏฺเฐปิ สตฺถวาเส.\csromanspacer{} สตฺตเมปิ สตฺถวาเส น อทฺทสํสุ ติณํ วา กฏฺฐํ วา อุทกํ วา.\csromanspacer{} ตญฺจ สตฺถํ อทฺทสํสุ อนยพฺยสนํ อาปนฺนํ.\csromanspacer{} เย จ ตสฺมิํ สตฺเถปิ อเหสุํ มนุสฺสา วา ปสู วา, เตสญฺจ อฏฺฐิกาเนว อทฺทสํสุ เตน ยกฺเขน อมนุสฺเสน ภกฺขิตานํ.}

\prose{อถ โข โส สตฺถวาโห สตฺถิเก อามนฺเตสิ “อยํ โข โภ สตฺโถ อนยพฺยสนํ อาปนฺโน, ยถา ตํ เตน พาเลน สตฺถวาเหน ปริณายเกน.\csromanspacer{} เตน หิ โภ ยานมฺหากํ สตฺเถ อปฺปสารานิ ปณิยานิ, ตานิ ฉฑฺเฑตฺวา, ยานิ อิมสฺมิํ สตฺเถ มหาสารานิ ปณิยานิ, ตานิ อาทิยถา”ติ. “เอวํ โภ”ติ โข เต สตฺถิกา ตสฺส สตฺถวาหสฺส ปฏิสฺสุตฺวา ยานิ สกสฺมิํ สตฺเถ อปฺปสารานิ ปณิยานิ, ตานิ ฉฑฺเฑตฺวา ยานิ ตสฺมิํ สตฺเถ มหาสารานิ ปณิยานิ, ตานิ อาทิยิตฺวา โสตฺถินา ตํ กนฺตารํ นิตฺถริํสุ, ยถา ตํ ปณฺฑิเตน สตฺถวาเหน ปริณายเกน.\csromanspacer{} เอวเมว โข ตฺวํ ราชญฺญ พาโล อพฺยตฺโต อนยพฺยสนํ อาปชฺชิสฺสสิ อโยนิโส ปรโลกํ คเวสนฺโต เสยฺยถาปิ โส ปุริโม สตฺถวาโห.\csromanspacer{} เยปิ ตว1โสตพฺพํ สทฺธาตพฺพํ\footnote{สทฺทหาตพฺพํ (อิ, ก)} มญฺญิสฺสนฺติ, เตปิ อนยพฺยสนํ อาปชฺชิสฺสนฺติ เสยฺยถาปิ เต สตฺถิกา.\csromanspacer{} ปฏินิสฺสชฺเชตํ}

\vspace{\tipitakaparskip}
\csromancenter{ปายาสิสุตฺต ราชญฺญ ปาปกํ ทิฏฺฐิคตํ, ปฏินิสฺสชฺเชตํ ราชญฺญ ปาปกํ ทิฏฺฐิคตํ, มา เต อโหสิ ทีฆรตฺตํ อหิตาย ทุกฺขายาติ.}
\csromanmatikaanchor{mtk.150}
\csromantocmarkref{section}{คูถภาริก-อุปมา}{mtk.150}
\bookmark[dest={mtk.150},level=1]{คูถภาริก-อุปมา}
\proseitem{431}{\csromanfolio{275}กิญฺจาปิ ภวํ กสฺสโป เอวมาห, อถ โข เนวาหํ สกฺโกมิ อิทํ ปาปกํ ทิฏฺฐิคตํ ปฏินิสฺสชฺชิตุํ.\csromanspacer{} ราชาปิ มํ ปเสนทิ โกสโล ชานาติ ติโรราชาโนปิ “ปายาสิ ราชญฺโญ เอวํวาที เอวํทิฏฺฐี ‘อิติปิ นตฺถิ ปโร โลโก ฯเปฯ วิปาโก’ติ”.\csromanspacer{} สจาหํ โภ กสฺสป อิทํ ปาปกํ ทิฏฺฐิคตํ ปฏินิสฺสชฺชิสฺสามิ, ภวิสฺสนฺติ เม วตฺตาโร “ยาว พาโล ปายาสิ ราชญฺโญ อพฺยตฺโต ทุคฺคหิตคาหี”ติ.\csromanspacer{} โกเปนปิ นํ หริสฺสามิ, มกฺเขนปิ นํ หริสฺสามิ, ปลาเสนปิ นํ หริสฺสามีติ.}

\vspace{\tipitakaparskip}
\csromancenter{\textbf{(12) คูถภาริก-อุปมา}}
\proseitem{432}{เตน หิ ราชญฺญ อุปมํ เต กริสฺสามิ, อุปมาย มิเธกจฺเจ วิญฺญู ปุริสา ภาสิตสฺส อตฺถํ อาชานนฺติ.\csromanspacer{} ภูตปุพฺพํ ราชญฺญ อญฺญตโร สูกรโปสโก ปุริโส สกมฺหา คามา อญฺญํ คามํ อคมาสิ, ตตฺถ อทฺทส ปหูตํ สุกฺขคูถํ ฉฑฺฑิตํ.\csromanspacer{} ทิสฺวานสฺส เอตทโหสิ “อยํ โข ปหูโต สุกฺขคูโถ ฉฑฺฑิโต, มม จ สูกรภตฺตํ\footnote{สูกรานํ ภกฺโข (สฺยา)}, ยํนูนาหํ อิโต สุกฺขคูถํ หเรยฺยนฺ”ติ.\csromanspacer{} โส อุตฺตราสงฺคํ ปตฺถริตฺวา ปหูตํ สุกฺขคูถํ อากิริตฺวา ภณฺฑิกํ พนฺธิตฺวา สีเส อุพฺพาเหตฺวา\footnote{อุจฺจาโรเปตฺวา (ก-สี, ก)} อคมาสิ.\csromanspacer{} ตสฺส อนฺตรามคฺเค มหา-อกาลเมโฆ ปาวสฺสิ.\csromanspacer{} โส อุคฺฆรนฺตํ ปคฺฆรนฺตํ ยาว อคฺคนขา คูเถน มกฺขิโต คูถภารํ อาทาย อคมาสิ.\csromanspacer{} ตเมนํ มนุสฺสา ทิสฺวา เอวมาหํสุ “กจฺจิโน ตฺวํ ภเณ อุมฺมตฺโต กจฺจิ วิเจโต, กถํ หิ นาม อุคฺฆรนฺตํ ปคฺฆรนฺตํ ยาว อคฺคนขา คูเถน มกฺขิโต คูถภารํ หริสฺสสี”ติ.\csromanspacer{} ตุเมฺห เขฺวตฺถ ภเณ อุมฺมตฺตา ตุเมฺห วิเจตา, ตถา หิ ปน เม สูกรภตฺตนฺติ.\csromanspacer{} เอวเมว โข ตฺวํ ราชญฺญ คูถภาริกูปโม\footnote{คูถหาริกูปโม (สี, อิ)} มญฺเญ ปฏิภาสิ.\csromanspacer{} ปฏินิสฺสชฺเชตํ ราชญฺญ ปาปกํ ทิฏฺฐิคตํ, ปฏินิสฺสชฺเชตํ ราชญฺญ ปาปกํ ทิฏฺฐิคตํ, มา เต อโหสิ ทีฆรตฺตํ อหิตาย ทุกฺขายาติ.}

\csromanmatikaanchor{mtk.151}
\csromantocmarkref{section}{อกฺขธุตฺตก-อุปมา}{mtk.151}
\bookmark[dest={mtk.151},level=1]{อกฺขธุตฺตก-อุปมา}
\proseitem{433}{กิญฺจาปิ ภวํ กสฺสโป เอวมาห, อถ โข เนวาหํ สกฺโกมิ อิทํ ปาปกํ ทิฏฺฐิคตํ ปฏินิสฺสชฺชิตุํ.\csromanspacer{} ราชาปิ มํ ปเสนทิ โกสโล ชานาติ ติโรราชาโนปิ “ปายาสิ ราชญฺโญ เอวํวาที เอวํทิฏฺฐี ‘อิติปิ \csromanfolio{276}นตฺถิ ปโร โลโก ฯเปฯ วิปาโก’ติ”.\csromanspacer{} สจาหํ โภ กสฺสป อิทํ ปาปกํ ทิฏฺฐิคตํ ปฏินิสฺสชฺชิสฺสามิ, ภวิสฺสนฺติ เม วตฺตาโร “ยาว พาโล ปายาสิ ราชญฺโญ อพฺยตฺโต ทุคฺคหิตคาหี”ติ.\csromanspacer{} โกเปนปิ นํ หริสฺสามิ, มกฺเขนปิ นํ หริสฺสามิ, ปลาเสนปิ นํ หริสฺสามีติ.}

\vspace{\tipitakaparskip}
\csromancenter{\textbf{(13) อกฺขธุตฺตก อุปมา}}
\proseitem{434}{เตน หิ ราชญฺญ อุปมํ เต กริสฺสามิ, อุปมาย มิเธกจฺเจ วิญฺญู ปุริสา ภาสิตสฺส อตฺถํ อาชานนฺติ.\csromanspacer{} ภูตปุพฺพํ ราชญฺญ เทฺว อกฺขธุตฺตา อกฺเขหิ ทิพฺพิํสุ.\csromanspacer{} เอโก อกฺขธุตฺโต อาคตาคตํ กลิํ คิลติ.\csromanspacer{} อทฺทสา โข ทุติโย อกฺขธุตฺโต ตํ อกฺขธุตฺตํ อาคตาคตํ กลิํ คิลนฺตํ, ทิสฺวา ตํ อกฺขธุตฺตํ เอตทโวจ “ตฺวํ โข สมฺม เอกนฺติเกน ชินาสิ, เทหิ เม สมฺม อกฺเข ปโชหิสฺสามี”ติ. “เอวํ สมฺมา”ติ โข โส อกฺขธุตฺโต ตสฺส อกฺขธุตฺตสฺส อกฺเข ปาทาสิ.\csromanspacer{} อถ โข โส อกฺขธุตฺโต อกฺเข วิเสน ปริภาเวตฺวา ตํ อกฺขธุตฺตํ เอตทโวจ “เอหิ โข สมฺม อกฺเขหิ ทิพฺพิสฺสามา”ติ. “เอวํ สมฺมา”ติ โข โส อกฺขธุตฺโต ตสฺส อกฺขธุตฺตสฺส ปจฺจสฺโสสิ.\csromanspacer{} ทุติยมฺปิ โข เต อกฺขธุตฺตา อกฺเขหิ ทิพฺพิํสุ.\csromanspacer{} ทุติยมฺปิ โข โส อกฺขธุตฺโต อาคตาคตํ กลิํ คิลติ.\csromanspacer{} อทฺทสา โข ทุติโย อกฺขธุตฺโต ตํ อกฺขธุตฺตํ ทุติยมฺปิ อาคตาคตํ กลิํ คิลนฺตํ, ทิสฺวา ตํ อกฺขธุตฺตํ เอตทโวจ\csromandash{}}

\csromangathameasure{“ลิตฺตํ ปรเมน เตชสา, คิลมกฺขํ ปุริโส น พุชฺฌติ. \\ คิล เร คิล ปาปธุตฺตก, ปจฺฉา เต กฏุกํ ภวิสฺสตี”ติ.}
\csromangathagroup{\csromangathastanza{“ลิตฺตํ ปรเมน เตชสา, คิลมกฺขํ ปุริโส น พุชฺฌติ. \\ คิล เร คิล ปาปธุตฺตก\footnote{คิลิ เร ปาปธุตฺตก (ก)}, ปจฺฉา เต กฏุกํ ภวิสฺสตี”ติ.}}

\prose{เอวเมว โข ตฺวํ ราชญฺญ อกฺขธุตฺตกูปโม มญฺเญ ปฏิภาสิ.\csromanspacer{} ปฏินิสฺสชฺเชตํ ราชญฺญ ปาปกํ ทิฏฺฐิคตํ, ปฏินิสฺสชฺเชตํ ราชญฺญ ปาปกํ ทิฏฺฐิคตํ, มา เต อโหสิ ทีฆรตฺตํ อหิตาย ทุกฺขายาติ.}

\proseitem{435}{กิญฺจาปิ ภวํ กสฺสโป เอวมาห, อถ โข เนวาหํ สกฺโกมิ อิทํ ปาปกํ ทิฏฺฐิคตํ ปฏินิสฺสชฺชิตุํ.\csromanspacer{} ราชาปิ มํ ปเสนทิ โกสโล ชานาติ ติโรราชาโนปิ “ปายาสิ ราชญฺโญ เอวํวาที เอวํทิฏฺฐี ‘อิติปิ นตฺถิ ปโร โลโก ฯเปฯ วิปาโก’ติ”.\csromanspacer{} สจาหํ โภ กสฺสป อิทํ ปาปกํ}

\vspace{\tipitakaparskip}
\csromancenter{ปายาสิสุตฺต ทิฏฺฐิคตํ ปฏินิสฺสชฺชิสฺสามิ, ภวิสฺสนฺติ เม วตฺตาโร “ยาว พาโล ปายาสิ ราชญฺโญ อพฺยตฺโต ทุคฺคหิตคาหี”ติ.\csromanspacer{} โกเปนปิ นํ หริสฺสามิ, มกฺเขนปิ นํ หริสฺสามิ, ปลาเสนปิ นํ หริสฺสามีติ.}
\csromancenter{\textbf{(14) สาณภาริก อุปมา}}
\csromanmatikaanchor{mtk.152}
\csromantocmarkref{section}{สาณภาริก-อุปมา}{mtk.152}
\bookmark[dest={mtk.152},level=1]{สาณภาริก-อุปมา}
\proseitem{436}{\csromanfolio{277}เตน หิ ราชญฺญ อุปมํ เต กริสฺสามิ, อุปมาย มิเธกจฺเจ วิญฺญู ปุริสา ภาสิตสฺส อตฺถํ อาชานนฺติ.\csromanspacer{} ภูตปุพฺพํ ราชญฺญ อญฺญตโร ชนปโท วุฏฺฐาสิ.\csromanspacer{} อถ โข สหายโก สหายกํ อามนฺเตสิ “อายาม สมฺม เยน โส ชนปโท เตนุปสงฺกมิสฺสาม, อปฺเปว นาเมตฺถ กิญฺจิ ธนํ อธิคจฺเฉยฺยามา”ติ. “เอวํ สมฺมา”ติ โข สหายโก สหายกสฺส ปจฺจสฺโสสิ.\csromanspacer{} เต เยน โส ชนปโท, เยน อญฺญตรํ คามปฏฺฏํ\footnote{คามปชฺชํ (สฺยา), คามปตฺตํ (สี)} เตนุปสงฺกมิํสุ, ตตฺถ อทฺทสํสุ ปหูตํ สาณํ ฉฑฺฑิตํ, ทิสฺวา สหายโก สหายกํ อามนฺเตสิ “อิทํ โข สมฺม ปหูตํ สาณํ ฉฑฺฑิตํ, เตน หิ สมฺม ตฺวญฺจ สาณภารํ พนฺธ, อหญฺจ สาณภารํ พนฺธิสฺสามิ, อุโภ สาณภารํ อาทาย คมิสฺสามา”ติ. “เอวํ สมฺมา”ติ โข สหายโก สหายกสฺส ปฏิสฺสุตฺวา สาณภารํ พนฺธิตฺวา เต อุโภ สาณภารํ อาทาย เยน อญฺญตรํ คามปฏฺฏํ เตนุปสงฺกมิํสุ.\csromanspacer{} ตตฺถ อทฺทสํสุ ปหูตํ สาณสุตฺตํ ฉฑฺฑิตํ, ทิสฺวา สหายโก สหายกํ อามนฺเตสิ “ยสฺส โข สมฺม อตฺถาย อิจฺเฉยฺยาม สาณํ, อิทํ ปหูตํ สาณสุตฺตํ ฉฑฺฑิตํ, เตน หิ สมฺม ตฺวญฺจ สาณภารํ ฉฑฺเฑหิ, อหญฺจ สาณภารํ ฉฑฺเฑสฺสามิ, อุโภ สาณสุตฺตภารํ อาทาย คมิสฺสามา”ติ.\csromanspacer{} อยํ โข เม สมฺม สาณภาโร ทูราภโต จ สุสนฺนทฺโธ จ, อลํ เม ตฺวํ ปชานาหีติ.\csromanspacer{} อถ โข โส สหายโก สาณภารํ ฉฑฺเฑตฺวา สาณสุตฺตภารํ อาทิยิ.}

\prose{เต เยน อญฺญตรํ คามปฏฺฏํ เตนุปสงฺกมิํสุ.\csromanspacer{} ตตฺถ อทฺทสํสุ ปหูตา สาณิโย ฉฑฺฑิตา, ทิสฺวา สหายโก สหายกํ อามนฺเตสิ “ยสฺส โข สมฺม อตฺถาย อิจฺเฉยฺยาม สาณํ วา สาณสุตฺตํ วา, อิมา ปหูตา สาณิโย ฉฑฺฑิตา, เตน หิ สมฺม ตฺวญฺจ สาณภารํ ฉฑฺเฑหิ, อหญฺจ สาณสุตฺตภารํ ฉฑฺเฑสฺสามิ, อุโภ สาณิภารํ อาทาย \csromanfolio{278}คมิสฺสามา”ติ.\csromanspacer{} อยํ โข เม สมฺม สาณภาโร ทูราภโต จ สุสนฺนทฺโธ จ, อลํ เม ตฺวํ ปชานาหีติ.\csromanspacer{} อถ โข โส สหายโก สาณสุตฺตภารํ ฉฑฺเฑตฺวา สาณิภารํ อาทิยิ.}

\prose{เต เยน อญฺญตรํ คามปฏฺฏํ เตนุปสงฺกมิํสุ.\csromanspacer{} ตตฺถ อทฺทสํสุ ปหูตํ โขมํ ฉฑฺฑิตํ, ทิสฺวา ฯเปฯ ปหูตํ โขมสุตฺตํ ฉฑฺฑิตํ, ทิสฺวา.\csromanspacer{} ปหูตํ โขมทุสฺสํ ฉฑฺฑิตํ, ทิสฺวา.\csromanspacer{} ปหูตํ กปฺปาสํ ฉฑฺฑิตํ, ทิสฺวา.\csromanspacer{} ปหูตํ กปฺปาสิกสุตฺตํ ฉฑฺฑิตํ, ทิสฺวา.\csromanspacer{} ปหูตํ กปฺปาสิกทุสฺสํ ฉฑฺฑิตํ, ทิสฺวา.\csromanspacer{} ปหูตํ อยํ\footnote{อยสํ (สฺยา)} ฉฑฺฑิตํ, ทิสฺวา.\csromanspacer{} ปหูตํ โลหํ ฉฑฺฑิตํ, ทิสฺวา.\csromanspacer{} ปหูตํ ติปุํ ฉฑฺฑิตํ, ทิสฺวา.\csromanspacer{} ปหูตํ สีสํ ฉฑฺฑิตํ, ทิสฺวา.\csromanspacer{} ปหูตํ สชฺฌํ\footnote{สชฺฌุํ (สี, สฺยา, อิ)} ฉฑฺฑิตํ, ทิสฺวา.\csromanspacer{} ปหูตํ สุวณฺณํ ฉฑฺฑิตํ, ทิสฺวา สหายโก สหายกํ อามนฺเตสิ “ยสฺส โข สมฺม อตฺถาย อิจฺเฉยฺยาม สาณํ วา สาณสุตฺตํ วา สาณิโย วา โขมํ วา โขมสุตฺตํ วา โขมทุสฺสํ วา กปฺปาสํ วา กปฺปาสิกสุตฺตํ วา กปฺปาสิกทุสฺสํ วา อยํ วา โลหํ วา ติปุํ วา สีสํ วา สชฺฌํ วา.\csromanspacer{} อิทํ ปหูตํ สุวณฺณํ ฉฑฺฑิตํ.\csromanspacer{} เตน หิ สมฺม ตฺวญฺจ สาณภารํ ฉฑฺเฑหิ, อหญฺจ สชฺฌภารํ\footnote{สชฺฌุภารํ (สี, สฺยา, อิ)} ฉฑฺเฑสฺสามิ, อุโภ สุวณฺณภารํ อาทาย คมิสฺสามา”ติ.\csromanspacer{} อยํ โข เม สมฺม สาณภาโร ทูราภโต จ สุสนฺนทฺโธ จ, อลํ เม ตฺวํ ปชานาหีติ.\csromanspacer{} อถ โข โส สหายโก สชฺฌภารํ ฉฑฺเฑตฺวา สุวณฺณภารํ อาทิยิ.}

\prose{เต เยน สโก คาโม เตนุปสงฺกมิํสุ, ตตฺถ โย โส สหายโก สาณภารํ อาทาย อคมาสิ, ตสฺส เนว มาตาปิตโร อภินนฺทิํสุ, น ปุตฺตทารา อภินนฺทิํสุ, น มิตฺตามจฺจา อภินนฺทิํสุ, น จ ตโตนิทานํ สุขํ โสมนสฺสํ อธิคจฺฉิ.\csromanspacer{} โย ปน โส สหายโก สุวณฺณภารํ อาทาย อคมาสิ, ตสฺส มาตาปิตโรปิ อภินนฺทิํสุ, ปุตฺตทาราปิ อภินนฺทิํสุ, มิตฺตามจฺจาปิ อภินนฺทิํสุ, ตโตนิทานญฺจ สุขํ โสมนสฺสํ อธิคจฺฉิ.\csromanspacer{} เอวเมว โข ตฺวํ ราชญฺญ สาณภาริกูปโม มญฺเญ ปฏิภาสิ.\csromanspacer{} ปฏินิสฺสชฺเชตํ ราชญฺญ ปาปกํ ทิฏฺฐิคตํ, ปฏินิสฺสชฺเชตํ ราชญฺญ ปาปกํ ทิฏฺฐิคตํ, มา เต อโหสิ ทีฆรตฺตํ อหิตาย ทุกฺขายาติ.}

\csromanmatikaanchor{mtk.153}
\csromantocmarkref{section}{สรณคมน}{mtk.153}
\bookmark[dest={mtk.153},level=1]{สรณคมน}
\csromanheader{1}{10. ปายาสิสุตฺต \textbf{สรณคมน}}
\proseitem{437}{\csromanfolio{279}ปุริเมเนว อหํ โอปมฺเมน โภโต กสฺสปสฺส อตฺตมโน อภิรทฺโธ.\csromanspacer{} อปิ จาหํ อิมานิ วิจิตฺรานิ ปญฺหาปฏิภานานิ โสตุกาโม เอวาหํ ภวนฺตํ กสฺสปํ ปจฺจนีกํ กาตพฺพํ อมญฺญิสฺสํ.\csromanspacer{} อภิกฺกนฺตํ โภ กสฺสป, อภิกฺกนฺตํ โภ กสฺสป, เสยฺยถาปิ โภ กสฺสป นิกฺกุชฺชิตํ วา อุกฺกุชฺเชยฺย, ปฏิจฺฉนฺนํ วา วิวเรยฺย, มูฬฺหสฺส วา มคฺคํ อาจิกฺเขยฺย, อนฺธกาเร วา เตลปชฺโชตํ ธาเรยฺย ‘จกฺขุมนฺโต รูปานิ ทกฺขนฺตี’ติ.\csromanspacer{} เอวเมวํ โภตา กสฺสเปน อเนกปริยาเยน ธมฺโม ปกาสิโต, เอสาหํ โภ กสฺสป ตํ ภวนฺตํ โคตมํ สรณํ คจฺฉามิ ธมฺมญฺจ ภิกฺขุสํฆญฺจ, อุปาสกํ มํ ภวํ กสฺสโป ธาเรตุ อชฺชตคฺเค ปาณุเปตํ สรณํ คตํ.}

\prose{อิจฺฉามิ จาหํ โภ กสฺสป มหายญฺญํ ยชิตุํ, อนุสาสตุ มํ ภวํ กสฺสโป ยํ มมสฺส ทีฆรตฺตํ หิตาย สุขายาติ.}

\csromanmatikaanchor{mtk.154}
\csromantocmarkref{section}{ยญฺญกถา}{mtk.154}
\bookmark[dest={mtk.154},level=1]{ยญฺญกถา}
\csromanheader{1}{\textbf{ยญฺญกถา}}
\proseitem{438}{ยถารูเป โข ราชญฺญ ยญฺเญ คาโว วา หญฺญนฺติ, อเชฬกา วา หญฺญนฺติ, กุกุฏสูกรา วา หญฺญนฺติ, วิวิธา วา ปาณา สํฆาตํ อาปชฺชนฺติ, ปฏิคฺคาหกา จ โหนฺติ มิจฺฉาทิฏฺฐี มิจฺฉาสงฺกปฺปา มิจฺฉาวาจา มิจฺฉากมฺมนฺตา มิจฺฉา-อาชีวา มิจฺฉาวายามา มิจฺฉาสตี มิจฺฉาสมาธี.\csromanspacer{} เอวรูโป โข ราชญฺญ ยญฺโญ น มหปฺผโล โหติ น มหานิสํโส น มหาชุติโก น มหาวิปฺผาโร.\csromanspacer{} เสยฺยถาปิ ราชญฺญ กสฺสโก พีชนงฺคลํ อาทาย วนํ ปวิเสยฺย, โส ตตฺถ ทุกฺเขตฺเต ทุพฺภูเม อวิหตขาณุกณฺฏเก พีชานิ ปติฏฺฐาเปยฺย ขณฺฑานิ ปูตีนิ วาตาตปหตานิ อสารทานิ อสุขสยิตานิ, เทโว จ น กาเลน กาลํ สมฺมาธารํ อนุปฺปเวจฺเฉยฺย.\csromanspacer{} อปิ นุ ตานิ พีชานิ วุทฺธิํ วิรูฬฺหิํ\footnote{วิรุฬฺหิํ (โมคฺคลฺลาเน)} เวปุลฺลํ อาปชฺเชยฺยุํ.\csromanspacer{} กสฺสโก วา วิปุลํ ผลํ อธิคจฺเฉยฺยาติ.\csromanspacer{} โน หิทํ\footnote{น เอวํ (สฺยา, ก)} โภ กสฺสป.\csromanspacer{} เอวเมว โข ราชญฺญ ยถารูเป ยญฺเญ คาโว วา หญฺญนฺติ, อเชฬกา วา หญฺญนฺติ, กุกฺกุฏสูกรา วา หญฺญนฺติ, วิวิธา วา ปาณา สํฆาตํ อาปชฺชนฺติ, ปฏิคฺคาหกา จ โหนฺติ มิจฺฉาทิฏฺฐี มิจฺฉาสงฺกปฺปา มิจฺฉาวาจา มิจฺฉากมฺมนฺตา มิจฺฉา-อาชีวา มิจฺฉาวายามา \csromanfolio{280}มิจฺฉาสตี มิจฺฉาสมาธี.\csromanspacer{} เอวรูโป โข ราชญฺญ ยญฺโญ น มหปฺผโล โหติ น มหานิสํโส น มหาชุติโก น มหาวิปฺผาโร.}

\prose{ยถารูเป จ โข ราชญฺญ ยญฺเญ เนว คาโว หญฺญนฺติ, น อเชฬกา หญฺญนฺติ, น กุกฺกุฏสูกรา หญฺญนฺติ, น วิวิธา วา ปาณา สํฆาตํ อาปชฺชนฺติ, ปฏิคฺคาหกา จ โหนฺติ สมฺมาทิฏฺฐี สมฺมาสงฺกปฺปา สมฺมาวาจา สมฺมากมฺมนฺตา สมฺมา-อาชีวา สมฺมาวายามา สมฺมาสตี สมฺมาสมาธี.\csromanspacer{} เอวรูโป โข ราชญฺญ ยญฺโญ มหปฺผโล โหติ มหานิสํโส มหาชุติโก มหาวิปฺผาโร.\csromanspacer{} เสยฺยถาปิ ราชญฺญ กสฺสโก พีชนงฺคลํ อาทาย วนํ ปวิเสยฺย, โส ตตฺถ สุเขตฺเต สุภูเม สุวิหตขาณุกณฺฏเก พีชานิ ปติฏฺฐเปยฺย อขณฺฑานิ อปูตีนิ อวาตาตปหตานิ สารทานิ สุขสยิตานิ.\csromanspacer{} เทโว จ กาเลน กาลํ สมฺมาธารํ อนุปฺปเวจฺเฉยฺย.\csromanspacer{} อปิ นุ ตานิ พีชานิ วุทฺธิํ วิรูฬฺหิํ เวปุลฺลํ อาปชฺเชยฺยุํ, กสฺสโก วา วิปุลํ ผลํ อธิคจฺเฉยฺยาติ.\csromanspacer{} เอวํ โภ กสฺสป.\csromanspacer{} เอวเมว โข ราชญฺญ ยถารูเป ยญฺเญ เนว คาโว หญฺญนฺติ, น อเชฬกา หญฺญนฺติ, น กุกฺกุฏสูกรา หญฺญนฺติ, น วิวิธา วา ปาณา สํฆาตํ อาปชฺชนฺติ, ปฏิคฺคาหกา จ โหนฺติ สมฺมาทิฏฺฐี สเมสงฺกปฺปา สเมวาจา สมฺมากมฺมนฺตา สมฺมา-อาชีวา สมฺมาวายามา สมฺมาสตี สเมสมาธี.\csromanspacer{} เอวรูโป โข ราชญฺญ ยญฺโญ มหปฺผโล โหติ มหานิสํโส มหาชุติโก มหาวิปฺผาโรติ.}

\csromanmatikaanchor{mtk.155}
\csromantocmarkref{section}{อุตฺตรมาณววตฺถุ}{mtk.155}
\bookmark[dest={mtk.155},level=1]{อุตฺตรมาณววตฺถุ}
\csromanheader{1}{\textbf{อุตฺตรมาณววตฺถุ}}
\proseitem{439}{อถ โข ปายาสิ ราชญฺโญ ทานํ ปฏฺฐเปสิ สมณพฺราหฺมณกปณทฺธิกวณิพฺพกยาจกานํ.\csromanspacer{} ตสฺมิํ โข ปน ทาเน เอวรูปํ โภชนํ ทียติ กณาชกํ พิลงฺคทุติยํ, โธรกานิ\footnote{เถรกานิ (สี, อิ), โจรกานิ (สฺยา)} จ วตฺถานิ คุฬวาลกานิ\footnote{คุฬคาฬกานิ (ก)}.\csromanspacer{} ตสฺมิํ โข ปน ทาเน อุตฺตโร นาม มาณโว วาวโฏ\footnote{พฺยาวโฏ (สี, อิ)} อโหสิ, โส ทานํ ทตฺวา เอวํ อนุทฺทิสติ “อิมินาหํ ทาเนน ปายาสิํ ราชญฺญเมว อิมสฺมิํ โลเก สมาคจฺฉิํ, มา ปรสฺมินฺ”ติ.\csromanspacer{} อสฺโสสิ โข ปายาสิ ราชญฺโญ “อุตฺตโร กิร มาณโว ทานํ ทตฺวา เอวํ อนุทฺทิสติ ‘อิมินาหํ ทาเนน ปายาสิํ ราชญฺญเมว อิมสฺมิํ โลเก สมาคจฺฉิํ, มา ปรสฺมินฺ’ติ”.}

\vspace{\tipitakaparskip}
\csromancenter{ปายาสิสุตฺต อถ โข ปายาสิ ราชญฺโญ อุตฺตรํ มาณวํ อามนฺตาเปตฺวา เอตทโวจ “สจฺจํ กิร ตฺวํ ตาต อุตฺตร ทานํ ทตฺวา เอวํ อนุทฺทิสสิ ‘อิมินาหํ ทาเนน ปายาสิํ ราชญฺญเมว อิมสฺมิํ โลเก สมาคจฺฉิํ, มา ปรสฺมินฺ’ติ”.\csromanspacer{} เอวํ โภ.\csromanspacer{} กิสฺส ปน ตฺวํ ตาต อุตฺตร ทานํ ทตฺวา เอวํ อนุทฺทิสสิ “อิมินาหํ ทาเนน ปายาสิํ ราชญฺญเมว อิมสฺมิํ โลเก สมาคจฺฉิํ, มา ปรสฺมินฺ”ติ.\csromanspacer{} นนุ มยํ ตาต อุตฺตร ปุญฺญตฺถิกา ทานสฺเสว ผลํ ปาฏิกงฺขิโนติ.\csromanspacer{} โภโต โข ทาเน เอวรูปํ โภชนํ ทียติ กณาชกํ พิลงฺคทุติยํ, ยํ ภวํ ปาทาปิ\footnote{ปาทาสิ (ก)} น อิจฺเฉยฺย สมฺผุสิตุํ\footnote{ฉุปิตุํ (อิ, ก)}, กุโต ภุญฺชิตุํ.\csromanspacer{} โธรกานิ จ วตฺถานิ คุฬวาลกานิ, ยานิ ภวํ ปาทาปิ1 น อิจฺเฉยฺย สมฺผุสิตุํ, กุโต ปริทหิตุํ.\csromanspacer{} ภวํ โข ปนมฺหากํ ปิโย มนาโป, กถํ มยํ มนาปํ อมนาเปน สํโยเชมาติ.\csromanspacer{} เตน หิ ตฺวํ ตาต อุตฺตร ยาทิสาหํ โภชนํ ภุญฺชามิ, ตาทิสํ โภชนํ ปฏฺฐเปหิ.\csromanspacer{} ยาทิสานิ จาหํ วตฺถานิ ปริทหามิ, ตาทิสานิ จ วตฺถานิ ปฏฺฐเปหีติ. “เอวํ โภ”ติ โข อุตฺตโร มาณโว ปายาสิสฺส ราชญฺญสฺส ปฏิสฺสุตฺวา ยาทิสํ โภชนํ ปายาสิ ราชญฺโญ ภุญฺชติ, ตาทิสํ โภชนํ ปฏฺฐเปสิ.\csromanspacer{} ยาทิสานิ จ วตฺถานิ ปายาสิ ราชญฺโญ ปริทหติ, ตาทิสานิ จ วตฺถานิ ปฏฺฐเปสิ.}
\proseitem{440}{\csromanfolio{281}อถ โข ปายาสิ ราชญฺโญ อสกฺกจฺจํ ทานํ ทตฺวา อสหตฺถา ทานํ ทตฺวา อจิตฺตีกตํ\footnote{อจิตฺติกตํ (ก)} ทานํ ทตฺวา อปวิทฺธํ ทานํ ทตฺวา กายสฺส เภทา ปรํ มรณา จาตุมหาราชิกานํ เทวานํ สหพฺยตํ อุปปชฺชิ สุญฺญํ เสรีสกํ วิมานํ.\csromanspacer{} โย ปน ตสฺส ทาเน วาวโฏ อโหสิ อุตฺตโร นาม มาณโว, โส สกฺกจฺจํ ทานํ ทตฺวา สหตฺถา ทานํ ทตฺวา จิตฺตีกตํ ทานํ ทตฺวา อนปวิทฺธํ ทานํ ทตฺวา กายสฺส เภทา ปรํ มรณา สุคติํ สคฺคํ โลกํ อุปปชฺชิ เทวานํ ตาวติํสานํ สหพฺยตํ.}

\csromanmatikaanchor{mtk.156}
\csromantocmarkref{section}{ปายาสิเทวปุตฺต}{mtk.156}
\bookmark[dest={mtk.156},level=1]{ปายาสิเทวปุตฺต}
\csromanheader{1}{\textbf{ปายาสิเทวปุตฺต}}
\proseitem{441}{เตน โข ปน สมเยน อายสฺมา ควมฺปติ อภิกฺขณํ สุญฺญํ เสรีสกํ วิมานํ ทิวาวิหารํ คจฺฉติ.\csromanspacer{} อถ โข ปายาสิ เทวปุตฺโต เยนายสฺมา ควมฺปติ เตนุปสงฺกมิ, อุปสงฺกมิตฺวา อายสฺมนฺตํ ควมฺปติํ \csromanfolio{282}อภิวาเทตฺวา เอกมนฺตํ อฏฺฐาสิ.\csromanspacer{} เอกมนฺตํ ฐิตํ โข ปายาสิํ เทวปุตฺตํ อายสฺมา ควมฺปติ เอตทโวจ “โกสิ ตฺวํ อาวุโส”ติ.\csromanspacer{} อหํ ภนฺเต ปายาสิ ราชญฺโญติ.\csromanspacer{} นนุ ตฺวํ อาวุโส เอวํทิฏฺฐิโก อโหสิ “อิติปิ นตฺถิ ปโร โลโก, นตฺถิ สตฺตา โอปปาติกา, นตฺถิ สุกตทุกฺกฏานํ กมฺมานํ ผลํ วิปาโก”ติ.\csromanspacer{} สจฺจาหํ ภนฺเต เอวํทิฏฺฐิโก อโหสิํ “อิติปิ นตฺถิ ปโร โลโก, นตฺถิ สตฺตา โอปปาติกา, นตฺถิ สุกตทุกฺกฏานํ กมฺมานํ ผลํ วิปาโก”ติ.\csromanspacer{} อปิ จาหํ อยฺเยน กุมารกสฺสเปน เอตสฺมา ปาปกา ทิฏฺฐิคตา วิเวจิโตติ.\csromanspacer{} โย ปน เต อาวุโส ทาเน วาวโฏ อโหสิ อุตฺตโร นาม มาณโว, โส กุหิํ อุปปนฺโนติ.\csromanspacer{} โย เม ภนฺเต ทาเน วาวโฏ อโหสิ อุตฺตโร นาม มาณโว, โส สกฺกจฺจํ ทานํ ทตฺวา สหตฺถา ทานํ ทตฺวา จิตฺตีกตํ ทานํ ทตฺวา อนปวิทฺธํ ทานํ ทตฺวา กายสฺส เภทา ปรํ มรณา สุคติํ สคฺคํ โลกํ อุปปนฺโน เทวานํ ตาวติํสานํ สหพฺยตํ.\csromanspacer{} อหํ ปน ภนฺเต อสกฺกจฺจํ ทานํ ทตฺวา อสหตฺถา ทานํ ทตฺวา อจิตฺตีกตํ ทานํ ทตฺวา อปวิทฺธํ ทานํ ทตฺวา กายสฺส เภทา ปรํ มรณา จาตุมหาราชิกานํ เทวานํ สหพฺยตํ อุปปนฺโน สุญฺญํ เสรีสกํ วิมานํ เอน หิ ภนฺเต ควมฺปติ มนุสฺสโลกํ คนฺตฺวา เอวมาโรเจหิ “สกฺกจฺจํ ทานํ เทถ, สหตฺถา ทานํ เทถ, จิตฺตีกตํ ทานํ เทถ, อนปวิทฺธํ ทานํ เทถ, ปายาสิ ราชญฺโญ อสกฺกจฺจํ ทานํ ทตฺวา อสหตฺถา ทานํ ทตฺวา อจิตฺตีกตํ ทานํ ทตฺวา อปวิทฺธํ ทานํ ทตฺวา กายสฺส เภทา ปรํ มรณา จาตุมหาราชิกานํ เทวานํ สหพฺยตํ อุปปนฺโน สุญฺญํ เสรีสกํ วิมานํ.\csromanspacer{} โย ปน ตสฺส ทาเน วาวโฏ อโหสิ อุตฺตโร นาม มาณโว, โส สกฺกจฺจํ ทานํ ทตฺวา สหตฺถา ทานํ ทตฺวา จิตฺตีกตํ ทานํ ทตฺวา อนปวิทฺธํ ทานํ ทตฺวา กายสฺส เภทา ปรํ มรณา สุคติํ สคฺคํ โลกํ อุปปนฺโน เทวานํ ตาวติํสานํ สหพฺยตนฺ”ติ.}

\prose{อถ โข อายสฺมา ควมฺปติ มนุสฺสโลกํ อาคนฺตฺวา เอวมาโรเจสิ “สกฺกจฺจํ ทานํ เทถ, สหตฺถา ทานํ เทถ, จิตฺตีกตํ ทานํ เทถ, อนปวิทฺธํ ทานํ เทถ.\csromanspacer{} ปายาสิ ราชญฺโญ อสกฺกจฺจํ ทานํ ทตฺวา อสหตฺถา ทานํ ทตฺวา อจิตฺตีกตํ ทานํ ทตฺวา อปวิทฺธํ ทานํ ทตฺวา กายสฺส เภทา ปรํ มรณา จาตุมหาราชิกานํ เทวานํ สหพฺยตํ อุปปนฺโน สุญฺญํ เสรีสกํ วิมานํ.\csromanspacer{} โย ปน ตสฺส ทาเน วาวโฏ อโหสิ}

\vspace{\tipitakaparskip}
\csromancenter{ปายาสิสุตฺต อุตฺตโร นาม มาณโว, โส สกฺกจฺจํ ทานํ ทตฺวา สหตฺถา ทานํ ทตฺวา จิตฺตีกตํ ทานํ ทตฺวา อนปวิทฺธํ ทานํ ทตฺวา กายสฺส เภทา ปรํ มรณา สุคติํ สคฺคํ โลกํ อุปปนฺโน เทวานํ ตาวติํสานํ สหพฺยตนฺ”ติ.}
\nitthitam{ปายาสิสุตฺตํ นิฏฺฐิตํ ทสมํ.}
\nitthitam{มหาวคฺโค นิฏฺฐิโต.}
\titlehead{\textbf{ตสฺสุทฺทานํ}}
\csromanmatikaanchor{mtk.157}
\csromantocmarkref{section}{อุทฺทานคาถา}{mtk.157}
\bookmark[dest={mtk.157},level=1]{อุทฺทานคาถา}
\csromangathameasure{มหาปทาน นิทานํ, นิพฺพานญฺจ สุทสฺสนํ. \\ ชนวสภ โควินฺทํ, สมยํ สกฺกปญฺหกํ. \\ มหาสติปฏฺฐานญฺจ, ปายาสิ ทสมํ ภเว.}
\csromangathagroup{\csromangathastanza{\csromanfolio{283}มหาปทาน นิทานํ, นิพฺพานญฺจ สุทสฺสนํ. \\ ชนวสภ โควินฺทํ, สมยํ สกฺกปญฺหกํ. \\ มหาสติปฏฺฐานญฺจ, ปายาสิ ทสมํ ภเว\footnote{สติปฏฺฐานปายาสิ, มหาวคฺคสฺส สงฺคโห (สี, อิ) สติปฏฺฐานปายาสิ, มหาวคฺโคติ วุจฺจตีติ (สฺยา)}.}}

\nitthitam{\textbf{มหาวคฺคปาฬิ นิฏฺฐิตา.}}
\orphannote{ยาวตา วฏฺฏํ วฏฺฏติ (ก-สี)}

\orphannote{เต (ก)}

